%=====================================================================
% ĐỀ THI THỬ TỐT NGHIỆP THPT 2026 -- MÔN VẬT LÍ
% Nguồn: Sở GD&ĐT Ninh Bình, khảo sát lần 3 (mã đề gốc 0201). Gõ lại theo template Eduzy.
%=====================================================================
\def\CheDo{solcolor}% loigiai | dethi
\def\tenchuong{Sở GD\&ĐT Ninh Bình -- Lần 3}
\def\thuonghieu{VẬT LÍ 12 -- Eduzy.com.vn}
\def\TachCauDS{1}
\def\TachCauLG{1}
%============================================================
%  PREAMBLE EDUZY v2 - dùng chung cho đề thi, bài học, chuyên đề (gói ex_test)
%  Khác bản preamble-eduzy.tex ở đầu trang, chân trang: dải nền nâu đồng tràn ngang,
%  logo + Eduzy.com.vn, ô chữ vàng kem cạnh xiên, chân trang "Trang (số)";
%  logo Eduzy chìm giữa mọi trang (tắt: \logochimtat, bật lại: \logochimbat).
%  Cách dùng ở file chính (xem mau-chuan.tex, mau-chuyen-de.tex):
%      \def\CheDo{loigiai}          % loigiai | dethi | solcolor | color | book
%      \def\tenchuong{<dòng chữ trên đầu trang>}
%      \def\thuonghieu{<chữ ở chân trang lẻ>}   % không khai báo thì dùng mặc định
%      \def\tacgia{<chữ ở chân trang chẵn>}      % không khai báo thì dùng mặc định
%      \def\khauhieu{<chữ cạnh logo ở đầu trang>}  % tùy chọn, mặc định Eduzy.com.vn
%      %============================================================
%  PREAMBLE EDUZY v2 - dùng chung cho đề thi, bài học, chuyên đề (gói ex_test)
%  Khác bản preamble-eduzy.tex ở đầu trang, chân trang: dải nền nâu đồng tràn ngang,
%  logo + Eduzy.com.vn, ô chữ vàng kem cạnh xiên, chân trang "Trang (số)";
%  logo Eduzy chìm giữa mọi trang (tắt: \logochimtat, bật lại: \logochimbat).
%  Cách dùng ở file chính (xem mau-chuan.tex, mau-chuyen-de.tex):
%      \def\CheDo{loigiai}          % loigiai | dethi | solcolor | color | book
%      \def\tenchuong{<dòng chữ trên đầu trang>}
%      \def\thuonghieu{<chữ ở chân trang lẻ>}   % không khai báo thì dùng mặc định
%      \def\tacgia{<chữ ở chân trang chẵn>}      % không khai báo thì dùng mặc định
%      \def\khauhieu{<chữ cạnh logo ở đầu trang>}  % tùy chọn, mặc định Eduzy.com.vn
%      %============================================================
%  PREAMBLE EDUZY v2 - dùng chung cho đề thi, bài học, chuyên đề (gói ex_test)
%  Khác bản preamble-eduzy.tex ở đầu trang, chân trang: dải nền nâu đồng tràn ngang,
%  logo + Eduzy.com.vn, ô chữ vàng kem cạnh xiên, chân trang "Trang (số)";
%  logo Eduzy chìm giữa mọi trang (tắt: \logochimtat, bật lại: \logochimbat).
%  Cách dùng ở file chính (xem mau-chuan.tex, mau-chuyen-de.tex):
%      \def\CheDo{loigiai}          % loigiai | dethi | solcolor | color | book
%      \def\tenchuong{<dòng chữ trên đầu trang>}
%      \def\thuonghieu{<chữ ở chân trang lẻ>}   % không khai báo thì dùng mặc định
%      \def\tacgia{<chữ ở chân trang chẵn>}      % không khai báo thì dùng mặc định
%      \def\khauhieu{<chữ cạnh logo ở đầu trang>}  % tùy chọn, mặc định Eduzy.com.vn
%      \input{preamble-eduzy-v2}
%      \TLdapso                      % tài liệu: \TLfull | \TLdapso | \TLan
%      \begin{document} ... \end{document}
%  Logo Eduzy vẽ bằng TikZ ngay trong file (\eduzylogo), không cần file ảnh.
%  Biên dịch: pdfLaTeX (máy có gói vietnam) hoặc XeLaTeX.
%============================================================
\documentclass[12pt,a4paper,twoside]{article}%<-- twoside: bat dau trang/chan trang chan-le
\usepackage{iftex}
\usepackage{amsmath,amssymb,tikz,mathrsfs,tkz-tab}
\usepackage{graphicx}
\usepackage{pgfplots}% đồ thị hàm số
\pgfplotsset{compat=1.17}
%\usepackage{mathptmx}
\usepackage{mathpazo}
\ifXeTeX% XeLaTeX: font chữ có đủ dấu tiếng Việt, giữ nguyên font toán mathpazo
	\usepackage[no-math]{fontspec}
	\setmainfont{TeX Gyre Pagella}
	\newfontfamily\EZqag{TeX Gyre Adventor}[NFSSFamily=qag,Ligatures=TeX]
\fi
\usepackage[many]{tcolorbox}
\usetikzlibrary{math,through,calc,intersections,angles,quotes,shapes,shapes.geometric,arrows,patterns,matrix,chains,arrows.meta,decorations.shapes,decorations.fractals,decorations.markings,shadows,decorations.pathreplacing}
\usepackage[top=1.6cm, bottom=1.8cm, left=1.8cm, right=1.8cm]{geometry}% v2: top 1.6cm (bản cũ 1.5cm) để chữ không sát dải đầu trang
\providecommand{\CheDo}{loigiai}% loigiai | dethi | solcolor | color | book
\expandafter\usepackage\expandafter[\CheDo]{ex_test}
% \tenchuong: khai báo ở file chính (trước \input) hoặc \renewcommand sau \input
%============================================================
%  [GỘP TỪ eduzy-de.tex] EDUZY - MÀU NÂU ĐỒNG / VÀNG KIM
%  + ĐẦU TRANG, CHÂN TRANG KIỂU DẢI NGANG (v2), ĐỐI XỨNG TRANG CHẴN / TRANG LẺ.
%  (không cần file eduzy-de.tex riêng nữa)
%  Khối này phải đặt SAU \usepackage[...]{ex_test} và SAU \def\tenchuong{...}
%============================================================
\usepackage{lastpage}
\usepackage{needspace}
%--- DÀN TRANG ĐỀ THI: tách câu (THÊM 05/11/2026; chỉ thêm lệnh mới, KHÔNG đổi hành vi mặc định) ---
%  Hai công tắc, mặc định tắt (= preamble chạy y như cũ). Bật bằng cách đặt TRƯỚC \input{preamble-eduzy-v2}
%  trong file đề (tools/dan_trang/dan_trang.py tự ghi khi đặt \NeedspaceCau):
%    \def\TachCauDS{1}   câu Đúng/Sai được tách: câu dẫn (kèm hình) ở cuối trang, bảng Phát biểu sang trang sau
%                        (câu dẫn và hình không tách nhau, bảng không tách giữa hai trang)
%    \def\TachCauLG{1}   bản lời giải (\CheDo khác dethi): phần đề của câu cao hơn 1/3 trang được phép tách
%  Lệnh dùng trong file đề (giá trị do tools/dan_trang/dan_trang.py đo ở bản dethi):
%    \NeedspaceCau[<cao câu dẫn + hình>]{<cao cả khối câu>}   ngay trên \begin{ex} (câu đầu phần: trên tiêu đề phần)
%        công tắc tắt: đúng bằng \Needspace{<cao cả khối câu>} như cũ;
%        \TachCauDS=1 và có [..] (câu Đúng/Sai được chọn tách): chỉ giữ câu dẫn + hình;
%        \TachCauLG=1 ở bản lời giải: đòi tối đa 1/3 \textheight.
%    \NeedspaceBang{<cao bảng Phát biểu>}%   ngay trước \choiceTFt: \TachCauDS=1 thì giữ trọn bảng trên một trang;
%        công tắc tắt: không làm gì.
\providecommand{\TachCauDS}{0}
\providecommand{\TachCauLG}{0}
\newlength{\EZnsCau}
\newcommand{\NeedspaceCau}[2][]{%
	\setlength{\EZnsCau}{#2}%
	\ifnum\TachCauDS=1 \if\relax\detokenize{#1}\relax\else\setlength{\EZnsCau}{#1}\fi\fi
	\ifnum\TachCauLG=1 \ifdefstring{\kindOPTN}{dethi}{}{\ifdim\EZnsCau>0.3333\textheight\setlength{\EZnsCau}{0.3333\textheight}\fi}\fi
	\Needspace{\EZnsCau}}
\newcommand{\NeedspaceBang}[1]{\ifnum\TachCauDS=1 \Needspace{#1}\fi}

%--- Thang màu: đậm -> nhạt
\definecolor{gold1}{RGB}{126,96,0}    % nâu đồng đậm  #7E6000
\definecolor{gold2}{RGB}{192,144,3}   % vàng kim      #C09003
\definecolor{gold3}{RGB}{255,215,99}  % vàng nhạt     #FFD763
\definecolor{gold4}{RGB}{255,231,153} % vàng rất nhạt #FFE799
\definecolor{gold5}{RGB}{255,243,205} % vàng kem      #FFF3CD
%--- Ghi đè tông xanh/đỏ cũ
\definecolor{myblue}{RGB}{126,96,0}
\colorlet{blue}{gold1}
\colorlet{cyan}{gold2}
\colorlet{teal}{gold2}

%--- Tên chương mặc định (file đề nên tự khai báo lại)
\providecommand{\tenchuong}{}
%--- Chân trang
\providecommand{\thuonghieu}{TOÁN 12 -- Eduzy.com.vn}
\providecommand{\tacgia}{Tô Gia Bảo -- 0353703783}

%--- Chữ cạnh logo ở đầu trang
\providecommand{\khauhieu}{Eduzy.com.vn}

%====================== THÔNG SỐ ĐẦU TRANG, CHÂN TRANG (v2) ======================
%  Dải nền đậm tràn ngang trang; ô chữ vàng kem có cạnh xiên, kèm 2 vạch xiên chuyển màu.
%  Đầu trang lẻ : logo + \khauhieu (trái)   | ô chữ \tenchuong (phải)
%  Chân trang lẻ: ô chữ \thuonghieu (trái)  | Trang (số) (phải)
%  Trang chẵn soi gương trang lẻ; chân trang chẵn ghi \tacgia.
%--- Màu: chỉ dùng thang vàng kim Eduzy
\colorlet{ezNen}{gold1}% dải nền đậm
\colorlet{ezVach1}{gold3}% vạch xiên sát dải nền (đậm hơn)
\colorlet{ezVach2}{gold4}% vạch xiên sát ô chữ (nhạt hơn)
\colorlet{ezO}{gold5}% ô chữ vàng kem
%--- Kích thước
\newlength{\EZdai}\setlength{\EZdai}{1.1cm}% chiều cao dải nền
\newlength{\EZo}\setlength{\EZo}{0.68cm}% chiều cao ô chữ
\newlength{\EZxien}\setlength{\EZxien}{0.55cm}% độ lệch ngang của cạnh xiên
\newlength{\EZvach}\setlength{\EZvach}{0.75cm}% bề rộng mỗi vạch xiên
\newlength{\EZkhe}\setlength{\EZkhe}{0.12cm}% khe trắng giữa vạch và ô chữ
\newlength{\EZdem}\setlength{\EZdem}{0.4cm}% khoảng đệm từ chữ tới cạnh xiên
\newlength{\EZle}\setlength{\EZle}{1.8cm}% lề trái/phải của chữ (bằng lề thân bài)
\newlength{\EZrongmin}\setlength{\EZrongmin}{4.5cm}% ô chữ rộng tối thiểu
\newlength{\EZrong}% bề rộng chữ trong ô (tính tự động)
%--- Font
\newcommand{\EZfontO}{\fontsize{10.5pt}{12pt}\selectfont\bfseries}% chữ trong ô
\newcommand{\EZfontKH}{\fontfamily{qag}\fontsize{11pt}{12pt}\selectfont\bfseries}% \khauhieu cạnh logo
\newcommand{\EZfontTrang}{\fontsize{11.5pt}{12pt}\selectfont\bfseries}% chữ "Trang" và số trang
%--- Đo bề rộng chữ trong ô
\newcommand{\EZdorong}[1]{%
	\settowidth{\EZrong}{\EZfontO #1}%
	\ifdim\EZrong<\EZrongmin\setlength{\EZrong}{\EZrongmin}\fi}
%--- Logo + chữ trắng \khauhieu (dựng sẵn một lần khi vào \begin{document})
\newsavebox{\EZlogobox}
\AtBeginDocument{\sbox{\EZlogobox}{%
		\raisebox{-0.5\height}{\eduzylogo[1cm]}\hspace{0.15cm}%
		\raisebox{-0.5\height}{\color{white}\EZfontKH\khauhieu}}}
%--- Hai vạch xiên + khe trắng. #1: điểm gốc ở mép dưới, #2: hướng (1 sang phải, -1 sang trái),
%    #3: hướng nghiêng (1 kiểu "/", -1 kiểu "\")
\newcommand{\EZvachxien}[3]{%
	\fill[white] (#1) --++({(#3)*\EZxien},\EZo) --++({(#2)*\EZkhe},0) --++({-(#3)*\EZxien},-\EZo) --cycle;
	\fill[ezVach2] ([xshift=(#2)*\EZkhe]#1) --++({(#3)*\EZxien},\EZo) --++({(#2)*\EZvach},0) --++({-(#3)*\EZxien},-\EZo) --cycle;
	\fill[ezVach1] ([xshift=(#2)*(\EZkhe+\EZvach)]#1) --++({(#3)*\EZxien},\EZo) --++({(#2)*\EZvach},0) --++({-(#3)*\EZxien},-\EZo) --cycle;
	\draw[white,line width=0.6pt] ([xshift=(#2)*(\EZkhe+\EZvach)]#1) --++({(#3)*\EZxien},\EZo);
}
%--- Chữ trong ô (căn giữa theo chiều cao ô, không lệch vì dấu tiếng Việt)
\tikzset{ezchu/.style={inner sep=0pt,text=ezNen,text height=0.72em,text depth=0.22em}}

%====================== LOGO CHÌM GIỮA TRANG ======================
%  Vẽ cùng đầu trang nên nằm dưới chữ của thân bài, có trên mọi trang.
%  Tắt: \logochimtat (từ trang đó trở đi), bật lại: \logochimbat.
%  Đổi cỡ, độ đậm: khai báo trước \input, ví dụ \def\EZchimcao{9cm} \def\EZchimdo{0.1}
\providecommand{\EZchimcao}{10cm}% chiều cao logo chìm
\providecommand{\EZchimdo}{0.12}% độ đậm (0 = trong suốt, 1 = rõ hẳn)
\newif\ifEZchim \EZchimtrue
\newcommand{\logochimtat}{\EZchimfalse}
\newcommand{\logochimbat}{\EZchimtrue}
%  Dựng sẵn logo mờ một lần khi vào \begin{document}; transparency group làm mờ cả logo
%  một lượt nên không lộ vệt ở chỗ các dải màu chồng nhau.
\newsavebox{\EZchimtho}
\newsavebox{\EZchimbox}
\AtBeginDocument{%
	\sbox{\EZchimtho}{\eduzylogo[\EZchimcao]}%
	\sbox{\EZchimbox}{\begin{tikzpicture}
			\begin{scope}[transparency group,opacity=\EZchimdo]
				\node[inner sep=0pt] {\usebox{\EZchimtho}};
			\end{scope}
	\end{tikzpicture}}}
\newcommand{\EZvechim}{%
	\ifEZchim
		\node[inner sep=0pt] at (current page.center) {\usebox{\EZchimbox}};
	\fi}

%====================== ĐẦU TRANG - TRANG LẺ ======================
%  (logo bên trái, ô chữ \tenchuong bên phải)
\newcommand{\myfancyheadOdd}{%
	\EZdorong{\tenchuong}%
	\begin{tikzpicture}[remember picture,overlay]
		\coordinate (T) at (current page.north west);
		\coordinate (P) at (current page.north east);
		%--- logo chìm giữa trang
		\EZvechim
		%--- dải nền
		\fill[ezNen] (T) rectangle ([yshift=-\EZdai]P);
		%--- ô chữ: X là góc dưới bên trái, cạnh trái xiên kiểu "/"
		\coordinate (X) at ([xshift=-\EZle-\EZrong-\EZdem-\EZxien,yshift=-\EZdai]P);
		\fill[ezO] (X) --++(\EZxien,\EZo) -- ([yshift=-\EZdai+\EZo]P) -- ([yshift=-\EZdai]P) --cycle;
		\EZvachxien{X}{-1}{1}
		\node[ezchu,anchor=east] at ([xshift=-\EZle,yshift=-\EZdai+0.5*\EZo]P) {\EZfontO\tenchuong};
		%--- logo + tên
		\node[inner sep=0pt,anchor=west] at ([xshift=\EZle,yshift=-0.5*\EZdai]T) {\usebox{\EZlogobox}};
	\end{tikzpicture}%
}

%====================== ĐẦU TRANG - TRANG CHẴN =====================
%  (soi gương: ô chữ \tenchuong bên trái, logo bên phải)
\newcommand{\myfancyheadEven}{%
	\EZdorong{\tenchuong}%
	\begin{tikzpicture}[remember picture,overlay]
		\coordinate (T) at (current page.north west);
		\coordinate (P) at (current page.north east);
		%--- logo chìm giữa trang
		\EZvechim
		%--- dải nền
		\fill[ezNen] (T) rectangle ([yshift=-\EZdai]P);
		%--- ô chữ: X là góc dưới bên phải, cạnh phải xiên kiểu "\"
		\coordinate (X) at ([xshift=\EZle+\EZrong+\EZdem+\EZxien,yshift=-\EZdai]T);
		\fill[ezO] (X) --++(-\EZxien,\EZo) -- ([yshift=-\EZdai+\EZo]T) -- ([yshift=-\EZdai]T) --cycle;
		\EZvachxien{X}{1}{-1}
		\node[ezchu,anchor=west] at ([xshift=\EZle,yshift=-\EZdai+0.5*\EZo]T) {\EZfontO\tenchuong};
		%--- logo + tên
		\node[inner sep=0pt,anchor=east] at ([xshift=-\EZle,yshift=-0.5*\EZdai]P) {\usebox{\EZlogobox}};
	\end{tikzpicture}%
}

%====================== CHÂN TRANG - TRANG LẺ ======================
%  (ô chữ \thuonghieu bên trái, "Trang (số)" bên phải)
\newcommand{\myfancyfootOdd}{%
	\EZdorong{\thuonghieu}%
	\begin{tikzpicture}[remember picture,overlay]
		\coordinate (T) at (current page.south west);
		\coordinate (P) at (current page.south east);
		%--- dải nền
		\fill[ezNen] (T) rectangle ([yshift=\EZdai]P);
		%--- ô chữ sát mép trên dải: X là góc dưới bên phải, cạnh phải xiên kiểu "/"
		\coordinate (X) at ([xshift=\EZle+\EZrong+\EZdem,yshift=\EZdai-\EZo]T);
		\fill[ezO] ([yshift=\EZdai-\EZo]T) -- (X) --++(\EZxien,\EZo) -- ([yshift=\EZdai]T) --cycle;
		\EZvachxien{X}{1}{1}
		\node[ezchu,anchor=west] at ([xshift=\EZle,yshift=\EZdai-0.5*\EZo]T) {\EZfontO\thuonghieu};
		%--- số trang
		\node[circle,draw=white,line width=0.9pt,text=white,inner sep=0pt,minimum size=0.62cm,anchor=east]
		(so) at ([xshift=-\EZle,yshift=0.5*\EZdai]P) {\EZfontTrang\thepage};
		\node[inner sep=0pt,text=white,anchor=base east] (tr) at ([xshift=-0.15cm]so.base west) {\EZfontTrang Trang};
	\end{tikzpicture}%
}

%====================== CHÂN TRANG - TRANG CHẴN ====================
%  (soi gương: "Trang (số)" bên trái, ô chữ \tacgia bên phải)
\newcommand{\myfancyfootEven}{%
	\EZdorong{\tacgia}%
	\begin{tikzpicture}[remember picture,overlay]
		\coordinate (T) at (current page.south west);
		\coordinate (P) at (current page.south east);
		%--- dải nền
		\fill[ezNen] (T) rectangle ([yshift=\EZdai]P);
		%--- ô chữ sát mép trên dải: X là góc dưới bên trái, cạnh trái xiên kiểu "\"
		\coordinate (X) at ([xshift=-\EZle-\EZrong-\EZdem,yshift=\EZdai-\EZo]P);
		\fill[ezO] ([yshift=\EZdai-\EZo]P) -- (X) --++(-\EZxien,\EZo) -- ([yshift=\EZdai]P) --cycle;
		\EZvachxien{X}{-1}{-1}
		\node[ezchu,anchor=east] at ([xshift=-\EZle,yshift=\EZdai-0.5*\EZo]P) {\EZfontO\tacgia};
		%--- số trang
		%    (vòng tròn ẩn so0 nằm giữa dải để lấy đường chân chữ, giống trang lẻ)
		\node[circle,inner sep=0pt,minimum size=0.62cm] (so0) at ([yshift=0.5*\EZdai]T) {\EZfontTrang\phantom{\thepage}};
		\node[inner sep=0pt,text=white,anchor=base west] (tr) at ($(so0.base -| T)+(\EZle,0)$) {\EZfontTrang Trang};
		\node[circle,draw=white,line width=0.9pt,text=white,inner sep=0pt,minimum size=0.62cm,anchor=base west]
		(so) at ([xshift=0.15cm]tr.base east) {\EZfontTrang\thepage};
	\end{tikzpicture}%
}

%--- Giữ tương thích với các file cũ đang gọi \myfancyhead, \myfancyfoot
\newcommand{\myfancyhead}{\myfancyheadOdd}
\newcommand{\myfancyfoot}{\myfancyfootOdd}

%======================= GẮN VÀO TRANG =============================
\pagestyle{fancy}
\fancyhf{}
\fancyhead[LO]{\myfancyheadOdd}
\fancyhead[RE]{\myfancyheadEven}
\fancyfoot[LO]{\myfancyfootOdd}
\fancyfoot[RE]{\myfancyfootEven}
\renewcommand{\headrulewidth}{0pt}
\renewcommand{\footrulewidth}{0pt}
\fancypagestyle{plain}{%
	\fancyhf{}
	\fancyhead[LO]{\myfancyheadOdd}
	\fancyhead[RE]{\myfancyheadEven}
	\fancyfoot[LO]{\myfancyfootOdd}
	\fancyfoot[RE]{\myfancyfootEven}
	\renewcommand{\headrulewidth}{0pt}
	\renewcommand{\footrulewidth}{0pt}
}
%=================== HẾT PHẦN GỘP TỪ eduzy-de.tex ===================
\newcommand{\indam}[1]{
	{\fontsize{12pt}{9pt}\sffamily \selectfont\bfseries\color{gold1} #1}
}
\newcommand{\ind}[1]{
	{\fontsize{14pt}{9pt}\fontfamily{qag} \selectfont\bfseries\color{gold1} #1}
}
\renewcommand{\FalseEX}{\stepcounter{dapan}\textbf{\selectfont\color{gold1}\Alph{dapan}.}\ignorespaces}
\renewtheorem{ex}{\color{gold1}Câu}
\renewcommand{\baselinestretch}{1.1}
\newcommand{\hoac}[1]{
	\left[\begin{aligned}#1\end{aligned}\right.}
\newcommand{\heva}[1]{
	\left\{\begin{aligned}#1\end{aligned}\right.}
\usepackage{esvect}
\def\vec{\vv}
\def\vt{\vv}
\def\overrightarrow{\vv}
%====================== LOGO EDUZY VẼ BẰNG TIKZ ======================
%  Đường viền dò từ logo gốc; dải màu vàng kim đo trực tiếp trên ảnh gốc
%  (nghiêng 21 độ, tô bằng 80 dải hẹp nên hiển thị giống nhau trên mọi trình xem PDF).
%  Dùng: \eduzylogo          -> cao 2.4cm (mặc định)
%        \eduzylogo[1.5cm]   -> cao 1.5cm
\definecolor{ezText}{RGB}{245,200,46}
\newcommand{\eduzylogo}[1][2.4cm]{%
	\resizebox{!}{#1}{%
		\begin{tikzpicture}[x=0.01cm,y=0.01cm]
			\path[use as bounding box] (0,0) rectangle (224,-224);
			% mũ tốt nghiệp + vương miện
			\begin{scope}
				\clip
				(119.4,-176.1)..controls(118.9,-176.1)and(116.8,-176)..(114.6,-176)
				..controls(101.4,-176)and(85.8,-175.9)..(83.7,-175.7)
				..controls(81.8,-175.6)and(81,-175.6)..(80.2,-175.7)
				..controls(78.8,-175.9)and(77.7,-175.9)..(76.2,-175.7)
				..controls(74.9,-175.5)and(70.8,-175.5)..(64.3,-175.7)
				..controls(59.1,-175.9)and(58.3,-175.7)..(58.6,-173.7)
				..controls(58.6,-172.8)and(59.6,-169.6)..(60.3,-168)
				..controls(60.8,-166.9)and(61.4,-165.8)..(62.3,-165)
				..controls(63.1,-164.2)and(64.3,-163.4)..(65,-163.2)
				..controls(65.5,-163.1)and(72.3,-163)..(76.6,-163.1)
				..controls(77.9,-163.1)and(80.1,-163.1)..(81.4,-163)
				..controls(82.8,-162.9)and(86.3,-162.9)..(89.2,-162.9)
				..controls(93,-162.9)and(94.9,-162.8)..(96,-162.7)
				..controls(97.2,-162.5)and(101,-162.5)..(114,-162.4)
				..controls(123,-162.3)and(134.9,-162.2)..(140.2,-162.1)
				..controls(150.1,-162)and(152.2,-162.1)..(153.6,-162.4)
				..controls(154.4,-162.6)and(157.1,-163.8)..(158.1,-164.4)
				..controls(158.8,-164.8)and(159.7,-165.7)..(159.9,-166.3)
				..controls(160,-166.5)and(160.2,-167.1)..(160.2,-167.6)
				..controls(160.3,-168)and(160.6,-169.4)..(160.9,-170.6)
				..controls(161.2,-171.7)and(161.5,-173.1)..(161.6,-173.6)
				..controls(161.8,-174.4)and(161.8,-174.5)..(161.5,-174.9)
				..controls(161.2,-175.6)and(160.8,-175.9)..(159.9,-176)
				..controls(159.2,-176.1)and(154.1,-176.2)..(151.5,-176.2)
				..controls(148.1,-176.1)and(122.6,-176.1)..(121.6,-176.2)
				..controls(120.8,-176.2)and(119.8,-176.2)..(119.4,-176.1)--cycle
				(130.4,-160)..controls(129,-160)and(119.8,-160)..(110.1,-159.9)
				..controls(100.3,-159.8)and(88.5,-159.7)..(83.9,-159.6)
				..controls(77.8,-159.5)and(75.3,-159.5)..(74.5,-159.6)
				..controls(73.9,-159.7)and(72.5,-159.8)..(71.3,-159.8)
				..controls(69.6,-159.8)and(69,-159.8)..(68.5,-159.6)
				..controls(67.6,-159.3)and(66.8,-158.6)..(66.2,-157.6)
				..controls(66,-157.3)and(65.6,-156.5)..(65.3,-156.1)
				..controls(64.6,-155.1)and(64.5,-154.9)..(63.7,-152.9)
				..controls(63.2,-151.6)and(62.9,-150.9)..(61.8,-149.6)
				..controls(61.5,-149.2)and(61.1,-148.5)..(60.9,-147.9)
				..controls(60.7,-147.4)and(60.3,-146.6)..(60.1,-146.1)
				..controls(59.9,-145.5)and(59.5,-144.9)..(59,-144.3)
				..controls(58.5,-143.7)and(58.2,-143.2)..(57.9,-142.5)
				..controls(57.7,-141.9)and(57.3,-141.1)..(57,-140.7)
				..controls(56.6,-140.3)and(56.1,-139.3)..(55.8,-138.5)
				..controls(55.5,-137.8)and(55.1,-136.8)..(54.9,-136.5)
				..controls(54.7,-136.1)and(54.3,-135.2)..(53.9,-134.5)
				..controls(53.6,-133.8)and(53,-132.7)..(52.6,-132)
				..controls(51.6,-130.5)and(51.4,-130.2)..(51.1,-129)
				..controls(50.6,-127.5)and(50.6,-127.4)..(49.7,-126.3)
				..controls(49.2,-125.7)and(48.6,-124.8)..(48.3,-124.2)
				..controls(48,-123.6)and(47.5,-122.7)..(47.3,-122.2)
				..controls(47,-121.8)and(46.5,-120.7)..(46.1,-119.8)
				..controls(45.8,-118.9)and(45.2,-117.8)..(44.9,-117.2)
				..controls(44.2,-116.1)and(43.3,-114.4)..(42.3,-112.2)
				..controls(41.8,-111.3)and(41.3,-110.4)..(41.1,-110.1)
				..controls(40.5,-109.4)and(40.3,-109)..(40,-107.8)
				..controls(39.8,-107)and(39.6,-106.5)..(39.2,-106)
				..controls(38.9,-105.6)and(38.4,-104.7)..(38,-104)
				..controls(37.7,-103.4)and(37.2,-102.5)..(37,-102.1)
				..controls(36.7,-101.7)and(36.3,-100.8)..(36,-100.1)
				..controls(35.6,-98.9)and(35.3,-98.4)..(34.6,-97.6)
				..controls(34.2,-97.1)and(34.2,-96.3)..(34.6,-96)
				..controls(35,-95.7)and(35.5,-95.9)..(36.2,-96.7)
				..controls(36.5,-97)and(37.2,-97.5)..(37.6,-97.7)..controls(38,-98)and(38.6,-98.4)..(39.1,-98.9)
				..controls(39.6,-99.4)and(40.3,-100)..(40.6,-100.3)
				..controls(41.3,-100.8)and(42.4,-101.8)..(43.7,-102.9)
				..controls(44,-103.2)and(44.6,-103.7)..(45,-103.9)
				..controls(45.4,-104.2)and(45.9,-104.6)..(46.1,-104.8)
				..controls(47.1,-105.9)and(48.8,-107.1)..(50.5,-107.9)
				..controls(51.8,-108.6)and(52.3,-109)..(53.4,-110)
				..controls(54.8,-111.2)and(55.6,-111.6)..(57,-112.1)
				..controls(57.5,-112.3)and(58.5,-112.8)..(59.4,-113.3)
				..controls(60.2,-113.7)and(61.3,-114.3)..(61.9,-114.5)
				..controls(62.4,-114.7)and(63.3,-115.1)..(63.8,-115.3)
				..controls(64.4,-115.6)and(65.4,-115.9)..(66.2,-116.1)
				..controls(67,-116.3)and(68.4,-116.6)..(69.2,-116.9)
				..controls(72.3,-117.9)and(72.8,-118)..(78,-118.1)
				..controls(84,-118.3)and(86.1,-118.1)..(88.1,-117.1)
				..controls(89.4,-116.5)and(93,-115)..(93.7,-114.9)
				..controls(94.8,-114.6)and(95.4,-114.2)..(95.9,-113.4)
				..controls(96.2,-113)and(96.7,-112.5)..(97.1,-112.2)
				..controls(98,-111.6)and(101.3,-108.4)..(102.1,-107.2)
				..controls(102.4,-106.8)and(102.9,-106.2)..(103.2,-105.8)
				..controls(103.5,-105.5)and(104,-104.8)..(104.2,-104.2)
				..controls(104.5,-103.6)and(105,-102.8)..(105.3,-102.4)
				..controls(105.6,-102)and(106.1,-101.2)..(106.5,-100.6)
				..controls(107.1,-99.5)and(108.6,-97.5)..(109.5,-96.5)
				..controls(109.9,-96.1)and(110.1,-96)..(110.4,-96)
				..controls(111.1,-96)and(111.6,-96.5)..(111.8,-97.6)
				..controls(112,-98.6)and(112.3,-99.3)..(112.8,-99.9)
				..controls(113.7,-100.9)and(115.2,-103.1)..(115.7,-104.2)
				..controls(116.1,-105)and(116.5,-105.5)..(116.9,-106)
				..controls(117.3,-106.4)and(117.8,-106.9)..(118,-107.3)
				..controls(118.2,-107.6)and(118.6,-108)..(119,-108.3)
				..controls(119.6,-108.8)and(121.1,-110.3)..(121.8,-111.2)
				..controls(122.1,-111.7)and(122.5,-112)..(123.4,-112.4)
				..controls(124.8,-113.2)and(125.8,-113.8)..(127,-114.6)
				..controls(127.5,-114.9)and(128.1,-115.3)..(128.2,-115.4)
				..controls(128.4,-115.5)and(129.3,-115.8)..(130.2,-116)
				..controls(131.1,-116.2)and(132.5,-116.6)..(133.3,-116.9)
				..controls(135.7,-117.9)and(136.1,-117.9)..(142.5,-117.8)
				..controls(145.8,-117.8)and(148.4,-117.7)..(148.8,-117.6)
				..controls(149.2,-117.6)and(150.2,-117.2)..(151,-116.8)
				..controls(152,-116.4)and(152.8,-116)..(153.6,-115.9)
				..controls(154.5,-115.7)and(155.2,-115.4)..(156.3,-114.9)
				..controls(157.1,-114.5)and(158.2,-114)..(158.7,-113.8)
				..controls(159.1,-113.7)and(159.8,-113.3)..(160.2,-113)
				..controls(160.6,-112.7)and(161.3,-112.4)..(161.7,-112.2)
				..controls(162.7,-111.9)and(163.2,-111.7)..(164.1,-111.1)
				..controls(164.5,-110.8)and(165.3,-110.4)..(165.9,-110)
				..controls(166.5,-109.7)and(167.3,-109.2)..(167.6,-108.9)
				..controls(167.9,-108.6)and(168.6,-108.2)..(169,-108)
				..controls(169.4,-107.8)and(170.2,-107.3)..(170.7,-106.9)
				..controls(171.2,-106.5)and(172.1,-105.8)..(172.7,-105.4)
				..controls(173.2,-105)and(174.1,-104.4)..(174.5,-104)
				..controls(174.9,-103.7)and(175.7,-103.1)..(176.2,-102.6)
				..controls(176.7,-102.2)and(177.4,-101.5)..(177.9,-101.2)
				..controls(178.3,-100.8)and(178.8,-100.4)..(179,-100.2)
				..controls(179.2,-100)and(179.6,-99.6)..(180,-99.3)
				..controls(180.4,-99.1)and(180.9,-98.6)..(181.2,-98.4)
				..controls(181.5,-98.1)and(182.1,-97.7)..(182.5,-97.5)
				..controls(183.2,-97.2)and(183.5,-96.9)..(184.3,-95.9)
				..controls(184.6,-95.6)and(184.6,-95.6)..(184.8,-95.7)
				..controls(185,-95.8)and(185.1,-96)..(185.1,-96.2)
				..controls(185.1,-96.5)and(185,-96.7)..(184.6,-97.2)
				..controls(184.1,-97.8)and(184,-97.9)..(183.9,-98.8)
				..controls(183.7,-100.3)and(183.3,-101.3)..(182.9,-101.8)
				..controls(182.6,-102.1)and(182.3,-102.6)..(182.2,-102.9)
				..controls(182,-103.3)and(181.6,-104)..(181.2,-104.6)
				..controls(180.9,-105.2)and(180.5,-106.1)..(180.4,-106.6)
				..controls(179.9,-108)and(179.7,-108.5)..(179,-109.4)
				..controls(178.6,-109.8)and(178.2,-110.5)..(177.9,-111)
				..controls(177.7,-111.5)and(177.3,-112.3)..(177,-112.7)
				..controls(176.6,-113.2)and(176.4,-113.7)..(176.3,-114.2)
				..controls(176.2,-114.7)and(175.9,-115.2)..(175.4,-116)
				..controls(175,-116.6)and(174.4,-117.7)..(174,-118.5)
				..controls(173.7,-119.2)and(173.3,-120.1)..(173,-120.5)
				..controls(172.8,-120.9)and(172.4,-121.8)..(172.1,-122.5)
				..controls(171.8,-123.2)and(171.3,-124.1)..(171,-124.6)
				..controls(170.8,-125.1)and(170.3,-126)..(169.9,-126.6)
				..controls(169.6,-127.3)and(169.2,-128.1)..(169,-128.3)
				..controls(168.8,-128.6)and(168.5,-129.3)..(168.3,-129.8)
				..controls(168.1,-130.3)and(167.8,-131.1)..(167.5,-131.6)
				..controls(166.9,-132.5)and(164.2,-137.9)..(164.2,-138.1)
				..controls(164.2,-138.1)and(163.9,-138.8)..(163.4,-139.7)
				..controls(162.9,-140.5)and(162.4,-141.5)..(162.3,-141.9)
				..controls(161.9,-143)and(161.6,-143.4)..(160.9,-144.1)
				..controls(160.4,-144.6)and(160.3,-144.8)..(160.1,-145.4)
				..controls(160,-145.8)and(159.6,-146.8)..(159.2,-147.8)
				..controls(158.7,-148.8)and(158.1,-150)..(157.9,-150.6)
				..controls(157.6,-151.1)and(157.2,-151.9)..(156.9,-152.4)
				..controls(156.5,-152.8)and(156.1,-153.7)..(155.8,-154.4)
				..controls(155.6,-155)and(155.2,-155.9)..(154.9,-156.3)
				..controls(154.7,-156.7)and(154.4,-157.4)..(154.2,-157.9)
				..controls(154.1,-158.4)and(153.8,-159)..(153.5,-159.3)
				..controls(152.8,-160.1)and(153,-160.1)..(142.4,-160.1)
				..controls(137.2,-160.1)and(131.8,-160.1)..(130.4,-160)--cycle
				(141.5,-103.7)..controls(140.8,-103)and(139.8,-102.1)..(139.2,-101.6)
				..controls(138.5,-101.2)and(137.9,-100.6)..(137.8,-100.4)
				..controls(137,-99.5)and(134,-96.5)..(133.2,-96)
				..controls(132.7,-95.7)and(131.8,-94.9)..(131.1,-94.3)
				..controls(129.4,-92.7)and(127.9,-91.6)..(126.6,-91.1)
				..controls(125.2,-90.5)and(124.1,-90)..(122.9,-89.1)
				..controls(122.1,-88.6)and(121.3,-88.2)..(119.9,-87.7)
				..controls(118.8,-87.3)and(117.5,-86.8)..(116.9,-86.5)
				..controls(114.3,-85.5)and(113.7,-85.4)..(109.7,-85.5)
				..controls(105.3,-85.7)and(103.9,-85.9)..(101.6,-87)
				..controls(100.7,-87.4)and(99.4,-87.8)..(97.9,-88)
				..controls(97,-88.2)and(96.6,-88.3)..(95.8,-88.9)
				..controls(95.4,-89.2)and(94.7,-89.6)..(94.2,-89.8)
				..controls(92.9,-90.3)and(92.3,-90.7)..(90.9,-92.1)
				..controls(90.2,-92.7)and(89.4,-93.3)..(89,-93.5)..controls(87,-94.5)and(85.2,-96)..(83.6,-97.9)
				..controls(83.3,-98.2)and(82.5,-98.9)..(81.8,-99.4)
				..controls(80.7,-100.1)and(80.5,-100.3)..(80,-101.1)
				..controls(78.7,-103.1)and(77.2,-103.9)..(76.4,-103)
				..controls(75.7,-102.2)and(75.7,-103.2)..(75.7,-83.6)
				..controls(75.6,-73.6)and(75.6,-64.8)..(75.5,-64)
				..controls(75.5,-63.2)and(75.5,-61.4)..(75.6,-59.8)--(75.7,-57.1)--(76,-56.8)
				..controls(76.8,-56.2)and(78,-56.3)..(79.8,-57.2)
				..controls(80.4,-57.5)and(81.3,-57.8)..(81.8,-57.9)
				..controls(82.3,-58.1)and(83,-58.4)..(83.6,-58.8)..controls(84.3,-59.2)and(85,-59.5)..(86,-59.8)
				..controls(86.8,-60.1)and(88,-60.6)..(88.7,-60.9)
				..controls(89.3,-61.2)and(90.6,-61.7)..(91.4,-62)
				..controls(92.2,-62.3)and(93.7,-62.9)..(94.6,-63.3)
				..controls(95.6,-63.7)and(96.8,-64.1)..(97.4,-64.3)
				..controls(98,-64.5)and(99,-64.9)..(99.7,-65.2)
				..controls(100.4,-65.6)and(101.5,-66)..(102.3,-66.2)
				..controls(104.2,-66.8)and(105,-67.1)..(106.2,-67.8)
				..controls(107.9,-68.7)and(110.2,-69.4)..(111.3,-69.2)
				..controls(112.5,-69)and(115,-68.2)..(115.6,-67.8)
				..controls(117.3,-66.7)and(117.6,-66.6)..(119.5,-66.2)
				..controls(121.9,-65.8)and(122.1,-65.7)..(123,-65.1)
				..controls(124.1,-64.4)and(124.4,-64.3)..(126,-63.9)
				..controls(127,-63.7)and(127.7,-63.4)..(128.4,-63.1)
				..controls(129,-62.8)and(129.7,-62.5)..(130.6,-62.4)
				..controls(132.1,-62.1)and(134,-61.5)..(135.2,-61)
				..controls(135.7,-60.8)and(136.7,-60.5)..(137.3,-60.3)
				..controls(138.4,-60)and(139.3,-59.7)..(141.1,-59)
				..controls(141.6,-58.7)and(142,-58.6)..(142.5,-58.6)
				..controls(143.2,-58.6)and(143.3,-58.7)..(143.6,-59)
				..controls(144,-59.5)and(144.1,-59.9)..(144,-60.8)
				..controls(144,-61.2)and(144,-71.1)..(144,-82.8)--(143.9,-104.1)--(143.7,-104.5)
				..controls(143.5,-104.7)and(143.3,-104.8)..(143,-104.9)
				..controls(142.6,-104.9)and(142.5,-104.8)..(141.5,-103.7)--cycle
				(27.8,-91.4)..controls(27.1,-90.9)and(26.3,-90)..(26.1,-89.5)
				..controls(26,-89.2)and(25.8,-88.6)..(25.7,-87.9)
				..controls(25.4,-86.2)and(25.6,-84.8)..(26.5,-83.5)
				..controls(27.4,-82.3)and(28.3,-82)..(29.4,-82.3)
				..controls(31.1,-82.8)and(32.1,-84.9)..(32,-87.4)
				..controls(31.9,-89.3)and(31.2,-90.6)..(29.9,-91.4)
				..controls(29.1,-91.9)and(28.6,-91.9)..(27.8,-91.4)--cycle
				(190.1,-91.3)..controls(189.7,-91.2)and(189.3,-90.9)..(189.1,-90.7)
				..controls(188.3,-89.8)and(187.6,-87.7)..(187.6,-86.2)
				..controls(187.6,-85.5)and(187.7,-85.2)..(188.2,-84.4)
				..controls(189.1,-82.6)and(190.1,-81.8)..(191.3,-82.1)
				..controls(192,-82.3)and(193.5,-83.7)..(193.8,-84.4)
				..controls(194.1,-85.1)and(194.2,-87.1)..(194,-88.3)
				..controls(193.8,-89.8)and(193.3,-90.5)..(192.3,-91.3)
				..controls(191.7,-91.7)and(190.9,-91.7)..(190.1,-91.3)--cycle
				(173.7,-83.4)..controls(172.2,-82.5)and(172,-82.5)..(170.6,-82.4)
				..controls(169.7,-82.4)and(168.8,-82.4)..(168.3,-82.5)
				..controls(167.2,-82.7)and(166.9,-82.7)..(166.6,-82.2)
				..controls(166.3,-81.8)and(166.3,-81)..(166.8,-80)
				..controls(167.2,-78.9)and(167.4,-77.4)..(167.7,-74.1)
				..controls(167.8,-72.8)and(167.9,-71.8)..(167.9,-71.7)
				..controls(168,-71.7)and(168,-70.2)..(168,-68.4)..controls(168,-65.4)and(168,-64.9)..(167.7,-63)
				..controls(167.3,-60.5)and(167.1,-59.9)..(166.4,-58.4)
				..controls(166.1,-57.7)and(165.8,-56.9)..(165.7,-56.6)
				..controls(165.7,-56.2)and(165.4,-55.6)..(165.2,-55.2)
				..controls(164.6,-53.8)and(161.9,-50.7)..(161,-50.3)
				..controls(159.8,-49.7)and(157.9,-50)..(155.2,-51.2)
				..controls(154.7,-51.4)and(153.4,-51.8)..(152.2,-52.1)
				..controls(150.8,-52.4)and(149.4,-52.9)..(148.6,-53.2)
				..controls(147.8,-53.6)and(146.7,-54)..(146,-54.2)
				..controls(145.4,-54.4)and(144.5,-54.7)..(144,-54.9)
				..controls(143.6,-55.2)and(142.7,-55.5)..(141.8,-55.7)
				..controls(141,-55.9)and(139.6,-56.3)..(138.7,-56.6)
				..controls(137.8,-57)and(136.3,-57.4)..(135.3,-57.6)
				..controls(133.8,-57.9)and(133.2,-58.1)..(132.1,-58.7)
				..controls(131.1,-59.2)and(130.3,-59.4)..(129.4,-59.6)
				..controls(128.7,-59.8)and(127.8,-60)..(127.4,-60.1)
				..controls(127,-60.3)and(126,-60.5)..(125.1,-60.7)
				..controls(123.2,-61.1)and(121.7,-61.5)..(120.2,-62.1)
				..controls(119.7,-62.3)and(118.1,-62.8)..(116.7,-63.2)
				..controls(114.6,-63.9)and(114,-64.1)..(113,-64.7)
				..controls(111.9,-65.4)and(111.8,-65.4)..(111,-65.5)
				..controls(109.9,-65.5)and(109.3,-65.4)..(107.7,-64.7)
				..controls(107,-64.4)and(106,-64)..(105.3,-63.9)
				..controls(103.5,-63.5)and(102.2,-63)..(101.2,-62.5)
				..controls(100.1,-61.8)and(99.3,-61.5)..(97.8,-61.1)
				..controls(97.2,-61)and(96.2,-60.7)..(95.7,-60.4)
				..controls(95.1,-60.1)and(94.2,-59.8)..(93.5,-59.6)
				..controls(92.9,-59.4)and(91.7,-59)..(91,-58.7)..controls(90.2,-58.4)and(88.9,-58)..(88,-57.7)
				..controls(87.2,-57.5)and(86.1,-57.1)..(85.5,-56.8)
				..controls(85,-56.6)and(84.1,-56.2)..(83.5,-56)..controls(82.9,-55.7)and(81.8,-55.2)..(81,-54.9)
				..controls(80.1,-54.5)and(79.1,-54.2)..(78.7,-54.1)
				..controls(77.9,-53.8)and(77.3,-53.6)..(76.2,-52.9)
				..controls(75.7,-52.6)and(75.2,-52.4)..(74.1,-52.2)
				..controls(72.5,-51.9)and(72.4,-51.8)..(71.3,-51.1)
				..controls(69.9,-50.1)and(69.4,-49.9)..(67.9,-49.8)
				..controls(67,-49.7)and(66.6,-49.6)..(66.3,-49.4)
				..controls(66.1,-49.3)and(65.7,-49)..(65.4,-48.8)
				..controls(65.1,-48.7)and(64.2,-48.4)..(63.6,-48.2)
				..controls(62.9,-48)and(61.7,-47.6)..(60.9,-47.2)
				..controls(60.1,-46.8)and(58.9,-46.4)..(58.2,-46.2)
				..controls(57.6,-46)and(56.8,-45.7)..(56.5,-45.5)
				..controls(54.9,-44.6)and(54.6,-44.5)..(53.2,-44.2)
				..controls(51.7,-43.8)and(51,-43.6)..(49.9,-42.9)
				..controls(49.5,-42.6)and(48.9,-42.4)..(47.5,-42)
				..controls(45.7,-41.6)and(45.6,-41.5)..(44.3,-40.6)
				..controls(43.5,-40.1)and(42.8,-39.6)..(42.7,-39.4)
				..controls(42.4,-38.7)and(43,-38.1)..(44.5,-37.6)..controls(44.9,-37.5)and(46,-37)..(46.9,-36.6)
				..controls(47.9,-36.1)and(48.8,-35.8)..(49.6,-35.6)
				..controls(50.5,-35.4)and(51.3,-35.2)..(52.4,-34.6)
				..controls(53.3,-34.1)and(54.5,-33.7)..(54.9,-33.5)
				..controls(55.4,-33.4)and(56.4,-33)..(57.2,-32.6)
				..controls(58.6,-32)and(60.1,-31.5)..(61.6,-31.2)
				..controls(62.3,-31.1)and(62.7,-31)..(63.2,-30.7)
				..controls(64.1,-30)and(64.3,-29.9)..(65.3,-29.7)
				..controls(65.7,-29.6)and(66.5,-29.2)..(67.1,-28.9)
				..controls(67.6,-28.6)and(68.8,-28.2)..(69.6,-27.9)
				..controls(70.5,-27.7)and(71.6,-27.3)..(72.1,-27)
				..controls(72.6,-26.7)and(73.8,-26.2)..(74.8,-25.9)
				..controls(75.7,-25.5)and(76.7,-25.1)..(77.1,-25)
				..controls(77.4,-24.8)and(78.5,-24.4)..(79.6,-24.1)
				..controls(80.6,-23.8)and(82,-23.2)..(82.8,-22.9)
				..controls(83.5,-22.5)and(84.8,-22)..(85.8,-21.8)
				..controls(86.7,-21.5)and(87.8,-21.1)..(88.1,-20.9)
				..controls(88.4,-20.7)and(89.3,-20.3)..(90.1,-20.1)
				..controls(90.8,-19.8)and(91.9,-19.4)..(92.4,-19.1)
				..controls(93.5,-18.5)and(94.4,-18.2)..(95.4,-18.1)
				..controls(95.7,-18)and(96.4,-17.8)..(96.9,-17.5)
				..controls(97.5,-17.3)and(98.4,-16.9)..(99,-16.6)
				..controls(99.7,-16.3)and(100.9,-15.9)..(101.9,-15.6)
				..controls(102.9,-15.3)and(104.4,-14.7)..(105.3,-14.3)
				..controls(108.1,-13.2)and(108.5,-13.1)..(110.1,-13.1)
				..controls(111.4,-13.1)and(111.5,-13.1)..(112.4,-13.5)
				..controls(113.2,-14)and(115,-14.8)..(116.6,-15.4)
				..controls(116.9,-15.5)and(117.9,-15.8)..(118.8,-16.1)
				..controls(120.9,-16.6)and(121.5,-16.9)..(122.6,-17.9)
				..controls(123,-18.3)and(123.7,-18.8)..(123.9,-18.9)
				..controls(124.6,-19.4)and(126,-19.8)..(127.9,-20.2)
				..controls(129.6,-20.5)and(129.8,-20.6)..(131.1,-21.4)
				..controls(131.5,-21.7)and(132.4,-22.1)..(133.1,-22.3)
				..controls(136.1,-23.2)and(138.3,-24.1)..(140.2,-25)
				..controls(141,-25.5)and(141.5,-25.6)..(142.5,-25.8)
				..controls(145.1,-26.3)and(146,-26.6)..(149.4,-28.1)
				..controls(151.3,-28.9)and(153.5,-29.8)..(154.3,-30.1)
				..controls(155.1,-30.4)and(156.2,-30.8)..(156.6,-31)
				..controls(157.1,-31.2)and(158.2,-31.6)..(159.1,-31.8)
				..controls(160.6,-32.2)and(160.9,-32.3)..(161.7,-32.9)
				..controls(162.4,-33.3)and(162.9,-33.6)..(163.7,-33.8)
				..controls(164.3,-34)and(164.8,-34.1)..(164.9,-34.1)
				..controls(165,-34.1)and(165.6,-34.4)..(166.4,-34.8)
				..controls(167.3,-35.2)and(168.1,-35.4)..(169,-35.6)
				..controls(170,-35.8)and(170.5,-36)..(171.7,-36.6)
				..controls(172.4,-37)and(173.3,-37.4)..(173.6,-37.4)
				..controls(175.6,-37.7)and(176.2,-38)..(176.7,-39.1)
				..controls(176.9,-39.6)and(176.9,-39.7)..(176.7,-40.4)
				..controls(176.4,-41.6)and(175.8,-42.7)..(175.2,-43.4)
				..controls(174.5,-44.1)and(172.6,-45.5)..(171.7,-45.8)
				..controls(170.6,-46.2)and(169.1,-46)..(168.2,-45.4)
				..controls(167.9,-45.2)and(167.3,-44.8)..(166.9,-44.4)
				..controls(166.4,-44.1)and(165.8,-43.6)..(165.4,-43.2)
				..controls(165.1,-42.9)and(164.4,-42.3)..(163.9,-41.9)
				..controls(163.3,-41.5)and(162.7,-40.9)..(162.4,-40.5)
				..controls(161.4,-39.2)and(159.9,-38)..(158.8,-37.8)
				..controls(158.4,-37.8)and(158.3,-37.7)..(158.2,-37.2)
				..controls(158,-36.8)and(157.4,-36.1)..(157.1,-36.1)
				..controls(157,-36.1)and(156.7,-36.1)..(156.5,-36)
				..controls(156.1,-35.9)and(155.9,-35.7)..(155.7,-35.2)
				..controls(155.3,-34.6)and(155.1,-34.4)..(154,-34.1)
				..controls(153.6,-34)and(153.2,-33.7)..(152.9,-33.5)
				..controls(152.4,-32.9)and(151.6,-32.4)..(151.3,-32.4)
				..controls(150.9,-32.4)and(150.4,-32.6)..(150.4,-32.9)
				..controls(150.4,-33.2)and(150.8,-33.4)..(151.4,-33.6)
				..controls(152.1,-33.7)and(152.8,-34.1)..(153.5,-34.9)
				..controls(154.1,-35.5)and(154.3,-35.6)..(155,-35.8)
				..controls(155.7,-36)and(155.7,-36)..(156,-36.6)
				..controls(156.3,-37.3)and(156.7,-37.7)..(157.5,-37.8)
				..controls(158.1,-37.9)and(158.2,-38)..(158.4,-38.8)
				..controls(158.5,-39.4)and(158.6,-39.6)..(159.6,-40.4)
				..controls(160.2,-41)and(160.9,-41.6)..(161.3,-41.8)
				..controls(161.7,-42)and(161.9,-42.3)..(162.2,-42.8)
				..controls(162.4,-43.2)and(162.8,-43.8)..(163.1,-44.1)
				..controls(163.9,-44.8)and(165.1,-46.7)..(165.7,-48)
				..controls(166.2,-49.3)and(166.4,-49.7)..(167.2,-50.5)
				..controls(167.5,-50.8)and(167.9,-51.4)..(168.1,-51.8)
				..controls(168.3,-52.2)and(168.8,-53)..(169.1,-53.6)
				..controls(169.6,-54.3)and(169.8,-54.9)..(169.9,-55.6)
				..controls(170.4,-57.8)and(170.5,-58.1)..(171.2,-59.5)
				..controls(171.6,-60.3)and(172,-61.5)..(172.2,-62.1)
				..controls(172.3,-62.8)and(172.7,-63.7)..(172.9,-64.1)
				..controls(173.2,-64.6)and(173.4,-65.3)..(173.4,-65.6)
				..controls(173.8,-68.8)and(174.2,-70.6)..(174.9,-72.6)
				..controls(175.5,-74.5)and(175.9,-77.1)..(176.1,-80.2)
				..controls(176.3,-83.4)and(175.5,-84.4)..(173.7,-83.4)--cycle;
				\definecolor{ezs}{RGB}{247,179,26}\fill[ezs](-5.7,-156.6)--(-3.4,-157.5)--(61.9,12.7)--(59.7,13.6)--cycle;
				\definecolor{ezs}{RGB}{247,179,22}\fill[ezs](-3.8,-157.3)--(-1.5,-158.2)--(63.9,12)--(61.6,12.9)--cycle;
				\definecolor{ezs}{RGB}{248,179,19}\fill[ezs](-1.8,-158.1)--(0.4,-158.9)--(65.8,11.2)--(63.5,12.1)--cycle;
				\definecolor{ezs}{RGB}{248,178,16}\fill[ezs](0.1,-158.8)--(2.4,-159.7)--(67.7,10.5)--(65.4,11.4)--cycle;
				\definecolor{ezs}{RGB}{249,178,15}\fill[ezs](2,-159.5)--(4.3,-160.4)--(69.6,9.8)--(67.3,10.7)--cycle;
				\definecolor{ezs}{RGB}{251,177,17}\fill[ezs](3.9,-160.3)--(6.2,-161.1)--(71.5,9)--(69.2,9.9)--cycle;
				\definecolor{ezs}{RGB}{252,179,18}\fill[ezs](5.8,-161)--(8.1,-161.9)--(73.4,8.3)--(71.1,9.2)--cycle;
				\definecolor{ezs}{RGB}{252,182,17}\fill[ezs](7.7,-161.7)--(10,-162.6)--(75.3,7.6)--(73,8.5)--cycle;
				\definecolor{ezs}{RGB}{252,185,17}\fill[ezs](9.6,-162.5)--(11.9,-163.4)--(77.2,6.8)--(75,7.7)--cycle;
				\definecolor{ezs}{RGB}{251,189,18}\fill[ezs](11.5,-163.2)--(13.8,-164.1)--(79.2,6.1)--(76.9,7)--cycle;
				\definecolor{ezs}{RGB}{251,191,22}\fill[ezs](13.4,-163.9)--(15.7,-164.8)--(81.1,5.4)--(78.8,6.2)--cycle;
				\definecolor{ezs}{RGB}{252,193,25}\fill[ezs](15.4,-164.7)--(17.6,-165.6)--(83,4.6)--(80.7,5.5)--cycle;
				\definecolor{ezs}{RGB}{253,195,27}\fill[ezs](17.3,-165.4)--(19.6,-166.3)--(84.9,3.9)--(82.6,4.8)--cycle;
				\definecolor{ezs}{RGB}{253,196,31}\fill[ezs](19.2,-166.1)--(21.5,-167)--(86.8,3.2)--(84.5,4)--cycle;
				\definecolor{ezs}{RGB}{254,199,35}\fill[ezs](21.1,-166.9)--(23.4,-167.8)--(88.7,2.4)--(86.4,3.3)--cycle;
				\definecolor{ezs}{RGB}{253,201,38}\fill[ezs](23,-167.6)--(25.3,-168.5)--(90.6,1.7)--(88.3,2.6)--cycle;
				\definecolor{ezs}{RGB}{253,203,40}\fill[ezs](24.9,-168.3)--(27.2,-169.2)--(92.5,1)--(90.3,1.8)--cycle;
				\definecolor{ezs}{RGB}{253,205,42}\fill[ezs](26.8,-169.1)--(29.1,-170)--(94.4,0.2)--(92.2,1.1)--cycle;
				\definecolor{ezs}{RGB}{253,207,45}\fill[ezs](28.7,-169.8)--(31,-170.7)--(96.4,-0.5)--(94.1,0.4)--cycle;
				\definecolor{ezs}{RGB}{253,210,49}\fill[ezs](30.7,-170.5)--(32.9,-171.4)--(98.3,-1.2)--(96,-0.4)--cycle;
				\definecolor{ezs}{RGB}{253,212,54}\fill[ezs](32.6,-171.3)--(34.9,-172.2)--(100.2,-2)--(97.9,-1.1)--cycle;
				\definecolor{ezs}{RGB}{253,214,59}\fill[ezs](34.5,-172)--(36.8,-172.9)--(102.1,-2.7)--(99.8,-1.8)--cycle;
				\definecolor{ezs}{RGB}{253,216,63}\fill[ezs](36.4,-172.7)--(38.7,-173.6)--(104,-3.4)--(101.7,-2.6)--cycle;
				\definecolor{ezs}{RGB}{253,218,65}\fill[ezs](38.3,-173.5)--(40.6,-174.4)--(105.9,-4.2)--(103.6,-3.3)--cycle;
				\definecolor{ezs}{RGB}{254,220,66}\fill[ezs](40.2,-174.2)--(42.5,-175.1)--(107.8,-4.9)--(105.5,-4)--cycle;
				\definecolor{ezs}{RGB}{254,222,68}\fill[ezs](42.1,-175)--(44.4,-175.8)--(109.7,-5.6)--(107.5,-4.8)--cycle;
				\definecolor{ezs}{RGB}{254,224,71}\fill[ezs](44,-175.7)--(46.3,-176.6)--(111.7,-6.4)--(109.4,-5.5)--cycle;
				\definecolor{ezs}{RGB}{254,226,75}\fill[ezs](46,-176.4)--(48.2,-177.3)--(113.6,-7.1)--(111.3,-6.2)--cycle;
				\definecolor{ezs}{RGB}{254,228,79}\fill[ezs](47.9,-177.2)--(50.1,-178)--(115.5,-7.8)--(113.2,-7)--cycle;
				\definecolor{ezs}{RGB}{254,230,82}\fill[ezs](49.8,-177.9)--(52.1,-178.8)--(117.4,-8.6)--(115.1,-7.7)--cycle;
				\definecolor{ezs}{RGB}{254,232,85}\fill[ezs](51.7,-178.6)--(54,-179.5)--(119.3,-9.3)--(117,-8.4)--cycle;
				\definecolor{ezs}{RGB}{254,234,88}\fill[ezs](53.6,-179.4)--(55.9,-180.2)--(121.2,-10)--(118.9,-9.2)--cycle;
				\definecolor{ezs}{RGB}{254,237,91}\fill[ezs](55.5,-180.1)--(57.8,-181)--(123.1,-10.8)--(120.8,-9.9)--cycle;
				\definecolor{ezs}{RGB}{254,239,94}\fill[ezs](57.4,-180.8)--(59.7,-181.7)--(125,-11.5)--(122.8,-10.6)--cycle;
				\definecolor{ezs}{RGB}{254,241,97}\fill[ezs](59.3,-181.6)--(61.6,-182.4)--(127,-12.2)--(124.7,-11.4)--cycle;
				\definecolor{ezs}{RGB}{254,243,101}\fill[ezs](61.2,-182.3)--(63.5,-183.2)--(128.9,-13)--(126.6,-12.1)--cycle;
				\definecolor{ezs}{RGB}{254,245,103}\fill[ezs](63.2,-183)--(65.4,-183.9)--(130.8,-13.7)--(128.5,-12.8)--cycle;
				\definecolor{ezs}{RGB}{254,245,105}\fill[ezs](65.1,-183.8)--(67.4,-184.6)--(132.7,-14.4)--(130.4,-13.6)--cycle;
				\definecolor{ezs}{RGB}{254,245,105}\fill[ezs](67,-184.5)--(69.3,-185.4)--(134.6,-15.2)--(132.3,-14.3)--cycle;
				\definecolor{ezs}{RGB}{254,244,103}\fill[ezs](68.9,-185.2)--(71.2,-186.1)--(136.5,-15.9)--(134.2,-15)--cycle;
				\definecolor{ezs}{RGB}{253,243,101}\fill[ezs](70.8,-186)--(73.1,-186.8)--(138.4,-16.6)--(136.1,-15.8)--cycle;
				\definecolor{ezs}{RGB}{253,241,98}\fill[ezs](72.7,-186.7)--(75,-187.6)--(140.3,-17.4)--(138,-16.5)--cycle;
				\definecolor{ezs}{RGB}{253,239,95}\fill[ezs](74.6,-187.4)--(76.9,-188.3)--(142.2,-18.1)--(140,-17.2)--cycle;
				\definecolor{ezs}{RGB}{253,236,92}\fill[ezs](76.5,-188.2)--(78.8,-189)--(144.2,-18.8)--(141.9,-18)--cycle;
				\definecolor{ezs}{RGB}{253,234,88}\fill[ezs](78.5,-188.9)--(80.7,-189.8)--(146.1,-19.6)--(143.8,-18.7)--cycle;
				\definecolor{ezs}{RGB}{253,232,85}\fill[ezs](80.4,-189.6)--(82.7,-190.5)--(148,-20.3)--(145.7,-19.4)--cycle;
				\definecolor{ezs}{RGB}{253,230,82}\fill[ezs](82.3,-190.4)--(84.6,-191.2)--(149.9,-21)--(147.6,-20.2)--cycle;
				\definecolor{ezs}{RGB}{253,228,78}\fill[ezs](84.2,-191.1)--(86.5,-192)--(151.8,-21.8)--(149.5,-20.9)--cycle;
				\definecolor{ezs}{RGB}{253,226,74}\fill[ezs](86.1,-191.8)--(88.4,-192.7)--(153.7,-22.5)--(151.4,-21.6)--cycle;
				\definecolor{ezs}{RGB}{253,224,70}\fill[ezs](88,-192.6)--(90.3,-193.4)--(155.6,-23.3)--(153.3,-22.4)--cycle;
				\definecolor{ezs}{RGB}{253,222,67}\fill[ezs](89.9,-193.3)--(92.2,-194.2)--(157.5,-24)--(155.3,-23.1)--cycle;
				\definecolor{ezs}{RGB}{253,220,65}\fill[ezs](91.8,-194)--(94.1,-194.9)--(159.5,-24.7)--(157.2,-23.8)--cycle;
				\definecolor{ezs}{RGB}{253,217,63}\fill[ezs](93.7,-194.8)--(96,-195.6)--(161.4,-25.5)--(159.1,-24.6)--cycle;
				\definecolor{ezs}{RGB}{253,215,60}\fill[ezs](95.7,-195.5)--(97.9,-196.4)--(163.3,-26.2)--(161,-25.3)--cycle;
				\definecolor{ezs}{RGB}{253,213,56}\fill[ezs](97.6,-196.2)--(99.9,-197.1)--(165.2,-26.9)--(162.9,-26)--cycle;
				\definecolor{ezs}{RGB}{252,211,53}\fill[ezs](99.5,-197)--(101.8,-197.8)--(167.1,-27.7)--(164.8,-26.8)--cycle;
				\definecolor{ezs}{RGB}{252,209,49}\fill[ezs](101.4,-197.7)--(103.7,-198.6)--(169,-28.4)--(166.7,-27.5)--cycle;
				\definecolor{ezs}{RGB}{252,207,45}\fill[ezs](103.3,-198.4)--(105.6,-199.3)--(170.9,-29.1)--(168.6,-28.2)--cycle;
				\definecolor{ezs}{RGB}{253,205,42}\fill[ezs](105.2,-199.2)--(107.5,-200)--(172.8,-29.9)--(170.6,-29)--cycle;
				\definecolor{ezs}{RGB}{253,202,40}\fill[ezs](107.1,-199.9)--(109.4,-200.8)--(174.7,-30.6)--(172.5,-29.7)--cycle;
				\definecolor{ezs}{RGB}{253,200,38}\fill[ezs](109,-200.6)--(111.3,-201.5)--(176.7,-31.3)--(174.4,-30.4)--cycle;
				\definecolor{ezs}{RGB}{253,198,35}\fill[ezs](111,-201.4)--(113.2,-202.2)--(178.6,-32.1)--(176.3,-31.2)--cycle;
				\definecolor{ezs}{RGB}{253,196,31}\fill[ezs](112.9,-202.1)--(115.2,-203)--(180.5,-32.8)--(178.2,-31.9)--cycle;
				\definecolor{ezs}{RGB}{253,194,28}\fill[ezs](114.8,-202.8)--(117.1,-203.7)--(182.4,-33.5)--(180.1,-32.6)--cycle;
				\definecolor{ezs}{RGB}{253,192,26}\fill[ezs](116.7,-203.6)--(119,-204.5)--(184.3,-34.3)--(182,-33.4)--cycle;
				\definecolor{ezs}{RGB}{253,189,24}\fill[ezs](118.6,-204.3)--(120.9,-205.2)--(186.2,-35)--(183.9,-34.1)--cycle;
				\definecolor{ezs}{RGB}{253,186,20}\fill[ezs](120.5,-205)--(122.8,-205.9)--(188.1,-35.7)--(185.8,-34.8)--cycle;
				\definecolor{ezs}{RGB}{253,184,17}\fill[ezs](122.4,-205.8)--(124.7,-206.7)--(190,-36.5)--(187.8,-35.6)--cycle;
				\definecolor{ezs}{RGB}{252,183,15}\fill[ezs](124.3,-206.5)--(126.6,-207.4)--(192,-37.2)--(189.7,-36.3)--cycle;
				\definecolor{ezs}{RGB}{252,181,14}\fill[ezs](126.3,-207.2)--(128.5,-208.1)--(193.9,-37.9)--(191.6,-37.1)--cycle;
				\definecolor{ezs}{RGB}{252,179,13}\fill[ezs](128.2,-208)--(130.4,-208.9)--(195.8,-38.7)--(193.5,-37.8)--cycle;
				\definecolor{ezs}{RGB}{252,178,12}\fill[ezs](130.1,-208.7)--(132.4,-209.6)--(197.7,-39.4)--(195.4,-38.5)--cycle;
				\definecolor{ezs}{RGB}{251,178,10}\fill[ezs](132,-209.4)--(134.3,-210.3)--(199.6,-40.1)--(197.3,-39.3)--cycle;
				\definecolor{ezs}{RGB}{251,178,10}\fill[ezs](133.9,-210.2)--(136.2,-211.1)--(201.5,-40.9)--(199.2,-40)--cycle;
				\definecolor{ezs}{RGB}{250,177,11}\fill[ezs](135.8,-210.9)--(138.1,-211.8)--(203.4,-41.6)--(201.1,-40.7)--cycle;
				\definecolor{ezs}{RGB}{250,178,11}\fill[ezs](137.7,-211.6)--(140,-212.5)--(205.3,-42.3)--(203.1,-41.5)--cycle;
				\definecolor{ezs}{RGB}{248,180,10}\fill[ezs](139.6,-212.4)--(141.9,-213.3)--(207.3,-43.1)--(205,-42.2)--cycle;
				\definecolor{ezs}{RGB}{247,181,9}\fill[ezs](141.5,-213.1)--(143.8,-214)--(209.2,-43.8)--(206.9,-42.9)--cycle;
				\definecolor{ezs}{RGB}{246,182,9}\fill[ezs](143.5,-213.8)--(145.7,-214.7)--(211.1,-44.5)--(208.8,-43.7)--cycle;
				\definecolor{ezs}{RGB}{246,182,10}\fill[ezs](145.4,-214.6)--(147.7,-215.5)--(213,-45.3)--(210.7,-44.4)--cycle;
			\end{scope}
			% chữ EDUZY (Times đậm, làm dày nét bằng 5 lớp lệch nhẹ; căn đúng khung chữ logo gốc)
			\foreach \dx/\dy in {-0.7/0,0.7/0,0/-0.3,0/0.3,0/0}
			\node[anchor=center,inner sep=0pt,outer sep=0pt,text=ezText,align=center] at (109.6+\dx,-193.6+\dy)
			{\resizebox{1.01cm}{0.22cm}{\fontencoding{T1}\fontfamily{ptm}\fontseries{b}\selectfont EDUZY}};
		\end{tikzpicture}%
	}%
}
%====================== THIẾT LẬP THÂN ĐỀ ======================
%--- Môi trường tự luận {bt}: ex_test 3.3.4 chưa khai báo sẵn nên tạo ở đây
%    (ex_test đã tự móc \fixshowans và \listenumerate cho bt)
\makeatletter
\@ifundefined{bt}{\newtheorem{bt}{\color{gold1}\namebt}}{}
\makeatother
\renewcommand{\sectionmark}[1]{\markright{\thesection\ #1}}
%--- Khung tiêu đề (logo + tối đa 4 dòng; dòng 3, 4 để trống {} thì bỏ qua)
%  Đề thi:    \khungtieude{THI THỬ THPT QUỐC GIA NĂM 2026}{Môn Toán -- Đề số 21}
%                         {(Đề thi gồm: ...)}{Thời gian làm bài: 90 phút (không kể thời gian phát đề).}
%  Chuyên đề: \khungtieude{CHUYÊN ĐỀ 3 -- DÃY SỐ VÀ CÔNG THỨC TRUY HỒI}
%                         {Bồi dưỡng học sinh giỏi Toán THPT}{...}{}
\newcommand{\khungtieude}[4]{%
	\begin{tcolorbox}[colframe=gold1,colback=gold5!45,boxrule=0.5mm]
		\begin{minipage}[c]{3cm}
			\centerline{\eduzylogo[2.4cm]}
		\end{minipage}\hspace{0.5cm}
		\begin{minipage}[c]{12.5cm}
			\vskip 0.15 cm
			\centerline{\ind{#1}}
			\vskip 0.15 cm
			\centerline{\textbf{#2}}
			\if\relax\detokenize{#3}\relax\else
				\vskip 0.15 cm
				\centerline{\textit{#3}}
			\fi
			\if\relax\detokenize{#4}\relax\else
				\vskip 0.15 cm
				\centerline{\textit{#4}}
			\fi
			\vskip 0.15 cm
		\end{minipage}
	\end{tcolorbox}
}
%====================== BỔ SUNG CHO BÀI HỌC, CHUYÊN ĐỀ ======================
\usepackage[shortlabels]{enumitem}
\usepackage{tabularx}
%--- Danh sách gọn: \begin{enumerate}[gon] ... \end{enumerate}
\SetEnumitemKey{gon}{leftmargin=*,itemsep=1pt,parsep=0pt,topsep=3pt}
%--- Khung lý thuyết: \begin{khung}{} ... \end{khung}; tiêu đề chỉ dùng khi thật ngắn ("Mẹo", "Lỗi thường gặp")
\newtcolorbox{khung}[1]{breakable,enhanced,colframe=gold1,colback=gold5!45,
	boxrule=0.4mm,fonttitle=\bfseries\sffamily,coltitle=white,colbacktitle=gold1,
	attach boxed title to top left={xshift=6mm,yshift=-3mm},
	boxed title style={arc=1mm,boxrule=0.3mm},title={#1}}
%--- Tiêu đề mục nhỏ trong một dạng: \muc{Dấu hiệu}, \muc{Phương pháp}, \muc{Bài tự luyện}
\makeatletter
\newcommand{\muc}[1]{\par\medskip\noindent\indam{#1}\par\nobreak\smallskip\@afterindentfalse\@afterheading}
\makeatother
%--- Ô công thức cần nhớ (căn giữa, viền vàng, vừa khít nội dung)
%     \begin{center}\congthuc{$Q(n)=(n-1)!$}\end{center}
\newtcbox{\congthuc}{on line,colframe=gold1,colback=gold5!45,boxrule=0.5pt,arc=2pt,
	left=8pt,right=8pt,top=4pt,bottom=4pt}
%--- Công tắc lời giải bài tự luyện (ví dụ mẫu luôn dùng \loigiai)
%     \loigiaiTL[{đáp số}]{lời giải đầy đủ}
%     \TLfull -> in lời giải đầy đủ;  \TLdapso -> chỉ in một dòng "Đáp số";  \TLan -> ẩn hoàn toàn
%     (\CheDo là dethi thì mọi lời giải và đáp số đều ẩn)
\makeatletter
\newif\ifTLhien   \TLhientrue
\newif\ifTLdapso  \TLdapsofalse
\newcommand{\TLfull}{\TLhientrue\TLdapsofalse}
\newcommand{\TLdapso}{\TLhientrue\TLdapsotrue}
\newcommand{\TLan}{\TLhienfalse}
\newcommand{\loigiaiTL}[2][]{%
	\ifTLhien
		\ifTLdapso
			\def\@tlarg{#1}%
			\ifx\@tlarg\@empty\else
				\ifdefstring{\kindOPTN}{dethi}{}{\par\noindent{\bfseries\color{gold1}Đáp số:} #1\par}%
			\fi
		\else
			\loigiai{#2}%
		\fi
	\fi}
\makeatother
\TLdapso% mặc định: chỉ in đáp số bài tự luyện
%====================== BẢNG ĐÚNG/SAI (\choiceTFt) ======================
%  Hàng tiêu đề nền nâu đồng chữ trắng (Phát biểu | Đ | S), các hàng tô xen kẽ kem / trắng,
%  đường kẻ vàng kim, nhãn a) b) c) d) màu nâu đồng. Bản có đáp án: cột được chọn có vòng tròn
%  nâu đồng với dấu tích trắng (câu đúng tích ở cột Đúng, câu sai tích ở cột Sai).
%  Tiêu đề ghi đủ "Đúng | Sai": \OPTN{viettat=0} sau \input (cột Đúng/Sai rộng ra, phát biểu hẹp lại).
%  Dùng bảng gốc của ex_test: \def\BangDSgoc{1} trước \input{preamble-eduzy-v2}.
\providecommand{\BangDSgoc}{0}
\ifnum\BangDSgoc=0
\usepackage{colortbl,pifont}
\colorlet{ezDSke}{gold2!60!gold1}% đường kẻ bảng (nét 0,4pt như ex_test để không đổi chiều cao bảng)
\colorlet{ezDSle}{gold5!60}% hàng a), c)
\colorlet{ezDSchan}{white}% hàng b), d)
%--- dấu chọn ở cột Đúng/Sai
%    (bề rộng 0, căn giữa ô: không làm đổi bề rộng cột của ex_test, không báo tràn lề)
\newcommand{\EZdauchon}{\makebox[0pt]{\tikz[baseline=-1ex]\node[circle,fill=gold1,inner sep=0pt,
	minimum size=1.2em,text=white,font=\small]{\ding{51}};}}
\def\tickT{\EZdauchon}\def\tickF{\EZdauchon}
\AtBeginDocument{\sbox{\ccTFTab}{\EZdauchon}}
%--- nhãn a) b) c) d) màu nâu đồng (ở bản lời giải, nhãn câu đúng vẫn đóng khung như ex_test)
\renewcommand{\FalseTF}{\stepcounter{dapan}\textbf{\color{gold1}\DapAnTF\sepTF}\ignorespaces}
%--- trước khi dựng bảng: kẻ màu vàng kim (bề rộng cột giữ nguyên như ex_test); sau bảng trả lại màu đen
\newcount\EZhangTF
\newcommand{\EZbatdauTF}{%
	\global\EZhangTF=0\relax
	\arrayrulecolor{ezDSke}}
\newcommand{\EZketthucTF}{\arrayrulecolor{black}}
%--- hàng tiêu đề nền nâu đồng, chữ trắng
\makeatletter
\newcommand{\EZo@TF}[1]{\centering\bfseries\color{white}#1}
\newcommand{\EZtieudeTF}{%
	\ifcase\cotDS
		\ifnum\cotSTT=1\relax
			\def\rowHeader{\hline\rowcolor{gold1}\EZo@TF{TT}&\EZo@TF{\phatbieu}\tabularnewline}%
		\else
			\def\rowHeader{\hline\rowcolor{gold1}\EZo@TF{\phatbieu}\tabularnewline}%
		\fi
	\or
		\ifnum\cotSTT=1\relax
			\def\rowHeader{\hline\rowcolor{gold1}\EZo@TF{TT}&\EZo@TF{\phatbieu}&\EZo@TF{\Dung}\tabularnewline}%
		\else
			\def\rowHeader{\hline\rowcolor{gold1}\EZo@TF{\phatbieu}&\EZo@TF{\Dung}\tabularnewline}%
		\fi
	\or
		\ifnum\cotSTT=1\relax
			\def\rowHeader{\hline\rowcolor{gold1}\EZo@TF{TT}&\EZo@TF{\phatbieu}&\EZo@TF{\Dung}&\EZo@TF{\Sai}\tabularnewline}%
		\else
			\def\rowHeader{\hline\rowcolor{gold1}\EZo@TF{\phatbieu}&\EZo@TF{\Dung}&\EZo@TF{\Sai}\tabularnewline}%
		\fi
	\fi}
%--- tô xen kẽ từng hàng (gọi ngay sau \hline đầu mỗi hàng)
\newcommand{\EZmauhangTF}{%
	\noalign{\global\advance\EZhangTF\@ne}%
	\ifodd\EZhangTF\rowcolor{ezDSle}\else\rowcolor{ezDSchan}\fi}
%--- gắn vào \choiceTFt / \writeChoiceTFt của ex_test
\xpretocmd{\choiceTFt}{\EZbatdauTF}{}{\typeout{[preamble-eduzy-v2] Không gắn được bảng Đúng/Sai (đầu bảng)}}
\xpatchcmd{\choiceTFt}{\ifnum\dongPB=0\def\rowHeader{}\fi}{\EZtieudeTF\ifnum\dongPB=0\def\rowHeader{}\fi}
	{}{\typeout{[preamble-eduzy-v2] Không gắn được bảng Đúng/Sai (hàng tiêu đề)}}
\xpatchcmd{\writeChoiceTFt}{\hline\xdef\dotEX{\sdotEX}}{\hline\EZmauhangTF\xdef\dotEX{\sdotEX}}
	{}{\typeout{[preamble-eduzy-v2] Không gắn được bảng Đúng/Sai (tô hàng)}}
\xpatchcmd{\writeChoiceTFt}{{\end{longtable}}}{{\end{longtable}\EZketthucTF}}
	{}{\typeout{[preamble-eduzy-v2] Không gắn được bảng Đúng/Sai (cuối longtable)}}
\xpatchcmd{\writeChoiceTFt}{{\end{xtabular}\end{center}\medskip\noindent}}{{\end{xtabular}\end{center}\EZketthucTF\medskip\noindent}}
	{}{\typeout{[preamble-eduzy-v2] Không gắn được bảng Đúng/Sai (cuối xtabular)}}
\fi
%=====================================================================
\renewcommand{\parallel}{\mathrel{/\!\!/}}

%      \TLdapso                      % tài liệu: \TLfull | \TLdapso | \TLan
%      \begin{document} ... \end{document}
%  Logo Eduzy vẽ bằng TikZ ngay trong file (\eduzylogo), không cần file ảnh.
%  Biên dịch: pdfLaTeX (máy có gói vietnam) hoặc XeLaTeX.
%============================================================
\documentclass[12pt,a4paper,twoside]{article}%<-- twoside: bat dau trang/chan trang chan-le
\usepackage{iftex}
\usepackage{amsmath,amssymb,tikz,mathrsfs,tkz-tab}
\usepackage{graphicx}
\usepackage{pgfplots}% đồ thị hàm số
\pgfplotsset{compat=1.17}
%\usepackage{mathptmx}
\usepackage{mathpazo}
\ifXeTeX% XeLaTeX: font chữ có đủ dấu tiếng Việt, giữ nguyên font toán mathpazo
	\usepackage[no-math]{fontspec}
	\setmainfont{TeX Gyre Pagella}
	\newfontfamily\EZqag{TeX Gyre Adventor}[NFSSFamily=qag,Ligatures=TeX]
\fi
\usepackage[many]{tcolorbox}
\usetikzlibrary{math,through,calc,intersections,angles,quotes,shapes,shapes.geometric,arrows,patterns,matrix,chains,arrows.meta,decorations.shapes,decorations.fractals,decorations.markings,shadows,decorations.pathreplacing}
\usepackage[top=1.6cm, bottom=1.8cm, left=1.8cm, right=1.8cm]{geometry}% v2: top 1.6cm (bản cũ 1.5cm) để chữ không sát dải đầu trang
\providecommand{\CheDo}{loigiai}% loigiai | dethi | solcolor | color | book
\expandafter\usepackage\expandafter[\CheDo]{ex_test}
% \tenchuong: khai báo ở file chính (trước \input) hoặc \renewcommand sau \input
%============================================================
%  [GỘP TỪ eduzy-de.tex] EDUZY - MÀU NÂU ĐỒNG / VÀNG KIM
%  + ĐẦU TRANG, CHÂN TRANG KIỂU DẢI NGANG (v2), ĐỐI XỨNG TRANG CHẴN / TRANG LẺ.
%  (không cần file eduzy-de.tex riêng nữa)
%  Khối này phải đặt SAU \usepackage[...]{ex_test} và SAU \def\tenchuong{...}
%============================================================
\usepackage{lastpage}
\usepackage{needspace}
%--- DÀN TRANG ĐỀ THI: tách câu (THÊM 05/11/2026; chỉ thêm lệnh mới, KHÔNG đổi hành vi mặc định) ---
%  Hai công tắc, mặc định tắt (= preamble chạy y như cũ). Bật bằng cách đặt TRƯỚC %============================================================
%  PREAMBLE EDUZY v2 - dùng chung cho đề thi, bài học, chuyên đề (gói ex_test)
%  Khác bản preamble-eduzy.tex ở đầu trang, chân trang: dải nền nâu đồng tràn ngang,
%  logo + Eduzy.com.vn, ô chữ vàng kem cạnh xiên, chân trang "Trang (số)";
%  logo Eduzy chìm giữa mọi trang (tắt: \logochimtat, bật lại: \logochimbat).
%  Cách dùng ở file chính (xem mau-chuan.tex, mau-chuyen-de.tex):
%      \def\CheDo{loigiai}          % loigiai | dethi | solcolor | color | book
%      \def\tenchuong{<dòng chữ trên đầu trang>}
%      \def\thuonghieu{<chữ ở chân trang lẻ>}   % không khai báo thì dùng mặc định
%      \def\tacgia{<chữ ở chân trang chẵn>}      % không khai báo thì dùng mặc định
%      \def\khauhieu{<chữ cạnh logo ở đầu trang>}  % tùy chọn, mặc định Eduzy.com.vn
%      \input{preamble-eduzy-v2}
%      \TLdapso                      % tài liệu: \TLfull | \TLdapso | \TLan
%      \begin{document} ... \end{document}
%  Logo Eduzy vẽ bằng TikZ ngay trong file (\eduzylogo), không cần file ảnh.
%  Biên dịch: pdfLaTeX (máy có gói vietnam) hoặc XeLaTeX.
%============================================================
\documentclass[12pt,a4paper,twoside]{article}%<-- twoside: bat dau trang/chan trang chan-le
\usepackage{iftex}
\usepackage{amsmath,amssymb,tikz,mathrsfs,tkz-tab}
\usepackage{graphicx}
\usepackage{pgfplots}% đồ thị hàm số
\pgfplotsset{compat=1.17}
%\usepackage{mathptmx}
\usepackage{mathpazo}
\ifXeTeX% XeLaTeX: font chữ có đủ dấu tiếng Việt, giữ nguyên font toán mathpazo
	\usepackage[no-math]{fontspec}
	\setmainfont{TeX Gyre Pagella}
	\newfontfamily\EZqag{TeX Gyre Adventor}[NFSSFamily=qag,Ligatures=TeX]
\fi
\usepackage[many]{tcolorbox}
\usetikzlibrary{math,through,calc,intersections,angles,quotes,shapes,shapes.geometric,arrows,patterns,matrix,chains,arrows.meta,decorations.shapes,decorations.fractals,decorations.markings,shadows,decorations.pathreplacing}
\usepackage[top=1.6cm, bottom=1.8cm, left=1.8cm, right=1.8cm]{geometry}% v2: top 1.6cm (bản cũ 1.5cm) để chữ không sát dải đầu trang
\providecommand{\CheDo}{loigiai}% loigiai | dethi | solcolor | color | book
\expandafter\usepackage\expandafter[\CheDo]{ex_test}
% \tenchuong: khai báo ở file chính (trước \input) hoặc \renewcommand sau \input
%============================================================
%  [GỘP TỪ eduzy-de.tex] EDUZY - MÀU NÂU ĐỒNG / VÀNG KIM
%  + ĐẦU TRANG, CHÂN TRANG KIỂU DẢI NGANG (v2), ĐỐI XỨNG TRANG CHẴN / TRANG LẺ.
%  (không cần file eduzy-de.tex riêng nữa)
%  Khối này phải đặt SAU \usepackage[...]{ex_test} và SAU \def\tenchuong{...}
%============================================================
\usepackage{lastpage}
\usepackage{needspace}
%--- DÀN TRANG ĐỀ THI: tách câu (THÊM 05/11/2026; chỉ thêm lệnh mới, KHÔNG đổi hành vi mặc định) ---
%  Hai công tắc, mặc định tắt (= preamble chạy y như cũ). Bật bằng cách đặt TRƯỚC \input{preamble-eduzy-v2}
%  trong file đề (tools/dan_trang/dan_trang.py tự ghi khi đặt \NeedspaceCau):
%    \def\TachCauDS{1}   câu Đúng/Sai được tách: câu dẫn (kèm hình) ở cuối trang, bảng Phát biểu sang trang sau
%                        (câu dẫn và hình không tách nhau, bảng không tách giữa hai trang)
%    \def\TachCauLG{1}   bản lời giải (\CheDo khác dethi): phần đề của câu cao hơn 1/3 trang được phép tách
%  Lệnh dùng trong file đề (giá trị do tools/dan_trang/dan_trang.py đo ở bản dethi):
%    \NeedspaceCau[<cao câu dẫn + hình>]{<cao cả khối câu>}   ngay trên \begin{ex} (câu đầu phần: trên tiêu đề phần)
%        công tắc tắt: đúng bằng \Needspace{<cao cả khối câu>} như cũ;
%        \TachCauDS=1 và có [..] (câu Đúng/Sai được chọn tách): chỉ giữ câu dẫn + hình;
%        \TachCauLG=1 ở bản lời giải: đòi tối đa 1/3 \textheight.
%    \NeedspaceBang{<cao bảng Phát biểu>}%   ngay trước \choiceTFt: \TachCauDS=1 thì giữ trọn bảng trên một trang;
%        công tắc tắt: không làm gì.
\providecommand{\TachCauDS}{0}
\providecommand{\TachCauLG}{0}
\newlength{\EZnsCau}
\newcommand{\NeedspaceCau}[2][]{%
	\setlength{\EZnsCau}{#2}%
	\ifnum\TachCauDS=1 \if\relax\detokenize{#1}\relax\else\setlength{\EZnsCau}{#1}\fi\fi
	\ifnum\TachCauLG=1 \ifdefstring{\kindOPTN}{dethi}{}{\ifdim\EZnsCau>0.3333\textheight\setlength{\EZnsCau}{0.3333\textheight}\fi}\fi
	\Needspace{\EZnsCau}}
\newcommand{\NeedspaceBang}[1]{\ifnum\TachCauDS=1 \Needspace{#1}\fi}

%--- Thang màu: đậm -> nhạt
\definecolor{gold1}{RGB}{126,96,0}    % nâu đồng đậm  #7E6000
\definecolor{gold2}{RGB}{192,144,3}   % vàng kim      #C09003
\definecolor{gold3}{RGB}{255,215,99}  % vàng nhạt     #FFD763
\definecolor{gold4}{RGB}{255,231,153} % vàng rất nhạt #FFE799
\definecolor{gold5}{RGB}{255,243,205} % vàng kem      #FFF3CD
%--- Ghi đè tông xanh/đỏ cũ
\definecolor{myblue}{RGB}{126,96,0}
\colorlet{blue}{gold1}
\colorlet{cyan}{gold2}
\colorlet{teal}{gold2}

%--- Tên chương mặc định (file đề nên tự khai báo lại)
\providecommand{\tenchuong}{}
%--- Chân trang
\providecommand{\thuonghieu}{TOÁN 12 -- Eduzy.com.vn}
\providecommand{\tacgia}{Tô Gia Bảo -- 0353703783}

%--- Chữ cạnh logo ở đầu trang
\providecommand{\khauhieu}{Eduzy.com.vn}

%====================== THÔNG SỐ ĐẦU TRANG, CHÂN TRANG (v2) ======================
%  Dải nền đậm tràn ngang trang; ô chữ vàng kem có cạnh xiên, kèm 2 vạch xiên chuyển màu.
%  Đầu trang lẻ : logo + \khauhieu (trái)   | ô chữ \tenchuong (phải)
%  Chân trang lẻ: ô chữ \thuonghieu (trái)  | Trang (số) (phải)
%  Trang chẵn soi gương trang lẻ; chân trang chẵn ghi \tacgia.
%--- Màu: chỉ dùng thang vàng kim Eduzy
\colorlet{ezNen}{gold1}% dải nền đậm
\colorlet{ezVach1}{gold3}% vạch xiên sát dải nền (đậm hơn)
\colorlet{ezVach2}{gold4}% vạch xiên sát ô chữ (nhạt hơn)
\colorlet{ezO}{gold5}% ô chữ vàng kem
%--- Kích thước
\newlength{\EZdai}\setlength{\EZdai}{1.1cm}% chiều cao dải nền
\newlength{\EZo}\setlength{\EZo}{0.68cm}% chiều cao ô chữ
\newlength{\EZxien}\setlength{\EZxien}{0.55cm}% độ lệch ngang của cạnh xiên
\newlength{\EZvach}\setlength{\EZvach}{0.75cm}% bề rộng mỗi vạch xiên
\newlength{\EZkhe}\setlength{\EZkhe}{0.12cm}% khe trắng giữa vạch và ô chữ
\newlength{\EZdem}\setlength{\EZdem}{0.4cm}% khoảng đệm từ chữ tới cạnh xiên
\newlength{\EZle}\setlength{\EZle}{1.8cm}% lề trái/phải của chữ (bằng lề thân bài)
\newlength{\EZrongmin}\setlength{\EZrongmin}{4.5cm}% ô chữ rộng tối thiểu
\newlength{\EZrong}% bề rộng chữ trong ô (tính tự động)
%--- Font
\newcommand{\EZfontO}{\fontsize{10.5pt}{12pt}\selectfont\bfseries}% chữ trong ô
\newcommand{\EZfontKH}{\fontfamily{qag}\fontsize{11pt}{12pt}\selectfont\bfseries}% \khauhieu cạnh logo
\newcommand{\EZfontTrang}{\fontsize{11.5pt}{12pt}\selectfont\bfseries}% chữ "Trang" và số trang
%--- Đo bề rộng chữ trong ô
\newcommand{\EZdorong}[1]{%
	\settowidth{\EZrong}{\EZfontO #1}%
	\ifdim\EZrong<\EZrongmin\setlength{\EZrong}{\EZrongmin}\fi}
%--- Logo + chữ trắng \khauhieu (dựng sẵn một lần khi vào \begin{document})
\newsavebox{\EZlogobox}
\AtBeginDocument{\sbox{\EZlogobox}{%
		\raisebox{-0.5\height}{\eduzylogo[1cm]}\hspace{0.15cm}%
		\raisebox{-0.5\height}{\color{white}\EZfontKH\khauhieu}}}
%--- Hai vạch xiên + khe trắng. #1: điểm gốc ở mép dưới, #2: hướng (1 sang phải, -1 sang trái),
%    #3: hướng nghiêng (1 kiểu "/", -1 kiểu "\")
\newcommand{\EZvachxien}[3]{%
	\fill[white] (#1) --++({(#3)*\EZxien},\EZo) --++({(#2)*\EZkhe},0) --++({-(#3)*\EZxien},-\EZo) --cycle;
	\fill[ezVach2] ([xshift=(#2)*\EZkhe]#1) --++({(#3)*\EZxien},\EZo) --++({(#2)*\EZvach},0) --++({-(#3)*\EZxien},-\EZo) --cycle;
	\fill[ezVach1] ([xshift=(#2)*(\EZkhe+\EZvach)]#1) --++({(#3)*\EZxien},\EZo) --++({(#2)*\EZvach},0) --++({-(#3)*\EZxien},-\EZo) --cycle;
	\draw[white,line width=0.6pt] ([xshift=(#2)*(\EZkhe+\EZvach)]#1) --++({(#3)*\EZxien},\EZo);
}
%--- Chữ trong ô (căn giữa theo chiều cao ô, không lệch vì dấu tiếng Việt)
\tikzset{ezchu/.style={inner sep=0pt,text=ezNen,text height=0.72em,text depth=0.22em}}

%====================== LOGO CHÌM GIỮA TRANG ======================
%  Vẽ cùng đầu trang nên nằm dưới chữ của thân bài, có trên mọi trang.
%  Tắt: \logochimtat (từ trang đó trở đi), bật lại: \logochimbat.
%  Đổi cỡ, độ đậm: khai báo trước \input, ví dụ \def\EZchimcao{9cm} \def\EZchimdo{0.1}
\providecommand{\EZchimcao}{10cm}% chiều cao logo chìm
\providecommand{\EZchimdo}{0.12}% độ đậm (0 = trong suốt, 1 = rõ hẳn)
\newif\ifEZchim \EZchimtrue
\newcommand{\logochimtat}{\EZchimfalse}
\newcommand{\logochimbat}{\EZchimtrue}
%  Dựng sẵn logo mờ một lần khi vào \begin{document}; transparency group làm mờ cả logo
%  một lượt nên không lộ vệt ở chỗ các dải màu chồng nhau.
\newsavebox{\EZchimtho}
\newsavebox{\EZchimbox}
\AtBeginDocument{%
	\sbox{\EZchimtho}{\eduzylogo[\EZchimcao]}%
	\sbox{\EZchimbox}{\begin{tikzpicture}
			\begin{scope}[transparency group,opacity=\EZchimdo]
				\node[inner sep=0pt] {\usebox{\EZchimtho}};
			\end{scope}
	\end{tikzpicture}}}
\newcommand{\EZvechim}{%
	\ifEZchim
		\node[inner sep=0pt] at (current page.center) {\usebox{\EZchimbox}};
	\fi}

%====================== ĐẦU TRANG - TRANG LẺ ======================
%  (logo bên trái, ô chữ \tenchuong bên phải)
\newcommand{\myfancyheadOdd}{%
	\EZdorong{\tenchuong}%
	\begin{tikzpicture}[remember picture,overlay]
		\coordinate (T) at (current page.north west);
		\coordinate (P) at (current page.north east);
		%--- logo chìm giữa trang
		\EZvechim
		%--- dải nền
		\fill[ezNen] (T) rectangle ([yshift=-\EZdai]P);
		%--- ô chữ: X là góc dưới bên trái, cạnh trái xiên kiểu "/"
		\coordinate (X) at ([xshift=-\EZle-\EZrong-\EZdem-\EZxien,yshift=-\EZdai]P);
		\fill[ezO] (X) --++(\EZxien,\EZo) -- ([yshift=-\EZdai+\EZo]P) -- ([yshift=-\EZdai]P) --cycle;
		\EZvachxien{X}{-1}{1}
		\node[ezchu,anchor=east] at ([xshift=-\EZle,yshift=-\EZdai+0.5*\EZo]P) {\EZfontO\tenchuong};
		%--- logo + tên
		\node[inner sep=0pt,anchor=west] at ([xshift=\EZle,yshift=-0.5*\EZdai]T) {\usebox{\EZlogobox}};
	\end{tikzpicture}%
}

%====================== ĐẦU TRANG - TRANG CHẴN =====================
%  (soi gương: ô chữ \tenchuong bên trái, logo bên phải)
\newcommand{\myfancyheadEven}{%
	\EZdorong{\tenchuong}%
	\begin{tikzpicture}[remember picture,overlay]
		\coordinate (T) at (current page.north west);
		\coordinate (P) at (current page.north east);
		%--- logo chìm giữa trang
		\EZvechim
		%--- dải nền
		\fill[ezNen] (T) rectangle ([yshift=-\EZdai]P);
		%--- ô chữ: X là góc dưới bên phải, cạnh phải xiên kiểu "\"
		\coordinate (X) at ([xshift=\EZle+\EZrong+\EZdem+\EZxien,yshift=-\EZdai]T);
		\fill[ezO] (X) --++(-\EZxien,\EZo) -- ([yshift=-\EZdai+\EZo]T) -- ([yshift=-\EZdai]T) --cycle;
		\EZvachxien{X}{1}{-1}
		\node[ezchu,anchor=west] at ([xshift=\EZle,yshift=-\EZdai+0.5*\EZo]T) {\EZfontO\tenchuong};
		%--- logo + tên
		\node[inner sep=0pt,anchor=east] at ([xshift=-\EZle,yshift=-0.5*\EZdai]P) {\usebox{\EZlogobox}};
	\end{tikzpicture}%
}

%====================== CHÂN TRANG - TRANG LẺ ======================
%  (ô chữ \thuonghieu bên trái, "Trang (số)" bên phải)
\newcommand{\myfancyfootOdd}{%
	\EZdorong{\thuonghieu}%
	\begin{tikzpicture}[remember picture,overlay]
		\coordinate (T) at (current page.south west);
		\coordinate (P) at (current page.south east);
		%--- dải nền
		\fill[ezNen] (T) rectangle ([yshift=\EZdai]P);
		%--- ô chữ sát mép trên dải: X là góc dưới bên phải, cạnh phải xiên kiểu "/"
		\coordinate (X) at ([xshift=\EZle+\EZrong+\EZdem,yshift=\EZdai-\EZo]T);
		\fill[ezO] ([yshift=\EZdai-\EZo]T) -- (X) --++(\EZxien,\EZo) -- ([yshift=\EZdai]T) --cycle;
		\EZvachxien{X}{1}{1}
		\node[ezchu,anchor=west] at ([xshift=\EZle,yshift=\EZdai-0.5*\EZo]T) {\EZfontO\thuonghieu};
		%--- số trang
		\node[circle,draw=white,line width=0.9pt,text=white,inner sep=0pt,minimum size=0.62cm,anchor=east]
		(so) at ([xshift=-\EZle,yshift=0.5*\EZdai]P) {\EZfontTrang\thepage};
		\node[inner sep=0pt,text=white,anchor=base east] (tr) at ([xshift=-0.15cm]so.base west) {\EZfontTrang Trang};
	\end{tikzpicture}%
}

%====================== CHÂN TRANG - TRANG CHẴN ====================
%  (soi gương: "Trang (số)" bên trái, ô chữ \tacgia bên phải)
\newcommand{\myfancyfootEven}{%
	\EZdorong{\tacgia}%
	\begin{tikzpicture}[remember picture,overlay]
		\coordinate (T) at (current page.south west);
		\coordinate (P) at (current page.south east);
		%--- dải nền
		\fill[ezNen] (T) rectangle ([yshift=\EZdai]P);
		%--- ô chữ sát mép trên dải: X là góc dưới bên trái, cạnh trái xiên kiểu "\"
		\coordinate (X) at ([xshift=-\EZle-\EZrong-\EZdem,yshift=\EZdai-\EZo]P);
		\fill[ezO] ([yshift=\EZdai-\EZo]P) -- (X) --++(-\EZxien,\EZo) -- ([yshift=\EZdai]P) --cycle;
		\EZvachxien{X}{-1}{-1}
		\node[ezchu,anchor=east] at ([xshift=-\EZle,yshift=\EZdai-0.5*\EZo]P) {\EZfontO\tacgia};
		%--- số trang
		%    (vòng tròn ẩn so0 nằm giữa dải để lấy đường chân chữ, giống trang lẻ)
		\node[circle,inner sep=0pt,minimum size=0.62cm] (so0) at ([yshift=0.5*\EZdai]T) {\EZfontTrang\phantom{\thepage}};
		\node[inner sep=0pt,text=white,anchor=base west] (tr) at ($(so0.base -| T)+(\EZle,0)$) {\EZfontTrang Trang};
		\node[circle,draw=white,line width=0.9pt,text=white,inner sep=0pt,minimum size=0.62cm,anchor=base west]
		(so) at ([xshift=0.15cm]tr.base east) {\EZfontTrang\thepage};
	\end{tikzpicture}%
}

%--- Giữ tương thích với các file cũ đang gọi \myfancyhead, \myfancyfoot
\newcommand{\myfancyhead}{\myfancyheadOdd}
\newcommand{\myfancyfoot}{\myfancyfootOdd}

%======================= GẮN VÀO TRANG =============================
\pagestyle{fancy}
\fancyhf{}
\fancyhead[LO]{\myfancyheadOdd}
\fancyhead[RE]{\myfancyheadEven}
\fancyfoot[LO]{\myfancyfootOdd}
\fancyfoot[RE]{\myfancyfootEven}
\renewcommand{\headrulewidth}{0pt}
\renewcommand{\footrulewidth}{0pt}
\fancypagestyle{plain}{%
	\fancyhf{}
	\fancyhead[LO]{\myfancyheadOdd}
	\fancyhead[RE]{\myfancyheadEven}
	\fancyfoot[LO]{\myfancyfootOdd}
	\fancyfoot[RE]{\myfancyfootEven}
	\renewcommand{\headrulewidth}{0pt}
	\renewcommand{\footrulewidth}{0pt}
}
%=================== HẾT PHẦN GỘP TỪ eduzy-de.tex ===================
\newcommand{\indam}[1]{
	{\fontsize{12pt}{9pt}\sffamily \selectfont\bfseries\color{gold1} #1}
}
\newcommand{\ind}[1]{
	{\fontsize{14pt}{9pt}\fontfamily{qag} \selectfont\bfseries\color{gold1} #1}
}
\renewcommand{\FalseEX}{\stepcounter{dapan}\textbf{\selectfont\color{gold1}\Alph{dapan}.}\ignorespaces}
\renewtheorem{ex}{\color{gold1}Câu}
\renewcommand{\baselinestretch}{1.1}
\newcommand{\hoac}[1]{
	\left[\begin{aligned}#1\end{aligned}\right.}
\newcommand{\heva}[1]{
	\left\{\begin{aligned}#1\end{aligned}\right.}
\usepackage{esvect}
\def\vec{\vv}
\def\vt{\vv}
\def\overrightarrow{\vv}
%====================== LOGO EDUZY VẼ BẰNG TIKZ ======================
%  Đường viền dò từ logo gốc; dải màu vàng kim đo trực tiếp trên ảnh gốc
%  (nghiêng 21 độ, tô bằng 80 dải hẹp nên hiển thị giống nhau trên mọi trình xem PDF).
%  Dùng: \eduzylogo          -> cao 2.4cm (mặc định)
%        \eduzylogo[1.5cm]   -> cao 1.5cm
\definecolor{ezText}{RGB}{245,200,46}
\newcommand{\eduzylogo}[1][2.4cm]{%
	\resizebox{!}{#1}{%
		\begin{tikzpicture}[x=0.01cm,y=0.01cm]
			\path[use as bounding box] (0,0) rectangle (224,-224);
			% mũ tốt nghiệp + vương miện
			\begin{scope}
				\clip
				(119.4,-176.1)..controls(118.9,-176.1)and(116.8,-176)..(114.6,-176)
				..controls(101.4,-176)and(85.8,-175.9)..(83.7,-175.7)
				..controls(81.8,-175.6)and(81,-175.6)..(80.2,-175.7)
				..controls(78.8,-175.9)and(77.7,-175.9)..(76.2,-175.7)
				..controls(74.9,-175.5)and(70.8,-175.5)..(64.3,-175.7)
				..controls(59.1,-175.9)and(58.3,-175.7)..(58.6,-173.7)
				..controls(58.6,-172.8)and(59.6,-169.6)..(60.3,-168)
				..controls(60.8,-166.9)and(61.4,-165.8)..(62.3,-165)
				..controls(63.1,-164.2)and(64.3,-163.4)..(65,-163.2)
				..controls(65.5,-163.1)and(72.3,-163)..(76.6,-163.1)
				..controls(77.9,-163.1)and(80.1,-163.1)..(81.4,-163)
				..controls(82.8,-162.9)and(86.3,-162.9)..(89.2,-162.9)
				..controls(93,-162.9)and(94.9,-162.8)..(96,-162.7)
				..controls(97.2,-162.5)and(101,-162.5)..(114,-162.4)
				..controls(123,-162.3)and(134.9,-162.2)..(140.2,-162.1)
				..controls(150.1,-162)and(152.2,-162.1)..(153.6,-162.4)
				..controls(154.4,-162.6)and(157.1,-163.8)..(158.1,-164.4)
				..controls(158.8,-164.8)and(159.7,-165.7)..(159.9,-166.3)
				..controls(160,-166.5)and(160.2,-167.1)..(160.2,-167.6)
				..controls(160.3,-168)and(160.6,-169.4)..(160.9,-170.6)
				..controls(161.2,-171.7)and(161.5,-173.1)..(161.6,-173.6)
				..controls(161.8,-174.4)and(161.8,-174.5)..(161.5,-174.9)
				..controls(161.2,-175.6)and(160.8,-175.9)..(159.9,-176)
				..controls(159.2,-176.1)and(154.1,-176.2)..(151.5,-176.2)
				..controls(148.1,-176.1)and(122.6,-176.1)..(121.6,-176.2)
				..controls(120.8,-176.2)and(119.8,-176.2)..(119.4,-176.1)--cycle
				(130.4,-160)..controls(129,-160)and(119.8,-160)..(110.1,-159.9)
				..controls(100.3,-159.8)and(88.5,-159.7)..(83.9,-159.6)
				..controls(77.8,-159.5)and(75.3,-159.5)..(74.5,-159.6)
				..controls(73.9,-159.7)and(72.5,-159.8)..(71.3,-159.8)
				..controls(69.6,-159.8)and(69,-159.8)..(68.5,-159.6)
				..controls(67.6,-159.3)and(66.8,-158.6)..(66.2,-157.6)
				..controls(66,-157.3)and(65.6,-156.5)..(65.3,-156.1)
				..controls(64.6,-155.1)and(64.5,-154.9)..(63.7,-152.9)
				..controls(63.2,-151.6)and(62.9,-150.9)..(61.8,-149.6)
				..controls(61.5,-149.2)and(61.1,-148.5)..(60.9,-147.9)
				..controls(60.7,-147.4)and(60.3,-146.6)..(60.1,-146.1)
				..controls(59.9,-145.5)and(59.5,-144.9)..(59,-144.3)
				..controls(58.5,-143.7)and(58.2,-143.2)..(57.9,-142.5)
				..controls(57.7,-141.9)and(57.3,-141.1)..(57,-140.7)
				..controls(56.6,-140.3)and(56.1,-139.3)..(55.8,-138.5)
				..controls(55.5,-137.8)and(55.1,-136.8)..(54.9,-136.5)
				..controls(54.7,-136.1)and(54.3,-135.2)..(53.9,-134.5)
				..controls(53.6,-133.8)and(53,-132.7)..(52.6,-132)
				..controls(51.6,-130.5)and(51.4,-130.2)..(51.1,-129)
				..controls(50.6,-127.5)and(50.6,-127.4)..(49.7,-126.3)
				..controls(49.2,-125.7)and(48.6,-124.8)..(48.3,-124.2)
				..controls(48,-123.6)and(47.5,-122.7)..(47.3,-122.2)
				..controls(47,-121.8)and(46.5,-120.7)..(46.1,-119.8)
				..controls(45.8,-118.9)and(45.2,-117.8)..(44.9,-117.2)
				..controls(44.2,-116.1)and(43.3,-114.4)..(42.3,-112.2)
				..controls(41.8,-111.3)and(41.3,-110.4)..(41.1,-110.1)
				..controls(40.5,-109.4)and(40.3,-109)..(40,-107.8)
				..controls(39.8,-107)and(39.6,-106.5)..(39.2,-106)
				..controls(38.9,-105.6)and(38.4,-104.7)..(38,-104)
				..controls(37.7,-103.4)and(37.2,-102.5)..(37,-102.1)
				..controls(36.7,-101.7)and(36.3,-100.8)..(36,-100.1)
				..controls(35.6,-98.9)and(35.3,-98.4)..(34.6,-97.6)
				..controls(34.2,-97.1)and(34.2,-96.3)..(34.6,-96)
				..controls(35,-95.7)and(35.5,-95.9)..(36.2,-96.7)
				..controls(36.5,-97)and(37.2,-97.5)..(37.6,-97.7)..controls(38,-98)and(38.6,-98.4)..(39.1,-98.9)
				..controls(39.6,-99.4)and(40.3,-100)..(40.6,-100.3)
				..controls(41.3,-100.8)and(42.4,-101.8)..(43.7,-102.9)
				..controls(44,-103.2)and(44.6,-103.7)..(45,-103.9)
				..controls(45.4,-104.2)and(45.9,-104.6)..(46.1,-104.8)
				..controls(47.1,-105.9)and(48.8,-107.1)..(50.5,-107.9)
				..controls(51.8,-108.6)and(52.3,-109)..(53.4,-110)
				..controls(54.8,-111.2)and(55.6,-111.6)..(57,-112.1)
				..controls(57.5,-112.3)and(58.5,-112.8)..(59.4,-113.3)
				..controls(60.2,-113.7)and(61.3,-114.3)..(61.9,-114.5)
				..controls(62.4,-114.7)and(63.3,-115.1)..(63.8,-115.3)
				..controls(64.4,-115.6)and(65.4,-115.9)..(66.2,-116.1)
				..controls(67,-116.3)and(68.4,-116.6)..(69.2,-116.9)
				..controls(72.3,-117.9)and(72.8,-118)..(78,-118.1)
				..controls(84,-118.3)and(86.1,-118.1)..(88.1,-117.1)
				..controls(89.4,-116.5)and(93,-115)..(93.7,-114.9)
				..controls(94.8,-114.6)and(95.4,-114.2)..(95.9,-113.4)
				..controls(96.2,-113)and(96.7,-112.5)..(97.1,-112.2)
				..controls(98,-111.6)and(101.3,-108.4)..(102.1,-107.2)
				..controls(102.4,-106.8)and(102.9,-106.2)..(103.2,-105.8)
				..controls(103.5,-105.5)and(104,-104.8)..(104.2,-104.2)
				..controls(104.5,-103.6)and(105,-102.8)..(105.3,-102.4)
				..controls(105.6,-102)and(106.1,-101.2)..(106.5,-100.6)
				..controls(107.1,-99.5)and(108.6,-97.5)..(109.5,-96.5)
				..controls(109.9,-96.1)and(110.1,-96)..(110.4,-96)
				..controls(111.1,-96)and(111.6,-96.5)..(111.8,-97.6)
				..controls(112,-98.6)and(112.3,-99.3)..(112.8,-99.9)
				..controls(113.7,-100.9)and(115.2,-103.1)..(115.7,-104.2)
				..controls(116.1,-105)and(116.5,-105.5)..(116.9,-106)
				..controls(117.3,-106.4)and(117.8,-106.9)..(118,-107.3)
				..controls(118.2,-107.6)and(118.6,-108)..(119,-108.3)
				..controls(119.6,-108.8)and(121.1,-110.3)..(121.8,-111.2)
				..controls(122.1,-111.7)and(122.5,-112)..(123.4,-112.4)
				..controls(124.8,-113.2)and(125.8,-113.8)..(127,-114.6)
				..controls(127.5,-114.9)and(128.1,-115.3)..(128.2,-115.4)
				..controls(128.4,-115.5)and(129.3,-115.8)..(130.2,-116)
				..controls(131.1,-116.2)and(132.5,-116.6)..(133.3,-116.9)
				..controls(135.7,-117.9)and(136.1,-117.9)..(142.5,-117.8)
				..controls(145.8,-117.8)and(148.4,-117.7)..(148.8,-117.6)
				..controls(149.2,-117.6)and(150.2,-117.2)..(151,-116.8)
				..controls(152,-116.4)and(152.8,-116)..(153.6,-115.9)
				..controls(154.5,-115.7)and(155.2,-115.4)..(156.3,-114.9)
				..controls(157.1,-114.5)and(158.2,-114)..(158.7,-113.8)
				..controls(159.1,-113.7)and(159.8,-113.3)..(160.2,-113)
				..controls(160.6,-112.7)and(161.3,-112.4)..(161.7,-112.2)
				..controls(162.7,-111.9)and(163.2,-111.7)..(164.1,-111.1)
				..controls(164.5,-110.8)and(165.3,-110.4)..(165.9,-110)
				..controls(166.5,-109.7)and(167.3,-109.2)..(167.6,-108.9)
				..controls(167.9,-108.6)and(168.6,-108.2)..(169,-108)
				..controls(169.4,-107.8)and(170.2,-107.3)..(170.7,-106.9)
				..controls(171.2,-106.5)and(172.1,-105.8)..(172.7,-105.4)
				..controls(173.2,-105)and(174.1,-104.4)..(174.5,-104)
				..controls(174.9,-103.7)and(175.7,-103.1)..(176.2,-102.6)
				..controls(176.7,-102.2)and(177.4,-101.5)..(177.9,-101.2)
				..controls(178.3,-100.8)and(178.8,-100.4)..(179,-100.2)
				..controls(179.2,-100)and(179.6,-99.6)..(180,-99.3)
				..controls(180.4,-99.1)and(180.9,-98.6)..(181.2,-98.4)
				..controls(181.5,-98.1)and(182.1,-97.7)..(182.5,-97.5)
				..controls(183.2,-97.2)and(183.5,-96.9)..(184.3,-95.9)
				..controls(184.6,-95.6)and(184.6,-95.6)..(184.8,-95.7)
				..controls(185,-95.8)and(185.1,-96)..(185.1,-96.2)
				..controls(185.1,-96.5)and(185,-96.7)..(184.6,-97.2)
				..controls(184.1,-97.8)and(184,-97.9)..(183.9,-98.8)
				..controls(183.7,-100.3)and(183.3,-101.3)..(182.9,-101.8)
				..controls(182.6,-102.1)and(182.3,-102.6)..(182.2,-102.9)
				..controls(182,-103.3)and(181.6,-104)..(181.2,-104.6)
				..controls(180.9,-105.2)and(180.5,-106.1)..(180.4,-106.6)
				..controls(179.9,-108)and(179.7,-108.5)..(179,-109.4)
				..controls(178.6,-109.8)and(178.2,-110.5)..(177.9,-111)
				..controls(177.7,-111.5)and(177.3,-112.3)..(177,-112.7)
				..controls(176.6,-113.2)and(176.4,-113.7)..(176.3,-114.2)
				..controls(176.2,-114.7)and(175.9,-115.2)..(175.4,-116)
				..controls(175,-116.6)and(174.4,-117.7)..(174,-118.5)
				..controls(173.7,-119.2)and(173.3,-120.1)..(173,-120.5)
				..controls(172.8,-120.9)and(172.4,-121.8)..(172.1,-122.5)
				..controls(171.8,-123.2)and(171.3,-124.1)..(171,-124.6)
				..controls(170.8,-125.1)and(170.3,-126)..(169.9,-126.6)
				..controls(169.6,-127.3)and(169.2,-128.1)..(169,-128.3)
				..controls(168.8,-128.6)and(168.5,-129.3)..(168.3,-129.8)
				..controls(168.1,-130.3)and(167.8,-131.1)..(167.5,-131.6)
				..controls(166.9,-132.5)and(164.2,-137.9)..(164.2,-138.1)
				..controls(164.2,-138.1)and(163.9,-138.8)..(163.4,-139.7)
				..controls(162.9,-140.5)and(162.4,-141.5)..(162.3,-141.9)
				..controls(161.9,-143)and(161.6,-143.4)..(160.9,-144.1)
				..controls(160.4,-144.6)and(160.3,-144.8)..(160.1,-145.4)
				..controls(160,-145.8)and(159.6,-146.8)..(159.2,-147.8)
				..controls(158.7,-148.8)and(158.1,-150)..(157.9,-150.6)
				..controls(157.6,-151.1)and(157.2,-151.9)..(156.9,-152.4)
				..controls(156.5,-152.8)and(156.1,-153.7)..(155.8,-154.4)
				..controls(155.6,-155)and(155.2,-155.9)..(154.9,-156.3)
				..controls(154.7,-156.7)and(154.4,-157.4)..(154.2,-157.9)
				..controls(154.1,-158.4)and(153.8,-159)..(153.5,-159.3)
				..controls(152.8,-160.1)and(153,-160.1)..(142.4,-160.1)
				..controls(137.2,-160.1)and(131.8,-160.1)..(130.4,-160)--cycle
				(141.5,-103.7)..controls(140.8,-103)and(139.8,-102.1)..(139.2,-101.6)
				..controls(138.5,-101.2)and(137.9,-100.6)..(137.8,-100.4)
				..controls(137,-99.5)and(134,-96.5)..(133.2,-96)
				..controls(132.7,-95.7)and(131.8,-94.9)..(131.1,-94.3)
				..controls(129.4,-92.7)and(127.9,-91.6)..(126.6,-91.1)
				..controls(125.2,-90.5)and(124.1,-90)..(122.9,-89.1)
				..controls(122.1,-88.6)and(121.3,-88.2)..(119.9,-87.7)
				..controls(118.8,-87.3)and(117.5,-86.8)..(116.9,-86.5)
				..controls(114.3,-85.5)and(113.7,-85.4)..(109.7,-85.5)
				..controls(105.3,-85.7)and(103.9,-85.9)..(101.6,-87)
				..controls(100.7,-87.4)and(99.4,-87.8)..(97.9,-88)
				..controls(97,-88.2)and(96.6,-88.3)..(95.8,-88.9)
				..controls(95.4,-89.2)and(94.7,-89.6)..(94.2,-89.8)
				..controls(92.9,-90.3)and(92.3,-90.7)..(90.9,-92.1)
				..controls(90.2,-92.7)and(89.4,-93.3)..(89,-93.5)..controls(87,-94.5)and(85.2,-96)..(83.6,-97.9)
				..controls(83.3,-98.2)and(82.5,-98.9)..(81.8,-99.4)
				..controls(80.7,-100.1)and(80.5,-100.3)..(80,-101.1)
				..controls(78.7,-103.1)and(77.2,-103.9)..(76.4,-103)
				..controls(75.7,-102.2)and(75.7,-103.2)..(75.7,-83.6)
				..controls(75.6,-73.6)and(75.6,-64.8)..(75.5,-64)
				..controls(75.5,-63.2)and(75.5,-61.4)..(75.6,-59.8)--(75.7,-57.1)--(76,-56.8)
				..controls(76.8,-56.2)and(78,-56.3)..(79.8,-57.2)
				..controls(80.4,-57.5)and(81.3,-57.8)..(81.8,-57.9)
				..controls(82.3,-58.1)and(83,-58.4)..(83.6,-58.8)..controls(84.3,-59.2)and(85,-59.5)..(86,-59.8)
				..controls(86.8,-60.1)and(88,-60.6)..(88.7,-60.9)
				..controls(89.3,-61.2)and(90.6,-61.7)..(91.4,-62)
				..controls(92.2,-62.3)and(93.7,-62.9)..(94.6,-63.3)
				..controls(95.6,-63.7)and(96.8,-64.1)..(97.4,-64.3)
				..controls(98,-64.5)and(99,-64.9)..(99.7,-65.2)
				..controls(100.4,-65.6)and(101.5,-66)..(102.3,-66.2)
				..controls(104.2,-66.8)and(105,-67.1)..(106.2,-67.8)
				..controls(107.9,-68.7)and(110.2,-69.4)..(111.3,-69.2)
				..controls(112.5,-69)and(115,-68.2)..(115.6,-67.8)
				..controls(117.3,-66.7)and(117.6,-66.6)..(119.5,-66.2)
				..controls(121.9,-65.8)and(122.1,-65.7)..(123,-65.1)
				..controls(124.1,-64.4)and(124.4,-64.3)..(126,-63.9)
				..controls(127,-63.7)and(127.7,-63.4)..(128.4,-63.1)
				..controls(129,-62.8)and(129.7,-62.5)..(130.6,-62.4)
				..controls(132.1,-62.1)and(134,-61.5)..(135.2,-61)
				..controls(135.7,-60.8)and(136.7,-60.5)..(137.3,-60.3)
				..controls(138.4,-60)and(139.3,-59.7)..(141.1,-59)
				..controls(141.6,-58.7)and(142,-58.6)..(142.5,-58.6)
				..controls(143.2,-58.6)and(143.3,-58.7)..(143.6,-59)
				..controls(144,-59.5)and(144.1,-59.9)..(144,-60.8)
				..controls(144,-61.2)and(144,-71.1)..(144,-82.8)--(143.9,-104.1)--(143.7,-104.5)
				..controls(143.5,-104.7)and(143.3,-104.8)..(143,-104.9)
				..controls(142.6,-104.9)and(142.5,-104.8)..(141.5,-103.7)--cycle
				(27.8,-91.4)..controls(27.1,-90.9)and(26.3,-90)..(26.1,-89.5)
				..controls(26,-89.2)and(25.8,-88.6)..(25.7,-87.9)
				..controls(25.4,-86.2)and(25.6,-84.8)..(26.5,-83.5)
				..controls(27.4,-82.3)and(28.3,-82)..(29.4,-82.3)
				..controls(31.1,-82.8)and(32.1,-84.9)..(32,-87.4)
				..controls(31.9,-89.3)and(31.2,-90.6)..(29.9,-91.4)
				..controls(29.1,-91.9)and(28.6,-91.9)..(27.8,-91.4)--cycle
				(190.1,-91.3)..controls(189.7,-91.2)and(189.3,-90.9)..(189.1,-90.7)
				..controls(188.3,-89.8)and(187.6,-87.7)..(187.6,-86.2)
				..controls(187.6,-85.5)and(187.7,-85.2)..(188.2,-84.4)
				..controls(189.1,-82.6)and(190.1,-81.8)..(191.3,-82.1)
				..controls(192,-82.3)and(193.5,-83.7)..(193.8,-84.4)
				..controls(194.1,-85.1)and(194.2,-87.1)..(194,-88.3)
				..controls(193.8,-89.8)and(193.3,-90.5)..(192.3,-91.3)
				..controls(191.7,-91.7)and(190.9,-91.7)..(190.1,-91.3)--cycle
				(173.7,-83.4)..controls(172.2,-82.5)and(172,-82.5)..(170.6,-82.4)
				..controls(169.7,-82.4)and(168.8,-82.4)..(168.3,-82.5)
				..controls(167.2,-82.7)and(166.9,-82.7)..(166.6,-82.2)
				..controls(166.3,-81.8)and(166.3,-81)..(166.8,-80)
				..controls(167.2,-78.9)and(167.4,-77.4)..(167.7,-74.1)
				..controls(167.8,-72.8)and(167.9,-71.8)..(167.9,-71.7)
				..controls(168,-71.7)and(168,-70.2)..(168,-68.4)..controls(168,-65.4)and(168,-64.9)..(167.7,-63)
				..controls(167.3,-60.5)and(167.1,-59.9)..(166.4,-58.4)
				..controls(166.1,-57.7)and(165.8,-56.9)..(165.7,-56.6)
				..controls(165.7,-56.2)and(165.4,-55.6)..(165.2,-55.2)
				..controls(164.6,-53.8)and(161.9,-50.7)..(161,-50.3)
				..controls(159.8,-49.7)and(157.9,-50)..(155.2,-51.2)
				..controls(154.7,-51.4)and(153.4,-51.8)..(152.2,-52.1)
				..controls(150.8,-52.4)and(149.4,-52.9)..(148.6,-53.2)
				..controls(147.8,-53.6)and(146.7,-54)..(146,-54.2)
				..controls(145.4,-54.4)and(144.5,-54.7)..(144,-54.9)
				..controls(143.6,-55.2)and(142.7,-55.5)..(141.8,-55.7)
				..controls(141,-55.9)and(139.6,-56.3)..(138.7,-56.6)
				..controls(137.8,-57)and(136.3,-57.4)..(135.3,-57.6)
				..controls(133.8,-57.9)and(133.2,-58.1)..(132.1,-58.7)
				..controls(131.1,-59.2)and(130.3,-59.4)..(129.4,-59.6)
				..controls(128.7,-59.8)and(127.8,-60)..(127.4,-60.1)
				..controls(127,-60.3)and(126,-60.5)..(125.1,-60.7)
				..controls(123.2,-61.1)and(121.7,-61.5)..(120.2,-62.1)
				..controls(119.7,-62.3)and(118.1,-62.8)..(116.7,-63.2)
				..controls(114.6,-63.9)and(114,-64.1)..(113,-64.7)
				..controls(111.9,-65.4)and(111.8,-65.4)..(111,-65.5)
				..controls(109.9,-65.5)and(109.3,-65.4)..(107.7,-64.7)
				..controls(107,-64.4)and(106,-64)..(105.3,-63.9)
				..controls(103.5,-63.5)and(102.2,-63)..(101.2,-62.5)
				..controls(100.1,-61.8)and(99.3,-61.5)..(97.8,-61.1)
				..controls(97.2,-61)and(96.2,-60.7)..(95.7,-60.4)
				..controls(95.1,-60.1)and(94.2,-59.8)..(93.5,-59.6)
				..controls(92.9,-59.4)and(91.7,-59)..(91,-58.7)..controls(90.2,-58.4)and(88.9,-58)..(88,-57.7)
				..controls(87.2,-57.5)and(86.1,-57.1)..(85.5,-56.8)
				..controls(85,-56.6)and(84.1,-56.2)..(83.5,-56)..controls(82.9,-55.7)and(81.8,-55.2)..(81,-54.9)
				..controls(80.1,-54.5)and(79.1,-54.2)..(78.7,-54.1)
				..controls(77.9,-53.8)and(77.3,-53.6)..(76.2,-52.9)
				..controls(75.7,-52.6)and(75.2,-52.4)..(74.1,-52.2)
				..controls(72.5,-51.9)and(72.4,-51.8)..(71.3,-51.1)
				..controls(69.9,-50.1)and(69.4,-49.9)..(67.9,-49.8)
				..controls(67,-49.7)and(66.6,-49.6)..(66.3,-49.4)
				..controls(66.1,-49.3)and(65.7,-49)..(65.4,-48.8)
				..controls(65.1,-48.7)and(64.2,-48.4)..(63.6,-48.2)
				..controls(62.9,-48)and(61.7,-47.6)..(60.9,-47.2)
				..controls(60.1,-46.8)and(58.9,-46.4)..(58.2,-46.2)
				..controls(57.6,-46)and(56.8,-45.7)..(56.5,-45.5)
				..controls(54.9,-44.6)and(54.6,-44.5)..(53.2,-44.2)
				..controls(51.7,-43.8)and(51,-43.6)..(49.9,-42.9)
				..controls(49.5,-42.6)and(48.9,-42.4)..(47.5,-42)
				..controls(45.7,-41.6)and(45.6,-41.5)..(44.3,-40.6)
				..controls(43.5,-40.1)and(42.8,-39.6)..(42.7,-39.4)
				..controls(42.4,-38.7)and(43,-38.1)..(44.5,-37.6)..controls(44.9,-37.5)and(46,-37)..(46.9,-36.6)
				..controls(47.9,-36.1)and(48.8,-35.8)..(49.6,-35.6)
				..controls(50.5,-35.4)and(51.3,-35.2)..(52.4,-34.6)
				..controls(53.3,-34.1)and(54.5,-33.7)..(54.9,-33.5)
				..controls(55.4,-33.4)and(56.4,-33)..(57.2,-32.6)
				..controls(58.6,-32)and(60.1,-31.5)..(61.6,-31.2)
				..controls(62.3,-31.1)and(62.7,-31)..(63.2,-30.7)
				..controls(64.1,-30)and(64.3,-29.9)..(65.3,-29.7)
				..controls(65.7,-29.6)and(66.5,-29.2)..(67.1,-28.9)
				..controls(67.6,-28.6)and(68.8,-28.2)..(69.6,-27.9)
				..controls(70.5,-27.7)and(71.6,-27.3)..(72.1,-27)
				..controls(72.6,-26.7)and(73.8,-26.2)..(74.8,-25.9)
				..controls(75.7,-25.5)and(76.7,-25.1)..(77.1,-25)
				..controls(77.4,-24.8)and(78.5,-24.4)..(79.6,-24.1)
				..controls(80.6,-23.8)and(82,-23.2)..(82.8,-22.9)
				..controls(83.5,-22.5)and(84.8,-22)..(85.8,-21.8)
				..controls(86.7,-21.5)and(87.8,-21.1)..(88.1,-20.9)
				..controls(88.4,-20.7)and(89.3,-20.3)..(90.1,-20.1)
				..controls(90.8,-19.8)and(91.9,-19.4)..(92.4,-19.1)
				..controls(93.5,-18.5)and(94.4,-18.2)..(95.4,-18.1)
				..controls(95.7,-18)and(96.4,-17.8)..(96.9,-17.5)
				..controls(97.5,-17.3)and(98.4,-16.9)..(99,-16.6)
				..controls(99.7,-16.3)and(100.9,-15.9)..(101.9,-15.6)
				..controls(102.9,-15.3)and(104.4,-14.7)..(105.3,-14.3)
				..controls(108.1,-13.2)and(108.5,-13.1)..(110.1,-13.1)
				..controls(111.4,-13.1)and(111.5,-13.1)..(112.4,-13.5)
				..controls(113.2,-14)and(115,-14.8)..(116.6,-15.4)
				..controls(116.9,-15.5)and(117.9,-15.8)..(118.8,-16.1)
				..controls(120.9,-16.6)and(121.5,-16.9)..(122.6,-17.9)
				..controls(123,-18.3)and(123.7,-18.8)..(123.9,-18.9)
				..controls(124.6,-19.4)and(126,-19.8)..(127.9,-20.2)
				..controls(129.6,-20.5)and(129.8,-20.6)..(131.1,-21.4)
				..controls(131.5,-21.7)and(132.4,-22.1)..(133.1,-22.3)
				..controls(136.1,-23.2)and(138.3,-24.1)..(140.2,-25)
				..controls(141,-25.5)and(141.5,-25.6)..(142.5,-25.8)
				..controls(145.1,-26.3)and(146,-26.6)..(149.4,-28.1)
				..controls(151.3,-28.9)and(153.5,-29.8)..(154.3,-30.1)
				..controls(155.1,-30.4)and(156.2,-30.8)..(156.6,-31)
				..controls(157.1,-31.2)and(158.2,-31.6)..(159.1,-31.8)
				..controls(160.6,-32.2)and(160.9,-32.3)..(161.7,-32.9)
				..controls(162.4,-33.3)and(162.9,-33.6)..(163.7,-33.8)
				..controls(164.3,-34)and(164.8,-34.1)..(164.9,-34.1)
				..controls(165,-34.1)and(165.6,-34.4)..(166.4,-34.8)
				..controls(167.3,-35.2)and(168.1,-35.4)..(169,-35.6)
				..controls(170,-35.8)and(170.5,-36)..(171.7,-36.6)
				..controls(172.4,-37)and(173.3,-37.4)..(173.6,-37.4)
				..controls(175.6,-37.7)and(176.2,-38)..(176.7,-39.1)
				..controls(176.9,-39.6)and(176.9,-39.7)..(176.7,-40.4)
				..controls(176.4,-41.6)and(175.8,-42.7)..(175.2,-43.4)
				..controls(174.5,-44.1)and(172.6,-45.5)..(171.7,-45.8)
				..controls(170.6,-46.2)and(169.1,-46)..(168.2,-45.4)
				..controls(167.9,-45.2)and(167.3,-44.8)..(166.9,-44.4)
				..controls(166.4,-44.1)and(165.8,-43.6)..(165.4,-43.2)
				..controls(165.1,-42.9)and(164.4,-42.3)..(163.9,-41.9)
				..controls(163.3,-41.5)and(162.7,-40.9)..(162.4,-40.5)
				..controls(161.4,-39.2)and(159.9,-38)..(158.8,-37.8)
				..controls(158.4,-37.8)and(158.3,-37.7)..(158.2,-37.2)
				..controls(158,-36.8)and(157.4,-36.1)..(157.1,-36.1)
				..controls(157,-36.1)and(156.7,-36.1)..(156.5,-36)
				..controls(156.1,-35.9)and(155.9,-35.7)..(155.7,-35.2)
				..controls(155.3,-34.6)and(155.1,-34.4)..(154,-34.1)
				..controls(153.6,-34)and(153.2,-33.7)..(152.9,-33.5)
				..controls(152.4,-32.9)and(151.6,-32.4)..(151.3,-32.4)
				..controls(150.9,-32.4)and(150.4,-32.6)..(150.4,-32.9)
				..controls(150.4,-33.2)and(150.8,-33.4)..(151.4,-33.6)
				..controls(152.1,-33.7)and(152.8,-34.1)..(153.5,-34.9)
				..controls(154.1,-35.5)and(154.3,-35.6)..(155,-35.8)
				..controls(155.7,-36)and(155.7,-36)..(156,-36.6)
				..controls(156.3,-37.3)and(156.7,-37.7)..(157.5,-37.8)
				..controls(158.1,-37.9)and(158.2,-38)..(158.4,-38.8)
				..controls(158.5,-39.4)and(158.6,-39.6)..(159.6,-40.4)
				..controls(160.2,-41)and(160.9,-41.6)..(161.3,-41.8)
				..controls(161.7,-42)and(161.9,-42.3)..(162.2,-42.8)
				..controls(162.4,-43.2)and(162.8,-43.8)..(163.1,-44.1)
				..controls(163.9,-44.8)and(165.1,-46.7)..(165.7,-48)
				..controls(166.2,-49.3)and(166.4,-49.7)..(167.2,-50.5)
				..controls(167.5,-50.8)and(167.9,-51.4)..(168.1,-51.8)
				..controls(168.3,-52.2)and(168.8,-53)..(169.1,-53.6)
				..controls(169.6,-54.3)and(169.8,-54.9)..(169.9,-55.6)
				..controls(170.4,-57.8)and(170.5,-58.1)..(171.2,-59.5)
				..controls(171.6,-60.3)and(172,-61.5)..(172.2,-62.1)
				..controls(172.3,-62.8)and(172.7,-63.7)..(172.9,-64.1)
				..controls(173.2,-64.6)and(173.4,-65.3)..(173.4,-65.6)
				..controls(173.8,-68.8)and(174.2,-70.6)..(174.9,-72.6)
				..controls(175.5,-74.5)and(175.9,-77.1)..(176.1,-80.2)
				..controls(176.3,-83.4)and(175.5,-84.4)..(173.7,-83.4)--cycle;
				\definecolor{ezs}{RGB}{247,179,26}\fill[ezs](-5.7,-156.6)--(-3.4,-157.5)--(61.9,12.7)--(59.7,13.6)--cycle;
				\definecolor{ezs}{RGB}{247,179,22}\fill[ezs](-3.8,-157.3)--(-1.5,-158.2)--(63.9,12)--(61.6,12.9)--cycle;
				\definecolor{ezs}{RGB}{248,179,19}\fill[ezs](-1.8,-158.1)--(0.4,-158.9)--(65.8,11.2)--(63.5,12.1)--cycle;
				\definecolor{ezs}{RGB}{248,178,16}\fill[ezs](0.1,-158.8)--(2.4,-159.7)--(67.7,10.5)--(65.4,11.4)--cycle;
				\definecolor{ezs}{RGB}{249,178,15}\fill[ezs](2,-159.5)--(4.3,-160.4)--(69.6,9.8)--(67.3,10.7)--cycle;
				\definecolor{ezs}{RGB}{251,177,17}\fill[ezs](3.9,-160.3)--(6.2,-161.1)--(71.5,9)--(69.2,9.9)--cycle;
				\definecolor{ezs}{RGB}{252,179,18}\fill[ezs](5.8,-161)--(8.1,-161.9)--(73.4,8.3)--(71.1,9.2)--cycle;
				\definecolor{ezs}{RGB}{252,182,17}\fill[ezs](7.7,-161.7)--(10,-162.6)--(75.3,7.6)--(73,8.5)--cycle;
				\definecolor{ezs}{RGB}{252,185,17}\fill[ezs](9.6,-162.5)--(11.9,-163.4)--(77.2,6.8)--(75,7.7)--cycle;
				\definecolor{ezs}{RGB}{251,189,18}\fill[ezs](11.5,-163.2)--(13.8,-164.1)--(79.2,6.1)--(76.9,7)--cycle;
				\definecolor{ezs}{RGB}{251,191,22}\fill[ezs](13.4,-163.9)--(15.7,-164.8)--(81.1,5.4)--(78.8,6.2)--cycle;
				\definecolor{ezs}{RGB}{252,193,25}\fill[ezs](15.4,-164.7)--(17.6,-165.6)--(83,4.6)--(80.7,5.5)--cycle;
				\definecolor{ezs}{RGB}{253,195,27}\fill[ezs](17.3,-165.4)--(19.6,-166.3)--(84.9,3.9)--(82.6,4.8)--cycle;
				\definecolor{ezs}{RGB}{253,196,31}\fill[ezs](19.2,-166.1)--(21.5,-167)--(86.8,3.2)--(84.5,4)--cycle;
				\definecolor{ezs}{RGB}{254,199,35}\fill[ezs](21.1,-166.9)--(23.4,-167.8)--(88.7,2.4)--(86.4,3.3)--cycle;
				\definecolor{ezs}{RGB}{253,201,38}\fill[ezs](23,-167.6)--(25.3,-168.5)--(90.6,1.7)--(88.3,2.6)--cycle;
				\definecolor{ezs}{RGB}{253,203,40}\fill[ezs](24.9,-168.3)--(27.2,-169.2)--(92.5,1)--(90.3,1.8)--cycle;
				\definecolor{ezs}{RGB}{253,205,42}\fill[ezs](26.8,-169.1)--(29.1,-170)--(94.4,0.2)--(92.2,1.1)--cycle;
				\definecolor{ezs}{RGB}{253,207,45}\fill[ezs](28.7,-169.8)--(31,-170.7)--(96.4,-0.5)--(94.1,0.4)--cycle;
				\definecolor{ezs}{RGB}{253,210,49}\fill[ezs](30.7,-170.5)--(32.9,-171.4)--(98.3,-1.2)--(96,-0.4)--cycle;
				\definecolor{ezs}{RGB}{253,212,54}\fill[ezs](32.6,-171.3)--(34.9,-172.2)--(100.2,-2)--(97.9,-1.1)--cycle;
				\definecolor{ezs}{RGB}{253,214,59}\fill[ezs](34.5,-172)--(36.8,-172.9)--(102.1,-2.7)--(99.8,-1.8)--cycle;
				\definecolor{ezs}{RGB}{253,216,63}\fill[ezs](36.4,-172.7)--(38.7,-173.6)--(104,-3.4)--(101.7,-2.6)--cycle;
				\definecolor{ezs}{RGB}{253,218,65}\fill[ezs](38.3,-173.5)--(40.6,-174.4)--(105.9,-4.2)--(103.6,-3.3)--cycle;
				\definecolor{ezs}{RGB}{254,220,66}\fill[ezs](40.2,-174.2)--(42.5,-175.1)--(107.8,-4.9)--(105.5,-4)--cycle;
				\definecolor{ezs}{RGB}{254,222,68}\fill[ezs](42.1,-175)--(44.4,-175.8)--(109.7,-5.6)--(107.5,-4.8)--cycle;
				\definecolor{ezs}{RGB}{254,224,71}\fill[ezs](44,-175.7)--(46.3,-176.6)--(111.7,-6.4)--(109.4,-5.5)--cycle;
				\definecolor{ezs}{RGB}{254,226,75}\fill[ezs](46,-176.4)--(48.2,-177.3)--(113.6,-7.1)--(111.3,-6.2)--cycle;
				\definecolor{ezs}{RGB}{254,228,79}\fill[ezs](47.9,-177.2)--(50.1,-178)--(115.5,-7.8)--(113.2,-7)--cycle;
				\definecolor{ezs}{RGB}{254,230,82}\fill[ezs](49.8,-177.9)--(52.1,-178.8)--(117.4,-8.6)--(115.1,-7.7)--cycle;
				\definecolor{ezs}{RGB}{254,232,85}\fill[ezs](51.7,-178.6)--(54,-179.5)--(119.3,-9.3)--(117,-8.4)--cycle;
				\definecolor{ezs}{RGB}{254,234,88}\fill[ezs](53.6,-179.4)--(55.9,-180.2)--(121.2,-10)--(118.9,-9.2)--cycle;
				\definecolor{ezs}{RGB}{254,237,91}\fill[ezs](55.5,-180.1)--(57.8,-181)--(123.1,-10.8)--(120.8,-9.9)--cycle;
				\definecolor{ezs}{RGB}{254,239,94}\fill[ezs](57.4,-180.8)--(59.7,-181.7)--(125,-11.5)--(122.8,-10.6)--cycle;
				\definecolor{ezs}{RGB}{254,241,97}\fill[ezs](59.3,-181.6)--(61.6,-182.4)--(127,-12.2)--(124.7,-11.4)--cycle;
				\definecolor{ezs}{RGB}{254,243,101}\fill[ezs](61.2,-182.3)--(63.5,-183.2)--(128.9,-13)--(126.6,-12.1)--cycle;
				\definecolor{ezs}{RGB}{254,245,103}\fill[ezs](63.2,-183)--(65.4,-183.9)--(130.8,-13.7)--(128.5,-12.8)--cycle;
				\definecolor{ezs}{RGB}{254,245,105}\fill[ezs](65.1,-183.8)--(67.4,-184.6)--(132.7,-14.4)--(130.4,-13.6)--cycle;
				\definecolor{ezs}{RGB}{254,245,105}\fill[ezs](67,-184.5)--(69.3,-185.4)--(134.6,-15.2)--(132.3,-14.3)--cycle;
				\definecolor{ezs}{RGB}{254,244,103}\fill[ezs](68.9,-185.2)--(71.2,-186.1)--(136.5,-15.9)--(134.2,-15)--cycle;
				\definecolor{ezs}{RGB}{253,243,101}\fill[ezs](70.8,-186)--(73.1,-186.8)--(138.4,-16.6)--(136.1,-15.8)--cycle;
				\definecolor{ezs}{RGB}{253,241,98}\fill[ezs](72.7,-186.7)--(75,-187.6)--(140.3,-17.4)--(138,-16.5)--cycle;
				\definecolor{ezs}{RGB}{253,239,95}\fill[ezs](74.6,-187.4)--(76.9,-188.3)--(142.2,-18.1)--(140,-17.2)--cycle;
				\definecolor{ezs}{RGB}{253,236,92}\fill[ezs](76.5,-188.2)--(78.8,-189)--(144.2,-18.8)--(141.9,-18)--cycle;
				\definecolor{ezs}{RGB}{253,234,88}\fill[ezs](78.5,-188.9)--(80.7,-189.8)--(146.1,-19.6)--(143.8,-18.7)--cycle;
				\definecolor{ezs}{RGB}{253,232,85}\fill[ezs](80.4,-189.6)--(82.7,-190.5)--(148,-20.3)--(145.7,-19.4)--cycle;
				\definecolor{ezs}{RGB}{253,230,82}\fill[ezs](82.3,-190.4)--(84.6,-191.2)--(149.9,-21)--(147.6,-20.2)--cycle;
				\definecolor{ezs}{RGB}{253,228,78}\fill[ezs](84.2,-191.1)--(86.5,-192)--(151.8,-21.8)--(149.5,-20.9)--cycle;
				\definecolor{ezs}{RGB}{253,226,74}\fill[ezs](86.1,-191.8)--(88.4,-192.7)--(153.7,-22.5)--(151.4,-21.6)--cycle;
				\definecolor{ezs}{RGB}{253,224,70}\fill[ezs](88,-192.6)--(90.3,-193.4)--(155.6,-23.3)--(153.3,-22.4)--cycle;
				\definecolor{ezs}{RGB}{253,222,67}\fill[ezs](89.9,-193.3)--(92.2,-194.2)--(157.5,-24)--(155.3,-23.1)--cycle;
				\definecolor{ezs}{RGB}{253,220,65}\fill[ezs](91.8,-194)--(94.1,-194.9)--(159.5,-24.7)--(157.2,-23.8)--cycle;
				\definecolor{ezs}{RGB}{253,217,63}\fill[ezs](93.7,-194.8)--(96,-195.6)--(161.4,-25.5)--(159.1,-24.6)--cycle;
				\definecolor{ezs}{RGB}{253,215,60}\fill[ezs](95.7,-195.5)--(97.9,-196.4)--(163.3,-26.2)--(161,-25.3)--cycle;
				\definecolor{ezs}{RGB}{253,213,56}\fill[ezs](97.6,-196.2)--(99.9,-197.1)--(165.2,-26.9)--(162.9,-26)--cycle;
				\definecolor{ezs}{RGB}{252,211,53}\fill[ezs](99.5,-197)--(101.8,-197.8)--(167.1,-27.7)--(164.8,-26.8)--cycle;
				\definecolor{ezs}{RGB}{252,209,49}\fill[ezs](101.4,-197.7)--(103.7,-198.6)--(169,-28.4)--(166.7,-27.5)--cycle;
				\definecolor{ezs}{RGB}{252,207,45}\fill[ezs](103.3,-198.4)--(105.6,-199.3)--(170.9,-29.1)--(168.6,-28.2)--cycle;
				\definecolor{ezs}{RGB}{253,205,42}\fill[ezs](105.2,-199.2)--(107.5,-200)--(172.8,-29.9)--(170.6,-29)--cycle;
				\definecolor{ezs}{RGB}{253,202,40}\fill[ezs](107.1,-199.9)--(109.4,-200.8)--(174.7,-30.6)--(172.5,-29.7)--cycle;
				\definecolor{ezs}{RGB}{253,200,38}\fill[ezs](109,-200.6)--(111.3,-201.5)--(176.7,-31.3)--(174.4,-30.4)--cycle;
				\definecolor{ezs}{RGB}{253,198,35}\fill[ezs](111,-201.4)--(113.2,-202.2)--(178.6,-32.1)--(176.3,-31.2)--cycle;
				\definecolor{ezs}{RGB}{253,196,31}\fill[ezs](112.9,-202.1)--(115.2,-203)--(180.5,-32.8)--(178.2,-31.9)--cycle;
				\definecolor{ezs}{RGB}{253,194,28}\fill[ezs](114.8,-202.8)--(117.1,-203.7)--(182.4,-33.5)--(180.1,-32.6)--cycle;
				\definecolor{ezs}{RGB}{253,192,26}\fill[ezs](116.7,-203.6)--(119,-204.5)--(184.3,-34.3)--(182,-33.4)--cycle;
				\definecolor{ezs}{RGB}{253,189,24}\fill[ezs](118.6,-204.3)--(120.9,-205.2)--(186.2,-35)--(183.9,-34.1)--cycle;
				\definecolor{ezs}{RGB}{253,186,20}\fill[ezs](120.5,-205)--(122.8,-205.9)--(188.1,-35.7)--(185.8,-34.8)--cycle;
				\definecolor{ezs}{RGB}{253,184,17}\fill[ezs](122.4,-205.8)--(124.7,-206.7)--(190,-36.5)--(187.8,-35.6)--cycle;
				\definecolor{ezs}{RGB}{252,183,15}\fill[ezs](124.3,-206.5)--(126.6,-207.4)--(192,-37.2)--(189.7,-36.3)--cycle;
				\definecolor{ezs}{RGB}{252,181,14}\fill[ezs](126.3,-207.2)--(128.5,-208.1)--(193.9,-37.9)--(191.6,-37.1)--cycle;
				\definecolor{ezs}{RGB}{252,179,13}\fill[ezs](128.2,-208)--(130.4,-208.9)--(195.8,-38.7)--(193.5,-37.8)--cycle;
				\definecolor{ezs}{RGB}{252,178,12}\fill[ezs](130.1,-208.7)--(132.4,-209.6)--(197.7,-39.4)--(195.4,-38.5)--cycle;
				\definecolor{ezs}{RGB}{251,178,10}\fill[ezs](132,-209.4)--(134.3,-210.3)--(199.6,-40.1)--(197.3,-39.3)--cycle;
				\definecolor{ezs}{RGB}{251,178,10}\fill[ezs](133.9,-210.2)--(136.2,-211.1)--(201.5,-40.9)--(199.2,-40)--cycle;
				\definecolor{ezs}{RGB}{250,177,11}\fill[ezs](135.8,-210.9)--(138.1,-211.8)--(203.4,-41.6)--(201.1,-40.7)--cycle;
				\definecolor{ezs}{RGB}{250,178,11}\fill[ezs](137.7,-211.6)--(140,-212.5)--(205.3,-42.3)--(203.1,-41.5)--cycle;
				\definecolor{ezs}{RGB}{248,180,10}\fill[ezs](139.6,-212.4)--(141.9,-213.3)--(207.3,-43.1)--(205,-42.2)--cycle;
				\definecolor{ezs}{RGB}{247,181,9}\fill[ezs](141.5,-213.1)--(143.8,-214)--(209.2,-43.8)--(206.9,-42.9)--cycle;
				\definecolor{ezs}{RGB}{246,182,9}\fill[ezs](143.5,-213.8)--(145.7,-214.7)--(211.1,-44.5)--(208.8,-43.7)--cycle;
				\definecolor{ezs}{RGB}{246,182,10}\fill[ezs](145.4,-214.6)--(147.7,-215.5)--(213,-45.3)--(210.7,-44.4)--cycle;
			\end{scope}
			% chữ EDUZY (Times đậm, làm dày nét bằng 5 lớp lệch nhẹ; căn đúng khung chữ logo gốc)
			\foreach \dx/\dy in {-0.7/0,0.7/0,0/-0.3,0/0.3,0/0}
			\node[anchor=center,inner sep=0pt,outer sep=0pt,text=ezText,align=center] at (109.6+\dx,-193.6+\dy)
			{\resizebox{1.01cm}{0.22cm}{\fontencoding{T1}\fontfamily{ptm}\fontseries{b}\selectfont EDUZY}};
		\end{tikzpicture}%
	}%
}
%====================== THIẾT LẬP THÂN ĐỀ ======================
%--- Môi trường tự luận {bt}: ex_test 3.3.4 chưa khai báo sẵn nên tạo ở đây
%    (ex_test đã tự móc \fixshowans và \listenumerate cho bt)
\makeatletter
\@ifundefined{bt}{\newtheorem{bt}{\color{gold1}\namebt}}{}
\makeatother
\renewcommand{\sectionmark}[1]{\markright{\thesection\ #1}}
%--- Khung tiêu đề (logo + tối đa 4 dòng; dòng 3, 4 để trống {} thì bỏ qua)
%  Đề thi:    \khungtieude{THI THỬ THPT QUỐC GIA NĂM 2026}{Môn Toán -- Đề số 21}
%                         {(Đề thi gồm: ...)}{Thời gian làm bài: 90 phút (không kể thời gian phát đề).}
%  Chuyên đề: \khungtieude{CHUYÊN ĐỀ 3 -- DÃY SỐ VÀ CÔNG THỨC TRUY HỒI}
%                         {Bồi dưỡng học sinh giỏi Toán THPT}{...}{}
\newcommand{\khungtieude}[4]{%
	\begin{tcolorbox}[colframe=gold1,colback=gold5!45,boxrule=0.5mm]
		\begin{minipage}[c]{3cm}
			\centerline{\eduzylogo[2.4cm]}
		\end{minipage}\hspace{0.5cm}
		\begin{minipage}[c]{12.5cm}
			\vskip 0.15 cm
			\centerline{\ind{#1}}
			\vskip 0.15 cm
			\centerline{\textbf{#2}}
			\if\relax\detokenize{#3}\relax\else
				\vskip 0.15 cm
				\centerline{\textit{#3}}
			\fi
			\if\relax\detokenize{#4}\relax\else
				\vskip 0.15 cm
				\centerline{\textit{#4}}
			\fi
			\vskip 0.15 cm
		\end{minipage}
	\end{tcolorbox}
}
%====================== BỔ SUNG CHO BÀI HỌC, CHUYÊN ĐỀ ======================
\usepackage[shortlabels]{enumitem}
\usepackage{tabularx}
%--- Danh sách gọn: \begin{enumerate}[gon] ... \end{enumerate}
\SetEnumitemKey{gon}{leftmargin=*,itemsep=1pt,parsep=0pt,topsep=3pt}
%--- Khung lý thuyết: \begin{khung}{} ... \end{khung}; tiêu đề chỉ dùng khi thật ngắn ("Mẹo", "Lỗi thường gặp")
\newtcolorbox{khung}[1]{breakable,enhanced,colframe=gold1,colback=gold5!45,
	boxrule=0.4mm,fonttitle=\bfseries\sffamily,coltitle=white,colbacktitle=gold1,
	attach boxed title to top left={xshift=6mm,yshift=-3mm},
	boxed title style={arc=1mm,boxrule=0.3mm},title={#1}}
%--- Tiêu đề mục nhỏ trong một dạng: \muc{Dấu hiệu}, \muc{Phương pháp}, \muc{Bài tự luyện}
\makeatletter
\newcommand{\muc}[1]{\par\medskip\noindent\indam{#1}\par\nobreak\smallskip\@afterindentfalse\@afterheading}
\makeatother
%--- Ô công thức cần nhớ (căn giữa, viền vàng, vừa khít nội dung)
%     \begin{center}\congthuc{$Q(n)=(n-1)!$}\end{center}
\newtcbox{\congthuc}{on line,colframe=gold1,colback=gold5!45,boxrule=0.5pt,arc=2pt,
	left=8pt,right=8pt,top=4pt,bottom=4pt}
%--- Công tắc lời giải bài tự luyện (ví dụ mẫu luôn dùng \loigiai)
%     \loigiaiTL[{đáp số}]{lời giải đầy đủ}
%     \TLfull -> in lời giải đầy đủ;  \TLdapso -> chỉ in một dòng "Đáp số";  \TLan -> ẩn hoàn toàn
%     (\CheDo là dethi thì mọi lời giải và đáp số đều ẩn)
\makeatletter
\newif\ifTLhien   \TLhientrue
\newif\ifTLdapso  \TLdapsofalse
\newcommand{\TLfull}{\TLhientrue\TLdapsofalse}
\newcommand{\TLdapso}{\TLhientrue\TLdapsotrue}
\newcommand{\TLan}{\TLhienfalse}
\newcommand{\loigiaiTL}[2][]{%
	\ifTLhien
		\ifTLdapso
			\def\@tlarg{#1}%
			\ifx\@tlarg\@empty\else
				\ifdefstring{\kindOPTN}{dethi}{}{\par\noindent{\bfseries\color{gold1}Đáp số:} #1\par}%
			\fi
		\else
			\loigiai{#2}%
		\fi
	\fi}
\makeatother
\TLdapso% mặc định: chỉ in đáp số bài tự luyện
%====================== BẢNG ĐÚNG/SAI (\choiceTFt) ======================
%  Hàng tiêu đề nền nâu đồng chữ trắng (Phát biểu | Đ | S), các hàng tô xen kẽ kem / trắng,
%  đường kẻ vàng kim, nhãn a) b) c) d) màu nâu đồng. Bản có đáp án: cột được chọn có vòng tròn
%  nâu đồng với dấu tích trắng (câu đúng tích ở cột Đúng, câu sai tích ở cột Sai).
%  Tiêu đề ghi đủ "Đúng | Sai": \OPTN{viettat=0} sau \input (cột Đúng/Sai rộng ra, phát biểu hẹp lại).
%  Dùng bảng gốc của ex_test: \def\BangDSgoc{1} trước \input{preamble-eduzy-v2}.
\providecommand{\BangDSgoc}{0}
\ifnum\BangDSgoc=0
\usepackage{colortbl,pifont}
\colorlet{ezDSke}{gold2!60!gold1}% đường kẻ bảng (nét 0,4pt như ex_test để không đổi chiều cao bảng)
\colorlet{ezDSle}{gold5!60}% hàng a), c)
\colorlet{ezDSchan}{white}% hàng b), d)
%--- dấu chọn ở cột Đúng/Sai
%    (bề rộng 0, căn giữa ô: không làm đổi bề rộng cột của ex_test, không báo tràn lề)
\newcommand{\EZdauchon}{\makebox[0pt]{\tikz[baseline=-1ex]\node[circle,fill=gold1,inner sep=0pt,
	minimum size=1.2em,text=white,font=\small]{\ding{51}};}}
\def\tickT{\EZdauchon}\def\tickF{\EZdauchon}
\AtBeginDocument{\sbox{\ccTFTab}{\EZdauchon}}
%--- nhãn a) b) c) d) màu nâu đồng (ở bản lời giải, nhãn câu đúng vẫn đóng khung như ex_test)
\renewcommand{\FalseTF}{\stepcounter{dapan}\textbf{\color{gold1}\DapAnTF\sepTF}\ignorespaces}
%--- trước khi dựng bảng: kẻ màu vàng kim (bề rộng cột giữ nguyên như ex_test); sau bảng trả lại màu đen
\newcount\EZhangTF
\newcommand{\EZbatdauTF}{%
	\global\EZhangTF=0\relax
	\arrayrulecolor{ezDSke}}
\newcommand{\EZketthucTF}{\arrayrulecolor{black}}
%--- hàng tiêu đề nền nâu đồng, chữ trắng
\makeatletter
\newcommand{\EZo@TF}[1]{\centering\bfseries\color{white}#1}
\newcommand{\EZtieudeTF}{%
	\ifcase\cotDS
		\ifnum\cotSTT=1\relax
			\def\rowHeader{\hline\rowcolor{gold1}\EZo@TF{TT}&\EZo@TF{\phatbieu}\tabularnewline}%
		\else
			\def\rowHeader{\hline\rowcolor{gold1}\EZo@TF{\phatbieu}\tabularnewline}%
		\fi
	\or
		\ifnum\cotSTT=1\relax
			\def\rowHeader{\hline\rowcolor{gold1}\EZo@TF{TT}&\EZo@TF{\phatbieu}&\EZo@TF{\Dung}\tabularnewline}%
		\else
			\def\rowHeader{\hline\rowcolor{gold1}\EZo@TF{\phatbieu}&\EZo@TF{\Dung}\tabularnewline}%
		\fi
	\or
		\ifnum\cotSTT=1\relax
			\def\rowHeader{\hline\rowcolor{gold1}\EZo@TF{TT}&\EZo@TF{\phatbieu}&\EZo@TF{\Dung}&\EZo@TF{\Sai}\tabularnewline}%
		\else
			\def\rowHeader{\hline\rowcolor{gold1}\EZo@TF{\phatbieu}&\EZo@TF{\Dung}&\EZo@TF{\Sai}\tabularnewline}%
		\fi
	\fi}
%--- tô xen kẽ từng hàng (gọi ngay sau \hline đầu mỗi hàng)
\newcommand{\EZmauhangTF}{%
	\noalign{\global\advance\EZhangTF\@ne}%
	\ifodd\EZhangTF\rowcolor{ezDSle}\else\rowcolor{ezDSchan}\fi}
%--- gắn vào \choiceTFt / \writeChoiceTFt của ex_test
\xpretocmd{\choiceTFt}{\EZbatdauTF}{}{\typeout{[preamble-eduzy-v2] Không gắn được bảng Đúng/Sai (đầu bảng)}}
\xpatchcmd{\choiceTFt}{\ifnum\dongPB=0\def\rowHeader{}\fi}{\EZtieudeTF\ifnum\dongPB=0\def\rowHeader{}\fi}
	{}{\typeout{[preamble-eduzy-v2] Không gắn được bảng Đúng/Sai (hàng tiêu đề)}}
\xpatchcmd{\writeChoiceTFt}{\hline\xdef\dotEX{\sdotEX}}{\hline\EZmauhangTF\xdef\dotEX{\sdotEX}}
	{}{\typeout{[preamble-eduzy-v2] Không gắn được bảng Đúng/Sai (tô hàng)}}
\xpatchcmd{\writeChoiceTFt}{{\end{longtable}}}{{\end{longtable}\EZketthucTF}}
	{}{\typeout{[preamble-eduzy-v2] Không gắn được bảng Đúng/Sai (cuối longtable)}}
\xpatchcmd{\writeChoiceTFt}{{\end{xtabular}\end{center}\medskip\noindent}}{{\end{xtabular}\end{center}\EZketthucTF\medskip\noindent}}
	{}{\typeout{[preamble-eduzy-v2] Không gắn được bảng Đúng/Sai (cuối xtabular)}}
\fi
%=====================================================================
\renewcommand{\parallel}{\mathrel{/\!\!/}}

%  trong file đề (tools/dan_trang/dan_trang.py tự ghi khi đặt \NeedspaceCau):
%    \def\TachCauDS{1}   câu Đúng/Sai được tách: câu dẫn (kèm hình) ở cuối trang, bảng Phát biểu sang trang sau
%                        (câu dẫn và hình không tách nhau, bảng không tách giữa hai trang)
%    \def\TachCauLG{1}   bản lời giải (\CheDo khác dethi): phần đề của câu cao hơn 1/3 trang được phép tách
%  Lệnh dùng trong file đề (giá trị do tools/dan_trang/dan_trang.py đo ở bản dethi):
%    \NeedspaceCau[<cao câu dẫn + hình>]{<cao cả khối câu>}   ngay trên \begin{ex} (câu đầu phần: trên tiêu đề phần)
%        công tắc tắt: đúng bằng \Needspace{<cao cả khối câu>} như cũ;
%        \TachCauDS=1 và có [..] (câu Đúng/Sai được chọn tách): chỉ giữ câu dẫn + hình;
%        \TachCauLG=1 ở bản lời giải: đòi tối đa 1/3 \textheight.
%    \NeedspaceBang{<cao bảng Phát biểu>}%   ngay trước \choiceTFt: \TachCauDS=1 thì giữ trọn bảng trên một trang;
%        công tắc tắt: không làm gì.
\providecommand{\TachCauDS}{0}
\providecommand{\TachCauLG}{0}
\newlength{\EZnsCau}
\newcommand{\NeedspaceCau}[2][]{%
	\setlength{\EZnsCau}{#2}%
	\ifnum\TachCauDS=1 \if\relax\detokenize{#1}\relax\else\setlength{\EZnsCau}{#1}\fi\fi
	\ifnum\TachCauLG=1 \ifdefstring{\kindOPTN}{dethi}{}{\ifdim\EZnsCau>0.3333\textheight\setlength{\EZnsCau}{0.3333\textheight}\fi}\fi
	\Needspace{\EZnsCau}}
\newcommand{\NeedspaceBang}[1]{\ifnum\TachCauDS=1 \Needspace{#1}\fi}

%--- Thang màu: đậm -> nhạt
\definecolor{gold1}{RGB}{126,96,0}    % nâu đồng đậm  #7E6000
\definecolor{gold2}{RGB}{192,144,3}   % vàng kim      #C09003
\definecolor{gold3}{RGB}{255,215,99}  % vàng nhạt     #FFD763
\definecolor{gold4}{RGB}{255,231,153} % vàng rất nhạt #FFE799
\definecolor{gold5}{RGB}{255,243,205} % vàng kem      #FFF3CD
%--- Ghi đè tông xanh/đỏ cũ
\definecolor{myblue}{RGB}{126,96,0}
\colorlet{blue}{gold1}
\colorlet{cyan}{gold2}
\colorlet{teal}{gold2}

%--- Tên chương mặc định (file đề nên tự khai báo lại)
\providecommand{\tenchuong}{}
%--- Chân trang
\providecommand{\thuonghieu}{TOÁN 12 -- Eduzy.com.vn}
\providecommand{\tacgia}{Tô Gia Bảo -- 0353703783}

%--- Chữ cạnh logo ở đầu trang
\providecommand{\khauhieu}{Eduzy.com.vn}

%====================== THÔNG SỐ ĐẦU TRANG, CHÂN TRANG (v2) ======================
%  Dải nền đậm tràn ngang trang; ô chữ vàng kem có cạnh xiên, kèm 2 vạch xiên chuyển màu.
%  Đầu trang lẻ : logo + \khauhieu (trái)   | ô chữ \tenchuong (phải)
%  Chân trang lẻ: ô chữ \thuonghieu (trái)  | Trang (số) (phải)
%  Trang chẵn soi gương trang lẻ; chân trang chẵn ghi \tacgia.
%--- Màu: chỉ dùng thang vàng kim Eduzy
\colorlet{ezNen}{gold1}% dải nền đậm
\colorlet{ezVach1}{gold3}% vạch xiên sát dải nền (đậm hơn)
\colorlet{ezVach2}{gold4}% vạch xiên sát ô chữ (nhạt hơn)
\colorlet{ezO}{gold5}% ô chữ vàng kem
%--- Kích thước
\newlength{\EZdai}\setlength{\EZdai}{1.1cm}% chiều cao dải nền
\newlength{\EZo}\setlength{\EZo}{0.68cm}% chiều cao ô chữ
\newlength{\EZxien}\setlength{\EZxien}{0.55cm}% độ lệch ngang của cạnh xiên
\newlength{\EZvach}\setlength{\EZvach}{0.75cm}% bề rộng mỗi vạch xiên
\newlength{\EZkhe}\setlength{\EZkhe}{0.12cm}% khe trắng giữa vạch và ô chữ
\newlength{\EZdem}\setlength{\EZdem}{0.4cm}% khoảng đệm từ chữ tới cạnh xiên
\newlength{\EZle}\setlength{\EZle}{1.8cm}% lề trái/phải của chữ (bằng lề thân bài)
\newlength{\EZrongmin}\setlength{\EZrongmin}{4.5cm}% ô chữ rộng tối thiểu
\newlength{\EZrong}% bề rộng chữ trong ô (tính tự động)
%--- Font
\newcommand{\EZfontO}{\fontsize{10.5pt}{12pt}\selectfont\bfseries}% chữ trong ô
\newcommand{\EZfontKH}{\fontfamily{qag}\fontsize{11pt}{12pt}\selectfont\bfseries}% \khauhieu cạnh logo
\newcommand{\EZfontTrang}{\fontsize{11.5pt}{12pt}\selectfont\bfseries}% chữ "Trang" và số trang
%--- Đo bề rộng chữ trong ô
\newcommand{\EZdorong}[1]{%
	\settowidth{\EZrong}{\EZfontO #1}%
	\ifdim\EZrong<\EZrongmin\setlength{\EZrong}{\EZrongmin}\fi}
%--- Logo + chữ trắng \khauhieu (dựng sẵn một lần khi vào \begin{document})
\newsavebox{\EZlogobox}
\AtBeginDocument{\sbox{\EZlogobox}{%
		\raisebox{-0.5\height}{\eduzylogo[1cm]}\hspace{0.15cm}%
		\raisebox{-0.5\height}{\color{white}\EZfontKH\khauhieu}}}
%--- Hai vạch xiên + khe trắng. #1: điểm gốc ở mép dưới, #2: hướng (1 sang phải, -1 sang trái),
%    #3: hướng nghiêng (1 kiểu "/", -1 kiểu "\")
\newcommand{\EZvachxien}[3]{%
	\fill[white] (#1) --++({(#3)*\EZxien},\EZo) --++({(#2)*\EZkhe},0) --++({-(#3)*\EZxien},-\EZo) --cycle;
	\fill[ezVach2] ([xshift=(#2)*\EZkhe]#1) --++({(#3)*\EZxien},\EZo) --++({(#2)*\EZvach},0) --++({-(#3)*\EZxien},-\EZo) --cycle;
	\fill[ezVach1] ([xshift=(#2)*(\EZkhe+\EZvach)]#1) --++({(#3)*\EZxien},\EZo) --++({(#2)*\EZvach},0) --++({-(#3)*\EZxien},-\EZo) --cycle;
	\draw[white,line width=0.6pt] ([xshift=(#2)*(\EZkhe+\EZvach)]#1) --++({(#3)*\EZxien},\EZo);
}
%--- Chữ trong ô (căn giữa theo chiều cao ô, không lệch vì dấu tiếng Việt)
\tikzset{ezchu/.style={inner sep=0pt,text=ezNen,text height=0.72em,text depth=0.22em}}

%====================== LOGO CHÌM GIỮA TRANG ======================
%  Vẽ cùng đầu trang nên nằm dưới chữ của thân bài, có trên mọi trang.
%  Tắt: \logochimtat (từ trang đó trở đi), bật lại: \logochimbat.
%  Đổi cỡ, độ đậm: khai báo trước \input, ví dụ \def\EZchimcao{9cm} \def\EZchimdo{0.1}
\providecommand{\EZchimcao}{10cm}% chiều cao logo chìm
\providecommand{\EZchimdo}{0.12}% độ đậm (0 = trong suốt, 1 = rõ hẳn)
\newif\ifEZchim \EZchimtrue
\newcommand{\logochimtat}{\EZchimfalse}
\newcommand{\logochimbat}{\EZchimtrue}
%  Dựng sẵn logo mờ một lần khi vào \begin{document}; transparency group làm mờ cả logo
%  một lượt nên không lộ vệt ở chỗ các dải màu chồng nhau.
\newsavebox{\EZchimtho}
\newsavebox{\EZchimbox}
\AtBeginDocument{%
	\sbox{\EZchimtho}{\eduzylogo[\EZchimcao]}%
	\sbox{\EZchimbox}{\begin{tikzpicture}
			\begin{scope}[transparency group,opacity=\EZchimdo]
				\node[inner sep=0pt] {\usebox{\EZchimtho}};
			\end{scope}
	\end{tikzpicture}}}
\newcommand{\EZvechim}{%
	\ifEZchim
		\node[inner sep=0pt] at (current page.center) {\usebox{\EZchimbox}};
	\fi}

%====================== ĐẦU TRANG - TRANG LẺ ======================
%  (logo bên trái, ô chữ \tenchuong bên phải)
\newcommand{\myfancyheadOdd}{%
	\EZdorong{\tenchuong}%
	\begin{tikzpicture}[remember picture,overlay]
		\coordinate (T) at (current page.north west);
		\coordinate (P) at (current page.north east);
		%--- logo chìm giữa trang
		\EZvechim
		%--- dải nền
		\fill[ezNen] (T) rectangle ([yshift=-\EZdai]P);
		%--- ô chữ: X là góc dưới bên trái, cạnh trái xiên kiểu "/"
		\coordinate (X) at ([xshift=-\EZle-\EZrong-\EZdem-\EZxien,yshift=-\EZdai]P);
		\fill[ezO] (X) --++(\EZxien,\EZo) -- ([yshift=-\EZdai+\EZo]P) -- ([yshift=-\EZdai]P) --cycle;
		\EZvachxien{X}{-1}{1}
		\node[ezchu,anchor=east] at ([xshift=-\EZle,yshift=-\EZdai+0.5*\EZo]P) {\EZfontO\tenchuong};
		%--- logo + tên
		\node[inner sep=0pt,anchor=west] at ([xshift=\EZle,yshift=-0.5*\EZdai]T) {\usebox{\EZlogobox}};
	\end{tikzpicture}%
}

%====================== ĐẦU TRANG - TRANG CHẴN =====================
%  (soi gương: ô chữ \tenchuong bên trái, logo bên phải)
\newcommand{\myfancyheadEven}{%
	\EZdorong{\tenchuong}%
	\begin{tikzpicture}[remember picture,overlay]
		\coordinate (T) at (current page.north west);
		\coordinate (P) at (current page.north east);
		%--- logo chìm giữa trang
		\EZvechim
		%--- dải nền
		\fill[ezNen] (T) rectangle ([yshift=-\EZdai]P);
		%--- ô chữ: X là góc dưới bên phải, cạnh phải xiên kiểu "\"
		\coordinate (X) at ([xshift=\EZle+\EZrong+\EZdem+\EZxien,yshift=-\EZdai]T);
		\fill[ezO] (X) --++(-\EZxien,\EZo) -- ([yshift=-\EZdai+\EZo]T) -- ([yshift=-\EZdai]T) --cycle;
		\EZvachxien{X}{1}{-1}
		\node[ezchu,anchor=west] at ([xshift=\EZle,yshift=-\EZdai+0.5*\EZo]T) {\EZfontO\tenchuong};
		%--- logo + tên
		\node[inner sep=0pt,anchor=east] at ([xshift=-\EZle,yshift=-0.5*\EZdai]P) {\usebox{\EZlogobox}};
	\end{tikzpicture}%
}

%====================== CHÂN TRANG - TRANG LẺ ======================
%  (ô chữ \thuonghieu bên trái, "Trang (số)" bên phải)
\newcommand{\myfancyfootOdd}{%
	\EZdorong{\thuonghieu}%
	\begin{tikzpicture}[remember picture,overlay]
		\coordinate (T) at (current page.south west);
		\coordinate (P) at (current page.south east);
		%--- dải nền
		\fill[ezNen] (T) rectangle ([yshift=\EZdai]P);
		%--- ô chữ sát mép trên dải: X là góc dưới bên phải, cạnh phải xiên kiểu "/"
		\coordinate (X) at ([xshift=\EZle+\EZrong+\EZdem,yshift=\EZdai-\EZo]T);
		\fill[ezO] ([yshift=\EZdai-\EZo]T) -- (X) --++(\EZxien,\EZo) -- ([yshift=\EZdai]T) --cycle;
		\EZvachxien{X}{1}{1}
		\node[ezchu,anchor=west] at ([xshift=\EZle,yshift=\EZdai-0.5*\EZo]T) {\EZfontO\thuonghieu};
		%--- số trang
		\node[circle,draw=white,line width=0.9pt,text=white,inner sep=0pt,minimum size=0.62cm,anchor=east]
		(so) at ([xshift=-\EZle,yshift=0.5*\EZdai]P) {\EZfontTrang\thepage};
		\node[inner sep=0pt,text=white,anchor=base east] (tr) at ([xshift=-0.15cm]so.base west) {\EZfontTrang Trang};
	\end{tikzpicture}%
}

%====================== CHÂN TRANG - TRANG CHẴN ====================
%  (soi gương: "Trang (số)" bên trái, ô chữ \tacgia bên phải)
\newcommand{\myfancyfootEven}{%
	\EZdorong{\tacgia}%
	\begin{tikzpicture}[remember picture,overlay]
		\coordinate (T) at (current page.south west);
		\coordinate (P) at (current page.south east);
		%--- dải nền
		\fill[ezNen] (T) rectangle ([yshift=\EZdai]P);
		%--- ô chữ sát mép trên dải: X là góc dưới bên trái, cạnh trái xiên kiểu "\"
		\coordinate (X) at ([xshift=-\EZle-\EZrong-\EZdem,yshift=\EZdai-\EZo]P);
		\fill[ezO] ([yshift=\EZdai-\EZo]P) -- (X) --++(-\EZxien,\EZo) -- ([yshift=\EZdai]P) --cycle;
		\EZvachxien{X}{-1}{-1}
		\node[ezchu,anchor=east] at ([xshift=-\EZle,yshift=\EZdai-0.5*\EZo]P) {\EZfontO\tacgia};
		%--- số trang
		%    (vòng tròn ẩn so0 nằm giữa dải để lấy đường chân chữ, giống trang lẻ)
		\node[circle,inner sep=0pt,minimum size=0.62cm] (so0) at ([yshift=0.5*\EZdai]T) {\EZfontTrang\phantom{\thepage}};
		\node[inner sep=0pt,text=white,anchor=base west] (tr) at ($(so0.base -| T)+(\EZle,0)$) {\EZfontTrang Trang};
		\node[circle,draw=white,line width=0.9pt,text=white,inner sep=0pt,minimum size=0.62cm,anchor=base west]
		(so) at ([xshift=0.15cm]tr.base east) {\EZfontTrang\thepage};
	\end{tikzpicture}%
}

%--- Giữ tương thích với các file cũ đang gọi \myfancyhead, \myfancyfoot
\newcommand{\myfancyhead}{\myfancyheadOdd}
\newcommand{\myfancyfoot}{\myfancyfootOdd}

%======================= GẮN VÀO TRANG =============================
\pagestyle{fancy}
\fancyhf{}
\fancyhead[LO]{\myfancyheadOdd}
\fancyhead[RE]{\myfancyheadEven}
\fancyfoot[LO]{\myfancyfootOdd}
\fancyfoot[RE]{\myfancyfootEven}
\renewcommand{\headrulewidth}{0pt}
\renewcommand{\footrulewidth}{0pt}
\fancypagestyle{plain}{%
	\fancyhf{}
	\fancyhead[LO]{\myfancyheadOdd}
	\fancyhead[RE]{\myfancyheadEven}
	\fancyfoot[LO]{\myfancyfootOdd}
	\fancyfoot[RE]{\myfancyfootEven}
	\renewcommand{\headrulewidth}{0pt}
	\renewcommand{\footrulewidth}{0pt}
}
%=================== HẾT PHẦN GỘP TỪ eduzy-de.tex ===================
\newcommand{\indam}[1]{
	{\fontsize{12pt}{9pt}\sffamily \selectfont\bfseries\color{gold1} #1}
}
\newcommand{\ind}[1]{
	{\fontsize{14pt}{9pt}\fontfamily{qag} \selectfont\bfseries\color{gold1} #1}
}
\renewcommand{\FalseEX}{\stepcounter{dapan}\textbf{\selectfont\color{gold1}\Alph{dapan}.}\ignorespaces}
\renewtheorem{ex}{\color{gold1}Câu}
\renewcommand{\baselinestretch}{1.1}
\newcommand{\hoac}[1]{
	\left[\begin{aligned}#1\end{aligned}\right.}
\newcommand{\heva}[1]{
	\left\{\begin{aligned}#1\end{aligned}\right.}
\usepackage{esvect}
\def\vec{\vv}
\def\vt{\vv}
\def\overrightarrow{\vv}
%====================== LOGO EDUZY VẼ BẰNG TIKZ ======================
%  Đường viền dò từ logo gốc; dải màu vàng kim đo trực tiếp trên ảnh gốc
%  (nghiêng 21 độ, tô bằng 80 dải hẹp nên hiển thị giống nhau trên mọi trình xem PDF).
%  Dùng: \eduzylogo          -> cao 2.4cm (mặc định)
%        \eduzylogo[1.5cm]   -> cao 1.5cm
\definecolor{ezText}{RGB}{245,200,46}
\newcommand{\eduzylogo}[1][2.4cm]{%
	\resizebox{!}{#1}{%
		\begin{tikzpicture}[x=0.01cm,y=0.01cm]
			\path[use as bounding box] (0,0) rectangle (224,-224);
			% mũ tốt nghiệp + vương miện
			\begin{scope}
				\clip
				(119.4,-176.1)..controls(118.9,-176.1)and(116.8,-176)..(114.6,-176)
				..controls(101.4,-176)and(85.8,-175.9)..(83.7,-175.7)
				..controls(81.8,-175.6)and(81,-175.6)..(80.2,-175.7)
				..controls(78.8,-175.9)and(77.7,-175.9)..(76.2,-175.7)
				..controls(74.9,-175.5)and(70.8,-175.5)..(64.3,-175.7)
				..controls(59.1,-175.9)and(58.3,-175.7)..(58.6,-173.7)
				..controls(58.6,-172.8)and(59.6,-169.6)..(60.3,-168)
				..controls(60.8,-166.9)and(61.4,-165.8)..(62.3,-165)
				..controls(63.1,-164.2)and(64.3,-163.4)..(65,-163.2)
				..controls(65.5,-163.1)and(72.3,-163)..(76.6,-163.1)
				..controls(77.9,-163.1)and(80.1,-163.1)..(81.4,-163)
				..controls(82.8,-162.9)and(86.3,-162.9)..(89.2,-162.9)
				..controls(93,-162.9)and(94.9,-162.8)..(96,-162.7)
				..controls(97.2,-162.5)and(101,-162.5)..(114,-162.4)
				..controls(123,-162.3)and(134.9,-162.2)..(140.2,-162.1)
				..controls(150.1,-162)and(152.2,-162.1)..(153.6,-162.4)
				..controls(154.4,-162.6)and(157.1,-163.8)..(158.1,-164.4)
				..controls(158.8,-164.8)and(159.7,-165.7)..(159.9,-166.3)
				..controls(160,-166.5)and(160.2,-167.1)..(160.2,-167.6)
				..controls(160.3,-168)and(160.6,-169.4)..(160.9,-170.6)
				..controls(161.2,-171.7)and(161.5,-173.1)..(161.6,-173.6)
				..controls(161.8,-174.4)and(161.8,-174.5)..(161.5,-174.9)
				..controls(161.2,-175.6)and(160.8,-175.9)..(159.9,-176)
				..controls(159.2,-176.1)and(154.1,-176.2)..(151.5,-176.2)
				..controls(148.1,-176.1)and(122.6,-176.1)..(121.6,-176.2)
				..controls(120.8,-176.2)and(119.8,-176.2)..(119.4,-176.1)--cycle
				(130.4,-160)..controls(129,-160)and(119.8,-160)..(110.1,-159.9)
				..controls(100.3,-159.8)and(88.5,-159.7)..(83.9,-159.6)
				..controls(77.8,-159.5)and(75.3,-159.5)..(74.5,-159.6)
				..controls(73.9,-159.7)and(72.5,-159.8)..(71.3,-159.8)
				..controls(69.6,-159.8)and(69,-159.8)..(68.5,-159.6)
				..controls(67.6,-159.3)and(66.8,-158.6)..(66.2,-157.6)
				..controls(66,-157.3)and(65.6,-156.5)..(65.3,-156.1)
				..controls(64.6,-155.1)and(64.5,-154.9)..(63.7,-152.9)
				..controls(63.2,-151.6)and(62.9,-150.9)..(61.8,-149.6)
				..controls(61.5,-149.2)and(61.1,-148.5)..(60.9,-147.9)
				..controls(60.7,-147.4)and(60.3,-146.6)..(60.1,-146.1)
				..controls(59.9,-145.5)and(59.5,-144.9)..(59,-144.3)
				..controls(58.5,-143.7)and(58.2,-143.2)..(57.9,-142.5)
				..controls(57.7,-141.9)and(57.3,-141.1)..(57,-140.7)
				..controls(56.6,-140.3)and(56.1,-139.3)..(55.8,-138.5)
				..controls(55.5,-137.8)and(55.1,-136.8)..(54.9,-136.5)
				..controls(54.7,-136.1)and(54.3,-135.2)..(53.9,-134.5)
				..controls(53.6,-133.8)and(53,-132.7)..(52.6,-132)
				..controls(51.6,-130.5)and(51.4,-130.2)..(51.1,-129)
				..controls(50.6,-127.5)and(50.6,-127.4)..(49.7,-126.3)
				..controls(49.2,-125.7)and(48.6,-124.8)..(48.3,-124.2)
				..controls(48,-123.6)and(47.5,-122.7)..(47.3,-122.2)
				..controls(47,-121.8)and(46.5,-120.7)..(46.1,-119.8)
				..controls(45.8,-118.9)and(45.2,-117.8)..(44.9,-117.2)
				..controls(44.2,-116.1)and(43.3,-114.4)..(42.3,-112.2)
				..controls(41.8,-111.3)and(41.3,-110.4)..(41.1,-110.1)
				..controls(40.5,-109.4)and(40.3,-109)..(40,-107.8)
				..controls(39.8,-107)and(39.6,-106.5)..(39.2,-106)
				..controls(38.9,-105.6)and(38.4,-104.7)..(38,-104)
				..controls(37.7,-103.4)and(37.2,-102.5)..(37,-102.1)
				..controls(36.7,-101.7)and(36.3,-100.8)..(36,-100.1)
				..controls(35.6,-98.9)and(35.3,-98.4)..(34.6,-97.6)
				..controls(34.2,-97.1)and(34.2,-96.3)..(34.6,-96)
				..controls(35,-95.7)and(35.5,-95.9)..(36.2,-96.7)
				..controls(36.5,-97)and(37.2,-97.5)..(37.6,-97.7)..controls(38,-98)and(38.6,-98.4)..(39.1,-98.9)
				..controls(39.6,-99.4)and(40.3,-100)..(40.6,-100.3)
				..controls(41.3,-100.8)and(42.4,-101.8)..(43.7,-102.9)
				..controls(44,-103.2)and(44.6,-103.7)..(45,-103.9)
				..controls(45.4,-104.2)and(45.9,-104.6)..(46.1,-104.8)
				..controls(47.1,-105.9)and(48.8,-107.1)..(50.5,-107.9)
				..controls(51.8,-108.6)and(52.3,-109)..(53.4,-110)
				..controls(54.8,-111.2)and(55.6,-111.6)..(57,-112.1)
				..controls(57.5,-112.3)and(58.5,-112.8)..(59.4,-113.3)
				..controls(60.2,-113.7)and(61.3,-114.3)..(61.9,-114.5)
				..controls(62.4,-114.7)and(63.3,-115.1)..(63.8,-115.3)
				..controls(64.4,-115.6)and(65.4,-115.9)..(66.2,-116.1)
				..controls(67,-116.3)and(68.4,-116.6)..(69.2,-116.9)
				..controls(72.3,-117.9)and(72.8,-118)..(78,-118.1)
				..controls(84,-118.3)and(86.1,-118.1)..(88.1,-117.1)
				..controls(89.4,-116.5)and(93,-115)..(93.7,-114.9)
				..controls(94.8,-114.6)and(95.4,-114.2)..(95.9,-113.4)
				..controls(96.2,-113)and(96.7,-112.5)..(97.1,-112.2)
				..controls(98,-111.6)and(101.3,-108.4)..(102.1,-107.2)
				..controls(102.4,-106.8)and(102.9,-106.2)..(103.2,-105.8)
				..controls(103.5,-105.5)and(104,-104.8)..(104.2,-104.2)
				..controls(104.5,-103.6)and(105,-102.8)..(105.3,-102.4)
				..controls(105.6,-102)and(106.1,-101.2)..(106.5,-100.6)
				..controls(107.1,-99.5)and(108.6,-97.5)..(109.5,-96.5)
				..controls(109.9,-96.1)and(110.1,-96)..(110.4,-96)
				..controls(111.1,-96)and(111.6,-96.5)..(111.8,-97.6)
				..controls(112,-98.6)and(112.3,-99.3)..(112.8,-99.9)
				..controls(113.7,-100.9)and(115.2,-103.1)..(115.7,-104.2)
				..controls(116.1,-105)and(116.5,-105.5)..(116.9,-106)
				..controls(117.3,-106.4)and(117.8,-106.9)..(118,-107.3)
				..controls(118.2,-107.6)and(118.6,-108)..(119,-108.3)
				..controls(119.6,-108.8)and(121.1,-110.3)..(121.8,-111.2)
				..controls(122.1,-111.7)and(122.5,-112)..(123.4,-112.4)
				..controls(124.8,-113.2)and(125.8,-113.8)..(127,-114.6)
				..controls(127.5,-114.9)and(128.1,-115.3)..(128.2,-115.4)
				..controls(128.4,-115.5)and(129.3,-115.8)..(130.2,-116)
				..controls(131.1,-116.2)and(132.5,-116.6)..(133.3,-116.9)
				..controls(135.7,-117.9)and(136.1,-117.9)..(142.5,-117.8)
				..controls(145.8,-117.8)and(148.4,-117.7)..(148.8,-117.6)
				..controls(149.2,-117.6)and(150.2,-117.2)..(151,-116.8)
				..controls(152,-116.4)and(152.8,-116)..(153.6,-115.9)
				..controls(154.5,-115.7)and(155.2,-115.4)..(156.3,-114.9)
				..controls(157.1,-114.5)and(158.2,-114)..(158.7,-113.8)
				..controls(159.1,-113.7)and(159.8,-113.3)..(160.2,-113)
				..controls(160.6,-112.7)and(161.3,-112.4)..(161.7,-112.2)
				..controls(162.7,-111.9)and(163.2,-111.7)..(164.1,-111.1)
				..controls(164.5,-110.8)and(165.3,-110.4)..(165.9,-110)
				..controls(166.5,-109.7)and(167.3,-109.2)..(167.6,-108.9)
				..controls(167.9,-108.6)and(168.6,-108.2)..(169,-108)
				..controls(169.4,-107.8)and(170.2,-107.3)..(170.7,-106.9)
				..controls(171.2,-106.5)and(172.1,-105.8)..(172.7,-105.4)
				..controls(173.2,-105)and(174.1,-104.4)..(174.5,-104)
				..controls(174.9,-103.7)and(175.7,-103.1)..(176.2,-102.6)
				..controls(176.7,-102.2)and(177.4,-101.5)..(177.9,-101.2)
				..controls(178.3,-100.8)and(178.8,-100.4)..(179,-100.2)
				..controls(179.2,-100)and(179.6,-99.6)..(180,-99.3)
				..controls(180.4,-99.1)and(180.9,-98.6)..(181.2,-98.4)
				..controls(181.5,-98.1)and(182.1,-97.7)..(182.5,-97.5)
				..controls(183.2,-97.2)and(183.5,-96.9)..(184.3,-95.9)
				..controls(184.6,-95.6)and(184.6,-95.6)..(184.8,-95.7)
				..controls(185,-95.8)and(185.1,-96)..(185.1,-96.2)
				..controls(185.1,-96.5)and(185,-96.7)..(184.6,-97.2)
				..controls(184.1,-97.8)and(184,-97.9)..(183.9,-98.8)
				..controls(183.7,-100.3)and(183.3,-101.3)..(182.9,-101.8)
				..controls(182.6,-102.1)and(182.3,-102.6)..(182.2,-102.9)
				..controls(182,-103.3)and(181.6,-104)..(181.2,-104.6)
				..controls(180.9,-105.2)and(180.5,-106.1)..(180.4,-106.6)
				..controls(179.9,-108)and(179.7,-108.5)..(179,-109.4)
				..controls(178.6,-109.8)and(178.2,-110.5)..(177.9,-111)
				..controls(177.7,-111.5)and(177.3,-112.3)..(177,-112.7)
				..controls(176.6,-113.2)and(176.4,-113.7)..(176.3,-114.2)
				..controls(176.2,-114.7)and(175.9,-115.2)..(175.4,-116)
				..controls(175,-116.6)and(174.4,-117.7)..(174,-118.5)
				..controls(173.7,-119.2)and(173.3,-120.1)..(173,-120.5)
				..controls(172.8,-120.9)and(172.4,-121.8)..(172.1,-122.5)
				..controls(171.8,-123.2)and(171.3,-124.1)..(171,-124.6)
				..controls(170.8,-125.1)and(170.3,-126)..(169.9,-126.6)
				..controls(169.6,-127.3)and(169.2,-128.1)..(169,-128.3)
				..controls(168.8,-128.6)and(168.5,-129.3)..(168.3,-129.8)
				..controls(168.1,-130.3)and(167.8,-131.1)..(167.5,-131.6)
				..controls(166.9,-132.5)and(164.2,-137.9)..(164.2,-138.1)
				..controls(164.2,-138.1)and(163.9,-138.8)..(163.4,-139.7)
				..controls(162.9,-140.5)and(162.4,-141.5)..(162.3,-141.9)
				..controls(161.9,-143)and(161.6,-143.4)..(160.9,-144.1)
				..controls(160.4,-144.6)and(160.3,-144.8)..(160.1,-145.4)
				..controls(160,-145.8)and(159.6,-146.8)..(159.2,-147.8)
				..controls(158.7,-148.8)and(158.1,-150)..(157.9,-150.6)
				..controls(157.6,-151.1)and(157.2,-151.9)..(156.9,-152.4)
				..controls(156.5,-152.8)and(156.1,-153.7)..(155.8,-154.4)
				..controls(155.6,-155)and(155.2,-155.9)..(154.9,-156.3)
				..controls(154.7,-156.7)and(154.4,-157.4)..(154.2,-157.9)
				..controls(154.1,-158.4)and(153.8,-159)..(153.5,-159.3)
				..controls(152.8,-160.1)and(153,-160.1)..(142.4,-160.1)
				..controls(137.2,-160.1)and(131.8,-160.1)..(130.4,-160)--cycle
				(141.5,-103.7)..controls(140.8,-103)and(139.8,-102.1)..(139.2,-101.6)
				..controls(138.5,-101.2)and(137.9,-100.6)..(137.8,-100.4)
				..controls(137,-99.5)and(134,-96.5)..(133.2,-96)
				..controls(132.7,-95.7)and(131.8,-94.9)..(131.1,-94.3)
				..controls(129.4,-92.7)and(127.9,-91.6)..(126.6,-91.1)
				..controls(125.2,-90.5)and(124.1,-90)..(122.9,-89.1)
				..controls(122.1,-88.6)and(121.3,-88.2)..(119.9,-87.7)
				..controls(118.8,-87.3)and(117.5,-86.8)..(116.9,-86.5)
				..controls(114.3,-85.5)and(113.7,-85.4)..(109.7,-85.5)
				..controls(105.3,-85.7)and(103.9,-85.9)..(101.6,-87)
				..controls(100.7,-87.4)and(99.4,-87.8)..(97.9,-88)
				..controls(97,-88.2)and(96.6,-88.3)..(95.8,-88.9)
				..controls(95.4,-89.2)and(94.7,-89.6)..(94.2,-89.8)
				..controls(92.9,-90.3)and(92.3,-90.7)..(90.9,-92.1)
				..controls(90.2,-92.7)and(89.4,-93.3)..(89,-93.5)..controls(87,-94.5)and(85.2,-96)..(83.6,-97.9)
				..controls(83.3,-98.2)and(82.5,-98.9)..(81.8,-99.4)
				..controls(80.7,-100.1)and(80.5,-100.3)..(80,-101.1)
				..controls(78.7,-103.1)and(77.2,-103.9)..(76.4,-103)
				..controls(75.7,-102.2)and(75.7,-103.2)..(75.7,-83.6)
				..controls(75.6,-73.6)and(75.6,-64.8)..(75.5,-64)
				..controls(75.5,-63.2)and(75.5,-61.4)..(75.6,-59.8)--(75.7,-57.1)--(76,-56.8)
				..controls(76.8,-56.2)and(78,-56.3)..(79.8,-57.2)
				..controls(80.4,-57.5)and(81.3,-57.8)..(81.8,-57.9)
				..controls(82.3,-58.1)and(83,-58.4)..(83.6,-58.8)..controls(84.3,-59.2)and(85,-59.5)..(86,-59.8)
				..controls(86.8,-60.1)and(88,-60.6)..(88.7,-60.9)
				..controls(89.3,-61.2)and(90.6,-61.7)..(91.4,-62)
				..controls(92.2,-62.3)and(93.7,-62.9)..(94.6,-63.3)
				..controls(95.6,-63.7)and(96.8,-64.1)..(97.4,-64.3)
				..controls(98,-64.5)and(99,-64.9)..(99.7,-65.2)
				..controls(100.4,-65.6)and(101.5,-66)..(102.3,-66.2)
				..controls(104.2,-66.8)and(105,-67.1)..(106.2,-67.8)
				..controls(107.9,-68.7)and(110.2,-69.4)..(111.3,-69.2)
				..controls(112.5,-69)and(115,-68.2)..(115.6,-67.8)
				..controls(117.3,-66.7)and(117.6,-66.6)..(119.5,-66.2)
				..controls(121.9,-65.8)and(122.1,-65.7)..(123,-65.1)
				..controls(124.1,-64.4)and(124.4,-64.3)..(126,-63.9)
				..controls(127,-63.7)and(127.7,-63.4)..(128.4,-63.1)
				..controls(129,-62.8)and(129.7,-62.5)..(130.6,-62.4)
				..controls(132.1,-62.1)and(134,-61.5)..(135.2,-61)
				..controls(135.7,-60.8)and(136.7,-60.5)..(137.3,-60.3)
				..controls(138.4,-60)and(139.3,-59.7)..(141.1,-59)
				..controls(141.6,-58.7)and(142,-58.6)..(142.5,-58.6)
				..controls(143.2,-58.6)and(143.3,-58.7)..(143.6,-59)
				..controls(144,-59.5)and(144.1,-59.9)..(144,-60.8)
				..controls(144,-61.2)and(144,-71.1)..(144,-82.8)--(143.9,-104.1)--(143.7,-104.5)
				..controls(143.5,-104.7)and(143.3,-104.8)..(143,-104.9)
				..controls(142.6,-104.9)and(142.5,-104.8)..(141.5,-103.7)--cycle
				(27.8,-91.4)..controls(27.1,-90.9)and(26.3,-90)..(26.1,-89.5)
				..controls(26,-89.2)and(25.8,-88.6)..(25.7,-87.9)
				..controls(25.4,-86.2)and(25.6,-84.8)..(26.5,-83.5)
				..controls(27.4,-82.3)and(28.3,-82)..(29.4,-82.3)
				..controls(31.1,-82.8)and(32.1,-84.9)..(32,-87.4)
				..controls(31.9,-89.3)and(31.2,-90.6)..(29.9,-91.4)
				..controls(29.1,-91.9)and(28.6,-91.9)..(27.8,-91.4)--cycle
				(190.1,-91.3)..controls(189.7,-91.2)and(189.3,-90.9)..(189.1,-90.7)
				..controls(188.3,-89.8)and(187.6,-87.7)..(187.6,-86.2)
				..controls(187.6,-85.5)and(187.7,-85.2)..(188.2,-84.4)
				..controls(189.1,-82.6)and(190.1,-81.8)..(191.3,-82.1)
				..controls(192,-82.3)and(193.5,-83.7)..(193.8,-84.4)
				..controls(194.1,-85.1)and(194.2,-87.1)..(194,-88.3)
				..controls(193.8,-89.8)and(193.3,-90.5)..(192.3,-91.3)
				..controls(191.7,-91.7)and(190.9,-91.7)..(190.1,-91.3)--cycle
				(173.7,-83.4)..controls(172.2,-82.5)and(172,-82.5)..(170.6,-82.4)
				..controls(169.7,-82.4)and(168.8,-82.4)..(168.3,-82.5)
				..controls(167.2,-82.7)and(166.9,-82.7)..(166.6,-82.2)
				..controls(166.3,-81.8)and(166.3,-81)..(166.8,-80)
				..controls(167.2,-78.9)and(167.4,-77.4)..(167.7,-74.1)
				..controls(167.8,-72.8)and(167.9,-71.8)..(167.9,-71.7)
				..controls(168,-71.7)and(168,-70.2)..(168,-68.4)..controls(168,-65.4)and(168,-64.9)..(167.7,-63)
				..controls(167.3,-60.5)and(167.1,-59.9)..(166.4,-58.4)
				..controls(166.1,-57.7)and(165.8,-56.9)..(165.7,-56.6)
				..controls(165.7,-56.2)and(165.4,-55.6)..(165.2,-55.2)
				..controls(164.6,-53.8)and(161.9,-50.7)..(161,-50.3)
				..controls(159.8,-49.7)and(157.9,-50)..(155.2,-51.2)
				..controls(154.7,-51.4)and(153.4,-51.8)..(152.2,-52.1)
				..controls(150.8,-52.4)and(149.4,-52.9)..(148.6,-53.2)
				..controls(147.8,-53.6)and(146.7,-54)..(146,-54.2)
				..controls(145.4,-54.4)and(144.5,-54.7)..(144,-54.9)
				..controls(143.6,-55.2)and(142.7,-55.5)..(141.8,-55.7)
				..controls(141,-55.9)and(139.6,-56.3)..(138.7,-56.6)
				..controls(137.8,-57)and(136.3,-57.4)..(135.3,-57.6)
				..controls(133.8,-57.9)and(133.2,-58.1)..(132.1,-58.7)
				..controls(131.1,-59.2)and(130.3,-59.4)..(129.4,-59.6)
				..controls(128.7,-59.8)and(127.8,-60)..(127.4,-60.1)
				..controls(127,-60.3)and(126,-60.5)..(125.1,-60.7)
				..controls(123.2,-61.1)and(121.7,-61.5)..(120.2,-62.1)
				..controls(119.7,-62.3)and(118.1,-62.8)..(116.7,-63.2)
				..controls(114.6,-63.9)and(114,-64.1)..(113,-64.7)
				..controls(111.9,-65.4)and(111.8,-65.4)..(111,-65.5)
				..controls(109.9,-65.5)and(109.3,-65.4)..(107.7,-64.7)
				..controls(107,-64.4)and(106,-64)..(105.3,-63.9)
				..controls(103.5,-63.5)and(102.2,-63)..(101.2,-62.5)
				..controls(100.1,-61.8)and(99.3,-61.5)..(97.8,-61.1)
				..controls(97.2,-61)and(96.2,-60.7)..(95.7,-60.4)
				..controls(95.1,-60.1)and(94.2,-59.8)..(93.5,-59.6)
				..controls(92.9,-59.4)and(91.7,-59)..(91,-58.7)..controls(90.2,-58.4)and(88.9,-58)..(88,-57.7)
				..controls(87.2,-57.5)and(86.1,-57.1)..(85.5,-56.8)
				..controls(85,-56.6)and(84.1,-56.2)..(83.5,-56)..controls(82.9,-55.7)and(81.8,-55.2)..(81,-54.9)
				..controls(80.1,-54.5)and(79.1,-54.2)..(78.7,-54.1)
				..controls(77.9,-53.8)and(77.3,-53.6)..(76.2,-52.9)
				..controls(75.7,-52.6)and(75.2,-52.4)..(74.1,-52.2)
				..controls(72.5,-51.9)and(72.4,-51.8)..(71.3,-51.1)
				..controls(69.9,-50.1)and(69.4,-49.9)..(67.9,-49.8)
				..controls(67,-49.7)and(66.6,-49.6)..(66.3,-49.4)
				..controls(66.1,-49.3)and(65.7,-49)..(65.4,-48.8)
				..controls(65.1,-48.7)and(64.2,-48.4)..(63.6,-48.2)
				..controls(62.9,-48)and(61.7,-47.6)..(60.9,-47.2)
				..controls(60.1,-46.8)and(58.9,-46.4)..(58.2,-46.2)
				..controls(57.6,-46)and(56.8,-45.7)..(56.5,-45.5)
				..controls(54.9,-44.6)and(54.6,-44.5)..(53.2,-44.2)
				..controls(51.7,-43.8)and(51,-43.6)..(49.9,-42.9)
				..controls(49.5,-42.6)and(48.9,-42.4)..(47.5,-42)
				..controls(45.7,-41.6)and(45.6,-41.5)..(44.3,-40.6)
				..controls(43.5,-40.1)and(42.8,-39.6)..(42.7,-39.4)
				..controls(42.4,-38.7)and(43,-38.1)..(44.5,-37.6)..controls(44.9,-37.5)and(46,-37)..(46.9,-36.6)
				..controls(47.9,-36.1)and(48.8,-35.8)..(49.6,-35.6)
				..controls(50.5,-35.4)and(51.3,-35.2)..(52.4,-34.6)
				..controls(53.3,-34.1)and(54.5,-33.7)..(54.9,-33.5)
				..controls(55.4,-33.4)and(56.4,-33)..(57.2,-32.6)
				..controls(58.6,-32)and(60.1,-31.5)..(61.6,-31.2)
				..controls(62.3,-31.1)and(62.7,-31)..(63.2,-30.7)
				..controls(64.1,-30)and(64.3,-29.9)..(65.3,-29.7)
				..controls(65.7,-29.6)and(66.5,-29.2)..(67.1,-28.9)
				..controls(67.6,-28.6)and(68.8,-28.2)..(69.6,-27.9)
				..controls(70.5,-27.7)and(71.6,-27.3)..(72.1,-27)
				..controls(72.6,-26.7)and(73.8,-26.2)..(74.8,-25.9)
				..controls(75.7,-25.5)and(76.7,-25.1)..(77.1,-25)
				..controls(77.4,-24.8)and(78.5,-24.4)..(79.6,-24.1)
				..controls(80.6,-23.8)and(82,-23.2)..(82.8,-22.9)
				..controls(83.5,-22.5)and(84.8,-22)..(85.8,-21.8)
				..controls(86.7,-21.5)and(87.8,-21.1)..(88.1,-20.9)
				..controls(88.4,-20.7)and(89.3,-20.3)..(90.1,-20.1)
				..controls(90.8,-19.8)and(91.9,-19.4)..(92.4,-19.1)
				..controls(93.5,-18.5)and(94.4,-18.2)..(95.4,-18.1)
				..controls(95.7,-18)and(96.4,-17.8)..(96.9,-17.5)
				..controls(97.5,-17.3)and(98.4,-16.9)..(99,-16.6)
				..controls(99.7,-16.3)and(100.9,-15.9)..(101.9,-15.6)
				..controls(102.9,-15.3)and(104.4,-14.7)..(105.3,-14.3)
				..controls(108.1,-13.2)and(108.5,-13.1)..(110.1,-13.1)
				..controls(111.4,-13.1)and(111.5,-13.1)..(112.4,-13.5)
				..controls(113.2,-14)and(115,-14.8)..(116.6,-15.4)
				..controls(116.9,-15.5)and(117.9,-15.8)..(118.8,-16.1)
				..controls(120.9,-16.6)and(121.5,-16.9)..(122.6,-17.9)
				..controls(123,-18.3)and(123.7,-18.8)..(123.9,-18.9)
				..controls(124.6,-19.4)and(126,-19.8)..(127.9,-20.2)
				..controls(129.6,-20.5)and(129.8,-20.6)..(131.1,-21.4)
				..controls(131.5,-21.7)and(132.4,-22.1)..(133.1,-22.3)
				..controls(136.1,-23.2)and(138.3,-24.1)..(140.2,-25)
				..controls(141,-25.5)and(141.5,-25.6)..(142.5,-25.8)
				..controls(145.1,-26.3)and(146,-26.6)..(149.4,-28.1)
				..controls(151.3,-28.9)and(153.5,-29.8)..(154.3,-30.1)
				..controls(155.1,-30.4)and(156.2,-30.8)..(156.6,-31)
				..controls(157.1,-31.2)and(158.2,-31.6)..(159.1,-31.8)
				..controls(160.6,-32.2)and(160.9,-32.3)..(161.7,-32.9)
				..controls(162.4,-33.3)and(162.9,-33.6)..(163.7,-33.8)
				..controls(164.3,-34)and(164.8,-34.1)..(164.9,-34.1)
				..controls(165,-34.1)and(165.6,-34.4)..(166.4,-34.8)
				..controls(167.3,-35.2)and(168.1,-35.4)..(169,-35.6)
				..controls(170,-35.8)and(170.5,-36)..(171.7,-36.6)
				..controls(172.4,-37)and(173.3,-37.4)..(173.6,-37.4)
				..controls(175.6,-37.7)and(176.2,-38)..(176.7,-39.1)
				..controls(176.9,-39.6)and(176.9,-39.7)..(176.7,-40.4)
				..controls(176.4,-41.6)and(175.8,-42.7)..(175.2,-43.4)
				..controls(174.5,-44.1)and(172.6,-45.5)..(171.7,-45.8)
				..controls(170.6,-46.2)and(169.1,-46)..(168.2,-45.4)
				..controls(167.9,-45.2)and(167.3,-44.8)..(166.9,-44.4)
				..controls(166.4,-44.1)and(165.8,-43.6)..(165.4,-43.2)
				..controls(165.1,-42.9)and(164.4,-42.3)..(163.9,-41.9)
				..controls(163.3,-41.5)and(162.7,-40.9)..(162.4,-40.5)
				..controls(161.4,-39.2)and(159.9,-38)..(158.8,-37.8)
				..controls(158.4,-37.8)and(158.3,-37.7)..(158.2,-37.2)
				..controls(158,-36.8)and(157.4,-36.1)..(157.1,-36.1)
				..controls(157,-36.1)and(156.7,-36.1)..(156.5,-36)
				..controls(156.1,-35.9)and(155.9,-35.7)..(155.7,-35.2)
				..controls(155.3,-34.6)and(155.1,-34.4)..(154,-34.1)
				..controls(153.6,-34)and(153.2,-33.7)..(152.9,-33.5)
				..controls(152.4,-32.9)and(151.6,-32.4)..(151.3,-32.4)
				..controls(150.9,-32.4)and(150.4,-32.6)..(150.4,-32.9)
				..controls(150.4,-33.2)and(150.8,-33.4)..(151.4,-33.6)
				..controls(152.1,-33.7)and(152.8,-34.1)..(153.5,-34.9)
				..controls(154.1,-35.5)and(154.3,-35.6)..(155,-35.8)
				..controls(155.7,-36)and(155.7,-36)..(156,-36.6)
				..controls(156.3,-37.3)and(156.7,-37.7)..(157.5,-37.8)
				..controls(158.1,-37.9)and(158.2,-38)..(158.4,-38.8)
				..controls(158.5,-39.4)and(158.6,-39.6)..(159.6,-40.4)
				..controls(160.2,-41)and(160.9,-41.6)..(161.3,-41.8)
				..controls(161.7,-42)and(161.9,-42.3)..(162.2,-42.8)
				..controls(162.4,-43.2)and(162.8,-43.8)..(163.1,-44.1)
				..controls(163.9,-44.8)and(165.1,-46.7)..(165.7,-48)
				..controls(166.2,-49.3)and(166.4,-49.7)..(167.2,-50.5)
				..controls(167.5,-50.8)and(167.9,-51.4)..(168.1,-51.8)
				..controls(168.3,-52.2)and(168.8,-53)..(169.1,-53.6)
				..controls(169.6,-54.3)and(169.8,-54.9)..(169.9,-55.6)
				..controls(170.4,-57.8)and(170.5,-58.1)..(171.2,-59.5)
				..controls(171.6,-60.3)and(172,-61.5)..(172.2,-62.1)
				..controls(172.3,-62.8)and(172.7,-63.7)..(172.9,-64.1)
				..controls(173.2,-64.6)and(173.4,-65.3)..(173.4,-65.6)
				..controls(173.8,-68.8)and(174.2,-70.6)..(174.9,-72.6)
				..controls(175.5,-74.5)and(175.9,-77.1)..(176.1,-80.2)
				..controls(176.3,-83.4)and(175.5,-84.4)..(173.7,-83.4)--cycle;
				\definecolor{ezs}{RGB}{247,179,26}\fill[ezs](-5.7,-156.6)--(-3.4,-157.5)--(61.9,12.7)--(59.7,13.6)--cycle;
				\definecolor{ezs}{RGB}{247,179,22}\fill[ezs](-3.8,-157.3)--(-1.5,-158.2)--(63.9,12)--(61.6,12.9)--cycle;
				\definecolor{ezs}{RGB}{248,179,19}\fill[ezs](-1.8,-158.1)--(0.4,-158.9)--(65.8,11.2)--(63.5,12.1)--cycle;
				\definecolor{ezs}{RGB}{248,178,16}\fill[ezs](0.1,-158.8)--(2.4,-159.7)--(67.7,10.5)--(65.4,11.4)--cycle;
				\definecolor{ezs}{RGB}{249,178,15}\fill[ezs](2,-159.5)--(4.3,-160.4)--(69.6,9.8)--(67.3,10.7)--cycle;
				\definecolor{ezs}{RGB}{251,177,17}\fill[ezs](3.9,-160.3)--(6.2,-161.1)--(71.5,9)--(69.2,9.9)--cycle;
				\definecolor{ezs}{RGB}{252,179,18}\fill[ezs](5.8,-161)--(8.1,-161.9)--(73.4,8.3)--(71.1,9.2)--cycle;
				\definecolor{ezs}{RGB}{252,182,17}\fill[ezs](7.7,-161.7)--(10,-162.6)--(75.3,7.6)--(73,8.5)--cycle;
				\definecolor{ezs}{RGB}{252,185,17}\fill[ezs](9.6,-162.5)--(11.9,-163.4)--(77.2,6.8)--(75,7.7)--cycle;
				\definecolor{ezs}{RGB}{251,189,18}\fill[ezs](11.5,-163.2)--(13.8,-164.1)--(79.2,6.1)--(76.9,7)--cycle;
				\definecolor{ezs}{RGB}{251,191,22}\fill[ezs](13.4,-163.9)--(15.7,-164.8)--(81.1,5.4)--(78.8,6.2)--cycle;
				\definecolor{ezs}{RGB}{252,193,25}\fill[ezs](15.4,-164.7)--(17.6,-165.6)--(83,4.6)--(80.7,5.5)--cycle;
				\definecolor{ezs}{RGB}{253,195,27}\fill[ezs](17.3,-165.4)--(19.6,-166.3)--(84.9,3.9)--(82.6,4.8)--cycle;
				\definecolor{ezs}{RGB}{253,196,31}\fill[ezs](19.2,-166.1)--(21.5,-167)--(86.8,3.2)--(84.5,4)--cycle;
				\definecolor{ezs}{RGB}{254,199,35}\fill[ezs](21.1,-166.9)--(23.4,-167.8)--(88.7,2.4)--(86.4,3.3)--cycle;
				\definecolor{ezs}{RGB}{253,201,38}\fill[ezs](23,-167.6)--(25.3,-168.5)--(90.6,1.7)--(88.3,2.6)--cycle;
				\definecolor{ezs}{RGB}{253,203,40}\fill[ezs](24.9,-168.3)--(27.2,-169.2)--(92.5,1)--(90.3,1.8)--cycle;
				\definecolor{ezs}{RGB}{253,205,42}\fill[ezs](26.8,-169.1)--(29.1,-170)--(94.4,0.2)--(92.2,1.1)--cycle;
				\definecolor{ezs}{RGB}{253,207,45}\fill[ezs](28.7,-169.8)--(31,-170.7)--(96.4,-0.5)--(94.1,0.4)--cycle;
				\definecolor{ezs}{RGB}{253,210,49}\fill[ezs](30.7,-170.5)--(32.9,-171.4)--(98.3,-1.2)--(96,-0.4)--cycle;
				\definecolor{ezs}{RGB}{253,212,54}\fill[ezs](32.6,-171.3)--(34.9,-172.2)--(100.2,-2)--(97.9,-1.1)--cycle;
				\definecolor{ezs}{RGB}{253,214,59}\fill[ezs](34.5,-172)--(36.8,-172.9)--(102.1,-2.7)--(99.8,-1.8)--cycle;
				\definecolor{ezs}{RGB}{253,216,63}\fill[ezs](36.4,-172.7)--(38.7,-173.6)--(104,-3.4)--(101.7,-2.6)--cycle;
				\definecolor{ezs}{RGB}{253,218,65}\fill[ezs](38.3,-173.5)--(40.6,-174.4)--(105.9,-4.2)--(103.6,-3.3)--cycle;
				\definecolor{ezs}{RGB}{254,220,66}\fill[ezs](40.2,-174.2)--(42.5,-175.1)--(107.8,-4.9)--(105.5,-4)--cycle;
				\definecolor{ezs}{RGB}{254,222,68}\fill[ezs](42.1,-175)--(44.4,-175.8)--(109.7,-5.6)--(107.5,-4.8)--cycle;
				\definecolor{ezs}{RGB}{254,224,71}\fill[ezs](44,-175.7)--(46.3,-176.6)--(111.7,-6.4)--(109.4,-5.5)--cycle;
				\definecolor{ezs}{RGB}{254,226,75}\fill[ezs](46,-176.4)--(48.2,-177.3)--(113.6,-7.1)--(111.3,-6.2)--cycle;
				\definecolor{ezs}{RGB}{254,228,79}\fill[ezs](47.9,-177.2)--(50.1,-178)--(115.5,-7.8)--(113.2,-7)--cycle;
				\definecolor{ezs}{RGB}{254,230,82}\fill[ezs](49.8,-177.9)--(52.1,-178.8)--(117.4,-8.6)--(115.1,-7.7)--cycle;
				\definecolor{ezs}{RGB}{254,232,85}\fill[ezs](51.7,-178.6)--(54,-179.5)--(119.3,-9.3)--(117,-8.4)--cycle;
				\definecolor{ezs}{RGB}{254,234,88}\fill[ezs](53.6,-179.4)--(55.9,-180.2)--(121.2,-10)--(118.9,-9.2)--cycle;
				\definecolor{ezs}{RGB}{254,237,91}\fill[ezs](55.5,-180.1)--(57.8,-181)--(123.1,-10.8)--(120.8,-9.9)--cycle;
				\definecolor{ezs}{RGB}{254,239,94}\fill[ezs](57.4,-180.8)--(59.7,-181.7)--(125,-11.5)--(122.8,-10.6)--cycle;
				\definecolor{ezs}{RGB}{254,241,97}\fill[ezs](59.3,-181.6)--(61.6,-182.4)--(127,-12.2)--(124.7,-11.4)--cycle;
				\definecolor{ezs}{RGB}{254,243,101}\fill[ezs](61.2,-182.3)--(63.5,-183.2)--(128.9,-13)--(126.6,-12.1)--cycle;
				\definecolor{ezs}{RGB}{254,245,103}\fill[ezs](63.2,-183)--(65.4,-183.9)--(130.8,-13.7)--(128.5,-12.8)--cycle;
				\definecolor{ezs}{RGB}{254,245,105}\fill[ezs](65.1,-183.8)--(67.4,-184.6)--(132.7,-14.4)--(130.4,-13.6)--cycle;
				\definecolor{ezs}{RGB}{254,245,105}\fill[ezs](67,-184.5)--(69.3,-185.4)--(134.6,-15.2)--(132.3,-14.3)--cycle;
				\definecolor{ezs}{RGB}{254,244,103}\fill[ezs](68.9,-185.2)--(71.2,-186.1)--(136.5,-15.9)--(134.2,-15)--cycle;
				\definecolor{ezs}{RGB}{253,243,101}\fill[ezs](70.8,-186)--(73.1,-186.8)--(138.4,-16.6)--(136.1,-15.8)--cycle;
				\definecolor{ezs}{RGB}{253,241,98}\fill[ezs](72.7,-186.7)--(75,-187.6)--(140.3,-17.4)--(138,-16.5)--cycle;
				\definecolor{ezs}{RGB}{253,239,95}\fill[ezs](74.6,-187.4)--(76.9,-188.3)--(142.2,-18.1)--(140,-17.2)--cycle;
				\definecolor{ezs}{RGB}{253,236,92}\fill[ezs](76.5,-188.2)--(78.8,-189)--(144.2,-18.8)--(141.9,-18)--cycle;
				\definecolor{ezs}{RGB}{253,234,88}\fill[ezs](78.5,-188.9)--(80.7,-189.8)--(146.1,-19.6)--(143.8,-18.7)--cycle;
				\definecolor{ezs}{RGB}{253,232,85}\fill[ezs](80.4,-189.6)--(82.7,-190.5)--(148,-20.3)--(145.7,-19.4)--cycle;
				\definecolor{ezs}{RGB}{253,230,82}\fill[ezs](82.3,-190.4)--(84.6,-191.2)--(149.9,-21)--(147.6,-20.2)--cycle;
				\definecolor{ezs}{RGB}{253,228,78}\fill[ezs](84.2,-191.1)--(86.5,-192)--(151.8,-21.8)--(149.5,-20.9)--cycle;
				\definecolor{ezs}{RGB}{253,226,74}\fill[ezs](86.1,-191.8)--(88.4,-192.7)--(153.7,-22.5)--(151.4,-21.6)--cycle;
				\definecolor{ezs}{RGB}{253,224,70}\fill[ezs](88,-192.6)--(90.3,-193.4)--(155.6,-23.3)--(153.3,-22.4)--cycle;
				\definecolor{ezs}{RGB}{253,222,67}\fill[ezs](89.9,-193.3)--(92.2,-194.2)--(157.5,-24)--(155.3,-23.1)--cycle;
				\definecolor{ezs}{RGB}{253,220,65}\fill[ezs](91.8,-194)--(94.1,-194.9)--(159.5,-24.7)--(157.2,-23.8)--cycle;
				\definecolor{ezs}{RGB}{253,217,63}\fill[ezs](93.7,-194.8)--(96,-195.6)--(161.4,-25.5)--(159.1,-24.6)--cycle;
				\definecolor{ezs}{RGB}{253,215,60}\fill[ezs](95.7,-195.5)--(97.9,-196.4)--(163.3,-26.2)--(161,-25.3)--cycle;
				\definecolor{ezs}{RGB}{253,213,56}\fill[ezs](97.6,-196.2)--(99.9,-197.1)--(165.2,-26.9)--(162.9,-26)--cycle;
				\definecolor{ezs}{RGB}{252,211,53}\fill[ezs](99.5,-197)--(101.8,-197.8)--(167.1,-27.7)--(164.8,-26.8)--cycle;
				\definecolor{ezs}{RGB}{252,209,49}\fill[ezs](101.4,-197.7)--(103.7,-198.6)--(169,-28.4)--(166.7,-27.5)--cycle;
				\definecolor{ezs}{RGB}{252,207,45}\fill[ezs](103.3,-198.4)--(105.6,-199.3)--(170.9,-29.1)--(168.6,-28.2)--cycle;
				\definecolor{ezs}{RGB}{253,205,42}\fill[ezs](105.2,-199.2)--(107.5,-200)--(172.8,-29.9)--(170.6,-29)--cycle;
				\definecolor{ezs}{RGB}{253,202,40}\fill[ezs](107.1,-199.9)--(109.4,-200.8)--(174.7,-30.6)--(172.5,-29.7)--cycle;
				\definecolor{ezs}{RGB}{253,200,38}\fill[ezs](109,-200.6)--(111.3,-201.5)--(176.7,-31.3)--(174.4,-30.4)--cycle;
				\definecolor{ezs}{RGB}{253,198,35}\fill[ezs](111,-201.4)--(113.2,-202.2)--(178.6,-32.1)--(176.3,-31.2)--cycle;
				\definecolor{ezs}{RGB}{253,196,31}\fill[ezs](112.9,-202.1)--(115.2,-203)--(180.5,-32.8)--(178.2,-31.9)--cycle;
				\definecolor{ezs}{RGB}{253,194,28}\fill[ezs](114.8,-202.8)--(117.1,-203.7)--(182.4,-33.5)--(180.1,-32.6)--cycle;
				\definecolor{ezs}{RGB}{253,192,26}\fill[ezs](116.7,-203.6)--(119,-204.5)--(184.3,-34.3)--(182,-33.4)--cycle;
				\definecolor{ezs}{RGB}{253,189,24}\fill[ezs](118.6,-204.3)--(120.9,-205.2)--(186.2,-35)--(183.9,-34.1)--cycle;
				\definecolor{ezs}{RGB}{253,186,20}\fill[ezs](120.5,-205)--(122.8,-205.9)--(188.1,-35.7)--(185.8,-34.8)--cycle;
				\definecolor{ezs}{RGB}{253,184,17}\fill[ezs](122.4,-205.8)--(124.7,-206.7)--(190,-36.5)--(187.8,-35.6)--cycle;
				\definecolor{ezs}{RGB}{252,183,15}\fill[ezs](124.3,-206.5)--(126.6,-207.4)--(192,-37.2)--(189.7,-36.3)--cycle;
				\definecolor{ezs}{RGB}{252,181,14}\fill[ezs](126.3,-207.2)--(128.5,-208.1)--(193.9,-37.9)--(191.6,-37.1)--cycle;
				\definecolor{ezs}{RGB}{252,179,13}\fill[ezs](128.2,-208)--(130.4,-208.9)--(195.8,-38.7)--(193.5,-37.8)--cycle;
				\definecolor{ezs}{RGB}{252,178,12}\fill[ezs](130.1,-208.7)--(132.4,-209.6)--(197.7,-39.4)--(195.4,-38.5)--cycle;
				\definecolor{ezs}{RGB}{251,178,10}\fill[ezs](132,-209.4)--(134.3,-210.3)--(199.6,-40.1)--(197.3,-39.3)--cycle;
				\definecolor{ezs}{RGB}{251,178,10}\fill[ezs](133.9,-210.2)--(136.2,-211.1)--(201.5,-40.9)--(199.2,-40)--cycle;
				\definecolor{ezs}{RGB}{250,177,11}\fill[ezs](135.8,-210.9)--(138.1,-211.8)--(203.4,-41.6)--(201.1,-40.7)--cycle;
				\definecolor{ezs}{RGB}{250,178,11}\fill[ezs](137.7,-211.6)--(140,-212.5)--(205.3,-42.3)--(203.1,-41.5)--cycle;
				\definecolor{ezs}{RGB}{248,180,10}\fill[ezs](139.6,-212.4)--(141.9,-213.3)--(207.3,-43.1)--(205,-42.2)--cycle;
				\definecolor{ezs}{RGB}{247,181,9}\fill[ezs](141.5,-213.1)--(143.8,-214)--(209.2,-43.8)--(206.9,-42.9)--cycle;
				\definecolor{ezs}{RGB}{246,182,9}\fill[ezs](143.5,-213.8)--(145.7,-214.7)--(211.1,-44.5)--(208.8,-43.7)--cycle;
				\definecolor{ezs}{RGB}{246,182,10}\fill[ezs](145.4,-214.6)--(147.7,-215.5)--(213,-45.3)--(210.7,-44.4)--cycle;
			\end{scope}
			% chữ EDUZY (Times đậm, làm dày nét bằng 5 lớp lệch nhẹ; căn đúng khung chữ logo gốc)
			\foreach \dx/\dy in {-0.7/0,0.7/0,0/-0.3,0/0.3,0/0}
			\node[anchor=center,inner sep=0pt,outer sep=0pt,text=ezText,align=center] at (109.6+\dx,-193.6+\dy)
			{\resizebox{1.01cm}{0.22cm}{\fontencoding{T1}\fontfamily{ptm}\fontseries{b}\selectfont EDUZY}};
		\end{tikzpicture}%
	}%
}
%====================== THIẾT LẬP THÂN ĐỀ ======================
%--- Môi trường tự luận {bt}: ex_test 3.3.4 chưa khai báo sẵn nên tạo ở đây
%    (ex_test đã tự móc \fixshowans và \listenumerate cho bt)
\makeatletter
\@ifundefined{bt}{\newtheorem{bt}{\color{gold1}\namebt}}{}
\makeatother
\renewcommand{\sectionmark}[1]{\markright{\thesection\ #1}}
%--- Khung tiêu đề (logo + tối đa 4 dòng; dòng 3, 4 để trống {} thì bỏ qua)
%  Đề thi:    \khungtieude{THI THỬ THPT QUỐC GIA NĂM 2026}{Môn Toán -- Đề số 21}
%                         {(Đề thi gồm: ...)}{Thời gian làm bài: 90 phút (không kể thời gian phát đề).}
%  Chuyên đề: \khungtieude{CHUYÊN ĐỀ 3 -- DÃY SỐ VÀ CÔNG THỨC TRUY HỒI}
%                         {Bồi dưỡng học sinh giỏi Toán THPT}{...}{}
\newcommand{\khungtieude}[4]{%
	\begin{tcolorbox}[colframe=gold1,colback=gold5!45,boxrule=0.5mm]
		\begin{minipage}[c]{3cm}
			\centerline{\eduzylogo[2.4cm]}
		\end{minipage}\hspace{0.5cm}
		\begin{minipage}[c]{12.5cm}
			\vskip 0.15 cm
			\centerline{\ind{#1}}
			\vskip 0.15 cm
			\centerline{\textbf{#2}}
			\if\relax\detokenize{#3}\relax\else
				\vskip 0.15 cm
				\centerline{\textit{#3}}
			\fi
			\if\relax\detokenize{#4}\relax\else
				\vskip 0.15 cm
				\centerline{\textit{#4}}
			\fi
			\vskip 0.15 cm
		\end{minipage}
	\end{tcolorbox}
}
%====================== BỔ SUNG CHO BÀI HỌC, CHUYÊN ĐỀ ======================
\usepackage[shortlabels]{enumitem}
\usepackage{tabularx}
%--- Danh sách gọn: \begin{enumerate}[gon] ... \end{enumerate}
\SetEnumitemKey{gon}{leftmargin=*,itemsep=1pt,parsep=0pt,topsep=3pt}
%--- Khung lý thuyết: \begin{khung}{} ... \end{khung}; tiêu đề chỉ dùng khi thật ngắn ("Mẹo", "Lỗi thường gặp")
\newtcolorbox{khung}[1]{breakable,enhanced,colframe=gold1,colback=gold5!45,
	boxrule=0.4mm,fonttitle=\bfseries\sffamily,coltitle=white,colbacktitle=gold1,
	attach boxed title to top left={xshift=6mm,yshift=-3mm},
	boxed title style={arc=1mm,boxrule=0.3mm},title={#1}}
%--- Tiêu đề mục nhỏ trong một dạng: \muc{Dấu hiệu}, \muc{Phương pháp}, \muc{Bài tự luyện}
\makeatletter
\newcommand{\muc}[1]{\par\medskip\noindent\indam{#1}\par\nobreak\smallskip\@afterindentfalse\@afterheading}
\makeatother
%--- Ô công thức cần nhớ (căn giữa, viền vàng, vừa khít nội dung)
%     \begin{center}\congthuc{$Q(n)=(n-1)!$}\end{center}
\newtcbox{\congthuc}{on line,colframe=gold1,colback=gold5!45,boxrule=0.5pt,arc=2pt,
	left=8pt,right=8pt,top=4pt,bottom=4pt}
%--- Công tắc lời giải bài tự luyện (ví dụ mẫu luôn dùng \loigiai)
%     \loigiaiTL[{đáp số}]{lời giải đầy đủ}
%     \TLfull -> in lời giải đầy đủ;  \TLdapso -> chỉ in một dòng "Đáp số";  \TLan -> ẩn hoàn toàn
%     (\CheDo là dethi thì mọi lời giải và đáp số đều ẩn)
\makeatletter
\newif\ifTLhien   \TLhientrue
\newif\ifTLdapso  \TLdapsofalse
\newcommand{\TLfull}{\TLhientrue\TLdapsofalse}
\newcommand{\TLdapso}{\TLhientrue\TLdapsotrue}
\newcommand{\TLan}{\TLhienfalse}
\newcommand{\loigiaiTL}[2][]{%
	\ifTLhien
		\ifTLdapso
			\def\@tlarg{#1}%
			\ifx\@tlarg\@empty\else
				\ifdefstring{\kindOPTN}{dethi}{}{\par\noindent{\bfseries\color{gold1}Đáp số:} #1\par}%
			\fi
		\else
			\loigiai{#2}%
		\fi
	\fi}
\makeatother
\TLdapso% mặc định: chỉ in đáp số bài tự luyện
%====================== BẢNG ĐÚNG/SAI (\choiceTFt) ======================
%  Hàng tiêu đề nền nâu đồng chữ trắng (Phát biểu | Đ | S), các hàng tô xen kẽ kem / trắng,
%  đường kẻ vàng kim, nhãn a) b) c) d) màu nâu đồng. Bản có đáp án: cột được chọn có vòng tròn
%  nâu đồng với dấu tích trắng (câu đúng tích ở cột Đúng, câu sai tích ở cột Sai).
%  Tiêu đề ghi đủ "Đúng | Sai": \OPTN{viettat=0} sau \input (cột Đúng/Sai rộng ra, phát biểu hẹp lại).
%  Dùng bảng gốc của ex_test: \def\BangDSgoc{1} trước %============================================================
%  PREAMBLE EDUZY v2 - dùng chung cho đề thi, bài học, chuyên đề (gói ex_test)
%  Khác bản preamble-eduzy.tex ở đầu trang, chân trang: dải nền nâu đồng tràn ngang,
%  logo + Eduzy.com.vn, ô chữ vàng kem cạnh xiên, chân trang "Trang (số)";
%  logo Eduzy chìm giữa mọi trang (tắt: \logochimtat, bật lại: \logochimbat).
%  Cách dùng ở file chính (xem mau-chuan.tex, mau-chuyen-de.tex):
%      \def\CheDo{loigiai}          % loigiai | dethi | solcolor | color | book
%      \def\tenchuong{<dòng chữ trên đầu trang>}
%      \def\thuonghieu{<chữ ở chân trang lẻ>}   % không khai báo thì dùng mặc định
%      \def\tacgia{<chữ ở chân trang chẵn>}      % không khai báo thì dùng mặc định
%      \def\khauhieu{<chữ cạnh logo ở đầu trang>}  % tùy chọn, mặc định Eduzy.com.vn
%      \input{preamble-eduzy-v2}
%      \TLdapso                      % tài liệu: \TLfull | \TLdapso | \TLan
%      \begin{document} ... \end{document}
%  Logo Eduzy vẽ bằng TikZ ngay trong file (\eduzylogo), không cần file ảnh.
%  Biên dịch: pdfLaTeX (máy có gói vietnam) hoặc XeLaTeX.
%============================================================
\documentclass[12pt,a4paper,twoside]{article}%<-- twoside: bat dau trang/chan trang chan-le
\usepackage{iftex}
\usepackage{amsmath,amssymb,tikz,mathrsfs,tkz-tab}
\usepackage{graphicx}
\usepackage{pgfplots}% đồ thị hàm số
\pgfplotsset{compat=1.17}
%\usepackage{mathptmx}
\usepackage{mathpazo}
\ifXeTeX% XeLaTeX: font chữ có đủ dấu tiếng Việt, giữ nguyên font toán mathpazo
	\usepackage[no-math]{fontspec}
	\setmainfont{TeX Gyre Pagella}
	\newfontfamily\EZqag{TeX Gyre Adventor}[NFSSFamily=qag,Ligatures=TeX]
\fi
\usepackage[many]{tcolorbox}
\usetikzlibrary{math,through,calc,intersections,angles,quotes,shapes,shapes.geometric,arrows,patterns,matrix,chains,arrows.meta,decorations.shapes,decorations.fractals,decorations.markings,shadows,decorations.pathreplacing}
\usepackage[top=1.6cm, bottom=1.8cm, left=1.8cm, right=1.8cm]{geometry}% v2: top 1.6cm (bản cũ 1.5cm) để chữ không sát dải đầu trang
\providecommand{\CheDo}{loigiai}% loigiai | dethi | solcolor | color | book
\expandafter\usepackage\expandafter[\CheDo]{ex_test}
% \tenchuong: khai báo ở file chính (trước \input) hoặc \renewcommand sau \input
%============================================================
%  [GỘP TỪ eduzy-de.tex] EDUZY - MÀU NÂU ĐỒNG / VÀNG KIM
%  + ĐẦU TRANG, CHÂN TRANG KIỂU DẢI NGANG (v2), ĐỐI XỨNG TRANG CHẴN / TRANG LẺ.
%  (không cần file eduzy-de.tex riêng nữa)
%  Khối này phải đặt SAU \usepackage[...]{ex_test} và SAU \def\tenchuong{...}
%============================================================
\usepackage{lastpage}
\usepackage{needspace}
%--- DÀN TRANG ĐỀ THI: tách câu (THÊM 05/11/2026; chỉ thêm lệnh mới, KHÔNG đổi hành vi mặc định) ---
%  Hai công tắc, mặc định tắt (= preamble chạy y như cũ). Bật bằng cách đặt TRƯỚC \input{preamble-eduzy-v2}
%  trong file đề (tools/dan_trang/dan_trang.py tự ghi khi đặt \NeedspaceCau):
%    \def\TachCauDS{1}   câu Đúng/Sai được tách: câu dẫn (kèm hình) ở cuối trang, bảng Phát biểu sang trang sau
%                        (câu dẫn và hình không tách nhau, bảng không tách giữa hai trang)
%    \def\TachCauLG{1}   bản lời giải (\CheDo khác dethi): phần đề của câu cao hơn 1/3 trang được phép tách
%  Lệnh dùng trong file đề (giá trị do tools/dan_trang/dan_trang.py đo ở bản dethi):
%    \NeedspaceCau[<cao câu dẫn + hình>]{<cao cả khối câu>}   ngay trên \begin{ex} (câu đầu phần: trên tiêu đề phần)
%        công tắc tắt: đúng bằng \Needspace{<cao cả khối câu>} như cũ;
%        \TachCauDS=1 và có [..] (câu Đúng/Sai được chọn tách): chỉ giữ câu dẫn + hình;
%        \TachCauLG=1 ở bản lời giải: đòi tối đa 1/3 \textheight.
%    \NeedspaceBang{<cao bảng Phát biểu>}%   ngay trước \choiceTFt: \TachCauDS=1 thì giữ trọn bảng trên một trang;
%        công tắc tắt: không làm gì.
\providecommand{\TachCauDS}{0}
\providecommand{\TachCauLG}{0}
\newlength{\EZnsCau}
\newcommand{\NeedspaceCau}[2][]{%
	\setlength{\EZnsCau}{#2}%
	\ifnum\TachCauDS=1 \if\relax\detokenize{#1}\relax\else\setlength{\EZnsCau}{#1}\fi\fi
	\ifnum\TachCauLG=1 \ifdefstring{\kindOPTN}{dethi}{}{\ifdim\EZnsCau>0.3333\textheight\setlength{\EZnsCau}{0.3333\textheight}\fi}\fi
	\Needspace{\EZnsCau}}
\newcommand{\NeedspaceBang}[1]{\ifnum\TachCauDS=1 \Needspace{#1}\fi}

%--- Thang màu: đậm -> nhạt
\definecolor{gold1}{RGB}{126,96,0}    % nâu đồng đậm  #7E6000
\definecolor{gold2}{RGB}{192,144,3}   % vàng kim      #C09003
\definecolor{gold3}{RGB}{255,215,99}  % vàng nhạt     #FFD763
\definecolor{gold4}{RGB}{255,231,153} % vàng rất nhạt #FFE799
\definecolor{gold5}{RGB}{255,243,205} % vàng kem      #FFF3CD
%--- Ghi đè tông xanh/đỏ cũ
\definecolor{myblue}{RGB}{126,96,0}
\colorlet{blue}{gold1}
\colorlet{cyan}{gold2}
\colorlet{teal}{gold2}

%--- Tên chương mặc định (file đề nên tự khai báo lại)
\providecommand{\tenchuong}{}
%--- Chân trang
\providecommand{\thuonghieu}{TOÁN 12 -- Eduzy.com.vn}
\providecommand{\tacgia}{Tô Gia Bảo -- 0353703783}

%--- Chữ cạnh logo ở đầu trang
\providecommand{\khauhieu}{Eduzy.com.vn}

%====================== THÔNG SỐ ĐẦU TRANG, CHÂN TRANG (v2) ======================
%  Dải nền đậm tràn ngang trang; ô chữ vàng kem có cạnh xiên, kèm 2 vạch xiên chuyển màu.
%  Đầu trang lẻ : logo + \khauhieu (trái)   | ô chữ \tenchuong (phải)
%  Chân trang lẻ: ô chữ \thuonghieu (trái)  | Trang (số) (phải)
%  Trang chẵn soi gương trang lẻ; chân trang chẵn ghi \tacgia.
%--- Màu: chỉ dùng thang vàng kim Eduzy
\colorlet{ezNen}{gold1}% dải nền đậm
\colorlet{ezVach1}{gold3}% vạch xiên sát dải nền (đậm hơn)
\colorlet{ezVach2}{gold4}% vạch xiên sát ô chữ (nhạt hơn)
\colorlet{ezO}{gold5}% ô chữ vàng kem
%--- Kích thước
\newlength{\EZdai}\setlength{\EZdai}{1.1cm}% chiều cao dải nền
\newlength{\EZo}\setlength{\EZo}{0.68cm}% chiều cao ô chữ
\newlength{\EZxien}\setlength{\EZxien}{0.55cm}% độ lệch ngang của cạnh xiên
\newlength{\EZvach}\setlength{\EZvach}{0.75cm}% bề rộng mỗi vạch xiên
\newlength{\EZkhe}\setlength{\EZkhe}{0.12cm}% khe trắng giữa vạch và ô chữ
\newlength{\EZdem}\setlength{\EZdem}{0.4cm}% khoảng đệm từ chữ tới cạnh xiên
\newlength{\EZle}\setlength{\EZle}{1.8cm}% lề trái/phải của chữ (bằng lề thân bài)
\newlength{\EZrongmin}\setlength{\EZrongmin}{4.5cm}% ô chữ rộng tối thiểu
\newlength{\EZrong}% bề rộng chữ trong ô (tính tự động)
%--- Font
\newcommand{\EZfontO}{\fontsize{10.5pt}{12pt}\selectfont\bfseries}% chữ trong ô
\newcommand{\EZfontKH}{\fontfamily{qag}\fontsize{11pt}{12pt}\selectfont\bfseries}% \khauhieu cạnh logo
\newcommand{\EZfontTrang}{\fontsize{11.5pt}{12pt}\selectfont\bfseries}% chữ "Trang" và số trang
%--- Đo bề rộng chữ trong ô
\newcommand{\EZdorong}[1]{%
	\settowidth{\EZrong}{\EZfontO #1}%
	\ifdim\EZrong<\EZrongmin\setlength{\EZrong}{\EZrongmin}\fi}
%--- Logo + chữ trắng \khauhieu (dựng sẵn một lần khi vào \begin{document})
\newsavebox{\EZlogobox}
\AtBeginDocument{\sbox{\EZlogobox}{%
		\raisebox{-0.5\height}{\eduzylogo[1cm]}\hspace{0.15cm}%
		\raisebox{-0.5\height}{\color{white}\EZfontKH\khauhieu}}}
%--- Hai vạch xiên + khe trắng. #1: điểm gốc ở mép dưới, #2: hướng (1 sang phải, -1 sang trái),
%    #3: hướng nghiêng (1 kiểu "/", -1 kiểu "\")
\newcommand{\EZvachxien}[3]{%
	\fill[white] (#1) --++({(#3)*\EZxien},\EZo) --++({(#2)*\EZkhe},0) --++({-(#3)*\EZxien},-\EZo) --cycle;
	\fill[ezVach2] ([xshift=(#2)*\EZkhe]#1) --++({(#3)*\EZxien},\EZo) --++({(#2)*\EZvach},0) --++({-(#3)*\EZxien},-\EZo) --cycle;
	\fill[ezVach1] ([xshift=(#2)*(\EZkhe+\EZvach)]#1) --++({(#3)*\EZxien},\EZo) --++({(#2)*\EZvach},0) --++({-(#3)*\EZxien},-\EZo) --cycle;
	\draw[white,line width=0.6pt] ([xshift=(#2)*(\EZkhe+\EZvach)]#1) --++({(#3)*\EZxien},\EZo);
}
%--- Chữ trong ô (căn giữa theo chiều cao ô, không lệch vì dấu tiếng Việt)
\tikzset{ezchu/.style={inner sep=0pt,text=ezNen,text height=0.72em,text depth=0.22em}}

%====================== LOGO CHÌM GIỮA TRANG ======================
%  Vẽ cùng đầu trang nên nằm dưới chữ của thân bài, có trên mọi trang.
%  Tắt: \logochimtat (từ trang đó trở đi), bật lại: \logochimbat.
%  Đổi cỡ, độ đậm: khai báo trước \input, ví dụ \def\EZchimcao{9cm} \def\EZchimdo{0.1}
\providecommand{\EZchimcao}{10cm}% chiều cao logo chìm
\providecommand{\EZchimdo}{0.12}% độ đậm (0 = trong suốt, 1 = rõ hẳn)
\newif\ifEZchim \EZchimtrue
\newcommand{\logochimtat}{\EZchimfalse}
\newcommand{\logochimbat}{\EZchimtrue}
%  Dựng sẵn logo mờ một lần khi vào \begin{document}; transparency group làm mờ cả logo
%  một lượt nên không lộ vệt ở chỗ các dải màu chồng nhau.
\newsavebox{\EZchimtho}
\newsavebox{\EZchimbox}
\AtBeginDocument{%
	\sbox{\EZchimtho}{\eduzylogo[\EZchimcao]}%
	\sbox{\EZchimbox}{\begin{tikzpicture}
			\begin{scope}[transparency group,opacity=\EZchimdo]
				\node[inner sep=0pt] {\usebox{\EZchimtho}};
			\end{scope}
	\end{tikzpicture}}}
\newcommand{\EZvechim}{%
	\ifEZchim
		\node[inner sep=0pt] at (current page.center) {\usebox{\EZchimbox}};
	\fi}

%====================== ĐẦU TRANG - TRANG LẺ ======================
%  (logo bên trái, ô chữ \tenchuong bên phải)
\newcommand{\myfancyheadOdd}{%
	\EZdorong{\tenchuong}%
	\begin{tikzpicture}[remember picture,overlay]
		\coordinate (T) at (current page.north west);
		\coordinate (P) at (current page.north east);
		%--- logo chìm giữa trang
		\EZvechim
		%--- dải nền
		\fill[ezNen] (T) rectangle ([yshift=-\EZdai]P);
		%--- ô chữ: X là góc dưới bên trái, cạnh trái xiên kiểu "/"
		\coordinate (X) at ([xshift=-\EZle-\EZrong-\EZdem-\EZxien,yshift=-\EZdai]P);
		\fill[ezO] (X) --++(\EZxien,\EZo) -- ([yshift=-\EZdai+\EZo]P) -- ([yshift=-\EZdai]P) --cycle;
		\EZvachxien{X}{-1}{1}
		\node[ezchu,anchor=east] at ([xshift=-\EZle,yshift=-\EZdai+0.5*\EZo]P) {\EZfontO\tenchuong};
		%--- logo + tên
		\node[inner sep=0pt,anchor=west] at ([xshift=\EZle,yshift=-0.5*\EZdai]T) {\usebox{\EZlogobox}};
	\end{tikzpicture}%
}

%====================== ĐẦU TRANG - TRANG CHẴN =====================
%  (soi gương: ô chữ \tenchuong bên trái, logo bên phải)
\newcommand{\myfancyheadEven}{%
	\EZdorong{\tenchuong}%
	\begin{tikzpicture}[remember picture,overlay]
		\coordinate (T) at (current page.north west);
		\coordinate (P) at (current page.north east);
		%--- logo chìm giữa trang
		\EZvechim
		%--- dải nền
		\fill[ezNen] (T) rectangle ([yshift=-\EZdai]P);
		%--- ô chữ: X là góc dưới bên phải, cạnh phải xiên kiểu "\"
		\coordinate (X) at ([xshift=\EZle+\EZrong+\EZdem+\EZxien,yshift=-\EZdai]T);
		\fill[ezO] (X) --++(-\EZxien,\EZo) -- ([yshift=-\EZdai+\EZo]T) -- ([yshift=-\EZdai]T) --cycle;
		\EZvachxien{X}{1}{-1}
		\node[ezchu,anchor=west] at ([xshift=\EZle,yshift=-\EZdai+0.5*\EZo]T) {\EZfontO\tenchuong};
		%--- logo + tên
		\node[inner sep=0pt,anchor=east] at ([xshift=-\EZle,yshift=-0.5*\EZdai]P) {\usebox{\EZlogobox}};
	\end{tikzpicture}%
}

%====================== CHÂN TRANG - TRANG LẺ ======================
%  (ô chữ \thuonghieu bên trái, "Trang (số)" bên phải)
\newcommand{\myfancyfootOdd}{%
	\EZdorong{\thuonghieu}%
	\begin{tikzpicture}[remember picture,overlay]
		\coordinate (T) at (current page.south west);
		\coordinate (P) at (current page.south east);
		%--- dải nền
		\fill[ezNen] (T) rectangle ([yshift=\EZdai]P);
		%--- ô chữ sát mép trên dải: X là góc dưới bên phải, cạnh phải xiên kiểu "/"
		\coordinate (X) at ([xshift=\EZle+\EZrong+\EZdem,yshift=\EZdai-\EZo]T);
		\fill[ezO] ([yshift=\EZdai-\EZo]T) -- (X) --++(\EZxien,\EZo) -- ([yshift=\EZdai]T) --cycle;
		\EZvachxien{X}{1}{1}
		\node[ezchu,anchor=west] at ([xshift=\EZle,yshift=\EZdai-0.5*\EZo]T) {\EZfontO\thuonghieu};
		%--- số trang
		\node[circle,draw=white,line width=0.9pt,text=white,inner sep=0pt,minimum size=0.62cm,anchor=east]
		(so) at ([xshift=-\EZle,yshift=0.5*\EZdai]P) {\EZfontTrang\thepage};
		\node[inner sep=0pt,text=white,anchor=base east] (tr) at ([xshift=-0.15cm]so.base west) {\EZfontTrang Trang};
	\end{tikzpicture}%
}

%====================== CHÂN TRANG - TRANG CHẴN ====================
%  (soi gương: "Trang (số)" bên trái, ô chữ \tacgia bên phải)
\newcommand{\myfancyfootEven}{%
	\EZdorong{\tacgia}%
	\begin{tikzpicture}[remember picture,overlay]
		\coordinate (T) at (current page.south west);
		\coordinate (P) at (current page.south east);
		%--- dải nền
		\fill[ezNen] (T) rectangle ([yshift=\EZdai]P);
		%--- ô chữ sát mép trên dải: X là góc dưới bên trái, cạnh trái xiên kiểu "\"
		\coordinate (X) at ([xshift=-\EZle-\EZrong-\EZdem,yshift=\EZdai-\EZo]P);
		\fill[ezO] ([yshift=\EZdai-\EZo]P) -- (X) --++(-\EZxien,\EZo) -- ([yshift=\EZdai]P) --cycle;
		\EZvachxien{X}{-1}{-1}
		\node[ezchu,anchor=east] at ([xshift=-\EZle,yshift=\EZdai-0.5*\EZo]P) {\EZfontO\tacgia};
		%--- số trang
		%    (vòng tròn ẩn so0 nằm giữa dải để lấy đường chân chữ, giống trang lẻ)
		\node[circle,inner sep=0pt,minimum size=0.62cm] (so0) at ([yshift=0.5*\EZdai]T) {\EZfontTrang\phantom{\thepage}};
		\node[inner sep=0pt,text=white,anchor=base west] (tr) at ($(so0.base -| T)+(\EZle,0)$) {\EZfontTrang Trang};
		\node[circle,draw=white,line width=0.9pt,text=white,inner sep=0pt,minimum size=0.62cm,anchor=base west]
		(so) at ([xshift=0.15cm]tr.base east) {\EZfontTrang\thepage};
	\end{tikzpicture}%
}

%--- Giữ tương thích với các file cũ đang gọi \myfancyhead, \myfancyfoot
\newcommand{\myfancyhead}{\myfancyheadOdd}
\newcommand{\myfancyfoot}{\myfancyfootOdd}

%======================= GẮN VÀO TRANG =============================
\pagestyle{fancy}
\fancyhf{}
\fancyhead[LO]{\myfancyheadOdd}
\fancyhead[RE]{\myfancyheadEven}
\fancyfoot[LO]{\myfancyfootOdd}
\fancyfoot[RE]{\myfancyfootEven}
\renewcommand{\headrulewidth}{0pt}
\renewcommand{\footrulewidth}{0pt}
\fancypagestyle{plain}{%
	\fancyhf{}
	\fancyhead[LO]{\myfancyheadOdd}
	\fancyhead[RE]{\myfancyheadEven}
	\fancyfoot[LO]{\myfancyfootOdd}
	\fancyfoot[RE]{\myfancyfootEven}
	\renewcommand{\headrulewidth}{0pt}
	\renewcommand{\footrulewidth}{0pt}
}
%=================== HẾT PHẦN GỘP TỪ eduzy-de.tex ===================
\newcommand{\indam}[1]{
	{\fontsize{12pt}{9pt}\sffamily \selectfont\bfseries\color{gold1} #1}
}
\newcommand{\ind}[1]{
	{\fontsize{14pt}{9pt}\fontfamily{qag} \selectfont\bfseries\color{gold1} #1}
}
\renewcommand{\FalseEX}{\stepcounter{dapan}\textbf{\selectfont\color{gold1}\Alph{dapan}.}\ignorespaces}
\renewtheorem{ex}{\color{gold1}Câu}
\renewcommand{\baselinestretch}{1.1}
\newcommand{\hoac}[1]{
	\left[\begin{aligned}#1\end{aligned}\right.}
\newcommand{\heva}[1]{
	\left\{\begin{aligned}#1\end{aligned}\right.}
\usepackage{esvect}
\def\vec{\vv}
\def\vt{\vv}
\def\overrightarrow{\vv}
%====================== LOGO EDUZY VẼ BẰNG TIKZ ======================
%  Đường viền dò từ logo gốc; dải màu vàng kim đo trực tiếp trên ảnh gốc
%  (nghiêng 21 độ, tô bằng 80 dải hẹp nên hiển thị giống nhau trên mọi trình xem PDF).
%  Dùng: \eduzylogo          -> cao 2.4cm (mặc định)
%        \eduzylogo[1.5cm]   -> cao 1.5cm
\definecolor{ezText}{RGB}{245,200,46}
\newcommand{\eduzylogo}[1][2.4cm]{%
	\resizebox{!}{#1}{%
		\begin{tikzpicture}[x=0.01cm,y=0.01cm]
			\path[use as bounding box] (0,0) rectangle (224,-224);
			% mũ tốt nghiệp + vương miện
			\begin{scope}
				\clip
				(119.4,-176.1)..controls(118.9,-176.1)and(116.8,-176)..(114.6,-176)
				..controls(101.4,-176)and(85.8,-175.9)..(83.7,-175.7)
				..controls(81.8,-175.6)and(81,-175.6)..(80.2,-175.7)
				..controls(78.8,-175.9)and(77.7,-175.9)..(76.2,-175.7)
				..controls(74.9,-175.5)and(70.8,-175.5)..(64.3,-175.7)
				..controls(59.1,-175.9)and(58.3,-175.7)..(58.6,-173.7)
				..controls(58.6,-172.8)and(59.6,-169.6)..(60.3,-168)
				..controls(60.8,-166.9)and(61.4,-165.8)..(62.3,-165)
				..controls(63.1,-164.2)and(64.3,-163.4)..(65,-163.2)
				..controls(65.5,-163.1)and(72.3,-163)..(76.6,-163.1)
				..controls(77.9,-163.1)and(80.1,-163.1)..(81.4,-163)
				..controls(82.8,-162.9)and(86.3,-162.9)..(89.2,-162.9)
				..controls(93,-162.9)and(94.9,-162.8)..(96,-162.7)
				..controls(97.2,-162.5)and(101,-162.5)..(114,-162.4)
				..controls(123,-162.3)and(134.9,-162.2)..(140.2,-162.1)
				..controls(150.1,-162)and(152.2,-162.1)..(153.6,-162.4)
				..controls(154.4,-162.6)and(157.1,-163.8)..(158.1,-164.4)
				..controls(158.8,-164.8)and(159.7,-165.7)..(159.9,-166.3)
				..controls(160,-166.5)and(160.2,-167.1)..(160.2,-167.6)
				..controls(160.3,-168)and(160.6,-169.4)..(160.9,-170.6)
				..controls(161.2,-171.7)and(161.5,-173.1)..(161.6,-173.6)
				..controls(161.8,-174.4)and(161.8,-174.5)..(161.5,-174.9)
				..controls(161.2,-175.6)and(160.8,-175.9)..(159.9,-176)
				..controls(159.2,-176.1)and(154.1,-176.2)..(151.5,-176.2)
				..controls(148.1,-176.1)and(122.6,-176.1)..(121.6,-176.2)
				..controls(120.8,-176.2)and(119.8,-176.2)..(119.4,-176.1)--cycle
				(130.4,-160)..controls(129,-160)and(119.8,-160)..(110.1,-159.9)
				..controls(100.3,-159.8)and(88.5,-159.7)..(83.9,-159.6)
				..controls(77.8,-159.5)and(75.3,-159.5)..(74.5,-159.6)
				..controls(73.9,-159.7)and(72.5,-159.8)..(71.3,-159.8)
				..controls(69.6,-159.8)and(69,-159.8)..(68.5,-159.6)
				..controls(67.6,-159.3)and(66.8,-158.6)..(66.2,-157.6)
				..controls(66,-157.3)and(65.6,-156.5)..(65.3,-156.1)
				..controls(64.6,-155.1)and(64.5,-154.9)..(63.7,-152.9)
				..controls(63.2,-151.6)and(62.9,-150.9)..(61.8,-149.6)
				..controls(61.5,-149.2)and(61.1,-148.5)..(60.9,-147.9)
				..controls(60.7,-147.4)and(60.3,-146.6)..(60.1,-146.1)
				..controls(59.9,-145.5)and(59.5,-144.9)..(59,-144.3)
				..controls(58.5,-143.7)and(58.2,-143.2)..(57.9,-142.5)
				..controls(57.7,-141.9)and(57.3,-141.1)..(57,-140.7)
				..controls(56.6,-140.3)and(56.1,-139.3)..(55.8,-138.5)
				..controls(55.5,-137.8)and(55.1,-136.8)..(54.9,-136.5)
				..controls(54.7,-136.1)and(54.3,-135.2)..(53.9,-134.5)
				..controls(53.6,-133.8)and(53,-132.7)..(52.6,-132)
				..controls(51.6,-130.5)and(51.4,-130.2)..(51.1,-129)
				..controls(50.6,-127.5)and(50.6,-127.4)..(49.7,-126.3)
				..controls(49.2,-125.7)and(48.6,-124.8)..(48.3,-124.2)
				..controls(48,-123.6)and(47.5,-122.7)..(47.3,-122.2)
				..controls(47,-121.8)and(46.5,-120.7)..(46.1,-119.8)
				..controls(45.8,-118.9)and(45.2,-117.8)..(44.9,-117.2)
				..controls(44.2,-116.1)and(43.3,-114.4)..(42.3,-112.2)
				..controls(41.8,-111.3)and(41.3,-110.4)..(41.1,-110.1)
				..controls(40.5,-109.4)and(40.3,-109)..(40,-107.8)
				..controls(39.8,-107)and(39.6,-106.5)..(39.2,-106)
				..controls(38.9,-105.6)and(38.4,-104.7)..(38,-104)
				..controls(37.7,-103.4)and(37.2,-102.5)..(37,-102.1)
				..controls(36.7,-101.7)and(36.3,-100.8)..(36,-100.1)
				..controls(35.6,-98.9)and(35.3,-98.4)..(34.6,-97.6)
				..controls(34.2,-97.1)and(34.2,-96.3)..(34.6,-96)
				..controls(35,-95.7)and(35.5,-95.9)..(36.2,-96.7)
				..controls(36.5,-97)and(37.2,-97.5)..(37.6,-97.7)..controls(38,-98)and(38.6,-98.4)..(39.1,-98.9)
				..controls(39.6,-99.4)and(40.3,-100)..(40.6,-100.3)
				..controls(41.3,-100.8)and(42.4,-101.8)..(43.7,-102.9)
				..controls(44,-103.2)and(44.6,-103.7)..(45,-103.9)
				..controls(45.4,-104.2)and(45.9,-104.6)..(46.1,-104.8)
				..controls(47.1,-105.9)and(48.8,-107.1)..(50.5,-107.9)
				..controls(51.8,-108.6)and(52.3,-109)..(53.4,-110)
				..controls(54.8,-111.2)and(55.6,-111.6)..(57,-112.1)
				..controls(57.5,-112.3)and(58.5,-112.8)..(59.4,-113.3)
				..controls(60.2,-113.7)and(61.3,-114.3)..(61.9,-114.5)
				..controls(62.4,-114.7)and(63.3,-115.1)..(63.8,-115.3)
				..controls(64.4,-115.6)and(65.4,-115.9)..(66.2,-116.1)
				..controls(67,-116.3)and(68.4,-116.6)..(69.2,-116.9)
				..controls(72.3,-117.9)and(72.8,-118)..(78,-118.1)
				..controls(84,-118.3)and(86.1,-118.1)..(88.1,-117.1)
				..controls(89.4,-116.5)and(93,-115)..(93.7,-114.9)
				..controls(94.8,-114.6)and(95.4,-114.2)..(95.9,-113.4)
				..controls(96.2,-113)and(96.7,-112.5)..(97.1,-112.2)
				..controls(98,-111.6)and(101.3,-108.4)..(102.1,-107.2)
				..controls(102.4,-106.8)and(102.9,-106.2)..(103.2,-105.8)
				..controls(103.5,-105.5)and(104,-104.8)..(104.2,-104.2)
				..controls(104.5,-103.6)and(105,-102.8)..(105.3,-102.4)
				..controls(105.6,-102)and(106.1,-101.2)..(106.5,-100.6)
				..controls(107.1,-99.5)and(108.6,-97.5)..(109.5,-96.5)
				..controls(109.9,-96.1)and(110.1,-96)..(110.4,-96)
				..controls(111.1,-96)and(111.6,-96.5)..(111.8,-97.6)
				..controls(112,-98.6)and(112.3,-99.3)..(112.8,-99.9)
				..controls(113.7,-100.9)and(115.2,-103.1)..(115.7,-104.2)
				..controls(116.1,-105)and(116.5,-105.5)..(116.9,-106)
				..controls(117.3,-106.4)and(117.8,-106.9)..(118,-107.3)
				..controls(118.2,-107.6)and(118.6,-108)..(119,-108.3)
				..controls(119.6,-108.8)and(121.1,-110.3)..(121.8,-111.2)
				..controls(122.1,-111.7)and(122.5,-112)..(123.4,-112.4)
				..controls(124.8,-113.2)and(125.8,-113.8)..(127,-114.6)
				..controls(127.5,-114.9)and(128.1,-115.3)..(128.2,-115.4)
				..controls(128.4,-115.5)and(129.3,-115.8)..(130.2,-116)
				..controls(131.1,-116.2)and(132.5,-116.6)..(133.3,-116.9)
				..controls(135.7,-117.9)and(136.1,-117.9)..(142.5,-117.8)
				..controls(145.8,-117.8)and(148.4,-117.7)..(148.8,-117.6)
				..controls(149.2,-117.6)and(150.2,-117.2)..(151,-116.8)
				..controls(152,-116.4)and(152.8,-116)..(153.6,-115.9)
				..controls(154.5,-115.7)and(155.2,-115.4)..(156.3,-114.9)
				..controls(157.1,-114.5)and(158.2,-114)..(158.7,-113.8)
				..controls(159.1,-113.7)and(159.8,-113.3)..(160.2,-113)
				..controls(160.6,-112.7)and(161.3,-112.4)..(161.7,-112.2)
				..controls(162.7,-111.9)and(163.2,-111.7)..(164.1,-111.1)
				..controls(164.5,-110.8)and(165.3,-110.4)..(165.9,-110)
				..controls(166.5,-109.7)and(167.3,-109.2)..(167.6,-108.9)
				..controls(167.9,-108.6)and(168.6,-108.2)..(169,-108)
				..controls(169.4,-107.8)and(170.2,-107.3)..(170.7,-106.9)
				..controls(171.2,-106.5)and(172.1,-105.8)..(172.7,-105.4)
				..controls(173.2,-105)and(174.1,-104.4)..(174.5,-104)
				..controls(174.9,-103.7)and(175.7,-103.1)..(176.2,-102.6)
				..controls(176.7,-102.2)and(177.4,-101.5)..(177.9,-101.2)
				..controls(178.3,-100.8)and(178.8,-100.4)..(179,-100.2)
				..controls(179.2,-100)and(179.6,-99.6)..(180,-99.3)
				..controls(180.4,-99.1)and(180.9,-98.6)..(181.2,-98.4)
				..controls(181.5,-98.1)and(182.1,-97.7)..(182.5,-97.5)
				..controls(183.2,-97.2)and(183.5,-96.9)..(184.3,-95.9)
				..controls(184.6,-95.6)and(184.6,-95.6)..(184.8,-95.7)
				..controls(185,-95.8)and(185.1,-96)..(185.1,-96.2)
				..controls(185.1,-96.5)and(185,-96.7)..(184.6,-97.2)
				..controls(184.1,-97.8)and(184,-97.9)..(183.9,-98.8)
				..controls(183.7,-100.3)and(183.3,-101.3)..(182.9,-101.8)
				..controls(182.6,-102.1)and(182.3,-102.6)..(182.2,-102.9)
				..controls(182,-103.3)and(181.6,-104)..(181.2,-104.6)
				..controls(180.9,-105.2)and(180.5,-106.1)..(180.4,-106.6)
				..controls(179.9,-108)and(179.7,-108.5)..(179,-109.4)
				..controls(178.6,-109.8)and(178.2,-110.5)..(177.9,-111)
				..controls(177.7,-111.5)and(177.3,-112.3)..(177,-112.7)
				..controls(176.6,-113.2)and(176.4,-113.7)..(176.3,-114.2)
				..controls(176.2,-114.7)and(175.9,-115.2)..(175.4,-116)
				..controls(175,-116.6)and(174.4,-117.7)..(174,-118.5)
				..controls(173.7,-119.2)and(173.3,-120.1)..(173,-120.5)
				..controls(172.8,-120.9)and(172.4,-121.8)..(172.1,-122.5)
				..controls(171.8,-123.2)and(171.3,-124.1)..(171,-124.6)
				..controls(170.8,-125.1)and(170.3,-126)..(169.9,-126.6)
				..controls(169.6,-127.3)and(169.2,-128.1)..(169,-128.3)
				..controls(168.8,-128.6)and(168.5,-129.3)..(168.3,-129.8)
				..controls(168.1,-130.3)and(167.8,-131.1)..(167.5,-131.6)
				..controls(166.9,-132.5)and(164.2,-137.9)..(164.2,-138.1)
				..controls(164.2,-138.1)and(163.9,-138.8)..(163.4,-139.7)
				..controls(162.9,-140.5)and(162.4,-141.5)..(162.3,-141.9)
				..controls(161.9,-143)and(161.6,-143.4)..(160.9,-144.1)
				..controls(160.4,-144.6)and(160.3,-144.8)..(160.1,-145.4)
				..controls(160,-145.8)and(159.6,-146.8)..(159.2,-147.8)
				..controls(158.7,-148.8)and(158.1,-150)..(157.9,-150.6)
				..controls(157.6,-151.1)and(157.2,-151.9)..(156.9,-152.4)
				..controls(156.5,-152.8)and(156.1,-153.7)..(155.8,-154.4)
				..controls(155.6,-155)and(155.2,-155.9)..(154.9,-156.3)
				..controls(154.7,-156.7)and(154.4,-157.4)..(154.2,-157.9)
				..controls(154.1,-158.4)and(153.8,-159)..(153.5,-159.3)
				..controls(152.8,-160.1)and(153,-160.1)..(142.4,-160.1)
				..controls(137.2,-160.1)and(131.8,-160.1)..(130.4,-160)--cycle
				(141.5,-103.7)..controls(140.8,-103)and(139.8,-102.1)..(139.2,-101.6)
				..controls(138.5,-101.2)and(137.9,-100.6)..(137.8,-100.4)
				..controls(137,-99.5)and(134,-96.5)..(133.2,-96)
				..controls(132.7,-95.7)and(131.8,-94.9)..(131.1,-94.3)
				..controls(129.4,-92.7)and(127.9,-91.6)..(126.6,-91.1)
				..controls(125.2,-90.5)and(124.1,-90)..(122.9,-89.1)
				..controls(122.1,-88.6)and(121.3,-88.2)..(119.9,-87.7)
				..controls(118.8,-87.3)and(117.5,-86.8)..(116.9,-86.5)
				..controls(114.3,-85.5)and(113.7,-85.4)..(109.7,-85.5)
				..controls(105.3,-85.7)and(103.9,-85.9)..(101.6,-87)
				..controls(100.7,-87.4)and(99.4,-87.8)..(97.9,-88)
				..controls(97,-88.2)and(96.6,-88.3)..(95.8,-88.9)
				..controls(95.4,-89.2)and(94.7,-89.6)..(94.2,-89.8)
				..controls(92.9,-90.3)and(92.3,-90.7)..(90.9,-92.1)
				..controls(90.2,-92.7)and(89.4,-93.3)..(89,-93.5)..controls(87,-94.5)and(85.2,-96)..(83.6,-97.9)
				..controls(83.3,-98.2)and(82.5,-98.9)..(81.8,-99.4)
				..controls(80.7,-100.1)and(80.5,-100.3)..(80,-101.1)
				..controls(78.7,-103.1)and(77.2,-103.9)..(76.4,-103)
				..controls(75.7,-102.2)and(75.7,-103.2)..(75.7,-83.6)
				..controls(75.6,-73.6)and(75.6,-64.8)..(75.5,-64)
				..controls(75.5,-63.2)and(75.5,-61.4)..(75.6,-59.8)--(75.7,-57.1)--(76,-56.8)
				..controls(76.8,-56.2)and(78,-56.3)..(79.8,-57.2)
				..controls(80.4,-57.5)and(81.3,-57.8)..(81.8,-57.9)
				..controls(82.3,-58.1)and(83,-58.4)..(83.6,-58.8)..controls(84.3,-59.2)and(85,-59.5)..(86,-59.8)
				..controls(86.8,-60.1)and(88,-60.6)..(88.7,-60.9)
				..controls(89.3,-61.2)and(90.6,-61.7)..(91.4,-62)
				..controls(92.2,-62.3)and(93.7,-62.9)..(94.6,-63.3)
				..controls(95.6,-63.7)and(96.8,-64.1)..(97.4,-64.3)
				..controls(98,-64.5)and(99,-64.9)..(99.7,-65.2)
				..controls(100.4,-65.6)and(101.5,-66)..(102.3,-66.2)
				..controls(104.2,-66.8)and(105,-67.1)..(106.2,-67.8)
				..controls(107.9,-68.7)and(110.2,-69.4)..(111.3,-69.2)
				..controls(112.5,-69)and(115,-68.2)..(115.6,-67.8)
				..controls(117.3,-66.7)and(117.6,-66.6)..(119.5,-66.2)
				..controls(121.9,-65.8)and(122.1,-65.7)..(123,-65.1)
				..controls(124.1,-64.4)and(124.4,-64.3)..(126,-63.9)
				..controls(127,-63.7)and(127.7,-63.4)..(128.4,-63.1)
				..controls(129,-62.8)and(129.7,-62.5)..(130.6,-62.4)
				..controls(132.1,-62.1)and(134,-61.5)..(135.2,-61)
				..controls(135.7,-60.8)and(136.7,-60.5)..(137.3,-60.3)
				..controls(138.4,-60)and(139.3,-59.7)..(141.1,-59)
				..controls(141.6,-58.7)and(142,-58.6)..(142.5,-58.6)
				..controls(143.2,-58.6)and(143.3,-58.7)..(143.6,-59)
				..controls(144,-59.5)and(144.1,-59.9)..(144,-60.8)
				..controls(144,-61.2)and(144,-71.1)..(144,-82.8)--(143.9,-104.1)--(143.7,-104.5)
				..controls(143.5,-104.7)and(143.3,-104.8)..(143,-104.9)
				..controls(142.6,-104.9)and(142.5,-104.8)..(141.5,-103.7)--cycle
				(27.8,-91.4)..controls(27.1,-90.9)and(26.3,-90)..(26.1,-89.5)
				..controls(26,-89.2)and(25.8,-88.6)..(25.7,-87.9)
				..controls(25.4,-86.2)and(25.6,-84.8)..(26.5,-83.5)
				..controls(27.4,-82.3)and(28.3,-82)..(29.4,-82.3)
				..controls(31.1,-82.8)and(32.1,-84.9)..(32,-87.4)
				..controls(31.9,-89.3)and(31.2,-90.6)..(29.9,-91.4)
				..controls(29.1,-91.9)and(28.6,-91.9)..(27.8,-91.4)--cycle
				(190.1,-91.3)..controls(189.7,-91.2)and(189.3,-90.9)..(189.1,-90.7)
				..controls(188.3,-89.8)and(187.6,-87.7)..(187.6,-86.2)
				..controls(187.6,-85.5)and(187.7,-85.2)..(188.2,-84.4)
				..controls(189.1,-82.6)and(190.1,-81.8)..(191.3,-82.1)
				..controls(192,-82.3)and(193.5,-83.7)..(193.8,-84.4)
				..controls(194.1,-85.1)and(194.2,-87.1)..(194,-88.3)
				..controls(193.8,-89.8)and(193.3,-90.5)..(192.3,-91.3)
				..controls(191.7,-91.7)and(190.9,-91.7)..(190.1,-91.3)--cycle
				(173.7,-83.4)..controls(172.2,-82.5)and(172,-82.5)..(170.6,-82.4)
				..controls(169.7,-82.4)and(168.8,-82.4)..(168.3,-82.5)
				..controls(167.2,-82.7)and(166.9,-82.7)..(166.6,-82.2)
				..controls(166.3,-81.8)and(166.3,-81)..(166.8,-80)
				..controls(167.2,-78.9)and(167.4,-77.4)..(167.7,-74.1)
				..controls(167.8,-72.8)and(167.9,-71.8)..(167.9,-71.7)
				..controls(168,-71.7)and(168,-70.2)..(168,-68.4)..controls(168,-65.4)and(168,-64.9)..(167.7,-63)
				..controls(167.3,-60.5)and(167.1,-59.9)..(166.4,-58.4)
				..controls(166.1,-57.7)and(165.8,-56.9)..(165.7,-56.6)
				..controls(165.7,-56.2)and(165.4,-55.6)..(165.2,-55.2)
				..controls(164.6,-53.8)and(161.9,-50.7)..(161,-50.3)
				..controls(159.8,-49.7)and(157.9,-50)..(155.2,-51.2)
				..controls(154.7,-51.4)and(153.4,-51.8)..(152.2,-52.1)
				..controls(150.8,-52.4)and(149.4,-52.9)..(148.6,-53.2)
				..controls(147.8,-53.6)and(146.7,-54)..(146,-54.2)
				..controls(145.4,-54.4)and(144.5,-54.7)..(144,-54.9)
				..controls(143.6,-55.2)and(142.7,-55.5)..(141.8,-55.7)
				..controls(141,-55.9)and(139.6,-56.3)..(138.7,-56.6)
				..controls(137.8,-57)and(136.3,-57.4)..(135.3,-57.6)
				..controls(133.8,-57.9)and(133.2,-58.1)..(132.1,-58.7)
				..controls(131.1,-59.2)and(130.3,-59.4)..(129.4,-59.6)
				..controls(128.7,-59.8)and(127.8,-60)..(127.4,-60.1)
				..controls(127,-60.3)and(126,-60.5)..(125.1,-60.7)
				..controls(123.2,-61.1)and(121.7,-61.5)..(120.2,-62.1)
				..controls(119.7,-62.3)and(118.1,-62.8)..(116.7,-63.2)
				..controls(114.6,-63.9)and(114,-64.1)..(113,-64.7)
				..controls(111.9,-65.4)and(111.8,-65.4)..(111,-65.5)
				..controls(109.9,-65.5)and(109.3,-65.4)..(107.7,-64.7)
				..controls(107,-64.4)and(106,-64)..(105.3,-63.9)
				..controls(103.5,-63.5)and(102.2,-63)..(101.2,-62.5)
				..controls(100.1,-61.8)and(99.3,-61.5)..(97.8,-61.1)
				..controls(97.2,-61)and(96.2,-60.7)..(95.7,-60.4)
				..controls(95.1,-60.1)and(94.2,-59.8)..(93.5,-59.6)
				..controls(92.9,-59.4)and(91.7,-59)..(91,-58.7)..controls(90.2,-58.4)and(88.9,-58)..(88,-57.7)
				..controls(87.2,-57.5)and(86.1,-57.1)..(85.5,-56.8)
				..controls(85,-56.6)and(84.1,-56.2)..(83.5,-56)..controls(82.9,-55.7)and(81.8,-55.2)..(81,-54.9)
				..controls(80.1,-54.5)and(79.1,-54.2)..(78.7,-54.1)
				..controls(77.9,-53.8)and(77.3,-53.6)..(76.2,-52.9)
				..controls(75.7,-52.6)and(75.2,-52.4)..(74.1,-52.2)
				..controls(72.5,-51.9)and(72.4,-51.8)..(71.3,-51.1)
				..controls(69.9,-50.1)and(69.4,-49.9)..(67.9,-49.8)
				..controls(67,-49.7)and(66.6,-49.6)..(66.3,-49.4)
				..controls(66.1,-49.3)and(65.7,-49)..(65.4,-48.8)
				..controls(65.1,-48.7)and(64.2,-48.4)..(63.6,-48.2)
				..controls(62.9,-48)and(61.7,-47.6)..(60.9,-47.2)
				..controls(60.1,-46.8)and(58.9,-46.4)..(58.2,-46.2)
				..controls(57.6,-46)and(56.8,-45.7)..(56.5,-45.5)
				..controls(54.9,-44.6)and(54.6,-44.5)..(53.2,-44.2)
				..controls(51.7,-43.8)and(51,-43.6)..(49.9,-42.9)
				..controls(49.5,-42.6)and(48.9,-42.4)..(47.5,-42)
				..controls(45.7,-41.6)and(45.6,-41.5)..(44.3,-40.6)
				..controls(43.5,-40.1)and(42.8,-39.6)..(42.7,-39.4)
				..controls(42.4,-38.7)and(43,-38.1)..(44.5,-37.6)..controls(44.9,-37.5)and(46,-37)..(46.9,-36.6)
				..controls(47.9,-36.1)and(48.8,-35.8)..(49.6,-35.6)
				..controls(50.5,-35.4)and(51.3,-35.2)..(52.4,-34.6)
				..controls(53.3,-34.1)and(54.5,-33.7)..(54.9,-33.5)
				..controls(55.4,-33.4)and(56.4,-33)..(57.2,-32.6)
				..controls(58.6,-32)and(60.1,-31.5)..(61.6,-31.2)
				..controls(62.3,-31.1)and(62.7,-31)..(63.2,-30.7)
				..controls(64.1,-30)and(64.3,-29.9)..(65.3,-29.7)
				..controls(65.7,-29.6)and(66.5,-29.2)..(67.1,-28.9)
				..controls(67.6,-28.6)and(68.8,-28.2)..(69.6,-27.9)
				..controls(70.5,-27.7)and(71.6,-27.3)..(72.1,-27)
				..controls(72.6,-26.7)and(73.8,-26.2)..(74.8,-25.9)
				..controls(75.7,-25.5)and(76.7,-25.1)..(77.1,-25)
				..controls(77.4,-24.8)and(78.5,-24.4)..(79.6,-24.1)
				..controls(80.6,-23.8)and(82,-23.2)..(82.8,-22.9)
				..controls(83.5,-22.5)and(84.8,-22)..(85.8,-21.8)
				..controls(86.7,-21.5)and(87.8,-21.1)..(88.1,-20.9)
				..controls(88.4,-20.7)and(89.3,-20.3)..(90.1,-20.1)
				..controls(90.8,-19.8)and(91.9,-19.4)..(92.4,-19.1)
				..controls(93.5,-18.5)and(94.4,-18.2)..(95.4,-18.1)
				..controls(95.7,-18)and(96.4,-17.8)..(96.9,-17.5)
				..controls(97.5,-17.3)and(98.4,-16.9)..(99,-16.6)
				..controls(99.7,-16.3)and(100.9,-15.9)..(101.9,-15.6)
				..controls(102.9,-15.3)and(104.4,-14.7)..(105.3,-14.3)
				..controls(108.1,-13.2)and(108.5,-13.1)..(110.1,-13.1)
				..controls(111.4,-13.1)and(111.5,-13.1)..(112.4,-13.5)
				..controls(113.2,-14)and(115,-14.8)..(116.6,-15.4)
				..controls(116.9,-15.5)and(117.9,-15.8)..(118.8,-16.1)
				..controls(120.9,-16.6)and(121.5,-16.9)..(122.6,-17.9)
				..controls(123,-18.3)and(123.7,-18.8)..(123.9,-18.9)
				..controls(124.6,-19.4)and(126,-19.8)..(127.9,-20.2)
				..controls(129.6,-20.5)and(129.8,-20.6)..(131.1,-21.4)
				..controls(131.5,-21.7)and(132.4,-22.1)..(133.1,-22.3)
				..controls(136.1,-23.2)and(138.3,-24.1)..(140.2,-25)
				..controls(141,-25.5)and(141.5,-25.6)..(142.5,-25.8)
				..controls(145.1,-26.3)and(146,-26.6)..(149.4,-28.1)
				..controls(151.3,-28.9)and(153.5,-29.8)..(154.3,-30.1)
				..controls(155.1,-30.4)and(156.2,-30.8)..(156.6,-31)
				..controls(157.1,-31.2)and(158.2,-31.6)..(159.1,-31.8)
				..controls(160.6,-32.2)and(160.9,-32.3)..(161.7,-32.9)
				..controls(162.4,-33.3)and(162.9,-33.6)..(163.7,-33.8)
				..controls(164.3,-34)and(164.8,-34.1)..(164.9,-34.1)
				..controls(165,-34.1)and(165.6,-34.4)..(166.4,-34.8)
				..controls(167.3,-35.2)and(168.1,-35.4)..(169,-35.6)
				..controls(170,-35.8)and(170.5,-36)..(171.7,-36.6)
				..controls(172.4,-37)and(173.3,-37.4)..(173.6,-37.4)
				..controls(175.6,-37.7)and(176.2,-38)..(176.7,-39.1)
				..controls(176.9,-39.6)and(176.9,-39.7)..(176.7,-40.4)
				..controls(176.4,-41.6)and(175.8,-42.7)..(175.2,-43.4)
				..controls(174.5,-44.1)and(172.6,-45.5)..(171.7,-45.8)
				..controls(170.6,-46.2)and(169.1,-46)..(168.2,-45.4)
				..controls(167.9,-45.2)and(167.3,-44.8)..(166.9,-44.4)
				..controls(166.4,-44.1)and(165.8,-43.6)..(165.4,-43.2)
				..controls(165.1,-42.9)and(164.4,-42.3)..(163.9,-41.9)
				..controls(163.3,-41.5)and(162.7,-40.9)..(162.4,-40.5)
				..controls(161.4,-39.2)and(159.9,-38)..(158.8,-37.8)
				..controls(158.4,-37.8)and(158.3,-37.7)..(158.2,-37.2)
				..controls(158,-36.8)and(157.4,-36.1)..(157.1,-36.1)
				..controls(157,-36.1)and(156.7,-36.1)..(156.5,-36)
				..controls(156.1,-35.9)and(155.9,-35.7)..(155.7,-35.2)
				..controls(155.3,-34.6)and(155.1,-34.4)..(154,-34.1)
				..controls(153.6,-34)and(153.2,-33.7)..(152.9,-33.5)
				..controls(152.4,-32.9)and(151.6,-32.4)..(151.3,-32.4)
				..controls(150.9,-32.4)and(150.4,-32.6)..(150.4,-32.9)
				..controls(150.4,-33.2)and(150.8,-33.4)..(151.4,-33.6)
				..controls(152.1,-33.7)and(152.8,-34.1)..(153.5,-34.9)
				..controls(154.1,-35.5)and(154.3,-35.6)..(155,-35.8)
				..controls(155.7,-36)and(155.7,-36)..(156,-36.6)
				..controls(156.3,-37.3)and(156.7,-37.7)..(157.5,-37.8)
				..controls(158.1,-37.9)and(158.2,-38)..(158.4,-38.8)
				..controls(158.5,-39.4)and(158.6,-39.6)..(159.6,-40.4)
				..controls(160.2,-41)and(160.9,-41.6)..(161.3,-41.8)
				..controls(161.7,-42)and(161.9,-42.3)..(162.2,-42.8)
				..controls(162.4,-43.2)and(162.8,-43.8)..(163.1,-44.1)
				..controls(163.9,-44.8)and(165.1,-46.7)..(165.7,-48)
				..controls(166.2,-49.3)and(166.4,-49.7)..(167.2,-50.5)
				..controls(167.5,-50.8)and(167.9,-51.4)..(168.1,-51.8)
				..controls(168.3,-52.2)and(168.8,-53)..(169.1,-53.6)
				..controls(169.6,-54.3)and(169.8,-54.9)..(169.9,-55.6)
				..controls(170.4,-57.8)and(170.5,-58.1)..(171.2,-59.5)
				..controls(171.6,-60.3)and(172,-61.5)..(172.2,-62.1)
				..controls(172.3,-62.8)and(172.7,-63.7)..(172.9,-64.1)
				..controls(173.2,-64.6)and(173.4,-65.3)..(173.4,-65.6)
				..controls(173.8,-68.8)and(174.2,-70.6)..(174.9,-72.6)
				..controls(175.5,-74.5)and(175.9,-77.1)..(176.1,-80.2)
				..controls(176.3,-83.4)and(175.5,-84.4)..(173.7,-83.4)--cycle;
				\definecolor{ezs}{RGB}{247,179,26}\fill[ezs](-5.7,-156.6)--(-3.4,-157.5)--(61.9,12.7)--(59.7,13.6)--cycle;
				\definecolor{ezs}{RGB}{247,179,22}\fill[ezs](-3.8,-157.3)--(-1.5,-158.2)--(63.9,12)--(61.6,12.9)--cycle;
				\definecolor{ezs}{RGB}{248,179,19}\fill[ezs](-1.8,-158.1)--(0.4,-158.9)--(65.8,11.2)--(63.5,12.1)--cycle;
				\definecolor{ezs}{RGB}{248,178,16}\fill[ezs](0.1,-158.8)--(2.4,-159.7)--(67.7,10.5)--(65.4,11.4)--cycle;
				\definecolor{ezs}{RGB}{249,178,15}\fill[ezs](2,-159.5)--(4.3,-160.4)--(69.6,9.8)--(67.3,10.7)--cycle;
				\definecolor{ezs}{RGB}{251,177,17}\fill[ezs](3.9,-160.3)--(6.2,-161.1)--(71.5,9)--(69.2,9.9)--cycle;
				\definecolor{ezs}{RGB}{252,179,18}\fill[ezs](5.8,-161)--(8.1,-161.9)--(73.4,8.3)--(71.1,9.2)--cycle;
				\definecolor{ezs}{RGB}{252,182,17}\fill[ezs](7.7,-161.7)--(10,-162.6)--(75.3,7.6)--(73,8.5)--cycle;
				\definecolor{ezs}{RGB}{252,185,17}\fill[ezs](9.6,-162.5)--(11.9,-163.4)--(77.2,6.8)--(75,7.7)--cycle;
				\definecolor{ezs}{RGB}{251,189,18}\fill[ezs](11.5,-163.2)--(13.8,-164.1)--(79.2,6.1)--(76.9,7)--cycle;
				\definecolor{ezs}{RGB}{251,191,22}\fill[ezs](13.4,-163.9)--(15.7,-164.8)--(81.1,5.4)--(78.8,6.2)--cycle;
				\definecolor{ezs}{RGB}{252,193,25}\fill[ezs](15.4,-164.7)--(17.6,-165.6)--(83,4.6)--(80.7,5.5)--cycle;
				\definecolor{ezs}{RGB}{253,195,27}\fill[ezs](17.3,-165.4)--(19.6,-166.3)--(84.9,3.9)--(82.6,4.8)--cycle;
				\definecolor{ezs}{RGB}{253,196,31}\fill[ezs](19.2,-166.1)--(21.5,-167)--(86.8,3.2)--(84.5,4)--cycle;
				\definecolor{ezs}{RGB}{254,199,35}\fill[ezs](21.1,-166.9)--(23.4,-167.8)--(88.7,2.4)--(86.4,3.3)--cycle;
				\definecolor{ezs}{RGB}{253,201,38}\fill[ezs](23,-167.6)--(25.3,-168.5)--(90.6,1.7)--(88.3,2.6)--cycle;
				\definecolor{ezs}{RGB}{253,203,40}\fill[ezs](24.9,-168.3)--(27.2,-169.2)--(92.5,1)--(90.3,1.8)--cycle;
				\definecolor{ezs}{RGB}{253,205,42}\fill[ezs](26.8,-169.1)--(29.1,-170)--(94.4,0.2)--(92.2,1.1)--cycle;
				\definecolor{ezs}{RGB}{253,207,45}\fill[ezs](28.7,-169.8)--(31,-170.7)--(96.4,-0.5)--(94.1,0.4)--cycle;
				\definecolor{ezs}{RGB}{253,210,49}\fill[ezs](30.7,-170.5)--(32.9,-171.4)--(98.3,-1.2)--(96,-0.4)--cycle;
				\definecolor{ezs}{RGB}{253,212,54}\fill[ezs](32.6,-171.3)--(34.9,-172.2)--(100.2,-2)--(97.9,-1.1)--cycle;
				\definecolor{ezs}{RGB}{253,214,59}\fill[ezs](34.5,-172)--(36.8,-172.9)--(102.1,-2.7)--(99.8,-1.8)--cycle;
				\definecolor{ezs}{RGB}{253,216,63}\fill[ezs](36.4,-172.7)--(38.7,-173.6)--(104,-3.4)--(101.7,-2.6)--cycle;
				\definecolor{ezs}{RGB}{253,218,65}\fill[ezs](38.3,-173.5)--(40.6,-174.4)--(105.9,-4.2)--(103.6,-3.3)--cycle;
				\definecolor{ezs}{RGB}{254,220,66}\fill[ezs](40.2,-174.2)--(42.5,-175.1)--(107.8,-4.9)--(105.5,-4)--cycle;
				\definecolor{ezs}{RGB}{254,222,68}\fill[ezs](42.1,-175)--(44.4,-175.8)--(109.7,-5.6)--(107.5,-4.8)--cycle;
				\definecolor{ezs}{RGB}{254,224,71}\fill[ezs](44,-175.7)--(46.3,-176.6)--(111.7,-6.4)--(109.4,-5.5)--cycle;
				\definecolor{ezs}{RGB}{254,226,75}\fill[ezs](46,-176.4)--(48.2,-177.3)--(113.6,-7.1)--(111.3,-6.2)--cycle;
				\definecolor{ezs}{RGB}{254,228,79}\fill[ezs](47.9,-177.2)--(50.1,-178)--(115.5,-7.8)--(113.2,-7)--cycle;
				\definecolor{ezs}{RGB}{254,230,82}\fill[ezs](49.8,-177.9)--(52.1,-178.8)--(117.4,-8.6)--(115.1,-7.7)--cycle;
				\definecolor{ezs}{RGB}{254,232,85}\fill[ezs](51.7,-178.6)--(54,-179.5)--(119.3,-9.3)--(117,-8.4)--cycle;
				\definecolor{ezs}{RGB}{254,234,88}\fill[ezs](53.6,-179.4)--(55.9,-180.2)--(121.2,-10)--(118.9,-9.2)--cycle;
				\definecolor{ezs}{RGB}{254,237,91}\fill[ezs](55.5,-180.1)--(57.8,-181)--(123.1,-10.8)--(120.8,-9.9)--cycle;
				\definecolor{ezs}{RGB}{254,239,94}\fill[ezs](57.4,-180.8)--(59.7,-181.7)--(125,-11.5)--(122.8,-10.6)--cycle;
				\definecolor{ezs}{RGB}{254,241,97}\fill[ezs](59.3,-181.6)--(61.6,-182.4)--(127,-12.2)--(124.7,-11.4)--cycle;
				\definecolor{ezs}{RGB}{254,243,101}\fill[ezs](61.2,-182.3)--(63.5,-183.2)--(128.9,-13)--(126.6,-12.1)--cycle;
				\definecolor{ezs}{RGB}{254,245,103}\fill[ezs](63.2,-183)--(65.4,-183.9)--(130.8,-13.7)--(128.5,-12.8)--cycle;
				\definecolor{ezs}{RGB}{254,245,105}\fill[ezs](65.1,-183.8)--(67.4,-184.6)--(132.7,-14.4)--(130.4,-13.6)--cycle;
				\definecolor{ezs}{RGB}{254,245,105}\fill[ezs](67,-184.5)--(69.3,-185.4)--(134.6,-15.2)--(132.3,-14.3)--cycle;
				\definecolor{ezs}{RGB}{254,244,103}\fill[ezs](68.9,-185.2)--(71.2,-186.1)--(136.5,-15.9)--(134.2,-15)--cycle;
				\definecolor{ezs}{RGB}{253,243,101}\fill[ezs](70.8,-186)--(73.1,-186.8)--(138.4,-16.6)--(136.1,-15.8)--cycle;
				\definecolor{ezs}{RGB}{253,241,98}\fill[ezs](72.7,-186.7)--(75,-187.6)--(140.3,-17.4)--(138,-16.5)--cycle;
				\definecolor{ezs}{RGB}{253,239,95}\fill[ezs](74.6,-187.4)--(76.9,-188.3)--(142.2,-18.1)--(140,-17.2)--cycle;
				\definecolor{ezs}{RGB}{253,236,92}\fill[ezs](76.5,-188.2)--(78.8,-189)--(144.2,-18.8)--(141.9,-18)--cycle;
				\definecolor{ezs}{RGB}{253,234,88}\fill[ezs](78.5,-188.9)--(80.7,-189.8)--(146.1,-19.6)--(143.8,-18.7)--cycle;
				\definecolor{ezs}{RGB}{253,232,85}\fill[ezs](80.4,-189.6)--(82.7,-190.5)--(148,-20.3)--(145.7,-19.4)--cycle;
				\definecolor{ezs}{RGB}{253,230,82}\fill[ezs](82.3,-190.4)--(84.6,-191.2)--(149.9,-21)--(147.6,-20.2)--cycle;
				\definecolor{ezs}{RGB}{253,228,78}\fill[ezs](84.2,-191.1)--(86.5,-192)--(151.8,-21.8)--(149.5,-20.9)--cycle;
				\definecolor{ezs}{RGB}{253,226,74}\fill[ezs](86.1,-191.8)--(88.4,-192.7)--(153.7,-22.5)--(151.4,-21.6)--cycle;
				\definecolor{ezs}{RGB}{253,224,70}\fill[ezs](88,-192.6)--(90.3,-193.4)--(155.6,-23.3)--(153.3,-22.4)--cycle;
				\definecolor{ezs}{RGB}{253,222,67}\fill[ezs](89.9,-193.3)--(92.2,-194.2)--(157.5,-24)--(155.3,-23.1)--cycle;
				\definecolor{ezs}{RGB}{253,220,65}\fill[ezs](91.8,-194)--(94.1,-194.9)--(159.5,-24.7)--(157.2,-23.8)--cycle;
				\definecolor{ezs}{RGB}{253,217,63}\fill[ezs](93.7,-194.8)--(96,-195.6)--(161.4,-25.5)--(159.1,-24.6)--cycle;
				\definecolor{ezs}{RGB}{253,215,60}\fill[ezs](95.7,-195.5)--(97.9,-196.4)--(163.3,-26.2)--(161,-25.3)--cycle;
				\definecolor{ezs}{RGB}{253,213,56}\fill[ezs](97.6,-196.2)--(99.9,-197.1)--(165.2,-26.9)--(162.9,-26)--cycle;
				\definecolor{ezs}{RGB}{252,211,53}\fill[ezs](99.5,-197)--(101.8,-197.8)--(167.1,-27.7)--(164.8,-26.8)--cycle;
				\definecolor{ezs}{RGB}{252,209,49}\fill[ezs](101.4,-197.7)--(103.7,-198.6)--(169,-28.4)--(166.7,-27.5)--cycle;
				\definecolor{ezs}{RGB}{252,207,45}\fill[ezs](103.3,-198.4)--(105.6,-199.3)--(170.9,-29.1)--(168.6,-28.2)--cycle;
				\definecolor{ezs}{RGB}{253,205,42}\fill[ezs](105.2,-199.2)--(107.5,-200)--(172.8,-29.9)--(170.6,-29)--cycle;
				\definecolor{ezs}{RGB}{253,202,40}\fill[ezs](107.1,-199.9)--(109.4,-200.8)--(174.7,-30.6)--(172.5,-29.7)--cycle;
				\definecolor{ezs}{RGB}{253,200,38}\fill[ezs](109,-200.6)--(111.3,-201.5)--(176.7,-31.3)--(174.4,-30.4)--cycle;
				\definecolor{ezs}{RGB}{253,198,35}\fill[ezs](111,-201.4)--(113.2,-202.2)--(178.6,-32.1)--(176.3,-31.2)--cycle;
				\definecolor{ezs}{RGB}{253,196,31}\fill[ezs](112.9,-202.1)--(115.2,-203)--(180.5,-32.8)--(178.2,-31.9)--cycle;
				\definecolor{ezs}{RGB}{253,194,28}\fill[ezs](114.8,-202.8)--(117.1,-203.7)--(182.4,-33.5)--(180.1,-32.6)--cycle;
				\definecolor{ezs}{RGB}{253,192,26}\fill[ezs](116.7,-203.6)--(119,-204.5)--(184.3,-34.3)--(182,-33.4)--cycle;
				\definecolor{ezs}{RGB}{253,189,24}\fill[ezs](118.6,-204.3)--(120.9,-205.2)--(186.2,-35)--(183.9,-34.1)--cycle;
				\definecolor{ezs}{RGB}{253,186,20}\fill[ezs](120.5,-205)--(122.8,-205.9)--(188.1,-35.7)--(185.8,-34.8)--cycle;
				\definecolor{ezs}{RGB}{253,184,17}\fill[ezs](122.4,-205.8)--(124.7,-206.7)--(190,-36.5)--(187.8,-35.6)--cycle;
				\definecolor{ezs}{RGB}{252,183,15}\fill[ezs](124.3,-206.5)--(126.6,-207.4)--(192,-37.2)--(189.7,-36.3)--cycle;
				\definecolor{ezs}{RGB}{252,181,14}\fill[ezs](126.3,-207.2)--(128.5,-208.1)--(193.9,-37.9)--(191.6,-37.1)--cycle;
				\definecolor{ezs}{RGB}{252,179,13}\fill[ezs](128.2,-208)--(130.4,-208.9)--(195.8,-38.7)--(193.5,-37.8)--cycle;
				\definecolor{ezs}{RGB}{252,178,12}\fill[ezs](130.1,-208.7)--(132.4,-209.6)--(197.7,-39.4)--(195.4,-38.5)--cycle;
				\definecolor{ezs}{RGB}{251,178,10}\fill[ezs](132,-209.4)--(134.3,-210.3)--(199.6,-40.1)--(197.3,-39.3)--cycle;
				\definecolor{ezs}{RGB}{251,178,10}\fill[ezs](133.9,-210.2)--(136.2,-211.1)--(201.5,-40.9)--(199.2,-40)--cycle;
				\definecolor{ezs}{RGB}{250,177,11}\fill[ezs](135.8,-210.9)--(138.1,-211.8)--(203.4,-41.6)--(201.1,-40.7)--cycle;
				\definecolor{ezs}{RGB}{250,178,11}\fill[ezs](137.7,-211.6)--(140,-212.5)--(205.3,-42.3)--(203.1,-41.5)--cycle;
				\definecolor{ezs}{RGB}{248,180,10}\fill[ezs](139.6,-212.4)--(141.9,-213.3)--(207.3,-43.1)--(205,-42.2)--cycle;
				\definecolor{ezs}{RGB}{247,181,9}\fill[ezs](141.5,-213.1)--(143.8,-214)--(209.2,-43.8)--(206.9,-42.9)--cycle;
				\definecolor{ezs}{RGB}{246,182,9}\fill[ezs](143.5,-213.8)--(145.7,-214.7)--(211.1,-44.5)--(208.8,-43.7)--cycle;
				\definecolor{ezs}{RGB}{246,182,10}\fill[ezs](145.4,-214.6)--(147.7,-215.5)--(213,-45.3)--(210.7,-44.4)--cycle;
			\end{scope}
			% chữ EDUZY (Times đậm, làm dày nét bằng 5 lớp lệch nhẹ; căn đúng khung chữ logo gốc)
			\foreach \dx/\dy in {-0.7/0,0.7/0,0/-0.3,0/0.3,0/0}
			\node[anchor=center,inner sep=0pt,outer sep=0pt,text=ezText,align=center] at (109.6+\dx,-193.6+\dy)
			{\resizebox{1.01cm}{0.22cm}{\fontencoding{T1}\fontfamily{ptm}\fontseries{b}\selectfont EDUZY}};
		\end{tikzpicture}%
	}%
}
%====================== THIẾT LẬP THÂN ĐỀ ======================
%--- Môi trường tự luận {bt}: ex_test 3.3.4 chưa khai báo sẵn nên tạo ở đây
%    (ex_test đã tự móc \fixshowans và \listenumerate cho bt)
\makeatletter
\@ifundefined{bt}{\newtheorem{bt}{\color{gold1}\namebt}}{}
\makeatother
\renewcommand{\sectionmark}[1]{\markright{\thesection\ #1}}
%--- Khung tiêu đề (logo + tối đa 4 dòng; dòng 3, 4 để trống {} thì bỏ qua)
%  Đề thi:    \khungtieude{THI THỬ THPT QUỐC GIA NĂM 2026}{Môn Toán -- Đề số 21}
%                         {(Đề thi gồm: ...)}{Thời gian làm bài: 90 phút (không kể thời gian phát đề).}
%  Chuyên đề: \khungtieude{CHUYÊN ĐỀ 3 -- DÃY SỐ VÀ CÔNG THỨC TRUY HỒI}
%                         {Bồi dưỡng học sinh giỏi Toán THPT}{...}{}
\newcommand{\khungtieude}[4]{%
	\begin{tcolorbox}[colframe=gold1,colback=gold5!45,boxrule=0.5mm]
		\begin{minipage}[c]{3cm}
			\centerline{\eduzylogo[2.4cm]}
		\end{minipage}\hspace{0.5cm}
		\begin{minipage}[c]{12.5cm}
			\vskip 0.15 cm
			\centerline{\ind{#1}}
			\vskip 0.15 cm
			\centerline{\textbf{#2}}
			\if\relax\detokenize{#3}\relax\else
				\vskip 0.15 cm
				\centerline{\textit{#3}}
			\fi
			\if\relax\detokenize{#4}\relax\else
				\vskip 0.15 cm
				\centerline{\textit{#4}}
			\fi
			\vskip 0.15 cm
		\end{minipage}
	\end{tcolorbox}
}
%====================== BỔ SUNG CHO BÀI HỌC, CHUYÊN ĐỀ ======================
\usepackage[shortlabels]{enumitem}
\usepackage{tabularx}
%--- Danh sách gọn: \begin{enumerate}[gon] ... \end{enumerate}
\SetEnumitemKey{gon}{leftmargin=*,itemsep=1pt,parsep=0pt,topsep=3pt}
%--- Khung lý thuyết: \begin{khung}{} ... \end{khung}; tiêu đề chỉ dùng khi thật ngắn ("Mẹo", "Lỗi thường gặp")
\newtcolorbox{khung}[1]{breakable,enhanced,colframe=gold1,colback=gold5!45,
	boxrule=0.4mm,fonttitle=\bfseries\sffamily,coltitle=white,colbacktitle=gold1,
	attach boxed title to top left={xshift=6mm,yshift=-3mm},
	boxed title style={arc=1mm,boxrule=0.3mm},title={#1}}
%--- Tiêu đề mục nhỏ trong một dạng: \muc{Dấu hiệu}, \muc{Phương pháp}, \muc{Bài tự luyện}
\makeatletter
\newcommand{\muc}[1]{\par\medskip\noindent\indam{#1}\par\nobreak\smallskip\@afterindentfalse\@afterheading}
\makeatother
%--- Ô công thức cần nhớ (căn giữa, viền vàng, vừa khít nội dung)
%     \begin{center}\congthuc{$Q(n)=(n-1)!$}\end{center}
\newtcbox{\congthuc}{on line,colframe=gold1,colback=gold5!45,boxrule=0.5pt,arc=2pt,
	left=8pt,right=8pt,top=4pt,bottom=4pt}
%--- Công tắc lời giải bài tự luyện (ví dụ mẫu luôn dùng \loigiai)
%     \loigiaiTL[{đáp số}]{lời giải đầy đủ}
%     \TLfull -> in lời giải đầy đủ;  \TLdapso -> chỉ in một dòng "Đáp số";  \TLan -> ẩn hoàn toàn
%     (\CheDo là dethi thì mọi lời giải và đáp số đều ẩn)
\makeatletter
\newif\ifTLhien   \TLhientrue
\newif\ifTLdapso  \TLdapsofalse
\newcommand{\TLfull}{\TLhientrue\TLdapsofalse}
\newcommand{\TLdapso}{\TLhientrue\TLdapsotrue}
\newcommand{\TLan}{\TLhienfalse}
\newcommand{\loigiaiTL}[2][]{%
	\ifTLhien
		\ifTLdapso
			\def\@tlarg{#1}%
			\ifx\@tlarg\@empty\else
				\ifdefstring{\kindOPTN}{dethi}{}{\par\noindent{\bfseries\color{gold1}Đáp số:} #1\par}%
			\fi
		\else
			\loigiai{#2}%
		\fi
	\fi}
\makeatother
\TLdapso% mặc định: chỉ in đáp số bài tự luyện
%====================== BẢNG ĐÚNG/SAI (\choiceTFt) ======================
%  Hàng tiêu đề nền nâu đồng chữ trắng (Phát biểu | Đ | S), các hàng tô xen kẽ kem / trắng,
%  đường kẻ vàng kim, nhãn a) b) c) d) màu nâu đồng. Bản có đáp án: cột được chọn có vòng tròn
%  nâu đồng với dấu tích trắng (câu đúng tích ở cột Đúng, câu sai tích ở cột Sai).
%  Tiêu đề ghi đủ "Đúng | Sai": \OPTN{viettat=0} sau \input (cột Đúng/Sai rộng ra, phát biểu hẹp lại).
%  Dùng bảng gốc của ex_test: \def\BangDSgoc{1} trước \input{preamble-eduzy-v2}.
\providecommand{\BangDSgoc}{0}
\ifnum\BangDSgoc=0
\usepackage{colortbl,pifont}
\colorlet{ezDSke}{gold2!60!gold1}% đường kẻ bảng (nét 0,4pt như ex_test để không đổi chiều cao bảng)
\colorlet{ezDSle}{gold5!60}% hàng a), c)
\colorlet{ezDSchan}{white}% hàng b), d)
%--- dấu chọn ở cột Đúng/Sai
%    (bề rộng 0, căn giữa ô: không làm đổi bề rộng cột của ex_test, không báo tràn lề)
\newcommand{\EZdauchon}{\makebox[0pt]{\tikz[baseline=-1ex]\node[circle,fill=gold1,inner sep=0pt,
	minimum size=1.2em,text=white,font=\small]{\ding{51}};}}
\def\tickT{\EZdauchon}\def\tickF{\EZdauchon}
\AtBeginDocument{\sbox{\ccTFTab}{\EZdauchon}}
%--- nhãn a) b) c) d) màu nâu đồng (ở bản lời giải, nhãn câu đúng vẫn đóng khung như ex_test)
\renewcommand{\FalseTF}{\stepcounter{dapan}\textbf{\color{gold1}\DapAnTF\sepTF}\ignorespaces}
%--- trước khi dựng bảng: kẻ màu vàng kim (bề rộng cột giữ nguyên như ex_test); sau bảng trả lại màu đen
\newcount\EZhangTF
\newcommand{\EZbatdauTF}{%
	\global\EZhangTF=0\relax
	\arrayrulecolor{ezDSke}}
\newcommand{\EZketthucTF}{\arrayrulecolor{black}}
%--- hàng tiêu đề nền nâu đồng, chữ trắng
\makeatletter
\newcommand{\EZo@TF}[1]{\centering\bfseries\color{white}#1}
\newcommand{\EZtieudeTF}{%
	\ifcase\cotDS
		\ifnum\cotSTT=1\relax
			\def\rowHeader{\hline\rowcolor{gold1}\EZo@TF{TT}&\EZo@TF{\phatbieu}\tabularnewline}%
		\else
			\def\rowHeader{\hline\rowcolor{gold1}\EZo@TF{\phatbieu}\tabularnewline}%
		\fi
	\or
		\ifnum\cotSTT=1\relax
			\def\rowHeader{\hline\rowcolor{gold1}\EZo@TF{TT}&\EZo@TF{\phatbieu}&\EZo@TF{\Dung}\tabularnewline}%
		\else
			\def\rowHeader{\hline\rowcolor{gold1}\EZo@TF{\phatbieu}&\EZo@TF{\Dung}\tabularnewline}%
		\fi
	\or
		\ifnum\cotSTT=1\relax
			\def\rowHeader{\hline\rowcolor{gold1}\EZo@TF{TT}&\EZo@TF{\phatbieu}&\EZo@TF{\Dung}&\EZo@TF{\Sai}\tabularnewline}%
		\else
			\def\rowHeader{\hline\rowcolor{gold1}\EZo@TF{\phatbieu}&\EZo@TF{\Dung}&\EZo@TF{\Sai}\tabularnewline}%
		\fi
	\fi}
%--- tô xen kẽ từng hàng (gọi ngay sau \hline đầu mỗi hàng)
\newcommand{\EZmauhangTF}{%
	\noalign{\global\advance\EZhangTF\@ne}%
	\ifodd\EZhangTF\rowcolor{ezDSle}\else\rowcolor{ezDSchan}\fi}
%--- gắn vào \choiceTFt / \writeChoiceTFt của ex_test
\xpretocmd{\choiceTFt}{\EZbatdauTF}{}{\typeout{[preamble-eduzy-v2] Không gắn được bảng Đúng/Sai (đầu bảng)}}
\xpatchcmd{\choiceTFt}{\ifnum\dongPB=0\def\rowHeader{}\fi}{\EZtieudeTF\ifnum\dongPB=0\def\rowHeader{}\fi}
	{}{\typeout{[preamble-eduzy-v2] Không gắn được bảng Đúng/Sai (hàng tiêu đề)}}
\xpatchcmd{\writeChoiceTFt}{\hline\xdef\dotEX{\sdotEX}}{\hline\EZmauhangTF\xdef\dotEX{\sdotEX}}
	{}{\typeout{[preamble-eduzy-v2] Không gắn được bảng Đúng/Sai (tô hàng)}}
\xpatchcmd{\writeChoiceTFt}{{\end{longtable}}}{{\end{longtable}\EZketthucTF}}
	{}{\typeout{[preamble-eduzy-v2] Không gắn được bảng Đúng/Sai (cuối longtable)}}
\xpatchcmd{\writeChoiceTFt}{{\end{xtabular}\end{center}\medskip\noindent}}{{\end{xtabular}\end{center}\EZketthucTF\medskip\noindent}}
	{}{\typeout{[preamble-eduzy-v2] Không gắn được bảng Đúng/Sai (cuối xtabular)}}
\fi
%=====================================================================
\renewcommand{\parallel}{\mathrel{/\!\!/}}
.
\providecommand{\BangDSgoc}{0}
\ifnum\BangDSgoc=0
\usepackage{colortbl,pifont}
\colorlet{ezDSke}{gold2!60!gold1}% đường kẻ bảng (nét 0,4pt như ex_test để không đổi chiều cao bảng)
\colorlet{ezDSle}{gold5!60}% hàng a), c)
\colorlet{ezDSchan}{white}% hàng b), d)
%--- dấu chọn ở cột Đúng/Sai
%    (bề rộng 0, căn giữa ô: không làm đổi bề rộng cột của ex_test, không báo tràn lề)
\newcommand{\EZdauchon}{\makebox[0pt]{\tikz[baseline=-1ex]\node[circle,fill=gold1,inner sep=0pt,
	minimum size=1.2em,text=white,font=\small]{\ding{51}};}}
\def\tickT{\EZdauchon}\def\tickF{\EZdauchon}
\AtBeginDocument{\sbox{\ccTFTab}{\EZdauchon}}
%--- nhãn a) b) c) d) màu nâu đồng (ở bản lời giải, nhãn câu đúng vẫn đóng khung như ex_test)
\renewcommand{\FalseTF}{\stepcounter{dapan}\textbf{\color{gold1}\DapAnTF\sepTF}\ignorespaces}
%--- trước khi dựng bảng: kẻ màu vàng kim (bề rộng cột giữ nguyên như ex_test); sau bảng trả lại màu đen
\newcount\EZhangTF
\newcommand{\EZbatdauTF}{%
	\global\EZhangTF=0\relax
	\arrayrulecolor{ezDSke}}
\newcommand{\EZketthucTF}{\arrayrulecolor{black}}
%--- hàng tiêu đề nền nâu đồng, chữ trắng
\makeatletter
\newcommand{\EZo@TF}[1]{\centering\bfseries\color{white}#1}
\newcommand{\EZtieudeTF}{%
	\ifcase\cotDS
		\ifnum\cotSTT=1\relax
			\def\rowHeader{\hline\rowcolor{gold1}\EZo@TF{TT}&\EZo@TF{\phatbieu}\tabularnewline}%
		\else
			\def\rowHeader{\hline\rowcolor{gold1}\EZo@TF{\phatbieu}\tabularnewline}%
		\fi
	\or
		\ifnum\cotSTT=1\relax
			\def\rowHeader{\hline\rowcolor{gold1}\EZo@TF{TT}&\EZo@TF{\phatbieu}&\EZo@TF{\Dung}\tabularnewline}%
		\else
			\def\rowHeader{\hline\rowcolor{gold1}\EZo@TF{\phatbieu}&\EZo@TF{\Dung}\tabularnewline}%
		\fi
	\or
		\ifnum\cotSTT=1\relax
			\def\rowHeader{\hline\rowcolor{gold1}\EZo@TF{TT}&\EZo@TF{\phatbieu}&\EZo@TF{\Dung}&\EZo@TF{\Sai}\tabularnewline}%
		\else
			\def\rowHeader{\hline\rowcolor{gold1}\EZo@TF{\phatbieu}&\EZo@TF{\Dung}&\EZo@TF{\Sai}\tabularnewline}%
		\fi
	\fi}
%--- tô xen kẽ từng hàng (gọi ngay sau \hline đầu mỗi hàng)
\newcommand{\EZmauhangTF}{%
	\noalign{\global\advance\EZhangTF\@ne}%
	\ifodd\EZhangTF\rowcolor{ezDSle}\else\rowcolor{ezDSchan}\fi}
%--- gắn vào \choiceTFt / \writeChoiceTFt của ex_test
\xpretocmd{\choiceTFt}{\EZbatdauTF}{}{\typeout{[preamble-eduzy-v2] Không gắn được bảng Đúng/Sai (đầu bảng)}}
\xpatchcmd{\choiceTFt}{\ifnum\dongPB=0\def\rowHeader{}\fi}{\EZtieudeTF\ifnum\dongPB=0\def\rowHeader{}\fi}
	{}{\typeout{[preamble-eduzy-v2] Không gắn được bảng Đúng/Sai (hàng tiêu đề)}}
\xpatchcmd{\writeChoiceTFt}{\hline\xdef\dotEX{\sdotEX}}{\hline\EZmauhangTF\xdef\dotEX{\sdotEX}}
	{}{\typeout{[preamble-eduzy-v2] Không gắn được bảng Đúng/Sai (tô hàng)}}
\xpatchcmd{\writeChoiceTFt}{{\end{longtable}}}{{\end{longtable}\EZketthucTF}}
	{}{\typeout{[preamble-eduzy-v2] Không gắn được bảng Đúng/Sai (cuối longtable)}}
\xpatchcmd{\writeChoiceTFt}{{\end{xtabular}\end{center}\medskip\noindent}}{{\end{xtabular}\end{center}\EZketthucTF\medskip\noindent}}
	{}{\typeout{[preamble-eduzy-v2] Không gắn được bảng Đúng/Sai (cuối xtabular)}}
\fi
%=====================================================================
\renewcommand{\parallel}{\mathrel{/\!\!/}}

%      \TLdapso                      % tài liệu: \TLfull | \TLdapso | \TLan
%      \begin{document} ... \end{document}
%  Logo Eduzy vẽ bằng TikZ ngay trong file (\eduzylogo), không cần file ảnh.
%  Biên dịch: pdfLaTeX (máy có gói vietnam) hoặc XeLaTeX.
%============================================================
\documentclass[12pt,a4paper,twoside]{article}%<-- twoside: bat dau trang/chan trang chan-le
\usepackage{iftex}
\usepackage{amsmath,amssymb,tikz,mathrsfs,tkz-tab}
\usepackage{graphicx}
\usepackage{pgfplots}% đồ thị hàm số
\pgfplotsset{compat=1.17}
%\usepackage{mathptmx}
\usepackage{mathpazo}
\ifXeTeX% XeLaTeX: font chữ có đủ dấu tiếng Việt, giữ nguyên font toán mathpazo
	\usepackage[no-math]{fontspec}
	\setmainfont{TeX Gyre Pagella}
	\newfontfamily\EZqag{TeX Gyre Adventor}[NFSSFamily=qag,Ligatures=TeX]
\fi
\usepackage[many]{tcolorbox}
\usetikzlibrary{math,through,calc,intersections,angles,quotes,shapes,shapes.geometric,arrows,patterns,matrix,chains,arrows.meta,decorations.shapes,decorations.fractals,decorations.markings,shadows,decorations.pathreplacing}
\usepackage[top=1.6cm, bottom=1.8cm, left=1.8cm, right=1.8cm]{geometry}% v2: top 1.6cm (bản cũ 1.5cm) để chữ không sát dải đầu trang
\providecommand{\CheDo}{loigiai}% loigiai | dethi | solcolor | color | book
\expandafter\usepackage\expandafter[\CheDo]{ex_test}
% \tenchuong: khai báo ở file chính (trước \input) hoặc \renewcommand sau \input
%============================================================
%  [GỘP TỪ eduzy-de.tex] EDUZY - MÀU NÂU ĐỒNG / VÀNG KIM
%  + ĐẦU TRANG, CHÂN TRANG KIỂU DẢI NGANG (v2), ĐỐI XỨNG TRANG CHẴN / TRANG LẺ.
%  (không cần file eduzy-de.tex riêng nữa)
%  Khối này phải đặt SAU \usepackage[...]{ex_test} và SAU \def\tenchuong{...}
%============================================================
\usepackage{lastpage}
\usepackage{needspace}
%--- DÀN TRANG ĐỀ THI: tách câu (THÊM 05/11/2026; chỉ thêm lệnh mới, KHÔNG đổi hành vi mặc định) ---
%  Hai công tắc, mặc định tắt (= preamble chạy y như cũ). Bật bằng cách đặt TRƯỚC %============================================================
%  PREAMBLE EDUZY v2 - dùng chung cho đề thi, bài học, chuyên đề (gói ex_test)
%  Khác bản preamble-eduzy.tex ở đầu trang, chân trang: dải nền nâu đồng tràn ngang,
%  logo + Eduzy.com.vn, ô chữ vàng kem cạnh xiên, chân trang "Trang (số)";
%  logo Eduzy chìm giữa mọi trang (tắt: \logochimtat, bật lại: \logochimbat).
%  Cách dùng ở file chính (xem mau-chuan.tex, mau-chuyen-de.tex):
%      \def\CheDo{loigiai}          % loigiai | dethi | solcolor | color | book
%      \def\tenchuong{<dòng chữ trên đầu trang>}
%      \def\thuonghieu{<chữ ở chân trang lẻ>}   % không khai báo thì dùng mặc định
%      \def\tacgia{<chữ ở chân trang chẵn>}      % không khai báo thì dùng mặc định
%      \def\khauhieu{<chữ cạnh logo ở đầu trang>}  % tùy chọn, mặc định Eduzy.com.vn
%      %============================================================
%  PREAMBLE EDUZY v2 - dùng chung cho đề thi, bài học, chuyên đề (gói ex_test)
%  Khác bản preamble-eduzy.tex ở đầu trang, chân trang: dải nền nâu đồng tràn ngang,
%  logo + Eduzy.com.vn, ô chữ vàng kem cạnh xiên, chân trang "Trang (số)";
%  logo Eduzy chìm giữa mọi trang (tắt: \logochimtat, bật lại: \logochimbat).
%  Cách dùng ở file chính (xem mau-chuan.tex, mau-chuyen-de.tex):
%      \def\CheDo{loigiai}          % loigiai | dethi | solcolor | color | book
%      \def\tenchuong{<dòng chữ trên đầu trang>}
%      \def\thuonghieu{<chữ ở chân trang lẻ>}   % không khai báo thì dùng mặc định
%      \def\tacgia{<chữ ở chân trang chẵn>}      % không khai báo thì dùng mặc định
%      \def\khauhieu{<chữ cạnh logo ở đầu trang>}  % tùy chọn, mặc định Eduzy.com.vn
%      \input{preamble-eduzy-v2}
%      \TLdapso                      % tài liệu: \TLfull | \TLdapso | \TLan
%      \begin{document} ... \end{document}
%  Logo Eduzy vẽ bằng TikZ ngay trong file (\eduzylogo), không cần file ảnh.
%  Biên dịch: pdfLaTeX (máy có gói vietnam) hoặc XeLaTeX.
%============================================================
\documentclass[12pt,a4paper,twoside]{article}%<-- twoside: bat dau trang/chan trang chan-le
\usepackage{iftex}
\usepackage{amsmath,amssymb,tikz,mathrsfs,tkz-tab}
\usepackage{graphicx}
\usepackage{pgfplots}% đồ thị hàm số
\pgfplotsset{compat=1.17}
%\usepackage{mathptmx}
\usepackage{mathpazo}
\ifXeTeX% XeLaTeX: font chữ có đủ dấu tiếng Việt, giữ nguyên font toán mathpazo
	\usepackage[no-math]{fontspec}
	\setmainfont{TeX Gyre Pagella}
	\newfontfamily\EZqag{TeX Gyre Adventor}[NFSSFamily=qag,Ligatures=TeX]
\fi
\usepackage[many]{tcolorbox}
\usetikzlibrary{math,through,calc,intersections,angles,quotes,shapes,shapes.geometric,arrows,patterns,matrix,chains,arrows.meta,decorations.shapes,decorations.fractals,decorations.markings,shadows,decorations.pathreplacing}
\usepackage[top=1.6cm, bottom=1.8cm, left=1.8cm, right=1.8cm]{geometry}% v2: top 1.6cm (bản cũ 1.5cm) để chữ không sát dải đầu trang
\providecommand{\CheDo}{loigiai}% loigiai | dethi | solcolor | color | book
\expandafter\usepackage\expandafter[\CheDo]{ex_test}
% \tenchuong: khai báo ở file chính (trước \input) hoặc \renewcommand sau \input
%============================================================
%  [GỘP TỪ eduzy-de.tex] EDUZY - MÀU NÂU ĐỒNG / VÀNG KIM
%  + ĐẦU TRANG, CHÂN TRANG KIỂU DẢI NGANG (v2), ĐỐI XỨNG TRANG CHẴN / TRANG LẺ.
%  (không cần file eduzy-de.tex riêng nữa)
%  Khối này phải đặt SAU \usepackage[...]{ex_test} và SAU \def\tenchuong{...}
%============================================================
\usepackage{lastpage}
\usepackage{needspace}
%--- DÀN TRANG ĐỀ THI: tách câu (THÊM 05/11/2026; chỉ thêm lệnh mới, KHÔNG đổi hành vi mặc định) ---
%  Hai công tắc, mặc định tắt (= preamble chạy y như cũ). Bật bằng cách đặt TRƯỚC \input{preamble-eduzy-v2}
%  trong file đề (tools/dan_trang/dan_trang.py tự ghi khi đặt \NeedspaceCau):
%    \def\TachCauDS{1}   câu Đúng/Sai được tách: câu dẫn (kèm hình) ở cuối trang, bảng Phát biểu sang trang sau
%                        (câu dẫn và hình không tách nhau, bảng không tách giữa hai trang)
%    \def\TachCauLG{1}   bản lời giải (\CheDo khác dethi): phần đề của câu cao hơn 1/3 trang được phép tách
%  Lệnh dùng trong file đề (giá trị do tools/dan_trang/dan_trang.py đo ở bản dethi):
%    \NeedspaceCau[<cao câu dẫn + hình>]{<cao cả khối câu>}   ngay trên \begin{ex} (câu đầu phần: trên tiêu đề phần)
%        công tắc tắt: đúng bằng \Needspace{<cao cả khối câu>} như cũ;
%        \TachCauDS=1 và có [..] (câu Đúng/Sai được chọn tách): chỉ giữ câu dẫn + hình;
%        \TachCauLG=1 ở bản lời giải: đòi tối đa 1/3 \textheight.
%    \NeedspaceBang{<cao bảng Phát biểu>}%   ngay trước \choiceTFt: \TachCauDS=1 thì giữ trọn bảng trên một trang;
%        công tắc tắt: không làm gì.
\providecommand{\TachCauDS}{0}
\providecommand{\TachCauLG}{0}
\newlength{\EZnsCau}
\newcommand{\NeedspaceCau}[2][]{%
	\setlength{\EZnsCau}{#2}%
	\ifnum\TachCauDS=1 \if\relax\detokenize{#1}\relax\else\setlength{\EZnsCau}{#1}\fi\fi
	\ifnum\TachCauLG=1 \ifdefstring{\kindOPTN}{dethi}{}{\ifdim\EZnsCau>0.3333\textheight\setlength{\EZnsCau}{0.3333\textheight}\fi}\fi
	\Needspace{\EZnsCau}}
\newcommand{\NeedspaceBang}[1]{\ifnum\TachCauDS=1 \Needspace{#1}\fi}

%--- Thang màu: đậm -> nhạt
\definecolor{gold1}{RGB}{126,96,0}    % nâu đồng đậm  #7E6000
\definecolor{gold2}{RGB}{192,144,3}   % vàng kim      #C09003
\definecolor{gold3}{RGB}{255,215,99}  % vàng nhạt     #FFD763
\definecolor{gold4}{RGB}{255,231,153} % vàng rất nhạt #FFE799
\definecolor{gold5}{RGB}{255,243,205} % vàng kem      #FFF3CD
%--- Ghi đè tông xanh/đỏ cũ
\definecolor{myblue}{RGB}{126,96,0}
\colorlet{blue}{gold1}
\colorlet{cyan}{gold2}
\colorlet{teal}{gold2}

%--- Tên chương mặc định (file đề nên tự khai báo lại)
\providecommand{\tenchuong}{}
%--- Chân trang
\providecommand{\thuonghieu}{TOÁN 12 -- Eduzy.com.vn}
\providecommand{\tacgia}{Tô Gia Bảo -- 0353703783}

%--- Chữ cạnh logo ở đầu trang
\providecommand{\khauhieu}{Eduzy.com.vn}

%====================== THÔNG SỐ ĐẦU TRANG, CHÂN TRANG (v2) ======================
%  Dải nền đậm tràn ngang trang; ô chữ vàng kem có cạnh xiên, kèm 2 vạch xiên chuyển màu.
%  Đầu trang lẻ : logo + \khauhieu (trái)   | ô chữ \tenchuong (phải)
%  Chân trang lẻ: ô chữ \thuonghieu (trái)  | Trang (số) (phải)
%  Trang chẵn soi gương trang lẻ; chân trang chẵn ghi \tacgia.
%--- Màu: chỉ dùng thang vàng kim Eduzy
\colorlet{ezNen}{gold1}% dải nền đậm
\colorlet{ezVach1}{gold3}% vạch xiên sát dải nền (đậm hơn)
\colorlet{ezVach2}{gold4}% vạch xiên sát ô chữ (nhạt hơn)
\colorlet{ezO}{gold5}% ô chữ vàng kem
%--- Kích thước
\newlength{\EZdai}\setlength{\EZdai}{1.1cm}% chiều cao dải nền
\newlength{\EZo}\setlength{\EZo}{0.68cm}% chiều cao ô chữ
\newlength{\EZxien}\setlength{\EZxien}{0.55cm}% độ lệch ngang của cạnh xiên
\newlength{\EZvach}\setlength{\EZvach}{0.75cm}% bề rộng mỗi vạch xiên
\newlength{\EZkhe}\setlength{\EZkhe}{0.12cm}% khe trắng giữa vạch và ô chữ
\newlength{\EZdem}\setlength{\EZdem}{0.4cm}% khoảng đệm từ chữ tới cạnh xiên
\newlength{\EZle}\setlength{\EZle}{1.8cm}% lề trái/phải của chữ (bằng lề thân bài)
\newlength{\EZrongmin}\setlength{\EZrongmin}{4.5cm}% ô chữ rộng tối thiểu
\newlength{\EZrong}% bề rộng chữ trong ô (tính tự động)
%--- Font
\newcommand{\EZfontO}{\fontsize{10.5pt}{12pt}\selectfont\bfseries}% chữ trong ô
\newcommand{\EZfontKH}{\fontfamily{qag}\fontsize{11pt}{12pt}\selectfont\bfseries}% \khauhieu cạnh logo
\newcommand{\EZfontTrang}{\fontsize{11.5pt}{12pt}\selectfont\bfseries}% chữ "Trang" và số trang
%--- Đo bề rộng chữ trong ô
\newcommand{\EZdorong}[1]{%
	\settowidth{\EZrong}{\EZfontO #1}%
	\ifdim\EZrong<\EZrongmin\setlength{\EZrong}{\EZrongmin}\fi}
%--- Logo + chữ trắng \khauhieu (dựng sẵn một lần khi vào \begin{document})
\newsavebox{\EZlogobox}
\AtBeginDocument{\sbox{\EZlogobox}{%
		\raisebox{-0.5\height}{\eduzylogo[1cm]}\hspace{0.15cm}%
		\raisebox{-0.5\height}{\color{white}\EZfontKH\khauhieu}}}
%--- Hai vạch xiên + khe trắng. #1: điểm gốc ở mép dưới, #2: hướng (1 sang phải, -1 sang trái),
%    #3: hướng nghiêng (1 kiểu "/", -1 kiểu "\")
\newcommand{\EZvachxien}[3]{%
	\fill[white] (#1) --++({(#3)*\EZxien},\EZo) --++({(#2)*\EZkhe},0) --++({-(#3)*\EZxien},-\EZo) --cycle;
	\fill[ezVach2] ([xshift=(#2)*\EZkhe]#1) --++({(#3)*\EZxien},\EZo) --++({(#2)*\EZvach},0) --++({-(#3)*\EZxien},-\EZo) --cycle;
	\fill[ezVach1] ([xshift=(#2)*(\EZkhe+\EZvach)]#1) --++({(#3)*\EZxien},\EZo) --++({(#2)*\EZvach},0) --++({-(#3)*\EZxien},-\EZo) --cycle;
	\draw[white,line width=0.6pt] ([xshift=(#2)*(\EZkhe+\EZvach)]#1) --++({(#3)*\EZxien},\EZo);
}
%--- Chữ trong ô (căn giữa theo chiều cao ô, không lệch vì dấu tiếng Việt)
\tikzset{ezchu/.style={inner sep=0pt,text=ezNen,text height=0.72em,text depth=0.22em}}

%====================== LOGO CHÌM GIỮA TRANG ======================
%  Vẽ cùng đầu trang nên nằm dưới chữ của thân bài, có trên mọi trang.
%  Tắt: \logochimtat (từ trang đó trở đi), bật lại: \logochimbat.
%  Đổi cỡ, độ đậm: khai báo trước \input, ví dụ \def\EZchimcao{9cm} \def\EZchimdo{0.1}
\providecommand{\EZchimcao}{10cm}% chiều cao logo chìm
\providecommand{\EZchimdo}{0.12}% độ đậm (0 = trong suốt, 1 = rõ hẳn)
\newif\ifEZchim \EZchimtrue
\newcommand{\logochimtat}{\EZchimfalse}
\newcommand{\logochimbat}{\EZchimtrue}
%  Dựng sẵn logo mờ một lần khi vào \begin{document}; transparency group làm mờ cả logo
%  một lượt nên không lộ vệt ở chỗ các dải màu chồng nhau.
\newsavebox{\EZchimtho}
\newsavebox{\EZchimbox}
\AtBeginDocument{%
	\sbox{\EZchimtho}{\eduzylogo[\EZchimcao]}%
	\sbox{\EZchimbox}{\begin{tikzpicture}
			\begin{scope}[transparency group,opacity=\EZchimdo]
				\node[inner sep=0pt] {\usebox{\EZchimtho}};
			\end{scope}
	\end{tikzpicture}}}
\newcommand{\EZvechim}{%
	\ifEZchim
		\node[inner sep=0pt] at (current page.center) {\usebox{\EZchimbox}};
	\fi}

%====================== ĐẦU TRANG - TRANG LẺ ======================
%  (logo bên trái, ô chữ \tenchuong bên phải)
\newcommand{\myfancyheadOdd}{%
	\EZdorong{\tenchuong}%
	\begin{tikzpicture}[remember picture,overlay]
		\coordinate (T) at (current page.north west);
		\coordinate (P) at (current page.north east);
		%--- logo chìm giữa trang
		\EZvechim
		%--- dải nền
		\fill[ezNen] (T) rectangle ([yshift=-\EZdai]P);
		%--- ô chữ: X là góc dưới bên trái, cạnh trái xiên kiểu "/"
		\coordinate (X) at ([xshift=-\EZle-\EZrong-\EZdem-\EZxien,yshift=-\EZdai]P);
		\fill[ezO] (X) --++(\EZxien,\EZo) -- ([yshift=-\EZdai+\EZo]P) -- ([yshift=-\EZdai]P) --cycle;
		\EZvachxien{X}{-1}{1}
		\node[ezchu,anchor=east] at ([xshift=-\EZle,yshift=-\EZdai+0.5*\EZo]P) {\EZfontO\tenchuong};
		%--- logo + tên
		\node[inner sep=0pt,anchor=west] at ([xshift=\EZle,yshift=-0.5*\EZdai]T) {\usebox{\EZlogobox}};
	\end{tikzpicture}%
}

%====================== ĐẦU TRANG - TRANG CHẴN =====================
%  (soi gương: ô chữ \tenchuong bên trái, logo bên phải)
\newcommand{\myfancyheadEven}{%
	\EZdorong{\tenchuong}%
	\begin{tikzpicture}[remember picture,overlay]
		\coordinate (T) at (current page.north west);
		\coordinate (P) at (current page.north east);
		%--- logo chìm giữa trang
		\EZvechim
		%--- dải nền
		\fill[ezNen] (T) rectangle ([yshift=-\EZdai]P);
		%--- ô chữ: X là góc dưới bên phải, cạnh phải xiên kiểu "\"
		\coordinate (X) at ([xshift=\EZle+\EZrong+\EZdem+\EZxien,yshift=-\EZdai]T);
		\fill[ezO] (X) --++(-\EZxien,\EZo) -- ([yshift=-\EZdai+\EZo]T) -- ([yshift=-\EZdai]T) --cycle;
		\EZvachxien{X}{1}{-1}
		\node[ezchu,anchor=west] at ([xshift=\EZle,yshift=-\EZdai+0.5*\EZo]T) {\EZfontO\tenchuong};
		%--- logo + tên
		\node[inner sep=0pt,anchor=east] at ([xshift=-\EZle,yshift=-0.5*\EZdai]P) {\usebox{\EZlogobox}};
	\end{tikzpicture}%
}

%====================== CHÂN TRANG - TRANG LẺ ======================
%  (ô chữ \thuonghieu bên trái, "Trang (số)" bên phải)
\newcommand{\myfancyfootOdd}{%
	\EZdorong{\thuonghieu}%
	\begin{tikzpicture}[remember picture,overlay]
		\coordinate (T) at (current page.south west);
		\coordinate (P) at (current page.south east);
		%--- dải nền
		\fill[ezNen] (T) rectangle ([yshift=\EZdai]P);
		%--- ô chữ sát mép trên dải: X là góc dưới bên phải, cạnh phải xiên kiểu "/"
		\coordinate (X) at ([xshift=\EZle+\EZrong+\EZdem,yshift=\EZdai-\EZo]T);
		\fill[ezO] ([yshift=\EZdai-\EZo]T) -- (X) --++(\EZxien,\EZo) -- ([yshift=\EZdai]T) --cycle;
		\EZvachxien{X}{1}{1}
		\node[ezchu,anchor=west] at ([xshift=\EZle,yshift=\EZdai-0.5*\EZo]T) {\EZfontO\thuonghieu};
		%--- số trang
		\node[circle,draw=white,line width=0.9pt,text=white,inner sep=0pt,minimum size=0.62cm,anchor=east]
		(so) at ([xshift=-\EZle,yshift=0.5*\EZdai]P) {\EZfontTrang\thepage};
		\node[inner sep=0pt,text=white,anchor=base east] (tr) at ([xshift=-0.15cm]so.base west) {\EZfontTrang Trang};
	\end{tikzpicture}%
}

%====================== CHÂN TRANG - TRANG CHẴN ====================
%  (soi gương: "Trang (số)" bên trái, ô chữ \tacgia bên phải)
\newcommand{\myfancyfootEven}{%
	\EZdorong{\tacgia}%
	\begin{tikzpicture}[remember picture,overlay]
		\coordinate (T) at (current page.south west);
		\coordinate (P) at (current page.south east);
		%--- dải nền
		\fill[ezNen] (T) rectangle ([yshift=\EZdai]P);
		%--- ô chữ sát mép trên dải: X là góc dưới bên trái, cạnh trái xiên kiểu "\"
		\coordinate (X) at ([xshift=-\EZle-\EZrong-\EZdem,yshift=\EZdai-\EZo]P);
		\fill[ezO] ([yshift=\EZdai-\EZo]P) -- (X) --++(-\EZxien,\EZo) -- ([yshift=\EZdai]P) --cycle;
		\EZvachxien{X}{-1}{-1}
		\node[ezchu,anchor=east] at ([xshift=-\EZle,yshift=\EZdai-0.5*\EZo]P) {\EZfontO\tacgia};
		%--- số trang
		%    (vòng tròn ẩn so0 nằm giữa dải để lấy đường chân chữ, giống trang lẻ)
		\node[circle,inner sep=0pt,minimum size=0.62cm] (so0) at ([yshift=0.5*\EZdai]T) {\EZfontTrang\phantom{\thepage}};
		\node[inner sep=0pt,text=white,anchor=base west] (tr) at ($(so0.base -| T)+(\EZle,0)$) {\EZfontTrang Trang};
		\node[circle,draw=white,line width=0.9pt,text=white,inner sep=0pt,minimum size=0.62cm,anchor=base west]
		(so) at ([xshift=0.15cm]tr.base east) {\EZfontTrang\thepage};
	\end{tikzpicture}%
}

%--- Giữ tương thích với các file cũ đang gọi \myfancyhead, \myfancyfoot
\newcommand{\myfancyhead}{\myfancyheadOdd}
\newcommand{\myfancyfoot}{\myfancyfootOdd}

%======================= GẮN VÀO TRANG =============================
\pagestyle{fancy}
\fancyhf{}
\fancyhead[LO]{\myfancyheadOdd}
\fancyhead[RE]{\myfancyheadEven}
\fancyfoot[LO]{\myfancyfootOdd}
\fancyfoot[RE]{\myfancyfootEven}
\renewcommand{\headrulewidth}{0pt}
\renewcommand{\footrulewidth}{0pt}
\fancypagestyle{plain}{%
	\fancyhf{}
	\fancyhead[LO]{\myfancyheadOdd}
	\fancyhead[RE]{\myfancyheadEven}
	\fancyfoot[LO]{\myfancyfootOdd}
	\fancyfoot[RE]{\myfancyfootEven}
	\renewcommand{\headrulewidth}{0pt}
	\renewcommand{\footrulewidth}{0pt}
}
%=================== HẾT PHẦN GỘP TỪ eduzy-de.tex ===================
\newcommand{\indam}[1]{
	{\fontsize{12pt}{9pt}\sffamily \selectfont\bfseries\color{gold1} #1}
}
\newcommand{\ind}[1]{
	{\fontsize{14pt}{9pt}\fontfamily{qag} \selectfont\bfseries\color{gold1} #1}
}
\renewcommand{\FalseEX}{\stepcounter{dapan}\textbf{\selectfont\color{gold1}\Alph{dapan}.}\ignorespaces}
\renewtheorem{ex}{\color{gold1}Câu}
\renewcommand{\baselinestretch}{1.1}
\newcommand{\hoac}[1]{
	\left[\begin{aligned}#1\end{aligned}\right.}
\newcommand{\heva}[1]{
	\left\{\begin{aligned}#1\end{aligned}\right.}
\usepackage{esvect}
\def\vec{\vv}
\def\vt{\vv}
\def\overrightarrow{\vv}
%====================== LOGO EDUZY VẼ BẰNG TIKZ ======================
%  Đường viền dò từ logo gốc; dải màu vàng kim đo trực tiếp trên ảnh gốc
%  (nghiêng 21 độ, tô bằng 80 dải hẹp nên hiển thị giống nhau trên mọi trình xem PDF).
%  Dùng: \eduzylogo          -> cao 2.4cm (mặc định)
%        \eduzylogo[1.5cm]   -> cao 1.5cm
\definecolor{ezText}{RGB}{245,200,46}
\newcommand{\eduzylogo}[1][2.4cm]{%
	\resizebox{!}{#1}{%
		\begin{tikzpicture}[x=0.01cm,y=0.01cm]
			\path[use as bounding box] (0,0) rectangle (224,-224);
			% mũ tốt nghiệp + vương miện
			\begin{scope}
				\clip
				(119.4,-176.1)..controls(118.9,-176.1)and(116.8,-176)..(114.6,-176)
				..controls(101.4,-176)and(85.8,-175.9)..(83.7,-175.7)
				..controls(81.8,-175.6)and(81,-175.6)..(80.2,-175.7)
				..controls(78.8,-175.9)and(77.7,-175.9)..(76.2,-175.7)
				..controls(74.9,-175.5)and(70.8,-175.5)..(64.3,-175.7)
				..controls(59.1,-175.9)and(58.3,-175.7)..(58.6,-173.7)
				..controls(58.6,-172.8)and(59.6,-169.6)..(60.3,-168)
				..controls(60.8,-166.9)and(61.4,-165.8)..(62.3,-165)
				..controls(63.1,-164.2)and(64.3,-163.4)..(65,-163.2)
				..controls(65.5,-163.1)and(72.3,-163)..(76.6,-163.1)
				..controls(77.9,-163.1)and(80.1,-163.1)..(81.4,-163)
				..controls(82.8,-162.9)and(86.3,-162.9)..(89.2,-162.9)
				..controls(93,-162.9)and(94.9,-162.8)..(96,-162.7)
				..controls(97.2,-162.5)and(101,-162.5)..(114,-162.4)
				..controls(123,-162.3)and(134.9,-162.2)..(140.2,-162.1)
				..controls(150.1,-162)and(152.2,-162.1)..(153.6,-162.4)
				..controls(154.4,-162.6)and(157.1,-163.8)..(158.1,-164.4)
				..controls(158.8,-164.8)and(159.7,-165.7)..(159.9,-166.3)
				..controls(160,-166.5)and(160.2,-167.1)..(160.2,-167.6)
				..controls(160.3,-168)and(160.6,-169.4)..(160.9,-170.6)
				..controls(161.2,-171.7)and(161.5,-173.1)..(161.6,-173.6)
				..controls(161.8,-174.4)and(161.8,-174.5)..(161.5,-174.9)
				..controls(161.2,-175.6)and(160.8,-175.9)..(159.9,-176)
				..controls(159.2,-176.1)and(154.1,-176.2)..(151.5,-176.2)
				..controls(148.1,-176.1)and(122.6,-176.1)..(121.6,-176.2)
				..controls(120.8,-176.2)and(119.8,-176.2)..(119.4,-176.1)--cycle
				(130.4,-160)..controls(129,-160)and(119.8,-160)..(110.1,-159.9)
				..controls(100.3,-159.8)and(88.5,-159.7)..(83.9,-159.6)
				..controls(77.8,-159.5)and(75.3,-159.5)..(74.5,-159.6)
				..controls(73.9,-159.7)and(72.5,-159.8)..(71.3,-159.8)
				..controls(69.6,-159.8)and(69,-159.8)..(68.5,-159.6)
				..controls(67.6,-159.3)and(66.8,-158.6)..(66.2,-157.6)
				..controls(66,-157.3)and(65.6,-156.5)..(65.3,-156.1)
				..controls(64.6,-155.1)and(64.5,-154.9)..(63.7,-152.9)
				..controls(63.2,-151.6)and(62.9,-150.9)..(61.8,-149.6)
				..controls(61.5,-149.2)and(61.1,-148.5)..(60.9,-147.9)
				..controls(60.7,-147.4)and(60.3,-146.6)..(60.1,-146.1)
				..controls(59.9,-145.5)and(59.5,-144.9)..(59,-144.3)
				..controls(58.5,-143.7)and(58.2,-143.2)..(57.9,-142.5)
				..controls(57.7,-141.9)and(57.3,-141.1)..(57,-140.7)
				..controls(56.6,-140.3)and(56.1,-139.3)..(55.8,-138.5)
				..controls(55.5,-137.8)and(55.1,-136.8)..(54.9,-136.5)
				..controls(54.7,-136.1)and(54.3,-135.2)..(53.9,-134.5)
				..controls(53.6,-133.8)and(53,-132.7)..(52.6,-132)
				..controls(51.6,-130.5)and(51.4,-130.2)..(51.1,-129)
				..controls(50.6,-127.5)and(50.6,-127.4)..(49.7,-126.3)
				..controls(49.2,-125.7)and(48.6,-124.8)..(48.3,-124.2)
				..controls(48,-123.6)and(47.5,-122.7)..(47.3,-122.2)
				..controls(47,-121.8)and(46.5,-120.7)..(46.1,-119.8)
				..controls(45.8,-118.9)and(45.2,-117.8)..(44.9,-117.2)
				..controls(44.2,-116.1)and(43.3,-114.4)..(42.3,-112.2)
				..controls(41.8,-111.3)and(41.3,-110.4)..(41.1,-110.1)
				..controls(40.5,-109.4)and(40.3,-109)..(40,-107.8)
				..controls(39.8,-107)and(39.6,-106.5)..(39.2,-106)
				..controls(38.9,-105.6)and(38.4,-104.7)..(38,-104)
				..controls(37.7,-103.4)and(37.2,-102.5)..(37,-102.1)
				..controls(36.7,-101.7)and(36.3,-100.8)..(36,-100.1)
				..controls(35.6,-98.9)and(35.3,-98.4)..(34.6,-97.6)
				..controls(34.2,-97.1)and(34.2,-96.3)..(34.6,-96)
				..controls(35,-95.7)and(35.5,-95.9)..(36.2,-96.7)
				..controls(36.5,-97)and(37.2,-97.5)..(37.6,-97.7)..controls(38,-98)and(38.6,-98.4)..(39.1,-98.9)
				..controls(39.6,-99.4)and(40.3,-100)..(40.6,-100.3)
				..controls(41.3,-100.8)and(42.4,-101.8)..(43.7,-102.9)
				..controls(44,-103.2)and(44.6,-103.7)..(45,-103.9)
				..controls(45.4,-104.2)and(45.9,-104.6)..(46.1,-104.8)
				..controls(47.1,-105.9)and(48.8,-107.1)..(50.5,-107.9)
				..controls(51.8,-108.6)and(52.3,-109)..(53.4,-110)
				..controls(54.8,-111.2)and(55.6,-111.6)..(57,-112.1)
				..controls(57.5,-112.3)and(58.5,-112.8)..(59.4,-113.3)
				..controls(60.2,-113.7)and(61.3,-114.3)..(61.9,-114.5)
				..controls(62.4,-114.7)and(63.3,-115.1)..(63.8,-115.3)
				..controls(64.4,-115.6)and(65.4,-115.9)..(66.2,-116.1)
				..controls(67,-116.3)and(68.4,-116.6)..(69.2,-116.9)
				..controls(72.3,-117.9)and(72.8,-118)..(78,-118.1)
				..controls(84,-118.3)and(86.1,-118.1)..(88.1,-117.1)
				..controls(89.4,-116.5)and(93,-115)..(93.7,-114.9)
				..controls(94.8,-114.6)and(95.4,-114.2)..(95.9,-113.4)
				..controls(96.2,-113)and(96.7,-112.5)..(97.1,-112.2)
				..controls(98,-111.6)and(101.3,-108.4)..(102.1,-107.2)
				..controls(102.4,-106.8)and(102.9,-106.2)..(103.2,-105.8)
				..controls(103.5,-105.5)and(104,-104.8)..(104.2,-104.2)
				..controls(104.5,-103.6)and(105,-102.8)..(105.3,-102.4)
				..controls(105.6,-102)and(106.1,-101.2)..(106.5,-100.6)
				..controls(107.1,-99.5)and(108.6,-97.5)..(109.5,-96.5)
				..controls(109.9,-96.1)and(110.1,-96)..(110.4,-96)
				..controls(111.1,-96)and(111.6,-96.5)..(111.8,-97.6)
				..controls(112,-98.6)and(112.3,-99.3)..(112.8,-99.9)
				..controls(113.7,-100.9)and(115.2,-103.1)..(115.7,-104.2)
				..controls(116.1,-105)and(116.5,-105.5)..(116.9,-106)
				..controls(117.3,-106.4)and(117.8,-106.9)..(118,-107.3)
				..controls(118.2,-107.6)and(118.6,-108)..(119,-108.3)
				..controls(119.6,-108.8)and(121.1,-110.3)..(121.8,-111.2)
				..controls(122.1,-111.7)and(122.5,-112)..(123.4,-112.4)
				..controls(124.8,-113.2)and(125.8,-113.8)..(127,-114.6)
				..controls(127.5,-114.9)and(128.1,-115.3)..(128.2,-115.4)
				..controls(128.4,-115.5)and(129.3,-115.8)..(130.2,-116)
				..controls(131.1,-116.2)and(132.5,-116.6)..(133.3,-116.9)
				..controls(135.7,-117.9)and(136.1,-117.9)..(142.5,-117.8)
				..controls(145.8,-117.8)and(148.4,-117.7)..(148.8,-117.6)
				..controls(149.2,-117.6)and(150.2,-117.2)..(151,-116.8)
				..controls(152,-116.4)and(152.8,-116)..(153.6,-115.9)
				..controls(154.5,-115.7)and(155.2,-115.4)..(156.3,-114.9)
				..controls(157.1,-114.5)and(158.2,-114)..(158.7,-113.8)
				..controls(159.1,-113.7)and(159.8,-113.3)..(160.2,-113)
				..controls(160.6,-112.7)and(161.3,-112.4)..(161.7,-112.2)
				..controls(162.7,-111.9)and(163.2,-111.7)..(164.1,-111.1)
				..controls(164.5,-110.8)and(165.3,-110.4)..(165.9,-110)
				..controls(166.5,-109.7)and(167.3,-109.2)..(167.6,-108.9)
				..controls(167.9,-108.6)and(168.6,-108.2)..(169,-108)
				..controls(169.4,-107.8)and(170.2,-107.3)..(170.7,-106.9)
				..controls(171.2,-106.5)and(172.1,-105.8)..(172.7,-105.4)
				..controls(173.2,-105)and(174.1,-104.4)..(174.5,-104)
				..controls(174.9,-103.7)and(175.7,-103.1)..(176.2,-102.6)
				..controls(176.7,-102.2)and(177.4,-101.5)..(177.9,-101.2)
				..controls(178.3,-100.8)and(178.8,-100.4)..(179,-100.2)
				..controls(179.2,-100)and(179.6,-99.6)..(180,-99.3)
				..controls(180.4,-99.1)and(180.9,-98.6)..(181.2,-98.4)
				..controls(181.5,-98.1)and(182.1,-97.7)..(182.5,-97.5)
				..controls(183.2,-97.2)and(183.5,-96.9)..(184.3,-95.9)
				..controls(184.6,-95.6)and(184.6,-95.6)..(184.8,-95.7)
				..controls(185,-95.8)and(185.1,-96)..(185.1,-96.2)
				..controls(185.1,-96.5)and(185,-96.7)..(184.6,-97.2)
				..controls(184.1,-97.8)and(184,-97.9)..(183.9,-98.8)
				..controls(183.7,-100.3)and(183.3,-101.3)..(182.9,-101.8)
				..controls(182.6,-102.1)and(182.3,-102.6)..(182.2,-102.9)
				..controls(182,-103.3)and(181.6,-104)..(181.2,-104.6)
				..controls(180.9,-105.2)and(180.5,-106.1)..(180.4,-106.6)
				..controls(179.9,-108)and(179.7,-108.5)..(179,-109.4)
				..controls(178.6,-109.8)and(178.2,-110.5)..(177.9,-111)
				..controls(177.7,-111.5)and(177.3,-112.3)..(177,-112.7)
				..controls(176.6,-113.2)and(176.4,-113.7)..(176.3,-114.2)
				..controls(176.2,-114.7)and(175.9,-115.2)..(175.4,-116)
				..controls(175,-116.6)and(174.4,-117.7)..(174,-118.5)
				..controls(173.7,-119.2)and(173.3,-120.1)..(173,-120.5)
				..controls(172.8,-120.9)and(172.4,-121.8)..(172.1,-122.5)
				..controls(171.8,-123.2)and(171.3,-124.1)..(171,-124.6)
				..controls(170.8,-125.1)and(170.3,-126)..(169.9,-126.6)
				..controls(169.6,-127.3)and(169.2,-128.1)..(169,-128.3)
				..controls(168.8,-128.6)and(168.5,-129.3)..(168.3,-129.8)
				..controls(168.1,-130.3)and(167.8,-131.1)..(167.5,-131.6)
				..controls(166.9,-132.5)and(164.2,-137.9)..(164.2,-138.1)
				..controls(164.2,-138.1)and(163.9,-138.8)..(163.4,-139.7)
				..controls(162.9,-140.5)and(162.4,-141.5)..(162.3,-141.9)
				..controls(161.9,-143)and(161.6,-143.4)..(160.9,-144.1)
				..controls(160.4,-144.6)and(160.3,-144.8)..(160.1,-145.4)
				..controls(160,-145.8)and(159.6,-146.8)..(159.2,-147.8)
				..controls(158.7,-148.8)and(158.1,-150)..(157.9,-150.6)
				..controls(157.6,-151.1)and(157.2,-151.9)..(156.9,-152.4)
				..controls(156.5,-152.8)and(156.1,-153.7)..(155.8,-154.4)
				..controls(155.6,-155)and(155.2,-155.9)..(154.9,-156.3)
				..controls(154.7,-156.7)and(154.4,-157.4)..(154.2,-157.9)
				..controls(154.1,-158.4)and(153.8,-159)..(153.5,-159.3)
				..controls(152.8,-160.1)and(153,-160.1)..(142.4,-160.1)
				..controls(137.2,-160.1)and(131.8,-160.1)..(130.4,-160)--cycle
				(141.5,-103.7)..controls(140.8,-103)and(139.8,-102.1)..(139.2,-101.6)
				..controls(138.5,-101.2)and(137.9,-100.6)..(137.8,-100.4)
				..controls(137,-99.5)and(134,-96.5)..(133.2,-96)
				..controls(132.7,-95.7)and(131.8,-94.9)..(131.1,-94.3)
				..controls(129.4,-92.7)and(127.9,-91.6)..(126.6,-91.1)
				..controls(125.2,-90.5)and(124.1,-90)..(122.9,-89.1)
				..controls(122.1,-88.6)and(121.3,-88.2)..(119.9,-87.7)
				..controls(118.8,-87.3)and(117.5,-86.8)..(116.9,-86.5)
				..controls(114.3,-85.5)and(113.7,-85.4)..(109.7,-85.5)
				..controls(105.3,-85.7)and(103.9,-85.9)..(101.6,-87)
				..controls(100.7,-87.4)and(99.4,-87.8)..(97.9,-88)
				..controls(97,-88.2)and(96.6,-88.3)..(95.8,-88.9)
				..controls(95.4,-89.2)and(94.7,-89.6)..(94.2,-89.8)
				..controls(92.9,-90.3)and(92.3,-90.7)..(90.9,-92.1)
				..controls(90.2,-92.7)and(89.4,-93.3)..(89,-93.5)..controls(87,-94.5)and(85.2,-96)..(83.6,-97.9)
				..controls(83.3,-98.2)and(82.5,-98.9)..(81.8,-99.4)
				..controls(80.7,-100.1)and(80.5,-100.3)..(80,-101.1)
				..controls(78.7,-103.1)and(77.2,-103.9)..(76.4,-103)
				..controls(75.7,-102.2)and(75.7,-103.2)..(75.7,-83.6)
				..controls(75.6,-73.6)and(75.6,-64.8)..(75.5,-64)
				..controls(75.5,-63.2)and(75.5,-61.4)..(75.6,-59.8)--(75.7,-57.1)--(76,-56.8)
				..controls(76.8,-56.2)and(78,-56.3)..(79.8,-57.2)
				..controls(80.4,-57.5)and(81.3,-57.8)..(81.8,-57.9)
				..controls(82.3,-58.1)and(83,-58.4)..(83.6,-58.8)..controls(84.3,-59.2)and(85,-59.5)..(86,-59.8)
				..controls(86.8,-60.1)and(88,-60.6)..(88.7,-60.9)
				..controls(89.3,-61.2)and(90.6,-61.7)..(91.4,-62)
				..controls(92.2,-62.3)and(93.7,-62.9)..(94.6,-63.3)
				..controls(95.6,-63.7)and(96.8,-64.1)..(97.4,-64.3)
				..controls(98,-64.5)and(99,-64.9)..(99.7,-65.2)
				..controls(100.4,-65.6)and(101.5,-66)..(102.3,-66.2)
				..controls(104.2,-66.8)and(105,-67.1)..(106.2,-67.8)
				..controls(107.9,-68.7)and(110.2,-69.4)..(111.3,-69.2)
				..controls(112.5,-69)and(115,-68.2)..(115.6,-67.8)
				..controls(117.3,-66.7)and(117.6,-66.6)..(119.5,-66.2)
				..controls(121.9,-65.8)and(122.1,-65.7)..(123,-65.1)
				..controls(124.1,-64.4)and(124.4,-64.3)..(126,-63.9)
				..controls(127,-63.7)and(127.7,-63.4)..(128.4,-63.1)
				..controls(129,-62.8)and(129.7,-62.5)..(130.6,-62.4)
				..controls(132.1,-62.1)and(134,-61.5)..(135.2,-61)
				..controls(135.7,-60.8)and(136.7,-60.5)..(137.3,-60.3)
				..controls(138.4,-60)and(139.3,-59.7)..(141.1,-59)
				..controls(141.6,-58.7)and(142,-58.6)..(142.5,-58.6)
				..controls(143.2,-58.6)and(143.3,-58.7)..(143.6,-59)
				..controls(144,-59.5)and(144.1,-59.9)..(144,-60.8)
				..controls(144,-61.2)and(144,-71.1)..(144,-82.8)--(143.9,-104.1)--(143.7,-104.5)
				..controls(143.5,-104.7)and(143.3,-104.8)..(143,-104.9)
				..controls(142.6,-104.9)and(142.5,-104.8)..(141.5,-103.7)--cycle
				(27.8,-91.4)..controls(27.1,-90.9)and(26.3,-90)..(26.1,-89.5)
				..controls(26,-89.2)and(25.8,-88.6)..(25.7,-87.9)
				..controls(25.4,-86.2)and(25.6,-84.8)..(26.5,-83.5)
				..controls(27.4,-82.3)and(28.3,-82)..(29.4,-82.3)
				..controls(31.1,-82.8)and(32.1,-84.9)..(32,-87.4)
				..controls(31.9,-89.3)and(31.2,-90.6)..(29.9,-91.4)
				..controls(29.1,-91.9)and(28.6,-91.9)..(27.8,-91.4)--cycle
				(190.1,-91.3)..controls(189.7,-91.2)and(189.3,-90.9)..(189.1,-90.7)
				..controls(188.3,-89.8)and(187.6,-87.7)..(187.6,-86.2)
				..controls(187.6,-85.5)and(187.7,-85.2)..(188.2,-84.4)
				..controls(189.1,-82.6)and(190.1,-81.8)..(191.3,-82.1)
				..controls(192,-82.3)and(193.5,-83.7)..(193.8,-84.4)
				..controls(194.1,-85.1)and(194.2,-87.1)..(194,-88.3)
				..controls(193.8,-89.8)and(193.3,-90.5)..(192.3,-91.3)
				..controls(191.7,-91.7)and(190.9,-91.7)..(190.1,-91.3)--cycle
				(173.7,-83.4)..controls(172.2,-82.5)and(172,-82.5)..(170.6,-82.4)
				..controls(169.7,-82.4)and(168.8,-82.4)..(168.3,-82.5)
				..controls(167.2,-82.7)and(166.9,-82.7)..(166.6,-82.2)
				..controls(166.3,-81.8)and(166.3,-81)..(166.8,-80)
				..controls(167.2,-78.9)and(167.4,-77.4)..(167.7,-74.1)
				..controls(167.8,-72.8)and(167.9,-71.8)..(167.9,-71.7)
				..controls(168,-71.7)and(168,-70.2)..(168,-68.4)..controls(168,-65.4)and(168,-64.9)..(167.7,-63)
				..controls(167.3,-60.5)and(167.1,-59.9)..(166.4,-58.4)
				..controls(166.1,-57.7)and(165.8,-56.9)..(165.7,-56.6)
				..controls(165.7,-56.2)and(165.4,-55.6)..(165.2,-55.2)
				..controls(164.6,-53.8)and(161.9,-50.7)..(161,-50.3)
				..controls(159.8,-49.7)and(157.9,-50)..(155.2,-51.2)
				..controls(154.7,-51.4)and(153.4,-51.8)..(152.2,-52.1)
				..controls(150.8,-52.4)and(149.4,-52.9)..(148.6,-53.2)
				..controls(147.8,-53.6)and(146.7,-54)..(146,-54.2)
				..controls(145.4,-54.4)and(144.5,-54.7)..(144,-54.9)
				..controls(143.6,-55.2)and(142.7,-55.5)..(141.8,-55.7)
				..controls(141,-55.9)and(139.6,-56.3)..(138.7,-56.6)
				..controls(137.8,-57)and(136.3,-57.4)..(135.3,-57.6)
				..controls(133.8,-57.9)and(133.2,-58.1)..(132.1,-58.7)
				..controls(131.1,-59.2)and(130.3,-59.4)..(129.4,-59.6)
				..controls(128.7,-59.8)and(127.8,-60)..(127.4,-60.1)
				..controls(127,-60.3)and(126,-60.5)..(125.1,-60.7)
				..controls(123.2,-61.1)and(121.7,-61.5)..(120.2,-62.1)
				..controls(119.7,-62.3)and(118.1,-62.8)..(116.7,-63.2)
				..controls(114.6,-63.9)and(114,-64.1)..(113,-64.7)
				..controls(111.9,-65.4)and(111.8,-65.4)..(111,-65.5)
				..controls(109.9,-65.5)and(109.3,-65.4)..(107.7,-64.7)
				..controls(107,-64.4)and(106,-64)..(105.3,-63.9)
				..controls(103.5,-63.5)and(102.2,-63)..(101.2,-62.5)
				..controls(100.1,-61.8)and(99.3,-61.5)..(97.8,-61.1)
				..controls(97.2,-61)and(96.2,-60.7)..(95.7,-60.4)
				..controls(95.1,-60.1)and(94.2,-59.8)..(93.5,-59.6)
				..controls(92.9,-59.4)and(91.7,-59)..(91,-58.7)..controls(90.2,-58.4)and(88.9,-58)..(88,-57.7)
				..controls(87.2,-57.5)and(86.1,-57.1)..(85.5,-56.8)
				..controls(85,-56.6)and(84.1,-56.2)..(83.5,-56)..controls(82.9,-55.7)and(81.8,-55.2)..(81,-54.9)
				..controls(80.1,-54.5)and(79.1,-54.2)..(78.7,-54.1)
				..controls(77.9,-53.8)and(77.3,-53.6)..(76.2,-52.9)
				..controls(75.7,-52.6)and(75.2,-52.4)..(74.1,-52.2)
				..controls(72.5,-51.9)and(72.4,-51.8)..(71.3,-51.1)
				..controls(69.9,-50.1)and(69.4,-49.9)..(67.9,-49.8)
				..controls(67,-49.7)and(66.6,-49.6)..(66.3,-49.4)
				..controls(66.1,-49.3)and(65.7,-49)..(65.4,-48.8)
				..controls(65.1,-48.7)and(64.2,-48.4)..(63.6,-48.2)
				..controls(62.9,-48)and(61.7,-47.6)..(60.9,-47.2)
				..controls(60.1,-46.8)and(58.9,-46.4)..(58.2,-46.2)
				..controls(57.6,-46)and(56.8,-45.7)..(56.5,-45.5)
				..controls(54.9,-44.6)and(54.6,-44.5)..(53.2,-44.2)
				..controls(51.7,-43.8)and(51,-43.6)..(49.9,-42.9)
				..controls(49.5,-42.6)and(48.9,-42.4)..(47.5,-42)
				..controls(45.7,-41.6)and(45.6,-41.5)..(44.3,-40.6)
				..controls(43.5,-40.1)and(42.8,-39.6)..(42.7,-39.4)
				..controls(42.4,-38.7)and(43,-38.1)..(44.5,-37.6)..controls(44.9,-37.5)and(46,-37)..(46.9,-36.6)
				..controls(47.9,-36.1)and(48.8,-35.8)..(49.6,-35.6)
				..controls(50.5,-35.4)and(51.3,-35.2)..(52.4,-34.6)
				..controls(53.3,-34.1)and(54.5,-33.7)..(54.9,-33.5)
				..controls(55.4,-33.4)and(56.4,-33)..(57.2,-32.6)
				..controls(58.6,-32)and(60.1,-31.5)..(61.6,-31.2)
				..controls(62.3,-31.1)and(62.7,-31)..(63.2,-30.7)
				..controls(64.1,-30)and(64.3,-29.9)..(65.3,-29.7)
				..controls(65.7,-29.6)and(66.5,-29.2)..(67.1,-28.9)
				..controls(67.6,-28.6)and(68.8,-28.2)..(69.6,-27.9)
				..controls(70.5,-27.7)and(71.6,-27.3)..(72.1,-27)
				..controls(72.6,-26.7)and(73.8,-26.2)..(74.8,-25.9)
				..controls(75.7,-25.5)and(76.7,-25.1)..(77.1,-25)
				..controls(77.4,-24.8)and(78.5,-24.4)..(79.6,-24.1)
				..controls(80.6,-23.8)and(82,-23.2)..(82.8,-22.9)
				..controls(83.5,-22.5)and(84.8,-22)..(85.8,-21.8)
				..controls(86.7,-21.5)and(87.8,-21.1)..(88.1,-20.9)
				..controls(88.4,-20.7)and(89.3,-20.3)..(90.1,-20.1)
				..controls(90.8,-19.8)and(91.9,-19.4)..(92.4,-19.1)
				..controls(93.5,-18.5)and(94.4,-18.2)..(95.4,-18.1)
				..controls(95.7,-18)and(96.4,-17.8)..(96.9,-17.5)
				..controls(97.5,-17.3)and(98.4,-16.9)..(99,-16.6)
				..controls(99.7,-16.3)and(100.9,-15.9)..(101.9,-15.6)
				..controls(102.9,-15.3)and(104.4,-14.7)..(105.3,-14.3)
				..controls(108.1,-13.2)and(108.5,-13.1)..(110.1,-13.1)
				..controls(111.4,-13.1)and(111.5,-13.1)..(112.4,-13.5)
				..controls(113.2,-14)and(115,-14.8)..(116.6,-15.4)
				..controls(116.9,-15.5)and(117.9,-15.8)..(118.8,-16.1)
				..controls(120.9,-16.6)and(121.5,-16.9)..(122.6,-17.9)
				..controls(123,-18.3)and(123.7,-18.8)..(123.9,-18.9)
				..controls(124.6,-19.4)and(126,-19.8)..(127.9,-20.2)
				..controls(129.6,-20.5)and(129.8,-20.6)..(131.1,-21.4)
				..controls(131.5,-21.7)and(132.4,-22.1)..(133.1,-22.3)
				..controls(136.1,-23.2)and(138.3,-24.1)..(140.2,-25)
				..controls(141,-25.5)and(141.5,-25.6)..(142.5,-25.8)
				..controls(145.1,-26.3)and(146,-26.6)..(149.4,-28.1)
				..controls(151.3,-28.9)and(153.5,-29.8)..(154.3,-30.1)
				..controls(155.1,-30.4)and(156.2,-30.8)..(156.6,-31)
				..controls(157.1,-31.2)and(158.2,-31.6)..(159.1,-31.8)
				..controls(160.6,-32.2)and(160.9,-32.3)..(161.7,-32.9)
				..controls(162.4,-33.3)and(162.9,-33.6)..(163.7,-33.8)
				..controls(164.3,-34)and(164.8,-34.1)..(164.9,-34.1)
				..controls(165,-34.1)and(165.6,-34.4)..(166.4,-34.8)
				..controls(167.3,-35.2)and(168.1,-35.4)..(169,-35.6)
				..controls(170,-35.8)and(170.5,-36)..(171.7,-36.6)
				..controls(172.4,-37)and(173.3,-37.4)..(173.6,-37.4)
				..controls(175.6,-37.7)and(176.2,-38)..(176.7,-39.1)
				..controls(176.9,-39.6)and(176.9,-39.7)..(176.7,-40.4)
				..controls(176.4,-41.6)and(175.8,-42.7)..(175.2,-43.4)
				..controls(174.5,-44.1)and(172.6,-45.5)..(171.7,-45.8)
				..controls(170.6,-46.2)and(169.1,-46)..(168.2,-45.4)
				..controls(167.9,-45.2)and(167.3,-44.8)..(166.9,-44.4)
				..controls(166.4,-44.1)and(165.8,-43.6)..(165.4,-43.2)
				..controls(165.1,-42.9)and(164.4,-42.3)..(163.9,-41.9)
				..controls(163.3,-41.5)and(162.7,-40.9)..(162.4,-40.5)
				..controls(161.4,-39.2)and(159.9,-38)..(158.8,-37.8)
				..controls(158.4,-37.8)and(158.3,-37.7)..(158.2,-37.2)
				..controls(158,-36.8)and(157.4,-36.1)..(157.1,-36.1)
				..controls(157,-36.1)and(156.7,-36.1)..(156.5,-36)
				..controls(156.1,-35.9)and(155.9,-35.7)..(155.7,-35.2)
				..controls(155.3,-34.6)and(155.1,-34.4)..(154,-34.1)
				..controls(153.6,-34)and(153.2,-33.7)..(152.9,-33.5)
				..controls(152.4,-32.9)and(151.6,-32.4)..(151.3,-32.4)
				..controls(150.9,-32.4)and(150.4,-32.6)..(150.4,-32.9)
				..controls(150.4,-33.2)and(150.8,-33.4)..(151.4,-33.6)
				..controls(152.1,-33.7)and(152.8,-34.1)..(153.5,-34.9)
				..controls(154.1,-35.5)and(154.3,-35.6)..(155,-35.8)
				..controls(155.7,-36)and(155.7,-36)..(156,-36.6)
				..controls(156.3,-37.3)and(156.7,-37.7)..(157.5,-37.8)
				..controls(158.1,-37.9)and(158.2,-38)..(158.4,-38.8)
				..controls(158.5,-39.4)and(158.6,-39.6)..(159.6,-40.4)
				..controls(160.2,-41)and(160.9,-41.6)..(161.3,-41.8)
				..controls(161.7,-42)and(161.9,-42.3)..(162.2,-42.8)
				..controls(162.4,-43.2)and(162.8,-43.8)..(163.1,-44.1)
				..controls(163.9,-44.8)and(165.1,-46.7)..(165.7,-48)
				..controls(166.2,-49.3)and(166.4,-49.7)..(167.2,-50.5)
				..controls(167.5,-50.8)and(167.9,-51.4)..(168.1,-51.8)
				..controls(168.3,-52.2)and(168.8,-53)..(169.1,-53.6)
				..controls(169.6,-54.3)and(169.8,-54.9)..(169.9,-55.6)
				..controls(170.4,-57.8)and(170.5,-58.1)..(171.2,-59.5)
				..controls(171.6,-60.3)and(172,-61.5)..(172.2,-62.1)
				..controls(172.3,-62.8)and(172.7,-63.7)..(172.9,-64.1)
				..controls(173.2,-64.6)and(173.4,-65.3)..(173.4,-65.6)
				..controls(173.8,-68.8)and(174.2,-70.6)..(174.9,-72.6)
				..controls(175.5,-74.5)and(175.9,-77.1)..(176.1,-80.2)
				..controls(176.3,-83.4)and(175.5,-84.4)..(173.7,-83.4)--cycle;
				\definecolor{ezs}{RGB}{247,179,26}\fill[ezs](-5.7,-156.6)--(-3.4,-157.5)--(61.9,12.7)--(59.7,13.6)--cycle;
				\definecolor{ezs}{RGB}{247,179,22}\fill[ezs](-3.8,-157.3)--(-1.5,-158.2)--(63.9,12)--(61.6,12.9)--cycle;
				\definecolor{ezs}{RGB}{248,179,19}\fill[ezs](-1.8,-158.1)--(0.4,-158.9)--(65.8,11.2)--(63.5,12.1)--cycle;
				\definecolor{ezs}{RGB}{248,178,16}\fill[ezs](0.1,-158.8)--(2.4,-159.7)--(67.7,10.5)--(65.4,11.4)--cycle;
				\definecolor{ezs}{RGB}{249,178,15}\fill[ezs](2,-159.5)--(4.3,-160.4)--(69.6,9.8)--(67.3,10.7)--cycle;
				\definecolor{ezs}{RGB}{251,177,17}\fill[ezs](3.9,-160.3)--(6.2,-161.1)--(71.5,9)--(69.2,9.9)--cycle;
				\definecolor{ezs}{RGB}{252,179,18}\fill[ezs](5.8,-161)--(8.1,-161.9)--(73.4,8.3)--(71.1,9.2)--cycle;
				\definecolor{ezs}{RGB}{252,182,17}\fill[ezs](7.7,-161.7)--(10,-162.6)--(75.3,7.6)--(73,8.5)--cycle;
				\definecolor{ezs}{RGB}{252,185,17}\fill[ezs](9.6,-162.5)--(11.9,-163.4)--(77.2,6.8)--(75,7.7)--cycle;
				\definecolor{ezs}{RGB}{251,189,18}\fill[ezs](11.5,-163.2)--(13.8,-164.1)--(79.2,6.1)--(76.9,7)--cycle;
				\definecolor{ezs}{RGB}{251,191,22}\fill[ezs](13.4,-163.9)--(15.7,-164.8)--(81.1,5.4)--(78.8,6.2)--cycle;
				\definecolor{ezs}{RGB}{252,193,25}\fill[ezs](15.4,-164.7)--(17.6,-165.6)--(83,4.6)--(80.7,5.5)--cycle;
				\definecolor{ezs}{RGB}{253,195,27}\fill[ezs](17.3,-165.4)--(19.6,-166.3)--(84.9,3.9)--(82.6,4.8)--cycle;
				\definecolor{ezs}{RGB}{253,196,31}\fill[ezs](19.2,-166.1)--(21.5,-167)--(86.8,3.2)--(84.5,4)--cycle;
				\definecolor{ezs}{RGB}{254,199,35}\fill[ezs](21.1,-166.9)--(23.4,-167.8)--(88.7,2.4)--(86.4,3.3)--cycle;
				\definecolor{ezs}{RGB}{253,201,38}\fill[ezs](23,-167.6)--(25.3,-168.5)--(90.6,1.7)--(88.3,2.6)--cycle;
				\definecolor{ezs}{RGB}{253,203,40}\fill[ezs](24.9,-168.3)--(27.2,-169.2)--(92.5,1)--(90.3,1.8)--cycle;
				\definecolor{ezs}{RGB}{253,205,42}\fill[ezs](26.8,-169.1)--(29.1,-170)--(94.4,0.2)--(92.2,1.1)--cycle;
				\definecolor{ezs}{RGB}{253,207,45}\fill[ezs](28.7,-169.8)--(31,-170.7)--(96.4,-0.5)--(94.1,0.4)--cycle;
				\definecolor{ezs}{RGB}{253,210,49}\fill[ezs](30.7,-170.5)--(32.9,-171.4)--(98.3,-1.2)--(96,-0.4)--cycle;
				\definecolor{ezs}{RGB}{253,212,54}\fill[ezs](32.6,-171.3)--(34.9,-172.2)--(100.2,-2)--(97.9,-1.1)--cycle;
				\definecolor{ezs}{RGB}{253,214,59}\fill[ezs](34.5,-172)--(36.8,-172.9)--(102.1,-2.7)--(99.8,-1.8)--cycle;
				\definecolor{ezs}{RGB}{253,216,63}\fill[ezs](36.4,-172.7)--(38.7,-173.6)--(104,-3.4)--(101.7,-2.6)--cycle;
				\definecolor{ezs}{RGB}{253,218,65}\fill[ezs](38.3,-173.5)--(40.6,-174.4)--(105.9,-4.2)--(103.6,-3.3)--cycle;
				\definecolor{ezs}{RGB}{254,220,66}\fill[ezs](40.2,-174.2)--(42.5,-175.1)--(107.8,-4.9)--(105.5,-4)--cycle;
				\definecolor{ezs}{RGB}{254,222,68}\fill[ezs](42.1,-175)--(44.4,-175.8)--(109.7,-5.6)--(107.5,-4.8)--cycle;
				\definecolor{ezs}{RGB}{254,224,71}\fill[ezs](44,-175.7)--(46.3,-176.6)--(111.7,-6.4)--(109.4,-5.5)--cycle;
				\definecolor{ezs}{RGB}{254,226,75}\fill[ezs](46,-176.4)--(48.2,-177.3)--(113.6,-7.1)--(111.3,-6.2)--cycle;
				\definecolor{ezs}{RGB}{254,228,79}\fill[ezs](47.9,-177.2)--(50.1,-178)--(115.5,-7.8)--(113.2,-7)--cycle;
				\definecolor{ezs}{RGB}{254,230,82}\fill[ezs](49.8,-177.9)--(52.1,-178.8)--(117.4,-8.6)--(115.1,-7.7)--cycle;
				\definecolor{ezs}{RGB}{254,232,85}\fill[ezs](51.7,-178.6)--(54,-179.5)--(119.3,-9.3)--(117,-8.4)--cycle;
				\definecolor{ezs}{RGB}{254,234,88}\fill[ezs](53.6,-179.4)--(55.9,-180.2)--(121.2,-10)--(118.9,-9.2)--cycle;
				\definecolor{ezs}{RGB}{254,237,91}\fill[ezs](55.5,-180.1)--(57.8,-181)--(123.1,-10.8)--(120.8,-9.9)--cycle;
				\definecolor{ezs}{RGB}{254,239,94}\fill[ezs](57.4,-180.8)--(59.7,-181.7)--(125,-11.5)--(122.8,-10.6)--cycle;
				\definecolor{ezs}{RGB}{254,241,97}\fill[ezs](59.3,-181.6)--(61.6,-182.4)--(127,-12.2)--(124.7,-11.4)--cycle;
				\definecolor{ezs}{RGB}{254,243,101}\fill[ezs](61.2,-182.3)--(63.5,-183.2)--(128.9,-13)--(126.6,-12.1)--cycle;
				\definecolor{ezs}{RGB}{254,245,103}\fill[ezs](63.2,-183)--(65.4,-183.9)--(130.8,-13.7)--(128.5,-12.8)--cycle;
				\definecolor{ezs}{RGB}{254,245,105}\fill[ezs](65.1,-183.8)--(67.4,-184.6)--(132.7,-14.4)--(130.4,-13.6)--cycle;
				\definecolor{ezs}{RGB}{254,245,105}\fill[ezs](67,-184.5)--(69.3,-185.4)--(134.6,-15.2)--(132.3,-14.3)--cycle;
				\definecolor{ezs}{RGB}{254,244,103}\fill[ezs](68.9,-185.2)--(71.2,-186.1)--(136.5,-15.9)--(134.2,-15)--cycle;
				\definecolor{ezs}{RGB}{253,243,101}\fill[ezs](70.8,-186)--(73.1,-186.8)--(138.4,-16.6)--(136.1,-15.8)--cycle;
				\definecolor{ezs}{RGB}{253,241,98}\fill[ezs](72.7,-186.7)--(75,-187.6)--(140.3,-17.4)--(138,-16.5)--cycle;
				\definecolor{ezs}{RGB}{253,239,95}\fill[ezs](74.6,-187.4)--(76.9,-188.3)--(142.2,-18.1)--(140,-17.2)--cycle;
				\definecolor{ezs}{RGB}{253,236,92}\fill[ezs](76.5,-188.2)--(78.8,-189)--(144.2,-18.8)--(141.9,-18)--cycle;
				\definecolor{ezs}{RGB}{253,234,88}\fill[ezs](78.5,-188.9)--(80.7,-189.8)--(146.1,-19.6)--(143.8,-18.7)--cycle;
				\definecolor{ezs}{RGB}{253,232,85}\fill[ezs](80.4,-189.6)--(82.7,-190.5)--(148,-20.3)--(145.7,-19.4)--cycle;
				\definecolor{ezs}{RGB}{253,230,82}\fill[ezs](82.3,-190.4)--(84.6,-191.2)--(149.9,-21)--(147.6,-20.2)--cycle;
				\definecolor{ezs}{RGB}{253,228,78}\fill[ezs](84.2,-191.1)--(86.5,-192)--(151.8,-21.8)--(149.5,-20.9)--cycle;
				\definecolor{ezs}{RGB}{253,226,74}\fill[ezs](86.1,-191.8)--(88.4,-192.7)--(153.7,-22.5)--(151.4,-21.6)--cycle;
				\definecolor{ezs}{RGB}{253,224,70}\fill[ezs](88,-192.6)--(90.3,-193.4)--(155.6,-23.3)--(153.3,-22.4)--cycle;
				\definecolor{ezs}{RGB}{253,222,67}\fill[ezs](89.9,-193.3)--(92.2,-194.2)--(157.5,-24)--(155.3,-23.1)--cycle;
				\definecolor{ezs}{RGB}{253,220,65}\fill[ezs](91.8,-194)--(94.1,-194.9)--(159.5,-24.7)--(157.2,-23.8)--cycle;
				\definecolor{ezs}{RGB}{253,217,63}\fill[ezs](93.7,-194.8)--(96,-195.6)--(161.4,-25.5)--(159.1,-24.6)--cycle;
				\definecolor{ezs}{RGB}{253,215,60}\fill[ezs](95.7,-195.5)--(97.9,-196.4)--(163.3,-26.2)--(161,-25.3)--cycle;
				\definecolor{ezs}{RGB}{253,213,56}\fill[ezs](97.6,-196.2)--(99.9,-197.1)--(165.2,-26.9)--(162.9,-26)--cycle;
				\definecolor{ezs}{RGB}{252,211,53}\fill[ezs](99.5,-197)--(101.8,-197.8)--(167.1,-27.7)--(164.8,-26.8)--cycle;
				\definecolor{ezs}{RGB}{252,209,49}\fill[ezs](101.4,-197.7)--(103.7,-198.6)--(169,-28.4)--(166.7,-27.5)--cycle;
				\definecolor{ezs}{RGB}{252,207,45}\fill[ezs](103.3,-198.4)--(105.6,-199.3)--(170.9,-29.1)--(168.6,-28.2)--cycle;
				\definecolor{ezs}{RGB}{253,205,42}\fill[ezs](105.2,-199.2)--(107.5,-200)--(172.8,-29.9)--(170.6,-29)--cycle;
				\definecolor{ezs}{RGB}{253,202,40}\fill[ezs](107.1,-199.9)--(109.4,-200.8)--(174.7,-30.6)--(172.5,-29.7)--cycle;
				\definecolor{ezs}{RGB}{253,200,38}\fill[ezs](109,-200.6)--(111.3,-201.5)--(176.7,-31.3)--(174.4,-30.4)--cycle;
				\definecolor{ezs}{RGB}{253,198,35}\fill[ezs](111,-201.4)--(113.2,-202.2)--(178.6,-32.1)--(176.3,-31.2)--cycle;
				\definecolor{ezs}{RGB}{253,196,31}\fill[ezs](112.9,-202.1)--(115.2,-203)--(180.5,-32.8)--(178.2,-31.9)--cycle;
				\definecolor{ezs}{RGB}{253,194,28}\fill[ezs](114.8,-202.8)--(117.1,-203.7)--(182.4,-33.5)--(180.1,-32.6)--cycle;
				\definecolor{ezs}{RGB}{253,192,26}\fill[ezs](116.7,-203.6)--(119,-204.5)--(184.3,-34.3)--(182,-33.4)--cycle;
				\definecolor{ezs}{RGB}{253,189,24}\fill[ezs](118.6,-204.3)--(120.9,-205.2)--(186.2,-35)--(183.9,-34.1)--cycle;
				\definecolor{ezs}{RGB}{253,186,20}\fill[ezs](120.5,-205)--(122.8,-205.9)--(188.1,-35.7)--(185.8,-34.8)--cycle;
				\definecolor{ezs}{RGB}{253,184,17}\fill[ezs](122.4,-205.8)--(124.7,-206.7)--(190,-36.5)--(187.8,-35.6)--cycle;
				\definecolor{ezs}{RGB}{252,183,15}\fill[ezs](124.3,-206.5)--(126.6,-207.4)--(192,-37.2)--(189.7,-36.3)--cycle;
				\definecolor{ezs}{RGB}{252,181,14}\fill[ezs](126.3,-207.2)--(128.5,-208.1)--(193.9,-37.9)--(191.6,-37.1)--cycle;
				\definecolor{ezs}{RGB}{252,179,13}\fill[ezs](128.2,-208)--(130.4,-208.9)--(195.8,-38.7)--(193.5,-37.8)--cycle;
				\definecolor{ezs}{RGB}{252,178,12}\fill[ezs](130.1,-208.7)--(132.4,-209.6)--(197.7,-39.4)--(195.4,-38.5)--cycle;
				\definecolor{ezs}{RGB}{251,178,10}\fill[ezs](132,-209.4)--(134.3,-210.3)--(199.6,-40.1)--(197.3,-39.3)--cycle;
				\definecolor{ezs}{RGB}{251,178,10}\fill[ezs](133.9,-210.2)--(136.2,-211.1)--(201.5,-40.9)--(199.2,-40)--cycle;
				\definecolor{ezs}{RGB}{250,177,11}\fill[ezs](135.8,-210.9)--(138.1,-211.8)--(203.4,-41.6)--(201.1,-40.7)--cycle;
				\definecolor{ezs}{RGB}{250,178,11}\fill[ezs](137.7,-211.6)--(140,-212.5)--(205.3,-42.3)--(203.1,-41.5)--cycle;
				\definecolor{ezs}{RGB}{248,180,10}\fill[ezs](139.6,-212.4)--(141.9,-213.3)--(207.3,-43.1)--(205,-42.2)--cycle;
				\definecolor{ezs}{RGB}{247,181,9}\fill[ezs](141.5,-213.1)--(143.8,-214)--(209.2,-43.8)--(206.9,-42.9)--cycle;
				\definecolor{ezs}{RGB}{246,182,9}\fill[ezs](143.5,-213.8)--(145.7,-214.7)--(211.1,-44.5)--(208.8,-43.7)--cycle;
				\definecolor{ezs}{RGB}{246,182,10}\fill[ezs](145.4,-214.6)--(147.7,-215.5)--(213,-45.3)--(210.7,-44.4)--cycle;
			\end{scope}
			% chữ EDUZY (Times đậm, làm dày nét bằng 5 lớp lệch nhẹ; căn đúng khung chữ logo gốc)
			\foreach \dx/\dy in {-0.7/0,0.7/0,0/-0.3,0/0.3,0/0}
			\node[anchor=center,inner sep=0pt,outer sep=0pt,text=ezText,align=center] at (109.6+\dx,-193.6+\dy)
			{\resizebox{1.01cm}{0.22cm}{\fontencoding{T1}\fontfamily{ptm}\fontseries{b}\selectfont EDUZY}};
		\end{tikzpicture}%
	}%
}
%====================== THIẾT LẬP THÂN ĐỀ ======================
%--- Môi trường tự luận {bt}: ex_test 3.3.4 chưa khai báo sẵn nên tạo ở đây
%    (ex_test đã tự móc \fixshowans và \listenumerate cho bt)
\makeatletter
\@ifundefined{bt}{\newtheorem{bt}{\color{gold1}\namebt}}{}
\makeatother
\renewcommand{\sectionmark}[1]{\markright{\thesection\ #1}}
%--- Khung tiêu đề (logo + tối đa 4 dòng; dòng 3, 4 để trống {} thì bỏ qua)
%  Đề thi:    \khungtieude{THI THỬ THPT QUỐC GIA NĂM 2026}{Môn Toán -- Đề số 21}
%                         {(Đề thi gồm: ...)}{Thời gian làm bài: 90 phút (không kể thời gian phát đề).}
%  Chuyên đề: \khungtieude{CHUYÊN ĐỀ 3 -- DÃY SỐ VÀ CÔNG THỨC TRUY HỒI}
%                         {Bồi dưỡng học sinh giỏi Toán THPT}{...}{}
\newcommand{\khungtieude}[4]{%
	\begin{tcolorbox}[colframe=gold1,colback=gold5!45,boxrule=0.5mm]
		\begin{minipage}[c]{3cm}
			\centerline{\eduzylogo[2.4cm]}
		\end{minipage}\hspace{0.5cm}
		\begin{minipage}[c]{12.5cm}
			\vskip 0.15 cm
			\centerline{\ind{#1}}
			\vskip 0.15 cm
			\centerline{\textbf{#2}}
			\if\relax\detokenize{#3}\relax\else
				\vskip 0.15 cm
				\centerline{\textit{#3}}
			\fi
			\if\relax\detokenize{#4}\relax\else
				\vskip 0.15 cm
				\centerline{\textit{#4}}
			\fi
			\vskip 0.15 cm
		\end{minipage}
	\end{tcolorbox}
}
%====================== BỔ SUNG CHO BÀI HỌC, CHUYÊN ĐỀ ======================
\usepackage[shortlabels]{enumitem}
\usepackage{tabularx}
%--- Danh sách gọn: \begin{enumerate}[gon] ... \end{enumerate}
\SetEnumitemKey{gon}{leftmargin=*,itemsep=1pt,parsep=0pt,topsep=3pt}
%--- Khung lý thuyết: \begin{khung}{} ... \end{khung}; tiêu đề chỉ dùng khi thật ngắn ("Mẹo", "Lỗi thường gặp")
\newtcolorbox{khung}[1]{breakable,enhanced,colframe=gold1,colback=gold5!45,
	boxrule=0.4mm,fonttitle=\bfseries\sffamily,coltitle=white,colbacktitle=gold1,
	attach boxed title to top left={xshift=6mm,yshift=-3mm},
	boxed title style={arc=1mm,boxrule=0.3mm},title={#1}}
%--- Tiêu đề mục nhỏ trong một dạng: \muc{Dấu hiệu}, \muc{Phương pháp}, \muc{Bài tự luyện}
\makeatletter
\newcommand{\muc}[1]{\par\medskip\noindent\indam{#1}\par\nobreak\smallskip\@afterindentfalse\@afterheading}
\makeatother
%--- Ô công thức cần nhớ (căn giữa, viền vàng, vừa khít nội dung)
%     \begin{center}\congthuc{$Q(n)=(n-1)!$}\end{center}
\newtcbox{\congthuc}{on line,colframe=gold1,colback=gold5!45,boxrule=0.5pt,arc=2pt,
	left=8pt,right=8pt,top=4pt,bottom=4pt}
%--- Công tắc lời giải bài tự luyện (ví dụ mẫu luôn dùng \loigiai)
%     \loigiaiTL[{đáp số}]{lời giải đầy đủ}
%     \TLfull -> in lời giải đầy đủ;  \TLdapso -> chỉ in một dòng "Đáp số";  \TLan -> ẩn hoàn toàn
%     (\CheDo là dethi thì mọi lời giải và đáp số đều ẩn)
\makeatletter
\newif\ifTLhien   \TLhientrue
\newif\ifTLdapso  \TLdapsofalse
\newcommand{\TLfull}{\TLhientrue\TLdapsofalse}
\newcommand{\TLdapso}{\TLhientrue\TLdapsotrue}
\newcommand{\TLan}{\TLhienfalse}
\newcommand{\loigiaiTL}[2][]{%
	\ifTLhien
		\ifTLdapso
			\def\@tlarg{#1}%
			\ifx\@tlarg\@empty\else
				\ifdefstring{\kindOPTN}{dethi}{}{\par\noindent{\bfseries\color{gold1}Đáp số:} #1\par}%
			\fi
		\else
			\loigiai{#2}%
		\fi
	\fi}
\makeatother
\TLdapso% mặc định: chỉ in đáp số bài tự luyện
%====================== BẢNG ĐÚNG/SAI (\choiceTFt) ======================
%  Hàng tiêu đề nền nâu đồng chữ trắng (Phát biểu | Đ | S), các hàng tô xen kẽ kem / trắng,
%  đường kẻ vàng kim, nhãn a) b) c) d) màu nâu đồng. Bản có đáp án: cột được chọn có vòng tròn
%  nâu đồng với dấu tích trắng (câu đúng tích ở cột Đúng, câu sai tích ở cột Sai).
%  Tiêu đề ghi đủ "Đúng | Sai": \OPTN{viettat=0} sau \input (cột Đúng/Sai rộng ra, phát biểu hẹp lại).
%  Dùng bảng gốc của ex_test: \def\BangDSgoc{1} trước \input{preamble-eduzy-v2}.
\providecommand{\BangDSgoc}{0}
\ifnum\BangDSgoc=0
\usepackage{colortbl,pifont}
\colorlet{ezDSke}{gold2!60!gold1}% đường kẻ bảng (nét 0,4pt như ex_test để không đổi chiều cao bảng)
\colorlet{ezDSle}{gold5!60}% hàng a), c)
\colorlet{ezDSchan}{white}% hàng b), d)
%--- dấu chọn ở cột Đúng/Sai
%    (bề rộng 0, căn giữa ô: không làm đổi bề rộng cột của ex_test, không báo tràn lề)
\newcommand{\EZdauchon}{\makebox[0pt]{\tikz[baseline=-1ex]\node[circle,fill=gold1,inner sep=0pt,
	minimum size=1.2em,text=white,font=\small]{\ding{51}};}}
\def\tickT{\EZdauchon}\def\tickF{\EZdauchon}
\AtBeginDocument{\sbox{\ccTFTab}{\EZdauchon}}
%--- nhãn a) b) c) d) màu nâu đồng (ở bản lời giải, nhãn câu đúng vẫn đóng khung như ex_test)
\renewcommand{\FalseTF}{\stepcounter{dapan}\textbf{\color{gold1}\DapAnTF\sepTF}\ignorespaces}
%--- trước khi dựng bảng: kẻ màu vàng kim (bề rộng cột giữ nguyên như ex_test); sau bảng trả lại màu đen
\newcount\EZhangTF
\newcommand{\EZbatdauTF}{%
	\global\EZhangTF=0\relax
	\arrayrulecolor{ezDSke}}
\newcommand{\EZketthucTF}{\arrayrulecolor{black}}
%--- hàng tiêu đề nền nâu đồng, chữ trắng
\makeatletter
\newcommand{\EZo@TF}[1]{\centering\bfseries\color{white}#1}
\newcommand{\EZtieudeTF}{%
	\ifcase\cotDS
		\ifnum\cotSTT=1\relax
			\def\rowHeader{\hline\rowcolor{gold1}\EZo@TF{TT}&\EZo@TF{\phatbieu}\tabularnewline}%
		\else
			\def\rowHeader{\hline\rowcolor{gold1}\EZo@TF{\phatbieu}\tabularnewline}%
		\fi
	\or
		\ifnum\cotSTT=1\relax
			\def\rowHeader{\hline\rowcolor{gold1}\EZo@TF{TT}&\EZo@TF{\phatbieu}&\EZo@TF{\Dung}\tabularnewline}%
		\else
			\def\rowHeader{\hline\rowcolor{gold1}\EZo@TF{\phatbieu}&\EZo@TF{\Dung}\tabularnewline}%
		\fi
	\or
		\ifnum\cotSTT=1\relax
			\def\rowHeader{\hline\rowcolor{gold1}\EZo@TF{TT}&\EZo@TF{\phatbieu}&\EZo@TF{\Dung}&\EZo@TF{\Sai}\tabularnewline}%
		\else
			\def\rowHeader{\hline\rowcolor{gold1}\EZo@TF{\phatbieu}&\EZo@TF{\Dung}&\EZo@TF{\Sai}\tabularnewline}%
		\fi
	\fi}
%--- tô xen kẽ từng hàng (gọi ngay sau \hline đầu mỗi hàng)
\newcommand{\EZmauhangTF}{%
	\noalign{\global\advance\EZhangTF\@ne}%
	\ifodd\EZhangTF\rowcolor{ezDSle}\else\rowcolor{ezDSchan}\fi}
%--- gắn vào \choiceTFt / \writeChoiceTFt của ex_test
\xpretocmd{\choiceTFt}{\EZbatdauTF}{}{\typeout{[preamble-eduzy-v2] Không gắn được bảng Đúng/Sai (đầu bảng)}}
\xpatchcmd{\choiceTFt}{\ifnum\dongPB=0\def\rowHeader{}\fi}{\EZtieudeTF\ifnum\dongPB=0\def\rowHeader{}\fi}
	{}{\typeout{[preamble-eduzy-v2] Không gắn được bảng Đúng/Sai (hàng tiêu đề)}}
\xpatchcmd{\writeChoiceTFt}{\hline\xdef\dotEX{\sdotEX}}{\hline\EZmauhangTF\xdef\dotEX{\sdotEX}}
	{}{\typeout{[preamble-eduzy-v2] Không gắn được bảng Đúng/Sai (tô hàng)}}
\xpatchcmd{\writeChoiceTFt}{{\end{longtable}}}{{\end{longtable}\EZketthucTF}}
	{}{\typeout{[preamble-eduzy-v2] Không gắn được bảng Đúng/Sai (cuối longtable)}}
\xpatchcmd{\writeChoiceTFt}{{\end{xtabular}\end{center}\medskip\noindent}}{{\end{xtabular}\end{center}\EZketthucTF\medskip\noindent}}
	{}{\typeout{[preamble-eduzy-v2] Không gắn được bảng Đúng/Sai (cuối xtabular)}}
\fi
%=====================================================================
\renewcommand{\parallel}{\mathrel{/\!\!/}}

%      \TLdapso                      % tài liệu: \TLfull | \TLdapso | \TLan
%      \begin{document} ... \end{document}
%  Logo Eduzy vẽ bằng TikZ ngay trong file (\eduzylogo), không cần file ảnh.
%  Biên dịch: pdfLaTeX (máy có gói vietnam) hoặc XeLaTeX.
%============================================================
\documentclass[12pt,a4paper,twoside]{article}%<-- twoside: bat dau trang/chan trang chan-le
\usepackage{iftex}
\usepackage{amsmath,amssymb,tikz,mathrsfs,tkz-tab}
\usepackage{graphicx}
\usepackage{pgfplots}% đồ thị hàm số
\pgfplotsset{compat=1.17}
%\usepackage{mathptmx}
\usepackage{mathpazo}
\ifXeTeX% XeLaTeX: font chữ có đủ dấu tiếng Việt, giữ nguyên font toán mathpazo
	\usepackage[no-math]{fontspec}
	\setmainfont{TeX Gyre Pagella}
	\newfontfamily\EZqag{TeX Gyre Adventor}[NFSSFamily=qag,Ligatures=TeX]
\fi
\usepackage[many]{tcolorbox}
\usetikzlibrary{math,through,calc,intersections,angles,quotes,shapes,shapes.geometric,arrows,patterns,matrix,chains,arrows.meta,decorations.shapes,decorations.fractals,decorations.markings,shadows,decorations.pathreplacing}
\usepackage[top=1.6cm, bottom=1.8cm, left=1.8cm, right=1.8cm]{geometry}% v2: top 1.6cm (bản cũ 1.5cm) để chữ không sát dải đầu trang
\providecommand{\CheDo}{loigiai}% loigiai | dethi | solcolor | color | book
\expandafter\usepackage\expandafter[\CheDo]{ex_test}
% \tenchuong: khai báo ở file chính (trước \input) hoặc \renewcommand sau \input
%============================================================
%  [GỘP TỪ eduzy-de.tex] EDUZY - MÀU NÂU ĐỒNG / VÀNG KIM
%  + ĐẦU TRANG, CHÂN TRANG KIỂU DẢI NGANG (v2), ĐỐI XỨNG TRANG CHẴN / TRANG LẺ.
%  (không cần file eduzy-de.tex riêng nữa)
%  Khối này phải đặt SAU \usepackage[...]{ex_test} và SAU \def\tenchuong{...}
%============================================================
\usepackage{lastpage}
\usepackage{needspace}
%--- DÀN TRANG ĐỀ THI: tách câu (THÊM 05/11/2026; chỉ thêm lệnh mới, KHÔNG đổi hành vi mặc định) ---
%  Hai công tắc, mặc định tắt (= preamble chạy y như cũ). Bật bằng cách đặt TRƯỚC %============================================================
%  PREAMBLE EDUZY v2 - dùng chung cho đề thi, bài học, chuyên đề (gói ex_test)
%  Khác bản preamble-eduzy.tex ở đầu trang, chân trang: dải nền nâu đồng tràn ngang,
%  logo + Eduzy.com.vn, ô chữ vàng kem cạnh xiên, chân trang "Trang (số)";
%  logo Eduzy chìm giữa mọi trang (tắt: \logochimtat, bật lại: \logochimbat).
%  Cách dùng ở file chính (xem mau-chuan.tex, mau-chuyen-de.tex):
%      \def\CheDo{loigiai}          % loigiai | dethi | solcolor | color | book
%      \def\tenchuong{<dòng chữ trên đầu trang>}
%      \def\thuonghieu{<chữ ở chân trang lẻ>}   % không khai báo thì dùng mặc định
%      \def\tacgia{<chữ ở chân trang chẵn>}      % không khai báo thì dùng mặc định
%      \def\khauhieu{<chữ cạnh logo ở đầu trang>}  % tùy chọn, mặc định Eduzy.com.vn
%      \input{preamble-eduzy-v2}
%      \TLdapso                      % tài liệu: \TLfull | \TLdapso | \TLan
%      \begin{document} ... \end{document}
%  Logo Eduzy vẽ bằng TikZ ngay trong file (\eduzylogo), không cần file ảnh.
%  Biên dịch: pdfLaTeX (máy có gói vietnam) hoặc XeLaTeX.
%============================================================
\documentclass[12pt,a4paper,twoside]{article}%<-- twoside: bat dau trang/chan trang chan-le
\usepackage{iftex}
\usepackage{amsmath,amssymb,tikz,mathrsfs,tkz-tab}
\usepackage{graphicx}
\usepackage{pgfplots}% đồ thị hàm số
\pgfplotsset{compat=1.17}
%\usepackage{mathptmx}
\usepackage{mathpazo}
\ifXeTeX% XeLaTeX: font chữ có đủ dấu tiếng Việt, giữ nguyên font toán mathpazo
	\usepackage[no-math]{fontspec}
	\setmainfont{TeX Gyre Pagella}
	\newfontfamily\EZqag{TeX Gyre Adventor}[NFSSFamily=qag,Ligatures=TeX]
\fi
\usepackage[many]{tcolorbox}
\usetikzlibrary{math,through,calc,intersections,angles,quotes,shapes,shapes.geometric,arrows,patterns,matrix,chains,arrows.meta,decorations.shapes,decorations.fractals,decorations.markings,shadows,decorations.pathreplacing}
\usepackage[top=1.6cm, bottom=1.8cm, left=1.8cm, right=1.8cm]{geometry}% v2: top 1.6cm (bản cũ 1.5cm) để chữ không sát dải đầu trang
\providecommand{\CheDo}{loigiai}% loigiai | dethi | solcolor | color | book
\expandafter\usepackage\expandafter[\CheDo]{ex_test}
% \tenchuong: khai báo ở file chính (trước \input) hoặc \renewcommand sau \input
%============================================================
%  [GỘP TỪ eduzy-de.tex] EDUZY - MÀU NÂU ĐỒNG / VÀNG KIM
%  + ĐẦU TRANG, CHÂN TRANG KIỂU DẢI NGANG (v2), ĐỐI XỨNG TRANG CHẴN / TRANG LẺ.
%  (không cần file eduzy-de.tex riêng nữa)
%  Khối này phải đặt SAU \usepackage[...]{ex_test} và SAU \def\tenchuong{...}
%============================================================
\usepackage{lastpage}
\usepackage{needspace}
%--- DÀN TRANG ĐỀ THI: tách câu (THÊM 05/11/2026; chỉ thêm lệnh mới, KHÔNG đổi hành vi mặc định) ---
%  Hai công tắc, mặc định tắt (= preamble chạy y như cũ). Bật bằng cách đặt TRƯỚC \input{preamble-eduzy-v2}
%  trong file đề (tools/dan_trang/dan_trang.py tự ghi khi đặt \NeedspaceCau):
%    \def\TachCauDS{1}   câu Đúng/Sai được tách: câu dẫn (kèm hình) ở cuối trang, bảng Phát biểu sang trang sau
%                        (câu dẫn và hình không tách nhau, bảng không tách giữa hai trang)
%    \def\TachCauLG{1}   bản lời giải (\CheDo khác dethi): phần đề của câu cao hơn 1/3 trang được phép tách
%  Lệnh dùng trong file đề (giá trị do tools/dan_trang/dan_trang.py đo ở bản dethi):
%    \NeedspaceCau[<cao câu dẫn + hình>]{<cao cả khối câu>}   ngay trên \begin{ex} (câu đầu phần: trên tiêu đề phần)
%        công tắc tắt: đúng bằng \Needspace{<cao cả khối câu>} như cũ;
%        \TachCauDS=1 và có [..] (câu Đúng/Sai được chọn tách): chỉ giữ câu dẫn + hình;
%        \TachCauLG=1 ở bản lời giải: đòi tối đa 1/3 \textheight.
%    \NeedspaceBang{<cao bảng Phát biểu>}%   ngay trước \choiceTFt: \TachCauDS=1 thì giữ trọn bảng trên một trang;
%        công tắc tắt: không làm gì.
\providecommand{\TachCauDS}{0}
\providecommand{\TachCauLG}{0}
\newlength{\EZnsCau}
\newcommand{\NeedspaceCau}[2][]{%
	\setlength{\EZnsCau}{#2}%
	\ifnum\TachCauDS=1 \if\relax\detokenize{#1}\relax\else\setlength{\EZnsCau}{#1}\fi\fi
	\ifnum\TachCauLG=1 \ifdefstring{\kindOPTN}{dethi}{}{\ifdim\EZnsCau>0.3333\textheight\setlength{\EZnsCau}{0.3333\textheight}\fi}\fi
	\Needspace{\EZnsCau}}
\newcommand{\NeedspaceBang}[1]{\ifnum\TachCauDS=1 \Needspace{#1}\fi}

%--- Thang màu: đậm -> nhạt
\definecolor{gold1}{RGB}{126,96,0}    % nâu đồng đậm  #7E6000
\definecolor{gold2}{RGB}{192,144,3}   % vàng kim      #C09003
\definecolor{gold3}{RGB}{255,215,99}  % vàng nhạt     #FFD763
\definecolor{gold4}{RGB}{255,231,153} % vàng rất nhạt #FFE799
\definecolor{gold5}{RGB}{255,243,205} % vàng kem      #FFF3CD
%--- Ghi đè tông xanh/đỏ cũ
\definecolor{myblue}{RGB}{126,96,0}
\colorlet{blue}{gold1}
\colorlet{cyan}{gold2}
\colorlet{teal}{gold2}

%--- Tên chương mặc định (file đề nên tự khai báo lại)
\providecommand{\tenchuong}{}
%--- Chân trang
\providecommand{\thuonghieu}{TOÁN 12 -- Eduzy.com.vn}
\providecommand{\tacgia}{Tô Gia Bảo -- 0353703783}

%--- Chữ cạnh logo ở đầu trang
\providecommand{\khauhieu}{Eduzy.com.vn}

%====================== THÔNG SỐ ĐẦU TRANG, CHÂN TRANG (v2) ======================
%  Dải nền đậm tràn ngang trang; ô chữ vàng kem có cạnh xiên, kèm 2 vạch xiên chuyển màu.
%  Đầu trang lẻ : logo + \khauhieu (trái)   | ô chữ \tenchuong (phải)
%  Chân trang lẻ: ô chữ \thuonghieu (trái)  | Trang (số) (phải)
%  Trang chẵn soi gương trang lẻ; chân trang chẵn ghi \tacgia.
%--- Màu: chỉ dùng thang vàng kim Eduzy
\colorlet{ezNen}{gold1}% dải nền đậm
\colorlet{ezVach1}{gold3}% vạch xiên sát dải nền (đậm hơn)
\colorlet{ezVach2}{gold4}% vạch xiên sát ô chữ (nhạt hơn)
\colorlet{ezO}{gold5}% ô chữ vàng kem
%--- Kích thước
\newlength{\EZdai}\setlength{\EZdai}{1.1cm}% chiều cao dải nền
\newlength{\EZo}\setlength{\EZo}{0.68cm}% chiều cao ô chữ
\newlength{\EZxien}\setlength{\EZxien}{0.55cm}% độ lệch ngang của cạnh xiên
\newlength{\EZvach}\setlength{\EZvach}{0.75cm}% bề rộng mỗi vạch xiên
\newlength{\EZkhe}\setlength{\EZkhe}{0.12cm}% khe trắng giữa vạch và ô chữ
\newlength{\EZdem}\setlength{\EZdem}{0.4cm}% khoảng đệm từ chữ tới cạnh xiên
\newlength{\EZle}\setlength{\EZle}{1.8cm}% lề trái/phải của chữ (bằng lề thân bài)
\newlength{\EZrongmin}\setlength{\EZrongmin}{4.5cm}% ô chữ rộng tối thiểu
\newlength{\EZrong}% bề rộng chữ trong ô (tính tự động)
%--- Font
\newcommand{\EZfontO}{\fontsize{10.5pt}{12pt}\selectfont\bfseries}% chữ trong ô
\newcommand{\EZfontKH}{\fontfamily{qag}\fontsize{11pt}{12pt}\selectfont\bfseries}% \khauhieu cạnh logo
\newcommand{\EZfontTrang}{\fontsize{11.5pt}{12pt}\selectfont\bfseries}% chữ "Trang" và số trang
%--- Đo bề rộng chữ trong ô
\newcommand{\EZdorong}[1]{%
	\settowidth{\EZrong}{\EZfontO #1}%
	\ifdim\EZrong<\EZrongmin\setlength{\EZrong}{\EZrongmin}\fi}
%--- Logo + chữ trắng \khauhieu (dựng sẵn một lần khi vào \begin{document})
\newsavebox{\EZlogobox}
\AtBeginDocument{\sbox{\EZlogobox}{%
		\raisebox{-0.5\height}{\eduzylogo[1cm]}\hspace{0.15cm}%
		\raisebox{-0.5\height}{\color{white}\EZfontKH\khauhieu}}}
%--- Hai vạch xiên + khe trắng. #1: điểm gốc ở mép dưới, #2: hướng (1 sang phải, -1 sang trái),
%    #3: hướng nghiêng (1 kiểu "/", -1 kiểu "\")
\newcommand{\EZvachxien}[3]{%
	\fill[white] (#1) --++({(#3)*\EZxien},\EZo) --++({(#2)*\EZkhe},0) --++({-(#3)*\EZxien},-\EZo) --cycle;
	\fill[ezVach2] ([xshift=(#2)*\EZkhe]#1) --++({(#3)*\EZxien},\EZo) --++({(#2)*\EZvach},0) --++({-(#3)*\EZxien},-\EZo) --cycle;
	\fill[ezVach1] ([xshift=(#2)*(\EZkhe+\EZvach)]#1) --++({(#3)*\EZxien},\EZo) --++({(#2)*\EZvach},0) --++({-(#3)*\EZxien},-\EZo) --cycle;
	\draw[white,line width=0.6pt] ([xshift=(#2)*(\EZkhe+\EZvach)]#1) --++({(#3)*\EZxien},\EZo);
}
%--- Chữ trong ô (căn giữa theo chiều cao ô, không lệch vì dấu tiếng Việt)
\tikzset{ezchu/.style={inner sep=0pt,text=ezNen,text height=0.72em,text depth=0.22em}}

%====================== LOGO CHÌM GIỮA TRANG ======================
%  Vẽ cùng đầu trang nên nằm dưới chữ của thân bài, có trên mọi trang.
%  Tắt: \logochimtat (từ trang đó trở đi), bật lại: \logochimbat.
%  Đổi cỡ, độ đậm: khai báo trước \input, ví dụ \def\EZchimcao{9cm} \def\EZchimdo{0.1}
\providecommand{\EZchimcao}{10cm}% chiều cao logo chìm
\providecommand{\EZchimdo}{0.12}% độ đậm (0 = trong suốt, 1 = rõ hẳn)
\newif\ifEZchim \EZchimtrue
\newcommand{\logochimtat}{\EZchimfalse}
\newcommand{\logochimbat}{\EZchimtrue}
%  Dựng sẵn logo mờ một lần khi vào \begin{document}; transparency group làm mờ cả logo
%  một lượt nên không lộ vệt ở chỗ các dải màu chồng nhau.
\newsavebox{\EZchimtho}
\newsavebox{\EZchimbox}
\AtBeginDocument{%
	\sbox{\EZchimtho}{\eduzylogo[\EZchimcao]}%
	\sbox{\EZchimbox}{\begin{tikzpicture}
			\begin{scope}[transparency group,opacity=\EZchimdo]
				\node[inner sep=0pt] {\usebox{\EZchimtho}};
			\end{scope}
	\end{tikzpicture}}}
\newcommand{\EZvechim}{%
	\ifEZchim
		\node[inner sep=0pt] at (current page.center) {\usebox{\EZchimbox}};
	\fi}

%====================== ĐẦU TRANG - TRANG LẺ ======================
%  (logo bên trái, ô chữ \tenchuong bên phải)
\newcommand{\myfancyheadOdd}{%
	\EZdorong{\tenchuong}%
	\begin{tikzpicture}[remember picture,overlay]
		\coordinate (T) at (current page.north west);
		\coordinate (P) at (current page.north east);
		%--- logo chìm giữa trang
		\EZvechim
		%--- dải nền
		\fill[ezNen] (T) rectangle ([yshift=-\EZdai]P);
		%--- ô chữ: X là góc dưới bên trái, cạnh trái xiên kiểu "/"
		\coordinate (X) at ([xshift=-\EZle-\EZrong-\EZdem-\EZxien,yshift=-\EZdai]P);
		\fill[ezO] (X) --++(\EZxien,\EZo) -- ([yshift=-\EZdai+\EZo]P) -- ([yshift=-\EZdai]P) --cycle;
		\EZvachxien{X}{-1}{1}
		\node[ezchu,anchor=east] at ([xshift=-\EZle,yshift=-\EZdai+0.5*\EZo]P) {\EZfontO\tenchuong};
		%--- logo + tên
		\node[inner sep=0pt,anchor=west] at ([xshift=\EZle,yshift=-0.5*\EZdai]T) {\usebox{\EZlogobox}};
	\end{tikzpicture}%
}

%====================== ĐẦU TRANG - TRANG CHẴN =====================
%  (soi gương: ô chữ \tenchuong bên trái, logo bên phải)
\newcommand{\myfancyheadEven}{%
	\EZdorong{\tenchuong}%
	\begin{tikzpicture}[remember picture,overlay]
		\coordinate (T) at (current page.north west);
		\coordinate (P) at (current page.north east);
		%--- logo chìm giữa trang
		\EZvechim
		%--- dải nền
		\fill[ezNen] (T) rectangle ([yshift=-\EZdai]P);
		%--- ô chữ: X là góc dưới bên phải, cạnh phải xiên kiểu "\"
		\coordinate (X) at ([xshift=\EZle+\EZrong+\EZdem+\EZxien,yshift=-\EZdai]T);
		\fill[ezO] (X) --++(-\EZxien,\EZo) -- ([yshift=-\EZdai+\EZo]T) -- ([yshift=-\EZdai]T) --cycle;
		\EZvachxien{X}{1}{-1}
		\node[ezchu,anchor=west] at ([xshift=\EZle,yshift=-\EZdai+0.5*\EZo]T) {\EZfontO\tenchuong};
		%--- logo + tên
		\node[inner sep=0pt,anchor=east] at ([xshift=-\EZle,yshift=-0.5*\EZdai]P) {\usebox{\EZlogobox}};
	\end{tikzpicture}%
}

%====================== CHÂN TRANG - TRANG LẺ ======================
%  (ô chữ \thuonghieu bên trái, "Trang (số)" bên phải)
\newcommand{\myfancyfootOdd}{%
	\EZdorong{\thuonghieu}%
	\begin{tikzpicture}[remember picture,overlay]
		\coordinate (T) at (current page.south west);
		\coordinate (P) at (current page.south east);
		%--- dải nền
		\fill[ezNen] (T) rectangle ([yshift=\EZdai]P);
		%--- ô chữ sát mép trên dải: X là góc dưới bên phải, cạnh phải xiên kiểu "/"
		\coordinate (X) at ([xshift=\EZle+\EZrong+\EZdem,yshift=\EZdai-\EZo]T);
		\fill[ezO] ([yshift=\EZdai-\EZo]T) -- (X) --++(\EZxien,\EZo) -- ([yshift=\EZdai]T) --cycle;
		\EZvachxien{X}{1}{1}
		\node[ezchu,anchor=west] at ([xshift=\EZle,yshift=\EZdai-0.5*\EZo]T) {\EZfontO\thuonghieu};
		%--- số trang
		\node[circle,draw=white,line width=0.9pt,text=white,inner sep=0pt,minimum size=0.62cm,anchor=east]
		(so) at ([xshift=-\EZle,yshift=0.5*\EZdai]P) {\EZfontTrang\thepage};
		\node[inner sep=0pt,text=white,anchor=base east] (tr) at ([xshift=-0.15cm]so.base west) {\EZfontTrang Trang};
	\end{tikzpicture}%
}

%====================== CHÂN TRANG - TRANG CHẴN ====================
%  (soi gương: "Trang (số)" bên trái, ô chữ \tacgia bên phải)
\newcommand{\myfancyfootEven}{%
	\EZdorong{\tacgia}%
	\begin{tikzpicture}[remember picture,overlay]
		\coordinate (T) at (current page.south west);
		\coordinate (P) at (current page.south east);
		%--- dải nền
		\fill[ezNen] (T) rectangle ([yshift=\EZdai]P);
		%--- ô chữ sát mép trên dải: X là góc dưới bên trái, cạnh trái xiên kiểu "\"
		\coordinate (X) at ([xshift=-\EZle-\EZrong-\EZdem,yshift=\EZdai-\EZo]P);
		\fill[ezO] ([yshift=\EZdai-\EZo]P) -- (X) --++(-\EZxien,\EZo) -- ([yshift=\EZdai]P) --cycle;
		\EZvachxien{X}{-1}{-1}
		\node[ezchu,anchor=east] at ([xshift=-\EZle,yshift=\EZdai-0.5*\EZo]P) {\EZfontO\tacgia};
		%--- số trang
		%    (vòng tròn ẩn so0 nằm giữa dải để lấy đường chân chữ, giống trang lẻ)
		\node[circle,inner sep=0pt,minimum size=0.62cm] (so0) at ([yshift=0.5*\EZdai]T) {\EZfontTrang\phantom{\thepage}};
		\node[inner sep=0pt,text=white,anchor=base west] (tr) at ($(so0.base -| T)+(\EZle,0)$) {\EZfontTrang Trang};
		\node[circle,draw=white,line width=0.9pt,text=white,inner sep=0pt,minimum size=0.62cm,anchor=base west]
		(so) at ([xshift=0.15cm]tr.base east) {\EZfontTrang\thepage};
	\end{tikzpicture}%
}

%--- Giữ tương thích với các file cũ đang gọi \myfancyhead, \myfancyfoot
\newcommand{\myfancyhead}{\myfancyheadOdd}
\newcommand{\myfancyfoot}{\myfancyfootOdd}

%======================= GẮN VÀO TRANG =============================
\pagestyle{fancy}
\fancyhf{}
\fancyhead[LO]{\myfancyheadOdd}
\fancyhead[RE]{\myfancyheadEven}
\fancyfoot[LO]{\myfancyfootOdd}
\fancyfoot[RE]{\myfancyfootEven}
\renewcommand{\headrulewidth}{0pt}
\renewcommand{\footrulewidth}{0pt}
\fancypagestyle{plain}{%
	\fancyhf{}
	\fancyhead[LO]{\myfancyheadOdd}
	\fancyhead[RE]{\myfancyheadEven}
	\fancyfoot[LO]{\myfancyfootOdd}
	\fancyfoot[RE]{\myfancyfootEven}
	\renewcommand{\headrulewidth}{0pt}
	\renewcommand{\footrulewidth}{0pt}
}
%=================== HẾT PHẦN GỘP TỪ eduzy-de.tex ===================
\newcommand{\indam}[1]{
	{\fontsize{12pt}{9pt}\sffamily \selectfont\bfseries\color{gold1} #1}
}
\newcommand{\ind}[1]{
	{\fontsize{14pt}{9pt}\fontfamily{qag} \selectfont\bfseries\color{gold1} #1}
}
\renewcommand{\FalseEX}{\stepcounter{dapan}\textbf{\selectfont\color{gold1}\Alph{dapan}.}\ignorespaces}
\renewtheorem{ex}{\color{gold1}Câu}
\renewcommand{\baselinestretch}{1.1}
\newcommand{\hoac}[1]{
	\left[\begin{aligned}#1\end{aligned}\right.}
\newcommand{\heva}[1]{
	\left\{\begin{aligned}#1\end{aligned}\right.}
\usepackage{esvect}
\def\vec{\vv}
\def\vt{\vv}
\def\overrightarrow{\vv}
%====================== LOGO EDUZY VẼ BẰNG TIKZ ======================
%  Đường viền dò từ logo gốc; dải màu vàng kim đo trực tiếp trên ảnh gốc
%  (nghiêng 21 độ, tô bằng 80 dải hẹp nên hiển thị giống nhau trên mọi trình xem PDF).
%  Dùng: \eduzylogo          -> cao 2.4cm (mặc định)
%        \eduzylogo[1.5cm]   -> cao 1.5cm
\definecolor{ezText}{RGB}{245,200,46}
\newcommand{\eduzylogo}[1][2.4cm]{%
	\resizebox{!}{#1}{%
		\begin{tikzpicture}[x=0.01cm,y=0.01cm]
			\path[use as bounding box] (0,0) rectangle (224,-224);
			% mũ tốt nghiệp + vương miện
			\begin{scope}
				\clip
				(119.4,-176.1)..controls(118.9,-176.1)and(116.8,-176)..(114.6,-176)
				..controls(101.4,-176)and(85.8,-175.9)..(83.7,-175.7)
				..controls(81.8,-175.6)and(81,-175.6)..(80.2,-175.7)
				..controls(78.8,-175.9)and(77.7,-175.9)..(76.2,-175.7)
				..controls(74.9,-175.5)and(70.8,-175.5)..(64.3,-175.7)
				..controls(59.1,-175.9)and(58.3,-175.7)..(58.6,-173.7)
				..controls(58.6,-172.8)and(59.6,-169.6)..(60.3,-168)
				..controls(60.8,-166.9)and(61.4,-165.8)..(62.3,-165)
				..controls(63.1,-164.2)and(64.3,-163.4)..(65,-163.2)
				..controls(65.5,-163.1)and(72.3,-163)..(76.6,-163.1)
				..controls(77.9,-163.1)and(80.1,-163.1)..(81.4,-163)
				..controls(82.8,-162.9)and(86.3,-162.9)..(89.2,-162.9)
				..controls(93,-162.9)and(94.9,-162.8)..(96,-162.7)
				..controls(97.2,-162.5)and(101,-162.5)..(114,-162.4)
				..controls(123,-162.3)and(134.9,-162.2)..(140.2,-162.1)
				..controls(150.1,-162)and(152.2,-162.1)..(153.6,-162.4)
				..controls(154.4,-162.6)and(157.1,-163.8)..(158.1,-164.4)
				..controls(158.8,-164.8)and(159.7,-165.7)..(159.9,-166.3)
				..controls(160,-166.5)and(160.2,-167.1)..(160.2,-167.6)
				..controls(160.3,-168)and(160.6,-169.4)..(160.9,-170.6)
				..controls(161.2,-171.7)and(161.5,-173.1)..(161.6,-173.6)
				..controls(161.8,-174.4)and(161.8,-174.5)..(161.5,-174.9)
				..controls(161.2,-175.6)and(160.8,-175.9)..(159.9,-176)
				..controls(159.2,-176.1)and(154.1,-176.2)..(151.5,-176.2)
				..controls(148.1,-176.1)and(122.6,-176.1)..(121.6,-176.2)
				..controls(120.8,-176.2)and(119.8,-176.2)..(119.4,-176.1)--cycle
				(130.4,-160)..controls(129,-160)and(119.8,-160)..(110.1,-159.9)
				..controls(100.3,-159.8)and(88.5,-159.7)..(83.9,-159.6)
				..controls(77.8,-159.5)and(75.3,-159.5)..(74.5,-159.6)
				..controls(73.9,-159.7)and(72.5,-159.8)..(71.3,-159.8)
				..controls(69.6,-159.8)and(69,-159.8)..(68.5,-159.6)
				..controls(67.6,-159.3)and(66.8,-158.6)..(66.2,-157.6)
				..controls(66,-157.3)and(65.6,-156.5)..(65.3,-156.1)
				..controls(64.6,-155.1)and(64.5,-154.9)..(63.7,-152.9)
				..controls(63.2,-151.6)and(62.9,-150.9)..(61.8,-149.6)
				..controls(61.5,-149.2)and(61.1,-148.5)..(60.9,-147.9)
				..controls(60.7,-147.4)and(60.3,-146.6)..(60.1,-146.1)
				..controls(59.9,-145.5)and(59.5,-144.9)..(59,-144.3)
				..controls(58.5,-143.7)and(58.2,-143.2)..(57.9,-142.5)
				..controls(57.7,-141.9)and(57.3,-141.1)..(57,-140.7)
				..controls(56.6,-140.3)and(56.1,-139.3)..(55.8,-138.5)
				..controls(55.5,-137.8)and(55.1,-136.8)..(54.9,-136.5)
				..controls(54.7,-136.1)and(54.3,-135.2)..(53.9,-134.5)
				..controls(53.6,-133.8)and(53,-132.7)..(52.6,-132)
				..controls(51.6,-130.5)and(51.4,-130.2)..(51.1,-129)
				..controls(50.6,-127.5)and(50.6,-127.4)..(49.7,-126.3)
				..controls(49.2,-125.7)and(48.6,-124.8)..(48.3,-124.2)
				..controls(48,-123.6)and(47.5,-122.7)..(47.3,-122.2)
				..controls(47,-121.8)and(46.5,-120.7)..(46.1,-119.8)
				..controls(45.8,-118.9)and(45.2,-117.8)..(44.9,-117.2)
				..controls(44.2,-116.1)and(43.3,-114.4)..(42.3,-112.2)
				..controls(41.8,-111.3)and(41.3,-110.4)..(41.1,-110.1)
				..controls(40.5,-109.4)and(40.3,-109)..(40,-107.8)
				..controls(39.8,-107)and(39.6,-106.5)..(39.2,-106)
				..controls(38.9,-105.6)and(38.4,-104.7)..(38,-104)
				..controls(37.7,-103.4)and(37.2,-102.5)..(37,-102.1)
				..controls(36.7,-101.7)and(36.3,-100.8)..(36,-100.1)
				..controls(35.6,-98.9)and(35.3,-98.4)..(34.6,-97.6)
				..controls(34.2,-97.1)and(34.2,-96.3)..(34.6,-96)
				..controls(35,-95.7)and(35.5,-95.9)..(36.2,-96.7)
				..controls(36.5,-97)and(37.2,-97.5)..(37.6,-97.7)..controls(38,-98)and(38.6,-98.4)..(39.1,-98.9)
				..controls(39.6,-99.4)and(40.3,-100)..(40.6,-100.3)
				..controls(41.3,-100.8)and(42.4,-101.8)..(43.7,-102.9)
				..controls(44,-103.2)and(44.6,-103.7)..(45,-103.9)
				..controls(45.4,-104.2)and(45.9,-104.6)..(46.1,-104.8)
				..controls(47.1,-105.9)and(48.8,-107.1)..(50.5,-107.9)
				..controls(51.8,-108.6)and(52.3,-109)..(53.4,-110)
				..controls(54.8,-111.2)and(55.6,-111.6)..(57,-112.1)
				..controls(57.5,-112.3)and(58.5,-112.8)..(59.4,-113.3)
				..controls(60.2,-113.7)and(61.3,-114.3)..(61.9,-114.5)
				..controls(62.4,-114.7)and(63.3,-115.1)..(63.8,-115.3)
				..controls(64.4,-115.6)and(65.4,-115.9)..(66.2,-116.1)
				..controls(67,-116.3)and(68.4,-116.6)..(69.2,-116.9)
				..controls(72.3,-117.9)and(72.8,-118)..(78,-118.1)
				..controls(84,-118.3)and(86.1,-118.1)..(88.1,-117.1)
				..controls(89.4,-116.5)and(93,-115)..(93.7,-114.9)
				..controls(94.8,-114.6)and(95.4,-114.2)..(95.9,-113.4)
				..controls(96.2,-113)and(96.7,-112.5)..(97.1,-112.2)
				..controls(98,-111.6)and(101.3,-108.4)..(102.1,-107.2)
				..controls(102.4,-106.8)and(102.9,-106.2)..(103.2,-105.8)
				..controls(103.5,-105.5)and(104,-104.8)..(104.2,-104.2)
				..controls(104.5,-103.6)and(105,-102.8)..(105.3,-102.4)
				..controls(105.6,-102)and(106.1,-101.2)..(106.5,-100.6)
				..controls(107.1,-99.5)and(108.6,-97.5)..(109.5,-96.5)
				..controls(109.9,-96.1)and(110.1,-96)..(110.4,-96)
				..controls(111.1,-96)and(111.6,-96.5)..(111.8,-97.6)
				..controls(112,-98.6)and(112.3,-99.3)..(112.8,-99.9)
				..controls(113.7,-100.9)and(115.2,-103.1)..(115.7,-104.2)
				..controls(116.1,-105)and(116.5,-105.5)..(116.9,-106)
				..controls(117.3,-106.4)and(117.8,-106.9)..(118,-107.3)
				..controls(118.2,-107.6)and(118.6,-108)..(119,-108.3)
				..controls(119.6,-108.8)and(121.1,-110.3)..(121.8,-111.2)
				..controls(122.1,-111.7)and(122.5,-112)..(123.4,-112.4)
				..controls(124.8,-113.2)and(125.8,-113.8)..(127,-114.6)
				..controls(127.5,-114.9)and(128.1,-115.3)..(128.2,-115.4)
				..controls(128.4,-115.5)and(129.3,-115.8)..(130.2,-116)
				..controls(131.1,-116.2)and(132.5,-116.6)..(133.3,-116.9)
				..controls(135.7,-117.9)and(136.1,-117.9)..(142.5,-117.8)
				..controls(145.8,-117.8)and(148.4,-117.7)..(148.8,-117.6)
				..controls(149.2,-117.6)and(150.2,-117.2)..(151,-116.8)
				..controls(152,-116.4)and(152.8,-116)..(153.6,-115.9)
				..controls(154.5,-115.7)and(155.2,-115.4)..(156.3,-114.9)
				..controls(157.1,-114.5)and(158.2,-114)..(158.7,-113.8)
				..controls(159.1,-113.7)and(159.8,-113.3)..(160.2,-113)
				..controls(160.6,-112.7)and(161.3,-112.4)..(161.7,-112.2)
				..controls(162.7,-111.9)and(163.2,-111.7)..(164.1,-111.1)
				..controls(164.5,-110.8)and(165.3,-110.4)..(165.9,-110)
				..controls(166.5,-109.7)and(167.3,-109.2)..(167.6,-108.9)
				..controls(167.9,-108.6)and(168.6,-108.2)..(169,-108)
				..controls(169.4,-107.8)and(170.2,-107.3)..(170.7,-106.9)
				..controls(171.2,-106.5)and(172.1,-105.8)..(172.7,-105.4)
				..controls(173.2,-105)and(174.1,-104.4)..(174.5,-104)
				..controls(174.9,-103.7)and(175.7,-103.1)..(176.2,-102.6)
				..controls(176.7,-102.2)and(177.4,-101.5)..(177.9,-101.2)
				..controls(178.3,-100.8)and(178.8,-100.4)..(179,-100.2)
				..controls(179.2,-100)and(179.6,-99.6)..(180,-99.3)
				..controls(180.4,-99.1)and(180.9,-98.6)..(181.2,-98.4)
				..controls(181.5,-98.1)and(182.1,-97.7)..(182.5,-97.5)
				..controls(183.2,-97.2)and(183.5,-96.9)..(184.3,-95.9)
				..controls(184.6,-95.6)and(184.6,-95.6)..(184.8,-95.7)
				..controls(185,-95.8)and(185.1,-96)..(185.1,-96.2)
				..controls(185.1,-96.5)and(185,-96.7)..(184.6,-97.2)
				..controls(184.1,-97.8)and(184,-97.9)..(183.9,-98.8)
				..controls(183.7,-100.3)and(183.3,-101.3)..(182.9,-101.8)
				..controls(182.6,-102.1)and(182.3,-102.6)..(182.2,-102.9)
				..controls(182,-103.3)and(181.6,-104)..(181.2,-104.6)
				..controls(180.9,-105.2)and(180.5,-106.1)..(180.4,-106.6)
				..controls(179.9,-108)and(179.7,-108.5)..(179,-109.4)
				..controls(178.6,-109.8)and(178.2,-110.5)..(177.9,-111)
				..controls(177.7,-111.5)and(177.3,-112.3)..(177,-112.7)
				..controls(176.6,-113.2)and(176.4,-113.7)..(176.3,-114.2)
				..controls(176.2,-114.7)and(175.9,-115.2)..(175.4,-116)
				..controls(175,-116.6)and(174.4,-117.7)..(174,-118.5)
				..controls(173.7,-119.2)and(173.3,-120.1)..(173,-120.5)
				..controls(172.8,-120.9)and(172.4,-121.8)..(172.1,-122.5)
				..controls(171.8,-123.2)and(171.3,-124.1)..(171,-124.6)
				..controls(170.8,-125.1)and(170.3,-126)..(169.9,-126.6)
				..controls(169.6,-127.3)and(169.2,-128.1)..(169,-128.3)
				..controls(168.8,-128.6)and(168.5,-129.3)..(168.3,-129.8)
				..controls(168.1,-130.3)and(167.8,-131.1)..(167.5,-131.6)
				..controls(166.9,-132.5)and(164.2,-137.9)..(164.2,-138.1)
				..controls(164.2,-138.1)and(163.9,-138.8)..(163.4,-139.7)
				..controls(162.9,-140.5)and(162.4,-141.5)..(162.3,-141.9)
				..controls(161.9,-143)and(161.6,-143.4)..(160.9,-144.1)
				..controls(160.4,-144.6)and(160.3,-144.8)..(160.1,-145.4)
				..controls(160,-145.8)and(159.6,-146.8)..(159.2,-147.8)
				..controls(158.7,-148.8)and(158.1,-150)..(157.9,-150.6)
				..controls(157.6,-151.1)and(157.2,-151.9)..(156.9,-152.4)
				..controls(156.5,-152.8)and(156.1,-153.7)..(155.8,-154.4)
				..controls(155.6,-155)and(155.2,-155.9)..(154.9,-156.3)
				..controls(154.7,-156.7)and(154.4,-157.4)..(154.2,-157.9)
				..controls(154.1,-158.4)and(153.8,-159)..(153.5,-159.3)
				..controls(152.8,-160.1)and(153,-160.1)..(142.4,-160.1)
				..controls(137.2,-160.1)and(131.8,-160.1)..(130.4,-160)--cycle
				(141.5,-103.7)..controls(140.8,-103)and(139.8,-102.1)..(139.2,-101.6)
				..controls(138.5,-101.2)and(137.9,-100.6)..(137.8,-100.4)
				..controls(137,-99.5)and(134,-96.5)..(133.2,-96)
				..controls(132.7,-95.7)and(131.8,-94.9)..(131.1,-94.3)
				..controls(129.4,-92.7)and(127.9,-91.6)..(126.6,-91.1)
				..controls(125.2,-90.5)and(124.1,-90)..(122.9,-89.1)
				..controls(122.1,-88.6)and(121.3,-88.2)..(119.9,-87.7)
				..controls(118.8,-87.3)and(117.5,-86.8)..(116.9,-86.5)
				..controls(114.3,-85.5)and(113.7,-85.4)..(109.7,-85.5)
				..controls(105.3,-85.7)and(103.9,-85.9)..(101.6,-87)
				..controls(100.7,-87.4)and(99.4,-87.8)..(97.9,-88)
				..controls(97,-88.2)and(96.6,-88.3)..(95.8,-88.9)
				..controls(95.4,-89.2)and(94.7,-89.6)..(94.2,-89.8)
				..controls(92.9,-90.3)and(92.3,-90.7)..(90.9,-92.1)
				..controls(90.2,-92.7)and(89.4,-93.3)..(89,-93.5)..controls(87,-94.5)and(85.2,-96)..(83.6,-97.9)
				..controls(83.3,-98.2)and(82.5,-98.9)..(81.8,-99.4)
				..controls(80.7,-100.1)and(80.5,-100.3)..(80,-101.1)
				..controls(78.7,-103.1)and(77.2,-103.9)..(76.4,-103)
				..controls(75.7,-102.2)and(75.7,-103.2)..(75.7,-83.6)
				..controls(75.6,-73.6)and(75.6,-64.8)..(75.5,-64)
				..controls(75.5,-63.2)and(75.5,-61.4)..(75.6,-59.8)--(75.7,-57.1)--(76,-56.8)
				..controls(76.8,-56.2)and(78,-56.3)..(79.8,-57.2)
				..controls(80.4,-57.5)and(81.3,-57.8)..(81.8,-57.9)
				..controls(82.3,-58.1)and(83,-58.4)..(83.6,-58.8)..controls(84.3,-59.2)and(85,-59.5)..(86,-59.8)
				..controls(86.8,-60.1)and(88,-60.6)..(88.7,-60.9)
				..controls(89.3,-61.2)and(90.6,-61.7)..(91.4,-62)
				..controls(92.2,-62.3)and(93.7,-62.9)..(94.6,-63.3)
				..controls(95.6,-63.7)and(96.8,-64.1)..(97.4,-64.3)
				..controls(98,-64.5)and(99,-64.9)..(99.7,-65.2)
				..controls(100.4,-65.6)and(101.5,-66)..(102.3,-66.2)
				..controls(104.2,-66.8)and(105,-67.1)..(106.2,-67.8)
				..controls(107.9,-68.7)and(110.2,-69.4)..(111.3,-69.2)
				..controls(112.5,-69)and(115,-68.2)..(115.6,-67.8)
				..controls(117.3,-66.7)and(117.6,-66.6)..(119.5,-66.2)
				..controls(121.9,-65.8)and(122.1,-65.7)..(123,-65.1)
				..controls(124.1,-64.4)and(124.4,-64.3)..(126,-63.9)
				..controls(127,-63.7)and(127.7,-63.4)..(128.4,-63.1)
				..controls(129,-62.8)and(129.7,-62.5)..(130.6,-62.4)
				..controls(132.1,-62.1)and(134,-61.5)..(135.2,-61)
				..controls(135.7,-60.8)and(136.7,-60.5)..(137.3,-60.3)
				..controls(138.4,-60)and(139.3,-59.7)..(141.1,-59)
				..controls(141.6,-58.7)and(142,-58.6)..(142.5,-58.6)
				..controls(143.2,-58.6)and(143.3,-58.7)..(143.6,-59)
				..controls(144,-59.5)and(144.1,-59.9)..(144,-60.8)
				..controls(144,-61.2)and(144,-71.1)..(144,-82.8)--(143.9,-104.1)--(143.7,-104.5)
				..controls(143.5,-104.7)and(143.3,-104.8)..(143,-104.9)
				..controls(142.6,-104.9)and(142.5,-104.8)..(141.5,-103.7)--cycle
				(27.8,-91.4)..controls(27.1,-90.9)and(26.3,-90)..(26.1,-89.5)
				..controls(26,-89.2)and(25.8,-88.6)..(25.7,-87.9)
				..controls(25.4,-86.2)and(25.6,-84.8)..(26.5,-83.5)
				..controls(27.4,-82.3)and(28.3,-82)..(29.4,-82.3)
				..controls(31.1,-82.8)and(32.1,-84.9)..(32,-87.4)
				..controls(31.9,-89.3)and(31.2,-90.6)..(29.9,-91.4)
				..controls(29.1,-91.9)and(28.6,-91.9)..(27.8,-91.4)--cycle
				(190.1,-91.3)..controls(189.7,-91.2)and(189.3,-90.9)..(189.1,-90.7)
				..controls(188.3,-89.8)and(187.6,-87.7)..(187.6,-86.2)
				..controls(187.6,-85.5)and(187.7,-85.2)..(188.2,-84.4)
				..controls(189.1,-82.6)and(190.1,-81.8)..(191.3,-82.1)
				..controls(192,-82.3)and(193.5,-83.7)..(193.8,-84.4)
				..controls(194.1,-85.1)and(194.2,-87.1)..(194,-88.3)
				..controls(193.8,-89.8)and(193.3,-90.5)..(192.3,-91.3)
				..controls(191.7,-91.7)and(190.9,-91.7)..(190.1,-91.3)--cycle
				(173.7,-83.4)..controls(172.2,-82.5)and(172,-82.5)..(170.6,-82.4)
				..controls(169.7,-82.4)and(168.8,-82.4)..(168.3,-82.5)
				..controls(167.2,-82.7)and(166.9,-82.7)..(166.6,-82.2)
				..controls(166.3,-81.8)and(166.3,-81)..(166.8,-80)
				..controls(167.2,-78.9)and(167.4,-77.4)..(167.7,-74.1)
				..controls(167.8,-72.8)and(167.9,-71.8)..(167.9,-71.7)
				..controls(168,-71.7)and(168,-70.2)..(168,-68.4)..controls(168,-65.4)and(168,-64.9)..(167.7,-63)
				..controls(167.3,-60.5)and(167.1,-59.9)..(166.4,-58.4)
				..controls(166.1,-57.7)and(165.8,-56.9)..(165.7,-56.6)
				..controls(165.7,-56.2)and(165.4,-55.6)..(165.2,-55.2)
				..controls(164.6,-53.8)and(161.9,-50.7)..(161,-50.3)
				..controls(159.8,-49.7)and(157.9,-50)..(155.2,-51.2)
				..controls(154.7,-51.4)and(153.4,-51.8)..(152.2,-52.1)
				..controls(150.8,-52.4)and(149.4,-52.9)..(148.6,-53.2)
				..controls(147.8,-53.6)and(146.7,-54)..(146,-54.2)
				..controls(145.4,-54.4)and(144.5,-54.7)..(144,-54.9)
				..controls(143.6,-55.2)and(142.7,-55.5)..(141.8,-55.7)
				..controls(141,-55.9)and(139.6,-56.3)..(138.7,-56.6)
				..controls(137.8,-57)and(136.3,-57.4)..(135.3,-57.6)
				..controls(133.8,-57.9)and(133.2,-58.1)..(132.1,-58.7)
				..controls(131.1,-59.2)and(130.3,-59.4)..(129.4,-59.6)
				..controls(128.7,-59.8)and(127.8,-60)..(127.4,-60.1)
				..controls(127,-60.3)and(126,-60.5)..(125.1,-60.7)
				..controls(123.2,-61.1)and(121.7,-61.5)..(120.2,-62.1)
				..controls(119.7,-62.3)and(118.1,-62.8)..(116.7,-63.2)
				..controls(114.6,-63.9)and(114,-64.1)..(113,-64.7)
				..controls(111.9,-65.4)and(111.8,-65.4)..(111,-65.5)
				..controls(109.9,-65.5)and(109.3,-65.4)..(107.7,-64.7)
				..controls(107,-64.4)and(106,-64)..(105.3,-63.9)
				..controls(103.5,-63.5)and(102.2,-63)..(101.2,-62.5)
				..controls(100.1,-61.8)and(99.3,-61.5)..(97.8,-61.1)
				..controls(97.2,-61)and(96.2,-60.7)..(95.7,-60.4)
				..controls(95.1,-60.1)and(94.2,-59.8)..(93.5,-59.6)
				..controls(92.9,-59.4)and(91.7,-59)..(91,-58.7)..controls(90.2,-58.4)and(88.9,-58)..(88,-57.7)
				..controls(87.2,-57.5)and(86.1,-57.1)..(85.5,-56.8)
				..controls(85,-56.6)and(84.1,-56.2)..(83.5,-56)..controls(82.9,-55.7)and(81.8,-55.2)..(81,-54.9)
				..controls(80.1,-54.5)and(79.1,-54.2)..(78.7,-54.1)
				..controls(77.9,-53.8)and(77.3,-53.6)..(76.2,-52.9)
				..controls(75.7,-52.6)and(75.2,-52.4)..(74.1,-52.2)
				..controls(72.5,-51.9)and(72.4,-51.8)..(71.3,-51.1)
				..controls(69.9,-50.1)and(69.4,-49.9)..(67.9,-49.8)
				..controls(67,-49.7)and(66.6,-49.6)..(66.3,-49.4)
				..controls(66.1,-49.3)and(65.7,-49)..(65.4,-48.8)
				..controls(65.1,-48.7)and(64.2,-48.4)..(63.6,-48.2)
				..controls(62.9,-48)and(61.7,-47.6)..(60.9,-47.2)
				..controls(60.1,-46.8)and(58.9,-46.4)..(58.2,-46.2)
				..controls(57.6,-46)and(56.8,-45.7)..(56.5,-45.5)
				..controls(54.9,-44.6)and(54.6,-44.5)..(53.2,-44.2)
				..controls(51.7,-43.8)and(51,-43.6)..(49.9,-42.9)
				..controls(49.5,-42.6)and(48.9,-42.4)..(47.5,-42)
				..controls(45.7,-41.6)and(45.6,-41.5)..(44.3,-40.6)
				..controls(43.5,-40.1)and(42.8,-39.6)..(42.7,-39.4)
				..controls(42.4,-38.7)and(43,-38.1)..(44.5,-37.6)..controls(44.9,-37.5)and(46,-37)..(46.9,-36.6)
				..controls(47.9,-36.1)and(48.8,-35.8)..(49.6,-35.6)
				..controls(50.5,-35.4)and(51.3,-35.2)..(52.4,-34.6)
				..controls(53.3,-34.1)and(54.5,-33.7)..(54.9,-33.5)
				..controls(55.4,-33.4)and(56.4,-33)..(57.2,-32.6)
				..controls(58.6,-32)and(60.1,-31.5)..(61.6,-31.2)
				..controls(62.3,-31.1)and(62.7,-31)..(63.2,-30.7)
				..controls(64.1,-30)and(64.3,-29.9)..(65.3,-29.7)
				..controls(65.7,-29.6)and(66.5,-29.2)..(67.1,-28.9)
				..controls(67.6,-28.6)and(68.8,-28.2)..(69.6,-27.9)
				..controls(70.5,-27.7)and(71.6,-27.3)..(72.1,-27)
				..controls(72.6,-26.7)and(73.8,-26.2)..(74.8,-25.9)
				..controls(75.7,-25.5)and(76.7,-25.1)..(77.1,-25)
				..controls(77.4,-24.8)and(78.5,-24.4)..(79.6,-24.1)
				..controls(80.6,-23.8)and(82,-23.2)..(82.8,-22.9)
				..controls(83.5,-22.5)and(84.8,-22)..(85.8,-21.8)
				..controls(86.7,-21.5)and(87.8,-21.1)..(88.1,-20.9)
				..controls(88.4,-20.7)and(89.3,-20.3)..(90.1,-20.1)
				..controls(90.8,-19.8)and(91.9,-19.4)..(92.4,-19.1)
				..controls(93.5,-18.5)and(94.4,-18.2)..(95.4,-18.1)
				..controls(95.7,-18)and(96.4,-17.8)..(96.9,-17.5)
				..controls(97.5,-17.3)and(98.4,-16.9)..(99,-16.6)
				..controls(99.7,-16.3)and(100.9,-15.9)..(101.9,-15.6)
				..controls(102.9,-15.3)and(104.4,-14.7)..(105.3,-14.3)
				..controls(108.1,-13.2)and(108.5,-13.1)..(110.1,-13.1)
				..controls(111.4,-13.1)and(111.5,-13.1)..(112.4,-13.5)
				..controls(113.2,-14)and(115,-14.8)..(116.6,-15.4)
				..controls(116.9,-15.5)and(117.9,-15.8)..(118.8,-16.1)
				..controls(120.9,-16.6)and(121.5,-16.9)..(122.6,-17.9)
				..controls(123,-18.3)and(123.7,-18.8)..(123.9,-18.9)
				..controls(124.6,-19.4)and(126,-19.8)..(127.9,-20.2)
				..controls(129.6,-20.5)and(129.8,-20.6)..(131.1,-21.4)
				..controls(131.5,-21.7)and(132.4,-22.1)..(133.1,-22.3)
				..controls(136.1,-23.2)and(138.3,-24.1)..(140.2,-25)
				..controls(141,-25.5)and(141.5,-25.6)..(142.5,-25.8)
				..controls(145.1,-26.3)and(146,-26.6)..(149.4,-28.1)
				..controls(151.3,-28.9)and(153.5,-29.8)..(154.3,-30.1)
				..controls(155.1,-30.4)and(156.2,-30.8)..(156.6,-31)
				..controls(157.1,-31.2)and(158.2,-31.6)..(159.1,-31.8)
				..controls(160.6,-32.2)and(160.9,-32.3)..(161.7,-32.9)
				..controls(162.4,-33.3)and(162.9,-33.6)..(163.7,-33.8)
				..controls(164.3,-34)and(164.8,-34.1)..(164.9,-34.1)
				..controls(165,-34.1)and(165.6,-34.4)..(166.4,-34.8)
				..controls(167.3,-35.2)and(168.1,-35.4)..(169,-35.6)
				..controls(170,-35.8)and(170.5,-36)..(171.7,-36.6)
				..controls(172.4,-37)and(173.3,-37.4)..(173.6,-37.4)
				..controls(175.6,-37.7)and(176.2,-38)..(176.7,-39.1)
				..controls(176.9,-39.6)and(176.9,-39.7)..(176.7,-40.4)
				..controls(176.4,-41.6)and(175.8,-42.7)..(175.2,-43.4)
				..controls(174.5,-44.1)and(172.6,-45.5)..(171.7,-45.8)
				..controls(170.6,-46.2)and(169.1,-46)..(168.2,-45.4)
				..controls(167.9,-45.2)and(167.3,-44.8)..(166.9,-44.4)
				..controls(166.4,-44.1)and(165.8,-43.6)..(165.4,-43.2)
				..controls(165.1,-42.9)and(164.4,-42.3)..(163.9,-41.9)
				..controls(163.3,-41.5)and(162.7,-40.9)..(162.4,-40.5)
				..controls(161.4,-39.2)and(159.9,-38)..(158.8,-37.8)
				..controls(158.4,-37.8)and(158.3,-37.7)..(158.2,-37.2)
				..controls(158,-36.8)and(157.4,-36.1)..(157.1,-36.1)
				..controls(157,-36.1)and(156.7,-36.1)..(156.5,-36)
				..controls(156.1,-35.9)and(155.9,-35.7)..(155.7,-35.2)
				..controls(155.3,-34.6)and(155.1,-34.4)..(154,-34.1)
				..controls(153.6,-34)and(153.2,-33.7)..(152.9,-33.5)
				..controls(152.4,-32.9)and(151.6,-32.4)..(151.3,-32.4)
				..controls(150.9,-32.4)and(150.4,-32.6)..(150.4,-32.9)
				..controls(150.4,-33.2)and(150.8,-33.4)..(151.4,-33.6)
				..controls(152.1,-33.7)and(152.8,-34.1)..(153.5,-34.9)
				..controls(154.1,-35.5)and(154.3,-35.6)..(155,-35.8)
				..controls(155.7,-36)and(155.7,-36)..(156,-36.6)
				..controls(156.3,-37.3)and(156.7,-37.7)..(157.5,-37.8)
				..controls(158.1,-37.9)and(158.2,-38)..(158.4,-38.8)
				..controls(158.5,-39.4)and(158.6,-39.6)..(159.6,-40.4)
				..controls(160.2,-41)and(160.9,-41.6)..(161.3,-41.8)
				..controls(161.7,-42)and(161.9,-42.3)..(162.2,-42.8)
				..controls(162.4,-43.2)and(162.8,-43.8)..(163.1,-44.1)
				..controls(163.9,-44.8)and(165.1,-46.7)..(165.7,-48)
				..controls(166.2,-49.3)and(166.4,-49.7)..(167.2,-50.5)
				..controls(167.5,-50.8)and(167.9,-51.4)..(168.1,-51.8)
				..controls(168.3,-52.2)and(168.8,-53)..(169.1,-53.6)
				..controls(169.6,-54.3)and(169.8,-54.9)..(169.9,-55.6)
				..controls(170.4,-57.8)and(170.5,-58.1)..(171.2,-59.5)
				..controls(171.6,-60.3)and(172,-61.5)..(172.2,-62.1)
				..controls(172.3,-62.8)and(172.7,-63.7)..(172.9,-64.1)
				..controls(173.2,-64.6)and(173.4,-65.3)..(173.4,-65.6)
				..controls(173.8,-68.8)and(174.2,-70.6)..(174.9,-72.6)
				..controls(175.5,-74.5)and(175.9,-77.1)..(176.1,-80.2)
				..controls(176.3,-83.4)and(175.5,-84.4)..(173.7,-83.4)--cycle;
				\definecolor{ezs}{RGB}{247,179,26}\fill[ezs](-5.7,-156.6)--(-3.4,-157.5)--(61.9,12.7)--(59.7,13.6)--cycle;
				\definecolor{ezs}{RGB}{247,179,22}\fill[ezs](-3.8,-157.3)--(-1.5,-158.2)--(63.9,12)--(61.6,12.9)--cycle;
				\definecolor{ezs}{RGB}{248,179,19}\fill[ezs](-1.8,-158.1)--(0.4,-158.9)--(65.8,11.2)--(63.5,12.1)--cycle;
				\definecolor{ezs}{RGB}{248,178,16}\fill[ezs](0.1,-158.8)--(2.4,-159.7)--(67.7,10.5)--(65.4,11.4)--cycle;
				\definecolor{ezs}{RGB}{249,178,15}\fill[ezs](2,-159.5)--(4.3,-160.4)--(69.6,9.8)--(67.3,10.7)--cycle;
				\definecolor{ezs}{RGB}{251,177,17}\fill[ezs](3.9,-160.3)--(6.2,-161.1)--(71.5,9)--(69.2,9.9)--cycle;
				\definecolor{ezs}{RGB}{252,179,18}\fill[ezs](5.8,-161)--(8.1,-161.9)--(73.4,8.3)--(71.1,9.2)--cycle;
				\definecolor{ezs}{RGB}{252,182,17}\fill[ezs](7.7,-161.7)--(10,-162.6)--(75.3,7.6)--(73,8.5)--cycle;
				\definecolor{ezs}{RGB}{252,185,17}\fill[ezs](9.6,-162.5)--(11.9,-163.4)--(77.2,6.8)--(75,7.7)--cycle;
				\definecolor{ezs}{RGB}{251,189,18}\fill[ezs](11.5,-163.2)--(13.8,-164.1)--(79.2,6.1)--(76.9,7)--cycle;
				\definecolor{ezs}{RGB}{251,191,22}\fill[ezs](13.4,-163.9)--(15.7,-164.8)--(81.1,5.4)--(78.8,6.2)--cycle;
				\definecolor{ezs}{RGB}{252,193,25}\fill[ezs](15.4,-164.7)--(17.6,-165.6)--(83,4.6)--(80.7,5.5)--cycle;
				\definecolor{ezs}{RGB}{253,195,27}\fill[ezs](17.3,-165.4)--(19.6,-166.3)--(84.9,3.9)--(82.6,4.8)--cycle;
				\definecolor{ezs}{RGB}{253,196,31}\fill[ezs](19.2,-166.1)--(21.5,-167)--(86.8,3.2)--(84.5,4)--cycle;
				\definecolor{ezs}{RGB}{254,199,35}\fill[ezs](21.1,-166.9)--(23.4,-167.8)--(88.7,2.4)--(86.4,3.3)--cycle;
				\definecolor{ezs}{RGB}{253,201,38}\fill[ezs](23,-167.6)--(25.3,-168.5)--(90.6,1.7)--(88.3,2.6)--cycle;
				\definecolor{ezs}{RGB}{253,203,40}\fill[ezs](24.9,-168.3)--(27.2,-169.2)--(92.5,1)--(90.3,1.8)--cycle;
				\definecolor{ezs}{RGB}{253,205,42}\fill[ezs](26.8,-169.1)--(29.1,-170)--(94.4,0.2)--(92.2,1.1)--cycle;
				\definecolor{ezs}{RGB}{253,207,45}\fill[ezs](28.7,-169.8)--(31,-170.7)--(96.4,-0.5)--(94.1,0.4)--cycle;
				\definecolor{ezs}{RGB}{253,210,49}\fill[ezs](30.7,-170.5)--(32.9,-171.4)--(98.3,-1.2)--(96,-0.4)--cycle;
				\definecolor{ezs}{RGB}{253,212,54}\fill[ezs](32.6,-171.3)--(34.9,-172.2)--(100.2,-2)--(97.9,-1.1)--cycle;
				\definecolor{ezs}{RGB}{253,214,59}\fill[ezs](34.5,-172)--(36.8,-172.9)--(102.1,-2.7)--(99.8,-1.8)--cycle;
				\definecolor{ezs}{RGB}{253,216,63}\fill[ezs](36.4,-172.7)--(38.7,-173.6)--(104,-3.4)--(101.7,-2.6)--cycle;
				\definecolor{ezs}{RGB}{253,218,65}\fill[ezs](38.3,-173.5)--(40.6,-174.4)--(105.9,-4.2)--(103.6,-3.3)--cycle;
				\definecolor{ezs}{RGB}{254,220,66}\fill[ezs](40.2,-174.2)--(42.5,-175.1)--(107.8,-4.9)--(105.5,-4)--cycle;
				\definecolor{ezs}{RGB}{254,222,68}\fill[ezs](42.1,-175)--(44.4,-175.8)--(109.7,-5.6)--(107.5,-4.8)--cycle;
				\definecolor{ezs}{RGB}{254,224,71}\fill[ezs](44,-175.7)--(46.3,-176.6)--(111.7,-6.4)--(109.4,-5.5)--cycle;
				\definecolor{ezs}{RGB}{254,226,75}\fill[ezs](46,-176.4)--(48.2,-177.3)--(113.6,-7.1)--(111.3,-6.2)--cycle;
				\definecolor{ezs}{RGB}{254,228,79}\fill[ezs](47.9,-177.2)--(50.1,-178)--(115.5,-7.8)--(113.2,-7)--cycle;
				\definecolor{ezs}{RGB}{254,230,82}\fill[ezs](49.8,-177.9)--(52.1,-178.8)--(117.4,-8.6)--(115.1,-7.7)--cycle;
				\definecolor{ezs}{RGB}{254,232,85}\fill[ezs](51.7,-178.6)--(54,-179.5)--(119.3,-9.3)--(117,-8.4)--cycle;
				\definecolor{ezs}{RGB}{254,234,88}\fill[ezs](53.6,-179.4)--(55.9,-180.2)--(121.2,-10)--(118.9,-9.2)--cycle;
				\definecolor{ezs}{RGB}{254,237,91}\fill[ezs](55.5,-180.1)--(57.8,-181)--(123.1,-10.8)--(120.8,-9.9)--cycle;
				\definecolor{ezs}{RGB}{254,239,94}\fill[ezs](57.4,-180.8)--(59.7,-181.7)--(125,-11.5)--(122.8,-10.6)--cycle;
				\definecolor{ezs}{RGB}{254,241,97}\fill[ezs](59.3,-181.6)--(61.6,-182.4)--(127,-12.2)--(124.7,-11.4)--cycle;
				\definecolor{ezs}{RGB}{254,243,101}\fill[ezs](61.2,-182.3)--(63.5,-183.2)--(128.9,-13)--(126.6,-12.1)--cycle;
				\definecolor{ezs}{RGB}{254,245,103}\fill[ezs](63.2,-183)--(65.4,-183.9)--(130.8,-13.7)--(128.5,-12.8)--cycle;
				\definecolor{ezs}{RGB}{254,245,105}\fill[ezs](65.1,-183.8)--(67.4,-184.6)--(132.7,-14.4)--(130.4,-13.6)--cycle;
				\definecolor{ezs}{RGB}{254,245,105}\fill[ezs](67,-184.5)--(69.3,-185.4)--(134.6,-15.2)--(132.3,-14.3)--cycle;
				\definecolor{ezs}{RGB}{254,244,103}\fill[ezs](68.9,-185.2)--(71.2,-186.1)--(136.5,-15.9)--(134.2,-15)--cycle;
				\definecolor{ezs}{RGB}{253,243,101}\fill[ezs](70.8,-186)--(73.1,-186.8)--(138.4,-16.6)--(136.1,-15.8)--cycle;
				\definecolor{ezs}{RGB}{253,241,98}\fill[ezs](72.7,-186.7)--(75,-187.6)--(140.3,-17.4)--(138,-16.5)--cycle;
				\definecolor{ezs}{RGB}{253,239,95}\fill[ezs](74.6,-187.4)--(76.9,-188.3)--(142.2,-18.1)--(140,-17.2)--cycle;
				\definecolor{ezs}{RGB}{253,236,92}\fill[ezs](76.5,-188.2)--(78.8,-189)--(144.2,-18.8)--(141.9,-18)--cycle;
				\definecolor{ezs}{RGB}{253,234,88}\fill[ezs](78.5,-188.9)--(80.7,-189.8)--(146.1,-19.6)--(143.8,-18.7)--cycle;
				\definecolor{ezs}{RGB}{253,232,85}\fill[ezs](80.4,-189.6)--(82.7,-190.5)--(148,-20.3)--(145.7,-19.4)--cycle;
				\definecolor{ezs}{RGB}{253,230,82}\fill[ezs](82.3,-190.4)--(84.6,-191.2)--(149.9,-21)--(147.6,-20.2)--cycle;
				\definecolor{ezs}{RGB}{253,228,78}\fill[ezs](84.2,-191.1)--(86.5,-192)--(151.8,-21.8)--(149.5,-20.9)--cycle;
				\definecolor{ezs}{RGB}{253,226,74}\fill[ezs](86.1,-191.8)--(88.4,-192.7)--(153.7,-22.5)--(151.4,-21.6)--cycle;
				\definecolor{ezs}{RGB}{253,224,70}\fill[ezs](88,-192.6)--(90.3,-193.4)--(155.6,-23.3)--(153.3,-22.4)--cycle;
				\definecolor{ezs}{RGB}{253,222,67}\fill[ezs](89.9,-193.3)--(92.2,-194.2)--(157.5,-24)--(155.3,-23.1)--cycle;
				\definecolor{ezs}{RGB}{253,220,65}\fill[ezs](91.8,-194)--(94.1,-194.9)--(159.5,-24.7)--(157.2,-23.8)--cycle;
				\definecolor{ezs}{RGB}{253,217,63}\fill[ezs](93.7,-194.8)--(96,-195.6)--(161.4,-25.5)--(159.1,-24.6)--cycle;
				\definecolor{ezs}{RGB}{253,215,60}\fill[ezs](95.7,-195.5)--(97.9,-196.4)--(163.3,-26.2)--(161,-25.3)--cycle;
				\definecolor{ezs}{RGB}{253,213,56}\fill[ezs](97.6,-196.2)--(99.9,-197.1)--(165.2,-26.9)--(162.9,-26)--cycle;
				\definecolor{ezs}{RGB}{252,211,53}\fill[ezs](99.5,-197)--(101.8,-197.8)--(167.1,-27.7)--(164.8,-26.8)--cycle;
				\definecolor{ezs}{RGB}{252,209,49}\fill[ezs](101.4,-197.7)--(103.7,-198.6)--(169,-28.4)--(166.7,-27.5)--cycle;
				\definecolor{ezs}{RGB}{252,207,45}\fill[ezs](103.3,-198.4)--(105.6,-199.3)--(170.9,-29.1)--(168.6,-28.2)--cycle;
				\definecolor{ezs}{RGB}{253,205,42}\fill[ezs](105.2,-199.2)--(107.5,-200)--(172.8,-29.9)--(170.6,-29)--cycle;
				\definecolor{ezs}{RGB}{253,202,40}\fill[ezs](107.1,-199.9)--(109.4,-200.8)--(174.7,-30.6)--(172.5,-29.7)--cycle;
				\definecolor{ezs}{RGB}{253,200,38}\fill[ezs](109,-200.6)--(111.3,-201.5)--(176.7,-31.3)--(174.4,-30.4)--cycle;
				\definecolor{ezs}{RGB}{253,198,35}\fill[ezs](111,-201.4)--(113.2,-202.2)--(178.6,-32.1)--(176.3,-31.2)--cycle;
				\definecolor{ezs}{RGB}{253,196,31}\fill[ezs](112.9,-202.1)--(115.2,-203)--(180.5,-32.8)--(178.2,-31.9)--cycle;
				\definecolor{ezs}{RGB}{253,194,28}\fill[ezs](114.8,-202.8)--(117.1,-203.7)--(182.4,-33.5)--(180.1,-32.6)--cycle;
				\definecolor{ezs}{RGB}{253,192,26}\fill[ezs](116.7,-203.6)--(119,-204.5)--(184.3,-34.3)--(182,-33.4)--cycle;
				\definecolor{ezs}{RGB}{253,189,24}\fill[ezs](118.6,-204.3)--(120.9,-205.2)--(186.2,-35)--(183.9,-34.1)--cycle;
				\definecolor{ezs}{RGB}{253,186,20}\fill[ezs](120.5,-205)--(122.8,-205.9)--(188.1,-35.7)--(185.8,-34.8)--cycle;
				\definecolor{ezs}{RGB}{253,184,17}\fill[ezs](122.4,-205.8)--(124.7,-206.7)--(190,-36.5)--(187.8,-35.6)--cycle;
				\definecolor{ezs}{RGB}{252,183,15}\fill[ezs](124.3,-206.5)--(126.6,-207.4)--(192,-37.2)--(189.7,-36.3)--cycle;
				\definecolor{ezs}{RGB}{252,181,14}\fill[ezs](126.3,-207.2)--(128.5,-208.1)--(193.9,-37.9)--(191.6,-37.1)--cycle;
				\definecolor{ezs}{RGB}{252,179,13}\fill[ezs](128.2,-208)--(130.4,-208.9)--(195.8,-38.7)--(193.5,-37.8)--cycle;
				\definecolor{ezs}{RGB}{252,178,12}\fill[ezs](130.1,-208.7)--(132.4,-209.6)--(197.7,-39.4)--(195.4,-38.5)--cycle;
				\definecolor{ezs}{RGB}{251,178,10}\fill[ezs](132,-209.4)--(134.3,-210.3)--(199.6,-40.1)--(197.3,-39.3)--cycle;
				\definecolor{ezs}{RGB}{251,178,10}\fill[ezs](133.9,-210.2)--(136.2,-211.1)--(201.5,-40.9)--(199.2,-40)--cycle;
				\definecolor{ezs}{RGB}{250,177,11}\fill[ezs](135.8,-210.9)--(138.1,-211.8)--(203.4,-41.6)--(201.1,-40.7)--cycle;
				\definecolor{ezs}{RGB}{250,178,11}\fill[ezs](137.7,-211.6)--(140,-212.5)--(205.3,-42.3)--(203.1,-41.5)--cycle;
				\definecolor{ezs}{RGB}{248,180,10}\fill[ezs](139.6,-212.4)--(141.9,-213.3)--(207.3,-43.1)--(205,-42.2)--cycle;
				\definecolor{ezs}{RGB}{247,181,9}\fill[ezs](141.5,-213.1)--(143.8,-214)--(209.2,-43.8)--(206.9,-42.9)--cycle;
				\definecolor{ezs}{RGB}{246,182,9}\fill[ezs](143.5,-213.8)--(145.7,-214.7)--(211.1,-44.5)--(208.8,-43.7)--cycle;
				\definecolor{ezs}{RGB}{246,182,10}\fill[ezs](145.4,-214.6)--(147.7,-215.5)--(213,-45.3)--(210.7,-44.4)--cycle;
			\end{scope}
			% chữ EDUZY (Times đậm, làm dày nét bằng 5 lớp lệch nhẹ; căn đúng khung chữ logo gốc)
			\foreach \dx/\dy in {-0.7/0,0.7/0,0/-0.3,0/0.3,0/0}
			\node[anchor=center,inner sep=0pt,outer sep=0pt,text=ezText,align=center] at (109.6+\dx,-193.6+\dy)
			{\resizebox{1.01cm}{0.22cm}{\fontencoding{T1}\fontfamily{ptm}\fontseries{b}\selectfont EDUZY}};
		\end{tikzpicture}%
	}%
}
%====================== THIẾT LẬP THÂN ĐỀ ======================
%--- Môi trường tự luận {bt}: ex_test 3.3.4 chưa khai báo sẵn nên tạo ở đây
%    (ex_test đã tự móc \fixshowans và \listenumerate cho bt)
\makeatletter
\@ifundefined{bt}{\newtheorem{bt}{\color{gold1}\namebt}}{}
\makeatother
\renewcommand{\sectionmark}[1]{\markright{\thesection\ #1}}
%--- Khung tiêu đề (logo + tối đa 4 dòng; dòng 3, 4 để trống {} thì bỏ qua)
%  Đề thi:    \khungtieude{THI THỬ THPT QUỐC GIA NĂM 2026}{Môn Toán -- Đề số 21}
%                         {(Đề thi gồm: ...)}{Thời gian làm bài: 90 phút (không kể thời gian phát đề).}
%  Chuyên đề: \khungtieude{CHUYÊN ĐỀ 3 -- DÃY SỐ VÀ CÔNG THỨC TRUY HỒI}
%                         {Bồi dưỡng học sinh giỏi Toán THPT}{...}{}
\newcommand{\khungtieude}[4]{%
	\begin{tcolorbox}[colframe=gold1,colback=gold5!45,boxrule=0.5mm]
		\begin{minipage}[c]{3cm}
			\centerline{\eduzylogo[2.4cm]}
		\end{minipage}\hspace{0.5cm}
		\begin{minipage}[c]{12.5cm}
			\vskip 0.15 cm
			\centerline{\ind{#1}}
			\vskip 0.15 cm
			\centerline{\textbf{#2}}
			\if\relax\detokenize{#3}\relax\else
				\vskip 0.15 cm
				\centerline{\textit{#3}}
			\fi
			\if\relax\detokenize{#4}\relax\else
				\vskip 0.15 cm
				\centerline{\textit{#4}}
			\fi
			\vskip 0.15 cm
		\end{minipage}
	\end{tcolorbox}
}
%====================== BỔ SUNG CHO BÀI HỌC, CHUYÊN ĐỀ ======================
\usepackage[shortlabels]{enumitem}
\usepackage{tabularx}
%--- Danh sách gọn: \begin{enumerate}[gon] ... \end{enumerate}
\SetEnumitemKey{gon}{leftmargin=*,itemsep=1pt,parsep=0pt,topsep=3pt}
%--- Khung lý thuyết: \begin{khung}{} ... \end{khung}; tiêu đề chỉ dùng khi thật ngắn ("Mẹo", "Lỗi thường gặp")
\newtcolorbox{khung}[1]{breakable,enhanced,colframe=gold1,colback=gold5!45,
	boxrule=0.4mm,fonttitle=\bfseries\sffamily,coltitle=white,colbacktitle=gold1,
	attach boxed title to top left={xshift=6mm,yshift=-3mm},
	boxed title style={arc=1mm,boxrule=0.3mm},title={#1}}
%--- Tiêu đề mục nhỏ trong một dạng: \muc{Dấu hiệu}, \muc{Phương pháp}, \muc{Bài tự luyện}
\makeatletter
\newcommand{\muc}[1]{\par\medskip\noindent\indam{#1}\par\nobreak\smallskip\@afterindentfalse\@afterheading}
\makeatother
%--- Ô công thức cần nhớ (căn giữa, viền vàng, vừa khít nội dung)
%     \begin{center}\congthuc{$Q(n)=(n-1)!$}\end{center}
\newtcbox{\congthuc}{on line,colframe=gold1,colback=gold5!45,boxrule=0.5pt,arc=2pt,
	left=8pt,right=8pt,top=4pt,bottom=4pt}
%--- Công tắc lời giải bài tự luyện (ví dụ mẫu luôn dùng \loigiai)
%     \loigiaiTL[{đáp số}]{lời giải đầy đủ}
%     \TLfull -> in lời giải đầy đủ;  \TLdapso -> chỉ in một dòng "Đáp số";  \TLan -> ẩn hoàn toàn
%     (\CheDo là dethi thì mọi lời giải và đáp số đều ẩn)
\makeatletter
\newif\ifTLhien   \TLhientrue
\newif\ifTLdapso  \TLdapsofalse
\newcommand{\TLfull}{\TLhientrue\TLdapsofalse}
\newcommand{\TLdapso}{\TLhientrue\TLdapsotrue}
\newcommand{\TLan}{\TLhienfalse}
\newcommand{\loigiaiTL}[2][]{%
	\ifTLhien
		\ifTLdapso
			\def\@tlarg{#1}%
			\ifx\@tlarg\@empty\else
				\ifdefstring{\kindOPTN}{dethi}{}{\par\noindent{\bfseries\color{gold1}Đáp số:} #1\par}%
			\fi
		\else
			\loigiai{#2}%
		\fi
	\fi}
\makeatother
\TLdapso% mặc định: chỉ in đáp số bài tự luyện
%====================== BẢNG ĐÚNG/SAI (\choiceTFt) ======================
%  Hàng tiêu đề nền nâu đồng chữ trắng (Phát biểu | Đ | S), các hàng tô xen kẽ kem / trắng,
%  đường kẻ vàng kim, nhãn a) b) c) d) màu nâu đồng. Bản có đáp án: cột được chọn có vòng tròn
%  nâu đồng với dấu tích trắng (câu đúng tích ở cột Đúng, câu sai tích ở cột Sai).
%  Tiêu đề ghi đủ "Đúng | Sai": \OPTN{viettat=0} sau \input (cột Đúng/Sai rộng ra, phát biểu hẹp lại).
%  Dùng bảng gốc của ex_test: \def\BangDSgoc{1} trước \input{preamble-eduzy-v2}.
\providecommand{\BangDSgoc}{0}
\ifnum\BangDSgoc=0
\usepackage{colortbl,pifont}
\colorlet{ezDSke}{gold2!60!gold1}% đường kẻ bảng (nét 0,4pt như ex_test để không đổi chiều cao bảng)
\colorlet{ezDSle}{gold5!60}% hàng a), c)
\colorlet{ezDSchan}{white}% hàng b), d)
%--- dấu chọn ở cột Đúng/Sai
%    (bề rộng 0, căn giữa ô: không làm đổi bề rộng cột của ex_test, không báo tràn lề)
\newcommand{\EZdauchon}{\makebox[0pt]{\tikz[baseline=-1ex]\node[circle,fill=gold1,inner sep=0pt,
	minimum size=1.2em,text=white,font=\small]{\ding{51}};}}
\def\tickT{\EZdauchon}\def\tickF{\EZdauchon}
\AtBeginDocument{\sbox{\ccTFTab}{\EZdauchon}}
%--- nhãn a) b) c) d) màu nâu đồng (ở bản lời giải, nhãn câu đúng vẫn đóng khung như ex_test)
\renewcommand{\FalseTF}{\stepcounter{dapan}\textbf{\color{gold1}\DapAnTF\sepTF}\ignorespaces}
%--- trước khi dựng bảng: kẻ màu vàng kim (bề rộng cột giữ nguyên như ex_test); sau bảng trả lại màu đen
\newcount\EZhangTF
\newcommand{\EZbatdauTF}{%
	\global\EZhangTF=0\relax
	\arrayrulecolor{ezDSke}}
\newcommand{\EZketthucTF}{\arrayrulecolor{black}}
%--- hàng tiêu đề nền nâu đồng, chữ trắng
\makeatletter
\newcommand{\EZo@TF}[1]{\centering\bfseries\color{white}#1}
\newcommand{\EZtieudeTF}{%
	\ifcase\cotDS
		\ifnum\cotSTT=1\relax
			\def\rowHeader{\hline\rowcolor{gold1}\EZo@TF{TT}&\EZo@TF{\phatbieu}\tabularnewline}%
		\else
			\def\rowHeader{\hline\rowcolor{gold1}\EZo@TF{\phatbieu}\tabularnewline}%
		\fi
	\or
		\ifnum\cotSTT=1\relax
			\def\rowHeader{\hline\rowcolor{gold1}\EZo@TF{TT}&\EZo@TF{\phatbieu}&\EZo@TF{\Dung}\tabularnewline}%
		\else
			\def\rowHeader{\hline\rowcolor{gold1}\EZo@TF{\phatbieu}&\EZo@TF{\Dung}\tabularnewline}%
		\fi
	\or
		\ifnum\cotSTT=1\relax
			\def\rowHeader{\hline\rowcolor{gold1}\EZo@TF{TT}&\EZo@TF{\phatbieu}&\EZo@TF{\Dung}&\EZo@TF{\Sai}\tabularnewline}%
		\else
			\def\rowHeader{\hline\rowcolor{gold1}\EZo@TF{\phatbieu}&\EZo@TF{\Dung}&\EZo@TF{\Sai}\tabularnewline}%
		\fi
	\fi}
%--- tô xen kẽ từng hàng (gọi ngay sau \hline đầu mỗi hàng)
\newcommand{\EZmauhangTF}{%
	\noalign{\global\advance\EZhangTF\@ne}%
	\ifodd\EZhangTF\rowcolor{ezDSle}\else\rowcolor{ezDSchan}\fi}
%--- gắn vào \choiceTFt / \writeChoiceTFt của ex_test
\xpretocmd{\choiceTFt}{\EZbatdauTF}{}{\typeout{[preamble-eduzy-v2] Không gắn được bảng Đúng/Sai (đầu bảng)}}
\xpatchcmd{\choiceTFt}{\ifnum\dongPB=0\def\rowHeader{}\fi}{\EZtieudeTF\ifnum\dongPB=0\def\rowHeader{}\fi}
	{}{\typeout{[preamble-eduzy-v2] Không gắn được bảng Đúng/Sai (hàng tiêu đề)}}
\xpatchcmd{\writeChoiceTFt}{\hline\xdef\dotEX{\sdotEX}}{\hline\EZmauhangTF\xdef\dotEX{\sdotEX}}
	{}{\typeout{[preamble-eduzy-v2] Không gắn được bảng Đúng/Sai (tô hàng)}}
\xpatchcmd{\writeChoiceTFt}{{\end{longtable}}}{{\end{longtable}\EZketthucTF}}
	{}{\typeout{[preamble-eduzy-v2] Không gắn được bảng Đúng/Sai (cuối longtable)}}
\xpatchcmd{\writeChoiceTFt}{{\end{xtabular}\end{center}\medskip\noindent}}{{\end{xtabular}\end{center}\EZketthucTF\medskip\noindent}}
	{}{\typeout{[preamble-eduzy-v2] Không gắn được bảng Đúng/Sai (cuối xtabular)}}
\fi
%=====================================================================
\renewcommand{\parallel}{\mathrel{/\!\!/}}

%  trong file đề (tools/dan_trang/dan_trang.py tự ghi khi đặt \NeedspaceCau):
%    \def\TachCauDS{1}   câu Đúng/Sai được tách: câu dẫn (kèm hình) ở cuối trang, bảng Phát biểu sang trang sau
%                        (câu dẫn và hình không tách nhau, bảng không tách giữa hai trang)
%    \def\TachCauLG{1}   bản lời giải (\CheDo khác dethi): phần đề của câu cao hơn 1/3 trang được phép tách
%  Lệnh dùng trong file đề (giá trị do tools/dan_trang/dan_trang.py đo ở bản dethi):
%    \NeedspaceCau[<cao câu dẫn + hình>]{<cao cả khối câu>}   ngay trên \begin{ex} (câu đầu phần: trên tiêu đề phần)
%        công tắc tắt: đúng bằng \Needspace{<cao cả khối câu>} như cũ;
%        \TachCauDS=1 và có [..] (câu Đúng/Sai được chọn tách): chỉ giữ câu dẫn + hình;
%        \TachCauLG=1 ở bản lời giải: đòi tối đa 1/3 \textheight.
%    \NeedspaceBang{<cao bảng Phát biểu>}%   ngay trước \choiceTFt: \TachCauDS=1 thì giữ trọn bảng trên một trang;
%        công tắc tắt: không làm gì.
\providecommand{\TachCauDS}{0}
\providecommand{\TachCauLG}{0}
\newlength{\EZnsCau}
\newcommand{\NeedspaceCau}[2][]{%
	\setlength{\EZnsCau}{#2}%
	\ifnum\TachCauDS=1 \if\relax\detokenize{#1}\relax\else\setlength{\EZnsCau}{#1}\fi\fi
	\ifnum\TachCauLG=1 \ifdefstring{\kindOPTN}{dethi}{}{\ifdim\EZnsCau>0.3333\textheight\setlength{\EZnsCau}{0.3333\textheight}\fi}\fi
	\Needspace{\EZnsCau}}
\newcommand{\NeedspaceBang}[1]{\ifnum\TachCauDS=1 \Needspace{#1}\fi}

%--- Thang màu: đậm -> nhạt
\definecolor{gold1}{RGB}{126,96,0}    % nâu đồng đậm  #7E6000
\definecolor{gold2}{RGB}{192,144,3}   % vàng kim      #C09003
\definecolor{gold3}{RGB}{255,215,99}  % vàng nhạt     #FFD763
\definecolor{gold4}{RGB}{255,231,153} % vàng rất nhạt #FFE799
\definecolor{gold5}{RGB}{255,243,205} % vàng kem      #FFF3CD
%--- Ghi đè tông xanh/đỏ cũ
\definecolor{myblue}{RGB}{126,96,0}
\colorlet{blue}{gold1}
\colorlet{cyan}{gold2}
\colorlet{teal}{gold2}

%--- Tên chương mặc định (file đề nên tự khai báo lại)
\providecommand{\tenchuong}{}
%--- Chân trang
\providecommand{\thuonghieu}{TOÁN 12 -- Eduzy.com.vn}
\providecommand{\tacgia}{Tô Gia Bảo -- 0353703783}

%--- Chữ cạnh logo ở đầu trang
\providecommand{\khauhieu}{Eduzy.com.vn}

%====================== THÔNG SỐ ĐẦU TRANG, CHÂN TRANG (v2) ======================
%  Dải nền đậm tràn ngang trang; ô chữ vàng kem có cạnh xiên, kèm 2 vạch xiên chuyển màu.
%  Đầu trang lẻ : logo + \khauhieu (trái)   | ô chữ \tenchuong (phải)
%  Chân trang lẻ: ô chữ \thuonghieu (trái)  | Trang (số) (phải)
%  Trang chẵn soi gương trang lẻ; chân trang chẵn ghi \tacgia.
%--- Màu: chỉ dùng thang vàng kim Eduzy
\colorlet{ezNen}{gold1}% dải nền đậm
\colorlet{ezVach1}{gold3}% vạch xiên sát dải nền (đậm hơn)
\colorlet{ezVach2}{gold4}% vạch xiên sát ô chữ (nhạt hơn)
\colorlet{ezO}{gold5}% ô chữ vàng kem
%--- Kích thước
\newlength{\EZdai}\setlength{\EZdai}{1.1cm}% chiều cao dải nền
\newlength{\EZo}\setlength{\EZo}{0.68cm}% chiều cao ô chữ
\newlength{\EZxien}\setlength{\EZxien}{0.55cm}% độ lệch ngang của cạnh xiên
\newlength{\EZvach}\setlength{\EZvach}{0.75cm}% bề rộng mỗi vạch xiên
\newlength{\EZkhe}\setlength{\EZkhe}{0.12cm}% khe trắng giữa vạch và ô chữ
\newlength{\EZdem}\setlength{\EZdem}{0.4cm}% khoảng đệm từ chữ tới cạnh xiên
\newlength{\EZle}\setlength{\EZle}{1.8cm}% lề trái/phải của chữ (bằng lề thân bài)
\newlength{\EZrongmin}\setlength{\EZrongmin}{4.5cm}% ô chữ rộng tối thiểu
\newlength{\EZrong}% bề rộng chữ trong ô (tính tự động)
%--- Font
\newcommand{\EZfontO}{\fontsize{10.5pt}{12pt}\selectfont\bfseries}% chữ trong ô
\newcommand{\EZfontKH}{\fontfamily{qag}\fontsize{11pt}{12pt}\selectfont\bfseries}% \khauhieu cạnh logo
\newcommand{\EZfontTrang}{\fontsize{11.5pt}{12pt}\selectfont\bfseries}% chữ "Trang" và số trang
%--- Đo bề rộng chữ trong ô
\newcommand{\EZdorong}[1]{%
	\settowidth{\EZrong}{\EZfontO #1}%
	\ifdim\EZrong<\EZrongmin\setlength{\EZrong}{\EZrongmin}\fi}
%--- Logo + chữ trắng \khauhieu (dựng sẵn một lần khi vào \begin{document})
\newsavebox{\EZlogobox}
\AtBeginDocument{\sbox{\EZlogobox}{%
		\raisebox{-0.5\height}{\eduzylogo[1cm]}\hspace{0.15cm}%
		\raisebox{-0.5\height}{\color{white}\EZfontKH\khauhieu}}}
%--- Hai vạch xiên + khe trắng. #1: điểm gốc ở mép dưới, #2: hướng (1 sang phải, -1 sang trái),
%    #3: hướng nghiêng (1 kiểu "/", -1 kiểu "\")
\newcommand{\EZvachxien}[3]{%
	\fill[white] (#1) --++({(#3)*\EZxien},\EZo) --++({(#2)*\EZkhe},0) --++({-(#3)*\EZxien},-\EZo) --cycle;
	\fill[ezVach2] ([xshift=(#2)*\EZkhe]#1) --++({(#3)*\EZxien},\EZo) --++({(#2)*\EZvach},0) --++({-(#3)*\EZxien},-\EZo) --cycle;
	\fill[ezVach1] ([xshift=(#2)*(\EZkhe+\EZvach)]#1) --++({(#3)*\EZxien},\EZo) --++({(#2)*\EZvach},0) --++({-(#3)*\EZxien},-\EZo) --cycle;
	\draw[white,line width=0.6pt] ([xshift=(#2)*(\EZkhe+\EZvach)]#1) --++({(#3)*\EZxien},\EZo);
}
%--- Chữ trong ô (căn giữa theo chiều cao ô, không lệch vì dấu tiếng Việt)
\tikzset{ezchu/.style={inner sep=0pt,text=ezNen,text height=0.72em,text depth=0.22em}}

%====================== LOGO CHÌM GIỮA TRANG ======================
%  Vẽ cùng đầu trang nên nằm dưới chữ của thân bài, có trên mọi trang.
%  Tắt: \logochimtat (từ trang đó trở đi), bật lại: \logochimbat.
%  Đổi cỡ, độ đậm: khai báo trước \input, ví dụ \def\EZchimcao{9cm} \def\EZchimdo{0.1}
\providecommand{\EZchimcao}{10cm}% chiều cao logo chìm
\providecommand{\EZchimdo}{0.12}% độ đậm (0 = trong suốt, 1 = rõ hẳn)
\newif\ifEZchim \EZchimtrue
\newcommand{\logochimtat}{\EZchimfalse}
\newcommand{\logochimbat}{\EZchimtrue}
%  Dựng sẵn logo mờ một lần khi vào \begin{document}; transparency group làm mờ cả logo
%  một lượt nên không lộ vệt ở chỗ các dải màu chồng nhau.
\newsavebox{\EZchimtho}
\newsavebox{\EZchimbox}
\AtBeginDocument{%
	\sbox{\EZchimtho}{\eduzylogo[\EZchimcao]}%
	\sbox{\EZchimbox}{\begin{tikzpicture}
			\begin{scope}[transparency group,opacity=\EZchimdo]
				\node[inner sep=0pt] {\usebox{\EZchimtho}};
			\end{scope}
	\end{tikzpicture}}}
\newcommand{\EZvechim}{%
	\ifEZchim
		\node[inner sep=0pt] at (current page.center) {\usebox{\EZchimbox}};
	\fi}

%====================== ĐẦU TRANG - TRANG LẺ ======================
%  (logo bên trái, ô chữ \tenchuong bên phải)
\newcommand{\myfancyheadOdd}{%
	\EZdorong{\tenchuong}%
	\begin{tikzpicture}[remember picture,overlay]
		\coordinate (T) at (current page.north west);
		\coordinate (P) at (current page.north east);
		%--- logo chìm giữa trang
		\EZvechim
		%--- dải nền
		\fill[ezNen] (T) rectangle ([yshift=-\EZdai]P);
		%--- ô chữ: X là góc dưới bên trái, cạnh trái xiên kiểu "/"
		\coordinate (X) at ([xshift=-\EZle-\EZrong-\EZdem-\EZxien,yshift=-\EZdai]P);
		\fill[ezO] (X) --++(\EZxien,\EZo) -- ([yshift=-\EZdai+\EZo]P) -- ([yshift=-\EZdai]P) --cycle;
		\EZvachxien{X}{-1}{1}
		\node[ezchu,anchor=east] at ([xshift=-\EZle,yshift=-\EZdai+0.5*\EZo]P) {\EZfontO\tenchuong};
		%--- logo + tên
		\node[inner sep=0pt,anchor=west] at ([xshift=\EZle,yshift=-0.5*\EZdai]T) {\usebox{\EZlogobox}};
	\end{tikzpicture}%
}

%====================== ĐẦU TRANG - TRANG CHẴN =====================
%  (soi gương: ô chữ \tenchuong bên trái, logo bên phải)
\newcommand{\myfancyheadEven}{%
	\EZdorong{\tenchuong}%
	\begin{tikzpicture}[remember picture,overlay]
		\coordinate (T) at (current page.north west);
		\coordinate (P) at (current page.north east);
		%--- logo chìm giữa trang
		\EZvechim
		%--- dải nền
		\fill[ezNen] (T) rectangle ([yshift=-\EZdai]P);
		%--- ô chữ: X là góc dưới bên phải, cạnh phải xiên kiểu "\"
		\coordinate (X) at ([xshift=\EZle+\EZrong+\EZdem+\EZxien,yshift=-\EZdai]T);
		\fill[ezO] (X) --++(-\EZxien,\EZo) -- ([yshift=-\EZdai+\EZo]T) -- ([yshift=-\EZdai]T) --cycle;
		\EZvachxien{X}{1}{-1}
		\node[ezchu,anchor=west] at ([xshift=\EZle,yshift=-\EZdai+0.5*\EZo]T) {\EZfontO\tenchuong};
		%--- logo + tên
		\node[inner sep=0pt,anchor=east] at ([xshift=-\EZle,yshift=-0.5*\EZdai]P) {\usebox{\EZlogobox}};
	\end{tikzpicture}%
}

%====================== CHÂN TRANG - TRANG LẺ ======================
%  (ô chữ \thuonghieu bên trái, "Trang (số)" bên phải)
\newcommand{\myfancyfootOdd}{%
	\EZdorong{\thuonghieu}%
	\begin{tikzpicture}[remember picture,overlay]
		\coordinate (T) at (current page.south west);
		\coordinate (P) at (current page.south east);
		%--- dải nền
		\fill[ezNen] (T) rectangle ([yshift=\EZdai]P);
		%--- ô chữ sát mép trên dải: X là góc dưới bên phải, cạnh phải xiên kiểu "/"
		\coordinate (X) at ([xshift=\EZle+\EZrong+\EZdem,yshift=\EZdai-\EZo]T);
		\fill[ezO] ([yshift=\EZdai-\EZo]T) -- (X) --++(\EZxien,\EZo) -- ([yshift=\EZdai]T) --cycle;
		\EZvachxien{X}{1}{1}
		\node[ezchu,anchor=west] at ([xshift=\EZle,yshift=\EZdai-0.5*\EZo]T) {\EZfontO\thuonghieu};
		%--- số trang
		\node[circle,draw=white,line width=0.9pt,text=white,inner sep=0pt,minimum size=0.62cm,anchor=east]
		(so) at ([xshift=-\EZle,yshift=0.5*\EZdai]P) {\EZfontTrang\thepage};
		\node[inner sep=0pt,text=white,anchor=base east] (tr) at ([xshift=-0.15cm]so.base west) {\EZfontTrang Trang};
	\end{tikzpicture}%
}

%====================== CHÂN TRANG - TRANG CHẴN ====================
%  (soi gương: "Trang (số)" bên trái, ô chữ \tacgia bên phải)
\newcommand{\myfancyfootEven}{%
	\EZdorong{\tacgia}%
	\begin{tikzpicture}[remember picture,overlay]
		\coordinate (T) at (current page.south west);
		\coordinate (P) at (current page.south east);
		%--- dải nền
		\fill[ezNen] (T) rectangle ([yshift=\EZdai]P);
		%--- ô chữ sát mép trên dải: X là góc dưới bên trái, cạnh trái xiên kiểu "\"
		\coordinate (X) at ([xshift=-\EZle-\EZrong-\EZdem,yshift=\EZdai-\EZo]P);
		\fill[ezO] ([yshift=\EZdai-\EZo]P) -- (X) --++(-\EZxien,\EZo) -- ([yshift=\EZdai]P) --cycle;
		\EZvachxien{X}{-1}{-1}
		\node[ezchu,anchor=east] at ([xshift=-\EZle,yshift=\EZdai-0.5*\EZo]P) {\EZfontO\tacgia};
		%--- số trang
		%    (vòng tròn ẩn so0 nằm giữa dải để lấy đường chân chữ, giống trang lẻ)
		\node[circle,inner sep=0pt,minimum size=0.62cm] (so0) at ([yshift=0.5*\EZdai]T) {\EZfontTrang\phantom{\thepage}};
		\node[inner sep=0pt,text=white,anchor=base west] (tr) at ($(so0.base -| T)+(\EZle,0)$) {\EZfontTrang Trang};
		\node[circle,draw=white,line width=0.9pt,text=white,inner sep=0pt,minimum size=0.62cm,anchor=base west]
		(so) at ([xshift=0.15cm]tr.base east) {\EZfontTrang\thepage};
	\end{tikzpicture}%
}

%--- Giữ tương thích với các file cũ đang gọi \myfancyhead, \myfancyfoot
\newcommand{\myfancyhead}{\myfancyheadOdd}
\newcommand{\myfancyfoot}{\myfancyfootOdd}

%======================= GẮN VÀO TRANG =============================
\pagestyle{fancy}
\fancyhf{}
\fancyhead[LO]{\myfancyheadOdd}
\fancyhead[RE]{\myfancyheadEven}
\fancyfoot[LO]{\myfancyfootOdd}
\fancyfoot[RE]{\myfancyfootEven}
\renewcommand{\headrulewidth}{0pt}
\renewcommand{\footrulewidth}{0pt}
\fancypagestyle{plain}{%
	\fancyhf{}
	\fancyhead[LO]{\myfancyheadOdd}
	\fancyhead[RE]{\myfancyheadEven}
	\fancyfoot[LO]{\myfancyfootOdd}
	\fancyfoot[RE]{\myfancyfootEven}
	\renewcommand{\headrulewidth}{0pt}
	\renewcommand{\footrulewidth}{0pt}
}
%=================== HẾT PHẦN GỘP TỪ eduzy-de.tex ===================
\newcommand{\indam}[1]{
	{\fontsize{12pt}{9pt}\sffamily \selectfont\bfseries\color{gold1} #1}
}
\newcommand{\ind}[1]{
	{\fontsize{14pt}{9pt}\fontfamily{qag} \selectfont\bfseries\color{gold1} #1}
}
\renewcommand{\FalseEX}{\stepcounter{dapan}\textbf{\selectfont\color{gold1}\Alph{dapan}.}\ignorespaces}
\renewtheorem{ex}{\color{gold1}Câu}
\renewcommand{\baselinestretch}{1.1}
\newcommand{\hoac}[1]{
	\left[\begin{aligned}#1\end{aligned}\right.}
\newcommand{\heva}[1]{
	\left\{\begin{aligned}#1\end{aligned}\right.}
\usepackage{esvect}
\def\vec{\vv}
\def\vt{\vv}
\def\overrightarrow{\vv}
%====================== LOGO EDUZY VẼ BẰNG TIKZ ======================
%  Đường viền dò từ logo gốc; dải màu vàng kim đo trực tiếp trên ảnh gốc
%  (nghiêng 21 độ, tô bằng 80 dải hẹp nên hiển thị giống nhau trên mọi trình xem PDF).
%  Dùng: \eduzylogo          -> cao 2.4cm (mặc định)
%        \eduzylogo[1.5cm]   -> cao 1.5cm
\definecolor{ezText}{RGB}{245,200,46}
\newcommand{\eduzylogo}[1][2.4cm]{%
	\resizebox{!}{#1}{%
		\begin{tikzpicture}[x=0.01cm,y=0.01cm]
			\path[use as bounding box] (0,0) rectangle (224,-224);
			% mũ tốt nghiệp + vương miện
			\begin{scope}
				\clip
				(119.4,-176.1)..controls(118.9,-176.1)and(116.8,-176)..(114.6,-176)
				..controls(101.4,-176)and(85.8,-175.9)..(83.7,-175.7)
				..controls(81.8,-175.6)and(81,-175.6)..(80.2,-175.7)
				..controls(78.8,-175.9)and(77.7,-175.9)..(76.2,-175.7)
				..controls(74.9,-175.5)and(70.8,-175.5)..(64.3,-175.7)
				..controls(59.1,-175.9)and(58.3,-175.7)..(58.6,-173.7)
				..controls(58.6,-172.8)and(59.6,-169.6)..(60.3,-168)
				..controls(60.8,-166.9)and(61.4,-165.8)..(62.3,-165)
				..controls(63.1,-164.2)and(64.3,-163.4)..(65,-163.2)
				..controls(65.5,-163.1)and(72.3,-163)..(76.6,-163.1)
				..controls(77.9,-163.1)and(80.1,-163.1)..(81.4,-163)
				..controls(82.8,-162.9)and(86.3,-162.9)..(89.2,-162.9)
				..controls(93,-162.9)and(94.9,-162.8)..(96,-162.7)
				..controls(97.2,-162.5)and(101,-162.5)..(114,-162.4)
				..controls(123,-162.3)and(134.9,-162.2)..(140.2,-162.1)
				..controls(150.1,-162)and(152.2,-162.1)..(153.6,-162.4)
				..controls(154.4,-162.6)and(157.1,-163.8)..(158.1,-164.4)
				..controls(158.8,-164.8)and(159.7,-165.7)..(159.9,-166.3)
				..controls(160,-166.5)and(160.2,-167.1)..(160.2,-167.6)
				..controls(160.3,-168)and(160.6,-169.4)..(160.9,-170.6)
				..controls(161.2,-171.7)and(161.5,-173.1)..(161.6,-173.6)
				..controls(161.8,-174.4)and(161.8,-174.5)..(161.5,-174.9)
				..controls(161.2,-175.6)and(160.8,-175.9)..(159.9,-176)
				..controls(159.2,-176.1)and(154.1,-176.2)..(151.5,-176.2)
				..controls(148.1,-176.1)and(122.6,-176.1)..(121.6,-176.2)
				..controls(120.8,-176.2)and(119.8,-176.2)..(119.4,-176.1)--cycle
				(130.4,-160)..controls(129,-160)and(119.8,-160)..(110.1,-159.9)
				..controls(100.3,-159.8)and(88.5,-159.7)..(83.9,-159.6)
				..controls(77.8,-159.5)and(75.3,-159.5)..(74.5,-159.6)
				..controls(73.9,-159.7)and(72.5,-159.8)..(71.3,-159.8)
				..controls(69.6,-159.8)and(69,-159.8)..(68.5,-159.6)
				..controls(67.6,-159.3)and(66.8,-158.6)..(66.2,-157.6)
				..controls(66,-157.3)and(65.6,-156.5)..(65.3,-156.1)
				..controls(64.6,-155.1)and(64.5,-154.9)..(63.7,-152.9)
				..controls(63.2,-151.6)and(62.9,-150.9)..(61.8,-149.6)
				..controls(61.5,-149.2)and(61.1,-148.5)..(60.9,-147.9)
				..controls(60.7,-147.4)and(60.3,-146.6)..(60.1,-146.1)
				..controls(59.9,-145.5)and(59.5,-144.9)..(59,-144.3)
				..controls(58.5,-143.7)and(58.2,-143.2)..(57.9,-142.5)
				..controls(57.7,-141.9)and(57.3,-141.1)..(57,-140.7)
				..controls(56.6,-140.3)and(56.1,-139.3)..(55.8,-138.5)
				..controls(55.5,-137.8)and(55.1,-136.8)..(54.9,-136.5)
				..controls(54.7,-136.1)and(54.3,-135.2)..(53.9,-134.5)
				..controls(53.6,-133.8)and(53,-132.7)..(52.6,-132)
				..controls(51.6,-130.5)and(51.4,-130.2)..(51.1,-129)
				..controls(50.6,-127.5)and(50.6,-127.4)..(49.7,-126.3)
				..controls(49.2,-125.7)and(48.6,-124.8)..(48.3,-124.2)
				..controls(48,-123.6)and(47.5,-122.7)..(47.3,-122.2)
				..controls(47,-121.8)and(46.5,-120.7)..(46.1,-119.8)
				..controls(45.8,-118.9)and(45.2,-117.8)..(44.9,-117.2)
				..controls(44.2,-116.1)and(43.3,-114.4)..(42.3,-112.2)
				..controls(41.8,-111.3)and(41.3,-110.4)..(41.1,-110.1)
				..controls(40.5,-109.4)and(40.3,-109)..(40,-107.8)
				..controls(39.8,-107)and(39.6,-106.5)..(39.2,-106)
				..controls(38.9,-105.6)and(38.4,-104.7)..(38,-104)
				..controls(37.7,-103.4)and(37.2,-102.5)..(37,-102.1)
				..controls(36.7,-101.7)and(36.3,-100.8)..(36,-100.1)
				..controls(35.6,-98.9)and(35.3,-98.4)..(34.6,-97.6)
				..controls(34.2,-97.1)and(34.2,-96.3)..(34.6,-96)
				..controls(35,-95.7)and(35.5,-95.9)..(36.2,-96.7)
				..controls(36.5,-97)and(37.2,-97.5)..(37.6,-97.7)..controls(38,-98)and(38.6,-98.4)..(39.1,-98.9)
				..controls(39.6,-99.4)and(40.3,-100)..(40.6,-100.3)
				..controls(41.3,-100.8)and(42.4,-101.8)..(43.7,-102.9)
				..controls(44,-103.2)and(44.6,-103.7)..(45,-103.9)
				..controls(45.4,-104.2)and(45.9,-104.6)..(46.1,-104.8)
				..controls(47.1,-105.9)and(48.8,-107.1)..(50.5,-107.9)
				..controls(51.8,-108.6)and(52.3,-109)..(53.4,-110)
				..controls(54.8,-111.2)and(55.6,-111.6)..(57,-112.1)
				..controls(57.5,-112.3)and(58.5,-112.8)..(59.4,-113.3)
				..controls(60.2,-113.7)and(61.3,-114.3)..(61.9,-114.5)
				..controls(62.4,-114.7)and(63.3,-115.1)..(63.8,-115.3)
				..controls(64.4,-115.6)and(65.4,-115.9)..(66.2,-116.1)
				..controls(67,-116.3)and(68.4,-116.6)..(69.2,-116.9)
				..controls(72.3,-117.9)and(72.8,-118)..(78,-118.1)
				..controls(84,-118.3)and(86.1,-118.1)..(88.1,-117.1)
				..controls(89.4,-116.5)and(93,-115)..(93.7,-114.9)
				..controls(94.8,-114.6)and(95.4,-114.2)..(95.9,-113.4)
				..controls(96.2,-113)and(96.7,-112.5)..(97.1,-112.2)
				..controls(98,-111.6)and(101.3,-108.4)..(102.1,-107.2)
				..controls(102.4,-106.8)and(102.9,-106.2)..(103.2,-105.8)
				..controls(103.5,-105.5)and(104,-104.8)..(104.2,-104.2)
				..controls(104.5,-103.6)and(105,-102.8)..(105.3,-102.4)
				..controls(105.6,-102)and(106.1,-101.2)..(106.5,-100.6)
				..controls(107.1,-99.5)and(108.6,-97.5)..(109.5,-96.5)
				..controls(109.9,-96.1)and(110.1,-96)..(110.4,-96)
				..controls(111.1,-96)and(111.6,-96.5)..(111.8,-97.6)
				..controls(112,-98.6)and(112.3,-99.3)..(112.8,-99.9)
				..controls(113.7,-100.9)and(115.2,-103.1)..(115.7,-104.2)
				..controls(116.1,-105)and(116.5,-105.5)..(116.9,-106)
				..controls(117.3,-106.4)and(117.8,-106.9)..(118,-107.3)
				..controls(118.2,-107.6)and(118.6,-108)..(119,-108.3)
				..controls(119.6,-108.8)and(121.1,-110.3)..(121.8,-111.2)
				..controls(122.1,-111.7)and(122.5,-112)..(123.4,-112.4)
				..controls(124.8,-113.2)and(125.8,-113.8)..(127,-114.6)
				..controls(127.5,-114.9)and(128.1,-115.3)..(128.2,-115.4)
				..controls(128.4,-115.5)and(129.3,-115.8)..(130.2,-116)
				..controls(131.1,-116.2)and(132.5,-116.6)..(133.3,-116.9)
				..controls(135.7,-117.9)and(136.1,-117.9)..(142.5,-117.8)
				..controls(145.8,-117.8)and(148.4,-117.7)..(148.8,-117.6)
				..controls(149.2,-117.6)and(150.2,-117.2)..(151,-116.8)
				..controls(152,-116.4)and(152.8,-116)..(153.6,-115.9)
				..controls(154.5,-115.7)and(155.2,-115.4)..(156.3,-114.9)
				..controls(157.1,-114.5)and(158.2,-114)..(158.7,-113.8)
				..controls(159.1,-113.7)and(159.8,-113.3)..(160.2,-113)
				..controls(160.6,-112.7)and(161.3,-112.4)..(161.7,-112.2)
				..controls(162.7,-111.9)and(163.2,-111.7)..(164.1,-111.1)
				..controls(164.5,-110.8)and(165.3,-110.4)..(165.9,-110)
				..controls(166.5,-109.7)and(167.3,-109.2)..(167.6,-108.9)
				..controls(167.9,-108.6)and(168.6,-108.2)..(169,-108)
				..controls(169.4,-107.8)and(170.2,-107.3)..(170.7,-106.9)
				..controls(171.2,-106.5)and(172.1,-105.8)..(172.7,-105.4)
				..controls(173.2,-105)and(174.1,-104.4)..(174.5,-104)
				..controls(174.9,-103.7)and(175.7,-103.1)..(176.2,-102.6)
				..controls(176.7,-102.2)and(177.4,-101.5)..(177.9,-101.2)
				..controls(178.3,-100.8)and(178.8,-100.4)..(179,-100.2)
				..controls(179.2,-100)and(179.6,-99.6)..(180,-99.3)
				..controls(180.4,-99.1)and(180.9,-98.6)..(181.2,-98.4)
				..controls(181.5,-98.1)and(182.1,-97.7)..(182.5,-97.5)
				..controls(183.2,-97.2)and(183.5,-96.9)..(184.3,-95.9)
				..controls(184.6,-95.6)and(184.6,-95.6)..(184.8,-95.7)
				..controls(185,-95.8)and(185.1,-96)..(185.1,-96.2)
				..controls(185.1,-96.5)and(185,-96.7)..(184.6,-97.2)
				..controls(184.1,-97.8)and(184,-97.9)..(183.9,-98.8)
				..controls(183.7,-100.3)and(183.3,-101.3)..(182.9,-101.8)
				..controls(182.6,-102.1)and(182.3,-102.6)..(182.2,-102.9)
				..controls(182,-103.3)and(181.6,-104)..(181.2,-104.6)
				..controls(180.9,-105.2)and(180.5,-106.1)..(180.4,-106.6)
				..controls(179.9,-108)and(179.7,-108.5)..(179,-109.4)
				..controls(178.6,-109.8)and(178.2,-110.5)..(177.9,-111)
				..controls(177.7,-111.5)and(177.3,-112.3)..(177,-112.7)
				..controls(176.6,-113.2)and(176.4,-113.7)..(176.3,-114.2)
				..controls(176.2,-114.7)and(175.9,-115.2)..(175.4,-116)
				..controls(175,-116.6)and(174.4,-117.7)..(174,-118.5)
				..controls(173.7,-119.2)and(173.3,-120.1)..(173,-120.5)
				..controls(172.8,-120.9)and(172.4,-121.8)..(172.1,-122.5)
				..controls(171.8,-123.2)and(171.3,-124.1)..(171,-124.6)
				..controls(170.8,-125.1)and(170.3,-126)..(169.9,-126.6)
				..controls(169.6,-127.3)and(169.2,-128.1)..(169,-128.3)
				..controls(168.8,-128.6)and(168.5,-129.3)..(168.3,-129.8)
				..controls(168.1,-130.3)and(167.8,-131.1)..(167.5,-131.6)
				..controls(166.9,-132.5)and(164.2,-137.9)..(164.2,-138.1)
				..controls(164.2,-138.1)and(163.9,-138.8)..(163.4,-139.7)
				..controls(162.9,-140.5)and(162.4,-141.5)..(162.3,-141.9)
				..controls(161.9,-143)and(161.6,-143.4)..(160.9,-144.1)
				..controls(160.4,-144.6)and(160.3,-144.8)..(160.1,-145.4)
				..controls(160,-145.8)and(159.6,-146.8)..(159.2,-147.8)
				..controls(158.7,-148.8)and(158.1,-150)..(157.9,-150.6)
				..controls(157.6,-151.1)and(157.2,-151.9)..(156.9,-152.4)
				..controls(156.5,-152.8)and(156.1,-153.7)..(155.8,-154.4)
				..controls(155.6,-155)and(155.2,-155.9)..(154.9,-156.3)
				..controls(154.7,-156.7)and(154.4,-157.4)..(154.2,-157.9)
				..controls(154.1,-158.4)and(153.8,-159)..(153.5,-159.3)
				..controls(152.8,-160.1)and(153,-160.1)..(142.4,-160.1)
				..controls(137.2,-160.1)and(131.8,-160.1)..(130.4,-160)--cycle
				(141.5,-103.7)..controls(140.8,-103)and(139.8,-102.1)..(139.2,-101.6)
				..controls(138.5,-101.2)and(137.9,-100.6)..(137.8,-100.4)
				..controls(137,-99.5)and(134,-96.5)..(133.2,-96)
				..controls(132.7,-95.7)and(131.8,-94.9)..(131.1,-94.3)
				..controls(129.4,-92.7)and(127.9,-91.6)..(126.6,-91.1)
				..controls(125.2,-90.5)and(124.1,-90)..(122.9,-89.1)
				..controls(122.1,-88.6)and(121.3,-88.2)..(119.9,-87.7)
				..controls(118.8,-87.3)and(117.5,-86.8)..(116.9,-86.5)
				..controls(114.3,-85.5)and(113.7,-85.4)..(109.7,-85.5)
				..controls(105.3,-85.7)and(103.9,-85.9)..(101.6,-87)
				..controls(100.7,-87.4)and(99.4,-87.8)..(97.9,-88)
				..controls(97,-88.2)and(96.6,-88.3)..(95.8,-88.9)
				..controls(95.4,-89.2)and(94.7,-89.6)..(94.2,-89.8)
				..controls(92.9,-90.3)and(92.3,-90.7)..(90.9,-92.1)
				..controls(90.2,-92.7)and(89.4,-93.3)..(89,-93.5)..controls(87,-94.5)and(85.2,-96)..(83.6,-97.9)
				..controls(83.3,-98.2)and(82.5,-98.9)..(81.8,-99.4)
				..controls(80.7,-100.1)and(80.5,-100.3)..(80,-101.1)
				..controls(78.7,-103.1)and(77.2,-103.9)..(76.4,-103)
				..controls(75.7,-102.2)and(75.7,-103.2)..(75.7,-83.6)
				..controls(75.6,-73.6)and(75.6,-64.8)..(75.5,-64)
				..controls(75.5,-63.2)and(75.5,-61.4)..(75.6,-59.8)--(75.7,-57.1)--(76,-56.8)
				..controls(76.8,-56.2)and(78,-56.3)..(79.8,-57.2)
				..controls(80.4,-57.5)and(81.3,-57.8)..(81.8,-57.9)
				..controls(82.3,-58.1)and(83,-58.4)..(83.6,-58.8)..controls(84.3,-59.2)and(85,-59.5)..(86,-59.8)
				..controls(86.8,-60.1)and(88,-60.6)..(88.7,-60.9)
				..controls(89.3,-61.2)and(90.6,-61.7)..(91.4,-62)
				..controls(92.2,-62.3)and(93.7,-62.9)..(94.6,-63.3)
				..controls(95.6,-63.7)and(96.8,-64.1)..(97.4,-64.3)
				..controls(98,-64.5)and(99,-64.9)..(99.7,-65.2)
				..controls(100.4,-65.6)and(101.5,-66)..(102.3,-66.2)
				..controls(104.2,-66.8)and(105,-67.1)..(106.2,-67.8)
				..controls(107.9,-68.7)and(110.2,-69.4)..(111.3,-69.2)
				..controls(112.5,-69)and(115,-68.2)..(115.6,-67.8)
				..controls(117.3,-66.7)and(117.6,-66.6)..(119.5,-66.2)
				..controls(121.9,-65.8)and(122.1,-65.7)..(123,-65.1)
				..controls(124.1,-64.4)and(124.4,-64.3)..(126,-63.9)
				..controls(127,-63.7)and(127.7,-63.4)..(128.4,-63.1)
				..controls(129,-62.8)and(129.7,-62.5)..(130.6,-62.4)
				..controls(132.1,-62.1)and(134,-61.5)..(135.2,-61)
				..controls(135.7,-60.8)and(136.7,-60.5)..(137.3,-60.3)
				..controls(138.4,-60)and(139.3,-59.7)..(141.1,-59)
				..controls(141.6,-58.7)and(142,-58.6)..(142.5,-58.6)
				..controls(143.2,-58.6)and(143.3,-58.7)..(143.6,-59)
				..controls(144,-59.5)and(144.1,-59.9)..(144,-60.8)
				..controls(144,-61.2)and(144,-71.1)..(144,-82.8)--(143.9,-104.1)--(143.7,-104.5)
				..controls(143.5,-104.7)and(143.3,-104.8)..(143,-104.9)
				..controls(142.6,-104.9)and(142.5,-104.8)..(141.5,-103.7)--cycle
				(27.8,-91.4)..controls(27.1,-90.9)and(26.3,-90)..(26.1,-89.5)
				..controls(26,-89.2)and(25.8,-88.6)..(25.7,-87.9)
				..controls(25.4,-86.2)and(25.6,-84.8)..(26.5,-83.5)
				..controls(27.4,-82.3)and(28.3,-82)..(29.4,-82.3)
				..controls(31.1,-82.8)and(32.1,-84.9)..(32,-87.4)
				..controls(31.9,-89.3)and(31.2,-90.6)..(29.9,-91.4)
				..controls(29.1,-91.9)and(28.6,-91.9)..(27.8,-91.4)--cycle
				(190.1,-91.3)..controls(189.7,-91.2)and(189.3,-90.9)..(189.1,-90.7)
				..controls(188.3,-89.8)and(187.6,-87.7)..(187.6,-86.2)
				..controls(187.6,-85.5)and(187.7,-85.2)..(188.2,-84.4)
				..controls(189.1,-82.6)and(190.1,-81.8)..(191.3,-82.1)
				..controls(192,-82.3)and(193.5,-83.7)..(193.8,-84.4)
				..controls(194.1,-85.1)and(194.2,-87.1)..(194,-88.3)
				..controls(193.8,-89.8)and(193.3,-90.5)..(192.3,-91.3)
				..controls(191.7,-91.7)and(190.9,-91.7)..(190.1,-91.3)--cycle
				(173.7,-83.4)..controls(172.2,-82.5)and(172,-82.5)..(170.6,-82.4)
				..controls(169.7,-82.4)and(168.8,-82.4)..(168.3,-82.5)
				..controls(167.2,-82.7)and(166.9,-82.7)..(166.6,-82.2)
				..controls(166.3,-81.8)and(166.3,-81)..(166.8,-80)
				..controls(167.2,-78.9)and(167.4,-77.4)..(167.7,-74.1)
				..controls(167.8,-72.8)and(167.9,-71.8)..(167.9,-71.7)
				..controls(168,-71.7)and(168,-70.2)..(168,-68.4)..controls(168,-65.4)and(168,-64.9)..(167.7,-63)
				..controls(167.3,-60.5)and(167.1,-59.9)..(166.4,-58.4)
				..controls(166.1,-57.7)and(165.8,-56.9)..(165.7,-56.6)
				..controls(165.7,-56.2)and(165.4,-55.6)..(165.2,-55.2)
				..controls(164.6,-53.8)and(161.9,-50.7)..(161,-50.3)
				..controls(159.8,-49.7)and(157.9,-50)..(155.2,-51.2)
				..controls(154.7,-51.4)and(153.4,-51.8)..(152.2,-52.1)
				..controls(150.8,-52.4)and(149.4,-52.9)..(148.6,-53.2)
				..controls(147.8,-53.6)and(146.7,-54)..(146,-54.2)
				..controls(145.4,-54.4)and(144.5,-54.7)..(144,-54.9)
				..controls(143.6,-55.2)and(142.7,-55.5)..(141.8,-55.7)
				..controls(141,-55.9)and(139.6,-56.3)..(138.7,-56.6)
				..controls(137.8,-57)and(136.3,-57.4)..(135.3,-57.6)
				..controls(133.8,-57.9)and(133.2,-58.1)..(132.1,-58.7)
				..controls(131.1,-59.2)and(130.3,-59.4)..(129.4,-59.6)
				..controls(128.7,-59.8)and(127.8,-60)..(127.4,-60.1)
				..controls(127,-60.3)and(126,-60.5)..(125.1,-60.7)
				..controls(123.2,-61.1)and(121.7,-61.5)..(120.2,-62.1)
				..controls(119.7,-62.3)and(118.1,-62.8)..(116.7,-63.2)
				..controls(114.6,-63.9)and(114,-64.1)..(113,-64.7)
				..controls(111.9,-65.4)and(111.8,-65.4)..(111,-65.5)
				..controls(109.9,-65.5)and(109.3,-65.4)..(107.7,-64.7)
				..controls(107,-64.4)and(106,-64)..(105.3,-63.9)
				..controls(103.5,-63.5)and(102.2,-63)..(101.2,-62.5)
				..controls(100.1,-61.8)and(99.3,-61.5)..(97.8,-61.1)
				..controls(97.2,-61)and(96.2,-60.7)..(95.7,-60.4)
				..controls(95.1,-60.1)and(94.2,-59.8)..(93.5,-59.6)
				..controls(92.9,-59.4)and(91.7,-59)..(91,-58.7)..controls(90.2,-58.4)and(88.9,-58)..(88,-57.7)
				..controls(87.2,-57.5)and(86.1,-57.1)..(85.5,-56.8)
				..controls(85,-56.6)and(84.1,-56.2)..(83.5,-56)..controls(82.9,-55.7)and(81.8,-55.2)..(81,-54.9)
				..controls(80.1,-54.5)and(79.1,-54.2)..(78.7,-54.1)
				..controls(77.9,-53.8)and(77.3,-53.6)..(76.2,-52.9)
				..controls(75.7,-52.6)and(75.2,-52.4)..(74.1,-52.2)
				..controls(72.5,-51.9)and(72.4,-51.8)..(71.3,-51.1)
				..controls(69.9,-50.1)and(69.4,-49.9)..(67.9,-49.8)
				..controls(67,-49.7)and(66.6,-49.6)..(66.3,-49.4)
				..controls(66.1,-49.3)and(65.7,-49)..(65.4,-48.8)
				..controls(65.1,-48.7)and(64.2,-48.4)..(63.6,-48.2)
				..controls(62.9,-48)and(61.7,-47.6)..(60.9,-47.2)
				..controls(60.1,-46.8)and(58.9,-46.4)..(58.2,-46.2)
				..controls(57.6,-46)and(56.8,-45.7)..(56.5,-45.5)
				..controls(54.9,-44.6)and(54.6,-44.5)..(53.2,-44.2)
				..controls(51.7,-43.8)and(51,-43.6)..(49.9,-42.9)
				..controls(49.5,-42.6)and(48.9,-42.4)..(47.5,-42)
				..controls(45.7,-41.6)and(45.6,-41.5)..(44.3,-40.6)
				..controls(43.5,-40.1)and(42.8,-39.6)..(42.7,-39.4)
				..controls(42.4,-38.7)and(43,-38.1)..(44.5,-37.6)..controls(44.9,-37.5)and(46,-37)..(46.9,-36.6)
				..controls(47.9,-36.1)and(48.8,-35.8)..(49.6,-35.6)
				..controls(50.5,-35.4)and(51.3,-35.2)..(52.4,-34.6)
				..controls(53.3,-34.1)and(54.5,-33.7)..(54.9,-33.5)
				..controls(55.4,-33.4)and(56.4,-33)..(57.2,-32.6)
				..controls(58.6,-32)and(60.1,-31.5)..(61.6,-31.2)
				..controls(62.3,-31.1)and(62.7,-31)..(63.2,-30.7)
				..controls(64.1,-30)and(64.3,-29.9)..(65.3,-29.7)
				..controls(65.7,-29.6)and(66.5,-29.2)..(67.1,-28.9)
				..controls(67.6,-28.6)and(68.8,-28.2)..(69.6,-27.9)
				..controls(70.5,-27.7)and(71.6,-27.3)..(72.1,-27)
				..controls(72.6,-26.7)and(73.8,-26.2)..(74.8,-25.9)
				..controls(75.7,-25.5)and(76.7,-25.1)..(77.1,-25)
				..controls(77.4,-24.8)and(78.5,-24.4)..(79.6,-24.1)
				..controls(80.6,-23.8)and(82,-23.2)..(82.8,-22.9)
				..controls(83.5,-22.5)and(84.8,-22)..(85.8,-21.8)
				..controls(86.7,-21.5)and(87.8,-21.1)..(88.1,-20.9)
				..controls(88.4,-20.7)and(89.3,-20.3)..(90.1,-20.1)
				..controls(90.8,-19.8)and(91.9,-19.4)..(92.4,-19.1)
				..controls(93.5,-18.5)and(94.4,-18.2)..(95.4,-18.1)
				..controls(95.7,-18)and(96.4,-17.8)..(96.9,-17.5)
				..controls(97.5,-17.3)and(98.4,-16.9)..(99,-16.6)
				..controls(99.7,-16.3)and(100.9,-15.9)..(101.9,-15.6)
				..controls(102.9,-15.3)and(104.4,-14.7)..(105.3,-14.3)
				..controls(108.1,-13.2)and(108.5,-13.1)..(110.1,-13.1)
				..controls(111.4,-13.1)and(111.5,-13.1)..(112.4,-13.5)
				..controls(113.2,-14)and(115,-14.8)..(116.6,-15.4)
				..controls(116.9,-15.5)and(117.9,-15.8)..(118.8,-16.1)
				..controls(120.9,-16.6)and(121.5,-16.9)..(122.6,-17.9)
				..controls(123,-18.3)and(123.7,-18.8)..(123.9,-18.9)
				..controls(124.6,-19.4)and(126,-19.8)..(127.9,-20.2)
				..controls(129.6,-20.5)and(129.8,-20.6)..(131.1,-21.4)
				..controls(131.5,-21.7)and(132.4,-22.1)..(133.1,-22.3)
				..controls(136.1,-23.2)and(138.3,-24.1)..(140.2,-25)
				..controls(141,-25.5)and(141.5,-25.6)..(142.5,-25.8)
				..controls(145.1,-26.3)and(146,-26.6)..(149.4,-28.1)
				..controls(151.3,-28.9)and(153.5,-29.8)..(154.3,-30.1)
				..controls(155.1,-30.4)and(156.2,-30.8)..(156.6,-31)
				..controls(157.1,-31.2)and(158.2,-31.6)..(159.1,-31.8)
				..controls(160.6,-32.2)and(160.9,-32.3)..(161.7,-32.9)
				..controls(162.4,-33.3)and(162.9,-33.6)..(163.7,-33.8)
				..controls(164.3,-34)and(164.8,-34.1)..(164.9,-34.1)
				..controls(165,-34.1)and(165.6,-34.4)..(166.4,-34.8)
				..controls(167.3,-35.2)and(168.1,-35.4)..(169,-35.6)
				..controls(170,-35.8)and(170.5,-36)..(171.7,-36.6)
				..controls(172.4,-37)and(173.3,-37.4)..(173.6,-37.4)
				..controls(175.6,-37.7)and(176.2,-38)..(176.7,-39.1)
				..controls(176.9,-39.6)and(176.9,-39.7)..(176.7,-40.4)
				..controls(176.4,-41.6)and(175.8,-42.7)..(175.2,-43.4)
				..controls(174.5,-44.1)and(172.6,-45.5)..(171.7,-45.8)
				..controls(170.6,-46.2)and(169.1,-46)..(168.2,-45.4)
				..controls(167.9,-45.2)and(167.3,-44.8)..(166.9,-44.4)
				..controls(166.4,-44.1)and(165.8,-43.6)..(165.4,-43.2)
				..controls(165.1,-42.9)and(164.4,-42.3)..(163.9,-41.9)
				..controls(163.3,-41.5)and(162.7,-40.9)..(162.4,-40.5)
				..controls(161.4,-39.2)and(159.9,-38)..(158.8,-37.8)
				..controls(158.4,-37.8)and(158.3,-37.7)..(158.2,-37.2)
				..controls(158,-36.8)and(157.4,-36.1)..(157.1,-36.1)
				..controls(157,-36.1)and(156.7,-36.1)..(156.5,-36)
				..controls(156.1,-35.9)and(155.9,-35.7)..(155.7,-35.2)
				..controls(155.3,-34.6)and(155.1,-34.4)..(154,-34.1)
				..controls(153.6,-34)and(153.2,-33.7)..(152.9,-33.5)
				..controls(152.4,-32.9)and(151.6,-32.4)..(151.3,-32.4)
				..controls(150.9,-32.4)and(150.4,-32.6)..(150.4,-32.9)
				..controls(150.4,-33.2)and(150.8,-33.4)..(151.4,-33.6)
				..controls(152.1,-33.7)and(152.8,-34.1)..(153.5,-34.9)
				..controls(154.1,-35.5)and(154.3,-35.6)..(155,-35.8)
				..controls(155.7,-36)and(155.7,-36)..(156,-36.6)
				..controls(156.3,-37.3)and(156.7,-37.7)..(157.5,-37.8)
				..controls(158.1,-37.9)and(158.2,-38)..(158.4,-38.8)
				..controls(158.5,-39.4)and(158.6,-39.6)..(159.6,-40.4)
				..controls(160.2,-41)and(160.9,-41.6)..(161.3,-41.8)
				..controls(161.7,-42)and(161.9,-42.3)..(162.2,-42.8)
				..controls(162.4,-43.2)and(162.8,-43.8)..(163.1,-44.1)
				..controls(163.9,-44.8)and(165.1,-46.7)..(165.7,-48)
				..controls(166.2,-49.3)and(166.4,-49.7)..(167.2,-50.5)
				..controls(167.5,-50.8)and(167.9,-51.4)..(168.1,-51.8)
				..controls(168.3,-52.2)and(168.8,-53)..(169.1,-53.6)
				..controls(169.6,-54.3)and(169.8,-54.9)..(169.9,-55.6)
				..controls(170.4,-57.8)and(170.5,-58.1)..(171.2,-59.5)
				..controls(171.6,-60.3)and(172,-61.5)..(172.2,-62.1)
				..controls(172.3,-62.8)and(172.7,-63.7)..(172.9,-64.1)
				..controls(173.2,-64.6)and(173.4,-65.3)..(173.4,-65.6)
				..controls(173.8,-68.8)and(174.2,-70.6)..(174.9,-72.6)
				..controls(175.5,-74.5)and(175.9,-77.1)..(176.1,-80.2)
				..controls(176.3,-83.4)and(175.5,-84.4)..(173.7,-83.4)--cycle;
				\definecolor{ezs}{RGB}{247,179,26}\fill[ezs](-5.7,-156.6)--(-3.4,-157.5)--(61.9,12.7)--(59.7,13.6)--cycle;
				\definecolor{ezs}{RGB}{247,179,22}\fill[ezs](-3.8,-157.3)--(-1.5,-158.2)--(63.9,12)--(61.6,12.9)--cycle;
				\definecolor{ezs}{RGB}{248,179,19}\fill[ezs](-1.8,-158.1)--(0.4,-158.9)--(65.8,11.2)--(63.5,12.1)--cycle;
				\definecolor{ezs}{RGB}{248,178,16}\fill[ezs](0.1,-158.8)--(2.4,-159.7)--(67.7,10.5)--(65.4,11.4)--cycle;
				\definecolor{ezs}{RGB}{249,178,15}\fill[ezs](2,-159.5)--(4.3,-160.4)--(69.6,9.8)--(67.3,10.7)--cycle;
				\definecolor{ezs}{RGB}{251,177,17}\fill[ezs](3.9,-160.3)--(6.2,-161.1)--(71.5,9)--(69.2,9.9)--cycle;
				\definecolor{ezs}{RGB}{252,179,18}\fill[ezs](5.8,-161)--(8.1,-161.9)--(73.4,8.3)--(71.1,9.2)--cycle;
				\definecolor{ezs}{RGB}{252,182,17}\fill[ezs](7.7,-161.7)--(10,-162.6)--(75.3,7.6)--(73,8.5)--cycle;
				\definecolor{ezs}{RGB}{252,185,17}\fill[ezs](9.6,-162.5)--(11.9,-163.4)--(77.2,6.8)--(75,7.7)--cycle;
				\definecolor{ezs}{RGB}{251,189,18}\fill[ezs](11.5,-163.2)--(13.8,-164.1)--(79.2,6.1)--(76.9,7)--cycle;
				\definecolor{ezs}{RGB}{251,191,22}\fill[ezs](13.4,-163.9)--(15.7,-164.8)--(81.1,5.4)--(78.8,6.2)--cycle;
				\definecolor{ezs}{RGB}{252,193,25}\fill[ezs](15.4,-164.7)--(17.6,-165.6)--(83,4.6)--(80.7,5.5)--cycle;
				\definecolor{ezs}{RGB}{253,195,27}\fill[ezs](17.3,-165.4)--(19.6,-166.3)--(84.9,3.9)--(82.6,4.8)--cycle;
				\definecolor{ezs}{RGB}{253,196,31}\fill[ezs](19.2,-166.1)--(21.5,-167)--(86.8,3.2)--(84.5,4)--cycle;
				\definecolor{ezs}{RGB}{254,199,35}\fill[ezs](21.1,-166.9)--(23.4,-167.8)--(88.7,2.4)--(86.4,3.3)--cycle;
				\definecolor{ezs}{RGB}{253,201,38}\fill[ezs](23,-167.6)--(25.3,-168.5)--(90.6,1.7)--(88.3,2.6)--cycle;
				\definecolor{ezs}{RGB}{253,203,40}\fill[ezs](24.9,-168.3)--(27.2,-169.2)--(92.5,1)--(90.3,1.8)--cycle;
				\definecolor{ezs}{RGB}{253,205,42}\fill[ezs](26.8,-169.1)--(29.1,-170)--(94.4,0.2)--(92.2,1.1)--cycle;
				\definecolor{ezs}{RGB}{253,207,45}\fill[ezs](28.7,-169.8)--(31,-170.7)--(96.4,-0.5)--(94.1,0.4)--cycle;
				\definecolor{ezs}{RGB}{253,210,49}\fill[ezs](30.7,-170.5)--(32.9,-171.4)--(98.3,-1.2)--(96,-0.4)--cycle;
				\definecolor{ezs}{RGB}{253,212,54}\fill[ezs](32.6,-171.3)--(34.9,-172.2)--(100.2,-2)--(97.9,-1.1)--cycle;
				\definecolor{ezs}{RGB}{253,214,59}\fill[ezs](34.5,-172)--(36.8,-172.9)--(102.1,-2.7)--(99.8,-1.8)--cycle;
				\definecolor{ezs}{RGB}{253,216,63}\fill[ezs](36.4,-172.7)--(38.7,-173.6)--(104,-3.4)--(101.7,-2.6)--cycle;
				\definecolor{ezs}{RGB}{253,218,65}\fill[ezs](38.3,-173.5)--(40.6,-174.4)--(105.9,-4.2)--(103.6,-3.3)--cycle;
				\definecolor{ezs}{RGB}{254,220,66}\fill[ezs](40.2,-174.2)--(42.5,-175.1)--(107.8,-4.9)--(105.5,-4)--cycle;
				\definecolor{ezs}{RGB}{254,222,68}\fill[ezs](42.1,-175)--(44.4,-175.8)--(109.7,-5.6)--(107.5,-4.8)--cycle;
				\definecolor{ezs}{RGB}{254,224,71}\fill[ezs](44,-175.7)--(46.3,-176.6)--(111.7,-6.4)--(109.4,-5.5)--cycle;
				\definecolor{ezs}{RGB}{254,226,75}\fill[ezs](46,-176.4)--(48.2,-177.3)--(113.6,-7.1)--(111.3,-6.2)--cycle;
				\definecolor{ezs}{RGB}{254,228,79}\fill[ezs](47.9,-177.2)--(50.1,-178)--(115.5,-7.8)--(113.2,-7)--cycle;
				\definecolor{ezs}{RGB}{254,230,82}\fill[ezs](49.8,-177.9)--(52.1,-178.8)--(117.4,-8.6)--(115.1,-7.7)--cycle;
				\definecolor{ezs}{RGB}{254,232,85}\fill[ezs](51.7,-178.6)--(54,-179.5)--(119.3,-9.3)--(117,-8.4)--cycle;
				\definecolor{ezs}{RGB}{254,234,88}\fill[ezs](53.6,-179.4)--(55.9,-180.2)--(121.2,-10)--(118.9,-9.2)--cycle;
				\definecolor{ezs}{RGB}{254,237,91}\fill[ezs](55.5,-180.1)--(57.8,-181)--(123.1,-10.8)--(120.8,-9.9)--cycle;
				\definecolor{ezs}{RGB}{254,239,94}\fill[ezs](57.4,-180.8)--(59.7,-181.7)--(125,-11.5)--(122.8,-10.6)--cycle;
				\definecolor{ezs}{RGB}{254,241,97}\fill[ezs](59.3,-181.6)--(61.6,-182.4)--(127,-12.2)--(124.7,-11.4)--cycle;
				\definecolor{ezs}{RGB}{254,243,101}\fill[ezs](61.2,-182.3)--(63.5,-183.2)--(128.9,-13)--(126.6,-12.1)--cycle;
				\definecolor{ezs}{RGB}{254,245,103}\fill[ezs](63.2,-183)--(65.4,-183.9)--(130.8,-13.7)--(128.5,-12.8)--cycle;
				\definecolor{ezs}{RGB}{254,245,105}\fill[ezs](65.1,-183.8)--(67.4,-184.6)--(132.7,-14.4)--(130.4,-13.6)--cycle;
				\definecolor{ezs}{RGB}{254,245,105}\fill[ezs](67,-184.5)--(69.3,-185.4)--(134.6,-15.2)--(132.3,-14.3)--cycle;
				\definecolor{ezs}{RGB}{254,244,103}\fill[ezs](68.9,-185.2)--(71.2,-186.1)--(136.5,-15.9)--(134.2,-15)--cycle;
				\definecolor{ezs}{RGB}{253,243,101}\fill[ezs](70.8,-186)--(73.1,-186.8)--(138.4,-16.6)--(136.1,-15.8)--cycle;
				\definecolor{ezs}{RGB}{253,241,98}\fill[ezs](72.7,-186.7)--(75,-187.6)--(140.3,-17.4)--(138,-16.5)--cycle;
				\definecolor{ezs}{RGB}{253,239,95}\fill[ezs](74.6,-187.4)--(76.9,-188.3)--(142.2,-18.1)--(140,-17.2)--cycle;
				\definecolor{ezs}{RGB}{253,236,92}\fill[ezs](76.5,-188.2)--(78.8,-189)--(144.2,-18.8)--(141.9,-18)--cycle;
				\definecolor{ezs}{RGB}{253,234,88}\fill[ezs](78.5,-188.9)--(80.7,-189.8)--(146.1,-19.6)--(143.8,-18.7)--cycle;
				\definecolor{ezs}{RGB}{253,232,85}\fill[ezs](80.4,-189.6)--(82.7,-190.5)--(148,-20.3)--(145.7,-19.4)--cycle;
				\definecolor{ezs}{RGB}{253,230,82}\fill[ezs](82.3,-190.4)--(84.6,-191.2)--(149.9,-21)--(147.6,-20.2)--cycle;
				\definecolor{ezs}{RGB}{253,228,78}\fill[ezs](84.2,-191.1)--(86.5,-192)--(151.8,-21.8)--(149.5,-20.9)--cycle;
				\definecolor{ezs}{RGB}{253,226,74}\fill[ezs](86.1,-191.8)--(88.4,-192.7)--(153.7,-22.5)--(151.4,-21.6)--cycle;
				\definecolor{ezs}{RGB}{253,224,70}\fill[ezs](88,-192.6)--(90.3,-193.4)--(155.6,-23.3)--(153.3,-22.4)--cycle;
				\definecolor{ezs}{RGB}{253,222,67}\fill[ezs](89.9,-193.3)--(92.2,-194.2)--(157.5,-24)--(155.3,-23.1)--cycle;
				\definecolor{ezs}{RGB}{253,220,65}\fill[ezs](91.8,-194)--(94.1,-194.9)--(159.5,-24.7)--(157.2,-23.8)--cycle;
				\definecolor{ezs}{RGB}{253,217,63}\fill[ezs](93.7,-194.8)--(96,-195.6)--(161.4,-25.5)--(159.1,-24.6)--cycle;
				\definecolor{ezs}{RGB}{253,215,60}\fill[ezs](95.7,-195.5)--(97.9,-196.4)--(163.3,-26.2)--(161,-25.3)--cycle;
				\definecolor{ezs}{RGB}{253,213,56}\fill[ezs](97.6,-196.2)--(99.9,-197.1)--(165.2,-26.9)--(162.9,-26)--cycle;
				\definecolor{ezs}{RGB}{252,211,53}\fill[ezs](99.5,-197)--(101.8,-197.8)--(167.1,-27.7)--(164.8,-26.8)--cycle;
				\definecolor{ezs}{RGB}{252,209,49}\fill[ezs](101.4,-197.7)--(103.7,-198.6)--(169,-28.4)--(166.7,-27.5)--cycle;
				\definecolor{ezs}{RGB}{252,207,45}\fill[ezs](103.3,-198.4)--(105.6,-199.3)--(170.9,-29.1)--(168.6,-28.2)--cycle;
				\definecolor{ezs}{RGB}{253,205,42}\fill[ezs](105.2,-199.2)--(107.5,-200)--(172.8,-29.9)--(170.6,-29)--cycle;
				\definecolor{ezs}{RGB}{253,202,40}\fill[ezs](107.1,-199.9)--(109.4,-200.8)--(174.7,-30.6)--(172.5,-29.7)--cycle;
				\definecolor{ezs}{RGB}{253,200,38}\fill[ezs](109,-200.6)--(111.3,-201.5)--(176.7,-31.3)--(174.4,-30.4)--cycle;
				\definecolor{ezs}{RGB}{253,198,35}\fill[ezs](111,-201.4)--(113.2,-202.2)--(178.6,-32.1)--(176.3,-31.2)--cycle;
				\definecolor{ezs}{RGB}{253,196,31}\fill[ezs](112.9,-202.1)--(115.2,-203)--(180.5,-32.8)--(178.2,-31.9)--cycle;
				\definecolor{ezs}{RGB}{253,194,28}\fill[ezs](114.8,-202.8)--(117.1,-203.7)--(182.4,-33.5)--(180.1,-32.6)--cycle;
				\definecolor{ezs}{RGB}{253,192,26}\fill[ezs](116.7,-203.6)--(119,-204.5)--(184.3,-34.3)--(182,-33.4)--cycle;
				\definecolor{ezs}{RGB}{253,189,24}\fill[ezs](118.6,-204.3)--(120.9,-205.2)--(186.2,-35)--(183.9,-34.1)--cycle;
				\definecolor{ezs}{RGB}{253,186,20}\fill[ezs](120.5,-205)--(122.8,-205.9)--(188.1,-35.7)--(185.8,-34.8)--cycle;
				\definecolor{ezs}{RGB}{253,184,17}\fill[ezs](122.4,-205.8)--(124.7,-206.7)--(190,-36.5)--(187.8,-35.6)--cycle;
				\definecolor{ezs}{RGB}{252,183,15}\fill[ezs](124.3,-206.5)--(126.6,-207.4)--(192,-37.2)--(189.7,-36.3)--cycle;
				\definecolor{ezs}{RGB}{252,181,14}\fill[ezs](126.3,-207.2)--(128.5,-208.1)--(193.9,-37.9)--(191.6,-37.1)--cycle;
				\definecolor{ezs}{RGB}{252,179,13}\fill[ezs](128.2,-208)--(130.4,-208.9)--(195.8,-38.7)--(193.5,-37.8)--cycle;
				\definecolor{ezs}{RGB}{252,178,12}\fill[ezs](130.1,-208.7)--(132.4,-209.6)--(197.7,-39.4)--(195.4,-38.5)--cycle;
				\definecolor{ezs}{RGB}{251,178,10}\fill[ezs](132,-209.4)--(134.3,-210.3)--(199.6,-40.1)--(197.3,-39.3)--cycle;
				\definecolor{ezs}{RGB}{251,178,10}\fill[ezs](133.9,-210.2)--(136.2,-211.1)--(201.5,-40.9)--(199.2,-40)--cycle;
				\definecolor{ezs}{RGB}{250,177,11}\fill[ezs](135.8,-210.9)--(138.1,-211.8)--(203.4,-41.6)--(201.1,-40.7)--cycle;
				\definecolor{ezs}{RGB}{250,178,11}\fill[ezs](137.7,-211.6)--(140,-212.5)--(205.3,-42.3)--(203.1,-41.5)--cycle;
				\definecolor{ezs}{RGB}{248,180,10}\fill[ezs](139.6,-212.4)--(141.9,-213.3)--(207.3,-43.1)--(205,-42.2)--cycle;
				\definecolor{ezs}{RGB}{247,181,9}\fill[ezs](141.5,-213.1)--(143.8,-214)--(209.2,-43.8)--(206.9,-42.9)--cycle;
				\definecolor{ezs}{RGB}{246,182,9}\fill[ezs](143.5,-213.8)--(145.7,-214.7)--(211.1,-44.5)--(208.8,-43.7)--cycle;
				\definecolor{ezs}{RGB}{246,182,10}\fill[ezs](145.4,-214.6)--(147.7,-215.5)--(213,-45.3)--(210.7,-44.4)--cycle;
			\end{scope}
			% chữ EDUZY (Times đậm, làm dày nét bằng 5 lớp lệch nhẹ; căn đúng khung chữ logo gốc)
			\foreach \dx/\dy in {-0.7/0,0.7/0,0/-0.3,0/0.3,0/0}
			\node[anchor=center,inner sep=0pt,outer sep=0pt,text=ezText,align=center] at (109.6+\dx,-193.6+\dy)
			{\resizebox{1.01cm}{0.22cm}{\fontencoding{T1}\fontfamily{ptm}\fontseries{b}\selectfont EDUZY}};
		\end{tikzpicture}%
	}%
}
%====================== THIẾT LẬP THÂN ĐỀ ======================
%--- Môi trường tự luận {bt}: ex_test 3.3.4 chưa khai báo sẵn nên tạo ở đây
%    (ex_test đã tự móc \fixshowans và \listenumerate cho bt)
\makeatletter
\@ifundefined{bt}{\newtheorem{bt}{\color{gold1}\namebt}}{}
\makeatother
\renewcommand{\sectionmark}[1]{\markright{\thesection\ #1}}
%--- Khung tiêu đề (logo + tối đa 4 dòng; dòng 3, 4 để trống {} thì bỏ qua)
%  Đề thi:    \khungtieude{THI THỬ THPT QUỐC GIA NĂM 2026}{Môn Toán -- Đề số 21}
%                         {(Đề thi gồm: ...)}{Thời gian làm bài: 90 phút (không kể thời gian phát đề).}
%  Chuyên đề: \khungtieude{CHUYÊN ĐỀ 3 -- DÃY SỐ VÀ CÔNG THỨC TRUY HỒI}
%                         {Bồi dưỡng học sinh giỏi Toán THPT}{...}{}
\newcommand{\khungtieude}[4]{%
	\begin{tcolorbox}[colframe=gold1,colback=gold5!45,boxrule=0.5mm]
		\begin{minipage}[c]{3cm}
			\centerline{\eduzylogo[2.4cm]}
		\end{minipage}\hspace{0.5cm}
		\begin{minipage}[c]{12.5cm}
			\vskip 0.15 cm
			\centerline{\ind{#1}}
			\vskip 0.15 cm
			\centerline{\textbf{#2}}
			\if\relax\detokenize{#3}\relax\else
				\vskip 0.15 cm
				\centerline{\textit{#3}}
			\fi
			\if\relax\detokenize{#4}\relax\else
				\vskip 0.15 cm
				\centerline{\textit{#4}}
			\fi
			\vskip 0.15 cm
		\end{minipage}
	\end{tcolorbox}
}
%====================== BỔ SUNG CHO BÀI HỌC, CHUYÊN ĐỀ ======================
\usepackage[shortlabels]{enumitem}
\usepackage{tabularx}
%--- Danh sách gọn: \begin{enumerate}[gon] ... \end{enumerate}
\SetEnumitemKey{gon}{leftmargin=*,itemsep=1pt,parsep=0pt,topsep=3pt}
%--- Khung lý thuyết: \begin{khung}{} ... \end{khung}; tiêu đề chỉ dùng khi thật ngắn ("Mẹo", "Lỗi thường gặp")
\newtcolorbox{khung}[1]{breakable,enhanced,colframe=gold1,colback=gold5!45,
	boxrule=0.4mm,fonttitle=\bfseries\sffamily,coltitle=white,colbacktitle=gold1,
	attach boxed title to top left={xshift=6mm,yshift=-3mm},
	boxed title style={arc=1mm,boxrule=0.3mm},title={#1}}
%--- Tiêu đề mục nhỏ trong một dạng: \muc{Dấu hiệu}, \muc{Phương pháp}, \muc{Bài tự luyện}
\makeatletter
\newcommand{\muc}[1]{\par\medskip\noindent\indam{#1}\par\nobreak\smallskip\@afterindentfalse\@afterheading}
\makeatother
%--- Ô công thức cần nhớ (căn giữa, viền vàng, vừa khít nội dung)
%     \begin{center}\congthuc{$Q(n)=(n-1)!$}\end{center}
\newtcbox{\congthuc}{on line,colframe=gold1,colback=gold5!45,boxrule=0.5pt,arc=2pt,
	left=8pt,right=8pt,top=4pt,bottom=4pt}
%--- Công tắc lời giải bài tự luyện (ví dụ mẫu luôn dùng \loigiai)
%     \loigiaiTL[{đáp số}]{lời giải đầy đủ}
%     \TLfull -> in lời giải đầy đủ;  \TLdapso -> chỉ in một dòng "Đáp số";  \TLan -> ẩn hoàn toàn
%     (\CheDo là dethi thì mọi lời giải và đáp số đều ẩn)
\makeatletter
\newif\ifTLhien   \TLhientrue
\newif\ifTLdapso  \TLdapsofalse
\newcommand{\TLfull}{\TLhientrue\TLdapsofalse}
\newcommand{\TLdapso}{\TLhientrue\TLdapsotrue}
\newcommand{\TLan}{\TLhienfalse}
\newcommand{\loigiaiTL}[2][]{%
	\ifTLhien
		\ifTLdapso
			\def\@tlarg{#1}%
			\ifx\@tlarg\@empty\else
				\ifdefstring{\kindOPTN}{dethi}{}{\par\noindent{\bfseries\color{gold1}Đáp số:} #1\par}%
			\fi
		\else
			\loigiai{#2}%
		\fi
	\fi}
\makeatother
\TLdapso% mặc định: chỉ in đáp số bài tự luyện
%====================== BẢNG ĐÚNG/SAI (\choiceTFt) ======================
%  Hàng tiêu đề nền nâu đồng chữ trắng (Phát biểu | Đ | S), các hàng tô xen kẽ kem / trắng,
%  đường kẻ vàng kim, nhãn a) b) c) d) màu nâu đồng. Bản có đáp án: cột được chọn có vòng tròn
%  nâu đồng với dấu tích trắng (câu đúng tích ở cột Đúng, câu sai tích ở cột Sai).
%  Tiêu đề ghi đủ "Đúng | Sai": \OPTN{viettat=0} sau \input (cột Đúng/Sai rộng ra, phát biểu hẹp lại).
%  Dùng bảng gốc của ex_test: \def\BangDSgoc{1} trước %============================================================
%  PREAMBLE EDUZY v2 - dùng chung cho đề thi, bài học, chuyên đề (gói ex_test)
%  Khác bản preamble-eduzy.tex ở đầu trang, chân trang: dải nền nâu đồng tràn ngang,
%  logo + Eduzy.com.vn, ô chữ vàng kem cạnh xiên, chân trang "Trang (số)";
%  logo Eduzy chìm giữa mọi trang (tắt: \logochimtat, bật lại: \logochimbat).
%  Cách dùng ở file chính (xem mau-chuan.tex, mau-chuyen-de.tex):
%      \def\CheDo{loigiai}          % loigiai | dethi | solcolor | color | book
%      \def\tenchuong{<dòng chữ trên đầu trang>}
%      \def\thuonghieu{<chữ ở chân trang lẻ>}   % không khai báo thì dùng mặc định
%      \def\tacgia{<chữ ở chân trang chẵn>}      % không khai báo thì dùng mặc định
%      \def\khauhieu{<chữ cạnh logo ở đầu trang>}  % tùy chọn, mặc định Eduzy.com.vn
%      \input{preamble-eduzy-v2}
%      \TLdapso                      % tài liệu: \TLfull | \TLdapso | \TLan
%      \begin{document} ... \end{document}
%  Logo Eduzy vẽ bằng TikZ ngay trong file (\eduzylogo), không cần file ảnh.
%  Biên dịch: pdfLaTeX (máy có gói vietnam) hoặc XeLaTeX.
%============================================================
\documentclass[12pt,a4paper,twoside]{article}%<-- twoside: bat dau trang/chan trang chan-le
\usepackage{iftex}
\usepackage{amsmath,amssymb,tikz,mathrsfs,tkz-tab}
\usepackage{graphicx}
\usepackage{pgfplots}% đồ thị hàm số
\pgfplotsset{compat=1.17}
%\usepackage{mathptmx}
\usepackage{mathpazo}
\ifXeTeX% XeLaTeX: font chữ có đủ dấu tiếng Việt, giữ nguyên font toán mathpazo
	\usepackage[no-math]{fontspec}
	\setmainfont{TeX Gyre Pagella}
	\newfontfamily\EZqag{TeX Gyre Adventor}[NFSSFamily=qag,Ligatures=TeX]
\fi
\usepackage[many]{tcolorbox}
\usetikzlibrary{math,through,calc,intersections,angles,quotes,shapes,shapes.geometric,arrows,patterns,matrix,chains,arrows.meta,decorations.shapes,decorations.fractals,decorations.markings,shadows,decorations.pathreplacing}
\usepackage[top=1.6cm, bottom=1.8cm, left=1.8cm, right=1.8cm]{geometry}% v2: top 1.6cm (bản cũ 1.5cm) để chữ không sát dải đầu trang
\providecommand{\CheDo}{loigiai}% loigiai | dethi | solcolor | color | book
\expandafter\usepackage\expandafter[\CheDo]{ex_test}
% \tenchuong: khai báo ở file chính (trước \input) hoặc \renewcommand sau \input
%============================================================
%  [GỘP TỪ eduzy-de.tex] EDUZY - MÀU NÂU ĐỒNG / VÀNG KIM
%  + ĐẦU TRANG, CHÂN TRANG KIỂU DẢI NGANG (v2), ĐỐI XỨNG TRANG CHẴN / TRANG LẺ.
%  (không cần file eduzy-de.tex riêng nữa)
%  Khối này phải đặt SAU \usepackage[...]{ex_test} và SAU \def\tenchuong{...}
%============================================================
\usepackage{lastpage}
\usepackage{needspace}
%--- DÀN TRANG ĐỀ THI: tách câu (THÊM 05/11/2026; chỉ thêm lệnh mới, KHÔNG đổi hành vi mặc định) ---
%  Hai công tắc, mặc định tắt (= preamble chạy y như cũ). Bật bằng cách đặt TRƯỚC \input{preamble-eduzy-v2}
%  trong file đề (tools/dan_trang/dan_trang.py tự ghi khi đặt \NeedspaceCau):
%    \def\TachCauDS{1}   câu Đúng/Sai được tách: câu dẫn (kèm hình) ở cuối trang, bảng Phát biểu sang trang sau
%                        (câu dẫn và hình không tách nhau, bảng không tách giữa hai trang)
%    \def\TachCauLG{1}   bản lời giải (\CheDo khác dethi): phần đề của câu cao hơn 1/3 trang được phép tách
%  Lệnh dùng trong file đề (giá trị do tools/dan_trang/dan_trang.py đo ở bản dethi):
%    \NeedspaceCau[<cao câu dẫn + hình>]{<cao cả khối câu>}   ngay trên \begin{ex} (câu đầu phần: trên tiêu đề phần)
%        công tắc tắt: đúng bằng \Needspace{<cao cả khối câu>} như cũ;
%        \TachCauDS=1 và có [..] (câu Đúng/Sai được chọn tách): chỉ giữ câu dẫn + hình;
%        \TachCauLG=1 ở bản lời giải: đòi tối đa 1/3 \textheight.
%    \NeedspaceBang{<cao bảng Phát biểu>}%   ngay trước \choiceTFt: \TachCauDS=1 thì giữ trọn bảng trên một trang;
%        công tắc tắt: không làm gì.
\providecommand{\TachCauDS}{0}
\providecommand{\TachCauLG}{0}
\newlength{\EZnsCau}
\newcommand{\NeedspaceCau}[2][]{%
	\setlength{\EZnsCau}{#2}%
	\ifnum\TachCauDS=1 \if\relax\detokenize{#1}\relax\else\setlength{\EZnsCau}{#1}\fi\fi
	\ifnum\TachCauLG=1 \ifdefstring{\kindOPTN}{dethi}{}{\ifdim\EZnsCau>0.3333\textheight\setlength{\EZnsCau}{0.3333\textheight}\fi}\fi
	\Needspace{\EZnsCau}}
\newcommand{\NeedspaceBang}[1]{\ifnum\TachCauDS=1 \Needspace{#1}\fi}

%--- Thang màu: đậm -> nhạt
\definecolor{gold1}{RGB}{126,96,0}    % nâu đồng đậm  #7E6000
\definecolor{gold2}{RGB}{192,144,3}   % vàng kim      #C09003
\definecolor{gold3}{RGB}{255,215,99}  % vàng nhạt     #FFD763
\definecolor{gold4}{RGB}{255,231,153} % vàng rất nhạt #FFE799
\definecolor{gold5}{RGB}{255,243,205} % vàng kem      #FFF3CD
%--- Ghi đè tông xanh/đỏ cũ
\definecolor{myblue}{RGB}{126,96,0}
\colorlet{blue}{gold1}
\colorlet{cyan}{gold2}
\colorlet{teal}{gold2}

%--- Tên chương mặc định (file đề nên tự khai báo lại)
\providecommand{\tenchuong}{}
%--- Chân trang
\providecommand{\thuonghieu}{TOÁN 12 -- Eduzy.com.vn}
\providecommand{\tacgia}{Tô Gia Bảo -- 0353703783}

%--- Chữ cạnh logo ở đầu trang
\providecommand{\khauhieu}{Eduzy.com.vn}

%====================== THÔNG SỐ ĐẦU TRANG, CHÂN TRANG (v2) ======================
%  Dải nền đậm tràn ngang trang; ô chữ vàng kem có cạnh xiên, kèm 2 vạch xiên chuyển màu.
%  Đầu trang lẻ : logo + \khauhieu (trái)   | ô chữ \tenchuong (phải)
%  Chân trang lẻ: ô chữ \thuonghieu (trái)  | Trang (số) (phải)
%  Trang chẵn soi gương trang lẻ; chân trang chẵn ghi \tacgia.
%--- Màu: chỉ dùng thang vàng kim Eduzy
\colorlet{ezNen}{gold1}% dải nền đậm
\colorlet{ezVach1}{gold3}% vạch xiên sát dải nền (đậm hơn)
\colorlet{ezVach2}{gold4}% vạch xiên sát ô chữ (nhạt hơn)
\colorlet{ezO}{gold5}% ô chữ vàng kem
%--- Kích thước
\newlength{\EZdai}\setlength{\EZdai}{1.1cm}% chiều cao dải nền
\newlength{\EZo}\setlength{\EZo}{0.68cm}% chiều cao ô chữ
\newlength{\EZxien}\setlength{\EZxien}{0.55cm}% độ lệch ngang của cạnh xiên
\newlength{\EZvach}\setlength{\EZvach}{0.75cm}% bề rộng mỗi vạch xiên
\newlength{\EZkhe}\setlength{\EZkhe}{0.12cm}% khe trắng giữa vạch và ô chữ
\newlength{\EZdem}\setlength{\EZdem}{0.4cm}% khoảng đệm từ chữ tới cạnh xiên
\newlength{\EZle}\setlength{\EZle}{1.8cm}% lề trái/phải của chữ (bằng lề thân bài)
\newlength{\EZrongmin}\setlength{\EZrongmin}{4.5cm}% ô chữ rộng tối thiểu
\newlength{\EZrong}% bề rộng chữ trong ô (tính tự động)
%--- Font
\newcommand{\EZfontO}{\fontsize{10.5pt}{12pt}\selectfont\bfseries}% chữ trong ô
\newcommand{\EZfontKH}{\fontfamily{qag}\fontsize{11pt}{12pt}\selectfont\bfseries}% \khauhieu cạnh logo
\newcommand{\EZfontTrang}{\fontsize{11.5pt}{12pt}\selectfont\bfseries}% chữ "Trang" và số trang
%--- Đo bề rộng chữ trong ô
\newcommand{\EZdorong}[1]{%
	\settowidth{\EZrong}{\EZfontO #1}%
	\ifdim\EZrong<\EZrongmin\setlength{\EZrong}{\EZrongmin}\fi}
%--- Logo + chữ trắng \khauhieu (dựng sẵn một lần khi vào \begin{document})
\newsavebox{\EZlogobox}
\AtBeginDocument{\sbox{\EZlogobox}{%
		\raisebox{-0.5\height}{\eduzylogo[1cm]}\hspace{0.15cm}%
		\raisebox{-0.5\height}{\color{white}\EZfontKH\khauhieu}}}
%--- Hai vạch xiên + khe trắng. #1: điểm gốc ở mép dưới, #2: hướng (1 sang phải, -1 sang trái),
%    #3: hướng nghiêng (1 kiểu "/", -1 kiểu "\")
\newcommand{\EZvachxien}[3]{%
	\fill[white] (#1) --++({(#3)*\EZxien},\EZo) --++({(#2)*\EZkhe},0) --++({-(#3)*\EZxien},-\EZo) --cycle;
	\fill[ezVach2] ([xshift=(#2)*\EZkhe]#1) --++({(#3)*\EZxien},\EZo) --++({(#2)*\EZvach},0) --++({-(#3)*\EZxien},-\EZo) --cycle;
	\fill[ezVach1] ([xshift=(#2)*(\EZkhe+\EZvach)]#1) --++({(#3)*\EZxien},\EZo) --++({(#2)*\EZvach},0) --++({-(#3)*\EZxien},-\EZo) --cycle;
	\draw[white,line width=0.6pt] ([xshift=(#2)*(\EZkhe+\EZvach)]#1) --++({(#3)*\EZxien},\EZo);
}
%--- Chữ trong ô (căn giữa theo chiều cao ô, không lệch vì dấu tiếng Việt)
\tikzset{ezchu/.style={inner sep=0pt,text=ezNen,text height=0.72em,text depth=0.22em}}

%====================== LOGO CHÌM GIỮA TRANG ======================
%  Vẽ cùng đầu trang nên nằm dưới chữ của thân bài, có trên mọi trang.
%  Tắt: \logochimtat (từ trang đó trở đi), bật lại: \logochimbat.
%  Đổi cỡ, độ đậm: khai báo trước \input, ví dụ \def\EZchimcao{9cm} \def\EZchimdo{0.1}
\providecommand{\EZchimcao}{10cm}% chiều cao logo chìm
\providecommand{\EZchimdo}{0.12}% độ đậm (0 = trong suốt, 1 = rõ hẳn)
\newif\ifEZchim \EZchimtrue
\newcommand{\logochimtat}{\EZchimfalse}
\newcommand{\logochimbat}{\EZchimtrue}
%  Dựng sẵn logo mờ một lần khi vào \begin{document}; transparency group làm mờ cả logo
%  một lượt nên không lộ vệt ở chỗ các dải màu chồng nhau.
\newsavebox{\EZchimtho}
\newsavebox{\EZchimbox}
\AtBeginDocument{%
	\sbox{\EZchimtho}{\eduzylogo[\EZchimcao]}%
	\sbox{\EZchimbox}{\begin{tikzpicture}
			\begin{scope}[transparency group,opacity=\EZchimdo]
				\node[inner sep=0pt] {\usebox{\EZchimtho}};
			\end{scope}
	\end{tikzpicture}}}
\newcommand{\EZvechim}{%
	\ifEZchim
		\node[inner sep=0pt] at (current page.center) {\usebox{\EZchimbox}};
	\fi}

%====================== ĐẦU TRANG - TRANG LẺ ======================
%  (logo bên trái, ô chữ \tenchuong bên phải)
\newcommand{\myfancyheadOdd}{%
	\EZdorong{\tenchuong}%
	\begin{tikzpicture}[remember picture,overlay]
		\coordinate (T) at (current page.north west);
		\coordinate (P) at (current page.north east);
		%--- logo chìm giữa trang
		\EZvechim
		%--- dải nền
		\fill[ezNen] (T) rectangle ([yshift=-\EZdai]P);
		%--- ô chữ: X là góc dưới bên trái, cạnh trái xiên kiểu "/"
		\coordinate (X) at ([xshift=-\EZle-\EZrong-\EZdem-\EZxien,yshift=-\EZdai]P);
		\fill[ezO] (X) --++(\EZxien,\EZo) -- ([yshift=-\EZdai+\EZo]P) -- ([yshift=-\EZdai]P) --cycle;
		\EZvachxien{X}{-1}{1}
		\node[ezchu,anchor=east] at ([xshift=-\EZle,yshift=-\EZdai+0.5*\EZo]P) {\EZfontO\tenchuong};
		%--- logo + tên
		\node[inner sep=0pt,anchor=west] at ([xshift=\EZle,yshift=-0.5*\EZdai]T) {\usebox{\EZlogobox}};
	\end{tikzpicture}%
}

%====================== ĐẦU TRANG - TRANG CHẴN =====================
%  (soi gương: ô chữ \tenchuong bên trái, logo bên phải)
\newcommand{\myfancyheadEven}{%
	\EZdorong{\tenchuong}%
	\begin{tikzpicture}[remember picture,overlay]
		\coordinate (T) at (current page.north west);
		\coordinate (P) at (current page.north east);
		%--- logo chìm giữa trang
		\EZvechim
		%--- dải nền
		\fill[ezNen] (T) rectangle ([yshift=-\EZdai]P);
		%--- ô chữ: X là góc dưới bên phải, cạnh phải xiên kiểu "\"
		\coordinate (X) at ([xshift=\EZle+\EZrong+\EZdem+\EZxien,yshift=-\EZdai]T);
		\fill[ezO] (X) --++(-\EZxien,\EZo) -- ([yshift=-\EZdai+\EZo]T) -- ([yshift=-\EZdai]T) --cycle;
		\EZvachxien{X}{1}{-1}
		\node[ezchu,anchor=west] at ([xshift=\EZle,yshift=-\EZdai+0.5*\EZo]T) {\EZfontO\tenchuong};
		%--- logo + tên
		\node[inner sep=0pt,anchor=east] at ([xshift=-\EZle,yshift=-0.5*\EZdai]P) {\usebox{\EZlogobox}};
	\end{tikzpicture}%
}

%====================== CHÂN TRANG - TRANG LẺ ======================
%  (ô chữ \thuonghieu bên trái, "Trang (số)" bên phải)
\newcommand{\myfancyfootOdd}{%
	\EZdorong{\thuonghieu}%
	\begin{tikzpicture}[remember picture,overlay]
		\coordinate (T) at (current page.south west);
		\coordinate (P) at (current page.south east);
		%--- dải nền
		\fill[ezNen] (T) rectangle ([yshift=\EZdai]P);
		%--- ô chữ sát mép trên dải: X là góc dưới bên phải, cạnh phải xiên kiểu "/"
		\coordinate (X) at ([xshift=\EZle+\EZrong+\EZdem,yshift=\EZdai-\EZo]T);
		\fill[ezO] ([yshift=\EZdai-\EZo]T) -- (X) --++(\EZxien,\EZo) -- ([yshift=\EZdai]T) --cycle;
		\EZvachxien{X}{1}{1}
		\node[ezchu,anchor=west] at ([xshift=\EZle,yshift=\EZdai-0.5*\EZo]T) {\EZfontO\thuonghieu};
		%--- số trang
		\node[circle,draw=white,line width=0.9pt,text=white,inner sep=0pt,minimum size=0.62cm,anchor=east]
		(so) at ([xshift=-\EZle,yshift=0.5*\EZdai]P) {\EZfontTrang\thepage};
		\node[inner sep=0pt,text=white,anchor=base east] (tr) at ([xshift=-0.15cm]so.base west) {\EZfontTrang Trang};
	\end{tikzpicture}%
}

%====================== CHÂN TRANG - TRANG CHẴN ====================
%  (soi gương: "Trang (số)" bên trái, ô chữ \tacgia bên phải)
\newcommand{\myfancyfootEven}{%
	\EZdorong{\tacgia}%
	\begin{tikzpicture}[remember picture,overlay]
		\coordinate (T) at (current page.south west);
		\coordinate (P) at (current page.south east);
		%--- dải nền
		\fill[ezNen] (T) rectangle ([yshift=\EZdai]P);
		%--- ô chữ sát mép trên dải: X là góc dưới bên trái, cạnh trái xiên kiểu "\"
		\coordinate (X) at ([xshift=-\EZle-\EZrong-\EZdem,yshift=\EZdai-\EZo]P);
		\fill[ezO] ([yshift=\EZdai-\EZo]P) -- (X) --++(-\EZxien,\EZo) -- ([yshift=\EZdai]P) --cycle;
		\EZvachxien{X}{-1}{-1}
		\node[ezchu,anchor=east] at ([xshift=-\EZle,yshift=\EZdai-0.5*\EZo]P) {\EZfontO\tacgia};
		%--- số trang
		%    (vòng tròn ẩn so0 nằm giữa dải để lấy đường chân chữ, giống trang lẻ)
		\node[circle,inner sep=0pt,minimum size=0.62cm] (so0) at ([yshift=0.5*\EZdai]T) {\EZfontTrang\phantom{\thepage}};
		\node[inner sep=0pt,text=white,anchor=base west] (tr) at ($(so0.base -| T)+(\EZle,0)$) {\EZfontTrang Trang};
		\node[circle,draw=white,line width=0.9pt,text=white,inner sep=0pt,minimum size=0.62cm,anchor=base west]
		(so) at ([xshift=0.15cm]tr.base east) {\EZfontTrang\thepage};
	\end{tikzpicture}%
}

%--- Giữ tương thích với các file cũ đang gọi \myfancyhead, \myfancyfoot
\newcommand{\myfancyhead}{\myfancyheadOdd}
\newcommand{\myfancyfoot}{\myfancyfootOdd}

%======================= GẮN VÀO TRANG =============================
\pagestyle{fancy}
\fancyhf{}
\fancyhead[LO]{\myfancyheadOdd}
\fancyhead[RE]{\myfancyheadEven}
\fancyfoot[LO]{\myfancyfootOdd}
\fancyfoot[RE]{\myfancyfootEven}
\renewcommand{\headrulewidth}{0pt}
\renewcommand{\footrulewidth}{0pt}
\fancypagestyle{plain}{%
	\fancyhf{}
	\fancyhead[LO]{\myfancyheadOdd}
	\fancyhead[RE]{\myfancyheadEven}
	\fancyfoot[LO]{\myfancyfootOdd}
	\fancyfoot[RE]{\myfancyfootEven}
	\renewcommand{\headrulewidth}{0pt}
	\renewcommand{\footrulewidth}{0pt}
}
%=================== HẾT PHẦN GỘP TỪ eduzy-de.tex ===================
\newcommand{\indam}[1]{
	{\fontsize{12pt}{9pt}\sffamily \selectfont\bfseries\color{gold1} #1}
}
\newcommand{\ind}[1]{
	{\fontsize{14pt}{9pt}\fontfamily{qag} \selectfont\bfseries\color{gold1} #1}
}
\renewcommand{\FalseEX}{\stepcounter{dapan}\textbf{\selectfont\color{gold1}\Alph{dapan}.}\ignorespaces}
\renewtheorem{ex}{\color{gold1}Câu}
\renewcommand{\baselinestretch}{1.1}
\newcommand{\hoac}[1]{
	\left[\begin{aligned}#1\end{aligned}\right.}
\newcommand{\heva}[1]{
	\left\{\begin{aligned}#1\end{aligned}\right.}
\usepackage{esvect}
\def\vec{\vv}
\def\vt{\vv}
\def\overrightarrow{\vv}
%====================== LOGO EDUZY VẼ BẰNG TIKZ ======================
%  Đường viền dò từ logo gốc; dải màu vàng kim đo trực tiếp trên ảnh gốc
%  (nghiêng 21 độ, tô bằng 80 dải hẹp nên hiển thị giống nhau trên mọi trình xem PDF).
%  Dùng: \eduzylogo          -> cao 2.4cm (mặc định)
%        \eduzylogo[1.5cm]   -> cao 1.5cm
\definecolor{ezText}{RGB}{245,200,46}
\newcommand{\eduzylogo}[1][2.4cm]{%
	\resizebox{!}{#1}{%
		\begin{tikzpicture}[x=0.01cm,y=0.01cm]
			\path[use as bounding box] (0,0) rectangle (224,-224);
			% mũ tốt nghiệp + vương miện
			\begin{scope}
				\clip
				(119.4,-176.1)..controls(118.9,-176.1)and(116.8,-176)..(114.6,-176)
				..controls(101.4,-176)and(85.8,-175.9)..(83.7,-175.7)
				..controls(81.8,-175.6)and(81,-175.6)..(80.2,-175.7)
				..controls(78.8,-175.9)and(77.7,-175.9)..(76.2,-175.7)
				..controls(74.9,-175.5)and(70.8,-175.5)..(64.3,-175.7)
				..controls(59.1,-175.9)and(58.3,-175.7)..(58.6,-173.7)
				..controls(58.6,-172.8)and(59.6,-169.6)..(60.3,-168)
				..controls(60.8,-166.9)and(61.4,-165.8)..(62.3,-165)
				..controls(63.1,-164.2)and(64.3,-163.4)..(65,-163.2)
				..controls(65.5,-163.1)and(72.3,-163)..(76.6,-163.1)
				..controls(77.9,-163.1)and(80.1,-163.1)..(81.4,-163)
				..controls(82.8,-162.9)and(86.3,-162.9)..(89.2,-162.9)
				..controls(93,-162.9)and(94.9,-162.8)..(96,-162.7)
				..controls(97.2,-162.5)and(101,-162.5)..(114,-162.4)
				..controls(123,-162.3)and(134.9,-162.2)..(140.2,-162.1)
				..controls(150.1,-162)and(152.2,-162.1)..(153.6,-162.4)
				..controls(154.4,-162.6)and(157.1,-163.8)..(158.1,-164.4)
				..controls(158.8,-164.8)and(159.7,-165.7)..(159.9,-166.3)
				..controls(160,-166.5)and(160.2,-167.1)..(160.2,-167.6)
				..controls(160.3,-168)and(160.6,-169.4)..(160.9,-170.6)
				..controls(161.2,-171.7)and(161.5,-173.1)..(161.6,-173.6)
				..controls(161.8,-174.4)and(161.8,-174.5)..(161.5,-174.9)
				..controls(161.2,-175.6)and(160.8,-175.9)..(159.9,-176)
				..controls(159.2,-176.1)and(154.1,-176.2)..(151.5,-176.2)
				..controls(148.1,-176.1)and(122.6,-176.1)..(121.6,-176.2)
				..controls(120.8,-176.2)and(119.8,-176.2)..(119.4,-176.1)--cycle
				(130.4,-160)..controls(129,-160)and(119.8,-160)..(110.1,-159.9)
				..controls(100.3,-159.8)and(88.5,-159.7)..(83.9,-159.6)
				..controls(77.8,-159.5)and(75.3,-159.5)..(74.5,-159.6)
				..controls(73.9,-159.7)and(72.5,-159.8)..(71.3,-159.8)
				..controls(69.6,-159.8)and(69,-159.8)..(68.5,-159.6)
				..controls(67.6,-159.3)and(66.8,-158.6)..(66.2,-157.6)
				..controls(66,-157.3)and(65.6,-156.5)..(65.3,-156.1)
				..controls(64.6,-155.1)and(64.5,-154.9)..(63.7,-152.9)
				..controls(63.2,-151.6)and(62.9,-150.9)..(61.8,-149.6)
				..controls(61.5,-149.2)and(61.1,-148.5)..(60.9,-147.9)
				..controls(60.7,-147.4)and(60.3,-146.6)..(60.1,-146.1)
				..controls(59.9,-145.5)and(59.5,-144.9)..(59,-144.3)
				..controls(58.5,-143.7)and(58.2,-143.2)..(57.9,-142.5)
				..controls(57.7,-141.9)and(57.3,-141.1)..(57,-140.7)
				..controls(56.6,-140.3)and(56.1,-139.3)..(55.8,-138.5)
				..controls(55.5,-137.8)and(55.1,-136.8)..(54.9,-136.5)
				..controls(54.7,-136.1)and(54.3,-135.2)..(53.9,-134.5)
				..controls(53.6,-133.8)and(53,-132.7)..(52.6,-132)
				..controls(51.6,-130.5)and(51.4,-130.2)..(51.1,-129)
				..controls(50.6,-127.5)and(50.6,-127.4)..(49.7,-126.3)
				..controls(49.2,-125.7)and(48.6,-124.8)..(48.3,-124.2)
				..controls(48,-123.6)and(47.5,-122.7)..(47.3,-122.2)
				..controls(47,-121.8)and(46.5,-120.7)..(46.1,-119.8)
				..controls(45.8,-118.9)and(45.2,-117.8)..(44.9,-117.2)
				..controls(44.2,-116.1)and(43.3,-114.4)..(42.3,-112.2)
				..controls(41.8,-111.3)and(41.3,-110.4)..(41.1,-110.1)
				..controls(40.5,-109.4)and(40.3,-109)..(40,-107.8)
				..controls(39.8,-107)and(39.6,-106.5)..(39.2,-106)
				..controls(38.9,-105.6)and(38.4,-104.7)..(38,-104)
				..controls(37.7,-103.4)and(37.2,-102.5)..(37,-102.1)
				..controls(36.7,-101.7)and(36.3,-100.8)..(36,-100.1)
				..controls(35.6,-98.9)and(35.3,-98.4)..(34.6,-97.6)
				..controls(34.2,-97.1)and(34.2,-96.3)..(34.6,-96)
				..controls(35,-95.7)and(35.5,-95.9)..(36.2,-96.7)
				..controls(36.5,-97)and(37.2,-97.5)..(37.6,-97.7)..controls(38,-98)and(38.6,-98.4)..(39.1,-98.9)
				..controls(39.6,-99.4)and(40.3,-100)..(40.6,-100.3)
				..controls(41.3,-100.8)and(42.4,-101.8)..(43.7,-102.9)
				..controls(44,-103.2)and(44.6,-103.7)..(45,-103.9)
				..controls(45.4,-104.2)and(45.9,-104.6)..(46.1,-104.8)
				..controls(47.1,-105.9)and(48.8,-107.1)..(50.5,-107.9)
				..controls(51.8,-108.6)and(52.3,-109)..(53.4,-110)
				..controls(54.8,-111.2)and(55.6,-111.6)..(57,-112.1)
				..controls(57.5,-112.3)and(58.5,-112.8)..(59.4,-113.3)
				..controls(60.2,-113.7)and(61.3,-114.3)..(61.9,-114.5)
				..controls(62.4,-114.7)and(63.3,-115.1)..(63.8,-115.3)
				..controls(64.4,-115.6)and(65.4,-115.9)..(66.2,-116.1)
				..controls(67,-116.3)and(68.4,-116.6)..(69.2,-116.9)
				..controls(72.3,-117.9)and(72.8,-118)..(78,-118.1)
				..controls(84,-118.3)and(86.1,-118.1)..(88.1,-117.1)
				..controls(89.4,-116.5)and(93,-115)..(93.7,-114.9)
				..controls(94.8,-114.6)and(95.4,-114.2)..(95.9,-113.4)
				..controls(96.2,-113)and(96.7,-112.5)..(97.1,-112.2)
				..controls(98,-111.6)and(101.3,-108.4)..(102.1,-107.2)
				..controls(102.4,-106.8)and(102.9,-106.2)..(103.2,-105.8)
				..controls(103.5,-105.5)and(104,-104.8)..(104.2,-104.2)
				..controls(104.5,-103.6)and(105,-102.8)..(105.3,-102.4)
				..controls(105.6,-102)and(106.1,-101.2)..(106.5,-100.6)
				..controls(107.1,-99.5)and(108.6,-97.5)..(109.5,-96.5)
				..controls(109.9,-96.1)and(110.1,-96)..(110.4,-96)
				..controls(111.1,-96)and(111.6,-96.5)..(111.8,-97.6)
				..controls(112,-98.6)and(112.3,-99.3)..(112.8,-99.9)
				..controls(113.7,-100.9)and(115.2,-103.1)..(115.7,-104.2)
				..controls(116.1,-105)and(116.5,-105.5)..(116.9,-106)
				..controls(117.3,-106.4)and(117.8,-106.9)..(118,-107.3)
				..controls(118.2,-107.6)and(118.6,-108)..(119,-108.3)
				..controls(119.6,-108.8)and(121.1,-110.3)..(121.8,-111.2)
				..controls(122.1,-111.7)and(122.5,-112)..(123.4,-112.4)
				..controls(124.8,-113.2)and(125.8,-113.8)..(127,-114.6)
				..controls(127.5,-114.9)and(128.1,-115.3)..(128.2,-115.4)
				..controls(128.4,-115.5)and(129.3,-115.8)..(130.2,-116)
				..controls(131.1,-116.2)and(132.5,-116.6)..(133.3,-116.9)
				..controls(135.7,-117.9)and(136.1,-117.9)..(142.5,-117.8)
				..controls(145.8,-117.8)and(148.4,-117.7)..(148.8,-117.6)
				..controls(149.2,-117.6)and(150.2,-117.2)..(151,-116.8)
				..controls(152,-116.4)and(152.8,-116)..(153.6,-115.9)
				..controls(154.5,-115.7)and(155.2,-115.4)..(156.3,-114.9)
				..controls(157.1,-114.5)and(158.2,-114)..(158.7,-113.8)
				..controls(159.1,-113.7)and(159.8,-113.3)..(160.2,-113)
				..controls(160.6,-112.7)and(161.3,-112.4)..(161.7,-112.2)
				..controls(162.7,-111.9)and(163.2,-111.7)..(164.1,-111.1)
				..controls(164.5,-110.8)and(165.3,-110.4)..(165.9,-110)
				..controls(166.5,-109.7)and(167.3,-109.2)..(167.6,-108.9)
				..controls(167.9,-108.6)and(168.6,-108.2)..(169,-108)
				..controls(169.4,-107.8)and(170.2,-107.3)..(170.7,-106.9)
				..controls(171.2,-106.5)and(172.1,-105.8)..(172.7,-105.4)
				..controls(173.2,-105)and(174.1,-104.4)..(174.5,-104)
				..controls(174.9,-103.7)and(175.7,-103.1)..(176.2,-102.6)
				..controls(176.7,-102.2)and(177.4,-101.5)..(177.9,-101.2)
				..controls(178.3,-100.8)and(178.8,-100.4)..(179,-100.2)
				..controls(179.2,-100)and(179.6,-99.6)..(180,-99.3)
				..controls(180.4,-99.1)and(180.9,-98.6)..(181.2,-98.4)
				..controls(181.5,-98.1)and(182.1,-97.7)..(182.5,-97.5)
				..controls(183.2,-97.2)and(183.5,-96.9)..(184.3,-95.9)
				..controls(184.6,-95.6)and(184.6,-95.6)..(184.8,-95.7)
				..controls(185,-95.8)and(185.1,-96)..(185.1,-96.2)
				..controls(185.1,-96.5)and(185,-96.7)..(184.6,-97.2)
				..controls(184.1,-97.8)and(184,-97.9)..(183.9,-98.8)
				..controls(183.7,-100.3)and(183.3,-101.3)..(182.9,-101.8)
				..controls(182.6,-102.1)and(182.3,-102.6)..(182.2,-102.9)
				..controls(182,-103.3)and(181.6,-104)..(181.2,-104.6)
				..controls(180.9,-105.2)and(180.5,-106.1)..(180.4,-106.6)
				..controls(179.9,-108)and(179.7,-108.5)..(179,-109.4)
				..controls(178.6,-109.8)and(178.2,-110.5)..(177.9,-111)
				..controls(177.7,-111.5)and(177.3,-112.3)..(177,-112.7)
				..controls(176.6,-113.2)and(176.4,-113.7)..(176.3,-114.2)
				..controls(176.2,-114.7)and(175.9,-115.2)..(175.4,-116)
				..controls(175,-116.6)and(174.4,-117.7)..(174,-118.5)
				..controls(173.7,-119.2)and(173.3,-120.1)..(173,-120.5)
				..controls(172.8,-120.9)and(172.4,-121.8)..(172.1,-122.5)
				..controls(171.8,-123.2)and(171.3,-124.1)..(171,-124.6)
				..controls(170.8,-125.1)and(170.3,-126)..(169.9,-126.6)
				..controls(169.6,-127.3)and(169.2,-128.1)..(169,-128.3)
				..controls(168.8,-128.6)and(168.5,-129.3)..(168.3,-129.8)
				..controls(168.1,-130.3)and(167.8,-131.1)..(167.5,-131.6)
				..controls(166.9,-132.5)and(164.2,-137.9)..(164.2,-138.1)
				..controls(164.2,-138.1)and(163.9,-138.8)..(163.4,-139.7)
				..controls(162.9,-140.5)and(162.4,-141.5)..(162.3,-141.9)
				..controls(161.9,-143)and(161.6,-143.4)..(160.9,-144.1)
				..controls(160.4,-144.6)and(160.3,-144.8)..(160.1,-145.4)
				..controls(160,-145.8)and(159.6,-146.8)..(159.2,-147.8)
				..controls(158.7,-148.8)and(158.1,-150)..(157.9,-150.6)
				..controls(157.6,-151.1)and(157.2,-151.9)..(156.9,-152.4)
				..controls(156.5,-152.8)and(156.1,-153.7)..(155.8,-154.4)
				..controls(155.6,-155)and(155.2,-155.9)..(154.9,-156.3)
				..controls(154.7,-156.7)and(154.4,-157.4)..(154.2,-157.9)
				..controls(154.1,-158.4)and(153.8,-159)..(153.5,-159.3)
				..controls(152.8,-160.1)and(153,-160.1)..(142.4,-160.1)
				..controls(137.2,-160.1)and(131.8,-160.1)..(130.4,-160)--cycle
				(141.5,-103.7)..controls(140.8,-103)and(139.8,-102.1)..(139.2,-101.6)
				..controls(138.5,-101.2)and(137.9,-100.6)..(137.8,-100.4)
				..controls(137,-99.5)and(134,-96.5)..(133.2,-96)
				..controls(132.7,-95.7)and(131.8,-94.9)..(131.1,-94.3)
				..controls(129.4,-92.7)and(127.9,-91.6)..(126.6,-91.1)
				..controls(125.2,-90.5)and(124.1,-90)..(122.9,-89.1)
				..controls(122.1,-88.6)and(121.3,-88.2)..(119.9,-87.7)
				..controls(118.8,-87.3)and(117.5,-86.8)..(116.9,-86.5)
				..controls(114.3,-85.5)and(113.7,-85.4)..(109.7,-85.5)
				..controls(105.3,-85.7)and(103.9,-85.9)..(101.6,-87)
				..controls(100.7,-87.4)and(99.4,-87.8)..(97.9,-88)
				..controls(97,-88.2)and(96.6,-88.3)..(95.8,-88.9)
				..controls(95.4,-89.2)and(94.7,-89.6)..(94.2,-89.8)
				..controls(92.9,-90.3)and(92.3,-90.7)..(90.9,-92.1)
				..controls(90.2,-92.7)and(89.4,-93.3)..(89,-93.5)..controls(87,-94.5)and(85.2,-96)..(83.6,-97.9)
				..controls(83.3,-98.2)and(82.5,-98.9)..(81.8,-99.4)
				..controls(80.7,-100.1)and(80.5,-100.3)..(80,-101.1)
				..controls(78.7,-103.1)and(77.2,-103.9)..(76.4,-103)
				..controls(75.7,-102.2)and(75.7,-103.2)..(75.7,-83.6)
				..controls(75.6,-73.6)and(75.6,-64.8)..(75.5,-64)
				..controls(75.5,-63.2)and(75.5,-61.4)..(75.6,-59.8)--(75.7,-57.1)--(76,-56.8)
				..controls(76.8,-56.2)and(78,-56.3)..(79.8,-57.2)
				..controls(80.4,-57.5)and(81.3,-57.8)..(81.8,-57.9)
				..controls(82.3,-58.1)and(83,-58.4)..(83.6,-58.8)..controls(84.3,-59.2)and(85,-59.5)..(86,-59.8)
				..controls(86.8,-60.1)and(88,-60.6)..(88.7,-60.9)
				..controls(89.3,-61.2)and(90.6,-61.7)..(91.4,-62)
				..controls(92.2,-62.3)and(93.7,-62.9)..(94.6,-63.3)
				..controls(95.6,-63.7)and(96.8,-64.1)..(97.4,-64.3)
				..controls(98,-64.5)and(99,-64.9)..(99.7,-65.2)
				..controls(100.4,-65.6)and(101.5,-66)..(102.3,-66.2)
				..controls(104.2,-66.8)and(105,-67.1)..(106.2,-67.8)
				..controls(107.9,-68.7)and(110.2,-69.4)..(111.3,-69.2)
				..controls(112.5,-69)and(115,-68.2)..(115.6,-67.8)
				..controls(117.3,-66.7)and(117.6,-66.6)..(119.5,-66.2)
				..controls(121.9,-65.8)and(122.1,-65.7)..(123,-65.1)
				..controls(124.1,-64.4)and(124.4,-64.3)..(126,-63.9)
				..controls(127,-63.7)and(127.7,-63.4)..(128.4,-63.1)
				..controls(129,-62.8)and(129.7,-62.5)..(130.6,-62.4)
				..controls(132.1,-62.1)and(134,-61.5)..(135.2,-61)
				..controls(135.7,-60.8)and(136.7,-60.5)..(137.3,-60.3)
				..controls(138.4,-60)and(139.3,-59.7)..(141.1,-59)
				..controls(141.6,-58.7)and(142,-58.6)..(142.5,-58.6)
				..controls(143.2,-58.6)and(143.3,-58.7)..(143.6,-59)
				..controls(144,-59.5)and(144.1,-59.9)..(144,-60.8)
				..controls(144,-61.2)and(144,-71.1)..(144,-82.8)--(143.9,-104.1)--(143.7,-104.5)
				..controls(143.5,-104.7)and(143.3,-104.8)..(143,-104.9)
				..controls(142.6,-104.9)and(142.5,-104.8)..(141.5,-103.7)--cycle
				(27.8,-91.4)..controls(27.1,-90.9)and(26.3,-90)..(26.1,-89.5)
				..controls(26,-89.2)and(25.8,-88.6)..(25.7,-87.9)
				..controls(25.4,-86.2)and(25.6,-84.8)..(26.5,-83.5)
				..controls(27.4,-82.3)and(28.3,-82)..(29.4,-82.3)
				..controls(31.1,-82.8)and(32.1,-84.9)..(32,-87.4)
				..controls(31.9,-89.3)and(31.2,-90.6)..(29.9,-91.4)
				..controls(29.1,-91.9)and(28.6,-91.9)..(27.8,-91.4)--cycle
				(190.1,-91.3)..controls(189.7,-91.2)and(189.3,-90.9)..(189.1,-90.7)
				..controls(188.3,-89.8)and(187.6,-87.7)..(187.6,-86.2)
				..controls(187.6,-85.5)and(187.7,-85.2)..(188.2,-84.4)
				..controls(189.1,-82.6)and(190.1,-81.8)..(191.3,-82.1)
				..controls(192,-82.3)and(193.5,-83.7)..(193.8,-84.4)
				..controls(194.1,-85.1)and(194.2,-87.1)..(194,-88.3)
				..controls(193.8,-89.8)and(193.3,-90.5)..(192.3,-91.3)
				..controls(191.7,-91.7)and(190.9,-91.7)..(190.1,-91.3)--cycle
				(173.7,-83.4)..controls(172.2,-82.5)and(172,-82.5)..(170.6,-82.4)
				..controls(169.7,-82.4)and(168.8,-82.4)..(168.3,-82.5)
				..controls(167.2,-82.7)and(166.9,-82.7)..(166.6,-82.2)
				..controls(166.3,-81.8)and(166.3,-81)..(166.8,-80)
				..controls(167.2,-78.9)and(167.4,-77.4)..(167.7,-74.1)
				..controls(167.8,-72.8)and(167.9,-71.8)..(167.9,-71.7)
				..controls(168,-71.7)and(168,-70.2)..(168,-68.4)..controls(168,-65.4)and(168,-64.9)..(167.7,-63)
				..controls(167.3,-60.5)and(167.1,-59.9)..(166.4,-58.4)
				..controls(166.1,-57.7)and(165.8,-56.9)..(165.7,-56.6)
				..controls(165.7,-56.2)and(165.4,-55.6)..(165.2,-55.2)
				..controls(164.6,-53.8)and(161.9,-50.7)..(161,-50.3)
				..controls(159.8,-49.7)and(157.9,-50)..(155.2,-51.2)
				..controls(154.7,-51.4)and(153.4,-51.8)..(152.2,-52.1)
				..controls(150.8,-52.4)and(149.4,-52.9)..(148.6,-53.2)
				..controls(147.8,-53.6)and(146.7,-54)..(146,-54.2)
				..controls(145.4,-54.4)and(144.5,-54.7)..(144,-54.9)
				..controls(143.6,-55.2)and(142.7,-55.5)..(141.8,-55.7)
				..controls(141,-55.9)and(139.6,-56.3)..(138.7,-56.6)
				..controls(137.8,-57)and(136.3,-57.4)..(135.3,-57.6)
				..controls(133.8,-57.9)and(133.2,-58.1)..(132.1,-58.7)
				..controls(131.1,-59.2)and(130.3,-59.4)..(129.4,-59.6)
				..controls(128.7,-59.8)and(127.8,-60)..(127.4,-60.1)
				..controls(127,-60.3)and(126,-60.5)..(125.1,-60.7)
				..controls(123.2,-61.1)and(121.7,-61.5)..(120.2,-62.1)
				..controls(119.7,-62.3)and(118.1,-62.8)..(116.7,-63.2)
				..controls(114.6,-63.9)and(114,-64.1)..(113,-64.7)
				..controls(111.9,-65.4)and(111.8,-65.4)..(111,-65.5)
				..controls(109.9,-65.5)and(109.3,-65.4)..(107.7,-64.7)
				..controls(107,-64.4)and(106,-64)..(105.3,-63.9)
				..controls(103.5,-63.5)and(102.2,-63)..(101.2,-62.5)
				..controls(100.1,-61.8)and(99.3,-61.5)..(97.8,-61.1)
				..controls(97.2,-61)and(96.2,-60.7)..(95.7,-60.4)
				..controls(95.1,-60.1)and(94.2,-59.8)..(93.5,-59.6)
				..controls(92.9,-59.4)and(91.7,-59)..(91,-58.7)..controls(90.2,-58.4)and(88.9,-58)..(88,-57.7)
				..controls(87.2,-57.5)and(86.1,-57.1)..(85.5,-56.8)
				..controls(85,-56.6)and(84.1,-56.2)..(83.5,-56)..controls(82.9,-55.7)and(81.8,-55.2)..(81,-54.9)
				..controls(80.1,-54.5)and(79.1,-54.2)..(78.7,-54.1)
				..controls(77.9,-53.8)and(77.3,-53.6)..(76.2,-52.9)
				..controls(75.7,-52.6)and(75.2,-52.4)..(74.1,-52.2)
				..controls(72.5,-51.9)and(72.4,-51.8)..(71.3,-51.1)
				..controls(69.9,-50.1)and(69.4,-49.9)..(67.9,-49.8)
				..controls(67,-49.7)and(66.6,-49.6)..(66.3,-49.4)
				..controls(66.1,-49.3)and(65.7,-49)..(65.4,-48.8)
				..controls(65.1,-48.7)and(64.2,-48.4)..(63.6,-48.2)
				..controls(62.9,-48)and(61.7,-47.6)..(60.9,-47.2)
				..controls(60.1,-46.8)and(58.9,-46.4)..(58.2,-46.2)
				..controls(57.6,-46)and(56.8,-45.7)..(56.5,-45.5)
				..controls(54.9,-44.6)and(54.6,-44.5)..(53.2,-44.2)
				..controls(51.7,-43.8)and(51,-43.6)..(49.9,-42.9)
				..controls(49.5,-42.6)and(48.9,-42.4)..(47.5,-42)
				..controls(45.7,-41.6)and(45.6,-41.5)..(44.3,-40.6)
				..controls(43.5,-40.1)and(42.8,-39.6)..(42.7,-39.4)
				..controls(42.4,-38.7)and(43,-38.1)..(44.5,-37.6)..controls(44.9,-37.5)and(46,-37)..(46.9,-36.6)
				..controls(47.9,-36.1)and(48.8,-35.8)..(49.6,-35.6)
				..controls(50.5,-35.4)and(51.3,-35.2)..(52.4,-34.6)
				..controls(53.3,-34.1)and(54.5,-33.7)..(54.9,-33.5)
				..controls(55.4,-33.4)and(56.4,-33)..(57.2,-32.6)
				..controls(58.6,-32)and(60.1,-31.5)..(61.6,-31.2)
				..controls(62.3,-31.1)and(62.7,-31)..(63.2,-30.7)
				..controls(64.1,-30)and(64.3,-29.9)..(65.3,-29.7)
				..controls(65.7,-29.6)and(66.5,-29.2)..(67.1,-28.9)
				..controls(67.6,-28.6)and(68.8,-28.2)..(69.6,-27.9)
				..controls(70.5,-27.7)and(71.6,-27.3)..(72.1,-27)
				..controls(72.6,-26.7)and(73.8,-26.2)..(74.8,-25.9)
				..controls(75.7,-25.5)and(76.7,-25.1)..(77.1,-25)
				..controls(77.4,-24.8)and(78.5,-24.4)..(79.6,-24.1)
				..controls(80.6,-23.8)and(82,-23.2)..(82.8,-22.9)
				..controls(83.5,-22.5)and(84.8,-22)..(85.8,-21.8)
				..controls(86.7,-21.5)and(87.8,-21.1)..(88.1,-20.9)
				..controls(88.4,-20.7)and(89.3,-20.3)..(90.1,-20.1)
				..controls(90.8,-19.8)and(91.9,-19.4)..(92.4,-19.1)
				..controls(93.5,-18.5)and(94.4,-18.2)..(95.4,-18.1)
				..controls(95.7,-18)and(96.4,-17.8)..(96.9,-17.5)
				..controls(97.5,-17.3)and(98.4,-16.9)..(99,-16.6)
				..controls(99.7,-16.3)and(100.9,-15.9)..(101.9,-15.6)
				..controls(102.9,-15.3)and(104.4,-14.7)..(105.3,-14.3)
				..controls(108.1,-13.2)and(108.5,-13.1)..(110.1,-13.1)
				..controls(111.4,-13.1)and(111.5,-13.1)..(112.4,-13.5)
				..controls(113.2,-14)and(115,-14.8)..(116.6,-15.4)
				..controls(116.9,-15.5)and(117.9,-15.8)..(118.8,-16.1)
				..controls(120.9,-16.6)and(121.5,-16.9)..(122.6,-17.9)
				..controls(123,-18.3)and(123.7,-18.8)..(123.9,-18.9)
				..controls(124.6,-19.4)and(126,-19.8)..(127.9,-20.2)
				..controls(129.6,-20.5)and(129.8,-20.6)..(131.1,-21.4)
				..controls(131.5,-21.7)and(132.4,-22.1)..(133.1,-22.3)
				..controls(136.1,-23.2)and(138.3,-24.1)..(140.2,-25)
				..controls(141,-25.5)and(141.5,-25.6)..(142.5,-25.8)
				..controls(145.1,-26.3)and(146,-26.6)..(149.4,-28.1)
				..controls(151.3,-28.9)and(153.5,-29.8)..(154.3,-30.1)
				..controls(155.1,-30.4)and(156.2,-30.8)..(156.6,-31)
				..controls(157.1,-31.2)and(158.2,-31.6)..(159.1,-31.8)
				..controls(160.6,-32.2)and(160.9,-32.3)..(161.7,-32.9)
				..controls(162.4,-33.3)and(162.9,-33.6)..(163.7,-33.8)
				..controls(164.3,-34)and(164.8,-34.1)..(164.9,-34.1)
				..controls(165,-34.1)and(165.6,-34.4)..(166.4,-34.8)
				..controls(167.3,-35.2)and(168.1,-35.4)..(169,-35.6)
				..controls(170,-35.8)and(170.5,-36)..(171.7,-36.6)
				..controls(172.4,-37)and(173.3,-37.4)..(173.6,-37.4)
				..controls(175.6,-37.7)and(176.2,-38)..(176.7,-39.1)
				..controls(176.9,-39.6)and(176.9,-39.7)..(176.7,-40.4)
				..controls(176.4,-41.6)and(175.8,-42.7)..(175.2,-43.4)
				..controls(174.5,-44.1)and(172.6,-45.5)..(171.7,-45.8)
				..controls(170.6,-46.2)and(169.1,-46)..(168.2,-45.4)
				..controls(167.9,-45.2)and(167.3,-44.8)..(166.9,-44.4)
				..controls(166.4,-44.1)and(165.8,-43.6)..(165.4,-43.2)
				..controls(165.1,-42.9)and(164.4,-42.3)..(163.9,-41.9)
				..controls(163.3,-41.5)and(162.7,-40.9)..(162.4,-40.5)
				..controls(161.4,-39.2)and(159.9,-38)..(158.8,-37.8)
				..controls(158.4,-37.8)and(158.3,-37.7)..(158.2,-37.2)
				..controls(158,-36.8)and(157.4,-36.1)..(157.1,-36.1)
				..controls(157,-36.1)and(156.7,-36.1)..(156.5,-36)
				..controls(156.1,-35.9)and(155.9,-35.7)..(155.7,-35.2)
				..controls(155.3,-34.6)and(155.1,-34.4)..(154,-34.1)
				..controls(153.6,-34)and(153.2,-33.7)..(152.9,-33.5)
				..controls(152.4,-32.9)and(151.6,-32.4)..(151.3,-32.4)
				..controls(150.9,-32.4)and(150.4,-32.6)..(150.4,-32.9)
				..controls(150.4,-33.2)and(150.8,-33.4)..(151.4,-33.6)
				..controls(152.1,-33.7)and(152.8,-34.1)..(153.5,-34.9)
				..controls(154.1,-35.5)and(154.3,-35.6)..(155,-35.8)
				..controls(155.7,-36)and(155.7,-36)..(156,-36.6)
				..controls(156.3,-37.3)and(156.7,-37.7)..(157.5,-37.8)
				..controls(158.1,-37.9)and(158.2,-38)..(158.4,-38.8)
				..controls(158.5,-39.4)and(158.6,-39.6)..(159.6,-40.4)
				..controls(160.2,-41)and(160.9,-41.6)..(161.3,-41.8)
				..controls(161.7,-42)and(161.9,-42.3)..(162.2,-42.8)
				..controls(162.4,-43.2)and(162.8,-43.8)..(163.1,-44.1)
				..controls(163.9,-44.8)and(165.1,-46.7)..(165.7,-48)
				..controls(166.2,-49.3)and(166.4,-49.7)..(167.2,-50.5)
				..controls(167.5,-50.8)and(167.9,-51.4)..(168.1,-51.8)
				..controls(168.3,-52.2)and(168.8,-53)..(169.1,-53.6)
				..controls(169.6,-54.3)and(169.8,-54.9)..(169.9,-55.6)
				..controls(170.4,-57.8)and(170.5,-58.1)..(171.2,-59.5)
				..controls(171.6,-60.3)and(172,-61.5)..(172.2,-62.1)
				..controls(172.3,-62.8)and(172.7,-63.7)..(172.9,-64.1)
				..controls(173.2,-64.6)and(173.4,-65.3)..(173.4,-65.6)
				..controls(173.8,-68.8)and(174.2,-70.6)..(174.9,-72.6)
				..controls(175.5,-74.5)and(175.9,-77.1)..(176.1,-80.2)
				..controls(176.3,-83.4)and(175.5,-84.4)..(173.7,-83.4)--cycle;
				\definecolor{ezs}{RGB}{247,179,26}\fill[ezs](-5.7,-156.6)--(-3.4,-157.5)--(61.9,12.7)--(59.7,13.6)--cycle;
				\definecolor{ezs}{RGB}{247,179,22}\fill[ezs](-3.8,-157.3)--(-1.5,-158.2)--(63.9,12)--(61.6,12.9)--cycle;
				\definecolor{ezs}{RGB}{248,179,19}\fill[ezs](-1.8,-158.1)--(0.4,-158.9)--(65.8,11.2)--(63.5,12.1)--cycle;
				\definecolor{ezs}{RGB}{248,178,16}\fill[ezs](0.1,-158.8)--(2.4,-159.7)--(67.7,10.5)--(65.4,11.4)--cycle;
				\definecolor{ezs}{RGB}{249,178,15}\fill[ezs](2,-159.5)--(4.3,-160.4)--(69.6,9.8)--(67.3,10.7)--cycle;
				\definecolor{ezs}{RGB}{251,177,17}\fill[ezs](3.9,-160.3)--(6.2,-161.1)--(71.5,9)--(69.2,9.9)--cycle;
				\definecolor{ezs}{RGB}{252,179,18}\fill[ezs](5.8,-161)--(8.1,-161.9)--(73.4,8.3)--(71.1,9.2)--cycle;
				\definecolor{ezs}{RGB}{252,182,17}\fill[ezs](7.7,-161.7)--(10,-162.6)--(75.3,7.6)--(73,8.5)--cycle;
				\definecolor{ezs}{RGB}{252,185,17}\fill[ezs](9.6,-162.5)--(11.9,-163.4)--(77.2,6.8)--(75,7.7)--cycle;
				\definecolor{ezs}{RGB}{251,189,18}\fill[ezs](11.5,-163.2)--(13.8,-164.1)--(79.2,6.1)--(76.9,7)--cycle;
				\definecolor{ezs}{RGB}{251,191,22}\fill[ezs](13.4,-163.9)--(15.7,-164.8)--(81.1,5.4)--(78.8,6.2)--cycle;
				\definecolor{ezs}{RGB}{252,193,25}\fill[ezs](15.4,-164.7)--(17.6,-165.6)--(83,4.6)--(80.7,5.5)--cycle;
				\definecolor{ezs}{RGB}{253,195,27}\fill[ezs](17.3,-165.4)--(19.6,-166.3)--(84.9,3.9)--(82.6,4.8)--cycle;
				\definecolor{ezs}{RGB}{253,196,31}\fill[ezs](19.2,-166.1)--(21.5,-167)--(86.8,3.2)--(84.5,4)--cycle;
				\definecolor{ezs}{RGB}{254,199,35}\fill[ezs](21.1,-166.9)--(23.4,-167.8)--(88.7,2.4)--(86.4,3.3)--cycle;
				\definecolor{ezs}{RGB}{253,201,38}\fill[ezs](23,-167.6)--(25.3,-168.5)--(90.6,1.7)--(88.3,2.6)--cycle;
				\definecolor{ezs}{RGB}{253,203,40}\fill[ezs](24.9,-168.3)--(27.2,-169.2)--(92.5,1)--(90.3,1.8)--cycle;
				\definecolor{ezs}{RGB}{253,205,42}\fill[ezs](26.8,-169.1)--(29.1,-170)--(94.4,0.2)--(92.2,1.1)--cycle;
				\definecolor{ezs}{RGB}{253,207,45}\fill[ezs](28.7,-169.8)--(31,-170.7)--(96.4,-0.5)--(94.1,0.4)--cycle;
				\definecolor{ezs}{RGB}{253,210,49}\fill[ezs](30.7,-170.5)--(32.9,-171.4)--(98.3,-1.2)--(96,-0.4)--cycle;
				\definecolor{ezs}{RGB}{253,212,54}\fill[ezs](32.6,-171.3)--(34.9,-172.2)--(100.2,-2)--(97.9,-1.1)--cycle;
				\definecolor{ezs}{RGB}{253,214,59}\fill[ezs](34.5,-172)--(36.8,-172.9)--(102.1,-2.7)--(99.8,-1.8)--cycle;
				\definecolor{ezs}{RGB}{253,216,63}\fill[ezs](36.4,-172.7)--(38.7,-173.6)--(104,-3.4)--(101.7,-2.6)--cycle;
				\definecolor{ezs}{RGB}{253,218,65}\fill[ezs](38.3,-173.5)--(40.6,-174.4)--(105.9,-4.2)--(103.6,-3.3)--cycle;
				\definecolor{ezs}{RGB}{254,220,66}\fill[ezs](40.2,-174.2)--(42.5,-175.1)--(107.8,-4.9)--(105.5,-4)--cycle;
				\definecolor{ezs}{RGB}{254,222,68}\fill[ezs](42.1,-175)--(44.4,-175.8)--(109.7,-5.6)--(107.5,-4.8)--cycle;
				\definecolor{ezs}{RGB}{254,224,71}\fill[ezs](44,-175.7)--(46.3,-176.6)--(111.7,-6.4)--(109.4,-5.5)--cycle;
				\definecolor{ezs}{RGB}{254,226,75}\fill[ezs](46,-176.4)--(48.2,-177.3)--(113.6,-7.1)--(111.3,-6.2)--cycle;
				\definecolor{ezs}{RGB}{254,228,79}\fill[ezs](47.9,-177.2)--(50.1,-178)--(115.5,-7.8)--(113.2,-7)--cycle;
				\definecolor{ezs}{RGB}{254,230,82}\fill[ezs](49.8,-177.9)--(52.1,-178.8)--(117.4,-8.6)--(115.1,-7.7)--cycle;
				\definecolor{ezs}{RGB}{254,232,85}\fill[ezs](51.7,-178.6)--(54,-179.5)--(119.3,-9.3)--(117,-8.4)--cycle;
				\definecolor{ezs}{RGB}{254,234,88}\fill[ezs](53.6,-179.4)--(55.9,-180.2)--(121.2,-10)--(118.9,-9.2)--cycle;
				\definecolor{ezs}{RGB}{254,237,91}\fill[ezs](55.5,-180.1)--(57.8,-181)--(123.1,-10.8)--(120.8,-9.9)--cycle;
				\definecolor{ezs}{RGB}{254,239,94}\fill[ezs](57.4,-180.8)--(59.7,-181.7)--(125,-11.5)--(122.8,-10.6)--cycle;
				\definecolor{ezs}{RGB}{254,241,97}\fill[ezs](59.3,-181.6)--(61.6,-182.4)--(127,-12.2)--(124.7,-11.4)--cycle;
				\definecolor{ezs}{RGB}{254,243,101}\fill[ezs](61.2,-182.3)--(63.5,-183.2)--(128.9,-13)--(126.6,-12.1)--cycle;
				\definecolor{ezs}{RGB}{254,245,103}\fill[ezs](63.2,-183)--(65.4,-183.9)--(130.8,-13.7)--(128.5,-12.8)--cycle;
				\definecolor{ezs}{RGB}{254,245,105}\fill[ezs](65.1,-183.8)--(67.4,-184.6)--(132.7,-14.4)--(130.4,-13.6)--cycle;
				\definecolor{ezs}{RGB}{254,245,105}\fill[ezs](67,-184.5)--(69.3,-185.4)--(134.6,-15.2)--(132.3,-14.3)--cycle;
				\definecolor{ezs}{RGB}{254,244,103}\fill[ezs](68.9,-185.2)--(71.2,-186.1)--(136.5,-15.9)--(134.2,-15)--cycle;
				\definecolor{ezs}{RGB}{253,243,101}\fill[ezs](70.8,-186)--(73.1,-186.8)--(138.4,-16.6)--(136.1,-15.8)--cycle;
				\definecolor{ezs}{RGB}{253,241,98}\fill[ezs](72.7,-186.7)--(75,-187.6)--(140.3,-17.4)--(138,-16.5)--cycle;
				\definecolor{ezs}{RGB}{253,239,95}\fill[ezs](74.6,-187.4)--(76.9,-188.3)--(142.2,-18.1)--(140,-17.2)--cycle;
				\definecolor{ezs}{RGB}{253,236,92}\fill[ezs](76.5,-188.2)--(78.8,-189)--(144.2,-18.8)--(141.9,-18)--cycle;
				\definecolor{ezs}{RGB}{253,234,88}\fill[ezs](78.5,-188.9)--(80.7,-189.8)--(146.1,-19.6)--(143.8,-18.7)--cycle;
				\definecolor{ezs}{RGB}{253,232,85}\fill[ezs](80.4,-189.6)--(82.7,-190.5)--(148,-20.3)--(145.7,-19.4)--cycle;
				\definecolor{ezs}{RGB}{253,230,82}\fill[ezs](82.3,-190.4)--(84.6,-191.2)--(149.9,-21)--(147.6,-20.2)--cycle;
				\definecolor{ezs}{RGB}{253,228,78}\fill[ezs](84.2,-191.1)--(86.5,-192)--(151.8,-21.8)--(149.5,-20.9)--cycle;
				\definecolor{ezs}{RGB}{253,226,74}\fill[ezs](86.1,-191.8)--(88.4,-192.7)--(153.7,-22.5)--(151.4,-21.6)--cycle;
				\definecolor{ezs}{RGB}{253,224,70}\fill[ezs](88,-192.6)--(90.3,-193.4)--(155.6,-23.3)--(153.3,-22.4)--cycle;
				\definecolor{ezs}{RGB}{253,222,67}\fill[ezs](89.9,-193.3)--(92.2,-194.2)--(157.5,-24)--(155.3,-23.1)--cycle;
				\definecolor{ezs}{RGB}{253,220,65}\fill[ezs](91.8,-194)--(94.1,-194.9)--(159.5,-24.7)--(157.2,-23.8)--cycle;
				\definecolor{ezs}{RGB}{253,217,63}\fill[ezs](93.7,-194.8)--(96,-195.6)--(161.4,-25.5)--(159.1,-24.6)--cycle;
				\definecolor{ezs}{RGB}{253,215,60}\fill[ezs](95.7,-195.5)--(97.9,-196.4)--(163.3,-26.2)--(161,-25.3)--cycle;
				\definecolor{ezs}{RGB}{253,213,56}\fill[ezs](97.6,-196.2)--(99.9,-197.1)--(165.2,-26.9)--(162.9,-26)--cycle;
				\definecolor{ezs}{RGB}{252,211,53}\fill[ezs](99.5,-197)--(101.8,-197.8)--(167.1,-27.7)--(164.8,-26.8)--cycle;
				\definecolor{ezs}{RGB}{252,209,49}\fill[ezs](101.4,-197.7)--(103.7,-198.6)--(169,-28.4)--(166.7,-27.5)--cycle;
				\definecolor{ezs}{RGB}{252,207,45}\fill[ezs](103.3,-198.4)--(105.6,-199.3)--(170.9,-29.1)--(168.6,-28.2)--cycle;
				\definecolor{ezs}{RGB}{253,205,42}\fill[ezs](105.2,-199.2)--(107.5,-200)--(172.8,-29.9)--(170.6,-29)--cycle;
				\definecolor{ezs}{RGB}{253,202,40}\fill[ezs](107.1,-199.9)--(109.4,-200.8)--(174.7,-30.6)--(172.5,-29.7)--cycle;
				\definecolor{ezs}{RGB}{253,200,38}\fill[ezs](109,-200.6)--(111.3,-201.5)--(176.7,-31.3)--(174.4,-30.4)--cycle;
				\definecolor{ezs}{RGB}{253,198,35}\fill[ezs](111,-201.4)--(113.2,-202.2)--(178.6,-32.1)--(176.3,-31.2)--cycle;
				\definecolor{ezs}{RGB}{253,196,31}\fill[ezs](112.9,-202.1)--(115.2,-203)--(180.5,-32.8)--(178.2,-31.9)--cycle;
				\definecolor{ezs}{RGB}{253,194,28}\fill[ezs](114.8,-202.8)--(117.1,-203.7)--(182.4,-33.5)--(180.1,-32.6)--cycle;
				\definecolor{ezs}{RGB}{253,192,26}\fill[ezs](116.7,-203.6)--(119,-204.5)--(184.3,-34.3)--(182,-33.4)--cycle;
				\definecolor{ezs}{RGB}{253,189,24}\fill[ezs](118.6,-204.3)--(120.9,-205.2)--(186.2,-35)--(183.9,-34.1)--cycle;
				\definecolor{ezs}{RGB}{253,186,20}\fill[ezs](120.5,-205)--(122.8,-205.9)--(188.1,-35.7)--(185.8,-34.8)--cycle;
				\definecolor{ezs}{RGB}{253,184,17}\fill[ezs](122.4,-205.8)--(124.7,-206.7)--(190,-36.5)--(187.8,-35.6)--cycle;
				\definecolor{ezs}{RGB}{252,183,15}\fill[ezs](124.3,-206.5)--(126.6,-207.4)--(192,-37.2)--(189.7,-36.3)--cycle;
				\definecolor{ezs}{RGB}{252,181,14}\fill[ezs](126.3,-207.2)--(128.5,-208.1)--(193.9,-37.9)--(191.6,-37.1)--cycle;
				\definecolor{ezs}{RGB}{252,179,13}\fill[ezs](128.2,-208)--(130.4,-208.9)--(195.8,-38.7)--(193.5,-37.8)--cycle;
				\definecolor{ezs}{RGB}{252,178,12}\fill[ezs](130.1,-208.7)--(132.4,-209.6)--(197.7,-39.4)--(195.4,-38.5)--cycle;
				\definecolor{ezs}{RGB}{251,178,10}\fill[ezs](132,-209.4)--(134.3,-210.3)--(199.6,-40.1)--(197.3,-39.3)--cycle;
				\definecolor{ezs}{RGB}{251,178,10}\fill[ezs](133.9,-210.2)--(136.2,-211.1)--(201.5,-40.9)--(199.2,-40)--cycle;
				\definecolor{ezs}{RGB}{250,177,11}\fill[ezs](135.8,-210.9)--(138.1,-211.8)--(203.4,-41.6)--(201.1,-40.7)--cycle;
				\definecolor{ezs}{RGB}{250,178,11}\fill[ezs](137.7,-211.6)--(140,-212.5)--(205.3,-42.3)--(203.1,-41.5)--cycle;
				\definecolor{ezs}{RGB}{248,180,10}\fill[ezs](139.6,-212.4)--(141.9,-213.3)--(207.3,-43.1)--(205,-42.2)--cycle;
				\definecolor{ezs}{RGB}{247,181,9}\fill[ezs](141.5,-213.1)--(143.8,-214)--(209.2,-43.8)--(206.9,-42.9)--cycle;
				\definecolor{ezs}{RGB}{246,182,9}\fill[ezs](143.5,-213.8)--(145.7,-214.7)--(211.1,-44.5)--(208.8,-43.7)--cycle;
				\definecolor{ezs}{RGB}{246,182,10}\fill[ezs](145.4,-214.6)--(147.7,-215.5)--(213,-45.3)--(210.7,-44.4)--cycle;
			\end{scope}
			% chữ EDUZY (Times đậm, làm dày nét bằng 5 lớp lệch nhẹ; căn đúng khung chữ logo gốc)
			\foreach \dx/\dy in {-0.7/0,0.7/0,0/-0.3,0/0.3,0/0}
			\node[anchor=center,inner sep=0pt,outer sep=0pt,text=ezText,align=center] at (109.6+\dx,-193.6+\dy)
			{\resizebox{1.01cm}{0.22cm}{\fontencoding{T1}\fontfamily{ptm}\fontseries{b}\selectfont EDUZY}};
		\end{tikzpicture}%
	}%
}
%====================== THIẾT LẬP THÂN ĐỀ ======================
%--- Môi trường tự luận {bt}: ex_test 3.3.4 chưa khai báo sẵn nên tạo ở đây
%    (ex_test đã tự móc \fixshowans và \listenumerate cho bt)
\makeatletter
\@ifundefined{bt}{\newtheorem{bt}{\color{gold1}\namebt}}{}
\makeatother
\renewcommand{\sectionmark}[1]{\markright{\thesection\ #1}}
%--- Khung tiêu đề (logo + tối đa 4 dòng; dòng 3, 4 để trống {} thì bỏ qua)
%  Đề thi:    \khungtieude{THI THỬ THPT QUỐC GIA NĂM 2026}{Môn Toán -- Đề số 21}
%                         {(Đề thi gồm: ...)}{Thời gian làm bài: 90 phút (không kể thời gian phát đề).}
%  Chuyên đề: \khungtieude{CHUYÊN ĐỀ 3 -- DÃY SỐ VÀ CÔNG THỨC TRUY HỒI}
%                         {Bồi dưỡng học sinh giỏi Toán THPT}{...}{}
\newcommand{\khungtieude}[4]{%
	\begin{tcolorbox}[colframe=gold1,colback=gold5!45,boxrule=0.5mm]
		\begin{minipage}[c]{3cm}
			\centerline{\eduzylogo[2.4cm]}
		\end{minipage}\hspace{0.5cm}
		\begin{minipage}[c]{12.5cm}
			\vskip 0.15 cm
			\centerline{\ind{#1}}
			\vskip 0.15 cm
			\centerline{\textbf{#2}}
			\if\relax\detokenize{#3}\relax\else
				\vskip 0.15 cm
				\centerline{\textit{#3}}
			\fi
			\if\relax\detokenize{#4}\relax\else
				\vskip 0.15 cm
				\centerline{\textit{#4}}
			\fi
			\vskip 0.15 cm
		\end{minipage}
	\end{tcolorbox}
}
%====================== BỔ SUNG CHO BÀI HỌC, CHUYÊN ĐỀ ======================
\usepackage[shortlabels]{enumitem}
\usepackage{tabularx}
%--- Danh sách gọn: \begin{enumerate}[gon] ... \end{enumerate}
\SetEnumitemKey{gon}{leftmargin=*,itemsep=1pt,parsep=0pt,topsep=3pt}
%--- Khung lý thuyết: \begin{khung}{} ... \end{khung}; tiêu đề chỉ dùng khi thật ngắn ("Mẹo", "Lỗi thường gặp")
\newtcolorbox{khung}[1]{breakable,enhanced,colframe=gold1,colback=gold5!45,
	boxrule=0.4mm,fonttitle=\bfseries\sffamily,coltitle=white,colbacktitle=gold1,
	attach boxed title to top left={xshift=6mm,yshift=-3mm},
	boxed title style={arc=1mm,boxrule=0.3mm},title={#1}}
%--- Tiêu đề mục nhỏ trong một dạng: \muc{Dấu hiệu}, \muc{Phương pháp}, \muc{Bài tự luyện}
\makeatletter
\newcommand{\muc}[1]{\par\medskip\noindent\indam{#1}\par\nobreak\smallskip\@afterindentfalse\@afterheading}
\makeatother
%--- Ô công thức cần nhớ (căn giữa, viền vàng, vừa khít nội dung)
%     \begin{center}\congthuc{$Q(n)=(n-1)!$}\end{center}
\newtcbox{\congthuc}{on line,colframe=gold1,colback=gold5!45,boxrule=0.5pt,arc=2pt,
	left=8pt,right=8pt,top=4pt,bottom=4pt}
%--- Công tắc lời giải bài tự luyện (ví dụ mẫu luôn dùng \loigiai)
%     \loigiaiTL[{đáp số}]{lời giải đầy đủ}
%     \TLfull -> in lời giải đầy đủ;  \TLdapso -> chỉ in một dòng "Đáp số";  \TLan -> ẩn hoàn toàn
%     (\CheDo là dethi thì mọi lời giải và đáp số đều ẩn)
\makeatletter
\newif\ifTLhien   \TLhientrue
\newif\ifTLdapso  \TLdapsofalse
\newcommand{\TLfull}{\TLhientrue\TLdapsofalse}
\newcommand{\TLdapso}{\TLhientrue\TLdapsotrue}
\newcommand{\TLan}{\TLhienfalse}
\newcommand{\loigiaiTL}[2][]{%
	\ifTLhien
		\ifTLdapso
			\def\@tlarg{#1}%
			\ifx\@tlarg\@empty\else
				\ifdefstring{\kindOPTN}{dethi}{}{\par\noindent{\bfseries\color{gold1}Đáp số:} #1\par}%
			\fi
		\else
			\loigiai{#2}%
		\fi
	\fi}
\makeatother
\TLdapso% mặc định: chỉ in đáp số bài tự luyện
%====================== BẢNG ĐÚNG/SAI (\choiceTFt) ======================
%  Hàng tiêu đề nền nâu đồng chữ trắng (Phát biểu | Đ | S), các hàng tô xen kẽ kem / trắng,
%  đường kẻ vàng kim, nhãn a) b) c) d) màu nâu đồng. Bản có đáp án: cột được chọn có vòng tròn
%  nâu đồng với dấu tích trắng (câu đúng tích ở cột Đúng, câu sai tích ở cột Sai).
%  Tiêu đề ghi đủ "Đúng | Sai": \OPTN{viettat=0} sau \input (cột Đúng/Sai rộng ra, phát biểu hẹp lại).
%  Dùng bảng gốc của ex_test: \def\BangDSgoc{1} trước \input{preamble-eduzy-v2}.
\providecommand{\BangDSgoc}{0}
\ifnum\BangDSgoc=0
\usepackage{colortbl,pifont}
\colorlet{ezDSke}{gold2!60!gold1}% đường kẻ bảng (nét 0,4pt như ex_test để không đổi chiều cao bảng)
\colorlet{ezDSle}{gold5!60}% hàng a), c)
\colorlet{ezDSchan}{white}% hàng b), d)
%--- dấu chọn ở cột Đúng/Sai
%    (bề rộng 0, căn giữa ô: không làm đổi bề rộng cột của ex_test, không báo tràn lề)
\newcommand{\EZdauchon}{\makebox[0pt]{\tikz[baseline=-1ex]\node[circle,fill=gold1,inner sep=0pt,
	minimum size=1.2em,text=white,font=\small]{\ding{51}};}}
\def\tickT{\EZdauchon}\def\tickF{\EZdauchon}
\AtBeginDocument{\sbox{\ccTFTab}{\EZdauchon}}
%--- nhãn a) b) c) d) màu nâu đồng (ở bản lời giải, nhãn câu đúng vẫn đóng khung như ex_test)
\renewcommand{\FalseTF}{\stepcounter{dapan}\textbf{\color{gold1}\DapAnTF\sepTF}\ignorespaces}
%--- trước khi dựng bảng: kẻ màu vàng kim (bề rộng cột giữ nguyên như ex_test); sau bảng trả lại màu đen
\newcount\EZhangTF
\newcommand{\EZbatdauTF}{%
	\global\EZhangTF=0\relax
	\arrayrulecolor{ezDSke}}
\newcommand{\EZketthucTF}{\arrayrulecolor{black}}
%--- hàng tiêu đề nền nâu đồng, chữ trắng
\makeatletter
\newcommand{\EZo@TF}[1]{\centering\bfseries\color{white}#1}
\newcommand{\EZtieudeTF}{%
	\ifcase\cotDS
		\ifnum\cotSTT=1\relax
			\def\rowHeader{\hline\rowcolor{gold1}\EZo@TF{TT}&\EZo@TF{\phatbieu}\tabularnewline}%
		\else
			\def\rowHeader{\hline\rowcolor{gold1}\EZo@TF{\phatbieu}\tabularnewline}%
		\fi
	\or
		\ifnum\cotSTT=1\relax
			\def\rowHeader{\hline\rowcolor{gold1}\EZo@TF{TT}&\EZo@TF{\phatbieu}&\EZo@TF{\Dung}\tabularnewline}%
		\else
			\def\rowHeader{\hline\rowcolor{gold1}\EZo@TF{\phatbieu}&\EZo@TF{\Dung}\tabularnewline}%
		\fi
	\or
		\ifnum\cotSTT=1\relax
			\def\rowHeader{\hline\rowcolor{gold1}\EZo@TF{TT}&\EZo@TF{\phatbieu}&\EZo@TF{\Dung}&\EZo@TF{\Sai}\tabularnewline}%
		\else
			\def\rowHeader{\hline\rowcolor{gold1}\EZo@TF{\phatbieu}&\EZo@TF{\Dung}&\EZo@TF{\Sai}\tabularnewline}%
		\fi
	\fi}
%--- tô xen kẽ từng hàng (gọi ngay sau \hline đầu mỗi hàng)
\newcommand{\EZmauhangTF}{%
	\noalign{\global\advance\EZhangTF\@ne}%
	\ifodd\EZhangTF\rowcolor{ezDSle}\else\rowcolor{ezDSchan}\fi}
%--- gắn vào \choiceTFt / \writeChoiceTFt của ex_test
\xpretocmd{\choiceTFt}{\EZbatdauTF}{}{\typeout{[preamble-eduzy-v2] Không gắn được bảng Đúng/Sai (đầu bảng)}}
\xpatchcmd{\choiceTFt}{\ifnum\dongPB=0\def\rowHeader{}\fi}{\EZtieudeTF\ifnum\dongPB=0\def\rowHeader{}\fi}
	{}{\typeout{[preamble-eduzy-v2] Không gắn được bảng Đúng/Sai (hàng tiêu đề)}}
\xpatchcmd{\writeChoiceTFt}{\hline\xdef\dotEX{\sdotEX}}{\hline\EZmauhangTF\xdef\dotEX{\sdotEX}}
	{}{\typeout{[preamble-eduzy-v2] Không gắn được bảng Đúng/Sai (tô hàng)}}
\xpatchcmd{\writeChoiceTFt}{{\end{longtable}}}{{\end{longtable}\EZketthucTF}}
	{}{\typeout{[preamble-eduzy-v2] Không gắn được bảng Đúng/Sai (cuối longtable)}}
\xpatchcmd{\writeChoiceTFt}{{\end{xtabular}\end{center}\medskip\noindent}}{{\end{xtabular}\end{center}\EZketthucTF\medskip\noindent}}
	{}{\typeout{[preamble-eduzy-v2] Không gắn được bảng Đúng/Sai (cuối xtabular)}}
\fi
%=====================================================================
\renewcommand{\parallel}{\mathrel{/\!\!/}}
.
\providecommand{\BangDSgoc}{0}
\ifnum\BangDSgoc=0
\usepackage{colortbl,pifont}
\colorlet{ezDSke}{gold2!60!gold1}% đường kẻ bảng (nét 0,4pt như ex_test để không đổi chiều cao bảng)
\colorlet{ezDSle}{gold5!60}% hàng a), c)
\colorlet{ezDSchan}{white}% hàng b), d)
%--- dấu chọn ở cột Đúng/Sai
%    (bề rộng 0, căn giữa ô: không làm đổi bề rộng cột của ex_test, không báo tràn lề)
\newcommand{\EZdauchon}{\makebox[0pt]{\tikz[baseline=-1ex]\node[circle,fill=gold1,inner sep=0pt,
	minimum size=1.2em,text=white,font=\small]{\ding{51}};}}
\def\tickT{\EZdauchon}\def\tickF{\EZdauchon}
\AtBeginDocument{\sbox{\ccTFTab}{\EZdauchon}}
%--- nhãn a) b) c) d) màu nâu đồng (ở bản lời giải, nhãn câu đúng vẫn đóng khung như ex_test)
\renewcommand{\FalseTF}{\stepcounter{dapan}\textbf{\color{gold1}\DapAnTF\sepTF}\ignorespaces}
%--- trước khi dựng bảng: kẻ màu vàng kim (bề rộng cột giữ nguyên như ex_test); sau bảng trả lại màu đen
\newcount\EZhangTF
\newcommand{\EZbatdauTF}{%
	\global\EZhangTF=0\relax
	\arrayrulecolor{ezDSke}}
\newcommand{\EZketthucTF}{\arrayrulecolor{black}}
%--- hàng tiêu đề nền nâu đồng, chữ trắng
\makeatletter
\newcommand{\EZo@TF}[1]{\centering\bfseries\color{white}#1}
\newcommand{\EZtieudeTF}{%
	\ifcase\cotDS
		\ifnum\cotSTT=1\relax
			\def\rowHeader{\hline\rowcolor{gold1}\EZo@TF{TT}&\EZo@TF{\phatbieu}\tabularnewline}%
		\else
			\def\rowHeader{\hline\rowcolor{gold1}\EZo@TF{\phatbieu}\tabularnewline}%
		\fi
	\or
		\ifnum\cotSTT=1\relax
			\def\rowHeader{\hline\rowcolor{gold1}\EZo@TF{TT}&\EZo@TF{\phatbieu}&\EZo@TF{\Dung}\tabularnewline}%
		\else
			\def\rowHeader{\hline\rowcolor{gold1}\EZo@TF{\phatbieu}&\EZo@TF{\Dung}\tabularnewline}%
		\fi
	\or
		\ifnum\cotSTT=1\relax
			\def\rowHeader{\hline\rowcolor{gold1}\EZo@TF{TT}&\EZo@TF{\phatbieu}&\EZo@TF{\Dung}&\EZo@TF{\Sai}\tabularnewline}%
		\else
			\def\rowHeader{\hline\rowcolor{gold1}\EZo@TF{\phatbieu}&\EZo@TF{\Dung}&\EZo@TF{\Sai}\tabularnewline}%
		\fi
	\fi}
%--- tô xen kẽ từng hàng (gọi ngay sau \hline đầu mỗi hàng)
\newcommand{\EZmauhangTF}{%
	\noalign{\global\advance\EZhangTF\@ne}%
	\ifodd\EZhangTF\rowcolor{ezDSle}\else\rowcolor{ezDSchan}\fi}
%--- gắn vào \choiceTFt / \writeChoiceTFt của ex_test
\xpretocmd{\choiceTFt}{\EZbatdauTF}{}{\typeout{[preamble-eduzy-v2] Không gắn được bảng Đúng/Sai (đầu bảng)}}
\xpatchcmd{\choiceTFt}{\ifnum\dongPB=0\def\rowHeader{}\fi}{\EZtieudeTF\ifnum\dongPB=0\def\rowHeader{}\fi}
	{}{\typeout{[preamble-eduzy-v2] Không gắn được bảng Đúng/Sai (hàng tiêu đề)}}
\xpatchcmd{\writeChoiceTFt}{\hline\xdef\dotEX{\sdotEX}}{\hline\EZmauhangTF\xdef\dotEX{\sdotEX}}
	{}{\typeout{[preamble-eduzy-v2] Không gắn được bảng Đúng/Sai (tô hàng)}}
\xpatchcmd{\writeChoiceTFt}{{\end{longtable}}}{{\end{longtable}\EZketthucTF}}
	{}{\typeout{[preamble-eduzy-v2] Không gắn được bảng Đúng/Sai (cuối longtable)}}
\xpatchcmd{\writeChoiceTFt}{{\end{xtabular}\end{center}\medskip\noindent}}{{\end{xtabular}\end{center}\EZketthucTF\medskip\noindent}}
	{}{\typeout{[preamble-eduzy-v2] Không gắn được bảng Đúng/Sai (cuối xtabular)}}
\fi
%=====================================================================
\renewcommand{\parallel}{\mathrel{/\!\!/}}

%  trong file đề (tools/dan_trang/dan_trang.py tự ghi khi đặt \NeedspaceCau):
%    \def\TachCauDS{1}   câu Đúng/Sai được tách: câu dẫn (kèm hình) ở cuối trang, bảng Phát biểu sang trang sau
%                        (câu dẫn và hình không tách nhau, bảng không tách giữa hai trang)
%    \def\TachCauLG{1}   bản lời giải (\CheDo khác dethi): phần đề của câu cao hơn 1/3 trang được phép tách
%  Lệnh dùng trong file đề (giá trị do tools/dan_trang/dan_trang.py đo ở bản dethi):
%    \NeedspaceCau[<cao câu dẫn + hình>]{<cao cả khối câu>}   ngay trên \begin{ex} (câu đầu phần: trên tiêu đề phần)
%        công tắc tắt: đúng bằng \Needspace{<cao cả khối câu>} như cũ;
%        \TachCauDS=1 và có [..] (câu Đúng/Sai được chọn tách): chỉ giữ câu dẫn + hình;
%        \TachCauLG=1 ở bản lời giải: đòi tối đa 1/3 \textheight.
%    \NeedspaceBang{<cao bảng Phát biểu>}%   ngay trước \choiceTFt: \TachCauDS=1 thì giữ trọn bảng trên một trang;
%        công tắc tắt: không làm gì.
\providecommand{\TachCauDS}{0}
\providecommand{\TachCauLG}{0}
\newlength{\EZnsCau}
\newcommand{\NeedspaceCau}[2][]{%
	\setlength{\EZnsCau}{#2}%
	\ifnum\TachCauDS=1 \if\relax\detokenize{#1}\relax\else\setlength{\EZnsCau}{#1}\fi\fi
	\ifnum\TachCauLG=1 \ifdefstring{\kindOPTN}{dethi}{}{\ifdim\EZnsCau>0.3333\textheight\setlength{\EZnsCau}{0.3333\textheight}\fi}\fi
	\Needspace{\EZnsCau}}
\newcommand{\NeedspaceBang}[1]{\ifnum\TachCauDS=1 \Needspace{#1}\fi}

%--- Thang màu: đậm -> nhạt
\definecolor{gold1}{RGB}{126,96,0}    % nâu đồng đậm  #7E6000
\definecolor{gold2}{RGB}{192,144,3}   % vàng kim      #C09003
\definecolor{gold3}{RGB}{255,215,99}  % vàng nhạt     #FFD763
\definecolor{gold4}{RGB}{255,231,153} % vàng rất nhạt #FFE799
\definecolor{gold5}{RGB}{255,243,205} % vàng kem      #FFF3CD
%--- Ghi đè tông xanh/đỏ cũ
\definecolor{myblue}{RGB}{126,96,0}
\colorlet{blue}{gold1}
\colorlet{cyan}{gold2}
\colorlet{teal}{gold2}

%--- Tên chương mặc định (file đề nên tự khai báo lại)
\providecommand{\tenchuong}{}
%--- Chân trang
\providecommand{\thuonghieu}{TOÁN 12 -- Eduzy.com.vn}
\providecommand{\tacgia}{Tô Gia Bảo -- 0353703783}

%--- Chữ cạnh logo ở đầu trang
\providecommand{\khauhieu}{Eduzy.com.vn}

%====================== THÔNG SỐ ĐẦU TRANG, CHÂN TRANG (v2) ======================
%  Dải nền đậm tràn ngang trang; ô chữ vàng kem có cạnh xiên, kèm 2 vạch xiên chuyển màu.
%  Đầu trang lẻ : logo + \khauhieu (trái)   | ô chữ \tenchuong (phải)
%  Chân trang lẻ: ô chữ \thuonghieu (trái)  | Trang (số) (phải)
%  Trang chẵn soi gương trang lẻ; chân trang chẵn ghi \tacgia.
%--- Màu: chỉ dùng thang vàng kim Eduzy
\colorlet{ezNen}{gold1}% dải nền đậm
\colorlet{ezVach1}{gold3}% vạch xiên sát dải nền (đậm hơn)
\colorlet{ezVach2}{gold4}% vạch xiên sát ô chữ (nhạt hơn)
\colorlet{ezO}{gold5}% ô chữ vàng kem
%--- Kích thước
\newlength{\EZdai}\setlength{\EZdai}{1.1cm}% chiều cao dải nền
\newlength{\EZo}\setlength{\EZo}{0.68cm}% chiều cao ô chữ
\newlength{\EZxien}\setlength{\EZxien}{0.55cm}% độ lệch ngang của cạnh xiên
\newlength{\EZvach}\setlength{\EZvach}{0.75cm}% bề rộng mỗi vạch xiên
\newlength{\EZkhe}\setlength{\EZkhe}{0.12cm}% khe trắng giữa vạch và ô chữ
\newlength{\EZdem}\setlength{\EZdem}{0.4cm}% khoảng đệm từ chữ tới cạnh xiên
\newlength{\EZle}\setlength{\EZle}{1.8cm}% lề trái/phải của chữ (bằng lề thân bài)
\newlength{\EZrongmin}\setlength{\EZrongmin}{4.5cm}% ô chữ rộng tối thiểu
\newlength{\EZrong}% bề rộng chữ trong ô (tính tự động)
%--- Font
\newcommand{\EZfontO}{\fontsize{10.5pt}{12pt}\selectfont\bfseries}% chữ trong ô
\newcommand{\EZfontKH}{\fontfamily{qag}\fontsize{11pt}{12pt}\selectfont\bfseries}% \khauhieu cạnh logo
\newcommand{\EZfontTrang}{\fontsize{11.5pt}{12pt}\selectfont\bfseries}% chữ "Trang" và số trang
%--- Đo bề rộng chữ trong ô
\newcommand{\EZdorong}[1]{%
	\settowidth{\EZrong}{\EZfontO #1}%
	\ifdim\EZrong<\EZrongmin\setlength{\EZrong}{\EZrongmin}\fi}
%--- Logo + chữ trắng \khauhieu (dựng sẵn một lần khi vào \begin{document})
\newsavebox{\EZlogobox}
\AtBeginDocument{\sbox{\EZlogobox}{%
		\raisebox{-0.5\height}{\eduzylogo[1cm]}\hspace{0.15cm}%
		\raisebox{-0.5\height}{\color{white}\EZfontKH\khauhieu}}}
%--- Hai vạch xiên + khe trắng. #1: điểm gốc ở mép dưới, #2: hướng (1 sang phải, -1 sang trái),
%    #3: hướng nghiêng (1 kiểu "/", -1 kiểu "\")
\newcommand{\EZvachxien}[3]{%
	\fill[white] (#1) --++({(#3)*\EZxien},\EZo) --++({(#2)*\EZkhe},0) --++({-(#3)*\EZxien},-\EZo) --cycle;
	\fill[ezVach2] ([xshift=(#2)*\EZkhe]#1) --++({(#3)*\EZxien},\EZo) --++({(#2)*\EZvach},0) --++({-(#3)*\EZxien},-\EZo) --cycle;
	\fill[ezVach1] ([xshift=(#2)*(\EZkhe+\EZvach)]#1) --++({(#3)*\EZxien},\EZo) --++({(#2)*\EZvach},0) --++({-(#3)*\EZxien},-\EZo) --cycle;
	\draw[white,line width=0.6pt] ([xshift=(#2)*(\EZkhe+\EZvach)]#1) --++({(#3)*\EZxien},\EZo);
}
%--- Chữ trong ô (căn giữa theo chiều cao ô, không lệch vì dấu tiếng Việt)
\tikzset{ezchu/.style={inner sep=0pt,text=ezNen,text height=0.72em,text depth=0.22em}}

%====================== LOGO CHÌM GIỮA TRANG ======================
%  Vẽ cùng đầu trang nên nằm dưới chữ của thân bài, có trên mọi trang.
%  Tắt: \logochimtat (từ trang đó trở đi), bật lại: \logochimbat.
%  Đổi cỡ, độ đậm: khai báo trước \input, ví dụ \def\EZchimcao{9cm} \def\EZchimdo{0.1}
\providecommand{\EZchimcao}{10cm}% chiều cao logo chìm
\providecommand{\EZchimdo}{0.12}% độ đậm (0 = trong suốt, 1 = rõ hẳn)
\newif\ifEZchim \EZchimtrue
\newcommand{\logochimtat}{\EZchimfalse}
\newcommand{\logochimbat}{\EZchimtrue}
%  Dựng sẵn logo mờ một lần khi vào \begin{document}; transparency group làm mờ cả logo
%  một lượt nên không lộ vệt ở chỗ các dải màu chồng nhau.
\newsavebox{\EZchimtho}
\newsavebox{\EZchimbox}
\AtBeginDocument{%
	\sbox{\EZchimtho}{\eduzylogo[\EZchimcao]}%
	\sbox{\EZchimbox}{\begin{tikzpicture}
			\begin{scope}[transparency group,opacity=\EZchimdo]
				\node[inner sep=0pt] {\usebox{\EZchimtho}};
			\end{scope}
	\end{tikzpicture}}}
\newcommand{\EZvechim}{%
	\ifEZchim
		\node[inner sep=0pt] at (current page.center) {\usebox{\EZchimbox}};
	\fi}

%====================== ĐẦU TRANG - TRANG LẺ ======================
%  (logo bên trái, ô chữ \tenchuong bên phải)
\newcommand{\myfancyheadOdd}{%
	\EZdorong{\tenchuong}%
	\begin{tikzpicture}[remember picture,overlay]
		\coordinate (T) at (current page.north west);
		\coordinate (P) at (current page.north east);
		%--- logo chìm giữa trang
		\EZvechim
		%--- dải nền
		\fill[ezNen] (T) rectangle ([yshift=-\EZdai]P);
		%--- ô chữ: X là góc dưới bên trái, cạnh trái xiên kiểu "/"
		\coordinate (X) at ([xshift=-\EZle-\EZrong-\EZdem-\EZxien,yshift=-\EZdai]P);
		\fill[ezO] (X) --++(\EZxien,\EZo) -- ([yshift=-\EZdai+\EZo]P) -- ([yshift=-\EZdai]P) --cycle;
		\EZvachxien{X}{-1}{1}
		\node[ezchu,anchor=east] at ([xshift=-\EZle,yshift=-\EZdai+0.5*\EZo]P) {\EZfontO\tenchuong};
		%--- logo + tên
		\node[inner sep=0pt,anchor=west] at ([xshift=\EZle,yshift=-0.5*\EZdai]T) {\usebox{\EZlogobox}};
	\end{tikzpicture}%
}

%====================== ĐẦU TRANG - TRANG CHẴN =====================
%  (soi gương: ô chữ \tenchuong bên trái, logo bên phải)
\newcommand{\myfancyheadEven}{%
	\EZdorong{\tenchuong}%
	\begin{tikzpicture}[remember picture,overlay]
		\coordinate (T) at (current page.north west);
		\coordinate (P) at (current page.north east);
		%--- logo chìm giữa trang
		\EZvechim
		%--- dải nền
		\fill[ezNen] (T) rectangle ([yshift=-\EZdai]P);
		%--- ô chữ: X là góc dưới bên phải, cạnh phải xiên kiểu "\"
		\coordinate (X) at ([xshift=\EZle+\EZrong+\EZdem+\EZxien,yshift=-\EZdai]T);
		\fill[ezO] (X) --++(-\EZxien,\EZo) -- ([yshift=-\EZdai+\EZo]T) -- ([yshift=-\EZdai]T) --cycle;
		\EZvachxien{X}{1}{-1}
		\node[ezchu,anchor=west] at ([xshift=\EZle,yshift=-\EZdai+0.5*\EZo]T) {\EZfontO\tenchuong};
		%--- logo + tên
		\node[inner sep=0pt,anchor=east] at ([xshift=-\EZle,yshift=-0.5*\EZdai]P) {\usebox{\EZlogobox}};
	\end{tikzpicture}%
}

%====================== CHÂN TRANG - TRANG LẺ ======================
%  (ô chữ \thuonghieu bên trái, "Trang (số)" bên phải)
\newcommand{\myfancyfootOdd}{%
	\EZdorong{\thuonghieu}%
	\begin{tikzpicture}[remember picture,overlay]
		\coordinate (T) at (current page.south west);
		\coordinate (P) at (current page.south east);
		%--- dải nền
		\fill[ezNen] (T) rectangle ([yshift=\EZdai]P);
		%--- ô chữ sát mép trên dải: X là góc dưới bên phải, cạnh phải xiên kiểu "/"
		\coordinate (X) at ([xshift=\EZle+\EZrong+\EZdem,yshift=\EZdai-\EZo]T);
		\fill[ezO] ([yshift=\EZdai-\EZo]T) -- (X) --++(\EZxien,\EZo) -- ([yshift=\EZdai]T) --cycle;
		\EZvachxien{X}{1}{1}
		\node[ezchu,anchor=west] at ([xshift=\EZle,yshift=\EZdai-0.5*\EZo]T) {\EZfontO\thuonghieu};
		%--- số trang
		\node[circle,draw=white,line width=0.9pt,text=white,inner sep=0pt,minimum size=0.62cm,anchor=east]
		(so) at ([xshift=-\EZle,yshift=0.5*\EZdai]P) {\EZfontTrang\thepage};
		\node[inner sep=0pt,text=white,anchor=base east] (tr) at ([xshift=-0.15cm]so.base west) {\EZfontTrang Trang};
	\end{tikzpicture}%
}

%====================== CHÂN TRANG - TRANG CHẴN ====================
%  (soi gương: "Trang (số)" bên trái, ô chữ \tacgia bên phải)
\newcommand{\myfancyfootEven}{%
	\EZdorong{\tacgia}%
	\begin{tikzpicture}[remember picture,overlay]
		\coordinate (T) at (current page.south west);
		\coordinate (P) at (current page.south east);
		%--- dải nền
		\fill[ezNen] (T) rectangle ([yshift=\EZdai]P);
		%--- ô chữ sát mép trên dải: X là góc dưới bên trái, cạnh trái xiên kiểu "\"
		\coordinate (X) at ([xshift=-\EZle-\EZrong-\EZdem,yshift=\EZdai-\EZo]P);
		\fill[ezO] ([yshift=\EZdai-\EZo]P) -- (X) --++(-\EZxien,\EZo) -- ([yshift=\EZdai]P) --cycle;
		\EZvachxien{X}{-1}{-1}
		\node[ezchu,anchor=east] at ([xshift=-\EZle,yshift=\EZdai-0.5*\EZo]P) {\EZfontO\tacgia};
		%--- số trang
		%    (vòng tròn ẩn so0 nằm giữa dải để lấy đường chân chữ, giống trang lẻ)
		\node[circle,inner sep=0pt,minimum size=0.62cm] (so0) at ([yshift=0.5*\EZdai]T) {\EZfontTrang\phantom{\thepage}};
		\node[inner sep=0pt,text=white,anchor=base west] (tr) at ($(so0.base -| T)+(\EZle,0)$) {\EZfontTrang Trang};
		\node[circle,draw=white,line width=0.9pt,text=white,inner sep=0pt,minimum size=0.62cm,anchor=base west]
		(so) at ([xshift=0.15cm]tr.base east) {\EZfontTrang\thepage};
	\end{tikzpicture}%
}

%--- Giữ tương thích với các file cũ đang gọi \myfancyhead, \myfancyfoot
\newcommand{\myfancyhead}{\myfancyheadOdd}
\newcommand{\myfancyfoot}{\myfancyfootOdd}

%======================= GẮN VÀO TRANG =============================
\pagestyle{fancy}
\fancyhf{}
\fancyhead[LO]{\myfancyheadOdd}
\fancyhead[RE]{\myfancyheadEven}
\fancyfoot[LO]{\myfancyfootOdd}
\fancyfoot[RE]{\myfancyfootEven}
\renewcommand{\headrulewidth}{0pt}
\renewcommand{\footrulewidth}{0pt}
\fancypagestyle{plain}{%
	\fancyhf{}
	\fancyhead[LO]{\myfancyheadOdd}
	\fancyhead[RE]{\myfancyheadEven}
	\fancyfoot[LO]{\myfancyfootOdd}
	\fancyfoot[RE]{\myfancyfootEven}
	\renewcommand{\headrulewidth}{0pt}
	\renewcommand{\footrulewidth}{0pt}
}
%=================== HẾT PHẦN GỘP TỪ eduzy-de.tex ===================
\newcommand{\indam}[1]{
	{\fontsize{12pt}{9pt}\sffamily \selectfont\bfseries\color{gold1} #1}
}
\newcommand{\ind}[1]{
	{\fontsize{14pt}{9pt}\fontfamily{qag} \selectfont\bfseries\color{gold1} #1}
}
\renewcommand{\FalseEX}{\stepcounter{dapan}\textbf{\selectfont\color{gold1}\Alph{dapan}.}\ignorespaces}
\renewtheorem{ex}{\color{gold1}Câu}
\renewcommand{\baselinestretch}{1.1}
\newcommand{\hoac}[1]{
	\left[\begin{aligned}#1\end{aligned}\right.}
\newcommand{\heva}[1]{
	\left\{\begin{aligned}#1\end{aligned}\right.}
\usepackage{esvect}
\def\vec{\vv}
\def\vt{\vv}
\def\overrightarrow{\vv}
%====================== LOGO EDUZY VẼ BẰNG TIKZ ======================
%  Đường viền dò từ logo gốc; dải màu vàng kim đo trực tiếp trên ảnh gốc
%  (nghiêng 21 độ, tô bằng 80 dải hẹp nên hiển thị giống nhau trên mọi trình xem PDF).
%  Dùng: \eduzylogo          -> cao 2.4cm (mặc định)
%        \eduzylogo[1.5cm]   -> cao 1.5cm
\definecolor{ezText}{RGB}{245,200,46}
\newcommand{\eduzylogo}[1][2.4cm]{%
	\resizebox{!}{#1}{%
		\begin{tikzpicture}[x=0.01cm,y=0.01cm]
			\path[use as bounding box] (0,0) rectangle (224,-224);
			% mũ tốt nghiệp + vương miện
			\begin{scope}
				\clip
				(119.4,-176.1)..controls(118.9,-176.1)and(116.8,-176)..(114.6,-176)
				..controls(101.4,-176)and(85.8,-175.9)..(83.7,-175.7)
				..controls(81.8,-175.6)and(81,-175.6)..(80.2,-175.7)
				..controls(78.8,-175.9)and(77.7,-175.9)..(76.2,-175.7)
				..controls(74.9,-175.5)and(70.8,-175.5)..(64.3,-175.7)
				..controls(59.1,-175.9)and(58.3,-175.7)..(58.6,-173.7)
				..controls(58.6,-172.8)and(59.6,-169.6)..(60.3,-168)
				..controls(60.8,-166.9)and(61.4,-165.8)..(62.3,-165)
				..controls(63.1,-164.2)and(64.3,-163.4)..(65,-163.2)
				..controls(65.5,-163.1)and(72.3,-163)..(76.6,-163.1)
				..controls(77.9,-163.1)and(80.1,-163.1)..(81.4,-163)
				..controls(82.8,-162.9)and(86.3,-162.9)..(89.2,-162.9)
				..controls(93,-162.9)and(94.9,-162.8)..(96,-162.7)
				..controls(97.2,-162.5)and(101,-162.5)..(114,-162.4)
				..controls(123,-162.3)and(134.9,-162.2)..(140.2,-162.1)
				..controls(150.1,-162)and(152.2,-162.1)..(153.6,-162.4)
				..controls(154.4,-162.6)and(157.1,-163.8)..(158.1,-164.4)
				..controls(158.8,-164.8)and(159.7,-165.7)..(159.9,-166.3)
				..controls(160,-166.5)and(160.2,-167.1)..(160.2,-167.6)
				..controls(160.3,-168)and(160.6,-169.4)..(160.9,-170.6)
				..controls(161.2,-171.7)and(161.5,-173.1)..(161.6,-173.6)
				..controls(161.8,-174.4)and(161.8,-174.5)..(161.5,-174.9)
				..controls(161.2,-175.6)and(160.8,-175.9)..(159.9,-176)
				..controls(159.2,-176.1)and(154.1,-176.2)..(151.5,-176.2)
				..controls(148.1,-176.1)and(122.6,-176.1)..(121.6,-176.2)
				..controls(120.8,-176.2)and(119.8,-176.2)..(119.4,-176.1)--cycle
				(130.4,-160)..controls(129,-160)and(119.8,-160)..(110.1,-159.9)
				..controls(100.3,-159.8)and(88.5,-159.7)..(83.9,-159.6)
				..controls(77.8,-159.5)and(75.3,-159.5)..(74.5,-159.6)
				..controls(73.9,-159.7)and(72.5,-159.8)..(71.3,-159.8)
				..controls(69.6,-159.8)and(69,-159.8)..(68.5,-159.6)
				..controls(67.6,-159.3)and(66.8,-158.6)..(66.2,-157.6)
				..controls(66,-157.3)and(65.6,-156.5)..(65.3,-156.1)
				..controls(64.6,-155.1)and(64.5,-154.9)..(63.7,-152.9)
				..controls(63.2,-151.6)and(62.9,-150.9)..(61.8,-149.6)
				..controls(61.5,-149.2)and(61.1,-148.5)..(60.9,-147.9)
				..controls(60.7,-147.4)and(60.3,-146.6)..(60.1,-146.1)
				..controls(59.9,-145.5)and(59.5,-144.9)..(59,-144.3)
				..controls(58.5,-143.7)and(58.2,-143.2)..(57.9,-142.5)
				..controls(57.7,-141.9)and(57.3,-141.1)..(57,-140.7)
				..controls(56.6,-140.3)and(56.1,-139.3)..(55.8,-138.5)
				..controls(55.5,-137.8)and(55.1,-136.8)..(54.9,-136.5)
				..controls(54.7,-136.1)and(54.3,-135.2)..(53.9,-134.5)
				..controls(53.6,-133.8)and(53,-132.7)..(52.6,-132)
				..controls(51.6,-130.5)and(51.4,-130.2)..(51.1,-129)
				..controls(50.6,-127.5)and(50.6,-127.4)..(49.7,-126.3)
				..controls(49.2,-125.7)and(48.6,-124.8)..(48.3,-124.2)
				..controls(48,-123.6)and(47.5,-122.7)..(47.3,-122.2)
				..controls(47,-121.8)and(46.5,-120.7)..(46.1,-119.8)
				..controls(45.8,-118.9)and(45.2,-117.8)..(44.9,-117.2)
				..controls(44.2,-116.1)and(43.3,-114.4)..(42.3,-112.2)
				..controls(41.8,-111.3)and(41.3,-110.4)..(41.1,-110.1)
				..controls(40.5,-109.4)and(40.3,-109)..(40,-107.8)
				..controls(39.8,-107)and(39.6,-106.5)..(39.2,-106)
				..controls(38.9,-105.6)and(38.4,-104.7)..(38,-104)
				..controls(37.7,-103.4)and(37.2,-102.5)..(37,-102.1)
				..controls(36.7,-101.7)and(36.3,-100.8)..(36,-100.1)
				..controls(35.6,-98.9)and(35.3,-98.4)..(34.6,-97.6)
				..controls(34.2,-97.1)and(34.2,-96.3)..(34.6,-96)
				..controls(35,-95.7)and(35.5,-95.9)..(36.2,-96.7)
				..controls(36.5,-97)and(37.2,-97.5)..(37.6,-97.7)..controls(38,-98)and(38.6,-98.4)..(39.1,-98.9)
				..controls(39.6,-99.4)and(40.3,-100)..(40.6,-100.3)
				..controls(41.3,-100.8)and(42.4,-101.8)..(43.7,-102.9)
				..controls(44,-103.2)and(44.6,-103.7)..(45,-103.9)
				..controls(45.4,-104.2)and(45.9,-104.6)..(46.1,-104.8)
				..controls(47.1,-105.9)and(48.8,-107.1)..(50.5,-107.9)
				..controls(51.8,-108.6)and(52.3,-109)..(53.4,-110)
				..controls(54.8,-111.2)and(55.6,-111.6)..(57,-112.1)
				..controls(57.5,-112.3)and(58.5,-112.8)..(59.4,-113.3)
				..controls(60.2,-113.7)and(61.3,-114.3)..(61.9,-114.5)
				..controls(62.4,-114.7)and(63.3,-115.1)..(63.8,-115.3)
				..controls(64.4,-115.6)and(65.4,-115.9)..(66.2,-116.1)
				..controls(67,-116.3)and(68.4,-116.6)..(69.2,-116.9)
				..controls(72.3,-117.9)and(72.8,-118)..(78,-118.1)
				..controls(84,-118.3)and(86.1,-118.1)..(88.1,-117.1)
				..controls(89.4,-116.5)and(93,-115)..(93.7,-114.9)
				..controls(94.8,-114.6)and(95.4,-114.2)..(95.9,-113.4)
				..controls(96.2,-113)and(96.7,-112.5)..(97.1,-112.2)
				..controls(98,-111.6)and(101.3,-108.4)..(102.1,-107.2)
				..controls(102.4,-106.8)and(102.9,-106.2)..(103.2,-105.8)
				..controls(103.5,-105.5)and(104,-104.8)..(104.2,-104.2)
				..controls(104.5,-103.6)and(105,-102.8)..(105.3,-102.4)
				..controls(105.6,-102)and(106.1,-101.2)..(106.5,-100.6)
				..controls(107.1,-99.5)and(108.6,-97.5)..(109.5,-96.5)
				..controls(109.9,-96.1)and(110.1,-96)..(110.4,-96)
				..controls(111.1,-96)and(111.6,-96.5)..(111.8,-97.6)
				..controls(112,-98.6)and(112.3,-99.3)..(112.8,-99.9)
				..controls(113.7,-100.9)and(115.2,-103.1)..(115.7,-104.2)
				..controls(116.1,-105)and(116.5,-105.5)..(116.9,-106)
				..controls(117.3,-106.4)and(117.8,-106.9)..(118,-107.3)
				..controls(118.2,-107.6)and(118.6,-108)..(119,-108.3)
				..controls(119.6,-108.8)and(121.1,-110.3)..(121.8,-111.2)
				..controls(122.1,-111.7)and(122.5,-112)..(123.4,-112.4)
				..controls(124.8,-113.2)and(125.8,-113.8)..(127,-114.6)
				..controls(127.5,-114.9)and(128.1,-115.3)..(128.2,-115.4)
				..controls(128.4,-115.5)and(129.3,-115.8)..(130.2,-116)
				..controls(131.1,-116.2)and(132.5,-116.6)..(133.3,-116.9)
				..controls(135.7,-117.9)and(136.1,-117.9)..(142.5,-117.8)
				..controls(145.8,-117.8)and(148.4,-117.7)..(148.8,-117.6)
				..controls(149.2,-117.6)and(150.2,-117.2)..(151,-116.8)
				..controls(152,-116.4)and(152.8,-116)..(153.6,-115.9)
				..controls(154.5,-115.7)and(155.2,-115.4)..(156.3,-114.9)
				..controls(157.1,-114.5)and(158.2,-114)..(158.7,-113.8)
				..controls(159.1,-113.7)and(159.8,-113.3)..(160.2,-113)
				..controls(160.6,-112.7)and(161.3,-112.4)..(161.7,-112.2)
				..controls(162.7,-111.9)and(163.2,-111.7)..(164.1,-111.1)
				..controls(164.5,-110.8)and(165.3,-110.4)..(165.9,-110)
				..controls(166.5,-109.7)and(167.3,-109.2)..(167.6,-108.9)
				..controls(167.9,-108.6)and(168.6,-108.2)..(169,-108)
				..controls(169.4,-107.8)and(170.2,-107.3)..(170.7,-106.9)
				..controls(171.2,-106.5)and(172.1,-105.8)..(172.7,-105.4)
				..controls(173.2,-105)and(174.1,-104.4)..(174.5,-104)
				..controls(174.9,-103.7)and(175.7,-103.1)..(176.2,-102.6)
				..controls(176.7,-102.2)and(177.4,-101.5)..(177.9,-101.2)
				..controls(178.3,-100.8)and(178.8,-100.4)..(179,-100.2)
				..controls(179.2,-100)and(179.6,-99.6)..(180,-99.3)
				..controls(180.4,-99.1)and(180.9,-98.6)..(181.2,-98.4)
				..controls(181.5,-98.1)and(182.1,-97.7)..(182.5,-97.5)
				..controls(183.2,-97.2)and(183.5,-96.9)..(184.3,-95.9)
				..controls(184.6,-95.6)and(184.6,-95.6)..(184.8,-95.7)
				..controls(185,-95.8)and(185.1,-96)..(185.1,-96.2)
				..controls(185.1,-96.5)and(185,-96.7)..(184.6,-97.2)
				..controls(184.1,-97.8)and(184,-97.9)..(183.9,-98.8)
				..controls(183.7,-100.3)and(183.3,-101.3)..(182.9,-101.8)
				..controls(182.6,-102.1)and(182.3,-102.6)..(182.2,-102.9)
				..controls(182,-103.3)and(181.6,-104)..(181.2,-104.6)
				..controls(180.9,-105.2)and(180.5,-106.1)..(180.4,-106.6)
				..controls(179.9,-108)and(179.7,-108.5)..(179,-109.4)
				..controls(178.6,-109.8)and(178.2,-110.5)..(177.9,-111)
				..controls(177.7,-111.5)and(177.3,-112.3)..(177,-112.7)
				..controls(176.6,-113.2)and(176.4,-113.7)..(176.3,-114.2)
				..controls(176.2,-114.7)and(175.9,-115.2)..(175.4,-116)
				..controls(175,-116.6)and(174.4,-117.7)..(174,-118.5)
				..controls(173.7,-119.2)and(173.3,-120.1)..(173,-120.5)
				..controls(172.8,-120.9)and(172.4,-121.8)..(172.1,-122.5)
				..controls(171.8,-123.2)and(171.3,-124.1)..(171,-124.6)
				..controls(170.8,-125.1)and(170.3,-126)..(169.9,-126.6)
				..controls(169.6,-127.3)and(169.2,-128.1)..(169,-128.3)
				..controls(168.8,-128.6)and(168.5,-129.3)..(168.3,-129.8)
				..controls(168.1,-130.3)and(167.8,-131.1)..(167.5,-131.6)
				..controls(166.9,-132.5)and(164.2,-137.9)..(164.2,-138.1)
				..controls(164.2,-138.1)and(163.9,-138.8)..(163.4,-139.7)
				..controls(162.9,-140.5)and(162.4,-141.5)..(162.3,-141.9)
				..controls(161.9,-143)and(161.6,-143.4)..(160.9,-144.1)
				..controls(160.4,-144.6)and(160.3,-144.8)..(160.1,-145.4)
				..controls(160,-145.8)and(159.6,-146.8)..(159.2,-147.8)
				..controls(158.7,-148.8)and(158.1,-150)..(157.9,-150.6)
				..controls(157.6,-151.1)and(157.2,-151.9)..(156.9,-152.4)
				..controls(156.5,-152.8)and(156.1,-153.7)..(155.8,-154.4)
				..controls(155.6,-155)and(155.2,-155.9)..(154.9,-156.3)
				..controls(154.7,-156.7)and(154.4,-157.4)..(154.2,-157.9)
				..controls(154.1,-158.4)and(153.8,-159)..(153.5,-159.3)
				..controls(152.8,-160.1)and(153,-160.1)..(142.4,-160.1)
				..controls(137.2,-160.1)and(131.8,-160.1)..(130.4,-160)--cycle
				(141.5,-103.7)..controls(140.8,-103)and(139.8,-102.1)..(139.2,-101.6)
				..controls(138.5,-101.2)and(137.9,-100.6)..(137.8,-100.4)
				..controls(137,-99.5)and(134,-96.5)..(133.2,-96)
				..controls(132.7,-95.7)and(131.8,-94.9)..(131.1,-94.3)
				..controls(129.4,-92.7)and(127.9,-91.6)..(126.6,-91.1)
				..controls(125.2,-90.5)and(124.1,-90)..(122.9,-89.1)
				..controls(122.1,-88.6)and(121.3,-88.2)..(119.9,-87.7)
				..controls(118.8,-87.3)and(117.5,-86.8)..(116.9,-86.5)
				..controls(114.3,-85.5)and(113.7,-85.4)..(109.7,-85.5)
				..controls(105.3,-85.7)and(103.9,-85.9)..(101.6,-87)
				..controls(100.7,-87.4)and(99.4,-87.8)..(97.9,-88)
				..controls(97,-88.2)and(96.6,-88.3)..(95.8,-88.9)
				..controls(95.4,-89.2)and(94.7,-89.6)..(94.2,-89.8)
				..controls(92.9,-90.3)and(92.3,-90.7)..(90.9,-92.1)
				..controls(90.2,-92.7)and(89.4,-93.3)..(89,-93.5)..controls(87,-94.5)and(85.2,-96)..(83.6,-97.9)
				..controls(83.3,-98.2)and(82.5,-98.9)..(81.8,-99.4)
				..controls(80.7,-100.1)and(80.5,-100.3)..(80,-101.1)
				..controls(78.7,-103.1)and(77.2,-103.9)..(76.4,-103)
				..controls(75.7,-102.2)and(75.7,-103.2)..(75.7,-83.6)
				..controls(75.6,-73.6)and(75.6,-64.8)..(75.5,-64)
				..controls(75.5,-63.2)and(75.5,-61.4)..(75.6,-59.8)--(75.7,-57.1)--(76,-56.8)
				..controls(76.8,-56.2)and(78,-56.3)..(79.8,-57.2)
				..controls(80.4,-57.5)and(81.3,-57.8)..(81.8,-57.9)
				..controls(82.3,-58.1)and(83,-58.4)..(83.6,-58.8)..controls(84.3,-59.2)and(85,-59.5)..(86,-59.8)
				..controls(86.8,-60.1)and(88,-60.6)..(88.7,-60.9)
				..controls(89.3,-61.2)and(90.6,-61.7)..(91.4,-62)
				..controls(92.2,-62.3)and(93.7,-62.9)..(94.6,-63.3)
				..controls(95.6,-63.7)and(96.8,-64.1)..(97.4,-64.3)
				..controls(98,-64.5)and(99,-64.9)..(99.7,-65.2)
				..controls(100.4,-65.6)and(101.5,-66)..(102.3,-66.2)
				..controls(104.2,-66.8)and(105,-67.1)..(106.2,-67.8)
				..controls(107.9,-68.7)and(110.2,-69.4)..(111.3,-69.2)
				..controls(112.5,-69)and(115,-68.2)..(115.6,-67.8)
				..controls(117.3,-66.7)and(117.6,-66.6)..(119.5,-66.2)
				..controls(121.9,-65.8)and(122.1,-65.7)..(123,-65.1)
				..controls(124.1,-64.4)and(124.4,-64.3)..(126,-63.9)
				..controls(127,-63.7)and(127.7,-63.4)..(128.4,-63.1)
				..controls(129,-62.8)and(129.7,-62.5)..(130.6,-62.4)
				..controls(132.1,-62.1)and(134,-61.5)..(135.2,-61)
				..controls(135.7,-60.8)and(136.7,-60.5)..(137.3,-60.3)
				..controls(138.4,-60)and(139.3,-59.7)..(141.1,-59)
				..controls(141.6,-58.7)and(142,-58.6)..(142.5,-58.6)
				..controls(143.2,-58.6)and(143.3,-58.7)..(143.6,-59)
				..controls(144,-59.5)and(144.1,-59.9)..(144,-60.8)
				..controls(144,-61.2)and(144,-71.1)..(144,-82.8)--(143.9,-104.1)--(143.7,-104.5)
				..controls(143.5,-104.7)and(143.3,-104.8)..(143,-104.9)
				..controls(142.6,-104.9)and(142.5,-104.8)..(141.5,-103.7)--cycle
				(27.8,-91.4)..controls(27.1,-90.9)and(26.3,-90)..(26.1,-89.5)
				..controls(26,-89.2)and(25.8,-88.6)..(25.7,-87.9)
				..controls(25.4,-86.2)and(25.6,-84.8)..(26.5,-83.5)
				..controls(27.4,-82.3)and(28.3,-82)..(29.4,-82.3)
				..controls(31.1,-82.8)and(32.1,-84.9)..(32,-87.4)
				..controls(31.9,-89.3)and(31.2,-90.6)..(29.9,-91.4)
				..controls(29.1,-91.9)and(28.6,-91.9)..(27.8,-91.4)--cycle
				(190.1,-91.3)..controls(189.7,-91.2)and(189.3,-90.9)..(189.1,-90.7)
				..controls(188.3,-89.8)and(187.6,-87.7)..(187.6,-86.2)
				..controls(187.6,-85.5)and(187.7,-85.2)..(188.2,-84.4)
				..controls(189.1,-82.6)and(190.1,-81.8)..(191.3,-82.1)
				..controls(192,-82.3)and(193.5,-83.7)..(193.8,-84.4)
				..controls(194.1,-85.1)and(194.2,-87.1)..(194,-88.3)
				..controls(193.8,-89.8)and(193.3,-90.5)..(192.3,-91.3)
				..controls(191.7,-91.7)and(190.9,-91.7)..(190.1,-91.3)--cycle
				(173.7,-83.4)..controls(172.2,-82.5)and(172,-82.5)..(170.6,-82.4)
				..controls(169.7,-82.4)and(168.8,-82.4)..(168.3,-82.5)
				..controls(167.2,-82.7)and(166.9,-82.7)..(166.6,-82.2)
				..controls(166.3,-81.8)and(166.3,-81)..(166.8,-80)
				..controls(167.2,-78.9)and(167.4,-77.4)..(167.7,-74.1)
				..controls(167.8,-72.8)and(167.9,-71.8)..(167.9,-71.7)
				..controls(168,-71.7)and(168,-70.2)..(168,-68.4)..controls(168,-65.4)and(168,-64.9)..(167.7,-63)
				..controls(167.3,-60.5)and(167.1,-59.9)..(166.4,-58.4)
				..controls(166.1,-57.7)and(165.8,-56.9)..(165.7,-56.6)
				..controls(165.7,-56.2)and(165.4,-55.6)..(165.2,-55.2)
				..controls(164.6,-53.8)and(161.9,-50.7)..(161,-50.3)
				..controls(159.8,-49.7)and(157.9,-50)..(155.2,-51.2)
				..controls(154.7,-51.4)and(153.4,-51.8)..(152.2,-52.1)
				..controls(150.8,-52.4)and(149.4,-52.9)..(148.6,-53.2)
				..controls(147.8,-53.6)and(146.7,-54)..(146,-54.2)
				..controls(145.4,-54.4)and(144.5,-54.7)..(144,-54.9)
				..controls(143.6,-55.2)and(142.7,-55.5)..(141.8,-55.7)
				..controls(141,-55.9)and(139.6,-56.3)..(138.7,-56.6)
				..controls(137.8,-57)and(136.3,-57.4)..(135.3,-57.6)
				..controls(133.8,-57.9)and(133.2,-58.1)..(132.1,-58.7)
				..controls(131.1,-59.2)and(130.3,-59.4)..(129.4,-59.6)
				..controls(128.7,-59.8)and(127.8,-60)..(127.4,-60.1)
				..controls(127,-60.3)and(126,-60.5)..(125.1,-60.7)
				..controls(123.2,-61.1)and(121.7,-61.5)..(120.2,-62.1)
				..controls(119.7,-62.3)and(118.1,-62.8)..(116.7,-63.2)
				..controls(114.6,-63.9)and(114,-64.1)..(113,-64.7)
				..controls(111.9,-65.4)and(111.8,-65.4)..(111,-65.5)
				..controls(109.9,-65.5)and(109.3,-65.4)..(107.7,-64.7)
				..controls(107,-64.4)and(106,-64)..(105.3,-63.9)
				..controls(103.5,-63.5)and(102.2,-63)..(101.2,-62.5)
				..controls(100.1,-61.8)and(99.3,-61.5)..(97.8,-61.1)
				..controls(97.2,-61)and(96.2,-60.7)..(95.7,-60.4)
				..controls(95.1,-60.1)and(94.2,-59.8)..(93.5,-59.6)
				..controls(92.9,-59.4)and(91.7,-59)..(91,-58.7)..controls(90.2,-58.4)and(88.9,-58)..(88,-57.7)
				..controls(87.2,-57.5)and(86.1,-57.1)..(85.5,-56.8)
				..controls(85,-56.6)and(84.1,-56.2)..(83.5,-56)..controls(82.9,-55.7)and(81.8,-55.2)..(81,-54.9)
				..controls(80.1,-54.5)and(79.1,-54.2)..(78.7,-54.1)
				..controls(77.9,-53.8)and(77.3,-53.6)..(76.2,-52.9)
				..controls(75.7,-52.6)and(75.2,-52.4)..(74.1,-52.2)
				..controls(72.5,-51.9)and(72.4,-51.8)..(71.3,-51.1)
				..controls(69.9,-50.1)and(69.4,-49.9)..(67.9,-49.8)
				..controls(67,-49.7)and(66.6,-49.6)..(66.3,-49.4)
				..controls(66.1,-49.3)and(65.7,-49)..(65.4,-48.8)
				..controls(65.1,-48.7)and(64.2,-48.4)..(63.6,-48.2)
				..controls(62.9,-48)and(61.7,-47.6)..(60.9,-47.2)
				..controls(60.1,-46.8)and(58.9,-46.4)..(58.2,-46.2)
				..controls(57.6,-46)and(56.8,-45.7)..(56.5,-45.5)
				..controls(54.9,-44.6)and(54.6,-44.5)..(53.2,-44.2)
				..controls(51.7,-43.8)and(51,-43.6)..(49.9,-42.9)
				..controls(49.5,-42.6)and(48.9,-42.4)..(47.5,-42)
				..controls(45.7,-41.6)and(45.6,-41.5)..(44.3,-40.6)
				..controls(43.5,-40.1)and(42.8,-39.6)..(42.7,-39.4)
				..controls(42.4,-38.7)and(43,-38.1)..(44.5,-37.6)..controls(44.9,-37.5)and(46,-37)..(46.9,-36.6)
				..controls(47.9,-36.1)and(48.8,-35.8)..(49.6,-35.6)
				..controls(50.5,-35.4)and(51.3,-35.2)..(52.4,-34.6)
				..controls(53.3,-34.1)and(54.5,-33.7)..(54.9,-33.5)
				..controls(55.4,-33.4)and(56.4,-33)..(57.2,-32.6)
				..controls(58.6,-32)and(60.1,-31.5)..(61.6,-31.2)
				..controls(62.3,-31.1)and(62.7,-31)..(63.2,-30.7)
				..controls(64.1,-30)and(64.3,-29.9)..(65.3,-29.7)
				..controls(65.7,-29.6)and(66.5,-29.2)..(67.1,-28.9)
				..controls(67.6,-28.6)and(68.8,-28.2)..(69.6,-27.9)
				..controls(70.5,-27.7)and(71.6,-27.3)..(72.1,-27)
				..controls(72.6,-26.7)and(73.8,-26.2)..(74.8,-25.9)
				..controls(75.7,-25.5)and(76.7,-25.1)..(77.1,-25)
				..controls(77.4,-24.8)and(78.5,-24.4)..(79.6,-24.1)
				..controls(80.6,-23.8)and(82,-23.2)..(82.8,-22.9)
				..controls(83.5,-22.5)and(84.8,-22)..(85.8,-21.8)
				..controls(86.7,-21.5)and(87.8,-21.1)..(88.1,-20.9)
				..controls(88.4,-20.7)and(89.3,-20.3)..(90.1,-20.1)
				..controls(90.8,-19.8)and(91.9,-19.4)..(92.4,-19.1)
				..controls(93.5,-18.5)and(94.4,-18.2)..(95.4,-18.1)
				..controls(95.7,-18)and(96.4,-17.8)..(96.9,-17.5)
				..controls(97.5,-17.3)and(98.4,-16.9)..(99,-16.6)
				..controls(99.7,-16.3)and(100.9,-15.9)..(101.9,-15.6)
				..controls(102.9,-15.3)and(104.4,-14.7)..(105.3,-14.3)
				..controls(108.1,-13.2)and(108.5,-13.1)..(110.1,-13.1)
				..controls(111.4,-13.1)and(111.5,-13.1)..(112.4,-13.5)
				..controls(113.2,-14)and(115,-14.8)..(116.6,-15.4)
				..controls(116.9,-15.5)and(117.9,-15.8)..(118.8,-16.1)
				..controls(120.9,-16.6)and(121.5,-16.9)..(122.6,-17.9)
				..controls(123,-18.3)and(123.7,-18.8)..(123.9,-18.9)
				..controls(124.6,-19.4)and(126,-19.8)..(127.9,-20.2)
				..controls(129.6,-20.5)and(129.8,-20.6)..(131.1,-21.4)
				..controls(131.5,-21.7)and(132.4,-22.1)..(133.1,-22.3)
				..controls(136.1,-23.2)and(138.3,-24.1)..(140.2,-25)
				..controls(141,-25.5)and(141.5,-25.6)..(142.5,-25.8)
				..controls(145.1,-26.3)and(146,-26.6)..(149.4,-28.1)
				..controls(151.3,-28.9)and(153.5,-29.8)..(154.3,-30.1)
				..controls(155.1,-30.4)and(156.2,-30.8)..(156.6,-31)
				..controls(157.1,-31.2)and(158.2,-31.6)..(159.1,-31.8)
				..controls(160.6,-32.2)and(160.9,-32.3)..(161.7,-32.9)
				..controls(162.4,-33.3)and(162.9,-33.6)..(163.7,-33.8)
				..controls(164.3,-34)and(164.8,-34.1)..(164.9,-34.1)
				..controls(165,-34.1)and(165.6,-34.4)..(166.4,-34.8)
				..controls(167.3,-35.2)and(168.1,-35.4)..(169,-35.6)
				..controls(170,-35.8)and(170.5,-36)..(171.7,-36.6)
				..controls(172.4,-37)and(173.3,-37.4)..(173.6,-37.4)
				..controls(175.6,-37.7)and(176.2,-38)..(176.7,-39.1)
				..controls(176.9,-39.6)and(176.9,-39.7)..(176.7,-40.4)
				..controls(176.4,-41.6)and(175.8,-42.7)..(175.2,-43.4)
				..controls(174.5,-44.1)and(172.6,-45.5)..(171.7,-45.8)
				..controls(170.6,-46.2)and(169.1,-46)..(168.2,-45.4)
				..controls(167.9,-45.2)and(167.3,-44.8)..(166.9,-44.4)
				..controls(166.4,-44.1)and(165.8,-43.6)..(165.4,-43.2)
				..controls(165.1,-42.9)and(164.4,-42.3)..(163.9,-41.9)
				..controls(163.3,-41.5)and(162.7,-40.9)..(162.4,-40.5)
				..controls(161.4,-39.2)and(159.9,-38)..(158.8,-37.8)
				..controls(158.4,-37.8)and(158.3,-37.7)..(158.2,-37.2)
				..controls(158,-36.8)and(157.4,-36.1)..(157.1,-36.1)
				..controls(157,-36.1)and(156.7,-36.1)..(156.5,-36)
				..controls(156.1,-35.9)and(155.9,-35.7)..(155.7,-35.2)
				..controls(155.3,-34.6)and(155.1,-34.4)..(154,-34.1)
				..controls(153.6,-34)and(153.2,-33.7)..(152.9,-33.5)
				..controls(152.4,-32.9)and(151.6,-32.4)..(151.3,-32.4)
				..controls(150.9,-32.4)and(150.4,-32.6)..(150.4,-32.9)
				..controls(150.4,-33.2)and(150.8,-33.4)..(151.4,-33.6)
				..controls(152.1,-33.7)and(152.8,-34.1)..(153.5,-34.9)
				..controls(154.1,-35.5)and(154.3,-35.6)..(155,-35.8)
				..controls(155.7,-36)and(155.7,-36)..(156,-36.6)
				..controls(156.3,-37.3)and(156.7,-37.7)..(157.5,-37.8)
				..controls(158.1,-37.9)and(158.2,-38)..(158.4,-38.8)
				..controls(158.5,-39.4)and(158.6,-39.6)..(159.6,-40.4)
				..controls(160.2,-41)and(160.9,-41.6)..(161.3,-41.8)
				..controls(161.7,-42)and(161.9,-42.3)..(162.2,-42.8)
				..controls(162.4,-43.2)and(162.8,-43.8)..(163.1,-44.1)
				..controls(163.9,-44.8)and(165.1,-46.7)..(165.7,-48)
				..controls(166.2,-49.3)and(166.4,-49.7)..(167.2,-50.5)
				..controls(167.5,-50.8)and(167.9,-51.4)..(168.1,-51.8)
				..controls(168.3,-52.2)and(168.8,-53)..(169.1,-53.6)
				..controls(169.6,-54.3)and(169.8,-54.9)..(169.9,-55.6)
				..controls(170.4,-57.8)and(170.5,-58.1)..(171.2,-59.5)
				..controls(171.6,-60.3)and(172,-61.5)..(172.2,-62.1)
				..controls(172.3,-62.8)and(172.7,-63.7)..(172.9,-64.1)
				..controls(173.2,-64.6)and(173.4,-65.3)..(173.4,-65.6)
				..controls(173.8,-68.8)and(174.2,-70.6)..(174.9,-72.6)
				..controls(175.5,-74.5)and(175.9,-77.1)..(176.1,-80.2)
				..controls(176.3,-83.4)and(175.5,-84.4)..(173.7,-83.4)--cycle;
				\definecolor{ezs}{RGB}{247,179,26}\fill[ezs](-5.7,-156.6)--(-3.4,-157.5)--(61.9,12.7)--(59.7,13.6)--cycle;
				\definecolor{ezs}{RGB}{247,179,22}\fill[ezs](-3.8,-157.3)--(-1.5,-158.2)--(63.9,12)--(61.6,12.9)--cycle;
				\definecolor{ezs}{RGB}{248,179,19}\fill[ezs](-1.8,-158.1)--(0.4,-158.9)--(65.8,11.2)--(63.5,12.1)--cycle;
				\definecolor{ezs}{RGB}{248,178,16}\fill[ezs](0.1,-158.8)--(2.4,-159.7)--(67.7,10.5)--(65.4,11.4)--cycle;
				\definecolor{ezs}{RGB}{249,178,15}\fill[ezs](2,-159.5)--(4.3,-160.4)--(69.6,9.8)--(67.3,10.7)--cycle;
				\definecolor{ezs}{RGB}{251,177,17}\fill[ezs](3.9,-160.3)--(6.2,-161.1)--(71.5,9)--(69.2,9.9)--cycle;
				\definecolor{ezs}{RGB}{252,179,18}\fill[ezs](5.8,-161)--(8.1,-161.9)--(73.4,8.3)--(71.1,9.2)--cycle;
				\definecolor{ezs}{RGB}{252,182,17}\fill[ezs](7.7,-161.7)--(10,-162.6)--(75.3,7.6)--(73,8.5)--cycle;
				\definecolor{ezs}{RGB}{252,185,17}\fill[ezs](9.6,-162.5)--(11.9,-163.4)--(77.2,6.8)--(75,7.7)--cycle;
				\definecolor{ezs}{RGB}{251,189,18}\fill[ezs](11.5,-163.2)--(13.8,-164.1)--(79.2,6.1)--(76.9,7)--cycle;
				\definecolor{ezs}{RGB}{251,191,22}\fill[ezs](13.4,-163.9)--(15.7,-164.8)--(81.1,5.4)--(78.8,6.2)--cycle;
				\definecolor{ezs}{RGB}{252,193,25}\fill[ezs](15.4,-164.7)--(17.6,-165.6)--(83,4.6)--(80.7,5.5)--cycle;
				\definecolor{ezs}{RGB}{253,195,27}\fill[ezs](17.3,-165.4)--(19.6,-166.3)--(84.9,3.9)--(82.6,4.8)--cycle;
				\definecolor{ezs}{RGB}{253,196,31}\fill[ezs](19.2,-166.1)--(21.5,-167)--(86.8,3.2)--(84.5,4)--cycle;
				\definecolor{ezs}{RGB}{254,199,35}\fill[ezs](21.1,-166.9)--(23.4,-167.8)--(88.7,2.4)--(86.4,3.3)--cycle;
				\definecolor{ezs}{RGB}{253,201,38}\fill[ezs](23,-167.6)--(25.3,-168.5)--(90.6,1.7)--(88.3,2.6)--cycle;
				\definecolor{ezs}{RGB}{253,203,40}\fill[ezs](24.9,-168.3)--(27.2,-169.2)--(92.5,1)--(90.3,1.8)--cycle;
				\definecolor{ezs}{RGB}{253,205,42}\fill[ezs](26.8,-169.1)--(29.1,-170)--(94.4,0.2)--(92.2,1.1)--cycle;
				\definecolor{ezs}{RGB}{253,207,45}\fill[ezs](28.7,-169.8)--(31,-170.7)--(96.4,-0.5)--(94.1,0.4)--cycle;
				\definecolor{ezs}{RGB}{253,210,49}\fill[ezs](30.7,-170.5)--(32.9,-171.4)--(98.3,-1.2)--(96,-0.4)--cycle;
				\definecolor{ezs}{RGB}{253,212,54}\fill[ezs](32.6,-171.3)--(34.9,-172.2)--(100.2,-2)--(97.9,-1.1)--cycle;
				\definecolor{ezs}{RGB}{253,214,59}\fill[ezs](34.5,-172)--(36.8,-172.9)--(102.1,-2.7)--(99.8,-1.8)--cycle;
				\definecolor{ezs}{RGB}{253,216,63}\fill[ezs](36.4,-172.7)--(38.7,-173.6)--(104,-3.4)--(101.7,-2.6)--cycle;
				\definecolor{ezs}{RGB}{253,218,65}\fill[ezs](38.3,-173.5)--(40.6,-174.4)--(105.9,-4.2)--(103.6,-3.3)--cycle;
				\definecolor{ezs}{RGB}{254,220,66}\fill[ezs](40.2,-174.2)--(42.5,-175.1)--(107.8,-4.9)--(105.5,-4)--cycle;
				\definecolor{ezs}{RGB}{254,222,68}\fill[ezs](42.1,-175)--(44.4,-175.8)--(109.7,-5.6)--(107.5,-4.8)--cycle;
				\definecolor{ezs}{RGB}{254,224,71}\fill[ezs](44,-175.7)--(46.3,-176.6)--(111.7,-6.4)--(109.4,-5.5)--cycle;
				\definecolor{ezs}{RGB}{254,226,75}\fill[ezs](46,-176.4)--(48.2,-177.3)--(113.6,-7.1)--(111.3,-6.2)--cycle;
				\definecolor{ezs}{RGB}{254,228,79}\fill[ezs](47.9,-177.2)--(50.1,-178)--(115.5,-7.8)--(113.2,-7)--cycle;
				\definecolor{ezs}{RGB}{254,230,82}\fill[ezs](49.8,-177.9)--(52.1,-178.8)--(117.4,-8.6)--(115.1,-7.7)--cycle;
				\definecolor{ezs}{RGB}{254,232,85}\fill[ezs](51.7,-178.6)--(54,-179.5)--(119.3,-9.3)--(117,-8.4)--cycle;
				\definecolor{ezs}{RGB}{254,234,88}\fill[ezs](53.6,-179.4)--(55.9,-180.2)--(121.2,-10)--(118.9,-9.2)--cycle;
				\definecolor{ezs}{RGB}{254,237,91}\fill[ezs](55.5,-180.1)--(57.8,-181)--(123.1,-10.8)--(120.8,-9.9)--cycle;
				\definecolor{ezs}{RGB}{254,239,94}\fill[ezs](57.4,-180.8)--(59.7,-181.7)--(125,-11.5)--(122.8,-10.6)--cycle;
				\definecolor{ezs}{RGB}{254,241,97}\fill[ezs](59.3,-181.6)--(61.6,-182.4)--(127,-12.2)--(124.7,-11.4)--cycle;
				\definecolor{ezs}{RGB}{254,243,101}\fill[ezs](61.2,-182.3)--(63.5,-183.2)--(128.9,-13)--(126.6,-12.1)--cycle;
				\definecolor{ezs}{RGB}{254,245,103}\fill[ezs](63.2,-183)--(65.4,-183.9)--(130.8,-13.7)--(128.5,-12.8)--cycle;
				\definecolor{ezs}{RGB}{254,245,105}\fill[ezs](65.1,-183.8)--(67.4,-184.6)--(132.7,-14.4)--(130.4,-13.6)--cycle;
				\definecolor{ezs}{RGB}{254,245,105}\fill[ezs](67,-184.5)--(69.3,-185.4)--(134.6,-15.2)--(132.3,-14.3)--cycle;
				\definecolor{ezs}{RGB}{254,244,103}\fill[ezs](68.9,-185.2)--(71.2,-186.1)--(136.5,-15.9)--(134.2,-15)--cycle;
				\definecolor{ezs}{RGB}{253,243,101}\fill[ezs](70.8,-186)--(73.1,-186.8)--(138.4,-16.6)--(136.1,-15.8)--cycle;
				\definecolor{ezs}{RGB}{253,241,98}\fill[ezs](72.7,-186.7)--(75,-187.6)--(140.3,-17.4)--(138,-16.5)--cycle;
				\definecolor{ezs}{RGB}{253,239,95}\fill[ezs](74.6,-187.4)--(76.9,-188.3)--(142.2,-18.1)--(140,-17.2)--cycle;
				\definecolor{ezs}{RGB}{253,236,92}\fill[ezs](76.5,-188.2)--(78.8,-189)--(144.2,-18.8)--(141.9,-18)--cycle;
				\definecolor{ezs}{RGB}{253,234,88}\fill[ezs](78.5,-188.9)--(80.7,-189.8)--(146.1,-19.6)--(143.8,-18.7)--cycle;
				\definecolor{ezs}{RGB}{253,232,85}\fill[ezs](80.4,-189.6)--(82.7,-190.5)--(148,-20.3)--(145.7,-19.4)--cycle;
				\definecolor{ezs}{RGB}{253,230,82}\fill[ezs](82.3,-190.4)--(84.6,-191.2)--(149.9,-21)--(147.6,-20.2)--cycle;
				\definecolor{ezs}{RGB}{253,228,78}\fill[ezs](84.2,-191.1)--(86.5,-192)--(151.8,-21.8)--(149.5,-20.9)--cycle;
				\definecolor{ezs}{RGB}{253,226,74}\fill[ezs](86.1,-191.8)--(88.4,-192.7)--(153.7,-22.5)--(151.4,-21.6)--cycle;
				\definecolor{ezs}{RGB}{253,224,70}\fill[ezs](88,-192.6)--(90.3,-193.4)--(155.6,-23.3)--(153.3,-22.4)--cycle;
				\definecolor{ezs}{RGB}{253,222,67}\fill[ezs](89.9,-193.3)--(92.2,-194.2)--(157.5,-24)--(155.3,-23.1)--cycle;
				\definecolor{ezs}{RGB}{253,220,65}\fill[ezs](91.8,-194)--(94.1,-194.9)--(159.5,-24.7)--(157.2,-23.8)--cycle;
				\definecolor{ezs}{RGB}{253,217,63}\fill[ezs](93.7,-194.8)--(96,-195.6)--(161.4,-25.5)--(159.1,-24.6)--cycle;
				\definecolor{ezs}{RGB}{253,215,60}\fill[ezs](95.7,-195.5)--(97.9,-196.4)--(163.3,-26.2)--(161,-25.3)--cycle;
				\definecolor{ezs}{RGB}{253,213,56}\fill[ezs](97.6,-196.2)--(99.9,-197.1)--(165.2,-26.9)--(162.9,-26)--cycle;
				\definecolor{ezs}{RGB}{252,211,53}\fill[ezs](99.5,-197)--(101.8,-197.8)--(167.1,-27.7)--(164.8,-26.8)--cycle;
				\definecolor{ezs}{RGB}{252,209,49}\fill[ezs](101.4,-197.7)--(103.7,-198.6)--(169,-28.4)--(166.7,-27.5)--cycle;
				\definecolor{ezs}{RGB}{252,207,45}\fill[ezs](103.3,-198.4)--(105.6,-199.3)--(170.9,-29.1)--(168.6,-28.2)--cycle;
				\definecolor{ezs}{RGB}{253,205,42}\fill[ezs](105.2,-199.2)--(107.5,-200)--(172.8,-29.9)--(170.6,-29)--cycle;
				\definecolor{ezs}{RGB}{253,202,40}\fill[ezs](107.1,-199.9)--(109.4,-200.8)--(174.7,-30.6)--(172.5,-29.7)--cycle;
				\definecolor{ezs}{RGB}{253,200,38}\fill[ezs](109,-200.6)--(111.3,-201.5)--(176.7,-31.3)--(174.4,-30.4)--cycle;
				\definecolor{ezs}{RGB}{253,198,35}\fill[ezs](111,-201.4)--(113.2,-202.2)--(178.6,-32.1)--(176.3,-31.2)--cycle;
				\definecolor{ezs}{RGB}{253,196,31}\fill[ezs](112.9,-202.1)--(115.2,-203)--(180.5,-32.8)--(178.2,-31.9)--cycle;
				\definecolor{ezs}{RGB}{253,194,28}\fill[ezs](114.8,-202.8)--(117.1,-203.7)--(182.4,-33.5)--(180.1,-32.6)--cycle;
				\definecolor{ezs}{RGB}{253,192,26}\fill[ezs](116.7,-203.6)--(119,-204.5)--(184.3,-34.3)--(182,-33.4)--cycle;
				\definecolor{ezs}{RGB}{253,189,24}\fill[ezs](118.6,-204.3)--(120.9,-205.2)--(186.2,-35)--(183.9,-34.1)--cycle;
				\definecolor{ezs}{RGB}{253,186,20}\fill[ezs](120.5,-205)--(122.8,-205.9)--(188.1,-35.7)--(185.8,-34.8)--cycle;
				\definecolor{ezs}{RGB}{253,184,17}\fill[ezs](122.4,-205.8)--(124.7,-206.7)--(190,-36.5)--(187.8,-35.6)--cycle;
				\definecolor{ezs}{RGB}{252,183,15}\fill[ezs](124.3,-206.5)--(126.6,-207.4)--(192,-37.2)--(189.7,-36.3)--cycle;
				\definecolor{ezs}{RGB}{252,181,14}\fill[ezs](126.3,-207.2)--(128.5,-208.1)--(193.9,-37.9)--(191.6,-37.1)--cycle;
				\definecolor{ezs}{RGB}{252,179,13}\fill[ezs](128.2,-208)--(130.4,-208.9)--(195.8,-38.7)--(193.5,-37.8)--cycle;
				\definecolor{ezs}{RGB}{252,178,12}\fill[ezs](130.1,-208.7)--(132.4,-209.6)--(197.7,-39.4)--(195.4,-38.5)--cycle;
				\definecolor{ezs}{RGB}{251,178,10}\fill[ezs](132,-209.4)--(134.3,-210.3)--(199.6,-40.1)--(197.3,-39.3)--cycle;
				\definecolor{ezs}{RGB}{251,178,10}\fill[ezs](133.9,-210.2)--(136.2,-211.1)--(201.5,-40.9)--(199.2,-40)--cycle;
				\definecolor{ezs}{RGB}{250,177,11}\fill[ezs](135.8,-210.9)--(138.1,-211.8)--(203.4,-41.6)--(201.1,-40.7)--cycle;
				\definecolor{ezs}{RGB}{250,178,11}\fill[ezs](137.7,-211.6)--(140,-212.5)--(205.3,-42.3)--(203.1,-41.5)--cycle;
				\definecolor{ezs}{RGB}{248,180,10}\fill[ezs](139.6,-212.4)--(141.9,-213.3)--(207.3,-43.1)--(205,-42.2)--cycle;
				\definecolor{ezs}{RGB}{247,181,9}\fill[ezs](141.5,-213.1)--(143.8,-214)--(209.2,-43.8)--(206.9,-42.9)--cycle;
				\definecolor{ezs}{RGB}{246,182,9}\fill[ezs](143.5,-213.8)--(145.7,-214.7)--(211.1,-44.5)--(208.8,-43.7)--cycle;
				\definecolor{ezs}{RGB}{246,182,10}\fill[ezs](145.4,-214.6)--(147.7,-215.5)--(213,-45.3)--(210.7,-44.4)--cycle;
			\end{scope}
			% chữ EDUZY (Times đậm, làm dày nét bằng 5 lớp lệch nhẹ; căn đúng khung chữ logo gốc)
			\foreach \dx/\dy in {-0.7/0,0.7/0,0/-0.3,0/0.3,0/0}
			\node[anchor=center,inner sep=0pt,outer sep=0pt,text=ezText,align=center] at (109.6+\dx,-193.6+\dy)
			{\resizebox{1.01cm}{0.22cm}{\fontencoding{T1}\fontfamily{ptm}\fontseries{b}\selectfont EDUZY}};
		\end{tikzpicture}%
	}%
}
%====================== THIẾT LẬP THÂN ĐỀ ======================
%--- Môi trường tự luận {bt}: ex_test 3.3.4 chưa khai báo sẵn nên tạo ở đây
%    (ex_test đã tự móc \fixshowans và \listenumerate cho bt)
\makeatletter
\@ifundefined{bt}{\newtheorem{bt}{\color{gold1}\namebt}}{}
\makeatother
\renewcommand{\sectionmark}[1]{\markright{\thesection\ #1}}
%--- Khung tiêu đề (logo + tối đa 4 dòng; dòng 3, 4 để trống {} thì bỏ qua)
%  Đề thi:    \khungtieude{THI THỬ THPT QUỐC GIA NĂM 2026}{Môn Toán -- Đề số 21}
%                         {(Đề thi gồm: ...)}{Thời gian làm bài: 90 phút (không kể thời gian phát đề).}
%  Chuyên đề: \khungtieude{CHUYÊN ĐỀ 3 -- DÃY SỐ VÀ CÔNG THỨC TRUY HỒI}
%                         {Bồi dưỡng học sinh giỏi Toán THPT}{...}{}
\newcommand{\khungtieude}[4]{%
	\begin{tcolorbox}[colframe=gold1,colback=gold5!45,boxrule=0.5mm]
		\begin{minipage}[c]{3cm}
			\centerline{\eduzylogo[2.4cm]}
		\end{minipage}\hspace{0.5cm}
		\begin{minipage}[c]{12.5cm}
			\vskip 0.15 cm
			\centerline{\ind{#1}}
			\vskip 0.15 cm
			\centerline{\textbf{#2}}
			\if\relax\detokenize{#3}\relax\else
				\vskip 0.15 cm
				\centerline{\textit{#3}}
			\fi
			\if\relax\detokenize{#4}\relax\else
				\vskip 0.15 cm
				\centerline{\textit{#4}}
			\fi
			\vskip 0.15 cm
		\end{minipage}
	\end{tcolorbox}
}
%====================== BỔ SUNG CHO BÀI HỌC, CHUYÊN ĐỀ ======================
\usepackage[shortlabels]{enumitem}
\usepackage{tabularx}
%--- Danh sách gọn: \begin{enumerate}[gon] ... \end{enumerate}
\SetEnumitemKey{gon}{leftmargin=*,itemsep=1pt,parsep=0pt,topsep=3pt}
%--- Khung lý thuyết: \begin{khung}{} ... \end{khung}; tiêu đề chỉ dùng khi thật ngắn ("Mẹo", "Lỗi thường gặp")
\newtcolorbox{khung}[1]{breakable,enhanced,colframe=gold1,colback=gold5!45,
	boxrule=0.4mm,fonttitle=\bfseries\sffamily,coltitle=white,colbacktitle=gold1,
	attach boxed title to top left={xshift=6mm,yshift=-3mm},
	boxed title style={arc=1mm,boxrule=0.3mm},title={#1}}
%--- Tiêu đề mục nhỏ trong một dạng: \muc{Dấu hiệu}, \muc{Phương pháp}, \muc{Bài tự luyện}
\makeatletter
\newcommand{\muc}[1]{\par\medskip\noindent\indam{#1}\par\nobreak\smallskip\@afterindentfalse\@afterheading}
\makeatother
%--- Ô công thức cần nhớ (căn giữa, viền vàng, vừa khít nội dung)
%     \begin{center}\congthuc{$Q(n)=(n-1)!$}\end{center}
\newtcbox{\congthuc}{on line,colframe=gold1,colback=gold5!45,boxrule=0.5pt,arc=2pt,
	left=8pt,right=8pt,top=4pt,bottom=4pt}
%--- Công tắc lời giải bài tự luyện (ví dụ mẫu luôn dùng \loigiai)
%     \loigiaiTL[{đáp số}]{lời giải đầy đủ}
%     \TLfull -> in lời giải đầy đủ;  \TLdapso -> chỉ in một dòng "Đáp số";  \TLan -> ẩn hoàn toàn
%     (\CheDo là dethi thì mọi lời giải và đáp số đều ẩn)
\makeatletter
\newif\ifTLhien   \TLhientrue
\newif\ifTLdapso  \TLdapsofalse
\newcommand{\TLfull}{\TLhientrue\TLdapsofalse}
\newcommand{\TLdapso}{\TLhientrue\TLdapsotrue}
\newcommand{\TLan}{\TLhienfalse}
\newcommand{\loigiaiTL}[2][]{%
	\ifTLhien
		\ifTLdapso
			\def\@tlarg{#1}%
			\ifx\@tlarg\@empty\else
				\ifdefstring{\kindOPTN}{dethi}{}{\par\noindent{\bfseries\color{gold1}Đáp số:} #1\par}%
			\fi
		\else
			\loigiai{#2}%
		\fi
	\fi}
\makeatother
\TLdapso% mặc định: chỉ in đáp số bài tự luyện
%====================== BẢNG ĐÚNG/SAI (\choiceTFt) ======================
%  Hàng tiêu đề nền nâu đồng chữ trắng (Phát biểu | Đ | S), các hàng tô xen kẽ kem / trắng,
%  đường kẻ vàng kim, nhãn a) b) c) d) màu nâu đồng. Bản có đáp án: cột được chọn có vòng tròn
%  nâu đồng với dấu tích trắng (câu đúng tích ở cột Đúng, câu sai tích ở cột Sai).
%  Tiêu đề ghi đủ "Đúng | Sai": \OPTN{viettat=0} sau \input (cột Đúng/Sai rộng ra, phát biểu hẹp lại).
%  Dùng bảng gốc của ex_test: \def\BangDSgoc{1} trước %============================================================
%  PREAMBLE EDUZY v2 - dùng chung cho đề thi, bài học, chuyên đề (gói ex_test)
%  Khác bản preamble-eduzy.tex ở đầu trang, chân trang: dải nền nâu đồng tràn ngang,
%  logo + Eduzy.com.vn, ô chữ vàng kem cạnh xiên, chân trang "Trang (số)";
%  logo Eduzy chìm giữa mọi trang (tắt: \logochimtat, bật lại: \logochimbat).
%  Cách dùng ở file chính (xem mau-chuan.tex, mau-chuyen-de.tex):
%      \def\CheDo{loigiai}          % loigiai | dethi | solcolor | color | book
%      \def\tenchuong{<dòng chữ trên đầu trang>}
%      \def\thuonghieu{<chữ ở chân trang lẻ>}   % không khai báo thì dùng mặc định
%      \def\tacgia{<chữ ở chân trang chẵn>}      % không khai báo thì dùng mặc định
%      \def\khauhieu{<chữ cạnh logo ở đầu trang>}  % tùy chọn, mặc định Eduzy.com.vn
%      %============================================================
%  PREAMBLE EDUZY v2 - dùng chung cho đề thi, bài học, chuyên đề (gói ex_test)
%  Khác bản preamble-eduzy.tex ở đầu trang, chân trang: dải nền nâu đồng tràn ngang,
%  logo + Eduzy.com.vn, ô chữ vàng kem cạnh xiên, chân trang "Trang (số)";
%  logo Eduzy chìm giữa mọi trang (tắt: \logochimtat, bật lại: \logochimbat).
%  Cách dùng ở file chính (xem mau-chuan.tex, mau-chuyen-de.tex):
%      \def\CheDo{loigiai}          % loigiai | dethi | solcolor | color | book
%      \def\tenchuong{<dòng chữ trên đầu trang>}
%      \def\thuonghieu{<chữ ở chân trang lẻ>}   % không khai báo thì dùng mặc định
%      \def\tacgia{<chữ ở chân trang chẵn>}      % không khai báo thì dùng mặc định
%      \def\khauhieu{<chữ cạnh logo ở đầu trang>}  % tùy chọn, mặc định Eduzy.com.vn
%      \input{preamble-eduzy-v2}
%      \TLdapso                      % tài liệu: \TLfull | \TLdapso | \TLan
%      \begin{document} ... \end{document}
%  Logo Eduzy vẽ bằng TikZ ngay trong file (\eduzylogo), không cần file ảnh.
%  Biên dịch: pdfLaTeX (máy có gói vietnam) hoặc XeLaTeX.
%============================================================
\documentclass[12pt,a4paper,twoside]{article}%<-- twoside: bat dau trang/chan trang chan-le
\usepackage{iftex}
\usepackage{amsmath,amssymb,tikz,mathrsfs,tkz-tab}
\usepackage{graphicx}
\usepackage{pgfplots}% đồ thị hàm số
\pgfplotsset{compat=1.17}
%\usepackage{mathptmx}
\usepackage{mathpazo}
\ifXeTeX% XeLaTeX: font chữ có đủ dấu tiếng Việt, giữ nguyên font toán mathpazo
	\usepackage[no-math]{fontspec}
	\setmainfont{TeX Gyre Pagella}
	\newfontfamily\EZqag{TeX Gyre Adventor}[NFSSFamily=qag,Ligatures=TeX]
\fi
\usepackage[many]{tcolorbox}
\usetikzlibrary{math,through,calc,intersections,angles,quotes,shapes,shapes.geometric,arrows,patterns,matrix,chains,arrows.meta,decorations.shapes,decorations.fractals,decorations.markings,shadows,decorations.pathreplacing}
\usepackage[top=1.6cm, bottom=1.8cm, left=1.8cm, right=1.8cm]{geometry}% v2: top 1.6cm (bản cũ 1.5cm) để chữ không sát dải đầu trang
\providecommand{\CheDo}{loigiai}% loigiai | dethi | solcolor | color | book
\expandafter\usepackage\expandafter[\CheDo]{ex_test}
% \tenchuong: khai báo ở file chính (trước \input) hoặc \renewcommand sau \input
%============================================================
%  [GỘP TỪ eduzy-de.tex] EDUZY - MÀU NÂU ĐỒNG / VÀNG KIM
%  + ĐẦU TRANG, CHÂN TRANG KIỂU DẢI NGANG (v2), ĐỐI XỨNG TRANG CHẴN / TRANG LẺ.
%  (không cần file eduzy-de.tex riêng nữa)
%  Khối này phải đặt SAU \usepackage[...]{ex_test} và SAU \def\tenchuong{...}
%============================================================
\usepackage{lastpage}
\usepackage{needspace}
%--- DÀN TRANG ĐỀ THI: tách câu (THÊM 05/11/2026; chỉ thêm lệnh mới, KHÔNG đổi hành vi mặc định) ---
%  Hai công tắc, mặc định tắt (= preamble chạy y như cũ). Bật bằng cách đặt TRƯỚC \input{preamble-eduzy-v2}
%  trong file đề (tools/dan_trang/dan_trang.py tự ghi khi đặt \NeedspaceCau):
%    \def\TachCauDS{1}   câu Đúng/Sai được tách: câu dẫn (kèm hình) ở cuối trang, bảng Phát biểu sang trang sau
%                        (câu dẫn và hình không tách nhau, bảng không tách giữa hai trang)
%    \def\TachCauLG{1}   bản lời giải (\CheDo khác dethi): phần đề của câu cao hơn 1/3 trang được phép tách
%  Lệnh dùng trong file đề (giá trị do tools/dan_trang/dan_trang.py đo ở bản dethi):
%    \NeedspaceCau[<cao câu dẫn + hình>]{<cao cả khối câu>}   ngay trên \begin{ex} (câu đầu phần: trên tiêu đề phần)
%        công tắc tắt: đúng bằng \Needspace{<cao cả khối câu>} như cũ;
%        \TachCauDS=1 và có [..] (câu Đúng/Sai được chọn tách): chỉ giữ câu dẫn + hình;
%        \TachCauLG=1 ở bản lời giải: đòi tối đa 1/3 \textheight.
%    \NeedspaceBang{<cao bảng Phát biểu>}%   ngay trước \choiceTFt: \TachCauDS=1 thì giữ trọn bảng trên một trang;
%        công tắc tắt: không làm gì.
\providecommand{\TachCauDS}{0}
\providecommand{\TachCauLG}{0}
\newlength{\EZnsCau}
\newcommand{\NeedspaceCau}[2][]{%
	\setlength{\EZnsCau}{#2}%
	\ifnum\TachCauDS=1 \if\relax\detokenize{#1}\relax\else\setlength{\EZnsCau}{#1}\fi\fi
	\ifnum\TachCauLG=1 \ifdefstring{\kindOPTN}{dethi}{}{\ifdim\EZnsCau>0.3333\textheight\setlength{\EZnsCau}{0.3333\textheight}\fi}\fi
	\Needspace{\EZnsCau}}
\newcommand{\NeedspaceBang}[1]{\ifnum\TachCauDS=1 \Needspace{#1}\fi}

%--- Thang màu: đậm -> nhạt
\definecolor{gold1}{RGB}{126,96,0}    % nâu đồng đậm  #7E6000
\definecolor{gold2}{RGB}{192,144,3}   % vàng kim      #C09003
\definecolor{gold3}{RGB}{255,215,99}  % vàng nhạt     #FFD763
\definecolor{gold4}{RGB}{255,231,153} % vàng rất nhạt #FFE799
\definecolor{gold5}{RGB}{255,243,205} % vàng kem      #FFF3CD
%--- Ghi đè tông xanh/đỏ cũ
\definecolor{myblue}{RGB}{126,96,0}
\colorlet{blue}{gold1}
\colorlet{cyan}{gold2}
\colorlet{teal}{gold2}

%--- Tên chương mặc định (file đề nên tự khai báo lại)
\providecommand{\tenchuong}{}
%--- Chân trang
\providecommand{\thuonghieu}{TOÁN 12 -- Eduzy.com.vn}
\providecommand{\tacgia}{Tô Gia Bảo -- 0353703783}

%--- Chữ cạnh logo ở đầu trang
\providecommand{\khauhieu}{Eduzy.com.vn}

%====================== THÔNG SỐ ĐẦU TRANG, CHÂN TRANG (v2) ======================
%  Dải nền đậm tràn ngang trang; ô chữ vàng kem có cạnh xiên, kèm 2 vạch xiên chuyển màu.
%  Đầu trang lẻ : logo + \khauhieu (trái)   | ô chữ \tenchuong (phải)
%  Chân trang lẻ: ô chữ \thuonghieu (trái)  | Trang (số) (phải)
%  Trang chẵn soi gương trang lẻ; chân trang chẵn ghi \tacgia.
%--- Màu: chỉ dùng thang vàng kim Eduzy
\colorlet{ezNen}{gold1}% dải nền đậm
\colorlet{ezVach1}{gold3}% vạch xiên sát dải nền (đậm hơn)
\colorlet{ezVach2}{gold4}% vạch xiên sát ô chữ (nhạt hơn)
\colorlet{ezO}{gold5}% ô chữ vàng kem
%--- Kích thước
\newlength{\EZdai}\setlength{\EZdai}{1.1cm}% chiều cao dải nền
\newlength{\EZo}\setlength{\EZo}{0.68cm}% chiều cao ô chữ
\newlength{\EZxien}\setlength{\EZxien}{0.55cm}% độ lệch ngang của cạnh xiên
\newlength{\EZvach}\setlength{\EZvach}{0.75cm}% bề rộng mỗi vạch xiên
\newlength{\EZkhe}\setlength{\EZkhe}{0.12cm}% khe trắng giữa vạch và ô chữ
\newlength{\EZdem}\setlength{\EZdem}{0.4cm}% khoảng đệm từ chữ tới cạnh xiên
\newlength{\EZle}\setlength{\EZle}{1.8cm}% lề trái/phải của chữ (bằng lề thân bài)
\newlength{\EZrongmin}\setlength{\EZrongmin}{4.5cm}% ô chữ rộng tối thiểu
\newlength{\EZrong}% bề rộng chữ trong ô (tính tự động)
%--- Font
\newcommand{\EZfontO}{\fontsize{10.5pt}{12pt}\selectfont\bfseries}% chữ trong ô
\newcommand{\EZfontKH}{\fontfamily{qag}\fontsize{11pt}{12pt}\selectfont\bfseries}% \khauhieu cạnh logo
\newcommand{\EZfontTrang}{\fontsize{11.5pt}{12pt}\selectfont\bfseries}% chữ "Trang" và số trang
%--- Đo bề rộng chữ trong ô
\newcommand{\EZdorong}[1]{%
	\settowidth{\EZrong}{\EZfontO #1}%
	\ifdim\EZrong<\EZrongmin\setlength{\EZrong}{\EZrongmin}\fi}
%--- Logo + chữ trắng \khauhieu (dựng sẵn một lần khi vào \begin{document})
\newsavebox{\EZlogobox}
\AtBeginDocument{\sbox{\EZlogobox}{%
		\raisebox{-0.5\height}{\eduzylogo[1cm]}\hspace{0.15cm}%
		\raisebox{-0.5\height}{\color{white}\EZfontKH\khauhieu}}}
%--- Hai vạch xiên + khe trắng. #1: điểm gốc ở mép dưới, #2: hướng (1 sang phải, -1 sang trái),
%    #3: hướng nghiêng (1 kiểu "/", -1 kiểu "\")
\newcommand{\EZvachxien}[3]{%
	\fill[white] (#1) --++({(#3)*\EZxien},\EZo) --++({(#2)*\EZkhe},0) --++({-(#3)*\EZxien},-\EZo) --cycle;
	\fill[ezVach2] ([xshift=(#2)*\EZkhe]#1) --++({(#3)*\EZxien},\EZo) --++({(#2)*\EZvach},0) --++({-(#3)*\EZxien},-\EZo) --cycle;
	\fill[ezVach1] ([xshift=(#2)*(\EZkhe+\EZvach)]#1) --++({(#3)*\EZxien},\EZo) --++({(#2)*\EZvach},0) --++({-(#3)*\EZxien},-\EZo) --cycle;
	\draw[white,line width=0.6pt] ([xshift=(#2)*(\EZkhe+\EZvach)]#1) --++({(#3)*\EZxien},\EZo);
}
%--- Chữ trong ô (căn giữa theo chiều cao ô, không lệch vì dấu tiếng Việt)
\tikzset{ezchu/.style={inner sep=0pt,text=ezNen,text height=0.72em,text depth=0.22em}}

%====================== LOGO CHÌM GIỮA TRANG ======================
%  Vẽ cùng đầu trang nên nằm dưới chữ của thân bài, có trên mọi trang.
%  Tắt: \logochimtat (từ trang đó trở đi), bật lại: \logochimbat.
%  Đổi cỡ, độ đậm: khai báo trước \input, ví dụ \def\EZchimcao{9cm} \def\EZchimdo{0.1}
\providecommand{\EZchimcao}{10cm}% chiều cao logo chìm
\providecommand{\EZchimdo}{0.12}% độ đậm (0 = trong suốt, 1 = rõ hẳn)
\newif\ifEZchim \EZchimtrue
\newcommand{\logochimtat}{\EZchimfalse}
\newcommand{\logochimbat}{\EZchimtrue}
%  Dựng sẵn logo mờ một lần khi vào \begin{document}; transparency group làm mờ cả logo
%  một lượt nên không lộ vệt ở chỗ các dải màu chồng nhau.
\newsavebox{\EZchimtho}
\newsavebox{\EZchimbox}
\AtBeginDocument{%
	\sbox{\EZchimtho}{\eduzylogo[\EZchimcao]}%
	\sbox{\EZchimbox}{\begin{tikzpicture}
			\begin{scope}[transparency group,opacity=\EZchimdo]
				\node[inner sep=0pt] {\usebox{\EZchimtho}};
			\end{scope}
	\end{tikzpicture}}}
\newcommand{\EZvechim}{%
	\ifEZchim
		\node[inner sep=0pt] at (current page.center) {\usebox{\EZchimbox}};
	\fi}

%====================== ĐẦU TRANG - TRANG LẺ ======================
%  (logo bên trái, ô chữ \tenchuong bên phải)
\newcommand{\myfancyheadOdd}{%
	\EZdorong{\tenchuong}%
	\begin{tikzpicture}[remember picture,overlay]
		\coordinate (T) at (current page.north west);
		\coordinate (P) at (current page.north east);
		%--- logo chìm giữa trang
		\EZvechim
		%--- dải nền
		\fill[ezNen] (T) rectangle ([yshift=-\EZdai]P);
		%--- ô chữ: X là góc dưới bên trái, cạnh trái xiên kiểu "/"
		\coordinate (X) at ([xshift=-\EZle-\EZrong-\EZdem-\EZxien,yshift=-\EZdai]P);
		\fill[ezO] (X) --++(\EZxien,\EZo) -- ([yshift=-\EZdai+\EZo]P) -- ([yshift=-\EZdai]P) --cycle;
		\EZvachxien{X}{-1}{1}
		\node[ezchu,anchor=east] at ([xshift=-\EZle,yshift=-\EZdai+0.5*\EZo]P) {\EZfontO\tenchuong};
		%--- logo + tên
		\node[inner sep=0pt,anchor=west] at ([xshift=\EZle,yshift=-0.5*\EZdai]T) {\usebox{\EZlogobox}};
	\end{tikzpicture}%
}

%====================== ĐẦU TRANG - TRANG CHẴN =====================
%  (soi gương: ô chữ \tenchuong bên trái, logo bên phải)
\newcommand{\myfancyheadEven}{%
	\EZdorong{\tenchuong}%
	\begin{tikzpicture}[remember picture,overlay]
		\coordinate (T) at (current page.north west);
		\coordinate (P) at (current page.north east);
		%--- logo chìm giữa trang
		\EZvechim
		%--- dải nền
		\fill[ezNen] (T) rectangle ([yshift=-\EZdai]P);
		%--- ô chữ: X là góc dưới bên phải, cạnh phải xiên kiểu "\"
		\coordinate (X) at ([xshift=\EZle+\EZrong+\EZdem+\EZxien,yshift=-\EZdai]T);
		\fill[ezO] (X) --++(-\EZxien,\EZo) -- ([yshift=-\EZdai+\EZo]T) -- ([yshift=-\EZdai]T) --cycle;
		\EZvachxien{X}{1}{-1}
		\node[ezchu,anchor=west] at ([xshift=\EZle,yshift=-\EZdai+0.5*\EZo]T) {\EZfontO\tenchuong};
		%--- logo + tên
		\node[inner sep=0pt,anchor=east] at ([xshift=-\EZle,yshift=-0.5*\EZdai]P) {\usebox{\EZlogobox}};
	\end{tikzpicture}%
}

%====================== CHÂN TRANG - TRANG LẺ ======================
%  (ô chữ \thuonghieu bên trái, "Trang (số)" bên phải)
\newcommand{\myfancyfootOdd}{%
	\EZdorong{\thuonghieu}%
	\begin{tikzpicture}[remember picture,overlay]
		\coordinate (T) at (current page.south west);
		\coordinate (P) at (current page.south east);
		%--- dải nền
		\fill[ezNen] (T) rectangle ([yshift=\EZdai]P);
		%--- ô chữ sát mép trên dải: X là góc dưới bên phải, cạnh phải xiên kiểu "/"
		\coordinate (X) at ([xshift=\EZle+\EZrong+\EZdem,yshift=\EZdai-\EZo]T);
		\fill[ezO] ([yshift=\EZdai-\EZo]T) -- (X) --++(\EZxien,\EZo) -- ([yshift=\EZdai]T) --cycle;
		\EZvachxien{X}{1}{1}
		\node[ezchu,anchor=west] at ([xshift=\EZle,yshift=\EZdai-0.5*\EZo]T) {\EZfontO\thuonghieu};
		%--- số trang
		\node[circle,draw=white,line width=0.9pt,text=white,inner sep=0pt,minimum size=0.62cm,anchor=east]
		(so) at ([xshift=-\EZle,yshift=0.5*\EZdai]P) {\EZfontTrang\thepage};
		\node[inner sep=0pt,text=white,anchor=base east] (tr) at ([xshift=-0.15cm]so.base west) {\EZfontTrang Trang};
	\end{tikzpicture}%
}

%====================== CHÂN TRANG - TRANG CHẴN ====================
%  (soi gương: "Trang (số)" bên trái, ô chữ \tacgia bên phải)
\newcommand{\myfancyfootEven}{%
	\EZdorong{\tacgia}%
	\begin{tikzpicture}[remember picture,overlay]
		\coordinate (T) at (current page.south west);
		\coordinate (P) at (current page.south east);
		%--- dải nền
		\fill[ezNen] (T) rectangle ([yshift=\EZdai]P);
		%--- ô chữ sát mép trên dải: X là góc dưới bên trái, cạnh trái xiên kiểu "\"
		\coordinate (X) at ([xshift=-\EZle-\EZrong-\EZdem,yshift=\EZdai-\EZo]P);
		\fill[ezO] ([yshift=\EZdai-\EZo]P) -- (X) --++(-\EZxien,\EZo) -- ([yshift=\EZdai]P) --cycle;
		\EZvachxien{X}{-1}{-1}
		\node[ezchu,anchor=east] at ([xshift=-\EZle,yshift=\EZdai-0.5*\EZo]P) {\EZfontO\tacgia};
		%--- số trang
		%    (vòng tròn ẩn so0 nằm giữa dải để lấy đường chân chữ, giống trang lẻ)
		\node[circle,inner sep=0pt,minimum size=0.62cm] (so0) at ([yshift=0.5*\EZdai]T) {\EZfontTrang\phantom{\thepage}};
		\node[inner sep=0pt,text=white,anchor=base west] (tr) at ($(so0.base -| T)+(\EZle,0)$) {\EZfontTrang Trang};
		\node[circle,draw=white,line width=0.9pt,text=white,inner sep=0pt,minimum size=0.62cm,anchor=base west]
		(so) at ([xshift=0.15cm]tr.base east) {\EZfontTrang\thepage};
	\end{tikzpicture}%
}

%--- Giữ tương thích với các file cũ đang gọi \myfancyhead, \myfancyfoot
\newcommand{\myfancyhead}{\myfancyheadOdd}
\newcommand{\myfancyfoot}{\myfancyfootOdd}

%======================= GẮN VÀO TRANG =============================
\pagestyle{fancy}
\fancyhf{}
\fancyhead[LO]{\myfancyheadOdd}
\fancyhead[RE]{\myfancyheadEven}
\fancyfoot[LO]{\myfancyfootOdd}
\fancyfoot[RE]{\myfancyfootEven}
\renewcommand{\headrulewidth}{0pt}
\renewcommand{\footrulewidth}{0pt}
\fancypagestyle{plain}{%
	\fancyhf{}
	\fancyhead[LO]{\myfancyheadOdd}
	\fancyhead[RE]{\myfancyheadEven}
	\fancyfoot[LO]{\myfancyfootOdd}
	\fancyfoot[RE]{\myfancyfootEven}
	\renewcommand{\headrulewidth}{0pt}
	\renewcommand{\footrulewidth}{0pt}
}
%=================== HẾT PHẦN GỘP TỪ eduzy-de.tex ===================
\newcommand{\indam}[1]{
	{\fontsize{12pt}{9pt}\sffamily \selectfont\bfseries\color{gold1} #1}
}
\newcommand{\ind}[1]{
	{\fontsize{14pt}{9pt}\fontfamily{qag} \selectfont\bfseries\color{gold1} #1}
}
\renewcommand{\FalseEX}{\stepcounter{dapan}\textbf{\selectfont\color{gold1}\Alph{dapan}.}\ignorespaces}
\renewtheorem{ex}{\color{gold1}Câu}
\renewcommand{\baselinestretch}{1.1}
\newcommand{\hoac}[1]{
	\left[\begin{aligned}#1\end{aligned}\right.}
\newcommand{\heva}[1]{
	\left\{\begin{aligned}#1\end{aligned}\right.}
\usepackage{esvect}
\def\vec{\vv}
\def\vt{\vv}
\def\overrightarrow{\vv}
%====================== LOGO EDUZY VẼ BẰNG TIKZ ======================
%  Đường viền dò từ logo gốc; dải màu vàng kim đo trực tiếp trên ảnh gốc
%  (nghiêng 21 độ, tô bằng 80 dải hẹp nên hiển thị giống nhau trên mọi trình xem PDF).
%  Dùng: \eduzylogo          -> cao 2.4cm (mặc định)
%        \eduzylogo[1.5cm]   -> cao 1.5cm
\definecolor{ezText}{RGB}{245,200,46}
\newcommand{\eduzylogo}[1][2.4cm]{%
	\resizebox{!}{#1}{%
		\begin{tikzpicture}[x=0.01cm,y=0.01cm]
			\path[use as bounding box] (0,0) rectangle (224,-224);
			% mũ tốt nghiệp + vương miện
			\begin{scope}
				\clip
				(119.4,-176.1)..controls(118.9,-176.1)and(116.8,-176)..(114.6,-176)
				..controls(101.4,-176)and(85.8,-175.9)..(83.7,-175.7)
				..controls(81.8,-175.6)and(81,-175.6)..(80.2,-175.7)
				..controls(78.8,-175.9)and(77.7,-175.9)..(76.2,-175.7)
				..controls(74.9,-175.5)and(70.8,-175.5)..(64.3,-175.7)
				..controls(59.1,-175.9)and(58.3,-175.7)..(58.6,-173.7)
				..controls(58.6,-172.8)and(59.6,-169.6)..(60.3,-168)
				..controls(60.8,-166.9)and(61.4,-165.8)..(62.3,-165)
				..controls(63.1,-164.2)and(64.3,-163.4)..(65,-163.2)
				..controls(65.5,-163.1)and(72.3,-163)..(76.6,-163.1)
				..controls(77.9,-163.1)and(80.1,-163.1)..(81.4,-163)
				..controls(82.8,-162.9)and(86.3,-162.9)..(89.2,-162.9)
				..controls(93,-162.9)and(94.9,-162.8)..(96,-162.7)
				..controls(97.2,-162.5)and(101,-162.5)..(114,-162.4)
				..controls(123,-162.3)and(134.9,-162.2)..(140.2,-162.1)
				..controls(150.1,-162)and(152.2,-162.1)..(153.6,-162.4)
				..controls(154.4,-162.6)and(157.1,-163.8)..(158.1,-164.4)
				..controls(158.8,-164.8)and(159.7,-165.7)..(159.9,-166.3)
				..controls(160,-166.5)and(160.2,-167.1)..(160.2,-167.6)
				..controls(160.3,-168)and(160.6,-169.4)..(160.9,-170.6)
				..controls(161.2,-171.7)and(161.5,-173.1)..(161.6,-173.6)
				..controls(161.8,-174.4)and(161.8,-174.5)..(161.5,-174.9)
				..controls(161.2,-175.6)and(160.8,-175.9)..(159.9,-176)
				..controls(159.2,-176.1)and(154.1,-176.2)..(151.5,-176.2)
				..controls(148.1,-176.1)and(122.6,-176.1)..(121.6,-176.2)
				..controls(120.8,-176.2)and(119.8,-176.2)..(119.4,-176.1)--cycle
				(130.4,-160)..controls(129,-160)and(119.8,-160)..(110.1,-159.9)
				..controls(100.3,-159.8)and(88.5,-159.7)..(83.9,-159.6)
				..controls(77.8,-159.5)and(75.3,-159.5)..(74.5,-159.6)
				..controls(73.9,-159.7)and(72.5,-159.8)..(71.3,-159.8)
				..controls(69.6,-159.8)and(69,-159.8)..(68.5,-159.6)
				..controls(67.6,-159.3)and(66.8,-158.6)..(66.2,-157.6)
				..controls(66,-157.3)and(65.6,-156.5)..(65.3,-156.1)
				..controls(64.6,-155.1)and(64.5,-154.9)..(63.7,-152.9)
				..controls(63.2,-151.6)and(62.9,-150.9)..(61.8,-149.6)
				..controls(61.5,-149.2)and(61.1,-148.5)..(60.9,-147.9)
				..controls(60.7,-147.4)and(60.3,-146.6)..(60.1,-146.1)
				..controls(59.9,-145.5)and(59.5,-144.9)..(59,-144.3)
				..controls(58.5,-143.7)and(58.2,-143.2)..(57.9,-142.5)
				..controls(57.7,-141.9)and(57.3,-141.1)..(57,-140.7)
				..controls(56.6,-140.3)and(56.1,-139.3)..(55.8,-138.5)
				..controls(55.5,-137.8)and(55.1,-136.8)..(54.9,-136.5)
				..controls(54.7,-136.1)and(54.3,-135.2)..(53.9,-134.5)
				..controls(53.6,-133.8)and(53,-132.7)..(52.6,-132)
				..controls(51.6,-130.5)and(51.4,-130.2)..(51.1,-129)
				..controls(50.6,-127.5)and(50.6,-127.4)..(49.7,-126.3)
				..controls(49.2,-125.7)and(48.6,-124.8)..(48.3,-124.2)
				..controls(48,-123.6)and(47.5,-122.7)..(47.3,-122.2)
				..controls(47,-121.8)and(46.5,-120.7)..(46.1,-119.8)
				..controls(45.8,-118.9)and(45.2,-117.8)..(44.9,-117.2)
				..controls(44.2,-116.1)and(43.3,-114.4)..(42.3,-112.2)
				..controls(41.8,-111.3)and(41.3,-110.4)..(41.1,-110.1)
				..controls(40.5,-109.4)and(40.3,-109)..(40,-107.8)
				..controls(39.8,-107)and(39.6,-106.5)..(39.2,-106)
				..controls(38.9,-105.6)and(38.4,-104.7)..(38,-104)
				..controls(37.7,-103.4)and(37.2,-102.5)..(37,-102.1)
				..controls(36.7,-101.7)and(36.3,-100.8)..(36,-100.1)
				..controls(35.6,-98.9)and(35.3,-98.4)..(34.6,-97.6)
				..controls(34.2,-97.1)and(34.2,-96.3)..(34.6,-96)
				..controls(35,-95.7)and(35.5,-95.9)..(36.2,-96.7)
				..controls(36.5,-97)and(37.2,-97.5)..(37.6,-97.7)..controls(38,-98)and(38.6,-98.4)..(39.1,-98.9)
				..controls(39.6,-99.4)and(40.3,-100)..(40.6,-100.3)
				..controls(41.3,-100.8)and(42.4,-101.8)..(43.7,-102.9)
				..controls(44,-103.2)and(44.6,-103.7)..(45,-103.9)
				..controls(45.4,-104.2)and(45.9,-104.6)..(46.1,-104.8)
				..controls(47.1,-105.9)and(48.8,-107.1)..(50.5,-107.9)
				..controls(51.8,-108.6)and(52.3,-109)..(53.4,-110)
				..controls(54.8,-111.2)and(55.6,-111.6)..(57,-112.1)
				..controls(57.5,-112.3)and(58.5,-112.8)..(59.4,-113.3)
				..controls(60.2,-113.7)and(61.3,-114.3)..(61.9,-114.5)
				..controls(62.4,-114.7)and(63.3,-115.1)..(63.8,-115.3)
				..controls(64.4,-115.6)and(65.4,-115.9)..(66.2,-116.1)
				..controls(67,-116.3)and(68.4,-116.6)..(69.2,-116.9)
				..controls(72.3,-117.9)and(72.8,-118)..(78,-118.1)
				..controls(84,-118.3)and(86.1,-118.1)..(88.1,-117.1)
				..controls(89.4,-116.5)and(93,-115)..(93.7,-114.9)
				..controls(94.8,-114.6)and(95.4,-114.2)..(95.9,-113.4)
				..controls(96.2,-113)and(96.7,-112.5)..(97.1,-112.2)
				..controls(98,-111.6)and(101.3,-108.4)..(102.1,-107.2)
				..controls(102.4,-106.8)and(102.9,-106.2)..(103.2,-105.8)
				..controls(103.5,-105.5)and(104,-104.8)..(104.2,-104.2)
				..controls(104.5,-103.6)and(105,-102.8)..(105.3,-102.4)
				..controls(105.6,-102)and(106.1,-101.2)..(106.5,-100.6)
				..controls(107.1,-99.5)and(108.6,-97.5)..(109.5,-96.5)
				..controls(109.9,-96.1)and(110.1,-96)..(110.4,-96)
				..controls(111.1,-96)and(111.6,-96.5)..(111.8,-97.6)
				..controls(112,-98.6)and(112.3,-99.3)..(112.8,-99.9)
				..controls(113.7,-100.9)and(115.2,-103.1)..(115.7,-104.2)
				..controls(116.1,-105)and(116.5,-105.5)..(116.9,-106)
				..controls(117.3,-106.4)and(117.8,-106.9)..(118,-107.3)
				..controls(118.2,-107.6)and(118.6,-108)..(119,-108.3)
				..controls(119.6,-108.8)and(121.1,-110.3)..(121.8,-111.2)
				..controls(122.1,-111.7)and(122.5,-112)..(123.4,-112.4)
				..controls(124.8,-113.2)and(125.8,-113.8)..(127,-114.6)
				..controls(127.5,-114.9)and(128.1,-115.3)..(128.2,-115.4)
				..controls(128.4,-115.5)and(129.3,-115.8)..(130.2,-116)
				..controls(131.1,-116.2)and(132.5,-116.6)..(133.3,-116.9)
				..controls(135.7,-117.9)and(136.1,-117.9)..(142.5,-117.8)
				..controls(145.8,-117.8)and(148.4,-117.7)..(148.8,-117.6)
				..controls(149.2,-117.6)and(150.2,-117.2)..(151,-116.8)
				..controls(152,-116.4)and(152.8,-116)..(153.6,-115.9)
				..controls(154.5,-115.7)and(155.2,-115.4)..(156.3,-114.9)
				..controls(157.1,-114.5)and(158.2,-114)..(158.7,-113.8)
				..controls(159.1,-113.7)and(159.8,-113.3)..(160.2,-113)
				..controls(160.6,-112.7)and(161.3,-112.4)..(161.7,-112.2)
				..controls(162.7,-111.9)and(163.2,-111.7)..(164.1,-111.1)
				..controls(164.5,-110.8)and(165.3,-110.4)..(165.9,-110)
				..controls(166.5,-109.7)and(167.3,-109.2)..(167.6,-108.9)
				..controls(167.9,-108.6)and(168.6,-108.2)..(169,-108)
				..controls(169.4,-107.8)and(170.2,-107.3)..(170.7,-106.9)
				..controls(171.2,-106.5)and(172.1,-105.8)..(172.7,-105.4)
				..controls(173.2,-105)and(174.1,-104.4)..(174.5,-104)
				..controls(174.9,-103.7)and(175.7,-103.1)..(176.2,-102.6)
				..controls(176.7,-102.2)and(177.4,-101.5)..(177.9,-101.2)
				..controls(178.3,-100.8)and(178.8,-100.4)..(179,-100.2)
				..controls(179.2,-100)and(179.6,-99.6)..(180,-99.3)
				..controls(180.4,-99.1)and(180.9,-98.6)..(181.2,-98.4)
				..controls(181.5,-98.1)and(182.1,-97.7)..(182.5,-97.5)
				..controls(183.2,-97.2)and(183.5,-96.9)..(184.3,-95.9)
				..controls(184.6,-95.6)and(184.6,-95.6)..(184.8,-95.7)
				..controls(185,-95.8)and(185.1,-96)..(185.1,-96.2)
				..controls(185.1,-96.5)and(185,-96.7)..(184.6,-97.2)
				..controls(184.1,-97.8)and(184,-97.9)..(183.9,-98.8)
				..controls(183.7,-100.3)and(183.3,-101.3)..(182.9,-101.8)
				..controls(182.6,-102.1)and(182.3,-102.6)..(182.2,-102.9)
				..controls(182,-103.3)and(181.6,-104)..(181.2,-104.6)
				..controls(180.9,-105.2)and(180.5,-106.1)..(180.4,-106.6)
				..controls(179.9,-108)and(179.7,-108.5)..(179,-109.4)
				..controls(178.6,-109.8)and(178.2,-110.5)..(177.9,-111)
				..controls(177.7,-111.5)and(177.3,-112.3)..(177,-112.7)
				..controls(176.6,-113.2)and(176.4,-113.7)..(176.3,-114.2)
				..controls(176.2,-114.7)and(175.9,-115.2)..(175.4,-116)
				..controls(175,-116.6)and(174.4,-117.7)..(174,-118.5)
				..controls(173.7,-119.2)and(173.3,-120.1)..(173,-120.5)
				..controls(172.8,-120.9)and(172.4,-121.8)..(172.1,-122.5)
				..controls(171.8,-123.2)and(171.3,-124.1)..(171,-124.6)
				..controls(170.8,-125.1)and(170.3,-126)..(169.9,-126.6)
				..controls(169.6,-127.3)and(169.2,-128.1)..(169,-128.3)
				..controls(168.8,-128.6)and(168.5,-129.3)..(168.3,-129.8)
				..controls(168.1,-130.3)and(167.8,-131.1)..(167.5,-131.6)
				..controls(166.9,-132.5)and(164.2,-137.9)..(164.2,-138.1)
				..controls(164.2,-138.1)and(163.9,-138.8)..(163.4,-139.7)
				..controls(162.9,-140.5)and(162.4,-141.5)..(162.3,-141.9)
				..controls(161.9,-143)and(161.6,-143.4)..(160.9,-144.1)
				..controls(160.4,-144.6)and(160.3,-144.8)..(160.1,-145.4)
				..controls(160,-145.8)and(159.6,-146.8)..(159.2,-147.8)
				..controls(158.7,-148.8)and(158.1,-150)..(157.9,-150.6)
				..controls(157.6,-151.1)and(157.2,-151.9)..(156.9,-152.4)
				..controls(156.5,-152.8)and(156.1,-153.7)..(155.8,-154.4)
				..controls(155.6,-155)and(155.2,-155.9)..(154.9,-156.3)
				..controls(154.7,-156.7)and(154.4,-157.4)..(154.2,-157.9)
				..controls(154.1,-158.4)and(153.8,-159)..(153.5,-159.3)
				..controls(152.8,-160.1)and(153,-160.1)..(142.4,-160.1)
				..controls(137.2,-160.1)and(131.8,-160.1)..(130.4,-160)--cycle
				(141.5,-103.7)..controls(140.8,-103)and(139.8,-102.1)..(139.2,-101.6)
				..controls(138.5,-101.2)and(137.9,-100.6)..(137.8,-100.4)
				..controls(137,-99.5)and(134,-96.5)..(133.2,-96)
				..controls(132.7,-95.7)and(131.8,-94.9)..(131.1,-94.3)
				..controls(129.4,-92.7)and(127.9,-91.6)..(126.6,-91.1)
				..controls(125.2,-90.5)and(124.1,-90)..(122.9,-89.1)
				..controls(122.1,-88.6)and(121.3,-88.2)..(119.9,-87.7)
				..controls(118.8,-87.3)and(117.5,-86.8)..(116.9,-86.5)
				..controls(114.3,-85.5)and(113.7,-85.4)..(109.7,-85.5)
				..controls(105.3,-85.7)and(103.9,-85.9)..(101.6,-87)
				..controls(100.7,-87.4)and(99.4,-87.8)..(97.9,-88)
				..controls(97,-88.2)and(96.6,-88.3)..(95.8,-88.9)
				..controls(95.4,-89.2)and(94.7,-89.6)..(94.2,-89.8)
				..controls(92.9,-90.3)and(92.3,-90.7)..(90.9,-92.1)
				..controls(90.2,-92.7)and(89.4,-93.3)..(89,-93.5)..controls(87,-94.5)and(85.2,-96)..(83.6,-97.9)
				..controls(83.3,-98.2)and(82.5,-98.9)..(81.8,-99.4)
				..controls(80.7,-100.1)and(80.5,-100.3)..(80,-101.1)
				..controls(78.7,-103.1)and(77.2,-103.9)..(76.4,-103)
				..controls(75.7,-102.2)and(75.7,-103.2)..(75.7,-83.6)
				..controls(75.6,-73.6)and(75.6,-64.8)..(75.5,-64)
				..controls(75.5,-63.2)and(75.5,-61.4)..(75.6,-59.8)--(75.7,-57.1)--(76,-56.8)
				..controls(76.8,-56.2)and(78,-56.3)..(79.8,-57.2)
				..controls(80.4,-57.5)and(81.3,-57.8)..(81.8,-57.9)
				..controls(82.3,-58.1)and(83,-58.4)..(83.6,-58.8)..controls(84.3,-59.2)and(85,-59.5)..(86,-59.8)
				..controls(86.8,-60.1)and(88,-60.6)..(88.7,-60.9)
				..controls(89.3,-61.2)and(90.6,-61.7)..(91.4,-62)
				..controls(92.2,-62.3)and(93.7,-62.9)..(94.6,-63.3)
				..controls(95.6,-63.7)and(96.8,-64.1)..(97.4,-64.3)
				..controls(98,-64.5)and(99,-64.9)..(99.7,-65.2)
				..controls(100.4,-65.6)and(101.5,-66)..(102.3,-66.2)
				..controls(104.2,-66.8)and(105,-67.1)..(106.2,-67.8)
				..controls(107.9,-68.7)and(110.2,-69.4)..(111.3,-69.2)
				..controls(112.5,-69)and(115,-68.2)..(115.6,-67.8)
				..controls(117.3,-66.7)and(117.6,-66.6)..(119.5,-66.2)
				..controls(121.9,-65.8)and(122.1,-65.7)..(123,-65.1)
				..controls(124.1,-64.4)and(124.4,-64.3)..(126,-63.9)
				..controls(127,-63.7)and(127.7,-63.4)..(128.4,-63.1)
				..controls(129,-62.8)and(129.7,-62.5)..(130.6,-62.4)
				..controls(132.1,-62.1)and(134,-61.5)..(135.2,-61)
				..controls(135.7,-60.8)and(136.7,-60.5)..(137.3,-60.3)
				..controls(138.4,-60)and(139.3,-59.7)..(141.1,-59)
				..controls(141.6,-58.7)and(142,-58.6)..(142.5,-58.6)
				..controls(143.2,-58.6)and(143.3,-58.7)..(143.6,-59)
				..controls(144,-59.5)and(144.1,-59.9)..(144,-60.8)
				..controls(144,-61.2)and(144,-71.1)..(144,-82.8)--(143.9,-104.1)--(143.7,-104.5)
				..controls(143.5,-104.7)and(143.3,-104.8)..(143,-104.9)
				..controls(142.6,-104.9)and(142.5,-104.8)..(141.5,-103.7)--cycle
				(27.8,-91.4)..controls(27.1,-90.9)and(26.3,-90)..(26.1,-89.5)
				..controls(26,-89.2)and(25.8,-88.6)..(25.7,-87.9)
				..controls(25.4,-86.2)and(25.6,-84.8)..(26.5,-83.5)
				..controls(27.4,-82.3)and(28.3,-82)..(29.4,-82.3)
				..controls(31.1,-82.8)and(32.1,-84.9)..(32,-87.4)
				..controls(31.9,-89.3)and(31.2,-90.6)..(29.9,-91.4)
				..controls(29.1,-91.9)and(28.6,-91.9)..(27.8,-91.4)--cycle
				(190.1,-91.3)..controls(189.7,-91.2)and(189.3,-90.9)..(189.1,-90.7)
				..controls(188.3,-89.8)and(187.6,-87.7)..(187.6,-86.2)
				..controls(187.6,-85.5)and(187.7,-85.2)..(188.2,-84.4)
				..controls(189.1,-82.6)and(190.1,-81.8)..(191.3,-82.1)
				..controls(192,-82.3)and(193.5,-83.7)..(193.8,-84.4)
				..controls(194.1,-85.1)and(194.2,-87.1)..(194,-88.3)
				..controls(193.8,-89.8)and(193.3,-90.5)..(192.3,-91.3)
				..controls(191.7,-91.7)and(190.9,-91.7)..(190.1,-91.3)--cycle
				(173.7,-83.4)..controls(172.2,-82.5)and(172,-82.5)..(170.6,-82.4)
				..controls(169.7,-82.4)and(168.8,-82.4)..(168.3,-82.5)
				..controls(167.2,-82.7)and(166.9,-82.7)..(166.6,-82.2)
				..controls(166.3,-81.8)and(166.3,-81)..(166.8,-80)
				..controls(167.2,-78.9)and(167.4,-77.4)..(167.7,-74.1)
				..controls(167.8,-72.8)and(167.9,-71.8)..(167.9,-71.7)
				..controls(168,-71.7)and(168,-70.2)..(168,-68.4)..controls(168,-65.4)and(168,-64.9)..(167.7,-63)
				..controls(167.3,-60.5)and(167.1,-59.9)..(166.4,-58.4)
				..controls(166.1,-57.7)and(165.8,-56.9)..(165.7,-56.6)
				..controls(165.7,-56.2)and(165.4,-55.6)..(165.2,-55.2)
				..controls(164.6,-53.8)and(161.9,-50.7)..(161,-50.3)
				..controls(159.8,-49.7)and(157.9,-50)..(155.2,-51.2)
				..controls(154.7,-51.4)and(153.4,-51.8)..(152.2,-52.1)
				..controls(150.8,-52.4)and(149.4,-52.9)..(148.6,-53.2)
				..controls(147.8,-53.6)and(146.7,-54)..(146,-54.2)
				..controls(145.4,-54.4)and(144.5,-54.7)..(144,-54.9)
				..controls(143.6,-55.2)and(142.7,-55.5)..(141.8,-55.7)
				..controls(141,-55.9)and(139.6,-56.3)..(138.7,-56.6)
				..controls(137.8,-57)and(136.3,-57.4)..(135.3,-57.6)
				..controls(133.8,-57.9)and(133.2,-58.1)..(132.1,-58.7)
				..controls(131.1,-59.2)and(130.3,-59.4)..(129.4,-59.6)
				..controls(128.7,-59.8)and(127.8,-60)..(127.4,-60.1)
				..controls(127,-60.3)and(126,-60.5)..(125.1,-60.7)
				..controls(123.2,-61.1)and(121.7,-61.5)..(120.2,-62.1)
				..controls(119.7,-62.3)and(118.1,-62.8)..(116.7,-63.2)
				..controls(114.6,-63.9)and(114,-64.1)..(113,-64.7)
				..controls(111.9,-65.4)and(111.8,-65.4)..(111,-65.5)
				..controls(109.9,-65.5)and(109.3,-65.4)..(107.7,-64.7)
				..controls(107,-64.4)and(106,-64)..(105.3,-63.9)
				..controls(103.5,-63.5)and(102.2,-63)..(101.2,-62.5)
				..controls(100.1,-61.8)and(99.3,-61.5)..(97.8,-61.1)
				..controls(97.2,-61)and(96.2,-60.7)..(95.7,-60.4)
				..controls(95.1,-60.1)and(94.2,-59.8)..(93.5,-59.6)
				..controls(92.9,-59.4)and(91.7,-59)..(91,-58.7)..controls(90.2,-58.4)and(88.9,-58)..(88,-57.7)
				..controls(87.2,-57.5)and(86.1,-57.1)..(85.5,-56.8)
				..controls(85,-56.6)and(84.1,-56.2)..(83.5,-56)..controls(82.9,-55.7)and(81.8,-55.2)..(81,-54.9)
				..controls(80.1,-54.5)and(79.1,-54.2)..(78.7,-54.1)
				..controls(77.9,-53.8)and(77.3,-53.6)..(76.2,-52.9)
				..controls(75.7,-52.6)and(75.2,-52.4)..(74.1,-52.2)
				..controls(72.5,-51.9)and(72.4,-51.8)..(71.3,-51.1)
				..controls(69.9,-50.1)and(69.4,-49.9)..(67.9,-49.8)
				..controls(67,-49.7)and(66.6,-49.6)..(66.3,-49.4)
				..controls(66.1,-49.3)and(65.7,-49)..(65.4,-48.8)
				..controls(65.1,-48.7)and(64.2,-48.4)..(63.6,-48.2)
				..controls(62.9,-48)and(61.7,-47.6)..(60.9,-47.2)
				..controls(60.1,-46.8)and(58.9,-46.4)..(58.2,-46.2)
				..controls(57.6,-46)and(56.8,-45.7)..(56.5,-45.5)
				..controls(54.9,-44.6)and(54.6,-44.5)..(53.2,-44.2)
				..controls(51.7,-43.8)and(51,-43.6)..(49.9,-42.9)
				..controls(49.5,-42.6)and(48.9,-42.4)..(47.5,-42)
				..controls(45.7,-41.6)and(45.6,-41.5)..(44.3,-40.6)
				..controls(43.5,-40.1)and(42.8,-39.6)..(42.7,-39.4)
				..controls(42.4,-38.7)and(43,-38.1)..(44.5,-37.6)..controls(44.9,-37.5)and(46,-37)..(46.9,-36.6)
				..controls(47.9,-36.1)and(48.8,-35.8)..(49.6,-35.6)
				..controls(50.5,-35.4)and(51.3,-35.2)..(52.4,-34.6)
				..controls(53.3,-34.1)and(54.5,-33.7)..(54.9,-33.5)
				..controls(55.4,-33.4)and(56.4,-33)..(57.2,-32.6)
				..controls(58.6,-32)and(60.1,-31.5)..(61.6,-31.2)
				..controls(62.3,-31.1)and(62.7,-31)..(63.2,-30.7)
				..controls(64.1,-30)and(64.3,-29.9)..(65.3,-29.7)
				..controls(65.7,-29.6)and(66.5,-29.2)..(67.1,-28.9)
				..controls(67.6,-28.6)and(68.8,-28.2)..(69.6,-27.9)
				..controls(70.5,-27.7)and(71.6,-27.3)..(72.1,-27)
				..controls(72.6,-26.7)and(73.8,-26.2)..(74.8,-25.9)
				..controls(75.7,-25.5)and(76.7,-25.1)..(77.1,-25)
				..controls(77.4,-24.8)and(78.5,-24.4)..(79.6,-24.1)
				..controls(80.6,-23.8)and(82,-23.2)..(82.8,-22.9)
				..controls(83.5,-22.5)and(84.8,-22)..(85.8,-21.8)
				..controls(86.7,-21.5)and(87.8,-21.1)..(88.1,-20.9)
				..controls(88.4,-20.7)and(89.3,-20.3)..(90.1,-20.1)
				..controls(90.8,-19.8)and(91.9,-19.4)..(92.4,-19.1)
				..controls(93.5,-18.5)and(94.4,-18.2)..(95.4,-18.1)
				..controls(95.7,-18)and(96.4,-17.8)..(96.9,-17.5)
				..controls(97.5,-17.3)and(98.4,-16.9)..(99,-16.6)
				..controls(99.7,-16.3)and(100.9,-15.9)..(101.9,-15.6)
				..controls(102.9,-15.3)and(104.4,-14.7)..(105.3,-14.3)
				..controls(108.1,-13.2)and(108.5,-13.1)..(110.1,-13.1)
				..controls(111.4,-13.1)and(111.5,-13.1)..(112.4,-13.5)
				..controls(113.2,-14)and(115,-14.8)..(116.6,-15.4)
				..controls(116.9,-15.5)and(117.9,-15.8)..(118.8,-16.1)
				..controls(120.9,-16.6)and(121.5,-16.9)..(122.6,-17.9)
				..controls(123,-18.3)and(123.7,-18.8)..(123.9,-18.9)
				..controls(124.6,-19.4)and(126,-19.8)..(127.9,-20.2)
				..controls(129.6,-20.5)and(129.8,-20.6)..(131.1,-21.4)
				..controls(131.5,-21.7)and(132.4,-22.1)..(133.1,-22.3)
				..controls(136.1,-23.2)and(138.3,-24.1)..(140.2,-25)
				..controls(141,-25.5)and(141.5,-25.6)..(142.5,-25.8)
				..controls(145.1,-26.3)and(146,-26.6)..(149.4,-28.1)
				..controls(151.3,-28.9)and(153.5,-29.8)..(154.3,-30.1)
				..controls(155.1,-30.4)and(156.2,-30.8)..(156.6,-31)
				..controls(157.1,-31.2)and(158.2,-31.6)..(159.1,-31.8)
				..controls(160.6,-32.2)and(160.9,-32.3)..(161.7,-32.9)
				..controls(162.4,-33.3)and(162.9,-33.6)..(163.7,-33.8)
				..controls(164.3,-34)and(164.8,-34.1)..(164.9,-34.1)
				..controls(165,-34.1)and(165.6,-34.4)..(166.4,-34.8)
				..controls(167.3,-35.2)and(168.1,-35.4)..(169,-35.6)
				..controls(170,-35.8)and(170.5,-36)..(171.7,-36.6)
				..controls(172.4,-37)and(173.3,-37.4)..(173.6,-37.4)
				..controls(175.6,-37.7)and(176.2,-38)..(176.7,-39.1)
				..controls(176.9,-39.6)and(176.9,-39.7)..(176.7,-40.4)
				..controls(176.4,-41.6)and(175.8,-42.7)..(175.2,-43.4)
				..controls(174.5,-44.1)and(172.6,-45.5)..(171.7,-45.8)
				..controls(170.6,-46.2)and(169.1,-46)..(168.2,-45.4)
				..controls(167.9,-45.2)and(167.3,-44.8)..(166.9,-44.4)
				..controls(166.4,-44.1)and(165.8,-43.6)..(165.4,-43.2)
				..controls(165.1,-42.9)and(164.4,-42.3)..(163.9,-41.9)
				..controls(163.3,-41.5)and(162.7,-40.9)..(162.4,-40.5)
				..controls(161.4,-39.2)and(159.9,-38)..(158.8,-37.8)
				..controls(158.4,-37.8)and(158.3,-37.7)..(158.2,-37.2)
				..controls(158,-36.8)and(157.4,-36.1)..(157.1,-36.1)
				..controls(157,-36.1)and(156.7,-36.1)..(156.5,-36)
				..controls(156.1,-35.9)and(155.9,-35.7)..(155.7,-35.2)
				..controls(155.3,-34.6)and(155.1,-34.4)..(154,-34.1)
				..controls(153.6,-34)and(153.2,-33.7)..(152.9,-33.5)
				..controls(152.4,-32.9)and(151.6,-32.4)..(151.3,-32.4)
				..controls(150.9,-32.4)and(150.4,-32.6)..(150.4,-32.9)
				..controls(150.4,-33.2)and(150.8,-33.4)..(151.4,-33.6)
				..controls(152.1,-33.7)and(152.8,-34.1)..(153.5,-34.9)
				..controls(154.1,-35.5)and(154.3,-35.6)..(155,-35.8)
				..controls(155.7,-36)and(155.7,-36)..(156,-36.6)
				..controls(156.3,-37.3)and(156.7,-37.7)..(157.5,-37.8)
				..controls(158.1,-37.9)and(158.2,-38)..(158.4,-38.8)
				..controls(158.5,-39.4)and(158.6,-39.6)..(159.6,-40.4)
				..controls(160.2,-41)and(160.9,-41.6)..(161.3,-41.8)
				..controls(161.7,-42)and(161.9,-42.3)..(162.2,-42.8)
				..controls(162.4,-43.2)and(162.8,-43.8)..(163.1,-44.1)
				..controls(163.9,-44.8)and(165.1,-46.7)..(165.7,-48)
				..controls(166.2,-49.3)and(166.4,-49.7)..(167.2,-50.5)
				..controls(167.5,-50.8)and(167.9,-51.4)..(168.1,-51.8)
				..controls(168.3,-52.2)and(168.8,-53)..(169.1,-53.6)
				..controls(169.6,-54.3)and(169.8,-54.9)..(169.9,-55.6)
				..controls(170.4,-57.8)and(170.5,-58.1)..(171.2,-59.5)
				..controls(171.6,-60.3)and(172,-61.5)..(172.2,-62.1)
				..controls(172.3,-62.8)and(172.7,-63.7)..(172.9,-64.1)
				..controls(173.2,-64.6)and(173.4,-65.3)..(173.4,-65.6)
				..controls(173.8,-68.8)and(174.2,-70.6)..(174.9,-72.6)
				..controls(175.5,-74.5)and(175.9,-77.1)..(176.1,-80.2)
				..controls(176.3,-83.4)and(175.5,-84.4)..(173.7,-83.4)--cycle;
				\definecolor{ezs}{RGB}{247,179,26}\fill[ezs](-5.7,-156.6)--(-3.4,-157.5)--(61.9,12.7)--(59.7,13.6)--cycle;
				\definecolor{ezs}{RGB}{247,179,22}\fill[ezs](-3.8,-157.3)--(-1.5,-158.2)--(63.9,12)--(61.6,12.9)--cycle;
				\definecolor{ezs}{RGB}{248,179,19}\fill[ezs](-1.8,-158.1)--(0.4,-158.9)--(65.8,11.2)--(63.5,12.1)--cycle;
				\definecolor{ezs}{RGB}{248,178,16}\fill[ezs](0.1,-158.8)--(2.4,-159.7)--(67.7,10.5)--(65.4,11.4)--cycle;
				\definecolor{ezs}{RGB}{249,178,15}\fill[ezs](2,-159.5)--(4.3,-160.4)--(69.6,9.8)--(67.3,10.7)--cycle;
				\definecolor{ezs}{RGB}{251,177,17}\fill[ezs](3.9,-160.3)--(6.2,-161.1)--(71.5,9)--(69.2,9.9)--cycle;
				\definecolor{ezs}{RGB}{252,179,18}\fill[ezs](5.8,-161)--(8.1,-161.9)--(73.4,8.3)--(71.1,9.2)--cycle;
				\definecolor{ezs}{RGB}{252,182,17}\fill[ezs](7.7,-161.7)--(10,-162.6)--(75.3,7.6)--(73,8.5)--cycle;
				\definecolor{ezs}{RGB}{252,185,17}\fill[ezs](9.6,-162.5)--(11.9,-163.4)--(77.2,6.8)--(75,7.7)--cycle;
				\definecolor{ezs}{RGB}{251,189,18}\fill[ezs](11.5,-163.2)--(13.8,-164.1)--(79.2,6.1)--(76.9,7)--cycle;
				\definecolor{ezs}{RGB}{251,191,22}\fill[ezs](13.4,-163.9)--(15.7,-164.8)--(81.1,5.4)--(78.8,6.2)--cycle;
				\definecolor{ezs}{RGB}{252,193,25}\fill[ezs](15.4,-164.7)--(17.6,-165.6)--(83,4.6)--(80.7,5.5)--cycle;
				\definecolor{ezs}{RGB}{253,195,27}\fill[ezs](17.3,-165.4)--(19.6,-166.3)--(84.9,3.9)--(82.6,4.8)--cycle;
				\definecolor{ezs}{RGB}{253,196,31}\fill[ezs](19.2,-166.1)--(21.5,-167)--(86.8,3.2)--(84.5,4)--cycle;
				\definecolor{ezs}{RGB}{254,199,35}\fill[ezs](21.1,-166.9)--(23.4,-167.8)--(88.7,2.4)--(86.4,3.3)--cycle;
				\definecolor{ezs}{RGB}{253,201,38}\fill[ezs](23,-167.6)--(25.3,-168.5)--(90.6,1.7)--(88.3,2.6)--cycle;
				\definecolor{ezs}{RGB}{253,203,40}\fill[ezs](24.9,-168.3)--(27.2,-169.2)--(92.5,1)--(90.3,1.8)--cycle;
				\definecolor{ezs}{RGB}{253,205,42}\fill[ezs](26.8,-169.1)--(29.1,-170)--(94.4,0.2)--(92.2,1.1)--cycle;
				\definecolor{ezs}{RGB}{253,207,45}\fill[ezs](28.7,-169.8)--(31,-170.7)--(96.4,-0.5)--(94.1,0.4)--cycle;
				\definecolor{ezs}{RGB}{253,210,49}\fill[ezs](30.7,-170.5)--(32.9,-171.4)--(98.3,-1.2)--(96,-0.4)--cycle;
				\definecolor{ezs}{RGB}{253,212,54}\fill[ezs](32.6,-171.3)--(34.9,-172.2)--(100.2,-2)--(97.9,-1.1)--cycle;
				\definecolor{ezs}{RGB}{253,214,59}\fill[ezs](34.5,-172)--(36.8,-172.9)--(102.1,-2.7)--(99.8,-1.8)--cycle;
				\definecolor{ezs}{RGB}{253,216,63}\fill[ezs](36.4,-172.7)--(38.7,-173.6)--(104,-3.4)--(101.7,-2.6)--cycle;
				\definecolor{ezs}{RGB}{253,218,65}\fill[ezs](38.3,-173.5)--(40.6,-174.4)--(105.9,-4.2)--(103.6,-3.3)--cycle;
				\definecolor{ezs}{RGB}{254,220,66}\fill[ezs](40.2,-174.2)--(42.5,-175.1)--(107.8,-4.9)--(105.5,-4)--cycle;
				\definecolor{ezs}{RGB}{254,222,68}\fill[ezs](42.1,-175)--(44.4,-175.8)--(109.7,-5.6)--(107.5,-4.8)--cycle;
				\definecolor{ezs}{RGB}{254,224,71}\fill[ezs](44,-175.7)--(46.3,-176.6)--(111.7,-6.4)--(109.4,-5.5)--cycle;
				\definecolor{ezs}{RGB}{254,226,75}\fill[ezs](46,-176.4)--(48.2,-177.3)--(113.6,-7.1)--(111.3,-6.2)--cycle;
				\definecolor{ezs}{RGB}{254,228,79}\fill[ezs](47.9,-177.2)--(50.1,-178)--(115.5,-7.8)--(113.2,-7)--cycle;
				\definecolor{ezs}{RGB}{254,230,82}\fill[ezs](49.8,-177.9)--(52.1,-178.8)--(117.4,-8.6)--(115.1,-7.7)--cycle;
				\definecolor{ezs}{RGB}{254,232,85}\fill[ezs](51.7,-178.6)--(54,-179.5)--(119.3,-9.3)--(117,-8.4)--cycle;
				\definecolor{ezs}{RGB}{254,234,88}\fill[ezs](53.6,-179.4)--(55.9,-180.2)--(121.2,-10)--(118.9,-9.2)--cycle;
				\definecolor{ezs}{RGB}{254,237,91}\fill[ezs](55.5,-180.1)--(57.8,-181)--(123.1,-10.8)--(120.8,-9.9)--cycle;
				\definecolor{ezs}{RGB}{254,239,94}\fill[ezs](57.4,-180.8)--(59.7,-181.7)--(125,-11.5)--(122.8,-10.6)--cycle;
				\definecolor{ezs}{RGB}{254,241,97}\fill[ezs](59.3,-181.6)--(61.6,-182.4)--(127,-12.2)--(124.7,-11.4)--cycle;
				\definecolor{ezs}{RGB}{254,243,101}\fill[ezs](61.2,-182.3)--(63.5,-183.2)--(128.9,-13)--(126.6,-12.1)--cycle;
				\definecolor{ezs}{RGB}{254,245,103}\fill[ezs](63.2,-183)--(65.4,-183.9)--(130.8,-13.7)--(128.5,-12.8)--cycle;
				\definecolor{ezs}{RGB}{254,245,105}\fill[ezs](65.1,-183.8)--(67.4,-184.6)--(132.7,-14.4)--(130.4,-13.6)--cycle;
				\definecolor{ezs}{RGB}{254,245,105}\fill[ezs](67,-184.5)--(69.3,-185.4)--(134.6,-15.2)--(132.3,-14.3)--cycle;
				\definecolor{ezs}{RGB}{254,244,103}\fill[ezs](68.9,-185.2)--(71.2,-186.1)--(136.5,-15.9)--(134.2,-15)--cycle;
				\definecolor{ezs}{RGB}{253,243,101}\fill[ezs](70.8,-186)--(73.1,-186.8)--(138.4,-16.6)--(136.1,-15.8)--cycle;
				\definecolor{ezs}{RGB}{253,241,98}\fill[ezs](72.7,-186.7)--(75,-187.6)--(140.3,-17.4)--(138,-16.5)--cycle;
				\definecolor{ezs}{RGB}{253,239,95}\fill[ezs](74.6,-187.4)--(76.9,-188.3)--(142.2,-18.1)--(140,-17.2)--cycle;
				\definecolor{ezs}{RGB}{253,236,92}\fill[ezs](76.5,-188.2)--(78.8,-189)--(144.2,-18.8)--(141.9,-18)--cycle;
				\definecolor{ezs}{RGB}{253,234,88}\fill[ezs](78.5,-188.9)--(80.7,-189.8)--(146.1,-19.6)--(143.8,-18.7)--cycle;
				\definecolor{ezs}{RGB}{253,232,85}\fill[ezs](80.4,-189.6)--(82.7,-190.5)--(148,-20.3)--(145.7,-19.4)--cycle;
				\definecolor{ezs}{RGB}{253,230,82}\fill[ezs](82.3,-190.4)--(84.6,-191.2)--(149.9,-21)--(147.6,-20.2)--cycle;
				\definecolor{ezs}{RGB}{253,228,78}\fill[ezs](84.2,-191.1)--(86.5,-192)--(151.8,-21.8)--(149.5,-20.9)--cycle;
				\definecolor{ezs}{RGB}{253,226,74}\fill[ezs](86.1,-191.8)--(88.4,-192.7)--(153.7,-22.5)--(151.4,-21.6)--cycle;
				\definecolor{ezs}{RGB}{253,224,70}\fill[ezs](88,-192.6)--(90.3,-193.4)--(155.6,-23.3)--(153.3,-22.4)--cycle;
				\definecolor{ezs}{RGB}{253,222,67}\fill[ezs](89.9,-193.3)--(92.2,-194.2)--(157.5,-24)--(155.3,-23.1)--cycle;
				\definecolor{ezs}{RGB}{253,220,65}\fill[ezs](91.8,-194)--(94.1,-194.9)--(159.5,-24.7)--(157.2,-23.8)--cycle;
				\definecolor{ezs}{RGB}{253,217,63}\fill[ezs](93.7,-194.8)--(96,-195.6)--(161.4,-25.5)--(159.1,-24.6)--cycle;
				\definecolor{ezs}{RGB}{253,215,60}\fill[ezs](95.7,-195.5)--(97.9,-196.4)--(163.3,-26.2)--(161,-25.3)--cycle;
				\definecolor{ezs}{RGB}{253,213,56}\fill[ezs](97.6,-196.2)--(99.9,-197.1)--(165.2,-26.9)--(162.9,-26)--cycle;
				\definecolor{ezs}{RGB}{252,211,53}\fill[ezs](99.5,-197)--(101.8,-197.8)--(167.1,-27.7)--(164.8,-26.8)--cycle;
				\definecolor{ezs}{RGB}{252,209,49}\fill[ezs](101.4,-197.7)--(103.7,-198.6)--(169,-28.4)--(166.7,-27.5)--cycle;
				\definecolor{ezs}{RGB}{252,207,45}\fill[ezs](103.3,-198.4)--(105.6,-199.3)--(170.9,-29.1)--(168.6,-28.2)--cycle;
				\definecolor{ezs}{RGB}{253,205,42}\fill[ezs](105.2,-199.2)--(107.5,-200)--(172.8,-29.9)--(170.6,-29)--cycle;
				\definecolor{ezs}{RGB}{253,202,40}\fill[ezs](107.1,-199.9)--(109.4,-200.8)--(174.7,-30.6)--(172.5,-29.7)--cycle;
				\definecolor{ezs}{RGB}{253,200,38}\fill[ezs](109,-200.6)--(111.3,-201.5)--(176.7,-31.3)--(174.4,-30.4)--cycle;
				\definecolor{ezs}{RGB}{253,198,35}\fill[ezs](111,-201.4)--(113.2,-202.2)--(178.6,-32.1)--(176.3,-31.2)--cycle;
				\definecolor{ezs}{RGB}{253,196,31}\fill[ezs](112.9,-202.1)--(115.2,-203)--(180.5,-32.8)--(178.2,-31.9)--cycle;
				\definecolor{ezs}{RGB}{253,194,28}\fill[ezs](114.8,-202.8)--(117.1,-203.7)--(182.4,-33.5)--(180.1,-32.6)--cycle;
				\definecolor{ezs}{RGB}{253,192,26}\fill[ezs](116.7,-203.6)--(119,-204.5)--(184.3,-34.3)--(182,-33.4)--cycle;
				\definecolor{ezs}{RGB}{253,189,24}\fill[ezs](118.6,-204.3)--(120.9,-205.2)--(186.2,-35)--(183.9,-34.1)--cycle;
				\definecolor{ezs}{RGB}{253,186,20}\fill[ezs](120.5,-205)--(122.8,-205.9)--(188.1,-35.7)--(185.8,-34.8)--cycle;
				\definecolor{ezs}{RGB}{253,184,17}\fill[ezs](122.4,-205.8)--(124.7,-206.7)--(190,-36.5)--(187.8,-35.6)--cycle;
				\definecolor{ezs}{RGB}{252,183,15}\fill[ezs](124.3,-206.5)--(126.6,-207.4)--(192,-37.2)--(189.7,-36.3)--cycle;
				\definecolor{ezs}{RGB}{252,181,14}\fill[ezs](126.3,-207.2)--(128.5,-208.1)--(193.9,-37.9)--(191.6,-37.1)--cycle;
				\definecolor{ezs}{RGB}{252,179,13}\fill[ezs](128.2,-208)--(130.4,-208.9)--(195.8,-38.7)--(193.5,-37.8)--cycle;
				\definecolor{ezs}{RGB}{252,178,12}\fill[ezs](130.1,-208.7)--(132.4,-209.6)--(197.7,-39.4)--(195.4,-38.5)--cycle;
				\definecolor{ezs}{RGB}{251,178,10}\fill[ezs](132,-209.4)--(134.3,-210.3)--(199.6,-40.1)--(197.3,-39.3)--cycle;
				\definecolor{ezs}{RGB}{251,178,10}\fill[ezs](133.9,-210.2)--(136.2,-211.1)--(201.5,-40.9)--(199.2,-40)--cycle;
				\definecolor{ezs}{RGB}{250,177,11}\fill[ezs](135.8,-210.9)--(138.1,-211.8)--(203.4,-41.6)--(201.1,-40.7)--cycle;
				\definecolor{ezs}{RGB}{250,178,11}\fill[ezs](137.7,-211.6)--(140,-212.5)--(205.3,-42.3)--(203.1,-41.5)--cycle;
				\definecolor{ezs}{RGB}{248,180,10}\fill[ezs](139.6,-212.4)--(141.9,-213.3)--(207.3,-43.1)--(205,-42.2)--cycle;
				\definecolor{ezs}{RGB}{247,181,9}\fill[ezs](141.5,-213.1)--(143.8,-214)--(209.2,-43.8)--(206.9,-42.9)--cycle;
				\definecolor{ezs}{RGB}{246,182,9}\fill[ezs](143.5,-213.8)--(145.7,-214.7)--(211.1,-44.5)--(208.8,-43.7)--cycle;
				\definecolor{ezs}{RGB}{246,182,10}\fill[ezs](145.4,-214.6)--(147.7,-215.5)--(213,-45.3)--(210.7,-44.4)--cycle;
			\end{scope}
			% chữ EDUZY (Times đậm, làm dày nét bằng 5 lớp lệch nhẹ; căn đúng khung chữ logo gốc)
			\foreach \dx/\dy in {-0.7/0,0.7/0,0/-0.3,0/0.3,0/0}
			\node[anchor=center,inner sep=0pt,outer sep=0pt,text=ezText,align=center] at (109.6+\dx,-193.6+\dy)
			{\resizebox{1.01cm}{0.22cm}{\fontencoding{T1}\fontfamily{ptm}\fontseries{b}\selectfont EDUZY}};
		\end{tikzpicture}%
	}%
}
%====================== THIẾT LẬP THÂN ĐỀ ======================
%--- Môi trường tự luận {bt}: ex_test 3.3.4 chưa khai báo sẵn nên tạo ở đây
%    (ex_test đã tự móc \fixshowans và \listenumerate cho bt)
\makeatletter
\@ifundefined{bt}{\newtheorem{bt}{\color{gold1}\namebt}}{}
\makeatother
\renewcommand{\sectionmark}[1]{\markright{\thesection\ #1}}
%--- Khung tiêu đề (logo + tối đa 4 dòng; dòng 3, 4 để trống {} thì bỏ qua)
%  Đề thi:    \khungtieude{THI THỬ THPT QUỐC GIA NĂM 2026}{Môn Toán -- Đề số 21}
%                         {(Đề thi gồm: ...)}{Thời gian làm bài: 90 phút (không kể thời gian phát đề).}
%  Chuyên đề: \khungtieude{CHUYÊN ĐỀ 3 -- DÃY SỐ VÀ CÔNG THỨC TRUY HỒI}
%                         {Bồi dưỡng học sinh giỏi Toán THPT}{...}{}
\newcommand{\khungtieude}[4]{%
	\begin{tcolorbox}[colframe=gold1,colback=gold5!45,boxrule=0.5mm]
		\begin{minipage}[c]{3cm}
			\centerline{\eduzylogo[2.4cm]}
		\end{minipage}\hspace{0.5cm}
		\begin{minipage}[c]{12.5cm}
			\vskip 0.15 cm
			\centerline{\ind{#1}}
			\vskip 0.15 cm
			\centerline{\textbf{#2}}
			\if\relax\detokenize{#3}\relax\else
				\vskip 0.15 cm
				\centerline{\textit{#3}}
			\fi
			\if\relax\detokenize{#4}\relax\else
				\vskip 0.15 cm
				\centerline{\textit{#4}}
			\fi
			\vskip 0.15 cm
		\end{minipage}
	\end{tcolorbox}
}
%====================== BỔ SUNG CHO BÀI HỌC, CHUYÊN ĐỀ ======================
\usepackage[shortlabels]{enumitem}
\usepackage{tabularx}
%--- Danh sách gọn: \begin{enumerate}[gon] ... \end{enumerate}
\SetEnumitemKey{gon}{leftmargin=*,itemsep=1pt,parsep=0pt,topsep=3pt}
%--- Khung lý thuyết: \begin{khung}{} ... \end{khung}; tiêu đề chỉ dùng khi thật ngắn ("Mẹo", "Lỗi thường gặp")
\newtcolorbox{khung}[1]{breakable,enhanced,colframe=gold1,colback=gold5!45,
	boxrule=0.4mm,fonttitle=\bfseries\sffamily,coltitle=white,colbacktitle=gold1,
	attach boxed title to top left={xshift=6mm,yshift=-3mm},
	boxed title style={arc=1mm,boxrule=0.3mm},title={#1}}
%--- Tiêu đề mục nhỏ trong một dạng: \muc{Dấu hiệu}, \muc{Phương pháp}, \muc{Bài tự luyện}
\makeatletter
\newcommand{\muc}[1]{\par\medskip\noindent\indam{#1}\par\nobreak\smallskip\@afterindentfalse\@afterheading}
\makeatother
%--- Ô công thức cần nhớ (căn giữa, viền vàng, vừa khít nội dung)
%     \begin{center}\congthuc{$Q(n)=(n-1)!$}\end{center}
\newtcbox{\congthuc}{on line,colframe=gold1,colback=gold5!45,boxrule=0.5pt,arc=2pt,
	left=8pt,right=8pt,top=4pt,bottom=4pt}
%--- Công tắc lời giải bài tự luyện (ví dụ mẫu luôn dùng \loigiai)
%     \loigiaiTL[{đáp số}]{lời giải đầy đủ}
%     \TLfull -> in lời giải đầy đủ;  \TLdapso -> chỉ in một dòng "Đáp số";  \TLan -> ẩn hoàn toàn
%     (\CheDo là dethi thì mọi lời giải và đáp số đều ẩn)
\makeatletter
\newif\ifTLhien   \TLhientrue
\newif\ifTLdapso  \TLdapsofalse
\newcommand{\TLfull}{\TLhientrue\TLdapsofalse}
\newcommand{\TLdapso}{\TLhientrue\TLdapsotrue}
\newcommand{\TLan}{\TLhienfalse}
\newcommand{\loigiaiTL}[2][]{%
	\ifTLhien
		\ifTLdapso
			\def\@tlarg{#1}%
			\ifx\@tlarg\@empty\else
				\ifdefstring{\kindOPTN}{dethi}{}{\par\noindent{\bfseries\color{gold1}Đáp số:} #1\par}%
			\fi
		\else
			\loigiai{#2}%
		\fi
	\fi}
\makeatother
\TLdapso% mặc định: chỉ in đáp số bài tự luyện
%====================== BẢNG ĐÚNG/SAI (\choiceTFt) ======================
%  Hàng tiêu đề nền nâu đồng chữ trắng (Phát biểu | Đ | S), các hàng tô xen kẽ kem / trắng,
%  đường kẻ vàng kim, nhãn a) b) c) d) màu nâu đồng. Bản có đáp án: cột được chọn có vòng tròn
%  nâu đồng với dấu tích trắng (câu đúng tích ở cột Đúng, câu sai tích ở cột Sai).
%  Tiêu đề ghi đủ "Đúng | Sai": \OPTN{viettat=0} sau \input (cột Đúng/Sai rộng ra, phát biểu hẹp lại).
%  Dùng bảng gốc của ex_test: \def\BangDSgoc{1} trước \input{preamble-eduzy-v2}.
\providecommand{\BangDSgoc}{0}
\ifnum\BangDSgoc=0
\usepackage{colortbl,pifont}
\colorlet{ezDSke}{gold2!60!gold1}% đường kẻ bảng (nét 0,4pt như ex_test để không đổi chiều cao bảng)
\colorlet{ezDSle}{gold5!60}% hàng a), c)
\colorlet{ezDSchan}{white}% hàng b), d)
%--- dấu chọn ở cột Đúng/Sai
%    (bề rộng 0, căn giữa ô: không làm đổi bề rộng cột của ex_test, không báo tràn lề)
\newcommand{\EZdauchon}{\makebox[0pt]{\tikz[baseline=-1ex]\node[circle,fill=gold1,inner sep=0pt,
	minimum size=1.2em,text=white,font=\small]{\ding{51}};}}
\def\tickT{\EZdauchon}\def\tickF{\EZdauchon}
\AtBeginDocument{\sbox{\ccTFTab}{\EZdauchon}}
%--- nhãn a) b) c) d) màu nâu đồng (ở bản lời giải, nhãn câu đúng vẫn đóng khung như ex_test)
\renewcommand{\FalseTF}{\stepcounter{dapan}\textbf{\color{gold1}\DapAnTF\sepTF}\ignorespaces}
%--- trước khi dựng bảng: kẻ màu vàng kim (bề rộng cột giữ nguyên như ex_test); sau bảng trả lại màu đen
\newcount\EZhangTF
\newcommand{\EZbatdauTF}{%
	\global\EZhangTF=0\relax
	\arrayrulecolor{ezDSke}}
\newcommand{\EZketthucTF}{\arrayrulecolor{black}}
%--- hàng tiêu đề nền nâu đồng, chữ trắng
\makeatletter
\newcommand{\EZo@TF}[1]{\centering\bfseries\color{white}#1}
\newcommand{\EZtieudeTF}{%
	\ifcase\cotDS
		\ifnum\cotSTT=1\relax
			\def\rowHeader{\hline\rowcolor{gold1}\EZo@TF{TT}&\EZo@TF{\phatbieu}\tabularnewline}%
		\else
			\def\rowHeader{\hline\rowcolor{gold1}\EZo@TF{\phatbieu}\tabularnewline}%
		\fi
	\or
		\ifnum\cotSTT=1\relax
			\def\rowHeader{\hline\rowcolor{gold1}\EZo@TF{TT}&\EZo@TF{\phatbieu}&\EZo@TF{\Dung}\tabularnewline}%
		\else
			\def\rowHeader{\hline\rowcolor{gold1}\EZo@TF{\phatbieu}&\EZo@TF{\Dung}\tabularnewline}%
		\fi
	\or
		\ifnum\cotSTT=1\relax
			\def\rowHeader{\hline\rowcolor{gold1}\EZo@TF{TT}&\EZo@TF{\phatbieu}&\EZo@TF{\Dung}&\EZo@TF{\Sai}\tabularnewline}%
		\else
			\def\rowHeader{\hline\rowcolor{gold1}\EZo@TF{\phatbieu}&\EZo@TF{\Dung}&\EZo@TF{\Sai}\tabularnewline}%
		\fi
	\fi}
%--- tô xen kẽ từng hàng (gọi ngay sau \hline đầu mỗi hàng)
\newcommand{\EZmauhangTF}{%
	\noalign{\global\advance\EZhangTF\@ne}%
	\ifodd\EZhangTF\rowcolor{ezDSle}\else\rowcolor{ezDSchan}\fi}
%--- gắn vào \choiceTFt / \writeChoiceTFt của ex_test
\xpretocmd{\choiceTFt}{\EZbatdauTF}{}{\typeout{[preamble-eduzy-v2] Không gắn được bảng Đúng/Sai (đầu bảng)}}
\xpatchcmd{\choiceTFt}{\ifnum\dongPB=0\def\rowHeader{}\fi}{\EZtieudeTF\ifnum\dongPB=0\def\rowHeader{}\fi}
	{}{\typeout{[preamble-eduzy-v2] Không gắn được bảng Đúng/Sai (hàng tiêu đề)}}
\xpatchcmd{\writeChoiceTFt}{\hline\xdef\dotEX{\sdotEX}}{\hline\EZmauhangTF\xdef\dotEX{\sdotEX}}
	{}{\typeout{[preamble-eduzy-v2] Không gắn được bảng Đúng/Sai (tô hàng)}}
\xpatchcmd{\writeChoiceTFt}{{\end{longtable}}}{{\end{longtable}\EZketthucTF}}
	{}{\typeout{[preamble-eduzy-v2] Không gắn được bảng Đúng/Sai (cuối longtable)}}
\xpatchcmd{\writeChoiceTFt}{{\end{xtabular}\end{center}\medskip\noindent}}{{\end{xtabular}\end{center}\EZketthucTF\medskip\noindent}}
	{}{\typeout{[preamble-eduzy-v2] Không gắn được bảng Đúng/Sai (cuối xtabular)}}
\fi
%=====================================================================
\renewcommand{\parallel}{\mathrel{/\!\!/}}

%      \TLdapso                      % tài liệu: \TLfull | \TLdapso | \TLan
%      \begin{document} ... \end{document}
%  Logo Eduzy vẽ bằng TikZ ngay trong file (\eduzylogo), không cần file ảnh.
%  Biên dịch: pdfLaTeX (máy có gói vietnam) hoặc XeLaTeX.
%============================================================
\documentclass[12pt,a4paper,twoside]{article}%<-- twoside: bat dau trang/chan trang chan-le
\usepackage{iftex}
\usepackage{amsmath,amssymb,tikz,mathrsfs,tkz-tab}
\usepackage{graphicx}
\usepackage{pgfplots}% đồ thị hàm số
\pgfplotsset{compat=1.17}
%\usepackage{mathptmx}
\usepackage{mathpazo}
\ifXeTeX% XeLaTeX: font chữ có đủ dấu tiếng Việt, giữ nguyên font toán mathpazo
	\usepackage[no-math]{fontspec}
	\setmainfont{TeX Gyre Pagella}
	\newfontfamily\EZqag{TeX Gyre Adventor}[NFSSFamily=qag,Ligatures=TeX]
\fi
\usepackage[many]{tcolorbox}
\usetikzlibrary{math,through,calc,intersections,angles,quotes,shapes,shapes.geometric,arrows,patterns,matrix,chains,arrows.meta,decorations.shapes,decorations.fractals,decorations.markings,shadows,decorations.pathreplacing}
\usepackage[top=1.6cm, bottom=1.8cm, left=1.8cm, right=1.8cm]{geometry}% v2: top 1.6cm (bản cũ 1.5cm) để chữ không sát dải đầu trang
\providecommand{\CheDo}{loigiai}% loigiai | dethi | solcolor | color | book
\expandafter\usepackage\expandafter[\CheDo]{ex_test}
% \tenchuong: khai báo ở file chính (trước \input) hoặc \renewcommand sau \input
%============================================================
%  [GỘP TỪ eduzy-de.tex] EDUZY - MÀU NÂU ĐỒNG / VÀNG KIM
%  + ĐẦU TRANG, CHÂN TRANG KIỂU DẢI NGANG (v2), ĐỐI XỨNG TRANG CHẴN / TRANG LẺ.
%  (không cần file eduzy-de.tex riêng nữa)
%  Khối này phải đặt SAU \usepackage[...]{ex_test} và SAU \def\tenchuong{...}
%============================================================
\usepackage{lastpage}
\usepackage{needspace}
%--- DÀN TRANG ĐỀ THI: tách câu (THÊM 05/11/2026; chỉ thêm lệnh mới, KHÔNG đổi hành vi mặc định) ---
%  Hai công tắc, mặc định tắt (= preamble chạy y như cũ). Bật bằng cách đặt TRƯỚC %============================================================
%  PREAMBLE EDUZY v2 - dùng chung cho đề thi, bài học, chuyên đề (gói ex_test)
%  Khác bản preamble-eduzy.tex ở đầu trang, chân trang: dải nền nâu đồng tràn ngang,
%  logo + Eduzy.com.vn, ô chữ vàng kem cạnh xiên, chân trang "Trang (số)";
%  logo Eduzy chìm giữa mọi trang (tắt: \logochimtat, bật lại: \logochimbat).
%  Cách dùng ở file chính (xem mau-chuan.tex, mau-chuyen-de.tex):
%      \def\CheDo{loigiai}          % loigiai | dethi | solcolor | color | book
%      \def\tenchuong{<dòng chữ trên đầu trang>}
%      \def\thuonghieu{<chữ ở chân trang lẻ>}   % không khai báo thì dùng mặc định
%      \def\tacgia{<chữ ở chân trang chẵn>}      % không khai báo thì dùng mặc định
%      \def\khauhieu{<chữ cạnh logo ở đầu trang>}  % tùy chọn, mặc định Eduzy.com.vn
%      \input{preamble-eduzy-v2}
%      \TLdapso                      % tài liệu: \TLfull | \TLdapso | \TLan
%      \begin{document} ... \end{document}
%  Logo Eduzy vẽ bằng TikZ ngay trong file (\eduzylogo), không cần file ảnh.
%  Biên dịch: pdfLaTeX (máy có gói vietnam) hoặc XeLaTeX.
%============================================================
\documentclass[12pt,a4paper,twoside]{article}%<-- twoside: bat dau trang/chan trang chan-le
\usepackage{iftex}
\usepackage{amsmath,amssymb,tikz,mathrsfs,tkz-tab}
\usepackage{graphicx}
\usepackage{pgfplots}% đồ thị hàm số
\pgfplotsset{compat=1.17}
%\usepackage{mathptmx}
\usepackage{mathpazo}
\ifXeTeX% XeLaTeX: font chữ có đủ dấu tiếng Việt, giữ nguyên font toán mathpazo
	\usepackage[no-math]{fontspec}
	\setmainfont{TeX Gyre Pagella}
	\newfontfamily\EZqag{TeX Gyre Adventor}[NFSSFamily=qag,Ligatures=TeX]
\fi
\usepackage[many]{tcolorbox}
\usetikzlibrary{math,through,calc,intersections,angles,quotes,shapes,shapes.geometric,arrows,patterns,matrix,chains,arrows.meta,decorations.shapes,decorations.fractals,decorations.markings,shadows,decorations.pathreplacing}
\usepackage[top=1.6cm, bottom=1.8cm, left=1.8cm, right=1.8cm]{geometry}% v2: top 1.6cm (bản cũ 1.5cm) để chữ không sát dải đầu trang
\providecommand{\CheDo}{loigiai}% loigiai | dethi | solcolor | color | book
\expandafter\usepackage\expandafter[\CheDo]{ex_test}
% \tenchuong: khai báo ở file chính (trước \input) hoặc \renewcommand sau \input
%============================================================
%  [GỘP TỪ eduzy-de.tex] EDUZY - MÀU NÂU ĐỒNG / VÀNG KIM
%  + ĐẦU TRANG, CHÂN TRANG KIỂU DẢI NGANG (v2), ĐỐI XỨNG TRANG CHẴN / TRANG LẺ.
%  (không cần file eduzy-de.tex riêng nữa)
%  Khối này phải đặt SAU \usepackage[...]{ex_test} và SAU \def\tenchuong{...}
%============================================================
\usepackage{lastpage}
\usepackage{needspace}
%--- DÀN TRANG ĐỀ THI: tách câu (THÊM 05/11/2026; chỉ thêm lệnh mới, KHÔNG đổi hành vi mặc định) ---
%  Hai công tắc, mặc định tắt (= preamble chạy y như cũ). Bật bằng cách đặt TRƯỚC \input{preamble-eduzy-v2}
%  trong file đề (tools/dan_trang/dan_trang.py tự ghi khi đặt \NeedspaceCau):
%    \def\TachCauDS{1}   câu Đúng/Sai được tách: câu dẫn (kèm hình) ở cuối trang, bảng Phát biểu sang trang sau
%                        (câu dẫn và hình không tách nhau, bảng không tách giữa hai trang)
%    \def\TachCauLG{1}   bản lời giải (\CheDo khác dethi): phần đề của câu cao hơn 1/3 trang được phép tách
%  Lệnh dùng trong file đề (giá trị do tools/dan_trang/dan_trang.py đo ở bản dethi):
%    \NeedspaceCau[<cao câu dẫn + hình>]{<cao cả khối câu>}   ngay trên \begin{ex} (câu đầu phần: trên tiêu đề phần)
%        công tắc tắt: đúng bằng \Needspace{<cao cả khối câu>} như cũ;
%        \TachCauDS=1 và có [..] (câu Đúng/Sai được chọn tách): chỉ giữ câu dẫn + hình;
%        \TachCauLG=1 ở bản lời giải: đòi tối đa 1/3 \textheight.
%    \NeedspaceBang{<cao bảng Phát biểu>}%   ngay trước \choiceTFt: \TachCauDS=1 thì giữ trọn bảng trên một trang;
%        công tắc tắt: không làm gì.
\providecommand{\TachCauDS}{0}
\providecommand{\TachCauLG}{0}
\newlength{\EZnsCau}
\newcommand{\NeedspaceCau}[2][]{%
	\setlength{\EZnsCau}{#2}%
	\ifnum\TachCauDS=1 \if\relax\detokenize{#1}\relax\else\setlength{\EZnsCau}{#1}\fi\fi
	\ifnum\TachCauLG=1 \ifdefstring{\kindOPTN}{dethi}{}{\ifdim\EZnsCau>0.3333\textheight\setlength{\EZnsCau}{0.3333\textheight}\fi}\fi
	\Needspace{\EZnsCau}}
\newcommand{\NeedspaceBang}[1]{\ifnum\TachCauDS=1 \Needspace{#1}\fi}

%--- Thang màu: đậm -> nhạt
\definecolor{gold1}{RGB}{126,96,0}    % nâu đồng đậm  #7E6000
\definecolor{gold2}{RGB}{192,144,3}   % vàng kim      #C09003
\definecolor{gold3}{RGB}{255,215,99}  % vàng nhạt     #FFD763
\definecolor{gold4}{RGB}{255,231,153} % vàng rất nhạt #FFE799
\definecolor{gold5}{RGB}{255,243,205} % vàng kem      #FFF3CD
%--- Ghi đè tông xanh/đỏ cũ
\definecolor{myblue}{RGB}{126,96,0}
\colorlet{blue}{gold1}
\colorlet{cyan}{gold2}
\colorlet{teal}{gold2}

%--- Tên chương mặc định (file đề nên tự khai báo lại)
\providecommand{\tenchuong}{}
%--- Chân trang
\providecommand{\thuonghieu}{TOÁN 12 -- Eduzy.com.vn}
\providecommand{\tacgia}{Tô Gia Bảo -- 0353703783}

%--- Chữ cạnh logo ở đầu trang
\providecommand{\khauhieu}{Eduzy.com.vn}

%====================== THÔNG SỐ ĐẦU TRANG, CHÂN TRANG (v2) ======================
%  Dải nền đậm tràn ngang trang; ô chữ vàng kem có cạnh xiên, kèm 2 vạch xiên chuyển màu.
%  Đầu trang lẻ : logo + \khauhieu (trái)   | ô chữ \tenchuong (phải)
%  Chân trang lẻ: ô chữ \thuonghieu (trái)  | Trang (số) (phải)
%  Trang chẵn soi gương trang lẻ; chân trang chẵn ghi \tacgia.
%--- Màu: chỉ dùng thang vàng kim Eduzy
\colorlet{ezNen}{gold1}% dải nền đậm
\colorlet{ezVach1}{gold3}% vạch xiên sát dải nền (đậm hơn)
\colorlet{ezVach2}{gold4}% vạch xiên sát ô chữ (nhạt hơn)
\colorlet{ezO}{gold5}% ô chữ vàng kem
%--- Kích thước
\newlength{\EZdai}\setlength{\EZdai}{1.1cm}% chiều cao dải nền
\newlength{\EZo}\setlength{\EZo}{0.68cm}% chiều cao ô chữ
\newlength{\EZxien}\setlength{\EZxien}{0.55cm}% độ lệch ngang của cạnh xiên
\newlength{\EZvach}\setlength{\EZvach}{0.75cm}% bề rộng mỗi vạch xiên
\newlength{\EZkhe}\setlength{\EZkhe}{0.12cm}% khe trắng giữa vạch và ô chữ
\newlength{\EZdem}\setlength{\EZdem}{0.4cm}% khoảng đệm từ chữ tới cạnh xiên
\newlength{\EZle}\setlength{\EZle}{1.8cm}% lề trái/phải của chữ (bằng lề thân bài)
\newlength{\EZrongmin}\setlength{\EZrongmin}{4.5cm}% ô chữ rộng tối thiểu
\newlength{\EZrong}% bề rộng chữ trong ô (tính tự động)
%--- Font
\newcommand{\EZfontO}{\fontsize{10.5pt}{12pt}\selectfont\bfseries}% chữ trong ô
\newcommand{\EZfontKH}{\fontfamily{qag}\fontsize{11pt}{12pt}\selectfont\bfseries}% \khauhieu cạnh logo
\newcommand{\EZfontTrang}{\fontsize{11.5pt}{12pt}\selectfont\bfseries}% chữ "Trang" và số trang
%--- Đo bề rộng chữ trong ô
\newcommand{\EZdorong}[1]{%
	\settowidth{\EZrong}{\EZfontO #1}%
	\ifdim\EZrong<\EZrongmin\setlength{\EZrong}{\EZrongmin}\fi}
%--- Logo + chữ trắng \khauhieu (dựng sẵn một lần khi vào \begin{document})
\newsavebox{\EZlogobox}
\AtBeginDocument{\sbox{\EZlogobox}{%
		\raisebox{-0.5\height}{\eduzylogo[1cm]}\hspace{0.15cm}%
		\raisebox{-0.5\height}{\color{white}\EZfontKH\khauhieu}}}
%--- Hai vạch xiên + khe trắng. #1: điểm gốc ở mép dưới, #2: hướng (1 sang phải, -1 sang trái),
%    #3: hướng nghiêng (1 kiểu "/", -1 kiểu "\")
\newcommand{\EZvachxien}[3]{%
	\fill[white] (#1) --++({(#3)*\EZxien},\EZo) --++({(#2)*\EZkhe},0) --++({-(#3)*\EZxien},-\EZo) --cycle;
	\fill[ezVach2] ([xshift=(#2)*\EZkhe]#1) --++({(#3)*\EZxien},\EZo) --++({(#2)*\EZvach},0) --++({-(#3)*\EZxien},-\EZo) --cycle;
	\fill[ezVach1] ([xshift=(#2)*(\EZkhe+\EZvach)]#1) --++({(#3)*\EZxien},\EZo) --++({(#2)*\EZvach},0) --++({-(#3)*\EZxien},-\EZo) --cycle;
	\draw[white,line width=0.6pt] ([xshift=(#2)*(\EZkhe+\EZvach)]#1) --++({(#3)*\EZxien},\EZo);
}
%--- Chữ trong ô (căn giữa theo chiều cao ô, không lệch vì dấu tiếng Việt)
\tikzset{ezchu/.style={inner sep=0pt,text=ezNen,text height=0.72em,text depth=0.22em}}

%====================== LOGO CHÌM GIỮA TRANG ======================
%  Vẽ cùng đầu trang nên nằm dưới chữ của thân bài, có trên mọi trang.
%  Tắt: \logochimtat (từ trang đó trở đi), bật lại: \logochimbat.
%  Đổi cỡ, độ đậm: khai báo trước \input, ví dụ \def\EZchimcao{9cm} \def\EZchimdo{0.1}
\providecommand{\EZchimcao}{10cm}% chiều cao logo chìm
\providecommand{\EZchimdo}{0.12}% độ đậm (0 = trong suốt, 1 = rõ hẳn)
\newif\ifEZchim \EZchimtrue
\newcommand{\logochimtat}{\EZchimfalse}
\newcommand{\logochimbat}{\EZchimtrue}
%  Dựng sẵn logo mờ một lần khi vào \begin{document}; transparency group làm mờ cả logo
%  một lượt nên không lộ vệt ở chỗ các dải màu chồng nhau.
\newsavebox{\EZchimtho}
\newsavebox{\EZchimbox}
\AtBeginDocument{%
	\sbox{\EZchimtho}{\eduzylogo[\EZchimcao]}%
	\sbox{\EZchimbox}{\begin{tikzpicture}
			\begin{scope}[transparency group,opacity=\EZchimdo]
				\node[inner sep=0pt] {\usebox{\EZchimtho}};
			\end{scope}
	\end{tikzpicture}}}
\newcommand{\EZvechim}{%
	\ifEZchim
		\node[inner sep=0pt] at (current page.center) {\usebox{\EZchimbox}};
	\fi}

%====================== ĐẦU TRANG - TRANG LẺ ======================
%  (logo bên trái, ô chữ \tenchuong bên phải)
\newcommand{\myfancyheadOdd}{%
	\EZdorong{\tenchuong}%
	\begin{tikzpicture}[remember picture,overlay]
		\coordinate (T) at (current page.north west);
		\coordinate (P) at (current page.north east);
		%--- logo chìm giữa trang
		\EZvechim
		%--- dải nền
		\fill[ezNen] (T) rectangle ([yshift=-\EZdai]P);
		%--- ô chữ: X là góc dưới bên trái, cạnh trái xiên kiểu "/"
		\coordinate (X) at ([xshift=-\EZle-\EZrong-\EZdem-\EZxien,yshift=-\EZdai]P);
		\fill[ezO] (X) --++(\EZxien,\EZo) -- ([yshift=-\EZdai+\EZo]P) -- ([yshift=-\EZdai]P) --cycle;
		\EZvachxien{X}{-1}{1}
		\node[ezchu,anchor=east] at ([xshift=-\EZle,yshift=-\EZdai+0.5*\EZo]P) {\EZfontO\tenchuong};
		%--- logo + tên
		\node[inner sep=0pt,anchor=west] at ([xshift=\EZle,yshift=-0.5*\EZdai]T) {\usebox{\EZlogobox}};
	\end{tikzpicture}%
}

%====================== ĐẦU TRANG - TRANG CHẴN =====================
%  (soi gương: ô chữ \tenchuong bên trái, logo bên phải)
\newcommand{\myfancyheadEven}{%
	\EZdorong{\tenchuong}%
	\begin{tikzpicture}[remember picture,overlay]
		\coordinate (T) at (current page.north west);
		\coordinate (P) at (current page.north east);
		%--- logo chìm giữa trang
		\EZvechim
		%--- dải nền
		\fill[ezNen] (T) rectangle ([yshift=-\EZdai]P);
		%--- ô chữ: X là góc dưới bên phải, cạnh phải xiên kiểu "\"
		\coordinate (X) at ([xshift=\EZle+\EZrong+\EZdem+\EZxien,yshift=-\EZdai]T);
		\fill[ezO] (X) --++(-\EZxien,\EZo) -- ([yshift=-\EZdai+\EZo]T) -- ([yshift=-\EZdai]T) --cycle;
		\EZvachxien{X}{1}{-1}
		\node[ezchu,anchor=west] at ([xshift=\EZle,yshift=-\EZdai+0.5*\EZo]T) {\EZfontO\tenchuong};
		%--- logo + tên
		\node[inner sep=0pt,anchor=east] at ([xshift=-\EZle,yshift=-0.5*\EZdai]P) {\usebox{\EZlogobox}};
	\end{tikzpicture}%
}

%====================== CHÂN TRANG - TRANG LẺ ======================
%  (ô chữ \thuonghieu bên trái, "Trang (số)" bên phải)
\newcommand{\myfancyfootOdd}{%
	\EZdorong{\thuonghieu}%
	\begin{tikzpicture}[remember picture,overlay]
		\coordinate (T) at (current page.south west);
		\coordinate (P) at (current page.south east);
		%--- dải nền
		\fill[ezNen] (T) rectangle ([yshift=\EZdai]P);
		%--- ô chữ sát mép trên dải: X là góc dưới bên phải, cạnh phải xiên kiểu "/"
		\coordinate (X) at ([xshift=\EZle+\EZrong+\EZdem,yshift=\EZdai-\EZo]T);
		\fill[ezO] ([yshift=\EZdai-\EZo]T) -- (X) --++(\EZxien,\EZo) -- ([yshift=\EZdai]T) --cycle;
		\EZvachxien{X}{1}{1}
		\node[ezchu,anchor=west] at ([xshift=\EZle,yshift=\EZdai-0.5*\EZo]T) {\EZfontO\thuonghieu};
		%--- số trang
		\node[circle,draw=white,line width=0.9pt,text=white,inner sep=0pt,minimum size=0.62cm,anchor=east]
		(so) at ([xshift=-\EZle,yshift=0.5*\EZdai]P) {\EZfontTrang\thepage};
		\node[inner sep=0pt,text=white,anchor=base east] (tr) at ([xshift=-0.15cm]so.base west) {\EZfontTrang Trang};
	\end{tikzpicture}%
}

%====================== CHÂN TRANG - TRANG CHẴN ====================
%  (soi gương: "Trang (số)" bên trái, ô chữ \tacgia bên phải)
\newcommand{\myfancyfootEven}{%
	\EZdorong{\tacgia}%
	\begin{tikzpicture}[remember picture,overlay]
		\coordinate (T) at (current page.south west);
		\coordinate (P) at (current page.south east);
		%--- dải nền
		\fill[ezNen] (T) rectangle ([yshift=\EZdai]P);
		%--- ô chữ sát mép trên dải: X là góc dưới bên trái, cạnh trái xiên kiểu "\"
		\coordinate (X) at ([xshift=-\EZle-\EZrong-\EZdem,yshift=\EZdai-\EZo]P);
		\fill[ezO] ([yshift=\EZdai-\EZo]P) -- (X) --++(-\EZxien,\EZo) -- ([yshift=\EZdai]P) --cycle;
		\EZvachxien{X}{-1}{-1}
		\node[ezchu,anchor=east] at ([xshift=-\EZle,yshift=\EZdai-0.5*\EZo]P) {\EZfontO\tacgia};
		%--- số trang
		%    (vòng tròn ẩn so0 nằm giữa dải để lấy đường chân chữ, giống trang lẻ)
		\node[circle,inner sep=0pt,minimum size=0.62cm] (so0) at ([yshift=0.5*\EZdai]T) {\EZfontTrang\phantom{\thepage}};
		\node[inner sep=0pt,text=white,anchor=base west] (tr) at ($(so0.base -| T)+(\EZle,0)$) {\EZfontTrang Trang};
		\node[circle,draw=white,line width=0.9pt,text=white,inner sep=0pt,minimum size=0.62cm,anchor=base west]
		(so) at ([xshift=0.15cm]tr.base east) {\EZfontTrang\thepage};
	\end{tikzpicture}%
}

%--- Giữ tương thích với các file cũ đang gọi \myfancyhead, \myfancyfoot
\newcommand{\myfancyhead}{\myfancyheadOdd}
\newcommand{\myfancyfoot}{\myfancyfootOdd}

%======================= GẮN VÀO TRANG =============================
\pagestyle{fancy}
\fancyhf{}
\fancyhead[LO]{\myfancyheadOdd}
\fancyhead[RE]{\myfancyheadEven}
\fancyfoot[LO]{\myfancyfootOdd}
\fancyfoot[RE]{\myfancyfootEven}
\renewcommand{\headrulewidth}{0pt}
\renewcommand{\footrulewidth}{0pt}
\fancypagestyle{plain}{%
	\fancyhf{}
	\fancyhead[LO]{\myfancyheadOdd}
	\fancyhead[RE]{\myfancyheadEven}
	\fancyfoot[LO]{\myfancyfootOdd}
	\fancyfoot[RE]{\myfancyfootEven}
	\renewcommand{\headrulewidth}{0pt}
	\renewcommand{\footrulewidth}{0pt}
}
%=================== HẾT PHẦN GỘP TỪ eduzy-de.tex ===================
\newcommand{\indam}[1]{
	{\fontsize{12pt}{9pt}\sffamily \selectfont\bfseries\color{gold1} #1}
}
\newcommand{\ind}[1]{
	{\fontsize{14pt}{9pt}\fontfamily{qag} \selectfont\bfseries\color{gold1} #1}
}
\renewcommand{\FalseEX}{\stepcounter{dapan}\textbf{\selectfont\color{gold1}\Alph{dapan}.}\ignorespaces}
\renewtheorem{ex}{\color{gold1}Câu}
\renewcommand{\baselinestretch}{1.1}
\newcommand{\hoac}[1]{
	\left[\begin{aligned}#1\end{aligned}\right.}
\newcommand{\heva}[1]{
	\left\{\begin{aligned}#1\end{aligned}\right.}
\usepackage{esvect}
\def\vec{\vv}
\def\vt{\vv}
\def\overrightarrow{\vv}
%====================== LOGO EDUZY VẼ BẰNG TIKZ ======================
%  Đường viền dò từ logo gốc; dải màu vàng kim đo trực tiếp trên ảnh gốc
%  (nghiêng 21 độ, tô bằng 80 dải hẹp nên hiển thị giống nhau trên mọi trình xem PDF).
%  Dùng: \eduzylogo          -> cao 2.4cm (mặc định)
%        \eduzylogo[1.5cm]   -> cao 1.5cm
\definecolor{ezText}{RGB}{245,200,46}
\newcommand{\eduzylogo}[1][2.4cm]{%
	\resizebox{!}{#1}{%
		\begin{tikzpicture}[x=0.01cm,y=0.01cm]
			\path[use as bounding box] (0,0) rectangle (224,-224);
			% mũ tốt nghiệp + vương miện
			\begin{scope}
				\clip
				(119.4,-176.1)..controls(118.9,-176.1)and(116.8,-176)..(114.6,-176)
				..controls(101.4,-176)and(85.8,-175.9)..(83.7,-175.7)
				..controls(81.8,-175.6)and(81,-175.6)..(80.2,-175.7)
				..controls(78.8,-175.9)and(77.7,-175.9)..(76.2,-175.7)
				..controls(74.9,-175.5)and(70.8,-175.5)..(64.3,-175.7)
				..controls(59.1,-175.9)and(58.3,-175.7)..(58.6,-173.7)
				..controls(58.6,-172.8)and(59.6,-169.6)..(60.3,-168)
				..controls(60.8,-166.9)and(61.4,-165.8)..(62.3,-165)
				..controls(63.1,-164.2)and(64.3,-163.4)..(65,-163.2)
				..controls(65.5,-163.1)and(72.3,-163)..(76.6,-163.1)
				..controls(77.9,-163.1)and(80.1,-163.1)..(81.4,-163)
				..controls(82.8,-162.9)and(86.3,-162.9)..(89.2,-162.9)
				..controls(93,-162.9)and(94.9,-162.8)..(96,-162.7)
				..controls(97.2,-162.5)and(101,-162.5)..(114,-162.4)
				..controls(123,-162.3)and(134.9,-162.2)..(140.2,-162.1)
				..controls(150.1,-162)and(152.2,-162.1)..(153.6,-162.4)
				..controls(154.4,-162.6)and(157.1,-163.8)..(158.1,-164.4)
				..controls(158.8,-164.8)and(159.7,-165.7)..(159.9,-166.3)
				..controls(160,-166.5)and(160.2,-167.1)..(160.2,-167.6)
				..controls(160.3,-168)and(160.6,-169.4)..(160.9,-170.6)
				..controls(161.2,-171.7)and(161.5,-173.1)..(161.6,-173.6)
				..controls(161.8,-174.4)and(161.8,-174.5)..(161.5,-174.9)
				..controls(161.2,-175.6)and(160.8,-175.9)..(159.9,-176)
				..controls(159.2,-176.1)and(154.1,-176.2)..(151.5,-176.2)
				..controls(148.1,-176.1)and(122.6,-176.1)..(121.6,-176.2)
				..controls(120.8,-176.2)and(119.8,-176.2)..(119.4,-176.1)--cycle
				(130.4,-160)..controls(129,-160)and(119.8,-160)..(110.1,-159.9)
				..controls(100.3,-159.8)and(88.5,-159.7)..(83.9,-159.6)
				..controls(77.8,-159.5)and(75.3,-159.5)..(74.5,-159.6)
				..controls(73.9,-159.7)and(72.5,-159.8)..(71.3,-159.8)
				..controls(69.6,-159.8)and(69,-159.8)..(68.5,-159.6)
				..controls(67.6,-159.3)and(66.8,-158.6)..(66.2,-157.6)
				..controls(66,-157.3)and(65.6,-156.5)..(65.3,-156.1)
				..controls(64.6,-155.1)and(64.5,-154.9)..(63.7,-152.9)
				..controls(63.2,-151.6)and(62.9,-150.9)..(61.8,-149.6)
				..controls(61.5,-149.2)and(61.1,-148.5)..(60.9,-147.9)
				..controls(60.7,-147.4)and(60.3,-146.6)..(60.1,-146.1)
				..controls(59.9,-145.5)and(59.5,-144.9)..(59,-144.3)
				..controls(58.5,-143.7)and(58.2,-143.2)..(57.9,-142.5)
				..controls(57.7,-141.9)and(57.3,-141.1)..(57,-140.7)
				..controls(56.6,-140.3)and(56.1,-139.3)..(55.8,-138.5)
				..controls(55.5,-137.8)and(55.1,-136.8)..(54.9,-136.5)
				..controls(54.7,-136.1)and(54.3,-135.2)..(53.9,-134.5)
				..controls(53.6,-133.8)and(53,-132.7)..(52.6,-132)
				..controls(51.6,-130.5)and(51.4,-130.2)..(51.1,-129)
				..controls(50.6,-127.5)and(50.6,-127.4)..(49.7,-126.3)
				..controls(49.2,-125.7)and(48.6,-124.8)..(48.3,-124.2)
				..controls(48,-123.6)and(47.5,-122.7)..(47.3,-122.2)
				..controls(47,-121.8)and(46.5,-120.7)..(46.1,-119.8)
				..controls(45.8,-118.9)and(45.2,-117.8)..(44.9,-117.2)
				..controls(44.2,-116.1)and(43.3,-114.4)..(42.3,-112.2)
				..controls(41.8,-111.3)and(41.3,-110.4)..(41.1,-110.1)
				..controls(40.5,-109.4)and(40.3,-109)..(40,-107.8)
				..controls(39.8,-107)and(39.6,-106.5)..(39.2,-106)
				..controls(38.9,-105.6)and(38.4,-104.7)..(38,-104)
				..controls(37.7,-103.4)and(37.2,-102.5)..(37,-102.1)
				..controls(36.7,-101.7)and(36.3,-100.8)..(36,-100.1)
				..controls(35.6,-98.9)and(35.3,-98.4)..(34.6,-97.6)
				..controls(34.2,-97.1)and(34.2,-96.3)..(34.6,-96)
				..controls(35,-95.7)and(35.5,-95.9)..(36.2,-96.7)
				..controls(36.5,-97)and(37.2,-97.5)..(37.6,-97.7)..controls(38,-98)and(38.6,-98.4)..(39.1,-98.9)
				..controls(39.6,-99.4)and(40.3,-100)..(40.6,-100.3)
				..controls(41.3,-100.8)and(42.4,-101.8)..(43.7,-102.9)
				..controls(44,-103.2)and(44.6,-103.7)..(45,-103.9)
				..controls(45.4,-104.2)and(45.9,-104.6)..(46.1,-104.8)
				..controls(47.1,-105.9)and(48.8,-107.1)..(50.5,-107.9)
				..controls(51.8,-108.6)and(52.3,-109)..(53.4,-110)
				..controls(54.8,-111.2)and(55.6,-111.6)..(57,-112.1)
				..controls(57.5,-112.3)and(58.5,-112.8)..(59.4,-113.3)
				..controls(60.2,-113.7)and(61.3,-114.3)..(61.9,-114.5)
				..controls(62.4,-114.7)and(63.3,-115.1)..(63.8,-115.3)
				..controls(64.4,-115.6)and(65.4,-115.9)..(66.2,-116.1)
				..controls(67,-116.3)and(68.4,-116.6)..(69.2,-116.9)
				..controls(72.3,-117.9)and(72.8,-118)..(78,-118.1)
				..controls(84,-118.3)and(86.1,-118.1)..(88.1,-117.1)
				..controls(89.4,-116.5)and(93,-115)..(93.7,-114.9)
				..controls(94.8,-114.6)and(95.4,-114.2)..(95.9,-113.4)
				..controls(96.2,-113)and(96.7,-112.5)..(97.1,-112.2)
				..controls(98,-111.6)and(101.3,-108.4)..(102.1,-107.2)
				..controls(102.4,-106.8)and(102.9,-106.2)..(103.2,-105.8)
				..controls(103.5,-105.5)and(104,-104.8)..(104.2,-104.2)
				..controls(104.5,-103.6)and(105,-102.8)..(105.3,-102.4)
				..controls(105.6,-102)and(106.1,-101.2)..(106.5,-100.6)
				..controls(107.1,-99.5)and(108.6,-97.5)..(109.5,-96.5)
				..controls(109.9,-96.1)and(110.1,-96)..(110.4,-96)
				..controls(111.1,-96)and(111.6,-96.5)..(111.8,-97.6)
				..controls(112,-98.6)and(112.3,-99.3)..(112.8,-99.9)
				..controls(113.7,-100.9)and(115.2,-103.1)..(115.7,-104.2)
				..controls(116.1,-105)and(116.5,-105.5)..(116.9,-106)
				..controls(117.3,-106.4)and(117.8,-106.9)..(118,-107.3)
				..controls(118.2,-107.6)and(118.6,-108)..(119,-108.3)
				..controls(119.6,-108.8)and(121.1,-110.3)..(121.8,-111.2)
				..controls(122.1,-111.7)and(122.5,-112)..(123.4,-112.4)
				..controls(124.8,-113.2)and(125.8,-113.8)..(127,-114.6)
				..controls(127.5,-114.9)and(128.1,-115.3)..(128.2,-115.4)
				..controls(128.4,-115.5)and(129.3,-115.8)..(130.2,-116)
				..controls(131.1,-116.2)and(132.5,-116.6)..(133.3,-116.9)
				..controls(135.7,-117.9)and(136.1,-117.9)..(142.5,-117.8)
				..controls(145.8,-117.8)and(148.4,-117.7)..(148.8,-117.6)
				..controls(149.2,-117.6)and(150.2,-117.2)..(151,-116.8)
				..controls(152,-116.4)and(152.8,-116)..(153.6,-115.9)
				..controls(154.5,-115.7)and(155.2,-115.4)..(156.3,-114.9)
				..controls(157.1,-114.5)and(158.2,-114)..(158.7,-113.8)
				..controls(159.1,-113.7)and(159.8,-113.3)..(160.2,-113)
				..controls(160.6,-112.7)and(161.3,-112.4)..(161.7,-112.2)
				..controls(162.7,-111.9)and(163.2,-111.7)..(164.1,-111.1)
				..controls(164.5,-110.8)and(165.3,-110.4)..(165.9,-110)
				..controls(166.5,-109.7)and(167.3,-109.2)..(167.6,-108.9)
				..controls(167.9,-108.6)and(168.6,-108.2)..(169,-108)
				..controls(169.4,-107.8)and(170.2,-107.3)..(170.7,-106.9)
				..controls(171.2,-106.5)and(172.1,-105.8)..(172.7,-105.4)
				..controls(173.2,-105)and(174.1,-104.4)..(174.5,-104)
				..controls(174.9,-103.7)and(175.7,-103.1)..(176.2,-102.6)
				..controls(176.7,-102.2)and(177.4,-101.5)..(177.9,-101.2)
				..controls(178.3,-100.8)and(178.8,-100.4)..(179,-100.2)
				..controls(179.2,-100)and(179.6,-99.6)..(180,-99.3)
				..controls(180.4,-99.1)and(180.9,-98.6)..(181.2,-98.4)
				..controls(181.5,-98.1)and(182.1,-97.7)..(182.5,-97.5)
				..controls(183.2,-97.2)and(183.5,-96.9)..(184.3,-95.9)
				..controls(184.6,-95.6)and(184.6,-95.6)..(184.8,-95.7)
				..controls(185,-95.8)and(185.1,-96)..(185.1,-96.2)
				..controls(185.1,-96.5)and(185,-96.7)..(184.6,-97.2)
				..controls(184.1,-97.8)and(184,-97.9)..(183.9,-98.8)
				..controls(183.7,-100.3)and(183.3,-101.3)..(182.9,-101.8)
				..controls(182.6,-102.1)and(182.3,-102.6)..(182.2,-102.9)
				..controls(182,-103.3)and(181.6,-104)..(181.2,-104.6)
				..controls(180.9,-105.2)and(180.5,-106.1)..(180.4,-106.6)
				..controls(179.9,-108)and(179.7,-108.5)..(179,-109.4)
				..controls(178.6,-109.8)and(178.2,-110.5)..(177.9,-111)
				..controls(177.7,-111.5)and(177.3,-112.3)..(177,-112.7)
				..controls(176.6,-113.2)and(176.4,-113.7)..(176.3,-114.2)
				..controls(176.2,-114.7)and(175.9,-115.2)..(175.4,-116)
				..controls(175,-116.6)and(174.4,-117.7)..(174,-118.5)
				..controls(173.7,-119.2)and(173.3,-120.1)..(173,-120.5)
				..controls(172.8,-120.9)and(172.4,-121.8)..(172.1,-122.5)
				..controls(171.8,-123.2)and(171.3,-124.1)..(171,-124.6)
				..controls(170.8,-125.1)and(170.3,-126)..(169.9,-126.6)
				..controls(169.6,-127.3)and(169.2,-128.1)..(169,-128.3)
				..controls(168.8,-128.6)and(168.5,-129.3)..(168.3,-129.8)
				..controls(168.1,-130.3)and(167.8,-131.1)..(167.5,-131.6)
				..controls(166.9,-132.5)and(164.2,-137.9)..(164.2,-138.1)
				..controls(164.2,-138.1)and(163.9,-138.8)..(163.4,-139.7)
				..controls(162.9,-140.5)and(162.4,-141.5)..(162.3,-141.9)
				..controls(161.9,-143)and(161.6,-143.4)..(160.9,-144.1)
				..controls(160.4,-144.6)and(160.3,-144.8)..(160.1,-145.4)
				..controls(160,-145.8)and(159.6,-146.8)..(159.2,-147.8)
				..controls(158.7,-148.8)and(158.1,-150)..(157.9,-150.6)
				..controls(157.6,-151.1)and(157.2,-151.9)..(156.9,-152.4)
				..controls(156.5,-152.8)and(156.1,-153.7)..(155.8,-154.4)
				..controls(155.6,-155)and(155.2,-155.9)..(154.9,-156.3)
				..controls(154.7,-156.7)and(154.4,-157.4)..(154.2,-157.9)
				..controls(154.1,-158.4)and(153.8,-159)..(153.5,-159.3)
				..controls(152.8,-160.1)and(153,-160.1)..(142.4,-160.1)
				..controls(137.2,-160.1)and(131.8,-160.1)..(130.4,-160)--cycle
				(141.5,-103.7)..controls(140.8,-103)and(139.8,-102.1)..(139.2,-101.6)
				..controls(138.5,-101.2)and(137.9,-100.6)..(137.8,-100.4)
				..controls(137,-99.5)and(134,-96.5)..(133.2,-96)
				..controls(132.7,-95.7)and(131.8,-94.9)..(131.1,-94.3)
				..controls(129.4,-92.7)and(127.9,-91.6)..(126.6,-91.1)
				..controls(125.2,-90.5)and(124.1,-90)..(122.9,-89.1)
				..controls(122.1,-88.6)and(121.3,-88.2)..(119.9,-87.7)
				..controls(118.8,-87.3)and(117.5,-86.8)..(116.9,-86.5)
				..controls(114.3,-85.5)and(113.7,-85.4)..(109.7,-85.5)
				..controls(105.3,-85.7)and(103.9,-85.9)..(101.6,-87)
				..controls(100.7,-87.4)and(99.4,-87.8)..(97.9,-88)
				..controls(97,-88.2)and(96.6,-88.3)..(95.8,-88.9)
				..controls(95.4,-89.2)and(94.7,-89.6)..(94.2,-89.8)
				..controls(92.9,-90.3)and(92.3,-90.7)..(90.9,-92.1)
				..controls(90.2,-92.7)and(89.4,-93.3)..(89,-93.5)..controls(87,-94.5)and(85.2,-96)..(83.6,-97.9)
				..controls(83.3,-98.2)and(82.5,-98.9)..(81.8,-99.4)
				..controls(80.7,-100.1)and(80.5,-100.3)..(80,-101.1)
				..controls(78.7,-103.1)and(77.2,-103.9)..(76.4,-103)
				..controls(75.7,-102.2)and(75.7,-103.2)..(75.7,-83.6)
				..controls(75.6,-73.6)and(75.6,-64.8)..(75.5,-64)
				..controls(75.5,-63.2)and(75.5,-61.4)..(75.6,-59.8)--(75.7,-57.1)--(76,-56.8)
				..controls(76.8,-56.2)and(78,-56.3)..(79.8,-57.2)
				..controls(80.4,-57.5)and(81.3,-57.8)..(81.8,-57.9)
				..controls(82.3,-58.1)and(83,-58.4)..(83.6,-58.8)..controls(84.3,-59.2)and(85,-59.5)..(86,-59.8)
				..controls(86.8,-60.1)and(88,-60.6)..(88.7,-60.9)
				..controls(89.3,-61.2)and(90.6,-61.7)..(91.4,-62)
				..controls(92.2,-62.3)and(93.7,-62.9)..(94.6,-63.3)
				..controls(95.6,-63.7)and(96.8,-64.1)..(97.4,-64.3)
				..controls(98,-64.5)and(99,-64.9)..(99.7,-65.2)
				..controls(100.4,-65.6)and(101.5,-66)..(102.3,-66.2)
				..controls(104.2,-66.8)and(105,-67.1)..(106.2,-67.8)
				..controls(107.9,-68.7)and(110.2,-69.4)..(111.3,-69.2)
				..controls(112.5,-69)and(115,-68.2)..(115.6,-67.8)
				..controls(117.3,-66.7)and(117.6,-66.6)..(119.5,-66.2)
				..controls(121.9,-65.8)and(122.1,-65.7)..(123,-65.1)
				..controls(124.1,-64.4)and(124.4,-64.3)..(126,-63.9)
				..controls(127,-63.7)and(127.7,-63.4)..(128.4,-63.1)
				..controls(129,-62.8)and(129.7,-62.5)..(130.6,-62.4)
				..controls(132.1,-62.1)and(134,-61.5)..(135.2,-61)
				..controls(135.7,-60.8)and(136.7,-60.5)..(137.3,-60.3)
				..controls(138.4,-60)and(139.3,-59.7)..(141.1,-59)
				..controls(141.6,-58.7)and(142,-58.6)..(142.5,-58.6)
				..controls(143.2,-58.6)and(143.3,-58.7)..(143.6,-59)
				..controls(144,-59.5)and(144.1,-59.9)..(144,-60.8)
				..controls(144,-61.2)and(144,-71.1)..(144,-82.8)--(143.9,-104.1)--(143.7,-104.5)
				..controls(143.5,-104.7)and(143.3,-104.8)..(143,-104.9)
				..controls(142.6,-104.9)and(142.5,-104.8)..(141.5,-103.7)--cycle
				(27.8,-91.4)..controls(27.1,-90.9)and(26.3,-90)..(26.1,-89.5)
				..controls(26,-89.2)and(25.8,-88.6)..(25.7,-87.9)
				..controls(25.4,-86.2)and(25.6,-84.8)..(26.5,-83.5)
				..controls(27.4,-82.3)and(28.3,-82)..(29.4,-82.3)
				..controls(31.1,-82.8)and(32.1,-84.9)..(32,-87.4)
				..controls(31.9,-89.3)and(31.2,-90.6)..(29.9,-91.4)
				..controls(29.1,-91.9)and(28.6,-91.9)..(27.8,-91.4)--cycle
				(190.1,-91.3)..controls(189.7,-91.2)and(189.3,-90.9)..(189.1,-90.7)
				..controls(188.3,-89.8)and(187.6,-87.7)..(187.6,-86.2)
				..controls(187.6,-85.5)and(187.7,-85.2)..(188.2,-84.4)
				..controls(189.1,-82.6)and(190.1,-81.8)..(191.3,-82.1)
				..controls(192,-82.3)and(193.5,-83.7)..(193.8,-84.4)
				..controls(194.1,-85.1)and(194.2,-87.1)..(194,-88.3)
				..controls(193.8,-89.8)and(193.3,-90.5)..(192.3,-91.3)
				..controls(191.7,-91.7)and(190.9,-91.7)..(190.1,-91.3)--cycle
				(173.7,-83.4)..controls(172.2,-82.5)and(172,-82.5)..(170.6,-82.4)
				..controls(169.7,-82.4)and(168.8,-82.4)..(168.3,-82.5)
				..controls(167.2,-82.7)and(166.9,-82.7)..(166.6,-82.2)
				..controls(166.3,-81.8)and(166.3,-81)..(166.8,-80)
				..controls(167.2,-78.9)and(167.4,-77.4)..(167.7,-74.1)
				..controls(167.8,-72.8)and(167.9,-71.8)..(167.9,-71.7)
				..controls(168,-71.7)and(168,-70.2)..(168,-68.4)..controls(168,-65.4)and(168,-64.9)..(167.7,-63)
				..controls(167.3,-60.5)and(167.1,-59.9)..(166.4,-58.4)
				..controls(166.1,-57.7)and(165.8,-56.9)..(165.7,-56.6)
				..controls(165.7,-56.2)and(165.4,-55.6)..(165.2,-55.2)
				..controls(164.6,-53.8)and(161.9,-50.7)..(161,-50.3)
				..controls(159.8,-49.7)and(157.9,-50)..(155.2,-51.2)
				..controls(154.7,-51.4)and(153.4,-51.8)..(152.2,-52.1)
				..controls(150.8,-52.4)and(149.4,-52.9)..(148.6,-53.2)
				..controls(147.8,-53.6)and(146.7,-54)..(146,-54.2)
				..controls(145.4,-54.4)and(144.5,-54.7)..(144,-54.9)
				..controls(143.6,-55.2)and(142.7,-55.5)..(141.8,-55.7)
				..controls(141,-55.9)and(139.6,-56.3)..(138.7,-56.6)
				..controls(137.8,-57)and(136.3,-57.4)..(135.3,-57.6)
				..controls(133.8,-57.9)and(133.2,-58.1)..(132.1,-58.7)
				..controls(131.1,-59.2)and(130.3,-59.4)..(129.4,-59.6)
				..controls(128.7,-59.8)and(127.8,-60)..(127.4,-60.1)
				..controls(127,-60.3)and(126,-60.5)..(125.1,-60.7)
				..controls(123.2,-61.1)and(121.7,-61.5)..(120.2,-62.1)
				..controls(119.7,-62.3)and(118.1,-62.8)..(116.7,-63.2)
				..controls(114.6,-63.9)and(114,-64.1)..(113,-64.7)
				..controls(111.9,-65.4)and(111.8,-65.4)..(111,-65.5)
				..controls(109.9,-65.5)and(109.3,-65.4)..(107.7,-64.7)
				..controls(107,-64.4)and(106,-64)..(105.3,-63.9)
				..controls(103.5,-63.5)and(102.2,-63)..(101.2,-62.5)
				..controls(100.1,-61.8)and(99.3,-61.5)..(97.8,-61.1)
				..controls(97.2,-61)and(96.2,-60.7)..(95.7,-60.4)
				..controls(95.1,-60.1)and(94.2,-59.8)..(93.5,-59.6)
				..controls(92.9,-59.4)and(91.7,-59)..(91,-58.7)..controls(90.2,-58.4)and(88.9,-58)..(88,-57.7)
				..controls(87.2,-57.5)and(86.1,-57.1)..(85.5,-56.8)
				..controls(85,-56.6)and(84.1,-56.2)..(83.5,-56)..controls(82.9,-55.7)and(81.8,-55.2)..(81,-54.9)
				..controls(80.1,-54.5)and(79.1,-54.2)..(78.7,-54.1)
				..controls(77.9,-53.8)and(77.3,-53.6)..(76.2,-52.9)
				..controls(75.7,-52.6)and(75.2,-52.4)..(74.1,-52.2)
				..controls(72.5,-51.9)and(72.4,-51.8)..(71.3,-51.1)
				..controls(69.9,-50.1)and(69.4,-49.9)..(67.9,-49.8)
				..controls(67,-49.7)and(66.6,-49.6)..(66.3,-49.4)
				..controls(66.1,-49.3)and(65.7,-49)..(65.4,-48.8)
				..controls(65.1,-48.7)and(64.2,-48.4)..(63.6,-48.2)
				..controls(62.9,-48)and(61.7,-47.6)..(60.9,-47.2)
				..controls(60.1,-46.8)and(58.9,-46.4)..(58.2,-46.2)
				..controls(57.6,-46)and(56.8,-45.7)..(56.5,-45.5)
				..controls(54.9,-44.6)and(54.6,-44.5)..(53.2,-44.2)
				..controls(51.7,-43.8)and(51,-43.6)..(49.9,-42.9)
				..controls(49.5,-42.6)and(48.9,-42.4)..(47.5,-42)
				..controls(45.7,-41.6)and(45.6,-41.5)..(44.3,-40.6)
				..controls(43.5,-40.1)and(42.8,-39.6)..(42.7,-39.4)
				..controls(42.4,-38.7)and(43,-38.1)..(44.5,-37.6)..controls(44.9,-37.5)and(46,-37)..(46.9,-36.6)
				..controls(47.9,-36.1)and(48.8,-35.8)..(49.6,-35.6)
				..controls(50.5,-35.4)and(51.3,-35.2)..(52.4,-34.6)
				..controls(53.3,-34.1)and(54.5,-33.7)..(54.9,-33.5)
				..controls(55.4,-33.4)and(56.4,-33)..(57.2,-32.6)
				..controls(58.6,-32)and(60.1,-31.5)..(61.6,-31.2)
				..controls(62.3,-31.1)and(62.7,-31)..(63.2,-30.7)
				..controls(64.1,-30)and(64.3,-29.9)..(65.3,-29.7)
				..controls(65.7,-29.6)and(66.5,-29.2)..(67.1,-28.9)
				..controls(67.6,-28.6)and(68.8,-28.2)..(69.6,-27.9)
				..controls(70.5,-27.7)and(71.6,-27.3)..(72.1,-27)
				..controls(72.6,-26.7)and(73.8,-26.2)..(74.8,-25.9)
				..controls(75.7,-25.5)and(76.7,-25.1)..(77.1,-25)
				..controls(77.4,-24.8)and(78.5,-24.4)..(79.6,-24.1)
				..controls(80.6,-23.8)and(82,-23.2)..(82.8,-22.9)
				..controls(83.5,-22.5)and(84.8,-22)..(85.8,-21.8)
				..controls(86.7,-21.5)and(87.8,-21.1)..(88.1,-20.9)
				..controls(88.4,-20.7)and(89.3,-20.3)..(90.1,-20.1)
				..controls(90.8,-19.8)and(91.9,-19.4)..(92.4,-19.1)
				..controls(93.5,-18.5)and(94.4,-18.2)..(95.4,-18.1)
				..controls(95.7,-18)and(96.4,-17.8)..(96.9,-17.5)
				..controls(97.5,-17.3)and(98.4,-16.9)..(99,-16.6)
				..controls(99.7,-16.3)and(100.9,-15.9)..(101.9,-15.6)
				..controls(102.9,-15.3)and(104.4,-14.7)..(105.3,-14.3)
				..controls(108.1,-13.2)and(108.5,-13.1)..(110.1,-13.1)
				..controls(111.4,-13.1)and(111.5,-13.1)..(112.4,-13.5)
				..controls(113.2,-14)and(115,-14.8)..(116.6,-15.4)
				..controls(116.9,-15.5)and(117.9,-15.8)..(118.8,-16.1)
				..controls(120.9,-16.6)and(121.5,-16.9)..(122.6,-17.9)
				..controls(123,-18.3)and(123.7,-18.8)..(123.9,-18.9)
				..controls(124.6,-19.4)and(126,-19.8)..(127.9,-20.2)
				..controls(129.6,-20.5)and(129.8,-20.6)..(131.1,-21.4)
				..controls(131.5,-21.7)and(132.4,-22.1)..(133.1,-22.3)
				..controls(136.1,-23.2)and(138.3,-24.1)..(140.2,-25)
				..controls(141,-25.5)and(141.5,-25.6)..(142.5,-25.8)
				..controls(145.1,-26.3)and(146,-26.6)..(149.4,-28.1)
				..controls(151.3,-28.9)and(153.5,-29.8)..(154.3,-30.1)
				..controls(155.1,-30.4)and(156.2,-30.8)..(156.6,-31)
				..controls(157.1,-31.2)and(158.2,-31.6)..(159.1,-31.8)
				..controls(160.6,-32.2)and(160.9,-32.3)..(161.7,-32.9)
				..controls(162.4,-33.3)and(162.9,-33.6)..(163.7,-33.8)
				..controls(164.3,-34)and(164.8,-34.1)..(164.9,-34.1)
				..controls(165,-34.1)and(165.6,-34.4)..(166.4,-34.8)
				..controls(167.3,-35.2)and(168.1,-35.4)..(169,-35.6)
				..controls(170,-35.8)and(170.5,-36)..(171.7,-36.6)
				..controls(172.4,-37)and(173.3,-37.4)..(173.6,-37.4)
				..controls(175.6,-37.7)and(176.2,-38)..(176.7,-39.1)
				..controls(176.9,-39.6)and(176.9,-39.7)..(176.7,-40.4)
				..controls(176.4,-41.6)and(175.8,-42.7)..(175.2,-43.4)
				..controls(174.5,-44.1)and(172.6,-45.5)..(171.7,-45.8)
				..controls(170.6,-46.2)and(169.1,-46)..(168.2,-45.4)
				..controls(167.9,-45.2)and(167.3,-44.8)..(166.9,-44.4)
				..controls(166.4,-44.1)and(165.8,-43.6)..(165.4,-43.2)
				..controls(165.1,-42.9)and(164.4,-42.3)..(163.9,-41.9)
				..controls(163.3,-41.5)and(162.7,-40.9)..(162.4,-40.5)
				..controls(161.4,-39.2)and(159.9,-38)..(158.8,-37.8)
				..controls(158.4,-37.8)and(158.3,-37.7)..(158.2,-37.2)
				..controls(158,-36.8)and(157.4,-36.1)..(157.1,-36.1)
				..controls(157,-36.1)and(156.7,-36.1)..(156.5,-36)
				..controls(156.1,-35.9)and(155.9,-35.7)..(155.7,-35.2)
				..controls(155.3,-34.6)and(155.1,-34.4)..(154,-34.1)
				..controls(153.6,-34)and(153.2,-33.7)..(152.9,-33.5)
				..controls(152.4,-32.9)and(151.6,-32.4)..(151.3,-32.4)
				..controls(150.9,-32.4)and(150.4,-32.6)..(150.4,-32.9)
				..controls(150.4,-33.2)and(150.8,-33.4)..(151.4,-33.6)
				..controls(152.1,-33.7)and(152.8,-34.1)..(153.5,-34.9)
				..controls(154.1,-35.5)and(154.3,-35.6)..(155,-35.8)
				..controls(155.7,-36)and(155.7,-36)..(156,-36.6)
				..controls(156.3,-37.3)and(156.7,-37.7)..(157.5,-37.8)
				..controls(158.1,-37.9)and(158.2,-38)..(158.4,-38.8)
				..controls(158.5,-39.4)and(158.6,-39.6)..(159.6,-40.4)
				..controls(160.2,-41)and(160.9,-41.6)..(161.3,-41.8)
				..controls(161.7,-42)and(161.9,-42.3)..(162.2,-42.8)
				..controls(162.4,-43.2)and(162.8,-43.8)..(163.1,-44.1)
				..controls(163.9,-44.8)and(165.1,-46.7)..(165.7,-48)
				..controls(166.2,-49.3)and(166.4,-49.7)..(167.2,-50.5)
				..controls(167.5,-50.8)and(167.9,-51.4)..(168.1,-51.8)
				..controls(168.3,-52.2)and(168.8,-53)..(169.1,-53.6)
				..controls(169.6,-54.3)and(169.8,-54.9)..(169.9,-55.6)
				..controls(170.4,-57.8)and(170.5,-58.1)..(171.2,-59.5)
				..controls(171.6,-60.3)and(172,-61.5)..(172.2,-62.1)
				..controls(172.3,-62.8)and(172.7,-63.7)..(172.9,-64.1)
				..controls(173.2,-64.6)and(173.4,-65.3)..(173.4,-65.6)
				..controls(173.8,-68.8)and(174.2,-70.6)..(174.9,-72.6)
				..controls(175.5,-74.5)and(175.9,-77.1)..(176.1,-80.2)
				..controls(176.3,-83.4)and(175.5,-84.4)..(173.7,-83.4)--cycle;
				\definecolor{ezs}{RGB}{247,179,26}\fill[ezs](-5.7,-156.6)--(-3.4,-157.5)--(61.9,12.7)--(59.7,13.6)--cycle;
				\definecolor{ezs}{RGB}{247,179,22}\fill[ezs](-3.8,-157.3)--(-1.5,-158.2)--(63.9,12)--(61.6,12.9)--cycle;
				\definecolor{ezs}{RGB}{248,179,19}\fill[ezs](-1.8,-158.1)--(0.4,-158.9)--(65.8,11.2)--(63.5,12.1)--cycle;
				\definecolor{ezs}{RGB}{248,178,16}\fill[ezs](0.1,-158.8)--(2.4,-159.7)--(67.7,10.5)--(65.4,11.4)--cycle;
				\definecolor{ezs}{RGB}{249,178,15}\fill[ezs](2,-159.5)--(4.3,-160.4)--(69.6,9.8)--(67.3,10.7)--cycle;
				\definecolor{ezs}{RGB}{251,177,17}\fill[ezs](3.9,-160.3)--(6.2,-161.1)--(71.5,9)--(69.2,9.9)--cycle;
				\definecolor{ezs}{RGB}{252,179,18}\fill[ezs](5.8,-161)--(8.1,-161.9)--(73.4,8.3)--(71.1,9.2)--cycle;
				\definecolor{ezs}{RGB}{252,182,17}\fill[ezs](7.7,-161.7)--(10,-162.6)--(75.3,7.6)--(73,8.5)--cycle;
				\definecolor{ezs}{RGB}{252,185,17}\fill[ezs](9.6,-162.5)--(11.9,-163.4)--(77.2,6.8)--(75,7.7)--cycle;
				\definecolor{ezs}{RGB}{251,189,18}\fill[ezs](11.5,-163.2)--(13.8,-164.1)--(79.2,6.1)--(76.9,7)--cycle;
				\definecolor{ezs}{RGB}{251,191,22}\fill[ezs](13.4,-163.9)--(15.7,-164.8)--(81.1,5.4)--(78.8,6.2)--cycle;
				\definecolor{ezs}{RGB}{252,193,25}\fill[ezs](15.4,-164.7)--(17.6,-165.6)--(83,4.6)--(80.7,5.5)--cycle;
				\definecolor{ezs}{RGB}{253,195,27}\fill[ezs](17.3,-165.4)--(19.6,-166.3)--(84.9,3.9)--(82.6,4.8)--cycle;
				\definecolor{ezs}{RGB}{253,196,31}\fill[ezs](19.2,-166.1)--(21.5,-167)--(86.8,3.2)--(84.5,4)--cycle;
				\definecolor{ezs}{RGB}{254,199,35}\fill[ezs](21.1,-166.9)--(23.4,-167.8)--(88.7,2.4)--(86.4,3.3)--cycle;
				\definecolor{ezs}{RGB}{253,201,38}\fill[ezs](23,-167.6)--(25.3,-168.5)--(90.6,1.7)--(88.3,2.6)--cycle;
				\definecolor{ezs}{RGB}{253,203,40}\fill[ezs](24.9,-168.3)--(27.2,-169.2)--(92.5,1)--(90.3,1.8)--cycle;
				\definecolor{ezs}{RGB}{253,205,42}\fill[ezs](26.8,-169.1)--(29.1,-170)--(94.4,0.2)--(92.2,1.1)--cycle;
				\definecolor{ezs}{RGB}{253,207,45}\fill[ezs](28.7,-169.8)--(31,-170.7)--(96.4,-0.5)--(94.1,0.4)--cycle;
				\definecolor{ezs}{RGB}{253,210,49}\fill[ezs](30.7,-170.5)--(32.9,-171.4)--(98.3,-1.2)--(96,-0.4)--cycle;
				\definecolor{ezs}{RGB}{253,212,54}\fill[ezs](32.6,-171.3)--(34.9,-172.2)--(100.2,-2)--(97.9,-1.1)--cycle;
				\definecolor{ezs}{RGB}{253,214,59}\fill[ezs](34.5,-172)--(36.8,-172.9)--(102.1,-2.7)--(99.8,-1.8)--cycle;
				\definecolor{ezs}{RGB}{253,216,63}\fill[ezs](36.4,-172.7)--(38.7,-173.6)--(104,-3.4)--(101.7,-2.6)--cycle;
				\definecolor{ezs}{RGB}{253,218,65}\fill[ezs](38.3,-173.5)--(40.6,-174.4)--(105.9,-4.2)--(103.6,-3.3)--cycle;
				\definecolor{ezs}{RGB}{254,220,66}\fill[ezs](40.2,-174.2)--(42.5,-175.1)--(107.8,-4.9)--(105.5,-4)--cycle;
				\definecolor{ezs}{RGB}{254,222,68}\fill[ezs](42.1,-175)--(44.4,-175.8)--(109.7,-5.6)--(107.5,-4.8)--cycle;
				\definecolor{ezs}{RGB}{254,224,71}\fill[ezs](44,-175.7)--(46.3,-176.6)--(111.7,-6.4)--(109.4,-5.5)--cycle;
				\definecolor{ezs}{RGB}{254,226,75}\fill[ezs](46,-176.4)--(48.2,-177.3)--(113.6,-7.1)--(111.3,-6.2)--cycle;
				\definecolor{ezs}{RGB}{254,228,79}\fill[ezs](47.9,-177.2)--(50.1,-178)--(115.5,-7.8)--(113.2,-7)--cycle;
				\definecolor{ezs}{RGB}{254,230,82}\fill[ezs](49.8,-177.9)--(52.1,-178.8)--(117.4,-8.6)--(115.1,-7.7)--cycle;
				\definecolor{ezs}{RGB}{254,232,85}\fill[ezs](51.7,-178.6)--(54,-179.5)--(119.3,-9.3)--(117,-8.4)--cycle;
				\definecolor{ezs}{RGB}{254,234,88}\fill[ezs](53.6,-179.4)--(55.9,-180.2)--(121.2,-10)--(118.9,-9.2)--cycle;
				\definecolor{ezs}{RGB}{254,237,91}\fill[ezs](55.5,-180.1)--(57.8,-181)--(123.1,-10.8)--(120.8,-9.9)--cycle;
				\definecolor{ezs}{RGB}{254,239,94}\fill[ezs](57.4,-180.8)--(59.7,-181.7)--(125,-11.5)--(122.8,-10.6)--cycle;
				\definecolor{ezs}{RGB}{254,241,97}\fill[ezs](59.3,-181.6)--(61.6,-182.4)--(127,-12.2)--(124.7,-11.4)--cycle;
				\definecolor{ezs}{RGB}{254,243,101}\fill[ezs](61.2,-182.3)--(63.5,-183.2)--(128.9,-13)--(126.6,-12.1)--cycle;
				\definecolor{ezs}{RGB}{254,245,103}\fill[ezs](63.2,-183)--(65.4,-183.9)--(130.8,-13.7)--(128.5,-12.8)--cycle;
				\definecolor{ezs}{RGB}{254,245,105}\fill[ezs](65.1,-183.8)--(67.4,-184.6)--(132.7,-14.4)--(130.4,-13.6)--cycle;
				\definecolor{ezs}{RGB}{254,245,105}\fill[ezs](67,-184.5)--(69.3,-185.4)--(134.6,-15.2)--(132.3,-14.3)--cycle;
				\definecolor{ezs}{RGB}{254,244,103}\fill[ezs](68.9,-185.2)--(71.2,-186.1)--(136.5,-15.9)--(134.2,-15)--cycle;
				\definecolor{ezs}{RGB}{253,243,101}\fill[ezs](70.8,-186)--(73.1,-186.8)--(138.4,-16.6)--(136.1,-15.8)--cycle;
				\definecolor{ezs}{RGB}{253,241,98}\fill[ezs](72.7,-186.7)--(75,-187.6)--(140.3,-17.4)--(138,-16.5)--cycle;
				\definecolor{ezs}{RGB}{253,239,95}\fill[ezs](74.6,-187.4)--(76.9,-188.3)--(142.2,-18.1)--(140,-17.2)--cycle;
				\definecolor{ezs}{RGB}{253,236,92}\fill[ezs](76.5,-188.2)--(78.8,-189)--(144.2,-18.8)--(141.9,-18)--cycle;
				\definecolor{ezs}{RGB}{253,234,88}\fill[ezs](78.5,-188.9)--(80.7,-189.8)--(146.1,-19.6)--(143.8,-18.7)--cycle;
				\definecolor{ezs}{RGB}{253,232,85}\fill[ezs](80.4,-189.6)--(82.7,-190.5)--(148,-20.3)--(145.7,-19.4)--cycle;
				\definecolor{ezs}{RGB}{253,230,82}\fill[ezs](82.3,-190.4)--(84.6,-191.2)--(149.9,-21)--(147.6,-20.2)--cycle;
				\definecolor{ezs}{RGB}{253,228,78}\fill[ezs](84.2,-191.1)--(86.5,-192)--(151.8,-21.8)--(149.5,-20.9)--cycle;
				\definecolor{ezs}{RGB}{253,226,74}\fill[ezs](86.1,-191.8)--(88.4,-192.7)--(153.7,-22.5)--(151.4,-21.6)--cycle;
				\definecolor{ezs}{RGB}{253,224,70}\fill[ezs](88,-192.6)--(90.3,-193.4)--(155.6,-23.3)--(153.3,-22.4)--cycle;
				\definecolor{ezs}{RGB}{253,222,67}\fill[ezs](89.9,-193.3)--(92.2,-194.2)--(157.5,-24)--(155.3,-23.1)--cycle;
				\definecolor{ezs}{RGB}{253,220,65}\fill[ezs](91.8,-194)--(94.1,-194.9)--(159.5,-24.7)--(157.2,-23.8)--cycle;
				\definecolor{ezs}{RGB}{253,217,63}\fill[ezs](93.7,-194.8)--(96,-195.6)--(161.4,-25.5)--(159.1,-24.6)--cycle;
				\definecolor{ezs}{RGB}{253,215,60}\fill[ezs](95.7,-195.5)--(97.9,-196.4)--(163.3,-26.2)--(161,-25.3)--cycle;
				\definecolor{ezs}{RGB}{253,213,56}\fill[ezs](97.6,-196.2)--(99.9,-197.1)--(165.2,-26.9)--(162.9,-26)--cycle;
				\definecolor{ezs}{RGB}{252,211,53}\fill[ezs](99.5,-197)--(101.8,-197.8)--(167.1,-27.7)--(164.8,-26.8)--cycle;
				\definecolor{ezs}{RGB}{252,209,49}\fill[ezs](101.4,-197.7)--(103.7,-198.6)--(169,-28.4)--(166.7,-27.5)--cycle;
				\definecolor{ezs}{RGB}{252,207,45}\fill[ezs](103.3,-198.4)--(105.6,-199.3)--(170.9,-29.1)--(168.6,-28.2)--cycle;
				\definecolor{ezs}{RGB}{253,205,42}\fill[ezs](105.2,-199.2)--(107.5,-200)--(172.8,-29.9)--(170.6,-29)--cycle;
				\definecolor{ezs}{RGB}{253,202,40}\fill[ezs](107.1,-199.9)--(109.4,-200.8)--(174.7,-30.6)--(172.5,-29.7)--cycle;
				\definecolor{ezs}{RGB}{253,200,38}\fill[ezs](109,-200.6)--(111.3,-201.5)--(176.7,-31.3)--(174.4,-30.4)--cycle;
				\definecolor{ezs}{RGB}{253,198,35}\fill[ezs](111,-201.4)--(113.2,-202.2)--(178.6,-32.1)--(176.3,-31.2)--cycle;
				\definecolor{ezs}{RGB}{253,196,31}\fill[ezs](112.9,-202.1)--(115.2,-203)--(180.5,-32.8)--(178.2,-31.9)--cycle;
				\definecolor{ezs}{RGB}{253,194,28}\fill[ezs](114.8,-202.8)--(117.1,-203.7)--(182.4,-33.5)--(180.1,-32.6)--cycle;
				\definecolor{ezs}{RGB}{253,192,26}\fill[ezs](116.7,-203.6)--(119,-204.5)--(184.3,-34.3)--(182,-33.4)--cycle;
				\definecolor{ezs}{RGB}{253,189,24}\fill[ezs](118.6,-204.3)--(120.9,-205.2)--(186.2,-35)--(183.9,-34.1)--cycle;
				\definecolor{ezs}{RGB}{253,186,20}\fill[ezs](120.5,-205)--(122.8,-205.9)--(188.1,-35.7)--(185.8,-34.8)--cycle;
				\definecolor{ezs}{RGB}{253,184,17}\fill[ezs](122.4,-205.8)--(124.7,-206.7)--(190,-36.5)--(187.8,-35.6)--cycle;
				\definecolor{ezs}{RGB}{252,183,15}\fill[ezs](124.3,-206.5)--(126.6,-207.4)--(192,-37.2)--(189.7,-36.3)--cycle;
				\definecolor{ezs}{RGB}{252,181,14}\fill[ezs](126.3,-207.2)--(128.5,-208.1)--(193.9,-37.9)--(191.6,-37.1)--cycle;
				\definecolor{ezs}{RGB}{252,179,13}\fill[ezs](128.2,-208)--(130.4,-208.9)--(195.8,-38.7)--(193.5,-37.8)--cycle;
				\definecolor{ezs}{RGB}{252,178,12}\fill[ezs](130.1,-208.7)--(132.4,-209.6)--(197.7,-39.4)--(195.4,-38.5)--cycle;
				\definecolor{ezs}{RGB}{251,178,10}\fill[ezs](132,-209.4)--(134.3,-210.3)--(199.6,-40.1)--(197.3,-39.3)--cycle;
				\definecolor{ezs}{RGB}{251,178,10}\fill[ezs](133.9,-210.2)--(136.2,-211.1)--(201.5,-40.9)--(199.2,-40)--cycle;
				\definecolor{ezs}{RGB}{250,177,11}\fill[ezs](135.8,-210.9)--(138.1,-211.8)--(203.4,-41.6)--(201.1,-40.7)--cycle;
				\definecolor{ezs}{RGB}{250,178,11}\fill[ezs](137.7,-211.6)--(140,-212.5)--(205.3,-42.3)--(203.1,-41.5)--cycle;
				\definecolor{ezs}{RGB}{248,180,10}\fill[ezs](139.6,-212.4)--(141.9,-213.3)--(207.3,-43.1)--(205,-42.2)--cycle;
				\definecolor{ezs}{RGB}{247,181,9}\fill[ezs](141.5,-213.1)--(143.8,-214)--(209.2,-43.8)--(206.9,-42.9)--cycle;
				\definecolor{ezs}{RGB}{246,182,9}\fill[ezs](143.5,-213.8)--(145.7,-214.7)--(211.1,-44.5)--(208.8,-43.7)--cycle;
				\definecolor{ezs}{RGB}{246,182,10}\fill[ezs](145.4,-214.6)--(147.7,-215.5)--(213,-45.3)--(210.7,-44.4)--cycle;
			\end{scope}
			% chữ EDUZY (Times đậm, làm dày nét bằng 5 lớp lệch nhẹ; căn đúng khung chữ logo gốc)
			\foreach \dx/\dy in {-0.7/0,0.7/0,0/-0.3,0/0.3,0/0}
			\node[anchor=center,inner sep=0pt,outer sep=0pt,text=ezText,align=center] at (109.6+\dx,-193.6+\dy)
			{\resizebox{1.01cm}{0.22cm}{\fontencoding{T1}\fontfamily{ptm}\fontseries{b}\selectfont EDUZY}};
		\end{tikzpicture}%
	}%
}
%====================== THIẾT LẬP THÂN ĐỀ ======================
%--- Môi trường tự luận {bt}: ex_test 3.3.4 chưa khai báo sẵn nên tạo ở đây
%    (ex_test đã tự móc \fixshowans và \listenumerate cho bt)
\makeatletter
\@ifundefined{bt}{\newtheorem{bt}{\color{gold1}\namebt}}{}
\makeatother
\renewcommand{\sectionmark}[1]{\markright{\thesection\ #1}}
%--- Khung tiêu đề (logo + tối đa 4 dòng; dòng 3, 4 để trống {} thì bỏ qua)
%  Đề thi:    \khungtieude{THI THỬ THPT QUỐC GIA NĂM 2026}{Môn Toán -- Đề số 21}
%                         {(Đề thi gồm: ...)}{Thời gian làm bài: 90 phút (không kể thời gian phát đề).}
%  Chuyên đề: \khungtieude{CHUYÊN ĐỀ 3 -- DÃY SỐ VÀ CÔNG THỨC TRUY HỒI}
%                         {Bồi dưỡng học sinh giỏi Toán THPT}{...}{}
\newcommand{\khungtieude}[4]{%
	\begin{tcolorbox}[colframe=gold1,colback=gold5!45,boxrule=0.5mm]
		\begin{minipage}[c]{3cm}
			\centerline{\eduzylogo[2.4cm]}
		\end{minipage}\hspace{0.5cm}
		\begin{minipage}[c]{12.5cm}
			\vskip 0.15 cm
			\centerline{\ind{#1}}
			\vskip 0.15 cm
			\centerline{\textbf{#2}}
			\if\relax\detokenize{#3}\relax\else
				\vskip 0.15 cm
				\centerline{\textit{#3}}
			\fi
			\if\relax\detokenize{#4}\relax\else
				\vskip 0.15 cm
				\centerline{\textit{#4}}
			\fi
			\vskip 0.15 cm
		\end{minipage}
	\end{tcolorbox}
}
%====================== BỔ SUNG CHO BÀI HỌC, CHUYÊN ĐỀ ======================
\usepackage[shortlabels]{enumitem}
\usepackage{tabularx}
%--- Danh sách gọn: \begin{enumerate}[gon] ... \end{enumerate}
\SetEnumitemKey{gon}{leftmargin=*,itemsep=1pt,parsep=0pt,topsep=3pt}
%--- Khung lý thuyết: \begin{khung}{} ... \end{khung}; tiêu đề chỉ dùng khi thật ngắn ("Mẹo", "Lỗi thường gặp")
\newtcolorbox{khung}[1]{breakable,enhanced,colframe=gold1,colback=gold5!45,
	boxrule=0.4mm,fonttitle=\bfseries\sffamily,coltitle=white,colbacktitle=gold1,
	attach boxed title to top left={xshift=6mm,yshift=-3mm},
	boxed title style={arc=1mm,boxrule=0.3mm},title={#1}}
%--- Tiêu đề mục nhỏ trong một dạng: \muc{Dấu hiệu}, \muc{Phương pháp}, \muc{Bài tự luyện}
\makeatletter
\newcommand{\muc}[1]{\par\medskip\noindent\indam{#1}\par\nobreak\smallskip\@afterindentfalse\@afterheading}
\makeatother
%--- Ô công thức cần nhớ (căn giữa, viền vàng, vừa khít nội dung)
%     \begin{center}\congthuc{$Q(n)=(n-1)!$}\end{center}
\newtcbox{\congthuc}{on line,colframe=gold1,colback=gold5!45,boxrule=0.5pt,arc=2pt,
	left=8pt,right=8pt,top=4pt,bottom=4pt}
%--- Công tắc lời giải bài tự luyện (ví dụ mẫu luôn dùng \loigiai)
%     \loigiaiTL[{đáp số}]{lời giải đầy đủ}
%     \TLfull -> in lời giải đầy đủ;  \TLdapso -> chỉ in một dòng "Đáp số";  \TLan -> ẩn hoàn toàn
%     (\CheDo là dethi thì mọi lời giải và đáp số đều ẩn)
\makeatletter
\newif\ifTLhien   \TLhientrue
\newif\ifTLdapso  \TLdapsofalse
\newcommand{\TLfull}{\TLhientrue\TLdapsofalse}
\newcommand{\TLdapso}{\TLhientrue\TLdapsotrue}
\newcommand{\TLan}{\TLhienfalse}
\newcommand{\loigiaiTL}[2][]{%
	\ifTLhien
		\ifTLdapso
			\def\@tlarg{#1}%
			\ifx\@tlarg\@empty\else
				\ifdefstring{\kindOPTN}{dethi}{}{\par\noindent{\bfseries\color{gold1}Đáp số:} #1\par}%
			\fi
		\else
			\loigiai{#2}%
		\fi
	\fi}
\makeatother
\TLdapso% mặc định: chỉ in đáp số bài tự luyện
%====================== BẢNG ĐÚNG/SAI (\choiceTFt) ======================
%  Hàng tiêu đề nền nâu đồng chữ trắng (Phát biểu | Đ | S), các hàng tô xen kẽ kem / trắng,
%  đường kẻ vàng kim, nhãn a) b) c) d) màu nâu đồng. Bản có đáp án: cột được chọn có vòng tròn
%  nâu đồng với dấu tích trắng (câu đúng tích ở cột Đúng, câu sai tích ở cột Sai).
%  Tiêu đề ghi đủ "Đúng | Sai": \OPTN{viettat=0} sau \input (cột Đúng/Sai rộng ra, phát biểu hẹp lại).
%  Dùng bảng gốc của ex_test: \def\BangDSgoc{1} trước \input{preamble-eduzy-v2}.
\providecommand{\BangDSgoc}{0}
\ifnum\BangDSgoc=0
\usepackage{colortbl,pifont}
\colorlet{ezDSke}{gold2!60!gold1}% đường kẻ bảng (nét 0,4pt như ex_test để không đổi chiều cao bảng)
\colorlet{ezDSle}{gold5!60}% hàng a), c)
\colorlet{ezDSchan}{white}% hàng b), d)
%--- dấu chọn ở cột Đúng/Sai
%    (bề rộng 0, căn giữa ô: không làm đổi bề rộng cột của ex_test, không báo tràn lề)
\newcommand{\EZdauchon}{\makebox[0pt]{\tikz[baseline=-1ex]\node[circle,fill=gold1,inner sep=0pt,
	minimum size=1.2em,text=white,font=\small]{\ding{51}};}}
\def\tickT{\EZdauchon}\def\tickF{\EZdauchon}
\AtBeginDocument{\sbox{\ccTFTab}{\EZdauchon}}
%--- nhãn a) b) c) d) màu nâu đồng (ở bản lời giải, nhãn câu đúng vẫn đóng khung như ex_test)
\renewcommand{\FalseTF}{\stepcounter{dapan}\textbf{\color{gold1}\DapAnTF\sepTF}\ignorespaces}
%--- trước khi dựng bảng: kẻ màu vàng kim (bề rộng cột giữ nguyên như ex_test); sau bảng trả lại màu đen
\newcount\EZhangTF
\newcommand{\EZbatdauTF}{%
	\global\EZhangTF=0\relax
	\arrayrulecolor{ezDSke}}
\newcommand{\EZketthucTF}{\arrayrulecolor{black}}
%--- hàng tiêu đề nền nâu đồng, chữ trắng
\makeatletter
\newcommand{\EZo@TF}[1]{\centering\bfseries\color{white}#1}
\newcommand{\EZtieudeTF}{%
	\ifcase\cotDS
		\ifnum\cotSTT=1\relax
			\def\rowHeader{\hline\rowcolor{gold1}\EZo@TF{TT}&\EZo@TF{\phatbieu}\tabularnewline}%
		\else
			\def\rowHeader{\hline\rowcolor{gold1}\EZo@TF{\phatbieu}\tabularnewline}%
		\fi
	\or
		\ifnum\cotSTT=1\relax
			\def\rowHeader{\hline\rowcolor{gold1}\EZo@TF{TT}&\EZo@TF{\phatbieu}&\EZo@TF{\Dung}\tabularnewline}%
		\else
			\def\rowHeader{\hline\rowcolor{gold1}\EZo@TF{\phatbieu}&\EZo@TF{\Dung}\tabularnewline}%
		\fi
	\or
		\ifnum\cotSTT=1\relax
			\def\rowHeader{\hline\rowcolor{gold1}\EZo@TF{TT}&\EZo@TF{\phatbieu}&\EZo@TF{\Dung}&\EZo@TF{\Sai}\tabularnewline}%
		\else
			\def\rowHeader{\hline\rowcolor{gold1}\EZo@TF{\phatbieu}&\EZo@TF{\Dung}&\EZo@TF{\Sai}\tabularnewline}%
		\fi
	\fi}
%--- tô xen kẽ từng hàng (gọi ngay sau \hline đầu mỗi hàng)
\newcommand{\EZmauhangTF}{%
	\noalign{\global\advance\EZhangTF\@ne}%
	\ifodd\EZhangTF\rowcolor{ezDSle}\else\rowcolor{ezDSchan}\fi}
%--- gắn vào \choiceTFt / \writeChoiceTFt của ex_test
\xpretocmd{\choiceTFt}{\EZbatdauTF}{}{\typeout{[preamble-eduzy-v2] Không gắn được bảng Đúng/Sai (đầu bảng)}}
\xpatchcmd{\choiceTFt}{\ifnum\dongPB=0\def\rowHeader{}\fi}{\EZtieudeTF\ifnum\dongPB=0\def\rowHeader{}\fi}
	{}{\typeout{[preamble-eduzy-v2] Không gắn được bảng Đúng/Sai (hàng tiêu đề)}}
\xpatchcmd{\writeChoiceTFt}{\hline\xdef\dotEX{\sdotEX}}{\hline\EZmauhangTF\xdef\dotEX{\sdotEX}}
	{}{\typeout{[preamble-eduzy-v2] Không gắn được bảng Đúng/Sai (tô hàng)}}
\xpatchcmd{\writeChoiceTFt}{{\end{longtable}}}{{\end{longtable}\EZketthucTF}}
	{}{\typeout{[preamble-eduzy-v2] Không gắn được bảng Đúng/Sai (cuối longtable)}}
\xpatchcmd{\writeChoiceTFt}{{\end{xtabular}\end{center}\medskip\noindent}}{{\end{xtabular}\end{center}\EZketthucTF\medskip\noindent}}
	{}{\typeout{[preamble-eduzy-v2] Không gắn được bảng Đúng/Sai (cuối xtabular)}}
\fi
%=====================================================================
\renewcommand{\parallel}{\mathrel{/\!\!/}}

%  trong file đề (tools/dan_trang/dan_trang.py tự ghi khi đặt \NeedspaceCau):
%    \def\TachCauDS{1}   câu Đúng/Sai được tách: câu dẫn (kèm hình) ở cuối trang, bảng Phát biểu sang trang sau
%                        (câu dẫn và hình không tách nhau, bảng không tách giữa hai trang)
%    \def\TachCauLG{1}   bản lời giải (\CheDo khác dethi): phần đề của câu cao hơn 1/3 trang được phép tách
%  Lệnh dùng trong file đề (giá trị do tools/dan_trang/dan_trang.py đo ở bản dethi):
%    \NeedspaceCau[<cao câu dẫn + hình>]{<cao cả khối câu>}   ngay trên \begin{ex} (câu đầu phần: trên tiêu đề phần)
%        công tắc tắt: đúng bằng \Needspace{<cao cả khối câu>} như cũ;
%        \TachCauDS=1 và có [..] (câu Đúng/Sai được chọn tách): chỉ giữ câu dẫn + hình;
%        \TachCauLG=1 ở bản lời giải: đòi tối đa 1/3 \textheight.
%    \NeedspaceBang{<cao bảng Phát biểu>}%   ngay trước \choiceTFt: \TachCauDS=1 thì giữ trọn bảng trên một trang;
%        công tắc tắt: không làm gì.
\providecommand{\TachCauDS}{0}
\providecommand{\TachCauLG}{0}
\newlength{\EZnsCau}
\newcommand{\NeedspaceCau}[2][]{%
	\setlength{\EZnsCau}{#2}%
	\ifnum\TachCauDS=1 \if\relax\detokenize{#1}\relax\else\setlength{\EZnsCau}{#1}\fi\fi
	\ifnum\TachCauLG=1 \ifdefstring{\kindOPTN}{dethi}{}{\ifdim\EZnsCau>0.3333\textheight\setlength{\EZnsCau}{0.3333\textheight}\fi}\fi
	\Needspace{\EZnsCau}}
\newcommand{\NeedspaceBang}[1]{\ifnum\TachCauDS=1 \Needspace{#1}\fi}

%--- Thang màu: đậm -> nhạt
\definecolor{gold1}{RGB}{126,96,0}    % nâu đồng đậm  #7E6000
\definecolor{gold2}{RGB}{192,144,3}   % vàng kim      #C09003
\definecolor{gold3}{RGB}{255,215,99}  % vàng nhạt     #FFD763
\definecolor{gold4}{RGB}{255,231,153} % vàng rất nhạt #FFE799
\definecolor{gold5}{RGB}{255,243,205} % vàng kem      #FFF3CD
%--- Ghi đè tông xanh/đỏ cũ
\definecolor{myblue}{RGB}{126,96,0}
\colorlet{blue}{gold1}
\colorlet{cyan}{gold2}
\colorlet{teal}{gold2}

%--- Tên chương mặc định (file đề nên tự khai báo lại)
\providecommand{\tenchuong}{}
%--- Chân trang
\providecommand{\thuonghieu}{TOÁN 12 -- Eduzy.com.vn}
\providecommand{\tacgia}{Tô Gia Bảo -- 0353703783}

%--- Chữ cạnh logo ở đầu trang
\providecommand{\khauhieu}{Eduzy.com.vn}

%====================== THÔNG SỐ ĐẦU TRANG, CHÂN TRANG (v2) ======================
%  Dải nền đậm tràn ngang trang; ô chữ vàng kem có cạnh xiên, kèm 2 vạch xiên chuyển màu.
%  Đầu trang lẻ : logo + \khauhieu (trái)   | ô chữ \tenchuong (phải)
%  Chân trang lẻ: ô chữ \thuonghieu (trái)  | Trang (số) (phải)
%  Trang chẵn soi gương trang lẻ; chân trang chẵn ghi \tacgia.
%--- Màu: chỉ dùng thang vàng kim Eduzy
\colorlet{ezNen}{gold1}% dải nền đậm
\colorlet{ezVach1}{gold3}% vạch xiên sát dải nền (đậm hơn)
\colorlet{ezVach2}{gold4}% vạch xiên sát ô chữ (nhạt hơn)
\colorlet{ezO}{gold5}% ô chữ vàng kem
%--- Kích thước
\newlength{\EZdai}\setlength{\EZdai}{1.1cm}% chiều cao dải nền
\newlength{\EZo}\setlength{\EZo}{0.68cm}% chiều cao ô chữ
\newlength{\EZxien}\setlength{\EZxien}{0.55cm}% độ lệch ngang của cạnh xiên
\newlength{\EZvach}\setlength{\EZvach}{0.75cm}% bề rộng mỗi vạch xiên
\newlength{\EZkhe}\setlength{\EZkhe}{0.12cm}% khe trắng giữa vạch và ô chữ
\newlength{\EZdem}\setlength{\EZdem}{0.4cm}% khoảng đệm từ chữ tới cạnh xiên
\newlength{\EZle}\setlength{\EZle}{1.8cm}% lề trái/phải của chữ (bằng lề thân bài)
\newlength{\EZrongmin}\setlength{\EZrongmin}{4.5cm}% ô chữ rộng tối thiểu
\newlength{\EZrong}% bề rộng chữ trong ô (tính tự động)
%--- Font
\newcommand{\EZfontO}{\fontsize{10.5pt}{12pt}\selectfont\bfseries}% chữ trong ô
\newcommand{\EZfontKH}{\fontfamily{qag}\fontsize{11pt}{12pt}\selectfont\bfseries}% \khauhieu cạnh logo
\newcommand{\EZfontTrang}{\fontsize{11.5pt}{12pt}\selectfont\bfseries}% chữ "Trang" và số trang
%--- Đo bề rộng chữ trong ô
\newcommand{\EZdorong}[1]{%
	\settowidth{\EZrong}{\EZfontO #1}%
	\ifdim\EZrong<\EZrongmin\setlength{\EZrong}{\EZrongmin}\fi}
%--- Logo + chữ trắng \khauhieu (dựng sẵn một lần khi vào \begin{document})
\newsavebox{\EZlogobox}
\AtBeginDocument{\sbox{\EZlogobox}{%
		\raisebox{-0.5\height}{\eduzylogo[1cm]}\hspace{0.15cm}%
		\raisebox{-0.5\height}{\color{white}\EZfontKH\khauhieu}}}
%--- Hai vạch xiên + khe trắng. #1: điểm gốc ở mép dưới, #2: hướng (1 sang phải, -1 sang trái),
%    #3: hướng nghiêng (1 kiểu "/", -1 kiểu "\")
\newcommand{\EZvachxien}[3]{%
	\fill[white] (#1) --++({(#3)*\EZxien},\EZo) --++({(#2)*\EZkhe},0) --++({-(#3)*\EZxien},-\EZo) --cycle;
	\fill[ezVach2] ([xshift=(#2)*\EZkhe]#1) --++({(#3)*\EZxien},\EZo) --++({(#2)*\EZvach},0) --++({-(#3)*\EZxien},-\EZo) --cycle;
	\fill[ezVach1] ([xshift=(#2)*(\EZkhe+\EZvach)]#1) --++({(#3)*\EZxien},\EZo) --++({(#2)*\EZvach},0) --++({-(#3)*\EZxien},-\EZo) --cycle;
	\draw[white,line width=0.6pt] ([xshift=(#2)*(\EZkhe+\EZvach)]#1) --++({(#3)*\EZxien},\EZo);
}
%--- Chữ trong ô (căn giữa theo chiều cao ô, không lệch vì dấu tiếng Việt)
\tikzset{ezchu/.style={inner sep=0pt,text=ezNen,text height=0.72em,text depth=0.22em}}

%====================== LOGO CHÌM GIỮA TRANG ======================
%  Vẽ cùng đầu trang nên nằm dưới chữ của thân bài, có trên mọi trang.
%  Tắt: \logochimtat (từ trang đó trở đi), bật lại: \logochimbat.
%  Đổi cỡ, độ đậm: khai báo trước \input, ví dụ \def\EZchimcao{9cm} \def\EZchimdo{0.1}
\providecommand{\EZchimcao}{10cm}% chiều cao logo chìm
\providecommand{\EZchimdo}{0.12}% độ đậm (0 = trong suốt, 1 = rõ hẳn)
\newif\ifEZchim \EZchimtrue
\newcommand{\logochimtat}{\EZchimfalse}
\newcommand{\logochimbat}{\EZchimtrue}
%  Dựng sẵn logo mờ một lần khi vào \begin{document}; transparency group làm mờ cả logo
%  một lượt nên không lộ vệt ở chỗ các dải màu chồng nhau.
\newsavebox{\EZchimtho}
\newsavebox{\EZchimbox}
\AtBeginDocument{%
	\sbox{\EZchimtho}{\eduzylogo[\EZchimcao]}%
	\sbox{\EZchimbox}{\begin{tikzpicture}
			\begin{scope}[transparency group,opacity=\EZchimdo]
				\node[inner sep=0pt] {\usebox{\EZchimtho}};
			\end{scope}
	\end{tikzpicture}}}
\newcommand{\EZvechim}{%
	\ifEZchim
		\node[inner sep=0pt] at (current page.center) {\usebox{\EZchimbox}};
	\fi}

%====================== ĐẦU TRANG - TRANG LẺ ======================
%  (logo bên trái, ô chữ \tenchuong bên phải)
\newcommand{\myfancyheadOdd}{%
	\EZdorong{\tenchuong}%
	\begin{tikzpicture}[remember picture,overlay]
		\coordinate (T) at (current page.north west);
		\coordinate (P) at (current page.north east);
		%--- logo chìm giữa trang
		\EZvechim
		%--- dải nền
		\fill[ezNen] (T) rectangle ([yshift=-\EZdai]P);
		%--- ô chữ: X là góc dưới bên trái, cạnh trái xiên kiểu "/"
		\coordinate (X) at ([xshift=-\EZle-\EZrong-\EZdem-\EZxien,yshift=-\EZdai]P);
		\fill[ezO] (X) --++(\EZxien,\EZo) -- ([yshift=-\EZdai+\EZo]P) -- ([yshift=-\EZdai]P) --cycle;
		\EZvachxien{X}{-1}{1}
		\node[ezchu,anchor=east] at ([xshift=-\EZle,yshift=-\EZdai+0.5*\EZo]P) {\EZfontO\tenchuong};
		%--- logo + tên
		\node[inner sep=0pt,anchor=west] at ([xshift=\EZle,yshift=-0.5*\EZdai]T) {\usebox{\EZlogobox}};
	\end{tikzpicture}%
}

%====================== ĐẦU TRANG - TRANG CHẴN =====================
%  (soi gương: ô chữ \tenchuong bên trái, logo bên phải)
\newcommand{\myfancyheadEven}{%
	\EZdorong{\tenchuong}%
	\begin{tikzpicture}[remember picture,overlay]
		\coordinate (T) at (current page.north west);
		\coordinate (P) at (current page.north east);
		%--- logo chìm giữa trang
		\EZvechim
		%--- dải nền
		\fill[ezNen] (T) rectangle ([yshift=-\EZdai]P);
		%--- ô chữ: X là góc dưới bên phải, cạnh phải xiên kiểu "\"
		\coordinate (X) at ([xshift=\EZle+\EZrong+\EZdem+\EZxien,yshift=-\EZdai]T);
		\fill[ezO] (X) --++(-\EZxien,\EZo) -- ([yshift=-\EZdai+\EZo]T) -- ([yshift=-\EZdai]T) --cycle;
		\EZvachxien{X}{1}{-1}
		\node[ezchu,anchor=west] at ([xshift=\EZle,yshift=-\EZdai+0.5*\EZo]T) {\EZfontO\tenchuong};
		%--- logo + tên
		\node[inner sep=0pt,anchor=east] at ([xshift=-\EZle,yshift=-0.5*\EZdai]P) {\usebox{\EZlogobox}};
	\end{tikzpicture}%
}

%====================== CHÂN TRANG - TRANG LẺ ======================
%  (ô chữ \thuonghieu bên trái, "Trang (số)" bên phải)
\newcommand{\myfancyfootOdd}{%
	\EZdorong{\thuonghieu}%
	\begin{tikzpicture}[remember picture,overlay]
		\coordinate (T) at (current page.south west);
		\coordinate (P) at (current page.south east);
		%--- dải nền
		\fill[ezNen] (T) rectangle ([yshift=\EZdai]P);
		%--- ô chữ sát mép trên dải: X là góc dưới bên phải, cạnh phải xiên kiểu "/"
		\coordinate (X) at ([xshift=\EZle+\EZrong+\EZdem,yshift=\EZdai-\EZo]T);
		\fill[ezO] ([yshift=\EZdai-\EZo]T) -- (X) --++(\EZxien,\EZo) -- ([yshift=\EZdai]T) --cycle;
		\EZvachxien{X}{1}{1}
		\node[ezchu,anchor=west] at ([xshift=\EZle,yshift=\EZdai-0.5*\EZo]T) {\EZfontO\thuonghieu};
		%--- số trang
		\node[circle,draw=white,line width=0.9pt,text=white,inner sep=0pt,minimum size=0.62cm,anchor=east]
		(so) at ([xshift=-\EZle,yshift=0.5*\EZdai]P) {\EZfontTrang\thepage};
		\node[inner sep=0pt,text=white,anchor=base east] (tr) at ([xshift=-0.15cm]so.base west) {\EZfontTrang Trang};
	\end{tikzpicture}%
}

%====================== CHÂN TRANG - TRANG CHẴN ====================
%  (soi gương: "Trang (số)" bên trái, ô chữ \tacgia bên phải)
\newcommand{\myfancyfootEven}{%
	\EZdorong{\tacgia}%
	\begin{tikzpicture}[remember picture,overlay]
		\coordinate (T) at (current page.south west);
		\coordinate (P) at (current page.south east);
		%--- dải nền
		\fill[ezNen] (T) rectangle ([yshift=\EZdai]P);
		%--- ô chữ sát mép trên dải: X là góc dưới bên trái, cạnh trái xiên kiểu "\"
		\coordinate (X) at ([xshift=-\EZle-\EZrong-\EZdem,yshift=\EZdai-\EZo]P);
		\fill[ezO] ([yshift=\EZdai-\EZo]P) -- (X) --++(-\EZxien,\EZo) -- ([yshift=\EZdai]P) --cycle;
		\EZvachxien{X}{-1}{-1}
		\node[ezchu,anchor=east] at ([xshift=-\EZle,yshift=\EZdai-0.5*\EZo]P) {\EZfontO\tacgia};
		%--- số trang
		%    (vòng tròn ẩn so0 nằm giữa dải để lấy đường chân chữ, giống trang lẻ)
		\node[circle,inner sep=0pt,minimum size=0.62cm] (so0) at ([yshift=0.5*\EZdai]T) {\EZfontTrang\phantom{\thepage}};
		\node[inner sep=0pt,text=white,anchor=base west] (tr) at ($(so0.base -| T)+(\EZle,0)$) {\EZfontTrang Trang};
		\node[circle,draw=white,line width=0.9pt,text=white,inner sep=0pt,minimum size=0.62cm,anchor=base west]
		(so) at ([xshift=0.15cm]tr.base east) {\EZfontTrang\thepage};
	\end{tikzpicture}%
}

%--- Giữ tương thích với các file cũ đang gọi \myfancyhead, \myfancyfoot
\newcommand{\myfancyhead}{\myfancyheadOdd}
\newcommand{\myfancyfoot}{\myfancyfootOdd}

%======================= GẮN VÀO TRANG =============================
\pagestyle{fancy}
\fancyhf{}
\fancyhead[LO]{\myfancyheadOdd}
\fancyhead[RE]{\myfancyheadEven}
\fancyfoot[LO]{\myfancyfootOdd}
\fancyfoot[RE]{\myfancyfootEven}
\renewcommand{\headrulewidth}{0pt}
\renewcommand{\footrulewidth}{0pt}
\fancypagestyle{plain}{%
	\fancyhf{}
	\fancyhead[LO]{\myfancyheadOdd}
	\fancyhead[RE]{\myfancyheadEven}
	\fancyfoot[LO]{\myfancyfootOdd}
	\fancyfoot[RE]{\myfancyfootEven}
	\renewcommand{\headrulewidth}{0pt}
	\renewcommand{\footrulewidth}{0pt}
}
%=================== HẾT PHẦN GỘP TỪ eduzy-de.tex ===================
\newcommand{\indam}[1]{
	{\fontsize{12pt}{9pt}\sffamily \selectfont\bfseries\color{gold1} #1}
}
\newcommand{\ind}[1]{
	{\fontsize{14pt}{9pt}\fontfamily{qag} \selectfont\bfseries\color{gold1} #1}
}
\renewcommand{\FalseEX}{\stepcounter{dapan}\textbf{\selectfont\color{gold1}\Alph{dapan}.}\ignorespaces}
\renewtheorem{ex}{\color{gold1}Câu}
\renewcommand{\baselinestretch}{1.1}
\newcommand{\hoac}[1]{
	\left[\begin{aligned}#1\end{aligned}\right.}
\newcommand{\heva}[1]{
	\left\{\begin{aligned}#1\end{aligned}\right.}
\usepackage{esvect}
\def\vec{\vv}
\def\vt{\vv}
\def\overrightarrow{\vv}
%====================== LOGO EDUZY VẼ BẰNG TIKZ ======================
%  Đường viền dò từ logo gốc; dải màu vàng kim đo trực tiếp trên ảnh gốc
%  (nghiêng 21 độ, tô bằng 80 dải hẹp nên hiển thị giống nhau trên mọi trình xem PDF).
%  Dùng: \eduzylogo          -> cao 2.4cm (mặc định)
%        \eduzylogo[1.5cm]   -> cao 1.5cm
\definecolor{ezText}{RGB}{245,200,46}
\newcommand{\eduzylogo}[1][2.4cm]{%
	\resizebox{!}{#1}{%
		\begin{tikzpicture}[x=0.01cm,y=0.01cm]
			\path[use as bounding box] (0,0) rectangle (224,-224);
			% mũ tốt nghiệp + vương miện
			\begin{scope}
				\clip
				(119.4,-176.1)..controls(118.9,-176.1)and(116.8,-176)..(114.6,-176)
				..controls(101.4,-176)and(85.8,-175.9)..(83.7,-175.7)
				..controls(81.8,-175.6)and(81,-175.6)..(80.2,-175.7)
				..controls(78.8,-175.9)and(77.7,-175.9)..(76.2,-175.7)
				..controls(74.9,-175.5)and(70.8,-175.5)..(64.3,-175.7)
				..controls(59.1,-175.9)and(58.3,-175.7)..(58.6,-173.7)
				..controls(58.6,-172.8)and(59.6,-169.6)..(60.3,-168)
				..controls(60.8,-166.9)and(61.4,-165.8)..(62.3,-165)
				..controls(63.1,-164.2)and(64.3,-163.4)..(65,-163.2)
				..controls(65.5,-163.1)and(72.3,-163)..(76.6,-163.1)
				..controls(77.9,-163.1)and(80.1,-163.1)..(81.4,-163)
				..controls(82.8,-162.9)and(86.3,-162.9)..(89.2,-162.9)
				..controls(93,-162.9)and(94.9,-162.8)..(96,-162.7)
				..controls(97.2,-162.5)and(101,-162.5)..(114,-162.4)
				..controls(123,-162.3)and(134.9,-162.2)..(140.2,-162.1)
				..controls(150.1,-162)and(152.2,-162.1)..(153.6,-162.4)
				..controls(154.4,-162.6)and(157.1,-163.8)..(158.1,-164.4)
				..controls(158.8,-164.8)and(159.7,-165.7)..(159.9,-166.3)
				..controls(160,-166.5)and(160.2,-167.1)..(160.2,-167.6)
				..controls(160.3,-168)and(160.6,-169.4)..(160.9,-170.6)
				..controls(161.2,-171.7)and(161.5,-173.1)..(161.6,-173.6)
				..controls(161.8,-174.4)and(161.8,-174.5)..(161.5,-174.9)
				..controls(161.2,-175.6)and(160.8,-175.9)..(159.9,-176)
				..controls(159.2,-176.1)and(154.1,-176.2)..(151.5,-176.2)
				..controls(148.1,-176.1)and(122.6,-176.1)..(121.6,-176.2)
				..controls(120.8,-176.2)and(119.8,-176.2)..(119.4,-176.1)--cycle
				(130.4,-160)..controls(129,-160)and(119.8,-160)..(110.1,-159.9)
				..controls(100.3,-159.8)and(88.5,-159.7)..(83.9,-159.6)
				..controls(77.8,-159.5)and(75.3,-159.5)..(74.5,-159.6)
				..controls(73.9,-159.7)and(72.5,-159.8)..(71.3,-159.8)
				..controls(69.6,-159.8)and(69,-159.8)..(68.5,-159.6)
				..controls(67.6,-159.3)and(66.8,-158.6)..(66.2,-157.6)
				..controls(66,-157.3)and(65.6,-156.5)..(65.3,-156.1)
				..controls(64.6,-155.1)and(64.5,-154.9)..(63.7,-152.9)
				..controls(63.2,-151.6)and(62.9,-150.9)..(61.8,-149.6)
				..controls(61.5,-149.2)and(61.1,-148.5)..(60.9,-147.9)
				..controls(60.7,-147.4)and(60.3,-146.6)..(60.1,-146.1)
				..controls(59.9,-145.5)and(59.5,-144.9)..(59,-144.3)
				..controls(58.5,-143.7)and(58.2,-143.2)..(57.9,-142.5)
				..controls(57.7,-141.9)and(57.3,-141.1)..(57,-140.7)
				..controls(56.6,-140.3)and(56.1,-139.3)..(55.8,-138.5)
				..controls(55.5,-137.8)and(55.1,-136.8)..(54.9,-136.5)
				..controls(54.7,-136.1)and(54.3,-135.2)..(53.9,-134.5)
				..controls(53.6,-133.8)and(53,-132.7)..(52.6,-132)
				..controls(51.6,-130.5)and(51.4,-130.2)..(51.1,-129)
				..controls(50.6,-127.5)and(50.6,-127.4)..(49.7,-126.3)
				..controls(49.2,-125.7)and(48.6,-124.8)..(48.3,-124.2)
				..controls(48,-123.6)and(47.5,-122.7)..(47.3,-122.2)
				..controls(47,-121.8)and(46.5,-120.7)..(46.1,-119.8)
				..controls(45.8,-118.9)and(45.2,-117.8)..(44.9,-117.2)
				..controls(44.2,-116.1)and(43.3,-114.4)..(42.3,-112.2)
				..controls(41.8,-111.3)and(41.3,-110.4)..(41.1,-110.1)
				..controls(40.5,-109.4)and(40.3,-109)..(40,-107.8)
				..controls(39.8,-107)and(39.6,-106.5)..(39.2,-106)
				..controls(38.9,-105.6)and(38.4,-104.7)..(38,-104)
				..controls(37.7,-103.4)and(37.2,-102.5)..(37,-102.1)
				..controls(36.7,-101.7)and(36.3,-100.8)..(36,-100.1)
				..controls(35.6,-98.9)and(35.3,-98.4)..(34.6,-97.6)
				..controls(34.2,-97.1)and(34.2,-96.3)..(34.6,-96)
				..controls(35,-95.7)and(35.5,-95.9)..(36.2,-96.7)
				..controls(36.5,-97)and(37.2,-97.5)..(37.6,-97.7)..controls(38,-98)and(38.6,-98.4)..(39.1,-98.9)
				..controls(39.6,-99.4)and(40.3,-100)..(40.6,-100.3)
				..controls(41.3,-100.8)and(42.4,-101.8)..(43.7,-102.9)
				..controls(44,-103.2)and(44.6,-103.7)..(45,-103.9)
				..controls(45.4,-104.2)and(45.9,-104.6)..(46.1,-104.8)
				..controls(47.1,-105.9)and(48.8,-107.1)..(50.5,-107.9)
				..controls(51.8,-108.6)and(52.3,-109)..(53.4,-110)
				..controls(54.8,-111.2)and(55.6,-111.6)..(57,-112.1)
				..controls(57.5,-112.3)and(58.5,-112.8)..(59.4,-113.3)
				..controls(60.2,-113.7)and(61.3,-114.3)..(61.9,-114.5)
				..controls(62.4,-114.7)and(63.3,-115.1)..(63.8,-115.3)
				..controls(64.4,-115.6)and(65.4,-115.9)..(66.2,-116.1)
				..controls(67,-116.3)and(68.4,-116.6)..(69.2,-116.9)
				..controls(72.3,-117.9)and(72.8,-118)..(78,-118.1)
				..controls(84,-118.3)and(86.1,-118.1)..(88.1,-117.1)
				..controls(89.4,-116.5)and(93,-115)..(93.7,-114.9)
				..controls(94.8,-114.6)and(95.4,-114.2)..(95.9,-113.4)
				..controls(96.2,-113)and(96.7,-112.5)..(97.1,-112.2)
				..controls(98,-111.6)and(101.3,-108.4)..(102.1,-107.2)
				..controls(102.4,-106.8)and(102.9,-106.2)..(103.2,-105.8)
				..controls(103.5,-105.5)and(104,-104.8)..(104.2,-104.2)
				..controls(104.5,-103.6)and(105,-102.8)..(105.3,-102.4)
				..controls(105.6,-102)and(106.1,-101.2)..(106.5,-100.6)
				..controls(107.1,-99.5)and(108.6,-97.5)..(109.5,-96.5)
				..controls(109.9,-96.1)and(110.1,-96)..(110.4,-96)
				..controls(111.1,-96)and(111.6,-96.5)..(111.8,-97.6)
				..controls(112,-98.6)and(112.3,-99.3)..(112.8,-99.9)
				..controls(113.7,-100.9)and(115.2,-103.1)..(115.7,-104.2)
				..controls(116.1,-105)and(116.5,-105.5)..(116.9,-106)
				..controls(117.3,-106.4)and(117.8,-106.9)..(118,-107.3)
				..controls(118.2,-107.6)and(118.6,-108)..(119,-108.3)
				..controls(119.6,-108.8)and(121.1,-110.3)..(121.8,-111.2)
				..controls(122.1,-111.7)and(122.5,-112)..(123.4,-112.4)
				..controls(124.8,-113.2)and(125.8,-113.8)..(127,-114.6)
				..controls(127.5,-114.9)and(128.1,-115.3)..(128.2,-115.4)
				..controls(128.4,-115.5)and(129.3,-115.8)..(130.2,-116)
				..controls(131.1,-116.2)and(132.5,-116.6)..(133.3,-116.9)
				..controls(135.7,-117.9)and(136.1,-117.9)..(142.5,-117.8)
				..controls(145.8,-117.8)and(148.4,-117.7)..(148.8,-117.6)
				..controls(149.2,-117.6)and(150.2,-117.2)..(151,-116.8)
				..controls(152,-116.4)and(152.8,-116)..(153.6,-115.9)
				..controls(154.5,-115.7)and(155.2,-115.4)..(156.3,-114.9)
				..controls(157.1,-114.5)and(158.2,-114)..(158.7,-113.8)
				..controls(159.1,-113.7)and(159.8,-113.3)..(160.2,-113)
				..controls(160.6,-112.7)and(161.3,-112.4)..(161.7,-112.2)
				..controls(162.7,-111.9)and(163.2,-111.7)..(164.1,-111.1)
				..controls(164.5,-110.8)and(165.3,-110.4)..(165.9,-110)
				..controls(166.5,-109.7)and(167.3,-109.2)..(167.6,-108.9)
				..controls(167.9,-108.6)and(168.6,-108.2)..(169,-108)
				..controls(169.4,-107.8)and(170.2,-107.3)..(170.7,-106.9)
				..controls(171.2,-106.5)and(172.1,-105.8)..(172.7,-105.4)
				..controls(173.2,-105)and(174.1,-104.4)..(174.5,-104)
				..controls(174.9,-103.7)and(175.7,-103.1)..(176.2,-102.6)
				..controls(176.7,-102.2)and(177.4,-101.5)..(177.9,-101.2)
				..controls(178.3,-100.8)and(178.8,-100.4)..(179,-100.2)
				..controls(179.2,-100)and(179.6,-99.6)..(180,-99.3)
				..controls(180.4,-99.1)and(180.9,-98.6)..(181.2,-98.4)
				..controls(181.5,-98.1)and(182.1,-97.7)..(182.5,-97.5)
				..controls(183.2,-97.2)and(183.5,-96.9)..(184.3,-95.9)
				..controls(184.6,-95.6)and(184.6,-95.6)..(184.8,-95.7)
				..controls(185,-95.8)and(185.1,-96)..(185.1,-96.2)
				..controls(185.1,-96.5)and(185,-96.7)..(184.6,-97.2)
				..controls(184.1,-97.8)and(184,-97.9)..(183.9,-98.8)
				..controls(183.7,-100.3)and(183.3,-101.3)..(182.9,-101.8)
				..controls(182.6,-102.1)and(182.3,-102.6)..(182.2,-102.9)
				..controls(182,-103.3)and(181.6,-104)..(181.2,-104.6)
				..controls(180.9,-105.2)and(180.5,-106.1)..(180.4,-106.6)
				..controls(179.9,-108)and(179.7,-108.5)..(179,-109.4)
				..controls(178.6,-109.8)and(178.2,-110.5)..(177.9,-111)
				..controls(177.7,-111.5)and(177.3,-112.3)..(177,-112.7)
				..controls(176.6,-113.2)and(176.4,-113.7)..(176.3,-114.2)
				..controls(176.2,-114.7)and(175.9,-115.2)..(175.4,-116)
				..controls(175,-116.6)and(174.4,-117.7)..(174,-118.5)
				..controls(173.7,-119.2)and(173.3,-120.1)..(173,-120.5)
				..controls(172.8,-120.9)and(172.4,-121.8)..(172.1,-122.5)
				..controls(171.8,-123.2)and(171.3,-124.1)..(171,-124.6)
				..controls(170.8,-125.1)and(170.3,-126)..(169.9,-126.6)
				..controls(169.6,-127.3)and(169.2,-128.1)..(169,-128.3)
				..controls(168.8,-128.6)and(168.5,-129.3)..(168.3,-129.8)
				..controls(168.1,-130.3)and(167.8,-131.1)..(167.5,-131.6)
				..controls(166.9,-132.5)and(164.2,-137.9)..(164.2,-138.1)
				..controls(164.2,-138.1)and(163.9,-138.8)..(163.4,-139.7)
				..controls(162.9,-140.5)and(162.4,-141.5)..(162.3,-141.9)
				..controls(161.9,-143)and(161.6,-143.4)..(160.9,-144.1)
				..controls(160.4,-144.6)and(160.3,-144.8)..(160.1,-145.4)
				..controls(160,-145.8)and(159.6,-146.8)..(159.2,-147.8)
				..controls(158.7,-148.8)and(158.1,-150)..(157.9,-150.6)
				..controls(157.6,-151.1)and(157.2,-151.9)..(156.9,-152.4)
				..controls(156.5,-152.8)and(156.1,-153.7)..(155.8,-154.4)
				..controls(155.6,-155)and(155.2,-155.9)..(154.9,-156.3)
				..controls(154.7,-156.7)and(154.4,-157.4)..(154.2,-157.9)
				..controls(154.1,-158.4)and(153.8,-159)..(153.5,-159.3)
				..controls(152.8,-160.1)and(153,-160.1)..(142.4,-160.1)
				..controls(137.2,-160.1)and(131.8,-160.1)..(130.4,-160)--cycle
				(141.5,-103.7)..controls(140.8,-103)and(139.8,-102.1)..(139.2,-101.6)
				..controls(138.5,-101.2)and(137.9,-100.6)..(137.8,-100.4)
				..controls(137,-99.5)and(134,-96.5)..(133.2,-96)
				..controls(132.7,-95.7)and(131.8,-94.9)..(131.1,-94.3)
				..controls(129.4,-92.7)and(127.9,-91.6)..(126.6,-91.1)
				..controls(125.2,-90.5)and(124.1,-90)..(122.9,-89.1)
				..controls(122.1,-88.6)and(121.3,-88.2)..(119.9,-87.7)
				..controls(118.8,-87.3)and(117.5,-86.8)..(116.9,-86.5)
				..controls(114.3,-85.5)and(113.7,-85.4)..(109.7,-85.5)
				..controls(105.3,-85.7)and(103.9,-85.9)..(101.6,-87)
				..controls(100.7,-87.4)and(99.4,-87.8)..(97.9,-88)
				..controls(97,-88.2)and(96.6,-88.3)..(95.8,-88.9)
				..controls(95.4,-89.2)and(94.7,-89.6)..(94.2,-89.8)
				..controls(92.9,-90.3)and(92.3,-90.7)..(90.9,-92.1)
				..controls(90.2,-92.7)and(89.4,-93.3)..(89,-93.5)..controls(87,-94.5)and(85.2,-96)..(83.6,-97.9)
				..controls(83.3,-98.2)and(82.5,-98.9)..(81.8,-99.4)
				..controls(80.7,-100.1)and(80.5,-100.3)..(80,-101.1)
				..controls(78.7,-103.1)and(77.2,-103.9)..(76.4,-103)
				..controls(75.7,-102.2)and(75.7,-103.2)..(75.7,-83.6)
				..controls(75.6,-73.6)and(75.6,-64.8)..(75.5,-64)
				..controls(75.5,-63.2)and(75.5,-61.4)..(75.6,-59.8)--(75.7,-57.1)--(76,-56.8)
				..controls(76.8,-56.2)and(78,-56.3)..(79.8,-57.2)
				..controls(80.4,-57.5)and(81.3,-57.8)..(81.8,-57.9)
				..controls(82.3,-58.1)and(83,-58.4)..(83.6,-58.8)..controls(84.3,-59.2)and(85,-59.5)..(86,-59.8)
				..controls(86.8,-60.1)and(88,-60.6)..(88.7,-60.9)
				..controls(89.3,-61.2)and(90.6,-61.7)..(91.4,-62)
				..controls(92.2,-62.3)and(93.7,-62.9)..(94.6,-63.3)
				..controls(95.6,-63.7)and(96.8,-64.1)..(97.4,-64.3)
				..controls(98,-64.5)and(99,-64.9)..(99.7,-65.2)
				..controls(100.4,-65.6)and(101.5,-66)..(102.3,-66.2)
				..controls(104.2,-66.8)and(105,-67.1)..(106.2,-67.8)
				..controls(107.9,-68.7)and(110.2,-69.4)..(111.3,-69.2)
				..controls(112.5,-69)and(115,-68.2)..(115.6,-67.8)
				..controls(117.3,-66.7)and(117.6,-66.6)..(119.5,-66.2)
				..controls(121.9,-65.8)and(122.1,-65.7)..(123,-65.1)
				..controls(124.1,-64.4)and(124.4,-64.3)..(126,-63.9)
				..controls(127,-63.7)and(127.7,-63.4)..(128.4,-63.1)
				..controls(129,-62.8)and(129.7,-62.5)..(130.6,-62.4)
				..controls(132.1,-62.1)and(134,-61.5)..(135.2,-61)
				..controls(135.7,-60.8)and(136.7,-60.5)..(137.3,-60.3)
				..controls(138.4,-60)and(139.3,-59.7)..(141.1,-59)
				..controls(141.6,-58.7)and(142,-58.6)..(142.5,-58.6)
				..controls(143.2,-58.6)and(143.3,-58.7)..(143.6,-59)
				..controls(144,-59.5)and(144.1,-59.9)..(144,-60.8)
				..controls(144,-61.2)and(144,-71.1)..(144,-82.8)--(143.9,-104.1)--(143.7,-104.5)
				..controls(143.5,-104.7)and(143.3,-104.8)..(143,-104.9)
				..controls(142.6,-104.9)and(142.5,-104.8)..(141.5,-103.7)--cycle
				(27.8,-91.4)..controls(27.1,-90.9)and(26.3,-90)..(26.1,-89.5)
				..controls(26,-89.2)and(25.8,-88.6)..(25.7,-87.9)
				..controls(25.4,-86.2)and(25.6,-84.8)..(26.5,-83.5)
				..controls(27.4,-82.3)and(28.3,-82)..(29.4,-82.3)
				..controls(31.1,-82.8)and(32.1,-84.9)..(32,-87.4)
				..controls(31.9,-89.3)and(31.2,-90.6)..(29.9,-91.4)
				..controls(29.1,-91.9)and(28.6,-91.9)..(27.8,-91.4)--cycle
				(190.1,-91.3)..controls(189.7,-91.2)and(189.3,-90.9)..(189.1,-90.7)
				..controls(188.3,-89.8)and(187.6,-87.7)..(187.6,-86.2)
				..controls(187.6,-85.5)and(187.7,-85.2)..(188.2,-84.4)
				..controls(189.1,-82.6)and(190.1,-81.8)..(191.3,-82.1)
				..controls(192,-82.3)and(193.5,-83.7)..(193.8,-84.4)
				..controls(194.1,-85.1)and(194.2,-87.1)..(194,-88.3)
				..controls(193.8,-89.8)and(193.3,-90.5)..(192.3,-91.3)
				..controls(191.7,-91.7)and(190.9,-91.7)..(190.1,-91.3)--cycle
				(173.7,-83.4)..controls(172.2,-82.5)and(172,-82.5)..(170.6,-82.4)
				..controls(169.7,-82.4)and(168.8,-82.4)..(168.3,-82.5)
				..controls(167.2,-82.7)and(166.9,-82.7)..(166.6,-82.2)
				..controls(166.3,-81.8)and(166.3,-81)..(166.8,-80)
				..controls(167.2,-78.9)and(167.4,-77.4)..(167.7,-74.1)
				..controls(167.8,-72.8)and(167.9,-71.8)..(167.9,-71.7)
				..controls(168,-71.7)and(168,-70.2)..(168,-68.4)..controls(168,-65.4)and(168,-64.9)..(167.7,-63)
				..controls(167.3,-60.5)and(167.1,-59.9)..(166.4,-58.4)
				..controls(166.1,-57.7)and(165.8,-56.9)..(165.7,-56.6)
				..controls(165.7,-56.2)and(165.4,-55.6)..(165.2,-55.2)
				..controls(164.6,-53.8)and(161.9,-50.7)..(161,-50.3)
				..controls(159.8,-49.7)and(157.9,-50)..(155.2,-51.2)
				..controls(154.7,-51.4)and(153.4,-51.8)..(152.2,-52.1)
				..controls(150.8,-52.4)and(149.4,-52.9)..(148.6,-53.2)
				..controls(147.8,-53.6)and(146.7,-54)..(146,-54.2)
				..controls(145.4,-54.4)and(144.5,-54.7)..(144,-54.9)
				..controls(143.6,-55.2)and(142.7,-55.5)..(141.8,-55.7)
				..controls(141,-55.9)and(139.6,-56.3)..(138.7,-56.6)
				..controls(137.8,-57)and(136.3,-57.4)..(135.3,-57.6)
				..controls(133.8,-57.9)and(133.2,-58.1)..(132.1,-58.7)
				..controls(131.1,-59.2)and(130.3,-59.4)..(129.4,-59.6)
				..controls(128.7,-59.8)and(127.8,-60)..(127.4,-60.1)
				..controls(127,-60.3)and(126,-60.5)..(125.1,-60.7)
				..controls(123.2,-61.1)and(121.7,-61.5)..(120.2,-62.1)
				..controls(119.7,-62.3)and(118.1,-62.8)..(116.7,-63.2)
				..controls(114.6,-63.9)and(114,-64.1)..(113,-64.7)
				..controls(111.9,-65.4)and(111.8,-65.4)..(111,-65.5)
				..controls(109.9,-65.5)and(109.3,-65.4)..(107.7,-64.7)
				..controls(107,-64.4)and(106,-64)..(105.3,-63.9)
				..controls(103.5,-63.5)and(102.2,-63)..(101.2,-62.5)
				..controls(100.1,-61.8)and(99.3,-61.5)..(97.8,-61.1)
				..controls(97.2,-61)and(96.2,-60.7)..(95.7,-60.4)
				..controls(95.1,-60.1)and(94.2,-59.8)..(93.5,-59.6)
				..controls(92.9,-59.4)and(91.7,-59)..(91,-58.7)..controls(90.2,-58.4)and(88.9,-58)..(88,-57.7)
				..controls(87.2,-57.5)and(86.1,-57.1)..(85.5,-56.8)
				..controls(85,-56.6)and(84.1,-56.2)..(83.5,-56)..controls(82.9,-55.7)and(81.8,-55.2)..(81,-54.9)
				..controls(80.1,-54.5)and(79.1,-54.2)..(78.7,-54.1)
				..controls(77.9,-53.8)and(77.3,-53.6)..(76.2,-52.9)
				..controls(75.7,-52.6)and(75.2,-52.4)..(74.1,-52.2)
				..controls(72.5,-51.9)and(72.4,-51.8)..(71.3,-51.1)
				..controls(69.9,-50.1)and(69.4,-49.9)..(67.9,-49.8)
				..controls(67,-49.7)and(66.6,-49.6)..(66.3,-49.4)
				..controls(66.1,-49.3)and(65.7,-49)..(65.4,-48.8)
				..controls(65.1,-48.7)and(64.2,-48.4)..(63.6,-48.2)
				..controls(62.9,-48)and(61.7,-47.6)..(60.9,-47.2)
				..controls(60.1,-46.8)and(58.9,-46.4)..(58.2,-46.2)
				..controls(57.6,-46)and(56.8,-45.7)..(56.5,-45.5)
				..controls(54.9,-44.6)and(54.6,-44.5)..(53.2,-44.2)
				..controls(51.7,-43.8)and(51,-43.6)..(49.9,-42.9)
				..controls(49.5,-42.6)and(48.9,-42.4)..(47.5,-42)
				..controls(45.7,-41.6)and(45.6,-41.5)..(44.3,-40.6)
				..controls(43.5,-40.1)and(42.8,-39.6)..(42.7,-39.4)
				..controls(42.4,-38.7)and(43,-38.1)..(44.5,-37.6)..controls(44.9,-37.5)and(46,-37)..(46.9,-36.6)
				..controls(47.9,-36.1)and(48.8,-35.8)..(49.6,-35.6)
				..controls(50.5,-35.4)and(51.3,-35.2)..(52.4,-34.6)
				..controls(53.3,-34.1)and(54.5,-33.7)..(54.9,-33.5)
				..controls(55.4,-33.4)and(56.4,-33)..(57.2,-32.6)
				..controls(58.6,-32)and(60.1,-31.5)..(61.6,-31.2)
				..controls(62.3,-31.1)and(62.7,-31)..(63.2,-30.7)
				..controls(64.1,-30)and(64.3,-29.9)..(65.3,-29.7)
				..controls(65.7,-29.6)and(66.5,-29.2)..(67.1,-28.9)
				..controls(67.6,-28.6)and(68.8,-28.2)..(69.6,-27.9)
				..controls(70.5,-27.7)and(71.6,-27.3)..(72.1,-27)
				..controls(72.6,-26.7)and(73.8,-26.2)..(74.8,-25.9)
				..controls(75.7,-25.5)and(76.7,-25.1)..(77.1,-25)
				..controls(77.4,-24.8)and(78.5,-24.4)..(79.6,-24.1)
				..controls(80.6,-23.8)and(82,-23.2)..(82.8,-22.9)
				..controls(83.5,-22.5)and(84.8,-22)..(85.8,-21.8)
				..controls(86.7,-21.5)and(87.8,-21.1)..(88.1,-20.9)
				..controls(88.4,-20.7)and(89.3,-20.3)..(90.1,-20.1)
				..controls(90.8,-19.8)and(91.9,-19.4)..(92.4,-19.1)
				..controls(93.5,-18.5)and(94.4,-18.2)..(95.4,-18.1)
				..controls(95.7,-18)and(96.4,-17.8)..(96.9,-17.5)
				..controls(97.5,-17.3)and(98.4,-16.9)..(99,-16.6)
				..controls(99.7,-16.3)and(100.9,-15.9)..(101.9,-15.6)
				..controls(102.9,-15.3)and(104.4,-14.7)..(105.3,-14.3)
				..controls(108.1,-13.2)and(108.5,-13.1)..(110.1,-13.1)
				..controls(111.4,-13.1)and(111.5,-13.1)..(112.4,-13.5)
				..controls(113.2,-14)and(115,-14.8)..(116.6,-15.4)
				..controls(116.9,-15.5)and(117.9,-15.8)..(118.8,-16.1)
				..controls(120.9,-16.6)and(121.5,-16.9)..(122.6,-17.9)
				..controls(123,-18.3)and(123.7,-18.8)..(123.9,-18.9)
				..controls(124.6,-19.4)and(126,-19.8)..(127.9,-20.2)
				..controls(129.6,-20.5)and(129.8,-20.6)..(131.1,-21.4)
				..controls(131.5,-21.7)and(132.4,-22.1)..(133.1,-22.3)
				..controls(136.1,-23.2)and(138.3,-24.1)..(140.2,-25)
				..controls(141,-25.5)and(141.5,-25.6)..(142.5,-25.8)
				..controls(145.1,-26.3)and(146,-26.6)..(149.4,-28.1)
				..controls(151.3,-28.9)and(153.5,-29.8)..(154.3,-30.1)
				..controls(155.1,-30.4)and(156.2,-30.8)..(156.6,-31)
				..controls(157.1,-31.2)and(158.2,-31.6)..(159.1,-31.8)
				..controls(160.6,-32.2)and(160.9,-32.3)..(161.7,-32.9)
				..controls(162.4,-33.3)and(162.9,-33.6)..(163.7,-33.8)
				..controls(164.3,-34)and(164.8,-34.1)..(164.9,-34.1)
				..controls(165,-34.1)and(165.6,-34.4)..(166.4,-34.8)
				..controls(167.3,-35.2)and(168.1,-35.4)..(169,-35.6)
				..controls(170,-35.8)and(170.5,-36)..(171.7,-36.6)
				..controls(172.4,-37)and(173.3,-37.4)..(173.6,-37.4)
				..controls(175.6,-37.7)and(176.2,-38)..(176.7,-39.1)
				..controls(176.9,-39.6)and(176.9,-39.7)..(176.7,-40.4)
				..controls(176.4,-41.6)and(175.8,-42.7)..(175.2,-43.4)
				..controls(174.5,-44.1)and(172.6,-45.5)..(171.7,-45.8)
				..controls(170.6,-46.2)and(169.1,-46)..(168.2,-45.4)
				..controls(167.9,-45.2)and(167.3,-44.8)..(166.9,-44.4)
				..controls(166.4,-44.1)and(165.8,-43.6)..(165.4,-43.2)
				..controls(165.1,-42.9)and(164.4,-42.3)..(163.9,-41.9)
				..controls(163.3,-41.5)and(162.7,-40.9)..(162.4,-40.5)
				..controls(161.4,-39.2)and(159.9,-38)..(158.8,-37.8)
				..controls(158.4,-37.8)and(158.3,-37.7)..(158.2,-37.2)
				..controls(158,-36.8)and(157.4,-36.1)..(157.1,-36.1)
				..controls(157,-36.1)and(156.7,-36.1)..(156.5,-36)
				..controls(156.1,-35.9)and(155.9,-35.7)..(155.7,-35.2)
				..controls(155.3,-34.6)and(155.1,-34.4)..(154,-34.1)
				..controls(153.6,-34)and(153.2,-33.7)..(152.9,-33.5)
				..controls(152.4,-32.9)and(151.6,-32.4)..(151.3,-32.4)
				..controls(150.9,-32.4)and(150.4,-32.6)..(150.4,-32.9)
				..controls(150.4,-33.2)and(150.8,-33.4)..(151.4,-33.6)
				..controls(152.1,-33.7)and(152.8,-34.1)..(153.5,-34.9)
				..controls(154.1,-35.5)and(154.3,-35.6)..(155,-35.8)
				..controls(155.7,-36)and(155.7,-36)..(156,-36.6)
				..controls(156.3,-37.3)and(156.7,-37.7)..(157.5,-37.8)
				..controls(158.1,-37.9)and(158.2,-38)..(158.4,-38.8)
				..controls(158.5,-39.4)and(158.6,-39.6)..(159.6,-40.4)
				..controls(160.2,-41)and(160.9,-41.6)..(161.3,-41.8)
				..controls(161.7,-42)and(161.9,-42.3)..(162.2,-42.8)
				..controls(162.4,-43.2)and(162.8,-43.8)..(163.1,-44.1)
				..controls(163.9,-44.8)and(165.1,-46.7)..(165.7,-48)
				..controls(166.2,-49.3)and(166.4,-49.7)..(167.2,-50.5)
				..controls(167.5,-50.8)and(167.9,-51.4)..(168.1,-51.8)
				..controls(168.3,-52.2)and(168.8,-53)..(169.1,-53.6)
				..controls(169.6,-54.3)and(169.8,-54.9)..(169.9,-55.6)
				..controls(170.4,-57.8)and(170.5,-58.1)..(171.2,-59.5)
				..controls(171.6,-60.3)and(172,-61.5)..(172.2,-62.1)
				..controls(172.3,-62.8)and(172.7,-63.7)..(172.9,-64.1)
				..controls(173.2,-64.6)and(173.4,-65.3)..(173.4,-65.6)
				..controls(173.8,-68.8)and(174.2,-70.6)..(174.9,-72.6)
				..controls(175.5,-74.5)and(175.9,-77.1)..(176.1,-80.2)
				..controls(176.3,-83.4)and(175.5,-84.4)..(173.7,-83.4)--cycle;
				\definecolor{ezs}{RGB}{247,179,26}\fill[ezs](-5.7,-156.6)--(-3.4,-157.5)--(61.9,12.7)--(59.7,13.6)--cycle;
				\definecolor{ezs}{RGB}{247,179,22}\fill[ezs](-3.8,-157.3)--(-1.5,-158.2)--(63.9,12)--(61.6,12.9)--cycle;
				\definecolor{ezs}{RGB}{248,179,19}\fill[ezs](-1.8,-158.1)--(0.4,-158.9)--(65.8,11.2)--(63.5,12.1)--cycle;
				\definecolor{ezs}{RGB}{248,178,16}\fill[ezs](0.1,-158.8)--(2.4,-159.7)--(67.7,10.5)--(65.4,11.4)--cycle;
				\definecolor{ezs}{RGB}{249,178,15}\fill[ezs](2,-159.5)--(4.3,-160.4)--(69.6,9.8)--(67.3,10.7)--cycle;
				\definecolor{ezs}{RGB}{251,177,17}\fill[ezs](3.9,-160.3)--(6.2,-161.1)--(71.5,9)--(69.2,9.9)--cycle;
				\definecolor{ezs}{RGB}{252,179,18}\fill[ezs](5.8,-161)--(8.1,-161.9)--(73.4,8.3)--(71.1,9.2)--cycle;
				\definecolor{ezs}{RGB}{252,182,17}\fill[ezs](7.7,-161.7)--(10,-162.6)--(75.3,7.6)--(73,8.5)--cycle;
				\definecolor{ezs}{RGB}{252,185,17}\fill[ezs](9.6,-162.5)--(11.9,-163.4)--(77.2,6.8)--(75,7.7)--cycle;
				\definecolor{ezs}{RGB}{251,189,18}\fill[ezs](11.5,-163.2)--(13.8,-164.1)--(79.2,6.1)--(76.9,7)--cycle;
				\definecolor{ezs}{RGB}{251,191,22}\fill[ezs](13.4,-163.9)--(15.7,-164.8)--(81.1,5.4)--(78.8,6.2)--cycle;
				\definecolor{ezs}{RGB}{252,193,25}\fill[ezs](15.4,-164.7)--(17.6,-165.6)--(83,4.6)--(80.7,5.5)--cycle;
				\definecolor{ezs}{RGB}{253,195,27}\fill[ezs](17.3,-165.4)--(19.6,-166.3)--(84.9,3.9)--(82.6,4.8)--cycle;
				\definecolor{ezs}{RGB}{253,196,31}\fill[ezs](19.2,-166.1)--(21.5,-167)--(86.8,3.2)--(84.5,4)--cycle;
				\definecolor{ezs}{RGB}{254,199,35}\fill[ezs](21.1,-166.9)--(23.4,-167.8)--(88.7,2.4)--(86.4,3.3)--cycle;
				\definecolor{ezs}{RGB}{253,201,38}\fill[ezs](23,-167.6)--(25.3,-168.5)--(90.6,1.7)--(88.3,2.6)--cycle;
				\definecolor{ezs}{RGB}{253,203,40}\fill[ezs](24.9,-168.3)--(27.2,-169.2)--(92.5,1)--(90.3,1.8)--cycle;
				\definecolor{ezs}{RGB}{253,205,42}\fill[ezs](26.8,-169.1)--(29.1,-170)--(94.4,0.2)--(92.2,1.1)--cycle;
				\definecolor{ezs}{RGB}{253,207,45}\fill[ezs](28.7,-169.8)--(31,-170.7)--(96.4,-0.5)--(94.1,0.4)--cycle;
				\definecolor{ezs}{RGB}{253,210,49}\fill[ezs](30.7,-170.5)--(32.9,-171.4)--(98.3,-1.2)--(96,-0.4)--cycle;
				\definecolor{ezs}{RGB}{253,212,54}\fill[ezs](32.6,-171.3)--(34.9,-172.2)--(100.2,-2)--(97.9,-1.1)--cycle;
				\definecolor{ezs}{RGB}{253,214,59}\fill[ezs](34.5,-172)--(36.8,-172.9)--(102.1,-2.7)--(99.8,-1.8)--cycle;
				\definecolor{ezs}{RGB}{253,216,63}\fill[ezs](36.4,-172.7)--(38.7,-173.6)--(104,-3.4)--(101.7,-2.6)--cycle;
				\definecolor{ezs}{RGB}{253,218,65}\fill[ezs](38.3,-173.5)--(40.6,-174.4)--(105.9,-4.2)--(103.6,-3.3)--cycle;
				\definecolor{ezs}{RGB}{254,220,66}\fill[ezs](40.2,-174.2)--(42.5,-175.1)--(107.8,-4.9)--(105.5,-4)--cycle;
				\definecolor{ezs}{RGB}{254,222,68}\fill[ezs](42.1,-175)--(44.4,-175.8)--(109.7,-5.6)--(107.5,-4.8)--cycle;
				\definecolor{ezs}{RGB}{254,224,71}\fill[ezs](44,-175.7)--(46.3,-176.6)--(111.7,-6.4)--(109.4,-5.5)--cycle;
				\definecolor{ezs}{RGB}{254,226,75}\fill[ezs](46,-176.4)--(48.2,-177.3)--(113.6,-7.1)--(111.3,-6.2)--cycle;
				\definecolor{ezs}{RGB}{254,228,79}\fill[ezs](47.9,-177.2)--(50.1,-178)--(115.5,-7.8)--(113.2,-7)--cycle;
				\definecolor{ezs}{RGB}{254,230,82}\fill[ezs](49.8,-177.9)--(52.1,-178.8)--(117.4,-8.6)--(115.1,-7.7)--cycle;
				\definecolor{ezs}{RGB}{254,232,85}\fill[ezs](51.7,-178.6)--(54,-179.5)--(119.3,-9.3)--(117,-8.4)--cycle;
				\definecolor{ezs}{RGB}{254,234,88}\fill[ezs](53.6,-179.4)--(55.9,-180.2)--(121.2,-10)--(118.9,-9.2)--cycle;
				\definecolor{ezs}{RGB}{254,237,91}\fill[ezs](55.5,-180.1)--(57.8,-181)--(123.1,-10.8)--(120.8,-9.9)--cycle;
				\definecolor{ezs}{RGB}{254,239,94}\fill[ezs](57.4,-180.8)--(59.7,-181.7)--(125,-11.5)--(122.8,-10.6)--cycle;
				\definecolor{ezs}{RGB}{254,241,97}\fill[ezs](59.3,-181.6)--(61.6,-182.4)--(127,-12.2)--(124.7,-11.4)--cycle;
				\definecolor{ezs}{RGB}{254,243,101}\fill[ezs](61.2,-182.3)--(63.5,-183.2)--(128.9,-13)--(126.6,-12.1)--cycle;
				\definecolor{ezs}{RGB}{254,245,103}\fill[ezs](63.2,-183)--(65.4,-183.9)--(130.8,-13.7)--(128.5,-12.8)--cycle;
				\definecolor{ezs}{RGB}{254,245,105}\fill[ezs](65.1,-183.8)--(67.4,-184.6)--(132.7,-14.4)--(130.4,-13.6)--cycle;
				\definecolor{ezs}{RGB}{254,245,105}\fill[ezs](67,-184.5)--(69.3,-185.4)--(134.6,-15.2)--(132.3,-14.3)--cycle;
				\definecolor{ezs}{RGB}{254,244,103}\fill[ezs](68.9,-185.2)--(71.2,-186.1)--(136.5,-15.9)--(134.2,-15)--cycle;
				\definecolor{ezs}{RGB}{253,243,101}\fill[ezs](70.8,-186)--(73.1,-186.8)--(138.4,-16.6)--(136.1,-15.8)--cycle;
				\definecolor{ezs}{RGB}{253,241,98}\fill[ezs](72.7,-186.7)--(75,-187.6)--(140.3,-17.4)--(138,-16.5)--cycle;
				\definecolor{ezs}{RGB}{253,239,95}\fill[ezs](74.6,-187.4)--(76.9,-188.3)--(142.2,-18.1)--(140,-17.2)--cycle;
				\definecolor{ezs}{RGB}{253,236,92}\fill[ezs](76.5,-188.2)--(78.8,-189)--(144.2,-18.8)--(141.9,-18)--cycle;
				\definecolor{ezs}{RGB}{253,234,88}\fill[ezs](78.5,-188.9)--(80.7,-189.8)--(146.1,-19.6)--(143.8,-18.7)--cycle;
				\definecolor{ezs}{RGB}{253,232,85}\fill[ezs](80.4,-189.6)--(82.7,-190.5)--(148,-20.3)--(145.7,-19.4)--cycle;
				\definecolor{ezs}{RGB}{253,230,82}\fill[ezs](82.3,-190.4)--(84.6,-191.2)--(149.9,-21)--(147.6,-20.2)--cycle;
				\definecolor{ezs}{RGB}{253,228,78}\fill[ezs](84.2,-191.1)--(86.5,-192)--(151.8,-21.8)--(149.5,-20.9)--cycle;
				\definecolor{ezs}{RGB}{253,226,74}\fill[ezs](86.1,-191.8)--(88.4,-192.7)--(153.7,-22.5)--(151.4,-21.6)--cycle;
				\definecolor{ezs}{RGB}{253,224,70}\fill[ezs](88,-192.6)--(90.3,-193.4)--(155.6,-23.3)--(153.3,-22.4)--cycle;
				\definecolor{ezs}{RGB}{253,222,67}\fill[ezs](89.9,-193.3)--(92.2,-194.2)--(157.5,-24)--(155.3,-23.1)--cycle;
				\definecolor{ezs}{RGB}{253,220,65}\fill[ezs](91.8,-194)--(94.1,-194.9)--(159.5,-24.7)--(157.2,-23.8)--cycle;
				\definecolor{ezs}{RGB}{253,217,63}\fill[ezs](93.7,-194.8)--(96,-195.6)--(161.4,-25.5)--(159.1,-24.6)--cycle;
				\definecolor{ezs}{RGB}{253,215,60}\fill[ezs](95.7,-195.5)--(97.9,-196.4)--(163.3,-26.2)--(161,-25.3)--cycle;
				\definecolor{ezs}{RGB}{253,213,56}\fill[ezs](97.6,-196.2)--(99.9,-197.1)--(165.2,-26.9)--(162.9,-26)--cycle;
				\definecolor{ezs}{RGB}{252,211,53}\fill[ezs](99.5,-197)--(101.8,-197.8)--(167.1,-27.7)--(164.8,-26.8)--cycle;
				\definecolor{ezs}{RGB}{252,209,49}\fill[ezs](101.4,-197.7)--(103.7,-198.6)--(169,-28.4)--(166.7,-27.5)--cycle;
				\definecolor{ezs}{RGB}{252,207,45}\fill[ezs](103.3,-198.4)--(105.6,-199.3)--(170.9,-29.1)--(168.6,-28.2)--cycle;
				\definecolor{ezs}{RGB}{253,205,42}\fill[ezs](105.2,-199.2)--(107.5,-200)--(172.8,-29.9)--(170.6,-29)--cycle;
				\definecolor{ezs}{RGB}{253,202,40}\fill[ezs](107.1,-199.9)--(109.4,-200.8)--(174.7,-30.6)--(172.5,-29.7)--cycle;
				\definecolor{ezs}{RGB}{253,200,38}\fill[ezs](109,-200.6)--(111.3,-201.5)--(176.7,-31.3)--(174.4,-30.4)--cycle;
				\definecolor{ezs}{RGB}{253,198,35}\fill[ezs](111,-201.4)--(113.2,-202.2)--(178.6,-32.1)--(176.3,-31.2)--cycle;
				\definecolor{ezs}{RGB}{253,196,31}\fill[ezs](112.9,-202.1)--(115.2,-203)--(180.5,-32.8)--(178.2,-31.9)--cycle;
				\definecolor{ezs}{RGB}{253,194,28}\fill[ezs](114.8,-202.8)--(117.1,-203.7)--(182.4,-33.5)--(180.1,-32.6)--cycle;
				\definecolor{ezs}{RGB}{253,192,26}\fill[ezs](116.7,-203.6)--(119,-204.5)--(184.3,-34.3)--(182,-33.4)--cycle;
				\definecolor{ezs}{RGB}{253,189,24}\fill[ezs](118.6,-204.3)--(120.9,-205.2)--(186.2,-35)--(183.9,-34.1)--cycle;
				\definecolor{ezs}{RGB}{253,186,20}\fill[ezs](120.5,-205)--(122.8,-205.9)--(188.1,-35.7)--(185.8,-34.8)--cycle;
				\definecolor{ezs}{RGB}{253,184,17}\fill[ezs](122.4,-205.8)--(124.7,-206.7)--(190,-36.5)--(187.8,-35.6)--cycle;
				\definecolor{ezs}{RGB}{252,183,15}\fill[ezs](124.3,-206.5)--(126.6,-207.4)--(192,-37.2)--(189.7,-36.3)--cycle;
				\definecolor{ezs}{RGB}{252,181,14}\fill[ezs](126.3,-207.2)--(128.5,-208.1)--(193.9,-37.9)--(191.6,-37.1)--cycle;
				\definecolor{ezs}{RGB}{252,179,13}\fill[ezs](128.2,-208)--(130.4,-208.9)--(195.8,-38.7)--(193.5,-37.8)--cycle;
				\definecolor{ezs}{RGB}{252,178,12}\fill[ezs](130.1,-208.7)--(132.4,-209.6)--(197.7,-39.4)--(195.4,-38.5)--cycle;
				\definecolor{ezs}{RGB}{251,178,10}\fill[ezs](132,-209.4)--(134.3,-210.3)--(199.6,-40.1)--(197.3,-39.3)--cycle;
				\definecolor{ezs}{RGB}{251,178,10}\fill[ezs](133.9,-210.2)--(136.2,-211.1)--(201.5,-40.9)--(199.2,-40)--cycle;
				\definecolor{ezs}{RGB}{250,177,11}\fill[ezs](135.8,-210.9)--(138.1,-211.8)--(203.4,-41.6)--(201.1,-40.7)--cycle;
				\definecolor{ezs}{RGB}{250,178,11}\fill[ezs](137.7,-211.6)--(140,-212.5)--(205.3,-42.3)--(203.1,-41.5)--cycle;
				\definecolor{ezs}{RGB}{248,180,10}\fill[ezs](139.6,-212.4)--(141.9,-213.3)--(207.3,-43.1)--(205,-42.2)--cycle;
				\definecolor{ezs}{RGB}{247,181,9}\fill[ezs](141.5,-213.1)--(143.8,-214)--(209.2,-43.8)--(206.9,-42.9)--cycle;
				\definecolor{ezs}{RGB}{246,182,9}\fill[ezs](143.5,-213.8)--(145.7,-214.7)--(211.1,-44.5)--(208.8,-43.7)--cycle;
				\definecolor{ezs}{RGB}{246,182,10}\fill[ezs](145.4,-214.6)--(147.7,-215.5)--(213,-45.3)--(210.7,-44.4)--cycle;
			\end{scope}
			% chữ EDUZY (Times đậm, làm dày nét bằng 5 lớp lệch nhẹ; căn đúng khung chữ logo gốc)
			\foreach \dx/\dy in {-0.7/0,0.7/0,0/-0.3,0/0.3,0/0}
			\node[anchor=center,inner sep=0pt,outer sep=0pt,text=ezText,align=center] at (109.6+\dx,-193.6+\dy)
			{\resizebox{1.01cm}{0.22cm}{\fontencoding{T1}\fontfamily{ptm}\fontseries{b}\selectfont EDUZY}};
		\end{tikzpicture}%
	}%
}
%====================== THIẾT LẬP THÂN ĐỀ ======================
%--- Môi trường tự luận {bt}: ex_test 3.3.4 chưa khai báo sẵn nên tạo ở đây
%    (ex_test đã tự móc \fixshowans và \listenumerate cho bt)
\makeatletter
\@ifundefined{bt}{\newtheorem{bt}{\color{gold1}\namebt}}{}
\makeatother
\renewcommand{\sectionmark}[1]{\markright{\thesection\ #1}}
%--- Khung tiêu đề (logo + tối đa 4 dòng; dòng 3, 4 để trống {} thì bỏ qua)
%  Đề thi:    \khungtieude{THI THỬ THPT QUỐC GIA NĂM 2026}{Môn Toán -- Đề số 21}
%                         {(Đề thi gồm: ...)}{Thời gian làm bài: 90 phút (không kể thời gian phát đề).}
%  Chuyên đề: \khungtieude{CHUYÊN ĐỀ 3 -- DÃY SỐ VÀ CÔNG THỨC TRUY HỒI}
%                         {Bồi dưỡng học sinh giỏi Toán THPT}{...}{}
\newcommand{\khungtieude}[4]{%
	\begin{tcolorbox}[colframe=gold1,colback=gold5!45,boxrule=0.5mm]
		\begin{minipage}[c]{3cm}
			\centerline{\eduzylogo[2.4cm]}
		\end{minipage}\hspace{0.5cm}
		\begin{minipage}[c]{12.5cm}
			\vskip 0.15 cm
			\centerline{\ind{#1}}
			\vskip 0.15 cm
			\centerline{\textbf{#2}}
			\if\relax\detokenize{#3}\relax\else
				\vskip 0.15 cm
				\centerline{\textit{#3}}
			\fi
			\if\relax\detokenize{#4}\relax\else
				\vskip 0.15 cm
				\centerline{\textit{#4}}
			\fi
			\vskip 0.15 cm
		\end{minipage}
	\end{tcolorbox}
}
%====================== BỔ SUNG CHO BÀI HỌC, CHUYÊN ĐỀ ======================
\usepackage[shortlabels]{enumitem}
\usepackage{tabularx}
%--- Danh sách gọn: \begin{enumerate}[gon] ... \end{enumerate}
\SetEnumitemKey{gon}{leftmargin=*,itemsep=1pt,parsep=0pt,topsep=3pt}
%--- Khung lý thuyết: \begin{khung}{} ... \end{khung}; tiêu đề chỉ dùng khi thật ngắn ("Mẹo", "Lỗi thường gặp")
\newtcolorbox{khung}[1]{breakable,enhanced,colframe=gold1,colback=gold5!45,
	boxrule=0.4mm,fonttitle=\bfseries\sffamily,coltitle=white,colbacktitle=gold1,
	attach boxed title to top left={xshift=6mm,yshift=-3mm},
	boxed title style={arc=1mm,boxrule=0.3mm},title={#1}}
%--- Tiêu đề mục nhỏ trong một dạng: \muc{Dấu hiệu}, \muc{Phương pháp}, \muc{Bài tự luyện}
\makeatletter
\newcommand{\muc}[1]{\par\medskip\noindent\indam{#1}\par\nobreak\smallskip\@afterindentfalse\@afterheading}
\makeatother
%--- Ô công thức cần nhớ (căn giữa, viền vàng, vừa khít nội dung)
%     \begin{center}\congthuc{$Q(n)=(n-1)!$}\end{center}
\newtcbox{\congthuc}{on line,colframe=gold1,colback=gold5!45,boxrule=0.5pt,arc=2pt,
	left=8pt,right=8pt,top=4pt,bottom=4pt}
%--- Công tắc lời giải bài tự luyện (ví dụ mẫu luôn dùng \loigiai)
%     \loigiaiTL[{đáp số}]{lời giải đầy đủ}
%     \TLfull -> in lời giải đầy đủ;  \TLdapso -> chỉ in một dòng "Đáp số";  \TLan -> ẩn hoàn toàn
%     (\CheDo là dethi thì mọi lời giải và đáp số đều ẩn)
\makeatletter
\newif\ifTLhien   \TLhientrue
\newif\ifTLdapso  \TLdapsofalse
\newcommand{\TLfull}{\TLhientrue\TLdapsofalse}
\newcommand{\TLdapso}{\TLhientrue\TLdapsotrue}
\newcommand{\TLan}{\TLhienfalse}
\newcommand{\loigiaiTL}[2][]{%
	\ifTLhien
		\ifTLdapso
			\def\@tlarg{#1}%
			\ifx\@tlarg\@empty\else
				\ifdefstring{\kindOPTN}{dethi}{}{\par\noindent{\bfseries\color{gold1}Đáp số:} #1\par}%
			\fi
		\else
			\loigiai{#2}%
		\fi
	\fi}
\makeatother
\TLdapso% mặc định: chỉ in đáp số bài tự luyện
%====================== BẢNG ĐÚNG/SAI (\choiceTFt) ======================
%  Hàng tiêu đề nền nâu đồng chữ trắng (Phát biểu | Đ | S), các hàng tô xen kẽ kem / trắng,
%  đường kẻ vàng kim, nhãn a) b) c) d) màu nâu đồng. Bản có đáp án: cột được chọn có vòng tròn
%  nâu đồng với dấu tích trắng (câu đúng tích ở cột Đúng, câu sai tích ở cột Sai).
%  Tiêu đề ghi đủ "Đúng | Sai": \OPTN{viettat=0} sau \input (cột Đúng/Sai rộng ra, phát biểu hẹp lại).
%  Dùng bảng gốc của ex_test: \def\BangDSgoc{1} trước %============================================================
%  PREAMBLE EDUZY v2 - dùng chung cho đề thi, bài học, chuyên đề (gói ex_test)
%  Khác bản preamble-eduzy.tex ở đầu trang, chân trang: dải nền nâu đồng tràn ngang,
%  logo + Eduzy.com.vn, ô chữ vàng kem cạnh xiên, chân trang "Trang (số)";
%  logo Eduzy chìm giữa mọi trang (tắt: \logochimtat, bật lại: \logochimbat).
%  Cách dùng ở file chính (xem mau-chuan.tex, mau-chuyen-de.tex):
%      \def\CheDo{loigiai}          % loigiai | dethi | solcolor | color | book
%      \def\tenchuong{<dòng chữ trên đầu trang>}
%      \def\thuonghieu{<chữ ở chân trang lẻ>}   % không khai báo thì dùng mặc định
%      \def\tacgia{<chữ ở chân trang chẵn>}      % không khai báo thì dùng mặc định
%      \def\khauhieu{<chữ cạnh logo ở đầu trang>}  % tùy chọn, mặc định Eduzy.com.vn
%      \input{preamble-eduzy-v2}
%      \TLdapso                      % tài liệu: \TLfull | \TLdapso | \TLan
%      \begin{document} ... \end{document}
%  Logo Eduzy vẽ bằng TikZ ngay trong file (\eduzylogo), không cần file ảnh.
%  Biên dịch: pdfLaTeX (máy có gói vietnam) hoặc XeLaTeX.
%============================================================
\documentclass[12pt,a4paper,twoside]{article}%<-- twoside: bat dau trang/chan trang chan-le
\usepackage{iftex}
\usepackage{amsmath,amssymb,tikz,mathrsfs,tkz-tab}
\usepackage{graphicx}
\usepackage{pgfplots}% đồ thị hàm số
\pgfplotsset{compat=1.17}
%\usepackage{mathptmx}
\usepackage{mathpazo}
\ifXeTeX% XeLaTeX: font chữ có đủ dấu tiếng Việt, giữ nguyên font toán mathpazo
	\usepackage[no-math]{fontspec}
	\setmainfont{TeX Gyre Pagella}
	\newfontfamily\EZqag{TeX Gyre Adventor}[NFSSFamily=qag,Ligatures=TeX]
\fi
\usepackage[many]{tcolorbox}
\usetikzlibrary{math,through,calc,intersections,angles,quotes,shapes,shapes.geometric,arrows,patterns,matrix,chains,arrows.meta,decorations.shapes,decorations.fractals,decorations.markings,shadows,decorations.pathreplacing}
\usepackage[top=1.6cm, bottom=1.8cm, left=1.8cm, right=1.8cm]{geometry}% v2: top 1.6cm (bản cũ 1.5cm) để chữ không sát dải đầu trang
\providecommand{\CheDo}{loigiai}% loigiai | dethi | solcolor | color | book
\expandafter\usepackage\expandafter[\CheDo]{ex_test}
% \tenchuong: khai báo ở file chính (trước \input) hoặc \renewcommand sau \input
%============================================================
%  [GỘP TỪ eduzy-de.tex] EDUZY - MÀU NÂU ĐỒNG / VÀNG KIM
%  + ĐẦU TRANG, CHÂN TRANG KIỂU DẢI NGANG (v2), ĐỐI XỨNG TRANG CHẴN / TRANG LẺ.
%  (không cần file eduzy-de.tex riêng nữa)
%  Khối này phải đặt SAU \usepackage[...]{ex_test} và SAU \def\tenchuong{...}
%============================================================
\usepackage{lastpage}
\usepackage{needspace}
%--- DÀN TRANG ĐỀ THI: tách câu (THÊM 05/11/2026; chỉ thêm lệnh mới, KHÔNG đổi hành vi mặc định) ---
%  Hai công tắc, mặc định tắt (= preamble chạy y như cũ). Bật bằng cách đặt TRƯỚC \input{preamble-eduzy-v2}
%  trong file đề (tools/dan_trang/dan_trang.py tự ghi khi đặt \NeedspaceCau):
%    \def\TachCauDS{1}   câu Đúng/Sai được tách: câu dẫn (kèm hình) ở cuối trang, bảng Phát biểu sang trang sau
%                        (câu dẫn và hình không tách nhau, bảng không tách giữa hai trang)
%    \def\TachCauLG{1}   bản lời giải (\CheDo khác dethi): phần đề của câu cao hơn 1/3 trang được phép tách
%  Lệnh dùng trong file đề (giá trị do tools/dan_trang/dan_trang.py đo ở bản dethi):
%    \NeedspaceCau[<cao câu dẫn + hình>]{<cao cả khối câu>}   ngay trên \begin{ex} (câu đầu phần: trên tiêu đề phần)
%        công tắc tắt: đúng bằng \Needspace{<cao cả khối câu>} như cũ;
%        \TachCauDS=1 và có [..] (câu Đúng/Sai được chọn tách): chỉ giữ câu dẫn + hình;
%        \TachCauLG=1 ở bản lời giải: đòi tối đa 1/3 \textheight.
%    \NeedspaceBang{<cao bảng Phát biểu>}%   ngay trước \choiceTFt: \TachCauDS=1 thì giữ trọn bảng trên một trang;
%        công tắc tắt: không làm gì.
\providecommand{\TachCauDS}{0}
\providecommand{\TachCauLG}{0}
\newlength{\EZnsCau}
\newcommand{\NeedspaceCau}[2][]{%
	\setlength{\EZnsCau}{#2}%
	\ifnum\TachCauDS=1 \if\relax\detokenize{#1}\relax\else\setlength{\EZnsCau}{#1}\fi\fi
	\ifnum\TachCauLG=1 \ifdefstring{\kindOPTN}{dethi}{}{\ifdim\EZnsCau>0.3333\textheight\setlength{\EZnsCau}{0.3333\textheight}\fi}\fi
	\Needspace{\EZnsCau}}
\newcommand{\NeedspaceBang}[1]{\ifnum\TachCauDS=1 \Needspace{#1}\fi}

%--- Thang màu: đậm -> nhạt
\definecolor{gold1}{RGB}{126,96,0}    % nâu đồng đậm  #7E6000
\definecolor{gold2}{RGB}{192,144,3}   % vàng kim      #C09003
\definecolor{gold3}{RGB}{255,215,99}  % vàng nhạt     #FFD763
\definecolor{gold4}{RGB}{255,231,153} % vàng rất nhạt #FFE799
\definecolor{gold5}{RGB}{255,243,205} % vàng kem      #FFF3CD
%--- Ghi đè tông xanh/đỏ cũ
\definecolor{myblue}{RGB}{126,96,0}
\colorlet{blue}{gold1}
\colorlet{cyan}{gold2}
\colorlet{teal}{gold2}

%--- Tên chương mặc định (file đề nên tự khai báo lại)
\providecommand{\tenchuong}{}
%--- Chân trang
\providecommand{\thuonghieu}{TOÁN 12 -- Eduzy.com.vn}
\providecommand{\tacgia}{Tô Gia Bảo -- 0353703783}

%--- Chữ cạnh logo ở đầu trang
\providecommand{\khauhieu}{Eduzy.com.vn}

%====================== THÔNG SỐ ĐẦU TRANG, CHÂN TRANG (v2) ======================
%  Dải nền đậm tràn ngang trang; ô chữ vàng kem có cạnh xiên, kèm 2 vạch xiên chuyển màu.
%  Đầu trang lẻ : logo + \khauhieu (trái)   | ô chữ \tenchuong (phải)
%  Chân trang lẻ: ô chữ \thuonghieu (trái)  | Trang (số) (phải)
%  Trang chẵn soi gương trang lẻ; chân trang chẵn ghi \tacgia.
%--- Màu: chỉ dùng thang vàng kim Eduzy
\colorlet{ezNen}{gold1}% dải nền đậm
\colorlet{ezVach1}{gold3}% vạch xiên sát dải nền (đậm hơn)
\colorlet{ezVach2}{gold4}% vạch xiên sát ô chữ (nhạt hơn)
\colorlet{ezO}{gold5}% ô chữ vàng kem
%--- Kích thước
\newlength{\EZdai}\setlength{\EZdai}{1.1cm}% chiều cao dải nền
\newlength{\EZo}\setlength{\EZo}{0.68cm}% chiều cao ô chữ
\newlength{\EZxien}\setlength{\EZxien}{0.55cm}% độ lệch ngang của cạnh xiên
\newlength{\EZvach}\setlength{\EZvach}{0.75cm}% bề rộng mỗi vạch xiên
\newlength{\EZkhe}\setlength{\EZkhe}{0.12cm}% khe trắng giữa vạch và ô chữ
\newlength{\EZdem}\setlength{\EZdem}{0.4cm}% khoảng đệm từ chữ tới cạnh xiên
\newlength{\EZle}\setlength{\EZle}{1.8cm}% lề trái/phải của chữ (bằng lề thân bài)
\newlength{\EZrongmin}\setlength{\EZrongmin}{4.5cm}% ô chữ rộng tối thiểu
\newlength{\EZrong}% bề rộng chữ trong ô (tính tự động)
%--- Font
\newcommand{\EZfontO}{\fontsize{10.5pt}{12pt}\selectfont\bfseries}% chữ trong ô
\newcommand{\EZfontKH}{\fontfamily{qag}\fontsize{11pt}{12pt}\selectfont\bfseries}% \khauhieu cạnh logo
\newcommand{\EZfontTrang}{\fontsize{11.5pt}{12pt}\selectfont\bfseries}% chữ "Trang" và số trang
%--- Đo bề rộng chữ trong ô
\newcommand{\EZdorong}[1]{%
	\settowidth{\EZrong}{\EZfontO #1}%
	\ifdim\EZrong<\EZrongmin\setlength{\EZrong}{\EZrongmin}\fi}
%--- Logo + chữ trắng \khauhieu (dựng sẵn một lần khi vào \begin{document})
\newsavebox{\EZlogobox}
\AtBeginDocument{\sbox{\EZlogobox}{%
		\raisebox{-0.5\height}{\eduzylogo[1cm]}\hspace{0.15cm}%
		\raisebox{-0.5\height}{\color{white}\EZfontKH\khauhieu}}}
%--- Hai vạch xiên + khe trắng. #1: điểm gốc ở mép dưới, #2: hướng (1 sang phải, -1 sang trái),
%    #3: hướng nghiêng (1 kiểu "/", -1 kiểu "\")
\newcommand{\EZvachxien}[3]{%
	\fill[white] (#1) --++({(#3)*\EZxien},\EZo) --++({(#2)*\EZkhe},0) --++({-(#3)*\EZxien},-\EZo) --cycle;
	\fill[ezVach2] ([xshift=(#2)*\EZkhe]#1) --++({(#3)*\EZxien},\EZo) --++({(#2)*\EZvach},0) --++({-(#3)*\EZxien},-\EZo) --cycle;
	\fill[ezVach1] ([xshift=(#2)*(\EZkhe+\EZvach)]#1) --++({(#3)*\EZxien},\EZo) --++({(#2)*\EZvach},0) --++({-(#3)*\EZxien},-\EZo) --cycle;
	\draw[white,line width=0.6pt] ([xshift=(#2)*(\EZkhe+\EZvach)]#1) --++({(#3)*\EZxien},\EZo);
}
%--- Chữ trong ô (căn giữa theo chiều cao ô, không lệch vì dấu tiếng Việt)
\tikzset{ezchu/.style={inner sep=0pt,text=ezNen,text height=0.72em,text depth=0.22em}}

%====================== LOGO CHÌM GIỮA TRANG ======================
%  Vẽ cùng đầu trang nên nằm dưới chữ của thân bài, có trên mọi trang.
%  Tắt: \logochimtat (từ trang đó trở đi), bật lại: \logochimbat.
%  Đổi cỡ, độ đậm: khai báo trước \input, ví dụ \def\EZchimcao{9cm} \def\EZchimdo{0.1}
\providecommand{\EZchimcao}{10cm}% chiều cao logo chìm
\providecommand{\EZchimdo}{0.12}% độ đậm (0 = trong suốt, 1 = rõ hẳn)
\newif\ifEZchim \EZchimtrue
\newcommand{\logochimtat}{\EZchimfalse}
\newcommand{\logochimbat}{\EZchimtrue}
%  Dựng sẵn logo mờ một lần khi vào \begin{document}; transparency group làm mờ cả logo
%  một lượt nên không lộ vệt ở chỗ các dải màu chồng nhau.
\newsavebox{\EZchimtho}
\newsavebox{\EZchimbox}
\AtBeginDocument{%
	\sbox{\EZchimtho}{\eduzylogo[\EZchimcao]}%
	\sbox{\EZchimbox}{\begin{tikzpicture}
			\begin{scope}[transparency group,opacity=\EZchimdo]
				\node[inner sep=0pt] {\usebox{\EZchimtho}};
			\end{scope}
	\end{tikzpicture}}}
\newcommand{\EZvechim}{%
	\ifEZchim
		\node[inner sep=0pt] at (current page.center) {\usebox{\EZchimbox}};
	\fi}

%====================== ĐẦU TRANG - TRANG LẺ ======================
%  (logo bên trái, ô chữ \tenchuong bên phải)
\newcommand{\myfancyheadOdd}{%
	\EZdorong{\tenchuong}%
	\begin{tikzpicture}[remember picture,overlay]
		\coordinate (T) at (current page.north west);
		\coordinate (P) at (current page.north east);
		%--- logo chìm giữa trang
		\EZvechim
		%--- dải nền
		\fill[ezNen] (T) rectangle ([yshift=-\EZdai]P);
		%--- ô chữ: X là góc dưới bên trái, cạnh trái xiên kiểu "/"
		\coordinate (X) at ([xshift=-\EZle-\EZrong-\EZdem-\EZxien,yshift=-\EZdai]P);
		\fill[ezO] (X) --++(\EZxien,\EZo) -- ([yshift=-\EZdai+\EZo]P) -- ([yshift=-\EZdai]P) --cycle;
		\EZvachxien{X}{-1}{1}
		\node[ezchu,anchor=east] at ([xshift=-\EZle,yshift=-\EZdai+0.5*\EZo]P) {\EZfontO\tenchuong};
		%--- logo + tên
		\node[inner sep=0pt,anchor=west] at ([xshift=\EZle,yshift=-0.5*\EZdai]T) {\usebox{\EZlogobox}};
	\end{tikzpicture}%
}

%====================== ĐẦU TRANG - TRANG CHẴN =====================
%  (soi gương: ô chữ \tenchuong bên trái, logo bên phải)
\newcommand{\myfancyheadEven}{%
	\EZdorong{\tenchuong}%
	\begin{tikzpicture}[remember picture,overlay]
		\coordinate (T) at (current page.north west);
		\coordinate (P) at (current page.north east);
		%--- logo chìm giữa trang
		\EZvechim
		%--- dải nền
		\fill[ezNen] (T) rectangle ([yshift=-\EZdai]P);
		%--- ô chữ: X là góc dưới bên phải, cạnh phải xiên kiểu "\"
		\coordinate (X) at ([xshift=\EZle+\EZrong+\EZdem+\EZxien,yshift=-\EZdai]T);
		\fill[ezO] (X) --++(-\EZxien,\EZo) -- ([yshift=-\EZdai+\EZo]T) -- ([yshift=-\EZdai]T) --cycle;
		\EZvachxien{X}{1}{-1}
		\node[ezchu,anchor=west] at ([xshift=\EZle,yshift=-\EZdai+0.5*\EZo]T) {\EZfontO\tenchuong};
		%--- logo + tên
		\node[inner sep=0pt,anchor=east] at ([xshift=-\EZle,yshift=-0.5*\EZdai]P) {\usebox{\EZlogobox}};
	\end{tikzpicture}%
}

%====================== CHÂN TRANG - TRANG LẺ ======================
%  (ô chữ \thuonghieu bên trái, "Trang (số)" bên phải)
\newcommand{\myfancyfootOdd}{%
	\EZdorong{\thuonghieu}%
	\begin{tikzpicture}[remember picture,overlay]
		\coordinate (T) at (current page.south west);
		\coordinate (P) at (current page.south east);
		%--- dải nền
		\fill[ezNen] (T) rectangle ([yshift=\EZdai]P);
		%--- ô chữ sát mép trên dải: X là góc dưới bên phải, cạnh phải xiên kiểu "/"
		\coordinate (X) at ([xshift=\EZle+\EZrong+\EZdem,yshift=\EZdai-\EZo]T);
		\fill[ezO] ([yshift=\EZdai-\EZo]T) -- (X) --++(\EZxien,\EZo) -- ([yshift=\EZdai]T) --cycle;
		\EZvachxien{X}{1}{1}
		\node[ezchu,anchor=west] at ([xshift=\EZle,yshift=\EZdai-0.5*\EZo]T) {\EZfontO\thuonghieu};
		%--- số trang
		\node[circle,draw=white,line width=0.9pt,text=white,inner sep=0pt,minimum size=0.62cm,anchor=east]
		(so) at ([xshift=-\EZle,yshift=0.5*\EZdai]P) {\EZfontTrang\thepage};
		\node[inner sep=0pt,text=white,anchor=base east] (tr) at ([xshift=-0.15cm]so.base west) {\EZfontTrang Trang};
	\end{tikzpicture}%
}

%====================== CHÂN TRANG - TRANG CHẴN ====================
%  (soi gương: "Trang (số)" bên trái, ô chữ \tacgia bên phải)
\newcommand{\myfancyfootEven}{%
	\EZdorong{\tacgia}%
	\begin{tikzpicture}[remember picture,overlay]
		\coordinate (T) at (current page.south west);
		\coordinate (P) at (current page.south east);
		%--- dải nền
		\fill[ezNen] (T) rectangle ([yshift=\EZdai]P);
		%--- ô chữ sát mép trên dải: X là góc dưới bên trái, cạnh trái xiên kiểu "\"
		\coordinate (X) at ([xshift=-\EZle-\EZrong-\EZdem,yshift=\EZdai-\EZo]P);
		\fill[ezO] ([yshift=\EZdai-\EZo]P) -- (X) --++(-\EZxien,\EZo) -- ([yshift=\EZdai]P) --cycle;
		\EZvachxien{X}{-1}{-1}
		\node[ezchu,anchor=east] at ([xshift=-\EZle,yshift=\EZdai-0.5*\EZo]P) {\EZfontO\tacgia};
		%--- số trang
		%    (vòng tròn ẩn so0 nằm giữa dải để lấy đường chân chữ, giống trang lẻ)
		\node[circle,inner sep=0pt,minimum size=0.62cm] (so0) at ([yshift=0.5*\EZdai]T) {\EZfontTrang\phantom{\thepage}};
		\node[inner sep=0pt,text=white,anchor=base west] (tr) at ($(so0.base -| T)+(\EZle,0)$) {\EZfontTrang Trang};
		\node[circle,draw=white,line width=0.9pt,text=white,inner sep=0pt,minimum size=0.62cm,anchor=base west]
		(so) at ([xshift=0.15cm]tr.base east) {\EZfontTrang\thepage};
	\end{tikzpicture}%
}

%--- Giữ tương thích với các file cũ đang gọi \myfancyhead, \myfancyfoot
\newcommand{\myfancyhead}{\myfancyheadOdd}
\newcommand{\myfancyfoot}{\myfancyfootOdd}

%======================= GẮN VÀO TRANG =============================
\pagestyle{fancy}
\fancyhf{}
\fancyhead[LO]{\myfancyheadOdd}
\fancyhead[RE]{\myfancyheadEven}
\fancyfoot[LO]{\myfancyfootOdd}
\fancyfoot[RE]{\myfancyfootEven}
\renewcommand{\headrulewidth}{0pt}
\renewcommand{\footrulewidth}{0pt}
\fancypagestyle{plain}{%
	\fancyhf{}
	\fancyhead[LO]{\myfancyheadOdd}
	\fancyhead[RE]{\myfancyheadEven}
	\fancyfoot[LO]{\myfancyfootOdd}
	\fancyfoot[RE]{\myfancyfootEven}
	\renewcommand{\headrulewidth}{0pt}
	\renewcommand{\footrulewidth}{0pt}
}
%=================== HẾT PHẦN GỘP TỪ eduzy-de.tex ===================
\newcommand{\indam}[1]{
	{\fontsize{12pt}{9pt}\sffamily \selectfont\bfseries\color{gold1} #1}
}
\newcommand{\ind}[1]{
	{\fontsize{14pt}{9pt}\fontfamily{qag} \selectfont\bfseries\color{gold1} #1}
}
\renewcommand{\FalseEX}{\stepcounter{dapan}\textbf{\selectfont\color{gold1}\Alph{dapan}.}\ignorespaces}
\renewtheorem{ex}{\color{gold1}Câu}
\renewcommand{\baselinestretch}{1.1}
\newcommand{\hoac}[1]{
	\left[\begin{aligned}#1\end{aligned}\right.}
\newcommand{\heva}[1]{
	\left\{\begin{aligned}#1\end{aligned}\right.}
\usepackage{esvect}
\def\vec{\vv}
\def\vt{\vv}
\def\overrightarrow{\vv}
%====================== LOGO EDUZY VẼ BẰNG TIKZ ======================
%  Đường viền dò từ logo gốc; dải màu vàng kim đo trực tiếp trên ảnh gốc
%  (nghiêng 21 độ, tô bằng 80 dải hẹp nên hiển thị giống nhau trên mọi trình xem PDF).
%  Dùng: \eduzylogo          -> cao 2.4cm (mặc định)
%        \eduzylogo[1.5cm]   -> cao 1.5cm
\definecolor{ezText}{RGB}{245,200,46}
\newcommand{\eduzylogo}[1][2.4cm]{%
	\resizebox{!}{#1}{%
		\begin{tikzpicture}[x=0.01cm,y=0.01cm]
			\path[use as bounding box] (0,0) rectangle (224,-224);
			% mũ tốt nghiệp + vương miện
			\begin{scope}
				\clip
				(119.4,-176.1)..controls(118.9,-176.1)and(116.8,-176)..(114.6,-176)
				..controls(101.4,-176)and(85.8,-175.9)..(83.7,-175.7)
				..controls(81.8,-175.6)and(81,-175.6)..(80.2,-175.7)
				..controls(78.8,-175.9)and(77.7,-175.9)..(76.2,-175.7)
				..controls(74.9,-175.5)and(70.8,-175.5)..(64.3,-175.7)
				..controls(59.1,-175.9)and(58.3,-175.7)..(58.6,-173.7)
				..controls(58.6,-172.8)and(59.6,-169.6)..(60.3,-168)
				..controls(60.8,-166.9)and(61.4,-165.8)..(62.3,-165)
				..controls(63.1,-164.2)and(64.3,-163.4)..(65,-163.2)
				..controls(65.5,-163.1)and(72.3,-163)..(76.6,-163.1)
				..controls(77.9,-163.1)and(80.1,-163.1)..(81.4,-163)
				..controls(82.8,-162.9)and(86.3,-162.9)..(89.2,-162.9)
				..controls(93,-162.9)and(94.9,-162.8)..(96,-162.7)
				..controls(97.2,-162.5)and(101,-162.5)..(114,-162.4)
				..controls(123,-162.3)and(134.9,-162.2)..(140.2,-162.1)
				..controls(150.1,-162)and(152.2,-162.1)..(153.6,-162.4)
				..controls(154.4,-162.6)and(157.1,-163.8)..(158.1,-164.4)
				..controls(158.8,-164.8)and(159.7,-165.7)..(159.9,-166.3)
				..controls(160,-166.5)and(160.2,-167.1)..(160.2,-167.6)
				..controls(160.3,-168)and(160.6,-169.4)..(160.9,-170.6)
				..controls(161.2,-171.7)and(161.5,-173.1)..(161.6,-173.6)
				..controls(161.8,-174.4)and(161.8,-174.5)..(161.5,-174.9)
				..controls(161.2,-175.6)and(160.8,-175.9)..(159.9,-176)
				..controls(159.2,-176.1)and(154.1,-176.2)..(151.5,-176.2)
				..controls(148.1,-176.1)and(122.6,-176.1)..(121.6,-176.2)
				..controls(120.8,-176.2)and(119.8,-176.2)..(119.4,-176.1)--cycle
				(130.4,-160)..controls(129,-160)and(119.8,-160)..(110.1,-159.9)
				..controls(100.3,-159.8)and(88.5,-159.7)..(83.9,-159.6)
				..controls(77.8,-159.5)and(75.3,-159.5)..(74.5,-159.6)
				..controls(73.9,-159.7)and(72.5,-159.8)..(71.3,-159.8)
				..controls(69.6,-159.8)and(69,-159.8)..(68.5,-159.6)
				..controls(67.6,-159.3)and(66.8,-158.6)..(66.2,-157.6)
				..controls(66,-157.3)and(65.6,-156.5)..(65.3,-156.1)
				..controls(64.6,-155.1)and(64.5,-154.9)..(63.7,-152.9)
				..controls(63.2,-151.6)and(62.9,-150.9)..(61.8,-149.6)
				..controls(61.5,-149.2)and(61.1,-148.5)..(60.9,-147.9)
				..controls(60.7,-147.4)and(60.3,-146.6)..(60.1,-146.1)
				..controls(59.9,-145.5)and(59.5,-144.9)..(59,-144.3)
				..controls(58.5,-143.7)and(58.2,-143.2)..(57.9,-142.5)
				..controls(57.7,-141.9)and(57.3,-141.1)..(57,-140.7)
				..controls(56.6,-140.3)and(56.1,-139.3)..(55.8,-138.5)
				..controls(55.5,-137.8)and(55.1,-136.8)..(54.9,-136.5)
				..controls(54.7,-136.1)and(54.3,-135.2)..(53.9,-134.5)
				..controls(53.6,-133.8)and(53,-132.7)..(52.6,-132)
				..controls(51.6,-130.5)and(51.4,-130.2)..(51.1,-129)
				..controls(50.6,-127.5)and(50.6,-127.4)..(49.7,-126.3)
				..controls(49.2,-125.7)and(48.6,-124.8)..(48.3,-124.2)
				..controls(48,-123.6)and(47.5,-122.7)..(47.3,-122.2)
				..controls(47,-121.8)and(46.5,-120.7)..(46.1,-119.8)
				..controls(45.8,-118.9)and(45.2,-117.8)..(44.9,-117.2)
				..controls(44.2,-116.1)and(43.3,-114.4)..(42.3,-112.2)
				..controls(41.8,-111.3)and(41.3,-110.4)..(41.1,-110.1)
				..controls(40.5,-109.4)and(40.3,-109)..(40,-107.8)
				..controls(39.8,-107)and(39.6,-106.5)..(39.2,-106)
				..controls(38.9,-105.6)and(38.4,-104.7)..(38,-104)
				..controls(37.7,-103.4)and(37.2,-102.5)..(37,-102.1)
				..controls(36.7,-101.7)and(36.3,-100.8)..(36,-100.1)
				..controls(35.6,-98.9)and(35.3,-98.4)..(34.6,-97.6)
				..controls(34.2,-97.1)and(34.2,-96.3)..(34.6,-96)
				..controls(35,-95.7)and(35.5,-95.9)..(36.2,-96.7)
				..controls(36.5,-97)and(37.2,-97.5)..(37.6,-97.7)..controls(38,-98)and(38.6,-98.4)..(39.1,-98.9)
				..controls(39.6,-99.4)and(40.3,-100)..(40.6,-100.3)
				..controls(41.3,-100.8)and(42.4,-101.8)..(43.7,-102.9)
				..controls(44,-103.2)and(44.6,-103.7)..(45,-103.9)
				..controls(45.4,-104.2)and(45.9,-104.6)..(46.1,-104.8)
				..controls(47.1,-105.9)and(48.8,-107.1)..(50.5,-107.9)
				..controls(51.8,-108.6)and(52.3,-109)..(53.4,-110)
				..controls(54.8,-111.2)and(55.6,-111.6)..(57,-112.1)
				..controls(57.5,-112.3)and(58.5,-112.8)..(59.4,-113.3)
				..controls(60.2,-113.7)and(61.3,-114.3)..(61.9,-114.5)
				..controls(62.4,-114.7)and(63.3,-115.1)..(63.8,-115.3)
				..controls(64.4,-115.6)and(65.4,-115.9)..(66.2,-116.1)
				..controls(67,-116.3)and(68.4,-116.6)..(69.2,-116.9)
				..controls(72.3,-117.9)and(72.8,-118)..(78,-118.1)
				..controls(84,-118.3)and(86.1,-118.1)..(88.1,-117.1)
				..controls(89.4,-116.5)and(93,-115)..(93.7,-114.9)
				..controls(94.8,-114.6)and(95.4,-114.2)..(95.9,-113.4)
				..controls(96.2,-113)and(96.7,-112.5)..(97.1,-112.2)
				..controls(98,-111.6)and(101.3,-108.4)..(102.1,-107.2)
				..controls(102.4,-106.8)and(102.9,-106.2)..(103.2,-105.8)
				..controls(103.5,-105.5)and(104,-104.8)..(104.2,-104.2)
				..controls(104.5,-103.6)and(105,-102.8)..(105.3,-102.4)
				..controls(105.6,-102)and(106.1,-101.2)..(106.5,-100.6)
				..controls(107.1,-99.5)and(108.6,-97.5)..(109.5,-96.5)
				..controls(109.9,-96.1)and(110.1,-96)..(110.4,-96)
				..controls(111.1,-96)and(111.6,-96.5)..(111.8,-97.6)
				..controls(112,-98.6)and(112.3,-99.3)..(112.8,-99.9)
				..controls(113.7,-100.9)and(115.2,-103.1)..(115.7,-104.2)
				..controls(116.1,-105)and(116.5,-105.5)..(116.9,-106)
				..controls(117.3,-106.4)and(117.8,-106.9)..(118,-107.3)
				..controls(118.2,-107.6)and(118.6,-108)..(119,-108.3)
				..controls(119.6,-108.8)and(121.1,-110.3)..(121.8,-111.2)
				..controls(122.1,-111.7)and(122.5,-112)..(123.4,-112.4)
				..controls(124.8,-113.2)and(125.8,-113.8)..(127,-114.6)
				..controls(127.5,-114.9)and(128.1,-115.3)..(128.2,-115.4)
				..controls(128.4,-115.5)and(129.3,-115.8)..(130.2,-116)
				..controls(131.1,-116.2)and(132.5,-116.6)..(133.3,-116.9)
				..controls(135.7,-117.9)and(136.1,-117.9)..(142.5,-117.8)
				..controls(145.8,-117.8)and(148.4,-117.7)..(148.8,-117.6)
				..controls(149.2,-117.6)and(150.2,-117.2)..(151,-116.8)
				..controls(152,-116.4)and(152.8,-116)..(153.6,-115.9)
				..controls(154.5,-115.7)and(155.2,-115.4)..(156.3,-114.9)
				..controls(157.1,-114.5)and(158.2,-114)..(158.7,-113.8)
				..controls(159.1,-113.7)and(159.8,-113.3)..(160.2,-113)
				..controls(160.6,-112.7)and(161.3,-112.4)..(161.7,-112.2)
				..controls(162.7,-111.9)and(163.2,-111.7)..(164.1,-111.1)
				..controls(164.5,-110.8)and(165.3,-110.4)..(165.9,-110)
				..controls(166.5,-109.7)and(167.3,-109.2)..(167.6,-108.9)
				..controls(167.9,-108.6)and(168.6,-108.2)..(169,-108)
				..controls(169.4,-107.8)and(170.2,-107.3)..(170.7,-106.9)
				..controls(171.2,-106.5)and(172.1,-105.8)..(172.7,-105.4)
				..controls(173.2,-105)and(174.1,-104.4)..(174.5,-104)
				..controls(174.9,-103.7)and(175.7,-103.1)..(176.2,-102.6)
				..controls(176.7,-102.2)and(177.4,-101.5)..(177.9,-101.2)
				..controls(178.3,-100.8)and(178.8,-100.4)..(179,-100.2)
				..controls(179.2,-100)and(179.6,-99.6)..(180,-99.3)
				..controls(180.4,-99.1)and(180.9,-98.6)..(181.2,-98.4)
				..controls(181.5,-98.1)and(182.1,-97.7)..(182.5,-97.5)
				..controls(183.2,-97.2)and(183.5,-96.9)..(184.3,-95.9)
				..controls(184.6,-95.6)and(184.6,-95.6)..(184.8,-95.7)
				..controls(185,-95.8)and(185.1,-96)..(185.1,-96.2)
				..controls(185.1,-96.5)and(185,-96.7)..(184.6,-97.2)
				..controls(184.1,-97.8)and(184,-97.9)..(183.9,-98.8)
				..controls(183.7,-100.3)and(183.3,-101.3)..(182.9,-101.8)
				..controls(182.6,-102.1)and(182.3,-102.6)..(182.2,-102.9)
				..controls(182,-103.3)and(181.6,-104)..(181.2,-104.6)
				..controls(180.9,-105.2)and(180.5,-106.1)..(180.4,-106.6)
				..controls(179.9,-108)and(179.7,-108.5)..(179,-109.4)
				..controls(178.6,-109.8)and(178.2,-110.5)..(177.9,-111)
				..controls(177.7,-111.5)and(177.3,-112.3)..(177,-112.7)
				..controls(176.6,-113.2)and(176.4,-113.7)..(176.3,-114.2)
				..controls(176.2,-114.7)and(175.9,-115.2)..(175.4,-116)
				..controls(175,-116.6)and(174.4,-117.7)..(174,-118.5)
				..controls(173.7,-119.2)and(173.3,-120.1)..(173,-120.5)
				..controls(172.8,-120.9)and(172.4,-121.8)..(172.1,-122.5)
				..controls(171.8,-123.2)and(171.3,-124.1)..(171,-124.6)
				..controls(170.8,-125.1)and(170.3,-126)..(169.9,-126.6)
				..controls(169.6,-127.3)and(169.2,-128.1)..(169,-128.3)
				..controls(168.8,-128.6)and(168.5,-129.3)..(168.3,-129.8)
				..controls(168.1,-130.3)and(167.8,-131.1)..(167.5,-131.6)
				..controls(166.9,-132.5)and(164.2,-137.9)..(164.2,-138.1)
				..controls(164.2,-138.1)and(163.9,-138.8)..(163.4,-139.7)
				..controls(162.9,-140.5)and(162.4,-141.5)..(162.3,-141.9)
				..controls(161.9,-143)and(161.6,-143.4)..(160.9,-144.1)
				..controls(160.4,-144.6)and(160.3,-144.8)..(160.1,-145.4)
				..controls(160,-145.8)and(159.6,-146.8)..(159.2,-147.8)
				..controls(158.7,-148.8)and(158.1,-150)..(157.9,-150.6)
				..controls(157.6,-151.1)and(157.2,-151.9)..(156.9,-152.4)
				..controls(156.5,-152.8)and(156.1,-153.7)..(155.8,-154.4)
				..controls(155.6,-155)and(155.2,-155.9)..(154.9,-156.3)
				..controls(154.7,-156.7)and(154.4,-157.4)..(154.2,-157.9)
				..controls(154.1,-158.4)and(153.8,-159)..(153.5,-159.3)
				..controls(152.8,-160.1)and(153,-160.1)..(142.4,-160.1)
				..controls(137.2,-160.1)and(131.8,-160.1)..(130.4,-160)--cycle
				(141.5,-103.7)..controls(140.8,-103)and(139.8,-102.1)..(139.2,-101.6)
				..controls(138.5,-101.2)and(137.9,-100.6)..(137.8,-100.4)
				..controls(137,-99.5)and(134,-96.5)..(133.2,-96)
				..controls(132.7,-95.7)and(131.8,-94.9)..(131.1,-94.3)
				..controls(129.4,-92.7)and(127.9,-91.6)..(126.6,-91.1)
				..controls(125.2,-90.5)and(124.1,-90)..(122.9,-89.1)
				..controls(122.1,-88.6)and(121.3,-88.2)..(119.9,-87.7)
				..controls(118.8,-87.3)and(117.5,-86.8)..(116.9,-86.5)
				..controls(114.3,-85.5)and(113.7,-85.4)..(109.7,-85.5)
				..controls(105.3,-85.7)and(103.9,-85.9)..(101.6,-87)
				..controls(100.7,-87.4)and(99.4,-87.8)..(97.9,-88)
				..controls(97,-88.2)and(96.6,-88.3)..(95.8,-88.9)
				..controls(95.4,-89.2)and(94.7,-89.6)..(94.2,-89.8)
				..controls(92.9,-90.3)and(92.3,-90.7)..(90.9,-92.1)
				..controls(90.2,-92.7)and(89.4,-93.3)..(89,-93.5)..controls(87,-94.5)and(85.2,-96)..(83.6,-97.9)
				..controls(83.3,-98.2)and(82.5,-98.9)..(81.8,-99.4)
				..controls(80.7,-100.1)and(80.5,-100.3)..(80,-101.1)
				..controls(78.7,-103.1)and(77.2,-103.9)..(76.4,-103)
				..controls(75.7,-102.2)and(75.7,-103.2)..(75.7,-83.6)
				..controls(75.6,-73.6)and(75.6,-64.8)..(75.5,-64)
				..controls(75.5,-63.2)and(75.5,-61.4)..(75.6,-59.8)--(75.7,-57.1)--(76,-56.8)
				..controls(76.8,-56.2)and(78,-56.3)..(79.8,-57.2)
				..controls(80.4,-57.5)and(81.3,-57.8)..(81.8,-57.9)
				..controls(82.3,-58.1)and(83,-58.4)..(83.6,-58.8)..controls(84.3,-59.2)and(85,-59.5)..(86,-59.8)
				..controls(86.8,-60.1)and(88,-60.6)..(88.7,-60.9)
				..controls(89.3,-61.2)and(90.6,-61.7)..(91.4,-62)
				..controls(92.2,-62.3)and(93.7,-62.9)..(94.6,-63.3)
				..controls(95.6,-63.7)and(96.8,-64.1)..(97.4,-64.3)
				..controls(98,-64.5)and(99,-64.9)..(99.7,-65.2)
				..controls(100.4,-65.6)and(101.5,-66)..(102.3,-66.2)
				..controls(104.2,-66.8)and(105,-67.1)..(106.2,-67.8)
				..controls(107.9,-68.7)and(110.2,-69.4)..(111.3,-69.2)
				..controls(112.5,-69)and(115,-68.2)..(115.6,-67.8)
				..controls(117.3,-66.7)and(117.6,-66.6)..(119.5,-66.2)
				..controls(121.9,-65.8)and(122.1,-65.7)..(123,-65.1)
				..controls(124.1,-64.4)and(124.4,-64.3)..(126,-63.9)
				..controls(127,-63.7)and(127.7,-63.4)..(128.4,-63.1)
				..controls(129,-62.8)and(129.7,-62.5)..(130.6,-62.4)
				..controls(132.1,-62.1)and(134,-61.5)..(135.2,-61)
				..controls(135.7,-60.8)and(136.7,-60.5)..(137.3,-60.3)
				..controls(138.4,-60)and(139.3,-59.7)..(141.1,-59)
				..controls(141.6,-58.7)and(142,-58.6)..(142.5,-58.6)
				..controls(143.2,-58.6)and(143.3,-58.7)..(143.6,-59)
				..controls(144,-59.5)and(144.1,-59.9)..(144,-60.8)
				..controls(144,-61.2)and(144,-71.1)..(144,-82.8)--(143.9,-104.1)--(143.7,-104.5)
				..controls(143.5,-104.7)and(143.3,-104.8)..(143,-104.9)
				..controls(142.6,-104.9)and(142.5,-104.8)..(141.5,-103.7)--cycle
				(27.8,-91.4)..controls(27.1,-90.9)and(26.3,-90)..(26.1,-89.5)
				..controls(26,-89.2)and(25.8,-88.6)..(25.7,-87.9)
				..controls(25.4,-86.2)and(25.6,-84.8)..(26.5,-83.5)
				..controls(27.4,-82.3)and(28.3,-82)..(29.4,-82.3)
				..controls(31.1,-82.8)and(32.1,-84.9)..(32,-87.4)
				..controls(31.9,-89.3)and(31.2,-90.6)..(29.9,-91.4)
				..controls(29.1,-91.9)and(28.6,-91.9)..(27.8,-91.4)--cycle
				(190.1,-91.3)..controls(189.7,-91.2)and(189.3,-90.9)..(189.1,-90.7)
				..controls(188.3,-89.8)and(187.6,-87.7)..(187.6,-86.2)
				..controls(187.6,-85.5)and(187.7,-85.2)..(188.2,-84.4)
				..controls(189.1,-82.6)and(190.1,-81.8)..(191.3,-82.1)
				..controls(192,-82.3)and(193.5,-83.7)..(193.8,-84.4)
				..controls(194.1,-85.1)and(194.2,-87.1)..(194,-88.3)
				..controls(193.8,-89.8)and(193.3,-90.5)..(192.3,-91.3)
				..controls(191.7,-91.7)and(190.9,-91.7)..(190.1,-91.3)--cycle
				(173.7,-83.4)..controls(172.2,-82.5)and(172,-82.5)..(170.6,-82.4)
				..controls(169.7,-82.4)and(168.8,-82.4)..(168.3,-82.5)
				..controls(167.2,-82.7)and(166.9,-82.7)..(166.6,-82.2)
				..controls(166.3,-81.8)and(166.3,-81)..(166.8,-80)
				..controls(167.2,-78.9)and(167.4,-77.4)..(167.7,-74.1)
				..controls(167.8,-72.8)and(167.9,-71.8)..(167.9,-71.7)
				..controls(168,-71.7)and(168,-70.2)..(168,-68.4)..controls(168,-65.4)and(168,-64.9)..(167.7,-63)
				..controls(167.3,-60.5)and(167.1,-59.9)..(166.4,-58.4)
				..controls(166.1,-57.7)and(165.8,-56.9)..(165.7,-56.6)
				..controls(165.7,-56.2)and(165.4,-55.6)..(165.2,-55.2)
				..controls(164.6,-53.8)and(161.9,-50.7)..(161,-50.3)
				..controls(159.8,-49.7)and(157.9,-50)..(155.2,-51.2)
				..controls(154.7,-51.4)and(153.4,-51.8)..(152.2,-52.1)
				..controls(150.8,-52.4)and(149.4,-52.9)..(148.6,-53.2)
				..controls(147.8,-53.6)and(146.7,-54)..(146,-54.2)
				..controls(145.4,-54.4)and(144.5,-54.7)..(144,-54.9)
				..controls(143.6,-55.2)and(142.7,-55.5)..(141.8,-55.7)
				..controls(141,-55.9)and(139.6,-56.3)..(138.7,-56.6)
				..controls(137.8,-57)and(136.3,-57.4)..(135.3,-57.6)
				..controls(133.8,-57.9)and(133.2,-58.1)..(132.1,-58.7)
				..controls(131.1,-59.2)and(130.3,-59.4)..(129.4,-59.6)
				..controls(128.7,-59.8)and(127.8,-60)..(127.4,-60.1)
				..controls(127,-60.3)and(126,-60.5)..(125.1,-60.7)
				..controls(123.2,-61.1)and(121.7,-61.5)..(120.2,-62.1)
				..controls(119.7,-62.3)and(118.1,-62.8)..(116.7,-63.2)
				..controls(114.6,-63.9)and(114,-64.1)..(113,-64.7)
				..controls(111.9,-65.4)and(111.8,-65.4)..(111,-65.5)
				..controls(109.9,-65.5)and(109.3,-65.4)..(107.7,-64.7)
				..controls(107,-64.4)and(106,-64)..(105.3,-63.9)
				..controls(103.5,-63.5)and(102.2,-63)..(101.2,-62.5)
				..controls(100.1,-61.8)and(99.3,-61.5)..(97.8,-61.1)
				..controls(97.2,-61)and(96.2,-60.7)..(95.7,-60.4)
				..controls(95.1,-60.1)and(94.2,-59.8)..(93.5,-59.6)
				..controls(92.9,-59.4)and(91.7,-59)..(91,-58.7)..controls(90.2,-58.4)and(88.9,-58)..(88,-57.7)
				..controls(87.2,-57.5)and(86.1,-57.1)..(85.5,-56.8)
				..controls(85,-56.6)and(84.1,-56.2)..(83.5,-56)..controls(82.9,-55.7)and(81.8,-55.2)..(81,-54.9)
				..controls(80.1,-54.5)and(79.1,-54.2)..(78.7,-54.1)
				..controls(77.9,-53.8)and(77.3,-53.6)..(76.2,-52.9)
				..controls(75.7,-52.6)and(75.2,-52.4)..(74.1,-52.2)
				..controls(72.5,-51.9)and(72.4,-51.8)..(71.3,-51.1)
				..controls(69.9,-50.1)and(69.4,-49.9)..(67.9,-49.8)
				..controls(67,-49.7)and(66.6,-49.6)..(66.3,-49.4)
				..controls(66.1,-49.3)and(65.7,-49)..(65.4,-48.8)
				..controls(65.1,-48.7)and(64.2,-48.4)..(63.6,-48.2)
				..controls(62.9,-48)and(61.7,-47.6)..(60.9,-47.2)
				..controls(60.1,-46.8)and(58.9,-46.4)..(58.2,-46.2)
				..controls(57.6,-46)and(56.8,-45.7)..(56.5,-45.5)
				..controls(54.9,-44.6)and(54.6,-44.5)..(53.2,-44.2)
				..controls(51.7,-43.8)and(51,-43.6)..(49.9,-42.9)
				..controls(49.5,-42.6)and(48.9,-42.4)..(47.5,-42)
				..controls(45.7,-41.6)and(45.6,-41.5)..(44.3,-40.6)
				..controls(43.5,-40.1)and(42.8,-39.6)..(42.7,-39.4)
				..controls(42.4,-38.7)and(43,-38.1)..(44.5,-37.6)..controls(44.9,-37.5)and(46,-37)..(46.9,-36.6)
				..controls(47.9,-36.1)and(48.8,-35.8)..(49.6,-35.6)
				..controls(50.5,-35.4)and(51.3,-35.2)..(52.4,-34.6)
				..controls(53.3,-34.1)and(54.5,-33.7)..(54.9,-33.5)
				..controls(55.4,-33.4)and(56.4,-33)..(57.2,-32.6)
				..controls(58.6,-32)and(60.1,-31.5)..(61.6,-31.2)
				..controls(62.3,-31.1)and(62.7,-31)..(63.2,-30.7)
				..controls(64.1,-30)and(64.3,-29.9)..(65.3,-29.7)
				..controls(65.7,-29.6)and(66.5,-29.2)..(67.1,-28.9)
				..controls(67.6,-28.6)and(68.8,-28.2)..(69.6,-27.9)
				..controls(70.5,-27.7)and(71.6,-27.3)..(72.1,-27)
				..controls(72.6,-26.7)and(73.8,-26.2)..(74.8,-25.9)
				..controls(75.7,-25.5)and(76.7,-25.1)..(77.1,-25)
				..controls(77.4,-24.8)and(78.5,-24.4)..(79.6,-24.1)
				..controls(80.6,-23.8)and(82,-23.2)..(82.8,-22.9)
				..controls(83.5,-22.5)and(84.8,-22)..(85.8,-21.8)
				..controls(86.7,-21.5)and(87.8,-21.1)..(88.1,-20.9)
				..controls(88.4,-20.7)and(89.3,-20.3)..(90.1,-20.1)
				..controls(90.8,-19.8)and(91.9,-19.4)..(92.4,-19.1)
				..controls(93.5,-18.5)and(94.4,-18.2)..(95.4,-18.1)
				..controls(95.7,-18)and(96.4,-17.8)..(96.9,-17.5)
				..controls(97.5,-17.3)and(98.4,-16.9)..(99,-16.6)
				..controls(99.7,-16.3)and(100.9,-15.9)..(101.9,-15.6)
				..controls(102.9,-15.3)and(104.4,-14.7)..(105.3,-14.3)
				..controls(108.1,-13.2)and(108.5,-13.1)..(110.1,-13.1)
				..controls(111.4,-13.1)and(111.5,-13.1)..(112.4,-13.5)
				..controls(113.2,-14)and(115,-14.8)..(116.6,-15.4)
				..controls(116.9,-15.5)and(117.9,-15.8)..(118.8,-16.1)
				..controls(120.9,-16.6)and(121.5,-16.9)..(122.6,-17.9)
				..controls(123,-18.3)and(123.7,-18.8)..(123.9,-18.9)
				..controls(124.6,-19.4)and(126,-19.8)..(127.9,-20.2)
				..controls(129.6,-20.5)and(129.8,-20.6)..(131.1,-21.4)
				..controls(131.5,-21.7)and(132.4,-22.1)..(133.1,-22.3)
				..controls(136.1,-23.2)and(138.3,-24.1)..(140.2,-25)
				..controls(141,-25.5)and(141.5,-25.6)..(142.5,-25.8)
				..controls(145.1,-26.3)and(146,-26.6)..(149.4,-28.1)
				..controls(151.3,-28.9)and(153.5,-29.8)..(154.3,-30.1)
				..controls(155.1,-30.4)and(156.2,-30.8)..(156.6,-31)
				..controls(157.1,-31.2)and(158.2,-31.6)..(159.1,-31.8)
				..controls(160.6,-32.2)and(160.9,-32.3)..(161.7,-32.9)
				..controls(162.4,-33.3)and(162.9,-33.6)..(163.7,-33.8)
				..controls(164.3,-34)and(164.8,-34.1)..(164.9,-34.1)
				..controls(165,-34.1)and(165.6,-34.4)..(166.4,-34.8)
				..controls(167.3,-35.2)and(168.1,-35.4)..(169,-35.6)
				..controls(170,-35.8)and(170.5,-36)..(171.7,-36.6)
				..controls(172.4,-37)and(173.3,-37.4)..(173.6,-37.4)
				..controls(175.6,-37.7)and(176.2,-38)..(176.7,-39.1)
				..controls(176.9,-39.6)and(176.9,-39.7)..(176.7,-40.4)
				..controls(176.4,-41.6)and(175.8,-42.7)..(175.2,-43.4)
				..controls(174.5,-44.1)and(172.6,-45.5)..(171.7,-45.8)
				..controls(170.6,-46.2)and(169.1,-46)..(168.2,-45.4)
				..controls(167.9,-45.2)and(167.3,-44.8)..(166.9,-44.4)
				..controls(166.4,-44.1)and(165.8,-43.6)..(165.4,-43.2)
				..controls(165.1,-42.9)and(164.4,-42.3)..(163.9,-41.9)
				..controls(163.3,-41.5)and(162.7,-40.9)..(162.4,-40.5)
				..controls(161.4,-39.2)and(159.9,-38)..(158.8,-37.8)
				..controls(158.4,-37.8)and(158.3,-37.7)..(158.2,-37.2)
				..controls(158,-36.8)and(157.4,-36.1)..(157.1,-36.1)
				..controls(157,-36.1)and(156.7,-36.1)..(156.5,-36)
				..controls(156.1,-35.9)and(155.9,-35.7)..(155.7,-35.2)
				..controls(155.3,-34.6)and(155.1,-34.4)..(154,-34.1)
				..controls(153.6,-34)and(153.2,-33.7)..(152.9,-33.5)
				..controls(152.4,-32.9)and(151.6,-32.4)..(151.3,-32.4)
				..controls(150.9,-32.4)and(150.4,-32.6)..(150.4,-32.9)
				..controls(150.4,-33.2)and(150.8,-33.4)..(151.4,-33.6)
				..controls(152.1,-33.7)and(152.8,-34.1)..(153.5,-34.9)
				..controls(154.1,-35.5)and(154.3,-35.6)..(155,-35.8)
				..controls(155.7,-36)and(155.7,-36)..(156,-36.6)
				..controls(156.3,-37.3)and(156.7,-37.7)..(157.5,-37.8)
				..controls(158.1,-37.9)and(158.2,-38)..(158.4,-38.8)
				..controls(158.5,-39.4)and(158.6,-39.6)..(159.6,-40.4)
				..controls(160.2,-41)and(160.9,-41.6)..(161.3,-41.8)
				..controls(161.7,-42)and(161.9,-42.3)..(162.2,-42.8)
				..controls(162.4,-43.2)and(162.8,-43.8)..(163.1,-44.1)
				..controls(163.9,-44.8)and(165.1,-46.7)..(165.7,-48)
				..controls(166.2,-49.3)and(166.4,-49.7)..(167.2,-50.5)
				..controls(167.5,-50.8)and(167.9,-51.4)..(168.1,-51.8)
				..controls(168.3,-52.2)and(168.8,-53)..(169.1,-53.6)
				..controls(169.6,-54.3)and(169.8,-54.9)..(169.9,-55.6)
				..controls(170.4,-57.8)and(170.5,-58.1)..(171.2,-59.5)
				..controls(171.6,-60.3)and(172,-61.5)..(172.2,-62.1)
				..controls(172.3,-62.8)and(172.7,-63.7)..(172.9,-64.1)
				..controls(173.2,-64.6)and(173.4,-65.3)..(173.4,-65.6)
				..controls(173.8,-68.8)and(174.2,-70.6)..(174.9,-72.6)
				..controls(175.5,-74.5)and(175.9,-77.1)..(176.1,-80.2)
				..controls(176.3,-83.4)and(175.5,-84.4)..(173.7,-83.4)--cycle;
				\definecolor{ezs}{RGB}{247,179,26}\fill[ezs](-5.7,-156.6)--(-3.4,-157.5)--(61.9,12.7)--(59.7,13.6)--cycle;
				\definecolor{ezs}{RGB}{247,179,22}\fill[ezs](-3.8,-157.3)--(-1.5,-158.2)--(63.9,12)--(61.6,12.9)--cycle;
				\definecolor{ezs}{RGB}{248,179,19}\fill[ezs](-1.8,-158.1)--(0.4,-158.9)--(65.8,11.2)--(63.5,12.1)--cycle;
				\definecolor{ezs}{RGB}{248,178,16}\fill[ezs](0.1,-158.8)--(2.4,-159.7)--(67.7,10.5)--(65.4,11.4)--cycle;
				\definecolor{ezs}{RGB}{249,178,15}\fill[ezs](2,-159.5)--(4.3,-160.4)--(69.6,9.8)--(67.3,10.7)--cycle;
				\definecolor{ezs}{RGB}{251,177,17}\fill[ezs](3.9,-160.3)--(6.2,-161.1)--(71.5,9)--(69.2,9.9)--cycle;
				\definecolor{ezs}{RGB}{252,179,18}\fill[ezs](5.8,-161)--(8.1,-161.9)--(73.4,8.3)--(71.1,9.2)--cycle;
				\definecolor{ezs}{RGB}{252,182,17}\fill[ezs](7.7,-161.7)--(10,-162.6)--(75.3,7.6)--(73,8.5)--cycle;
				\definecolor{ezs}{RGB}{252,185,17}\fill[ezs](9.6,-162.5)--(11.9,-163.4)--(77.2,6.8)--(75,7.7)--cycle;
				\definecolor{ezs}{RGB}{251,189,18}\fill[ezs](11.5,-163.2)--(13.8,-164.1)--(79.2,6.1)--(76.9,7)--cycle;
				\definecolor{ezs}{RGB}{251,191,22}\fill[ezs](13.4,-163.9)--(15.7,-164.8)--(81.1,5.4)--(78.8,6.2)--cycle;
				\definecolor{ezs}{RGB}{252,193,25}\fill[ezs](15.4,-164.7)--(17.6,-165.6)--(83,4.6)--(80.7,5.5)--cycle;
				\definecolor{ezs}{RGB}{253,195,27}\fill[ezs](17.3,-165.4)--(19.6,-166.3)--(84.9,3.9)--(82.6,4.8)--cycle;
				\definecolor{ezs}{RGB}{253,196,31}\fill[ezs](19.2,-166.1)--(21.5,-167)--(86.8,3.2)--(84.5,4)--cycle;
				\definecolor{ezs}{RGB}{254,199,35}\fill[ezs](21.1,-166.9)--(23.4,-167.8)--(88.7,2.4)--(86.4,3.3)--cycle;
				\definecolor{ezs}{RGB}{253,201,38}\fill[ezs](23,-167.6)--(25.3,-168.5)--(90.6,1.7)--(88.3,2.6)--cycle;
				\definecolor{ezs}{RGB}{253,203,40}\fill[ezs](24.9,-168.3)--(27.2,-169.2)--(92.5,1)--(90.3,1.8)--cycle;
				\definecolor{ezs}{RGB}{253,205,42}\fill[ezs](26.8,-169.1)--(29.1,-170)--(94.4,0.2)--(92.2,1.1)--cycle;
				\definecolor{ezs}{RGB}{253,207,45}\fill[ezs](28.7,-169.8)--(31,-170.7)--(96.4,-0.5)--(94.1,0.4)--cycle;
				\definecolor{ezs}{RGB}{253,210,49}\fill[ezs](30.7,-170.5)--(32.9,-171.4)--(98.3,-1.2)--(96,-0.4)--cycle;
				\definecolor{ezs}{RGB}{253,212,54}\fill[ezs](32.6,-171.3)--(34.9,-172.2)--(100.2,-2)--(97.9,-1.1)--cycle;
				\definecolor{ezs}{RGB}{253,214,59}\fill[ezs](34.5,-172)--(36.8,-172.9)--(102.1,-2.7)--(99.8,-1.8)--cycle;
				\definecolor{ezs}{RGB}{253,216,63}\fill[ezs](36.4,-172.7)--(38.7,-173.6)--(104,-3.4)--(101.7,-2.6)--cycle;
				\definecolor{ezs}{RGB}{253,218,65}\fill[ezs](38.3,-173.5)--(40.6,-174.4)--(105.9,-4.2)--(103.6,-3.3)--cycle;
				\definecolor{ezs}{RGB}{254,220,66}\fill[ezs](40.2,-174.2)--(42.5,-175.1)--(107.8,-4.9)--(105.5,-4)--cycle;
				\definecolor{ezs}{RGB}{254,222,68}\fill[ezs](42.1,-175)--(44.4,-175.8)--(109.7,-5.6)--(107.5,-4.8)--cycle;
				\definecolor{ezs}{RGB}{254,224,71}\fill[ezs](44,-175.7)--(46.3,-176.6)--(111.7,-6.4)--(109.4,-5.5)--cycle;
				\definecolor{ezs}{RGB}{254,226,75}\fill[ezs](46,-176.4)--(48.2,-177.3)--(113.6,-7.1)--(111.3,-6.2)--cycle;
				\definecolor{ezs}{RGB}{254,228,79}\fill[ezs](47.9,-177.2)--(50.1,-178)--(115.5,-7.8)--(113.2,-7)--cycle;
				\definecolor{ezs}{RGB}{254,230,82}\fill[ezs](49.8,-177.9)--(52.1,-178.8)--(117.4,-8.6)--(115.1,-7.7)--cycle;
				\definecolor{ezs}{RGB}{254,232,85}\fill[ezs](51.7,-178.6)--(54,-179.5)--(119.3,-9.3)--(117,-8.4)--cycle;
				\definecolor{ezs}{RGB}{254,234,88}\fill[ezs](53.6,-179.4)--(55.9,-180.2)--(121.2,-10)--(118.9,-9.2)--cycle;
				\definecolor{ezs}{RGB}{254,237,91}\fill[ezs](55.5,-180.1)--(57.8,-181)--(123.1,-10.8)--(120.8,-9.9)--cycle;
				\definecolor{ezs}{RGB}{254,239,94}\fill[ezs](57.4,-180.8)--(59.7,-181.7)--(125,-11.5)--(122.8,-10.6)--cycle;
				\definecolor{ezs}{RGB}{254,241,97}\fill[ezs](59.3,-181.6)--(61.6,-182.4)--(127,-12.2)--(124.7,-11.4)--cycle;
				\definecolor{ezs}{RGB}{254,243,101}\fill[ezs](61.2,-182.3)--(63.5,-183.2)--(128.9,-13)--(126.6,-12.1)--cycle;
				\definecolor{ezs}{RGB}{254,245,103}\fill[ezs](63.2,-183)--(65.4,-183.9)--(130.8,-13.7)--(128.5,-12.8)--cycle;
				\definecolor{ezs}{RGB}{254,245,105}\fill[ezs](65.1,-183.8)--(67.4,-184.6)--(132.7,-14.4)--(130.4,-13.6)--cycle;
				\definecolor{ezs}{RGB}{254,245,105}\fill[ezs](67,-184.5)--(69.3,-185.4)--(134.6,-15.2)--(132.3,-14.3)--cycle;
				\definecolor{ezs}{RGB}{254,244,103}\fill[ezs](68.9,-185.2)--(71.2,-186.1)--(136.5,-15.9)--(134.2,-15)--cycle;
				\definecolor{ezs}{RGB}{253,243,101}\fill[ezs](70.8,-186)--(73.1,-186.8)--(138.4,-16.6)--(136.1,-15.8)--cycle;
				\definecolor{ezs}{RGB}{253,241,98}\fill[ezs](72.7,-186.7)--(75,-187.6)--(140.3,-17.4)--(138,-16.5)--cycle;
				\definecolor{ezs}{RGB}{253,239,95}\fill[ezs](74.6,-187.4)--(76.9,-188.3)--(142.2,-18.1)--(140,-17.2)--cycle;
				\definecolor{ezs}{RGB}{253,236,92}\fill[ezs](76.5,-188.2)--(78.8,-189)--(144.2,-18.8)--(141.9,-18)--cycle;
				\definecolor{ezs}{RGB}{253,234,88}\fill[ezs](78.5,-188.9)--(80.7,-189.8)--(146.1,-19.6)--(143.8,-18.7)--cycle;
				\definecolor{ezs}{RGB}{253,232,85}\fill[ezs](80.4,-189.6)--(82.7,-190.5)--(148,-20.3)--(145.7,-19.4)--cycle;
				\definecolor{ezs}{RGB}{253,230,82}\fill[ezs](82.3,-190.4)--(84.6,-191.2)--(149.9,-21)--(147.6,-20.2)--cycle;
				\definecolor{ezs}{RGB}{253,228,78}\fill[ezs](84.2,-191.1)--(86.5,-192)--(151.8,-21.8)--(149.5,-20.9)--cycle;
				\definecolor{ezs}{RGB}{253,226,74}\fill[ezs](86.1,-191.8)--(88.4,-192.7)--(153.7,-22.5)--(151.4,-21.6)--cycle;
				\definecolor{ezs}{RGB}{253,224,70}\fill[ezs](88,-192.6)--(90.3,-193.4)--(155.6,-23.3)--(153.3,-22.4)--cycle;
				\definecolor{ezs}{RGB}{253,222,67}\fill[ezs](89.9,-193.3)--(92.2,-194.2)--(157.5,-24)--(155.3,-23.1)--cycle;
				\definecolor{ezs}{RGB}{253,220,65}\fill[ezs](91.8,-194)--(94.1,-194.9)--(159.5,-24.7)--(157.2,-23.8)--cycle;
				\definecolor{ezs}{RGB}{253,217,63}\fill[ezs](93.7,-194.8)--(96,-195.6)--(161.4,-25.5)--(159.1,-24.6)--cycle;
				\definecolor{ezs}{RGB}{253,215,60}\fill[ezs](95.7,-195.5)--(97.9,-196.4)--(163.3,-26.2)--(161,-25.3)--cycle;
				\definecolor{ezs}{RGB}{253,213,56}\fill[ezs](97.6,-196.2)--(99.9,-197.1)--(165.2,-26.9)--(162.9,-26)--cycle;
				\definecolor{ezs}{RGB}{252,211,53}\fill[ezs](99.5,-197)--(101.8,-197.8)--(167.1,-27.7)--(164.8,-26.8)--cycle;
				\definecolor{ezs}{RGB}{252,209,49}\fill[ezs](101.4,-197.7)--(103.7,-198.6)--(169,-28.4)--(166.7,-27.5)--cycle;
				\definecolor{ezs}{RGB}{252,207,45}\fill[ezs](103.3,-198.4)--(105.6,-199.3)--(170.9,-29.1)--(168.6,-28.2)--cycle;
				\definecolor{ezs}{RGB}{253,205,42}\fill[ezs](105.2,-199.2)--(107.5,-200)--(172.8,-29.9)--(170.6,-29)--cycle;
				\definecolor{ezs}{RGB}{253,202,40}\fill[ezs](107.1,-199.9)--(109.4,-200.8)--(174.7,-30.6)--(172.5,-29.7)--cycle;
				\definecolor{ezs}{RGB}{253,200,38}\fill[ezs](109,-200.6)--(111.3,-201.5)--(176.7,-31.3)--(174.4,-30.4)--cycle;
				\definecolor{ezs}{RGB}{253,198,35}\fill[ezs](111,-201.4)--(113.2,-202.2)--(178.6,-32.1)--(176.3,-31.2)--cycle;
				\definecolor{ezs}{RGB}{253,196,31}\fill[ezs](112.9,-202.1)--(115.2,-203)--(180.5,-32.8)--(178.2,-31.9)--cycle;
				\definecolor{ezs}{RGB}{253,194,28}\fill[ezs](114.8,-202.8)--(117.1,-203.7)--(182.4,-33.5)--(180.1,-32.6)--cycle;
				\definecolor{ezs}{RGB}{253,192,26}\fill[ezs](116.7,-203.6)--(119,-204.5)--(184.3,-34.3)--(182,-33.4)--cycle;
				\definecolor{ezs}{RGB}{253,189,24}\fill[ezs](118.6,-204.3)--(120.9,-205.2)--(186.2,-35)--(183.9,-34.1)--cycle;
				\definecolor{ezs}{RGB}{253,186,20}\fill[ezs](120.5,-205)--(122.8,-205.9)--(188.1,-35.7)--(185.8,-34.8)--cycle;
				\definecolor{ezs}{RGB}{253,184,17}\fill[ezs](122.4,-205.8)--(124.7,-206.7)--(190,-36.5)--(187.8,-35.6)--cycle;
				\definecolor{ezs}{RGB}{252,183,15}\fill[ezs](124.3,-206.5)--(126.6,-207.4)--(192,-37.2)--(189.7,-36.3)--cycle;
				\definecolor{ezs}{RGB}{252,181,14}\fill[ezs](126.3,-207.2)--(128.5,-208.1)--(193.9,-37.9)--(191.6,-37.1)--cycle;
				\definecolor{ezs}{RGB}{252,179,13}\fill[ezs](128.2,-208)--(130.4,-208.9)--(195.8,-38.7)--(193.5,-37.8)--cycle;
				\definecolor{ezs}{RGB}{252,178,12}\fill[ezs](130.1,-208.7)--(132.4,-209.6)--(197.7,-39.4)--(195.4,-38.5)--cycle;
				\definecolor{ezs}{RGB}{251,178,10}\fill[ezs](132,-209.4)--(134.3,-210.3)--(199.6,-40.1)--(197.3,-39.3)--cycle;
				\definecolor{ezs}{RGB}{251,178,10}\fill[ezs](133.9,-210.2)--(136.2,-211.1)--(201.5,-40.9)--(199.2,-40)--cycle;
				\definecolor{ezs}{RGB}{250,177,11}\fill[ezs](135.8,-210.9)--(138.1,-211.8)--(203.4,-41.6)--(201.1,-40.7)--cycle;
				\definecolor{ezs}{RGB}{250,178,11}\fill[ezs](137.7,-211.6)--(140,-212.5)--(205.3,-42.3)--(203.1,-41.5)--cycle;
				\definecolor{ezs}{RGB}{248,180,10}\fill[ezs](139.6,-212.4)--(141.9,-213.3)--(207.3,-43.1)--(205,-42.2)--cycle;
				\definecolor{ezs}{RGB}{247,181,9}\fill[ezs](141.5,-213.1)--(143.8,-214)--(209.2,-43.8)--(206.9,-42.9)--cycle;
				\definecolor{ezs}{RGB}{246,182,9}\fill[ezs](143.5,-213.8)--(145.7,-214.7)--(211.1,-44.5)--(208.8,-43.7)--cycle;
				\definecolor{ezs}{RGB}{246,182,10}\fill[ezs](145.4,-214.6)--(147.7,-215.5)--(213,-45.3)--(210.7,-44.4)--cycle;
			\end{scope}
			% chữ EDUZY (Times đậm, làm dày nét bằng 5 lớp lệch nhẹ; căn đúng khung chữ logo gốc)
			\foreach \dx/\dy in {-0.7/0,0.7/0,0/-0.3,0/0.3,0/0}
			\node[anchor=center,inner sep=0pt,outer sep=0pt,text=ezText,align=center] at (109.6+\dx,-193.6+\dy)
			{\resizebox{1.01cm}{0.22cm}{\fontencoding{T1}\fontfamily{ptm}\fontseries{b}\selectfont EDUZY}};
		\end{tikzpicture}%
	}%
}
%====================== THIẾT LẬP THÂN ĐỀ ======================
%--- Môi trường tự luận {bt}: ex_test 3.3.4 chưa khai báo sẵn nên tạo ở đây
%    (ex_test đã tự móc \fixshowans và \listenumerate cho bt)
\makeatletter
\@ifundefined{bt}{\newtheorem{bt}{\color{gold1}\namebt}}{}
\makeatother
\renewcommand{\sectionmark}[1]{\markright{\thesection\ #1}}
%--- Khung tiêu đề (logo + tối đa 4 dòng; dòng 3, 4 để trống {} thì bỏ qua)
%  Đề thi:    \khungtieude{THI THỬ THPT QUỐC GIA NĂM 2026}{Môn Toán -- Đề số 21}
%                         {(Đề thi gồm: ...)}{Thời gian làm bài: 90 phút (không kể thời gian phát đề).}
%  Chuyên đề: \khungtieude{CHUYÊN ĐỀ 3 -- DÃY SỐ VÀ CÔNG THỨC TRUY HỒI}
%                         {Bồi dưỡng học sinh giỏi Toán THPT}{...}{}
\newcommand{\khungtieude}[4]{%
	\begin{tcolorbox}[colframe=gold1,colback=gold5!45,boxrule=0.5mm]
		\begin{minipage}[c]{3cm}
			\centerline{\eduzylogo[2.4cm]}
		\end{minipage}\hspace{0.5cm}
		\begin{minipage}[c]{12.5cm}
			\vskip 0.15 cm
			\centerline{\ind{#1}}
			\vskip 0.15 cm
			\centerline{\textbf{#2}}
			\if\relax\detokenize{#3}\relax\else
				\vskip 0.15 cm
				\centerline{\textit{#3}}
			\fi
			\if\relax\detokenize{#4}\relax\else
				\vskip 0.15 cm
				\centerline{\textit{#4}}
			\fi
			\vskip 0.15 cm
		\end{minipage}
	\end{tcolorbox}
}
%====================== BỔ SUNG CHO BÀI HỌC, CHUYÊN ĐỀ ======================
\usepackage[shortlabels]{enumitem}
\usepackage{tabularx}
%--- Danh sách gọn: \begin{enumerate}[gon] ... \end{enumerate}
\SetEnumitemKey{gon}{leftmargin=*,itemsep=1pt,parsep=0pt,topsep=3pt}
%--- Khung lý thuyết: \begin{khung}{} ... \end{khung}; tiêu đề chỉ dùng khi thật ngắn ("Mẹo", "Lỗi thường gặp")
\newtcolorbox{khung}[1]{breakable,enhanced,colframe=gold1,colback=gold5!45,
	boxrule=0.4mm,fonttitle=\bfseries\sffamily,coltitle=white,colbacktitle=gold1,
	attach boxed title to top left={xshift=6mm,yshift=-3mm},
	boxed title style={arc=1mm,boxrule=0.3mm},title={#1}}
%--- Tiêu đề mục nhỏ trong một dạng: \muc{Dấu hiệu}, \muc{Phương pháp}, \muc{Bài tự luyện}
\makeatletter
\newcommand{\muc}[1]{\par\medskip\noindent\indam{#1}\par\nobreak\smallskip\@afterindentfalse\@afterheading}
\makeatother
%--- Ô công thức cần nhớ (căn giữa, viền vàng, vừa khít nội dung)
%     \begin{center}\congthuc{$Q(n)=(n-1)!$}\end{center}
\newtcbox{\congthuc}{on line,colframe=gold1,colback=gold5!45,boxrule=0.5pt,arc=2pt,
	left=8pt,right=8pt,top=4pt,bottom=4pt}
%--- Công tắc lời giải bài tự luyện (ví dụ mẫu luôn dùng \loigiai)
%     \loigiaiTL[{đáp số}]{lời giải đầy đủ}
%     \TLfull -> in lời giải đầy đủ;  \TLdapso -> chỉ in một dòng "Đáp số";  \TLan -> ẩn hoàn toàn
%     (\CheDo là dethi thì mọi lời giải và đáp số đều ẩn)
\makeatletter
\newif\ifTLhien   \TLhientrue
\newif\ifTLdapso  \TLdapsofalse
\newcommand{\TLfull}{\TLhientrue\TLdapsofalse}
\newcommand{\TLdapso}{\TLhientrue\TLdapsotrue}
\newcommand{\TLan}{\TLhienfalse}
\newcommand{\loigiaiTL}[2][]{%
	\ifTLhien
		\ifTLdapso
			\def\@tlarg{#1}%
			\ifx\@tlarg\@empty\else
				\ifdefstring{\kindOPTN}{dethi}{}{\par\noindent{\bfseries\color{gold1}Đáp số:} #1\par}%
			\fi
		\else
			\loigiai{#2}%
		\fi
	\fi}
\makeatother
\TLdapso% mặc định: chỉ in đáp số bài tự luyện
%====================== BẢNG ĐÚNG/SAI (\choiceTFt) ======================
%  Hàng tiêu đề nền nâu đồng chữ trắng (Phát biểu | Đ | S), các hàng tô xen kẽ kem / trắng,
%  đường kẻ vàng kim, nhãn a) b) c) d) màu nâu đồng. Bản có đáp án: cột được chọn có vòng tròn
%  nâu đồng với dấu tích trắng (câu đúng tích ở cột Đúng, câu sai tích ở cột Sai).
%  Tiêu đề ghi đủ "Đúng | Sai": \OPTN{viettat=0} sau \input (cột Đúng/Sai rộng ra, phát biểu hẹp lại).
%  Dùng bảng gốc của ex_test: \def\BangDSgoc{1} trước \input{preamble-eduzy-v2}.
\providecommand{\BangDSgoc}{0}
\ifnum\BangDSgoc=0
\usepackage{colortbl,pifont}
\colorlet{ezDSke}{gold2!60!gold1}% đường kẻ bảng (nét 0,4pt như ex_test để không đổi chiều cao bảng)
\colorlet{ezDSle}{gold5!60}% hàng a), c)
\colorlet{ezDSchan}{white}% hàng b), d)
%--- dấu chọn ở cột Đúng/Sai
%    (bề rộng 0, căn giữa ô: không làm đổi bề rộng cột của ex_test, không báo tràn lề)
\newcommand{\EZdauchon}{\makebox[0pt]{\tikz[baseline=-1ex]\node[circle,fill=gold1,inner sep=0pt,
	minimum size=1.2em,text=white,font=\small]{\ding{51}};}}
\def\tickT{\EZdauchon}\def\tickF{\EZdauchon}
\AtBeginDocument{\sbox{\ccTFTab}{\EZdauchon}}
%--- nhãn a) b) c) d) màu nâu đồng (ở bản lời giải, nhãn câu đúng vẫn đóng khung như ex_test)
\renewcommand{\FalseTF}{\stepcounter{dapan}\textbf{\color{gold1}\DapAnTF\sepTF}\ignorespaces}
%--- trước khi dựng bảng: kẻ màu vàng kim (bề rộng cột giữ nguyên như ex_test); sau bảng trả lại màu đen
\newcount\EZhangTF
\newcommand{\EZbatdauTF}{%
	\global\EZhangTF=0\relax
	\arrayrulecolor{ezDSke}}
\newcommand{\EZketthucTF}{\arrayrulecolor{black}}
%--- hàng tiêu đề nền nâu đồng, chữ trắng
\makeatletter
\newcommand{\EZo@TF}[1]{\centering\bfseries\color{white}#1}
\newcommand{\EZtieudeTF}{%
	\ifcase\cotDS
		\ifnum\cotSTT=1\relax
			\def\rowHeader{\hline\rowcolor{gold1}\EZo@TF{TT}&\EZo@TF{\phatbieu}\tabularnewline}%
		\else
			\def\rowHeader{\hline\rowcolor{gold1}\EZo@TF{\phatbieu}\tabularnewline}%
		\fi
	\or
		\ifnum\cotSTT=1\relax
			\def\rowHeader{\hline\rowcolor{gold1}\EZo@TF{TT}&\EZo@TF{\phatbieu}&\EZo@TF{\Dung}\tabularnewline}%
		\else
			\def\rowHeader{\hline\rowcolor{gold1}\EZo@TF{\phatbieu}&\EZo@TF{\Dung}\tabularnewline}%
		\fi
	\or
		\ifnum\cotSTT=1\relax
			\def\rowHeader{\hline\rowcolor{gold1}\EZo@TF{TT}&\EZo@TF{\phatbieu}&\EZo@TF{\Dung}&\EZo@TF{\Sai}\tabularnewline}%
		\else
			\def\rowHeader{\hline\rowcolor{gold1}\EZo@TF{\phatbieu}&\EZo@TF{\Dung}&\EZo@TF{\Sai}\tabularnewline}%
		\fi
	\fi}
%--- tô xen kẽ từng hàng (gọi ngay sau \hline đầu mỗi hàng)
\newcommand{\EZmauhangTF}{%
	\noalign{\global\advance\EZhangTF\@ne}%
	\ifodd\EZhangTF\rowcolor{ezDSle}\else\rowcolor{ezDSchan}\fi}
%--- gắn vào \choiceTFt / \writeChoiceTFt của ex_test
\xpretocmd{\choiceTFt}{\EZbatdauTF}{}{\typeout{[preamble-eduzy-v2] Không gắn được bảng Đúng/Sai (đầu bảng)}}
\xpatchcmd{\choiceTFt}{\ifnum\dongPB=0\def\rowHeader{}\fi}{\EZtieudeTF\ifnum\dongPB=0\def\rowHeader{}\fi}
	{}{\typeout{[preamble-eduzy-v2] Không gắn được bảng Đúng/Sai (hàng tiêu đề)}}
\xpatchcmd{\writeChoiceTFt}{\hline\xdef\dotEX{\sdotEX}}{\hline\EZmauhangTF\xdef\dotEX{\sdotEX}}
	{}{\typeout{[preamble-eduzy-v2] Không gắn được bảng Đúng/Sai (tô hàng)}}
\xpatchcmd{\writeChoiceTFt}{{\end{longtable}}}{{\end{longtable}\EZketthucTF}}
	{}{\typeout{[preamble-eduzy-v2] Không gắn được bảng Đúng/Sai (cuối longtable)}}
\xpatchcmd{\writeChoiceTFt}{{\end{xtabular}\end{center}\medskip\noindent}}{{\end{xtabular}\end{center}\EZketthucTF\medskip\noindent}}
	{}{\typeout{[preamble-eduzy-v2] Không gắn được bảng Đúng/Sai (cuối xtabular)}}
\fi
%=====================================================================
\renewcommand{\parallel}{\mathrel{/\!\!/}}
.
\providecommand{\BangDSgoc}{0}
\ifnum\BangDSgoc=0
\usepackage{colortbl,pifont}
\colorlet{ezDSke}{gold2!60!gold1}% đường kẻ bảng (nét 0,4pt như ex_test để không đổi chiều cao bảng)
\colorlet{ezDSle}{gold5!60}% hàng a), c)
\colorlet{ezDSchan}{white}% hàng b), d)
%--- dấu chọn ở cột Đúng/Sai
%    (bề rộng 0, căn giữa ô: không làm đổi bề rộng cột của ex_test, không báo tràn lề)
\newcommand{\EZdauchon}{\makebox[0pt]{\tikz[baseline=-1ex]\node[circle,fill=gold1,inner sep=0pt,
	minimum size=1.2em,text=white,font=\small]{\ding{51}};}}
\def\tickT{\EZdauchon}\def\tickF{\EZdauchon}
\AtBeginDocument{\sbox{\ccTFTab}{\EZdauchon}}
%--- nhãn a) b) c) d) màu nâu đồng (ở bản lời giải, nhãn câu đúng vẫn đóng khung như ex_test)
\renewcommand{\FalseTF}{\stepcounter{dapan}\textbf{\color{gold1}\DapAnTF\sepTF}\ignorespaces}
%--- trước khi dựng bảng: kẻ màu vàng kim (bề rộng cột giữ nguyên như ex_test); sau bảng trả lại màu đen
\newcount\EZhangTF
\newcommand{\EZbatdauTF}{%
	\global\EZhangTF=0\relax
	\arrayrulecolor{ezDSke}}
\newcommand{\EZketthucTF}{\arrayrulecolor{black}}
%--- hàng tiêu đề nền nâu đồng, chữ trắng
\makeatletter
\newcommand{\EZo@TF}[1]{\centering\bfseries\color{white}#1}
\newcommand{\EZtieudeTF}{%
	\ifcase\cotDS
		\ifnum\cotSTT=1\relax
			\def\rowHeader{\hline\rowcolor{gold1}\EZo@TF{TT}&\EZo@TF{\phatbieu}\tabularnewline}%
		\else
			\def\rowHeader{\hline\rowcolor{gold1}\EZo@TF{\phatbieu}\tabularnewline}%
		\fi
	\or
		\ifnum\cotSTT=1\relax
			\def\rowHeader{\hline\rowcolor{gold1}\EZo@TF{TT}&\EZo@TF{\phatbieu}&\EZo@TF{\Dung}\tabularnewline}%
		\else
			\def\rowHeader{\hline\rowcolor{gold1}\EZo@TF{\phatbieu}&\EZo@TF{\Dung}\tabularnewline}%
		\fi
	\or
		\ifnum\cotSTT=1\relax
			\def\rowHeader{\hline\rowcolor{gold1}\EZo@TF{TT}&\EZo@TF{\phatbieu}&\EZo@TF{\Dung}&\EZo@TF{\Sai}\tabularnewline}%
		\else
			\def\rowHeader{\hline\rowcolor{gold1}\EZo@TF{\phatbieu}&\EZo@TF{\Dung}&\EZo@TF{\Sai}\tabularnewline}%
		\fi
	\fi}
%--- tô xen kẽ từng hàng (gọi ngay sau \hline đầu mỗi hàng)
\newcommand{\EZmauhangTF}{%
	\noalign{\global\advance\EZhangTF\@ne}%
	\ifodd\EZhangTF\rowcolor{ezDSle}\else\rowcolor{ezDSchan}\fi}
%--- gắn vào \choiceTFt / \writeChoiceTFt của ex_test
\xpretocmd{\choiceTFt}{\EZbatdauTF}{}{\typeout{[preamble-eduzy-v2] Không gắn được bảng Đúng/Sai (đầu bảng)}}
\xpatchcmd{\choiceTFt}{\ifnum\dongPB=0\def\rowHeader{}\fi}{\EZtieudeTF\ifnum\dongPB=0\def\rowHeader{}\fi}
	{}{\typeout{[preamble-eduzy-v2] Không gắn được bảng Đúng/Sai (hàng tiêu đề)}}
\xpatchcmd{\writeChoiceTFt}{\hline\xdef\dotEX{\sdotEX}}{\hline\EZmauhangTF\xdef\dotEX{\sdotEX}}
	{}{\typeout{[preamble-eduzy-v2] Không gắn được bảng Đúng/Sai (tô hàng)}}
\xpatchcmd{\writeChoiceTFt}{{\end{longtable}}}{{\end{longtable}\EZketthucTF}}
	{}{\typeout{[preamble-eduzy-v2] Không gắn được bảng Đúng/Sai (cuối longtable)}}
\xpatchcmd{\writeChoiceTFt}{{\end{xtabular}\end{center}\medskip\noindent}}{{\end{xtabular}\end{center}\EZketthucTF\medskip\noindent}}
	{}{\typeout{[preamble-eduzy-v2] Không gắn được bảng Đúng/Sai (cuối xtabular)}}
\fi
%=====================================================================
\renewcommand{\parallel}{\mathrel{/\!\!/}}
.
\providecommand{\BangDSgoc}{0}
\ifnum\BangDSgoc=0
\usepackage{colortbl,pifont}
\colorlet{ezDSke}{gold2!60!gold1}% đường kẻ bảng (nét 0,4pt như ex_test để không đổi chiều cao bảng)
\colorlet{ezDSle}{gold5!60}% hàng a), c)
\colorlet{ezDSchan}{white}% hàng b), d)
%--- dấu chọn ở cột Đúng/Sai
%    (bề rộng 0, căn giữa ô: không làm đổi bề rộng cột của ex_test, không báo tràn lề)
\newcommand{\EZdauchon}{\makebox[0pt]{\tikz[baseline=-1ex]\node[circle,fill=gold1,inner sep=0pt,
	minimum size=1.2em,text=white,font=\small]{\ding{51}};}}
\def\tickT{\EZdauchon}\def\tickF{\EZdauchon}
\AtBeginDocument{\sbox{\ccTFTab}{\EZdauchon}}
%--- nhãn a) b) c) d) màu nâu đồng (ở bản lời giải, nhãn câu đúng vẫn đóng khung như ex_test)
\renewcommand{\FalseTF}{\stepcounter{dapan}\textbf{\color{gold1}\DapAnTF\sepTF}\ignorespaces}
%--- trước khi dựng bảng: kẻ màu vàng kim (bề rộng cột giữ nguyên như ex_test); sau bảng trả lại màu đen
\newcount\EZhangTF
\newcommand{\EZbatdauTF}{%
	\global\EZhangTF=0\relax
	\arrayrulecolor{ezDSke}}
\newcommand{\EZketthucTF}{\arrayrulecolor{black}}
%--- hàng tiêu đề nền nâu đồng, chữ trắng
\makeatletter
\newcommand{\EZo@TF}[1]{\centering\bfseries\color{white}#1}
\newcommand{\EZtieudeTF}{%
	\ifcase\cotDS
		\ifnum\cotSTT=1\relax
			\def\rowHeader{\hline\rowcolor{gold1}\EZo@TF{TT}&\EZo@TF{\phatbieu}\tabularnewline}%
		\else
			\def\rowHeader{\hline\rowcolor{gold1}\EZo@TF{\phatbieu}\tabularnewline}%
		\fi
	\or
		\ifnum\cotSTT=1\relax
			\def\rowHeader{\hline\rowcolor{gold1}\EZo@TF{TT}&\EZo@TF{\phatbieu}&\EZo@TF{\Dung}\tabularnewline}%
		\else
			\def\rowHeader{\hline\rowcolor{gold1}\EZo@TF{\phatbieu}&\EZo@TF{\Dung}\tabularnewline}%
		\fi
	\or
		\ifnum\cotSTT=1\relax
			\def\rowHeader{\hline\rowcolor{gold1}\EZo@TF{TT}&\EZo@TF{\phatbieu}&\EZo@TF{\Dung}&\EZo@TF{\Sai}\tabularnewline}%
		\else
			\def\rowHeader{\hline\rowcolor{gold1}\EZo@TF{\phatbieu}&\EZo@TF{\Dung}&\EZo@TF{\Sai}\tabularnewline}%
		\fi
	\fi}
%--- tô xen kẽ từng hàng (gọi ngay sau \hline đầu mỗi hàng)
\newcommand{\EZmauhangTF}{%
	\noalign{\global\advance\EZhangTF\@ne}%
	\ifodd\EZhangTF\rowcolor{ezDSle}\else\rowcolor{ezDSchan}\fi}
%--- gắn vào \choiceTFt / \writeChoiceTFt của ex_test
\xpretocmd{\choiceTFt}{\EZbatdauTF}{}{\typeout{[preamble-eduzy-v2] Không gắn được bảng Đúng/Sai (đầu bảng)}}
\xpatchcmd{\choiceTFt}{\ifnum\dongPB=0\def\rowHeader{}\fi}{\EZtieudeTF\ifnum\dongPB=0\def\rowHeader{}\fi}
	{}{\typeout{[preamble-eduzy-v2] Không gắn được bảng Đúng/Sai (hàng tiêu đề)}}
\xpatchcmd{\writeChoiceTFt}{\hline\xdef\dotEX{\sdotEX}}{\hline\EZmauhangTF\xdef\dotEX{\sdotEX}}
	{}{\typeout{[preamble-eduzy-v2] Không gắn được bảng Đúng/Sai (tô hàng)}}
\xpatchcmd{\writeChoiceTFt}{{\end{longtable}}}{{\end{longtable}\EZketthucTF}}
	{}{\typeout{[preamble-eduzy-v2] Không gắn được bảng Đúng/Sai (cuối longtable)}}
\xpatchcmd{\writeChoiceTFt}{{\end{xtabular}\end{center}\medskip\noindent}}{{\end{xtabular}\end{center}\EZketthucTF\medskip\noindent}}
	{}{\typeout{[preamble-eduzy-v2] Không gắn được bảng Đúng/Sai (cuối xtabular)}}
\fi
%=====================================================================
\renewcommand{\parallel}{\mathrel{/\!\!/}}

%=====================================================================
% preamble-vatli-hinh.tex -- BỘ MACRO HÌNH VẬT LÍ (dùng chung cho mọi đề Vật lí)
% Gọi ngay sau preamble chính:
%     %============================================================
%  PREAMBLE EDUZY v2 - dùng chung cho đề thi, bài học, chuyên đề (gói ex_test)
%  Khác bản preamble-eduzy.tex ở đầu trang, chân trang: dải nền nâu đồng tràn ngang,
%  logo + Eduzy.com.vn, ô chữ vàng kem cạnh xiên, chân trang "Trang (số)";
%  logo Eduzy chìm giữa mọi trang (tắt: \logochimtat, bật lại: \logochimbat).
%  Cách dùng ở file chính (xem mau-chuan.tex, mau-chuyen-de.tex):
%      \def\CheDo{loigiai}          % loigiai | dethi | solcolor | color | book
%      \def\tenchuong{<dòng chữ trên đầu trang>}
%      \def\thuonghieu{<chữ ở chân trang lẻ>}   % không khai báo thì dùng mặc định
%      \def\tacgia{<chữ ở chân trang chẵn>}      % không khai báo thì dùng mặc định
%      \def\khauhieu{<chữ cạnh logo ở đầu trang>}  % tùy chọn, mặc định Eduzy.com.vn
%      %============================================================
%  PREAMBLE EDUZY v2 - dùng chung cho đề thi, bài học, chuyên đề (gói ex_test)
%  Khác bản preamble-eduzy.tex ở đầu trang, chân trang: dải nền nâu đồng tràn ngang,
%  logo + Eduzy.com.vn, ô chữ vàng kem cạnh xiên, chân trang "Trang (số)";
%  logo Eduzy chìm giữa mọi trang (tắt: \logochimtat, bật lại: \logochimbat).
%  Cách dùng ở file chính (xem mau-chuan.tex, mau-chuyen-de.tex):
%      \def\CheDo{loigiai}          % loigiai | dethi | solcolor | color | book
%      \def\tenchuong{<dòng chữ trên đầu trang>}
%      \def\thuonghieu{<chữ ở chân trang lẻ>}   % không khai báo thì dùng mặc định
%      \def\tacgia{<chữ ở chân trang chẵn>}      % không khai báo thì dùng mặc định
%      \def\khauhieu{<chữ cạnh logo ở đầu trang>}  % tùy chọn, mặc định Eduzy.com.vn
%      \input{preamble-eduzy-v2}
%      \TLdapso                      % tài liệu: \TLfull | \TLdapso | \TLan
%      \begin{document} ... \end{document}
%  Logo Eduzy vẽ bằng TikZ ngay trong file (\eduzylogo), không cần file ảnh.
%  Biên dịch: pdfLaTeX (máy có gói vietnam) hoặc XeLaTeX.
%============================================================
\documentclass[12pt,a4paper,twoside]{article}%<-- twoside: bat dau trang/chan trang chan-le
\usepackage{iftex}
\usepackage{amsmath,amssymb,tikz,mathrsfs,tkz-tab}
\usepackage{graphicx}
\usepackage{pgfplots}% đồ thị hàm số
\pgfplotsset{compat=1.17}
%\usepackage{mathptmx}
\usepackage{mathpazo}
\ifXeTeX% XeLaTeX: font chữ có đủ dấu tiếng Việt, giữ nguyên font toán mathpazo
	\usepackage[no-math]{fontspec}
	\setmainfont{TeX Gyre Pagella}
	\newfontfamily\EZqag{TeX Gyre Adventor}[NFSSFamily=qag,Ligatures=TeX]
\fi
\usepackage[many]{tcolorbox}
\usetikzlibrary{math,through,calc,intersections,angles,quotes,shapes,shapes.geometric,arrows,patterns,matrix,chains,arrows.meta,decorations.shapes,decorations.fractals,decorations.markings,shadows,decorations.pathreplacing}
\usepackage[top=1.6cm, bottom=1.8cm, left=1.8cm, right=1.8cm]{geometry}% v2: top 1.6cm (bản cũ 1.5cm) để chữ không sát dải đầu trang
\providecommand{\CheDo}{loigiai}% loigiai | dethi | solcolor | color | book
\expandafter\usepackage\expandafter[\CheDo]{ex_test}
% \tenchuong: khai báo ở file chính (trước \input) hoặc \renewcommand sau \input
%============================================================
%  [GỘP TỪ eduzy-de.tex] EDUZY - MÀU NÂU ĐỒNG / VÀNG KIM
%  + ĐẦU TRANG, CHÂN TRANG KIỂU DẢI NGANG (v2), ĐỐI XỨNG TRANG CHẴN / TRANG LẺ.
%  (không cần file eduzy-de.tex riêng nữa)
%  Khối này phải đặt SAU \usepackage[...]{ex_test} và SAU \def\tenchuong{...}
%============================================================
\usepackage{lastpage}
\usepackage{needspace}
%--- DÀN TRANG ĐỀ THI: tách câu (THÊM 05/11/2026; chỉ thêm lệnh mới, KHÔNG đổi hành vi mặc định) ---
%  Hai công tắc, mặc định tắt (= preamble chạy y như cũ). Bật bằng cách đặt TRƯỚC \input{preamble-eduzy-v2}
%  trong file đề (tools/dan_trang/dan_trang.py tự ghi khi đặt \NeedspaceCau):
%    \def\TachCauDS{1}   câu Đúng/Sai được tách: câu dẫn (kèm hình) ở cuối trang, bảng Phát biểu sang trang sau
%                        (câu dẫn và hình không tách nhau, bảng không tách giữa hai trang)
%    \def\TachCauLG{1}   bản lời giải (\CheDo khác dethi): phần đề của câu cao hơn 1/3 trang được phép tách
%  Lệnh dùng trong file đề (giá trị do tools/dan_trang/dan_trang.py đo ở bản dethi):
%    \NeedspaceCau[<cao câu dẫn + hình>]{<cao cả khối câu>}   ngay trên \begin{ex} (câu đầu phần: trên tiêu đề phần)
%        công tắc tắt: đúng bằng \Needspace{<cao cả khối câu>} như cũ;
%        \TachCauDS=1 và có [..] (câu Đúng/Sai được chọn tách): chỉ giữ câu dẫn + hình;
%        \TachCauLG=1 ở bản lời giải: đòi tối đa 1/3 \textheight.
%    \NeedspaceBang{<cao bảng Phát biểu>}%   ngay trước \choiceTFt: \TachCauDS=1 thì giữ trọn bảng trên một trang;
%        công tắc tắt: không làm gì.
\providecommand{\TachCauDS}{0}
\providecommand{\TachCauLG}{0}
\newlength{\EZnsCau}
\newcommand{\NeedspaceCau}[2][]{%
	\setlength{\EZnsCau}{#2}%
	\ifnum\TachCauDS=1 \if\relax\detokenize{#1}\relax\else\setlength{\EZnsCau}{#1}\fi\fi
	\ifnum\TachCauLG=1 \ifdefstring{\kindOPTN}{dethi}{}{\ifdim\EZnsCau>0.3333\textheight\setlength{\EZnsCau}{0.3333\textheight}\fi}\fi
	\Needspace{\EZnsCau}}
\newcommand{\NeedspaceBang}[1]{\ifnum\TachCauDS=1 \Needspace{#1}\fi}

%--- Thang màu: đậm -> nhạt
\definecolor{gold1}{RGB}{126,96,0}    % nâu đồng đậm  #7E6000
\definecolor{gold2}{RGB}{192,144,3}   % vàng kim      #C09003
\definecolor{gold3}{RGB}{255,215,99}  % vàng nhạt     #FFD763
\definecolor{gold4}{RGB}{255,231,153} % vàng rất nhạt #FFE799
\definecolor{gold5}{RGB}{255,243,205} % vàng kem      #FFF3CD
%--- Ghi đè tông xanh/đỏ cũ
\definecolor{myblue}{RGB}{126,96,0}
\colorlet{blue}{gold1}
\colorlet{cyan}{gold2}
\colorlet{teal}{gold2}

%--- Tên chương mặc định (file đề nên tự khai báo lại)
\providecommand{\tenchuong}{}
%--- Chân trang
\providecommand{\thuonghieu}{TOÁN 12 -- Eduzy.com.vn}
\providecommand{\tacgia}{Tô Gia Bảo -- 0353703783}

%--- Chữ cạnh logo ở đầu trang
\providecommand{\khauhieu}{Eduzy.com.vn}

%====================== THÔNG SỐ ĐẦU TRANG, CHÂN TRANG (v2) ======================
%  Dải nền đậm tràn ngang trang; ô chữ vàng kem có cạnh xiên, kèm 2 vạch xiên chuyển màu.
%  Đầu trang lẻ : logo + \khauhieu (trái)   | ô chữ \tenchuong (phải)
%  Chân trang lẻ: ô chữ \thuonghieu (trái)  | Trang (số) (phải)
%  Trang chẵn soi gương trang lẻ; chân trang chẵn ghi \tacgia.
%--- Màu: chỉ dùng thang vàng kim Eduzy
\colorlet{ezNen}{gold1}% dải nền đậm
\colorlet{ezVach1}{gold3}% vạch xiên sát dải nền (đậm hơn)
\colorlet{ezVach2}{gold4}% vạch xiên sát ô chữ (nhạt hơn)
\colorlet{ezO}{gold5}% ô chữ vàng kem
%--- Kích thước
\newlength{\EZdai}\setlength{\EZdai}{1.1cm}% chiều cao dải nền
\newlength{\EZo}\setlength{\EZo}{0.68cm}% chiều cao ô chữ
\newlength{\EZxien}\setlength{\EZxien}{0.55cm}% độ lệch ngang của cạnh xiên
\newlength{\EZvach}\setlength{\EZvach}{0.75cm}% bề rộng mỗi vạch xiên
\newlength{\EZkhe}\setlength{\EZkhe}{0.12cm}% khe trắng giữa vạch và ô chữ
\newlength{\EZdem}\setlength{\EZdem}{0.4cm}% khoảng đệm từ chữ tới cạnh xiên
\newlength{\EZle}\setlength{\EZle}{1.8cm}% lề trái/phải của chữ (bằng lề thân bài)
\newlength{\EZrongmin}\setlength{\EZrongmin}{4.5cm}% ô chữ rộng tối thiểu
\newlength{\EZrong}% bề rộng chữ trong ô (tính tự động)
%--- Font
\newcommand{\EZfontO}{\fontsize{10.5pt}{12pt}\selectfont\bfseries}% chữ trong ô
\newcommand{\EZfontKH}{\fontfamily{qag}\fontsize{11pt}{12pt}\selectfont\bfseries}% \khauhieu cạnh logo
\newcommand{\EZfontTrang}{\fontsize{11.5pt}{12pt}\selectfont\bfseries}% chữ "Trang" và số trang
%--- Đo bề rộng chữ trong ô
\newcommand{\EZdorong}[1]{%
	\settowidth{\EZrong}{\EZfontO #1}%
	\ifdim\EZrong<\EZrongmin\setlength{\EZrong}{\EZrongmin}\fi}
%--- Logo + chữ trắng \khauhieu (dựng sẵn một lần khi vào \begin{document})
\newsavebox{\EZlogobox}
\AtBeginDocument{\sbox{\EZlogobox}{%
		\raisebox{-0.5\height}{\eduzylogo[1cm]}\hspace{0.15cm}%
		\raisebox{-0.5\height}{\color{white}\EZfontKH\khauhieu}}}
%--- Hai vạch xiên + khe trắng. #1: điểm gốc ở mép dưới, #2: hướng (1 sang phải, -1 sang trái),
%    #3: hướng nghiêng (1 kiểu "/", -1 kiểu "\")
\newcommand{\EZvachxien}[3]{%
	\fill[white] (#1) --++({(#3)*\EZxien},\EZo) --++({(#2)*\EZkhe},0) --++({-(#3)*\EZxien},-\EZo) --cycle;
	\fill[ezVach2] ([xshift=(#2)*\EZkhe]#1) --++({(#3)*\EZxien},\EZo) --++({(#2)*\EZvach},0) --++({-(#3)*\EZxien},-\EZo) --cycle;
	\fill[ezVach1] ([xshift=(#2)*(\EZkhe+\EZvach)]#1) --++({(#3)*\EZxien},\EZo) --++({(#2)*\EZvach},0) --++({-(#3)*\EZxien},-\EZo) --cycle;
	\draw[white,line width=0.6pt] ([xshift=(#2)*(\EZkhe+\EZvach)]#1) --++({(#3)*\EZxien},\EZo);
}
%--- Chữ trong ô (căn giữa theo chiều cao ô, không lệch vì dấu tiếng Việt)
\tikzset{ezchu/.style={inner sep=0pt,text=ezNen,text height=0.72em,text depth=0.22em}}

%====================== LOGO CHÌM GIỮA TRANG ======================
%  Vẽ cùng đầu trang nên nằm dưới chữ của thân bài, có trên mọi trang.
%  Tắt: \logochimtat (từ trang đó trở đi), bật lại: \logochimbat.
%  Đổi cỡ, độ đậm: khai báo trước \input, ví dụ \def\EZchimcao{9cm} \def\EZchimdo{0.1}
\providecommand{\EZchimcao}{10cm}% chiều cao logo chìm
\providecommand{\EZchimdo}{0.12}% độ đậm (0 = trong suốt, 1 = rõ hẳn)
\newif\ifEZchim \EZchimtrue
\newcommand{\logochimtat}{\EZchimfalse}
\newcommand{\logochimbat}{\EZchimtrue}
%  Dựng sẵn logo mờ một lần khi vào \begin{document}; transparency group làm mờ cả logo
%  một lượt nên không lộ vệt ở chỗ các dải màu chồng nhau.
\newsavebox{\EZchimtho}
\newsavebox{\EZchimbox}
\AtBeginDocument{%
	\sbox{\EZchimtho}{\eduzylogo[\EZchimcao]}%
	\sbox{\EZchimbox}{\begin{tikzpicture}
			\begin{scope}[transparency group,opacity=\EZchimdo]
				\node[inner sep=0pt] {\usebox{\EZchimtho}};
			\end{scope}
	\end{tikzpicture}}}
\newcommand{\EZvechim}{%
	\ifEZchim
		\node[inner sep=0pt] at (current page.center) {\usebox{\EZchimbox}};
	\fi}

%====================== ĐẦU TRANG - TRANG LẺ ======================
%  (logo bên trái, ô chữ \tenchuong bên phải)
\newcommand{\myfancyheadOdd}{%
	\EZdorong{\tenchuong}%
	\begin{tikzpicture}[remember picture,overlay]
		\coordinate (T) at (current page.north west);
		\coordinate (P) at (current page.north east);
		%--- logo chìm giữa trang
		\EZvechim
		%--- dải nền
		\fill[ezNen] (T) rectangle ([yshift=-\EZdai]P);
		%--- ô chữ: X là góc dưới bên trái, cạnh trái xiên kiểu "/"
		\coordinate (X) at ([xshift=-\EZle-\EZrong-\EZdem-\EZxien,yshift=-\EZdai]P);
		\fill[ezO] (X) --++(\EZxien,\EZo) -- ([yshift=-\EZdai+\EZo]P) -- ([yshift=-\EZdai]P) --cycle;
		\EZvachxien{X}{-1}{1}
		\node[ezchu,anchor=east] at ([xshift=-\EZle,yshift=-\EZdai+0.5*\EZo]P) {\EZfontO\tenchuong};
		%--- logo + tên
		\node[inner sep=0pt,anchor=west] at ([xshift=\EZle,yshift=-0.5*\EZdai]T) {\usebox{\EZlogobox}};
	\end{tikzpicture}%
}

%====================== ĐẦU TRANG - TRANG CHẴN =====================
%  (soi gương: ô chữ \tenchuong bên trái, logo bên phải)
\newcommand{\myfancyheadEven}{%
	\EZdorong{\tenchuong}%
	\begin{tikzpicture}[remember picture,overlay]
		\coordinate (T) at (current page.north west);
		\coordinate (P) at (current page.north east);
		%--- logo chìm giữa trang
		\EZvechim
		%--- dải nền
		\fill[ezNen] (T) rectangle ([yshift=-\EZdai]P);
		%--- ô chữ: X là góc dưới bên phải, cạnh phải xiên kiểu "\"
		\coordinate (X) at ([xshift=\EZle+\EZrong+\EZdem+\EZxien,yshift=-\EZdai]T);
		\fill[ezO] (X) --++(-\EZxien,\EZo) -- ([yshift=-\EZdai+\EZo]T) -- ([yshift=-\EZdai]T) --cycle;
		\EZvachxien{X}{1}{-1}
		\node[ezchu,anchor=west] at ([xshift=\EZle,yshift=-\EZdai+0.5*\EZo]T) {\EZfontO\tenchuong};
		%--- logo + tên
		\node[inner sep=0pt,anchor=east] at ([xshift=-\EZle,yshift=-0.5*\EZdai]P) {\usebox{\EZlogobox}};
	\end{tikzpicture}%
}

%====================== CHÂN TRANG - TRANG LẺ ======================
%  (ô chữ \thuonghieu bên trái, "Trang (số)" bên phải)
\newcommand{\myfancyfootOdd}{%
	\EZdorong{\thuonghieu}%
	\begin{tikzpicture}[remember picture,overlay]
		\coordinate (T) at (current page.south west);
		\coordinate (P) at (current page.south east);
		%--- dải nền
		\fill[ezNen] (T) rectangle ([yshift=\EZdai]P);
		%--- ô chữ sát mép trên dải: X là góc dưới bên phải, cạnh phải xiên kiểu "/"
		\coordinate (X) at ([xshift=\EZle+\EZrong+\EZdem,yshift=\EZdai-\EZo]T);
		\fill[ezO] ([yshift=\EZdai-\EZo]T) -- (X) --++(\EZxien,\EZo) -- ([yshift=\EZdai]T) --cycle;
		\EZvachxien{X}{1}{1}
		\node[ezchu,anchor=west] at ([xshift=\EZle,yshift=\EZdai-0.5*\EZo]T) {\EZfontO\thuonghieu};
		%--- số trang
		\node[circle,draw=white,line width=0.9pt,text=white,inner sep=0pt,minimum size=0.62cm,anchor=east]
		(so) at ([xshift=-\EZle,yshift=0.5*\EZdai]P) {\EZfontTrang\thepage};
		\node[inner sep=0pt,text=white,anchor=base east] (tr) at ([xshift=-0.15cm]so.base west) {\EZfontTrang Trang};
	\end{tikzpicture}%
}

%====================== CHÂN TRANG - TRANG CHẴN ====================
%  (soi gương: "Trang (số)" bên trái, ô chữ \tacgia bên phải)
\newcommand{\myfancyfootEven}{%
	\EZdorong{\tacgia}%
	\begin{tikzpicture}[remember picture,overlay]
		\coordinate (T) at (current page.south west);
		\coordinate (P) at (current page.south east);
		%--- dải nền
		\fill[ezNen] (T) rectangle ([yshift=\EZdai]P);
		%--- ô chữ sát mép trên dải: X là góc dưới bên trái, cạnh trái xiên kiểu "\"
		\coordinate (X) at ([xshift=-\EZle-\EZrong-\EZdem,yshift=\EZdai-\EZo]P);
		\fill[ezO] ([yshift=\EZdai-\EZo]P) -- (X) --++(-\EZxien,\EZo) -- ([yshift=\EZdai]P) --cycle;
		\EZvachxien{X}{-1}{-1}
		\node[ezchu,anchor=east] at ([xshift=-\EZle,yshift=\EZdai-0.5*\EZo]P) {\EZfontO\tacgia};
		%--- số trang
		%    (vòng tròn ẩn so0 nằm giữa dải để lấy đường chân chữ, giống trang lẻ)
		\node[circle,inner sep=0pt,minimum size=0.62cm] (so0) at ([yshift=0.5*\EZdai]T) {\EZfontTrang\phantom{\thepage}};
		\node[inner sep=0pt,text=white,anchor=base west] (tr) at ($(so0.base -| T)+(\EZle,0)$) {\EZfontTrang Trang};
		\node[circle,draw=white,line width=0.9pt,text=white,inner sep=0pt,minimum size=0.62cm,anchor=base west]
		(so) at ([xshift=0.15cm]tr.base east) {\EZfontTrang\thepage};
	\end{tikzpicture}%
}

%--- Giữ tương thích với các file cũ đang gọi \myfancyhead, \myfancyfoot
\newcommand{\myfancyhead}{\myfancyheadOdd}
\newcommand{\myfancyfoot}{\myfancyfootOdd}

%======================= GẮN VÀO TRANG =============================
\pagestyle{fancy}
\fancyhf{}
\fancyhead[LO]{\myfancyheadOdd}
\fancyhead[RE]{\myfancyheadEven}
\fancyfoot[LO]{\myfancyfootOdd}
\fancyfoot[RE]{\myfancyfootEven}
\renewcommand{\headrulewidth}{0pt}
\renewcommand{\footrulewidth}{0pt}
\fancypagestyle{plain}{%
	\fancyhf{}
	\fancyhead[LO]{\myfancyheadOdd}
	\fancyhead[RE]{\myfancyheadEven}
	\fancyfoot[LO]{\myfancyfootOdd}
	\fancyfoot[RE]{\myfancyfootEven}
	\renewcommand{\headrulewidth}{0pt}
	\renewcommand{\footrulewidth}{0pt}
}
%=================== HẾT PHẦN GỘP TỪ eduzy-de.tex ===================
\newcommand{\indam}[1]{
	{\fontsize{12pt}{9pt}\sffamily \selectfont\bfseries\color{gold1} #1}
}
\newcommand{\ind}[1]{
	{\fontsize{14pt}{9pt}\fontfamily{qag} \selectfont\bfseries\color{gold1} #1}
}
\renewcommand{\FalseEX}{\stepcounter{dapan}\textbf{\selectfont\color{gold1}\Alph{dapan}.}\ignorespaces}
\renewtheorem{ex}{\color{gold1}Câu}
\renewcommand{\baselinestretch}{1.1}
\newcommand{\hoac}[1]{
	\left[\begin{aligned}#1\end{aligned}\right.}
\newcommand{\heva}[1]{
	\left\{\begin{aligned}#1\end{aligned}\right.}
\usepackage{esvect}
\def\vec{\vv}
\def\vt{\vv}
\def\overrightarrow{\vv}
%====================== LOGO EDUZY VẼ BẰNG TIKZ ======================
%  Đường viền dò từ logo gốc; dải màu vàng kim đo trực tiếp trên ảnh gốc
%  (nghiêng 21 độ, tô bằng 80 dải hẹp nên hiển thị giống nhau trên mọi trình xem PDF).
%  Dùng: \eduzylogo          -> cao 2.4cm (mặc định)
%        \eduzylogo[1.5cm]   -> cao 1.5cm
\definecolor{ezText}{RGB}{245,200,46}
\newcommand{\eduzylogo}[1][2.4cm]{%
	\resizebox{!}{#1}{%
		\begin{tikzpicture}[x=0.01cm,y=0.01cm]
			\path[use as bounding box] (0,0) rectangle (224,-224);
			% mũ tốt nghiệp + vương miện
			\begin{scope}
				\clip
				(119.4,-176.1)..controls(118.9,-176.1)and(116.8,-176)..(114.6,-176)
				..controls(101.4,-176)and(85.8,-175.9)..(83.7,-175.7)
				..controls(81.8,-175.6)and(81,-175.6)..(80.2,-175.7)
				..controls(78.8,-175.9)and(77.7,-175.9)..(76.2,-175.7)
				..controls(74.9,-175.5)and(70.8,-175.5)..(64.3,-175.7)
				..controls(59.1,-175.9)and(58.3,-175.7)..(58.6,-173.7)
				..controls(58.6,-172.8)and(59.6,-169.6)..(60.3,-168)
				..controls(60.8,-166.9)and(61.4,-165.8)..(62.3,-165)
				..controls(63.1,-164.2)and(64.3,-163.4)..(65,-163.2)
				..controls(65.5,-163.1)and(72.3,-163)..(76.6,-163.1)
				..controls(77.9,-163.1)and(80.1,-163.1)..(81.4,-163)
				..controls(82.8,-162.9)and(86.3,-162.9)..(89.2,-162.9)
				..controls(93,-162.9)and(94.9,-162.8)..(96,-162.7)
				..controls(97.2,-162.5)and(101,-162.5)..(114,-162.4)
				..controls(123,-162.3)and(134.9,-162.2)..(140.2,-162.1)
				..controls(150.1,-162)and(152.2,-162.1)..(153.6,-162.4)
				..controls(154.4,-162.6)and(157.1,-163.8)..(158.1,-164.4)
				..controls(158.8,-164.8)and(159.7,-165.7)..(159.9,-166.3)
				..controls(160,-166.5)and(160.2,-167.1)..(160.2,-167.6)
				..controls(160.3,-168)and(160.6,-169.4)..(160.9,-170.6)
				..controls(161.2,-171.7)and(161.5,-173.1)..(161.6,-173.6)
				..controls(161.8,-174.4)and(161.8,-174.5)..(161.5,-174.9)
				..controls(161.2,-175.6)and(160.8,-175.9)..(159.9,-176)
				..controls(159.2,-176.1)and(154.1,-176.2)..(151.5,-176.2)
				..controls(148.1,-176.1)and(122.6,-176.1)..(121.6,-176.2)
				..controls(120.8,-176.2)and(119.8,-176.2)..(119.4,-176.1)--cycle
				(130.4,-160)..controls(129,-160)and(119.8,-160)..(110.1,-159.9)
				..controls(100.3,-159.8)and(88.5,-159.7)..(83.9,-159.6)
				..controls(77.8,-159.5)and(75.3,-159.5)..(74.5,-159.6)
				..controls(73.9,-159.7)and(72.5,-159.8)..(71.3,-159.8)
				..controls(69.6,-159.8)and(69,-159.8)..(68.5,-159.6)
				..controls(67.6,-159.3)and(66.8,-158.6)..(66.2,-157.6)
				..controls(66,-157.3)and(65.6,-156.5)..(65.3,-156.1)
				..controls(64.6,-155.1)and(64.5,-154.9)..(63.7,-152.9)
				..controls(63.2,-151.6)and(62.9,-150.9)..(61.8,-149.6)
				..controls(61.5,-149.2)and(61.1,-148.5)..(60.9,-147.9)
				..controls(60.7,-147.4)and(60.3,-146.6)..(60.1,-146.1)
				..controls(59.9,-145.5)and(59.5,-144.9)..(59,-144.3)
				..controls(58.5,-143.7)and(58.2,-143.2)..(57.9,-142.5)
				..controls(57.7,-141.9)and(57.3,-141.1)..(57,-140.7)
				..controls(56.6,-140.3)and(56.1,-139.3)..(55.8,-138.5)
				..controls(55.5,-137.8)and(55.1,-136.8)..(54.9,-136.5)
				..controls(54.7,-136.1)and(54.3,-135.2)..(53.9,-134.5)
				..controls(53.6,-133.8)and(53,-132.7)..(52.6,-132)
				..controls(51.6,-130.5)and(51.4,-130.2)..(51.1,-129)
				..controls(50.6,-127.5)and(50.6,-127.4)..(49.7,-126.3)
				..controls(49.2,-125.7)and(48.6,-124.8)..(48.3,-124.2)
				..controls(48,-123.6)and(47.5,-122.7)..(47.3,-122.2)
				..controls(47,-121.8)and(46.5,-120.7)..(46.1,-119.8)
				..controls(45.8,-118.9)and(45.2,-117.8)..(44.9,-117.2)
				..controls(44.2,-116.1)and(43.3,-114.4)..(42.3,-112.2)
				..controls(41.8,-111.3)and(41.3,-110.4)..(41.1,-110.1)
				..controls(40.5,-109.4)and(40.3,-109)..(40,-107.8)
				..controls(39.8,-107)and(39.6,-106.5)..(39.2,-106)
				..controls(38.9,-105.6)and(38.4,-104.7)..(38,-104)
				..controls(37.7,-103.4)and(37.2,-102.5)..(37,-102.1)
				..controls(36.7,-101.7)and(36.3,-100.8)..(36,-100.1)
				..controls(35.6,-98.9)and(35.3,-98.4)..(34.6,-97.6)
				..controls(34.2,-97.1)and(34.2,-96.3)..(34.6,-96)
				..controls(35,-95.7)and(35.5,-95.9)..(36.2,-96.7)
				..controls(36.5,-97)and(37.2,-97.5)..(37.6,-97.7)..controls(38,-98)and(38.6,-98.4)..(39.1,-98.9)
				..controls(39.6,-99.4)and(40.3,-100)..(40.6,-100.3)
				..controls(41.3,-100.8)and(42.4,-101.8)..(43.7,-102.9)
				..controls(44,-103.2)and(44.6,-103.7)..(45,-103.9)
				..controls(45.4,-104.2)and(45.9,-104.6)..(46.1,-104.8)
				..controls(47.1,-105.9)and(48.8,-107.1)..(50.5,-107.9)
				..controls(51.8,-108.6)and(52.3,-109)..(53.4,-110)
				..controls(54.8,-111.2)and(55.6,-111.6)..(57,-112.1)
				..controls(57.5,-112.3)and(58.5,-112.8)..(59.4,-113.3)
				..controls(60.2,-113.7)and(61.3,-114.3)..(61.9,-114.5)
				..controls(62.4,-114.7)and(63.3,-115.1)..(63.8,-115.3)
				..controls(64.4,-115.6)and(65.4,-115.9)..(66.2,-116.1)
				..controls(67,-116.3)and(68.4,-116.6)..(69.2,-116.9)
				..controls(72.3,-117.9)and(72.8,-118)..(78,-118.1)
				..controls(84,-118.3)and(86.1,-118.1)..(88.1,-117.1)
				..controls(89.4,-116.5)and(93,-115)..(93.7,-114.9)
				..controls(94.8,-114.6)and(95.4,-114.2)..(95.9,-113.4)
				..controls(96.2,-113)and(96.7,-112.5)..(97.1,-112.2)
				..controls(98,-111.6)and(101.3,-108.4)..(102.1,-107.2)
				..controls(102.4,-106.8)and(102.9,-106.2)..(103.2,-105.8)
				..controls(103.5,-105.5)and(104,-104.8)..(104.2,-104.2)
				..controls(104.5,-103.6)and(105,-102.8)..(105.3,-102.4)
				..controls(105.6,-102)and(106.1,-101.2)..(106.5,-100.6)
				..controls(107.1,-99.5)and(108.6,-97.5)..(109.5,-96.5)
				..controls(109.9,-96.1)and(110.1,-96)..(110.4,-96)
				..controls(111.1,-96)and(111.6,-96.5)..(111.8,-97.6)
				..controls(112,-98.6)and(112.3,-99.3)..(112.8,-99.9)
				..controls(113.7,-100.9)and(115.2,-103.1)..(115.7,-104.2)
				..controls(116.1,-105)and(116.5,-105.5)..(116.9,-106)
				..controls(117.3,-106.4)and(117.8,-106.9)..(118,-107.3)
				..controls(118.2,-107.6)and(118.6,-108)..(119,-108.3)
				..controls(119.6,-108.8)and(121.1,-110.3)..(121.8,-111.2)
				..controls(122.1,-111.7)and(122.5,-112)..(123.4,-112.4)
				..controls(124.8,-113.2)and(125.8,-113.8)..(127,-114.6)
				..controls(127.5,-114.9)and(128.1,-115.3)..(128.2,-115.4)
				..controls(128.4,-115.5)and(129.3,-115.8)..(130.2,-116)
				..controls(131.1,-116.2)and(132.5,-116.6)..(133.3,-116.9)
				..controls(135.7,-117.9)and(136.1,-117.9)..(142.5,-117.8)
				..controls(145.8,-117.8)and(148.4,-117.7)..(148.8,-117.6)
				..controls(149.2,-117.6)and(150.2,-117.2)..(151,-116.8)
				..controls(152,-116.4)and(152.8,-116)..(153.6,-115.9)
				..controls(154.5,-115.7)and(155.2,-115.4)..(156.3,-114.9)
				..controls(157.1,-114.5)and(158.2,-114)..(158.7,-113.8)
				..controls(159.1,-113.7)and(159.8,-113.3)..(160.2,-113)
				..controls(160.6,-112.7)and(161.3,-112.4)..(161.7,-112.2)
				..controls(162.7,-111.9)and(163.2,-111.7)..(164.1,-111.1)
				..controls(164.5,-110.8)and(165.3,-110.4)..(165.9,-110)
				..controls(166.5,-109.7)and(167.3,-109.2)..(167.6,-108.9)
				..controls(167.9,-108.6)and(168.6,-108.2)..(169,-108)
				..controls(169.4,-107.8)and(170.2,-107.3)..(170.7,-106.9)
				..controls(171.2,-106.5)and(172.1,-105.8)..(172.7,-105.4)
				..controls(173.2,-105)and(174.1,-104.4)..(174.5,-104)
				..controls(174.9,-103.7)and(175.7,-103.1)..(176.2,-102.6)
				..controls(176.7,-102.2)and(177.4,-101.5)..(177.9,-101.2)
				..controls(178.3,-100.8)and(178.8,-100.4)..(179,-100.2)
				..controls(179.2,-100)and(179.6,-99.6)..(180,-99.3)
				..controls(180.4,-99.1)and(180.9,-98.6)..(181.2,-98.4)
				..controls(181.5,-98.1)and(182.1,-97.7)..(182.5,-97.5)
				..controls(183.2,-97.2)and(183.5,-96.9)..(184.3,-95.9)
				..controls(184.6,-95.6)and(184.6,-95.6)..(184.8,-95.7)
				..controls(185,-95.8)and(185.1,-96)..(185.1,-96.2)
				..controls(185.1,-96.5)and(185,-96.7)..(184.6,-97.2)
				..controls(184.1,-97.8)and(184,-97.9)..(183.9,-98.8)
				..controls(183.7,-100.3)and(183.3,-101.3)..(182.9,-101.8)
				..controls(182.6,-102.1)and(182.3,-102.6)..(182.2,-102.9)
				..controls(182,-103.3)and(181.6,-104)..(181.2,-104.6)
				..controls(180.9,-105.2)and(180.5,-106.1)..(180.4,-106.6)
				..controls(179.9,-108)and(179.7,-108.5)..(179,-109.4)
				..controls(178.6,-109.8)and(178.2,-110.5)..(177.9,-111)
				..controls(177.7,-111.5)and(177.3,-112.3)..(177,-112.7)
				..controls(176.6,-113.2)and(176.4,-113.7)..(176.3,-114.2)
				..controls(176.2,-114.7)and(175.9,-115.2)..(175.4,-116)
				..controls(175,-116.6)and(174.4,-117.7)..(174,-118.5)
				..controls(173.7,-119.2)and(173.3,-120.1)..(173,-120.5)
				..controls(172.8,-120.9)and(172.4,-121.8)..(172.1,-122.5)
				..controls(171.8,-123.2)and(171.3,-124.1)..(171,-124.6)
				..controls(170.8,-125.1)and(170.3,-126)..(169.9,-126.6)
				..controls(169.6,-127.3)and(169.2,-128.1)..(169,-128.3)
				..controls(168.8,-128.6)and(168.5,-129.3)..(168.3,-129.8)
				..controls(168.1,-130.3)and(167.8,-131.1)..(167.5,-131.6)
				..controls(166.9,-132.5)and(164.2,-137.9)..(164.2,-138.1)
				..controls(164.2,-138.1)and(163.9,-138.8)..(163.4,-139.7)
				..controls(162.9,-140.5)and(162.4,-141.5)..(162.3,-141.9)
				..controls(161.9,-143)and(161.6,-143.4)..(160.9,-144.1)
				..controls(160.4,-144.6)and(160.3,-144.8)..(160.1,-145.4)
				..controls(160,-145.8)and(159.6,-146.8)..(159.2,-147.8)
				..controls(158.7,-148.8)and(158.1,-150)..(157.9,-150.6)
				..controls(157.6,-151.1)and(157.2,-151.9)..(156.9,-152.4)
				..controls(156.5,-152.8)and(156.1,-153.7)..(155.8,-154.4)
				..controls(155.6,-155)and(155.2,-155.9)..(154.9,-156.3)
				..controls(154.7,-156.7)and(154.4,-157.4)..(154.2,-157.9)
				..controls(154.1,-158.4)and(153.8,-159)..(153.5,-159.3)
				..controls(152.8,-160.1)and(153,-160.1)..(142.4,-160.1)
				..controls(137.2,-160.1)and(131.8,-160.1)..(130.4,-160)--cycle
				(141.5,-103.7)..controls(140.8,-103)and(139.8,-102.1)..(139.2,-101.6)
				..controls(138.5,-101.2)and(137.9,-100.6)..(137.8,-100.4)
				..controls(137,-99.5)and(134,-96.5)..(133.2,-96)
				..controls(132.7,-95.7)and(131.8,-94.9)..(131.1,-94.3)
				..controls(129.4,-92.7)and(127.9,-91.6)..(126.6,-91.1)
				..controls(125.2,-90.5)and(124.1,-90)..(122.9,-89.1)
				..controls(122.1,-88.6)and(121.3,-88.2)..(119.9,-87.7)
				..controls(118.8,-87.3)and(117.5,-86.8)..(116.9,-86.5)
				..controls(114.3,-85.5)and(113.7,-85.4)..(109.7,-85.5)
				..controls(105.3,-85.7)and(103.9,-85.9)..(101.6,-87)
				..controls(100.7,-87.4)and(99.4,-87.8)..(97.9,-88)
				..controls(97,-88.2)and(96.6,-88.3)..(95.8,-88.9)
				..controls(95.4,-89.2)and(94.7,-89.6)..(94.2,-89.8)
				..controls(92.9,-90.3)and(92.3,-90.7)..(90.9,-92.1)
				..controls(90.2,-92.7)and(89.4,-93.3)..(89,-93.5)..controls(87,-94.5)and(85.2,-96)..(83.6,-97.9)
				..controls(83.3,-98.2)and(82.5,-98.9)..(81.8,-99.4)
				..controls(80.7,-100.1)and(80.5,-100.3)..(80,-101.1)
				..controls(78.7,-103.1)and(77.2,-103.9)..(76.4,-103)
				..controls(75.7,-102.2)and(75.7,-103.2)..(75.7,-83.6)
				..controls(75.6,-73.6)and(75.6,-64.8)..(75.5,-64)
				..controls(75.5,-63.2)and(75.5,-61.4)..(75.6,-59.8)--(75.7,-57.1)--(76,-56.8)
				..controls(76.8,-56.2)and(78,-56.3)..(79.8,-57.2)
				..controls(80.4,-57.5)and(81.3,-57.8)..(81.8,-57.9)
				..controls(82.3,-58.1)and(83,-58.4)..(83.6,-58.8)..controls(84.3,-59.2)and(85,-59.5)..(86,-59.8)
				..controls(86.8,-60.1)and(88,-60.6)..(88.7,-60.9)
				..controls(89.3,-61.2)and(90.6,-61.7)..(91.4,-62)
				..controls(92.2,-62.3)and(93.7,-62.9)..(94.6,-63.3)
				..controls(95.6,-63.7)and(96.8,-64.1)..(97.4,-64.3)
				..controls(98,-64.5)and(99,-64.9)..(99.7,-65.2)
				..controls(100.4,-65.6)and(101.5,-66)..(102.3,-66.2)
				..controls(104.2,-66.8)and(105,-67.1)..(106.2,-67.8)
				..controls(107.9,-68.7)and(110.2,-69.4)..(111.3,-69.2)
				..controls(112.5,-69)and(115,-68.2)..(115.6,-67.8)
				..controls(117.3,-66.7)and(117.6,-66.6)..(119.5,-66.2)
				..controls(121.9,-65.8)and(122.1,-65.7)..(123,-65.1)
				..controls(124.1,-64.4)and(124.4,-64.3)..(126,-63.9)
				..controls(127,-63.7)and(127.7,-63.4)..(128.4,-63.1)
				..controls(129,-62.8)and(129.7,-62.5)..(130.6,-62.4)
				..controls(132.1,-62.1)and(134,-61.5)..(135.2,-61)
				..controls(135.7,-60.8)and(136.7,-60.5)..(137.3,-60.3)
				..controls(138.4,-60)and(139.3,-59.7)..(141.1,-59)
				..controls(141.6,-58.7)and(142,-58.6)..(142.5,-58.6)
				..controls(143.2,-58.6)and(143.3,-58.7)..(143.6,-59)
				..controls(144,-59.5)and(144.1,-59.9)..(144,-60.8)
				..controls(144,-61.2)and(144,-71.1)..(144,-82.8)--(143.9,-104.1)--(143.7,-104.5)
				..controls(143.5,-104.7)and(143.3,-104.8)..(143,-104.9)
				..controls(142.6,-104.9)and(142.5,-104.8)..(141.5,-103.7)--cycle
				(27.8,-91.4)..controls(27.1,-90.9)and(26.3,-90)..(26.1,-89.5)
				..controls(26,-89.2)and(25.8,-88.6)..(25.7,-87.9)
				..controls(25.4,-86.2)and(25.6,-84.8)..(26.5,-83.5)
				..controls(27.4,-82.3)and(28.3,-82)..(29.4,-82.3)
				..controls(31.1,-82.8)and(32.1,-84.9)..(32,-87.4)
				..controls(31.9,-89.3)and(31.2,-90.6)..(29.9,-91.4)
				..controls(29.1,-91.9)and(28.6,-91.9)..(27.8,-91.4)--cycle
				(190.1,-91.3)..controls(189.7,-91.2)and(189.3,-90.9)..(189.1,-90.7)
				..controls(188.3,-89.8)and(187.6,-87.7)..(187.6,-86.2)
				..controls(187.6,-85.5)and(187.7,-85.2)..(188.2,-84.4)
				..controls(189.1,-82.6)and(190.1,-81.8)..(191.3,-82.1)
				..controls(192,-82.3)and(193.5,-83.7)..(193.8,-84.4)
				..controls(194.1,-85.1)and(194.2,-87.1)..(194,-88.3)
				..controls(193.8,-89.8)and(193.3,-90.5)..(192.3,-91.3)
				..controls(191.7,-91.7)and(190.9,-91.7)..(190.1,-91.3)--cycle
				(173.7,-83.4)..controls(172.2,-82.5)and(172,-82.5)..(170.6,-82.4)
				..controls(169.7,-82.4)and(168.8,-82.4)..(168.3,-82.5)
				..controls(167.2,-82.7)and(166.9,-82.7)..(166.6,-82.2)
				..controls(166.3,-81.8)and(166.3,-81)..(166.8,-80)
				..controls(167.2,-78.9)and(167.4,-77.4)..(167.7,-74.1)
				..controls(167.8,-72.8)and(167.9,-71.8)..(167.9,-71.7)
				..controls(168,-71.7)and(168,-70.2)..(168,-68.4)..controls(168,-65.4)and(168,-64.9)..(167.7,-63)
				..controls(167.3,-60.5)and(167.1,-59.9)..(166.4,-58.4)
				..controls(166.1,-57.7)and(165.8,-56.9)..(165.7,-56.6)
				..controls(165.7,-56.2)and(165.4,-55.6)..(165.2,-55.2)
				..controls(164.6,-53.8)and(161.9,-50.7)..(161,-50.3)
				..controls(159.8,-49.7)and(157.9,-50)..(155.2,-51.2)
				..controls(154.7,-51.4)and(153.4,-51.8)..(152.2,-52.1)
				..controls(150.8,-52.4)and(149.4,-52.9)..(148.6,-53.2)
				..controls(147.8,-53.6)and(146.7,-54)..(146,-54.2)
				..controls(145.4,-54.4)and(144.5,-54.7)..(144,-54.9)
				..controls(143.6,-55.2)and(142.7,-55.5)..(141.8,-55.7)
				..controls(141,-55.9)and(139.6,-56.3)..(138.7,-56.6)
				..controls(137.8,-57)and(136.3,-57.4)..(135.3,-57.6)
				..controls(133.8,-57.9)and(133.2,-58.1)..(132.1,-58.7)
				..controls(131.1,-59.2)and(130.3,-59.4)..(129.4,-59.6)
				..controls(128.7,-59.8)and(127.8,-60)..(127.4,-60.1)
				..controls(127,-60.3)and(126,-60.5)..(125.1,-60.7)
				..controls(123.2,-61.1)and(121.7,-61.5)..(120.2,-62.1)
				..controls(119.7,-62.3)and(118.1,-62.8)..(116.7,-63.2)
				..controls(114.6,-63.9)and(114,-64.1)..(113,-64.7)
				..controls(111.9,-65.4)and(111.8,-65.4)..(111,-65.5)
				..controls(109.9,-65.5)and(109.3,-65.4)..(107.7,-64.7)
				..controls(107,-64.4)and(106,-64)..(105.3,-63.9)
				..controls(103.5,-63.5)and(102.2,-63)..(101.2,-62.5)
				..controls(100.1,-61.8)and(99.3,-61.5)..(97.8,-61.1)
				..controls(97.2,-61)and(96.2,-60.7)..(95.7,-60.4)
				..controls(95.1,-60.1)and(94.2,-59.8)..(93.5,-59.6)
				..controls(92.9,-59.4)and(91.7,-59)..(91,-58.7)..controls(90.2,-58.4)and(88.9,-58)..(88,-57.7)
				..controls(87.2,-57.5)and(86.1,-57.1)..(85.5,-56.8)
				..controls(85,-56.6)and(84.1,-56.2)..(83.5,-56)..controls(82.9,-55.7)and(81.8,-55.2)..(81,-54.9)
				..controls(80.1,-54.5)and(79.1,-54.2)..(78.7,-54.1)
				..controls(77.9,-53.8)and(77.3,-53.6)..(76.2,-52.9)
				..controls(75.7,-52.6)and(75.2,-52.4)..(74.1,-52.2)
				..controls(72.5,-51.9)and(72.4,-51.8)..(71.3,-51.1)
				..controls(69.9,-50.1)and(69.4,-49.9)..(67.9,-49.8)
				..controls(67,-49.7)and(66.6,-49.6)..(66.3,-49.4)
				..controls(66.1,-49.3)and(65.7,-49)..(65.4,-48.8)
				..controls(65.1,-48.7)and(64.2,-48.4)..(63.6,-48.2)
				..controls(62.9,-48)and(61.7,-47.6)..(60.9,-47.2)
				..controls(60.1,-46.8)and(58.9,-46.4)..(58.2,-46.2)
				..controls(57.6,-46)and(56.8,-45.7)..(56.5,-45.5)
				..controls(54.9,-44.6)and(54.6,-44.5)..(53.2,-44.2)
				..controls(51.7,-43.8)and(51,-43.6)..(49.9,-42.9)
				..controls(49.5,-42.6)and(48.9,-42.4)..(47.5,-42)
				..controls(45.7,-41.6)and(45.6,-41.5)..(44.3,-40.6)
				..controls(43.5,-40.1)and(42.8,-39.6)..(42.7,-39.4)
				..controls(42.4,-38.7)and(43,-38.1)..(44.5,-37.6)..controls(44.9,-37.5)and(46,-37)..(46.9,-36.6)
				..controls(47.9,-36.1)and(48.8,-35.8)..(49.6,-35.6)
				..controls(50.5,-35.4)and(51.3,-35.2)..(52.4,-34.6)
				..controls(53.3,-34.1)and(54.5,-33.7)..(54.9,-33.5)
				..controls(55.4,-33.4)and(56.4,-33)..(57.2,-32.6)
				..controls(58.6,-32)and(60.1,-31.5)..(61.6,-31.2)
				..controls(62.3,-31.1)and(62.7,-31)..(63.2,-30.7)
				..controls(64.1,-30)and(64.3,-29.9)..(65.3,-29.7)
				..controls(65.7,-29.6)and(66.5,-29.2)..(67.1,-28.9)
				..controls(67.6,-28.6)and(68.8,-28.2)..(69.6,-27.9)
				..controls(70.5,-27.7)and(71.6,-27.3)..(72.1,-27)
				..controls(72.6,-26.7)and(73.8,-26.2)..(74.8,-25.9)
				..controls(75.7,-25.5)and(76.7,-25.1)..(77.1,-25)
				..controls(77.4,-24.8)and(78.5,-24.4)..(79.6,-24.1)
				..controls(80.6,-23.8)and(82,-23.2)..(82.8,-22.9)
				..controls(83.5,-22.5)and(84.8,-22)..(85.8,-21.8)
				..controls(86.7,-21.5)and(87.8,-21.1)..(88.1,-20.9)
				..controls(88.4,-20.7)and(89.3,-20.3)..(90.1,-20.1)
				..controls(90.8,-19.8)and(91.9,-19.4)..(92.4,-19.1)
				..controls(93.5,-18.5)and(94.4,-18.2)..(95.4,-18.1)
				..controls(95.7,-18)and(96.4,-17.8)..(96.9,-17.5)
				..controls(97.5,-17.3)and(98.4,-16.9)..(99,-16.6)
				..controls(99.7,-16.3)and(100.9,-15.9)..(101.9,-15.6)
				..controls(102.9,-15.3)and(104.4,-14.7)..(105.3,-14.3)
				..controls(108.1,-13.2)and(108.5,-13.1)..(110.1,-13.1)
				..controls(111.4,-13.1)and(111.5,-13.1)..(112.4,-13.5)
				..controls(113.2,-14)and(115,-14.8)..(116.6,-15.4)
				..controls(116.9,-15.5)and(117.9,-15.8)..(118.8,-16.1)
				..controls(120.9,-16.6)and(121.5,-16.9)..(122.6,-17.9)
				..controls(123,-18.3)and(123.7,-18.8)..(123.9,-18.9)
				..controls(124.6,-19.4)and(126,-19.8)..(127.9,-20.2)
				..controls(129.6,-20.5)and(129.8,-20.6)..(131.1,-21.4)
				..controls(131.5,-21.7)and(132.4,-22.1)..(133.1,-22.3)
				..controls(136.1,-23.2)and(138.3,-24.1)..(140.2,-25)
				..controls(141,-25.5)and(141.5,-25.6)..(142.5,-25.8)
				..controls(145.1,-26.3)and(146,-26.6)..(149.4,-28.1)
				..controls(151.3,-28.9)and(153.5,-29.8)..(154.3,-30.1)
				..controls(155.1,-30.4)and(156.2,-30.8)..(156.6,-31)
				..controls(157.1,-31.2)and(158.2,-31.6)..(159.1,-31.8)
				..controls(160.6,-32.2)and(160.9,-32.3)..(161.7,-32.9)
				..controls(162.4,-33.3)and(162.9,-33.6)..(163.7,-33.8)
				..controls(164.3,-34)and(164.8,-34.1)..(164.9,-34.1)
				..controls(165,-34.1)and(165.6,-34.4)..(166.4,-34.8)
				..controls(167.3,-35.2)and(168.1,-35.4)..(169,-35.6)
				..controls(170,-35.8)and(170.5,-36)..(171.7,-36.6)
				..controls(172.4,-37)and(173.3,-37.4)..(173.6,-37.4)
				..controls(175.6,-37.7)and(176.2,-38)..(176.7,-39.1)
				..controls(176.9,-39.6)and(176.9,-39.7)..(176.7,-40.4)
				..controls(176.4,-41.6)and(175.8,-42.7)..(175.2,-43.4)
				..controls(174.5,-44.1)and(172.6,-45.5)..(171.7,-45.8)
				..controls(170.6,-46.2)and(169.1,-46)..(168.2,-45.4)
				..controls(167.9,-45.2)and(167.3,-44.8)..(166.9,-44.4)
				..controls(166.4,-44.1)and(165.8,-43.6)..(165.4,-43.2)
				..controls(165.1,-42.9)and(164.4,-42.3)..(163.9,-41.9)
				..controls(163.3,-41.5)and(162.7,-40.9)..(162.4,-40.5)
				..controls(161.4,-39.2)and(159.9,-38)..(158.8,-37.8)
				..controls(158.4,-37.8)and(158.3,-37.7)..(158.2,-37.2)
				..controls(158,-36.8)and(157.4,-36.1)..(157.1,-36.1)
				..controls(157,-36.1)and(156.7,-36.1)..(156.5,-36)
				..controls(156.1,-35.9)and(155.9,-35.7)..(155.7,-35.2)
				..controls(155.3,-34.6)and(155.1,-34.4)..(154,-34.1)
				..controls(153.6,-34)and(153.2,-33.7)..(152.9,-33.5)
				..controls(152.4,-32.9)and(151.6,-32.4)..(151.3,-32.4)
				..controls(150.9,-32.4)and(150.4,-32.6)..(150.4,-32.9)
				..controls(150.4,-33.2)and(150.8,-33.4)..(151.4,-33.6)
				..controls(152.1,-33.7)and(152.8,-34.1)..(153.5,-34.9)
				..controls(154.1,-35.5)and(154.3,-35.6)..(155,-35.8)
				..controls(155.7,-36)and(155.7,-36)..(156,-36.6)
				..controls(156.3,-37.3)and(156.7,-37.7)..(157.5,-37.8)
				..controls(158.1,-37.9)and(158.2,-38)..(158.4,-38.8)
				..controls(158.5,-39.4)and(158.6,-39.6)..(159.6,-40.4)
				..controls(160.2,-41)and(160.9,-41.6)..(161.3,-41.8)
				..controls(161.7,-42)and(161.9,-42.3)..(162.2,-42.8)
				..controls(162.4,-43.2)and(162.8,-43.8)..(163.1,-44.1)
				..controls(163.9,-44.8)and(165.1,-46.7)..(165.7,-48)
				..controls(166.2,-49.3)and(166.4,-49.7)..(167.2,-50.5)
				..controls(167.5,-50.8)and(167.9,-51.4)..(168.1,-51.8)
				..controls(168.3,-52.2)and(168.8,-53)..(169.1,-53.6)
				..controls(169.6,-54.3)and(169.8,-54.9)..(169.9,-55.6)
				..controls(170.4,-57.8)and(170.5,-58.1)..(171.2,-59.5)
				..controls(171.6,-60.3)and(172,-61.5)..(172.2,-62.1)
				..controls(172.3,-62.8)and(172.7,-63.7)..(172.9,-64.1)
				..controls(173.2,-64.6)and(173.4,-65.3)..(173.4,-65.6)
				..controls(173.8,-68.8)and(174.2,-70.6)..(174.9,-72.6)
				..controls(175.5,-74.5)and(175.9,-77.1)..(176.1,-80.2)
				..controls(176.3,-83.4)and(175.5,-84.4)..(173.7,-83.4)--cycle;
				\definecolor{ezs}{RGB}{247,179,26}\fill[ezs](-5.7,-156.6)--(-3.4,-157.5)--(61.9,12.7)--(59.7,13.6)--cycle;
				\definecolor{ezs}{RGB}{247,179,22}\fill[ezs](-3.8,-157.3)--(-1.5,-158.2)--(63.9,12)--(61.6,12.9)--cycle;
				\definecolor{ezs}{RGB}{248,179,19}\fill[ezs](-1.8,-158.1)--(0.4,-158.9)--(65.8,11.2)--(63.5,12.1)--cycle;
				\definecolor{ezs}{RGB}{248,178,16}\fill[ezs](0.1,-158.8)--(2.4,-159.7)--(67.7,10.5)--(65.4,11.4)--cycle;
				\definecolor{ezs}{RGB}{249,178,15}\fill[ezs](2,-159.5)--(4.3,-160.4)--(69.6,9.8)--(67.3,10.7)--cycle;
				\definecolor{ezs}{RGB}{251,177,17}\fill[ezs](3.9,-160.3)--(6.2,-161.1)--(71.5,9)--(69.2,9.9)--cycle;
				\definecolor{ezs}{RGB}{252,179,18}\fill[ezs](5.8,-161)--(8.1,-161.9)--(73.4,8.3)--(71.1,9.2)--cycle;
				\definecolor{ezs}{RGB}{252,182,17}\fill[ezs](7.7,-161.7)--(10,-162.6)--(75.3,7.6)--(73,8.5)--cycle;
				\definecolor{ezs}{RGB}{252,185,17}\fill[ezs](9.6,-162.5)--(11.9,-163.4)--(77.2,6.8)--(75,7.7)--cycle;
				\definecolor{ezs}{RGB}{251,189,18}\fill[ezs](11.5,-163.2)--(13.8,-164.1)--(79.2,6.1)--(76.9,7)--cycle;
				\definecolor{ezs}{RGB}{251,191,22}\fill[ezs](13.4,-163.9)--(15.7,-164.8)--(81.1,5.4)--(78.8,6.2)--cycle;
				\definecolor{ezs}{RGB}{252,193,25}\fill[ezs](15.4,-164.7)--(17.6,-165.6)--(83,4.6)--(80.7,5.5)--cycle;
				\definecolor{ezs}{RGB}{253,195,27}\fill[ezs](17.3,-165.4)--(19.6,-166.3)--(84.9,3.9)--(82.6,4.8)--cycle;
				\definecolor{ezs}{RGB}{253,196,31}\fill[ezs](19.2,-166.1)--(21.5,-167)--(86.8,3.2)--(84.5,4)--cycle;
				\definecolor{ezs}{RGB}{254,199,35}\fill[ezs](21.1,-166.9)--(23.4,-167.8)--(88.7,2.4)--(86.4,3.3)--cycle;
				\definecolor{ezs}{RGB}{253,201,38}\fill[ezs](23,-167.6)--(25.3,-168.5)--(90.6,1.7)--(88.3,2.6)--cycle;
				\definecolor{ezs}{RGB}{253,203,40}\fill[ezs](24.9,-168.3)--(27.2,-169.2)--(92.5,1)--(90.3,1.8)--cycle;
				\definecolor{ezs}{RGB}{253,205,42}\fill[ezs](26.8,-169.1)--(29.1,-170)--(94.4,0.2)--(92.2,1.1)--cycle;
				\definecolor{ezs}{RGB}{253,207,45}\fill[ezs](28.7,-169.8)--(31,-170.7)--(96.4,-0.5)--(94.1,0.4)--cycle;
				\definecolor{ezs}{RGB}{253,210,49}\fill[ezs](30.7,-170.5)--(32.9,-171.4)--(98.3,-1.2)--(96,-0.4)--cycle;
				\definecolor{ezs}{RGB}{253,212,54}\fill[ezs](32.6,-171.3)--(34.9,-172.2)--(100.2,-2)--(97.9,-1.1)--cycle;
				\definecolor{ezs}{RGB}{253,214,59}\fill[ezs](34.5,-172)--(36.8,-172.9)--(102.1,-2.7)--(99.8,-1.8)--cycle;
				\definecolor{ezs}{RGB}{253,216,63}\fill[ezs](36.4,-172.7)--(38.7,-173.6)--(104,-3.4)--(101.7,-2.6)--cycle;
				\definecolor{ezs}{RGB}{253,218,65}\fill[ezs](38.3,-173.5)--(40.6,-174.4)--(105.9,-4.2)--(103.6,-3.3)--cycle;
				\definecolor{ezs}{RGB}{254,220,66}\fill[ezs](40.2,-174.2)--(42.5,-175.1)--(107.8,-4.9)--(105.5,-4)--cycle;
				\definecolor{ezs}{RGB}{254,222,68}\fill[ezs](42.1,-175)--(44.4,-175.8)--(109.7,-5.6)--(107.5,-4.8)--cycle;
				\definecolor{ezs}{RGB}{254,224,71}\fill[ezs](44,-175.7)--(46.3,-176.6)--(111.7,-6.4)--(109.4,-5.5)--cycle;
				\definecolor{ezs}{RGB}{254,226,75}\fill[ezs](46,-176.4)--(48.2,-177.3)--(113.6,-7.1)--(111.3,-6.2)--cycle;
				\definecolor{ezs}{RGB}{254,228,79}\fill[ezs](47.9,-177.2)--(50.1,-178)--(115.5,-7.8)--(113.2,-7)--cycle;
				\definecolor{ezs}{RGB}{254,230,82}\fill[ezs](49.8,-177.9)--(52.1,-178.8)--(117.4,-8.6)--(115.1,-7.7)--cycle;
				\definecolor{ezs}{RGB}{254,232,85}\fill[ezs](51.7,-178.6)--(54,-179.5)--(119.3,-9.3)--(117,-8.4)--cycle;
				\definecolor{ezs}{RGB}{254,234,88}\fill[ezs](53.6,-179.4)--(55.9,-180.2)--(121.2,-10)--(118.9,-9.2)--cycle;
				\definecolor{ezs}{RGB}{254,237,91}\fill[ezs](55.5,-180.1)--(57.8,-181)--(123.1,-10.8)--(120.8,-9.9)--cycle;
				\definecolor{ezs}{RGB}{254,239,94}\fill[ezs](57.4,-180.8)--(59.7,-181.7)--(125,-11.5)--(122.8,-10.6)--cycle;
				\definecolor{ezs}{RGB}{254,241,97}\fill[ezs](59.3,-181.6)--(61.6,-182.4)--(127,-12.2)--(124.7,-11.4)--cycle;
				\definecolor{ezs}{RGB}{254,243,101}\fill[ezs](61.2,-182.3)--(63.5,-183.2)--(128.9,-13)--(126.6,-12.1)--cycle;
				\definecolor{ezs}{RGB}{254,245,103}\fill[ezs](63.2,-183)--(65.4,-183.9)--(130.8,-13.7)--(128.5,-12.8)--cycle;
				\definecolor{ezs}{RGB}{254,245,105}\fill[ezs](65.1,-183.8)--(67.4,-184.6)--(132.7,-14.4)--(130.4,-13.6)--cycle;
				\definecolor{ezs}{RGB}{254,245,105}\fill[ezs](67,-184.5)--(69.3,-185.4)--(134.6,-15.2)--(132.3,-14.3)--cycle;
				\definecolor{ezs}{RGB}{254,244,103}\fill[ezs](68.9,-185.2)--(71.2,-186.1)--(136.5,-15.9)--(134.2,-15)--cycle;
				\definecolor{ezs}{RGB}{253,243,101}\fill[ezs](70.8,-186)--(73.1,-186.8)--(138.4,-16.6)--(136.1,-15.8)--cycle;
				\definecolor{ezs}{RGB}{253,241,98}\fill[ezs](72.7,-186.7)--(75,-187.6)--(140.3,-17.4)--(138,-16.5)--cycle;
				\definecolor{ezs}{RGB}{253,239,95}\fill[ezs](74.6,-187.4)--(76.9,-188.3)--(142.2,-18.1)--(140,-17.2)--cycle;
				\definecolor{ezs}{RGB}{253,236,92}\fill[ezs](76.5,-188.2)--(78.8,-189)--(144.2,-18.8)--(141.9,-18)--cycle;
				\definecolor{ezs}{RGB}{253,234,88}\fill[ezs](78.5,-188.9)--(80.7,-189.8)--(146.1,-19.6)--(143.8,-18.7)--cycle;
				\definecolor{ezs}{RGB}{253,232,85}\fill[ezs](80.4,-189.6)--(82.7,-190.5)--(148,-20.3)--(145.7,-19.4)--cycle;
				\definecolor{ezs}{RGB}{253,230,82}\fill[ezs](82.3,-190.4)--(84.6,-191.2)--(149.9,-21)--(147.6,-20.2)--cycle;
				\definecolor{ezs}{RGB}{253,228,78}\fill[ezs](84.2,-191.1)--(86.5,-192)--(151.8,-21.8)--(149.5,-20.9)--cycle;
				\definecolor{ezs}{RGB}{253,226,74}\fill[ezs](86.1,-191.8)--(88.4,-192.7)--(153.7,-22.5)--(151.4,-21.6)--cycle;
				\definecolor{ezs}{RGB}{253,224,70}\fill[ezs](88,-192.6)--(90.3,-193.4)--(155.6,-23.3)--(153.3,-22.4)--cycle;
				\definecolor{ezs}{RGB}{253,222,67}\fill[ezs](89.9,-193.3)--(92.2,-194.2)--(157.5,-24)--(155.3,-23.1)--cycle;
				\definecolor{ezs}{RGB}{253,220,65}\fill[ezs](91.8,-194)--(94.1,-194.9)--(159.5,-24.7)--(157.2,-23.8)--cycle;
				\definecolor{ezs}{RGB}{253,217,63}\fill[ezs](93.7,-194.8)--(96,-195.6)--(161.4,-25.5)--(159.1,-24.6)--cycle;
				\definecolor{ezs}{RGB}{253,215,60}\fill[ezs](95.7,-195.5)--(97.9,-196.4)--(163.3,-26.2)--(161,-25.3)--cycle;
				\definecolor{ezs}{RGB}{253,213,56}\fill[ezs](97.6,-196.2)--(99.9,-197.1)--(165.2,-26.9)--(162.9,-26)--cycle;
				\definecolor{ezs}{RGB}{252,211,53}\fill[ezs](99.5,-197)--(101.8,-197.8)--(167.1,-27.7)--(164.8,-26.8)--cycle;
				\definecolor{ezs}{RGB}{252,209,49}\fill[ezs](101.4,-197.7)--(103.7,-198.6)--(169,-28.4)--(166.7,-27.5)--cycle;
				\definecolor{ezs}{RGB}{252,207,45}\fill[ezs](103.3,-198.4)--(105.6,-199.3)--(170.9,-29.1)--(168.6,-28.2)--cycle;
				\definecolor{ezs}{RGB}{253,205,42}\fill[ezs](105.2,-199.2)--(107.5,-200)--(172.8,-29.9)--(170.6,-29)--cycle;
				\definecolor{ezs}{RGB}{253,202,40}\fill[ezs](107.1,-199.9)--(109.4,-200.8)--(174.7,-30.6)--(172.5,-29.7)--cycle;
				\definecolor{ezs}{RGB}{253,200,38}\fill[ezs](109,-200.6)--(111.3,-201.5)--(176.7,-31.3)--(174.4,-30.4)--cycle;
				\definecolor{ezs}{RGB}{253,198,35}\fill[ezs](111,-201.4)--(113.2,-202.2)--(178.6,-32.1)--(176.3,-31.2)--cycle;
				\definecolor{ezs}{RGB}{253,196,31}\fill[ezs](112.9,-202.1)--(115.2,-203)--(180.5,-32.8)--(178.2,-31.9)--cycle;
				\definecolor{ezs}{RGB}{253,194,28}\fill[ezs](114.8,-202.8)--(117.1,-203.7)--(182.4,-33.5)--(180.1,-32.6)--cycle;
				\definecolor{ezs}{RGB}{253,192,26}\fill[ezs](116.7,-203.6)--(119,-204.5)--(184.3,-34.3)--(182,-33.4)--cycle;
				\definecolor{ezs}{RGB}{253,189,24}\fill[ezs](118.6,-204.3)--(120.9,-205.2)--(186.2,-35)--(183.9,-34.1)--cycle;
				\definecolor{ezs}{RGB}{253,186,20}\fill[ezs](120.5,-205)--(122.8,-205.9)--(188.1,-35.7)--(185.8,-34.8)--cycle;
				\definecolor{ezs}{RGB}{253,184,17}\fill[ezs](122.4,-205.8)--(124.7,-206.7)--(190,-36.5)--(187.8,-35.6)--cycle;
				\definecolor{ezs}{RGB}{252,183,15}\fill[ezs](124.3,-206.5)--(126.6,-207.4)--(192,-37.2)--(189.7,-36.3)--cycle;
				\definecolor{ezs}{RGB}{252,181,14}\fill[ezs](126.3,-207.2)--(128.5,-208.1)--(193.9,-37.9)--(191.6,-37.1)--cycle;
				\definecolor{ezs}{RGB}{252,179,13}\fill[ezs](128.2,-208)--(130.4,-208.9)--(195.8,-38.7)--(193.5,-37.8)--cycle;
				\definecolor{ezs}{RGB}{252,178,12}\fill[ezs](130.1,-208.7)--(132.4,-209.6)--(197.7,-39.4)--(195.4,-38.5)--cycle;
				\definecolor{ezs}{RGB}{251,178,10}\fill[ezs](132,-209.4)--(134.3,-210.3)--(199.6,-40.1)--(197.3,-39.3)--cycle;
				\definecolor{ezs}{RGB}{251,178,10}\fill[ezs](133.9,-210.2)--(136.2,-211.1)--(201.5,-40.9)--(199.2,-40)--cycle;
				\definecolor{ezs}{RGB}{250,177,11}\fill[ezs](135.8,-210.9)--(138.1,-211.8)--(203.4,-41.6)--(201.1,-40.7)--cycle;
				\definecolor{ezs}{RGB}{250,178,11}\fill[ezs](137.7,-211.6)--(140,-212.5)--(205.3,-42.3)--(203.1,-41.5)--cycle;
				\definecolor{ezs}{RGB}{248,180,10}\fill[ezs](139.6,-212.4)--(141.9,-213.3)--(207.3,-43.1)--(205,-42.2)--cycle;
				\definecolor{ezs}{RGB}{247,181,9}\fill[ezs](141.5,-213.1)--(143.8,-214)--(209.2,-43.8)--(206.9,-42.9)--cycle;
				\definecolor{ezs}{RGB}{246,182,9}\fill[ezs](143.5,-213.8)--(145.7,-214.7)--(211.1,-44.5)--(208.8,-43.7)--cycle;
				\definecolor{ezs}{RGB}{246,182,10}\fill[ezs](145.4,-214.6)--(147.7,-215.5)--(213,-45.3)--(210.7,-44.4)--cycle;
			\end{scope}
			% chữ EDUZY (Times đậm, làm dày nét bằng 5 lớp lệch nhẹ; căn đúng khung chữ logo gốc)
			\foreach \dx/\dy in {-0.7/0,0.7/0,0/-0.3,0/0.3,0/0}
			\node[anchor=center,inner sep=0pt,outer sep=0pt,text=ezText,align=center] at (109.6+\dx,-193.6+\dy)
			{\resizebox{1.01cm}{0.22cm}{\fontencoding{T1}\fontfamily{ptm}\fontseries{b}\selectfont EDUZY}};
		\end{tikzpicture}%
	}%
}
%====================== THIẾT LẬP THÂN ĐỀ ======================
%--- Môi trường tự luận {bt}: ex_test 3.3.4 chưa khai báo sẵn nên tạo ở đây
%    (ex_test đã tự móc \fixshowans và \listenumerate cho bt)
\makeatletter
\@ifundefined{bt}{\newtheorem{bt}{\color{gold1}\namebt}}{}
\makeatother
\renewcommand{\sectionmark}[1]{\markright{\thesection\ #1}}
%--- Khung tiêu đề (logo + tối đa 4 dòng; dòng 3, 4 để trống {} thì bỏ qua)
%  Đề thi:    \khungtieude{THI THỬ THPT QUỐC GIA NĂM 2026}{Môn Toán -- Đề số 21}
%                         {(Đề thi gồm: ...)}{Thời gian làm bài: 90 phút (không kể thời gian phát đề).}
%  Chuyên đề: \khungtieude{CHUYÊN ĐỀ 3 -- DÃY SỐ VÀ CÔNG THỨC TRUY HỒI}
%                         {Bồi dưỡng học sinh giỏi Toán THPT}{...}{}
\newcommand{\khungtieude}[4]{%
	\begin{tcolorbox}[colframe=gold1,colback=gold5!45,boxrule=0.5mm]
		\begin{minipage}[c]{3cm}
			\centerline{\eduzylogo[2.4cm]}
		\end{minipage}\hspace{0.5cm}
		\begin{minipage}[c]{12.5cm}
			\vskip 0.15 cm
			\centerline{\ind{#1}}
			\vskip 0.15 cm
			\centerline{\textbf{#2}}
			\if\relax\detokenize{#3}\relax\else
				\vskip 0.15 cm
				\centerline{\textit{#3}}
			\fi
			\if\relax\detokenize{#4}\relax\else
				\vskip 0.15 cm
				\centerline{\textit{#4}}
			\fi
			\vskip 0.15 cm
		\end{minipage}
	\end{tcolorbox}
}
%====================== BỔ SUNG CHO BÀI HỌC, CHUYÊN ĐỀ ======================
\usepackage[shortlabels]{enumitem}
\usepackage{tabularx}
%--- Danh sách gọn: \begin{enumerate}[gon] ... \end{enumerate}
\SetEnumitemKey{gon}{leftmargin=*,itemsep=1pt,parsep=0pt,topsep=3pt}
%--- Khung lý thuyết: \begin{khung}{} ... \end{khung}; tiêu đề chỉ dùng khi thật ngắn ("Mẹo", "Lỗi thường gặp")
\newtcolorbox{khung}[1]{breakable,enhanced,colframe=gold1,colback=gold5!45,
	boxrule=0.4mm,fonttitle=\bfseries\sffamily,coltitle=white,colbacktitle=gold1,
	attach boxed title to top left={xshift=6mm,yshift=-3mm},
	boxed title style={arc=1mm,boxrule=0.3mm},title={#1}}
%--- Tiêu đề mục nhỏ trong một dạng: \muc{Dấu hiệu}, \muc{Phương pháp}, \muc{Bài tự luyện}
\makeatletter
\newcommand{\muc}[1]{\par\medskip\noindent\indam{#1}\par\nobreak\smallskip\@afterindentfalse\@afterheading}
\makeatother
%--- Ô công thức cần nhớ (căn giữa, viền vàng, vừa khít nội dung)
%     \begin{center}\congthuc{$Q(n)=(n-1)!$}\end{center}
\newtcbox{\congthuc}{on line,colframe=gold1,colback=gold5!45,boxrule=0.5pt,arc=2pt,
	left=8pt,right=8pt,top=4pt,bottom=4pt}
%--- Công tắc lời giải bài tự luyện (ví dụ mẫu luôn dùng \loigiai)
%     \loigiaiTL[{đáp số}]{lời giải đầy đủ}
%     \TLfull -> in lời giải đầy đủ;  \TLdapso -> chỉ in một dòng "Đáp số";  \TLan -> ẩn hoàn toàn
%     (\CheDo là dethi thì mọi lời giải và đáp số đều ẩn)
\makeatletter
\newif\ifTLhien   \TLhientrue
\newif\ifTLdapso  \TLdapsofalse
\newcommand{\TLfull}{\TLhientrue\TLdapsofalse}
\newcommand{\TLdapso}{\TLhientrue\TLdapsotrue}
\newcommand{\TLan}{\TLhienfalse}
\newcommand{\loigiaiTL}[2][]{%
	\ifTLhien
		\ifTLdapso
			\def\@tlarg{#1}%
			\ifx\@tlarg\@empty\else
				\ifdefstring{\kindOPTN}{dethi}{}{\par\noindent{\bfseries\color{gold1}Đáp số:} #1\par}%
			\fi
		\else
			\loigiai{#2}%
		\fi
	\fi}
\makeatother
\TLdapso% mặc định: chỉ in đáp số bài tự luyện
%====================== BẢNG ĐÚNG/SAI (\choiceTFt) ======================
%  Hàng tiêu đề nền nâu đồng chữ trắng (Phát biểu | Đ | S), các hàng tô xen kẽ kem / trắng,
%  đường kẻ vàng kim, nhãn a) b) c) d) màu nâu đồng. Bản có đáp án: cột được chọn có vòng tròn
%  nâu đồng với dấu tích trắng (câu đúng tích ở cột Đúng, câu sai tích ở cột Sai).
%  Tiêu đề ghi đủ "Đúng | Sai": \OPTN{viettat=0} sau \input (cột Đúng/Sai rộng ra, phát biểu hẹp lại).
%  Dùng bảng gốc của ex_test: \def\BangDSgoc{1} trước \input{preamble-eduzy-v2}.
\providecommand{\BangDSgoc}{0}
\ifnum\BangDSgoc=0
\usepackage{colortbl,pifont}
\colorlet{ezDSke}{gold2!60!gold1}% đường kẻ bảng (nét 0,4pt như ex_test để không đổi chiều cao bảng)
\colorlet{ezDSle}{gold5!60}% hàng a), c)
\colorlet{ezDSchan}{white}% hàng b), d)
%--- dấu chọn ở cột Đúng/Sai
%    (bề rộng 0, căn giữa ô: không làm đổi bề rộng cột của ex_test, không báo tràn lề)
\newcommand{\EZdauchon}{\makebox[0pt]{\tikz[baseline=-1ex]\node[circle,fill=gold1,inner sep=0pt,
	minimum size=1.2em,text=white,font=\small]{\ding{51}};}}
\def\tickT{\EZdauchon}\def\tickF{\EZdauchon}
\AtBeginDocument{\sbox{\ccTFTab}{\EZdauchon}}
%--- nhãn a) b) c) d) màu nâu đồng (ở bản lời giải, nhãn câu đúng vẫn đóng khung như ex_test)
\renewcommand{\FalseTF}{\stepcounter{dapan}\textbf{\color{gold1}\DapAnTF\sepTF}\ignorespaces}
%--- trước khi dựng bảng: kẻ màu vàng kim (bề rộng cột giữ nguyên như ex_test); sau bảng trả lại màu đen
\newcount\EZhangTF
\newcommand{\EZbatdauTF}{%
	\global\EZhangTF=0\relax
	\arrayrulecolor{ezDSke}}
\newcommand{\EZketthucTF}{\arrayrulecolor{black}}
%--- hàng tiêu đề nền nâu đồng, chữ trắng
\makeatletter
\newcommand{\EZo@TF}[1]{\centering\bfseries\color{white}#1}
\newcommand{\EZtieudeTF}{%
	\ifcase\cotDS
		\ifnum\cotSTT=1\relax
			\def\rowHeader{\hline\rowcolor{gold1}\EZo@TF{TT}&\EZo@TF{\phatbieu}\tabularnewline}%
		\else
			\def\rowHeader{\hline\rowcolor{gold1}\EZo@TF{\phatbieu}\tabularnewline}%
		\fi
	\or
		\ifnum\cotSTT=1\relax
			\def\rowHeader{\hline\rowcolor{gold1}\EZo@TF{TT}&\EZo@TF{\phatbieu}&\EZo@TF{\Dung}\tabularnewline}%
		\else
			\def\rowHeader{\hline\rowcolor{gold1}\EZo@TF{\phatbieu}&\EZo@TF{\Dung}\tabularnewline}%
		\fi
	\or
		\ifnum\cotSTT=1\relax
			\def\rowHeader{\hline\rowcolor{gold1}\EZo@TF{TT}&\EZo@TF{\phatbieu}&\EZo@TF{\Dung}&\EZo@TF{\Sai}\tabularnewline}%
		\else
			\def\rowHeader{\hline\rowcolor{gold1}\EZo@TF{\phatbieu}&\EZo@TF{\Dung}&\EZo@TF{\Sai}\tabularnewline}%
		\fi
	\fi}
%--- tô xen kẽ từng hàng (gọi ngay sau \hline đầu mỗi hàng)
\newcommand{\EZmauhangTF}{%
	\noalign{\global\advance\EZhangTF\@ne}%
	\ifodd\EZhangTF\rowcolor{ezDSle}\else\rowcolor{ezDSchan}\fi}
%--- gắn vào \choiceTFt / \writeChoiceTFt của ex_test
\xpretocmd{\choiceTFt}{\EZbatdauTF}{}{\typeout{[preamble-eduzy-v2] Không gắn được bảng Đúng/Sai (đầu bảng)}}
\xpatchcmd{\choiceTFt}{\ifnum\dongPB=0\def\rowHeader{}\fi}{\EZtieudeTF\ifnum\dongPB=0\def\rowHeader{}\fi}
	{}{\typeout{[preamble-eduzy-v2] Không gắn được bảng Đúng/Sai (hàng tiêu đề)}}
\xpatchcmd{\writeChoiceTFt}{\hline\xdef\dotEX{\sdotEX}}{\hline\EZmauhangTF\xdef\dotEX{\sdotEX}}
	{}{\typeout{[preamble-eduzy-v2] Không gắn được bảng Đúng/Sai (tô hàng)}}
\xpatchcmd{\writeChoiceTFt}{{\end{longtable}}}{{\end{longtable}\EZketthucTF}}
	{}{\typeout{[preamble-eduzy-v2] Không gắn được bảng Đúng/Sai (cuối longtable)}}
\xpatchcmd{\writeChoiceTFt}{{\end{xtabular}\end{center}\medskip\noindent}}{{\end{xtabular}\end{center}\EZketthucTF\medskip\noindent}}
	{}{\typeout{[preamble-eduzy-v2] Không gắn được bảng Đúng/Sai (cuối xtabular)}}
\fi
%=====================================================================
\renewcommand{\parallel}{\mathrel{/\!\!/}}

%      \TLdapso                      % tài liệu: \TLfull | \TLdapso | \TLan
%      \begin{document} ... \end{document}
%  Logo Eduzy vẽ bằng TikZ ngay trong file (\eduzylogo), không cần file ảnh.
%  Biên dịch: pdfLaTeX (máy có gói vietnam) hoặc XeLaTeX.
%============================================================
\documentclass[12pt,a4paper,twoside]{article}%<-- twoside: bat dau trang/chan trang chan-le
\usepackage{iftex}
\usepackage{amsmath,amssymb,tikz,mathrsfs,tkz-tab}
\usepackage{graphicx}
\usepackage{pgfplots}% đồ thị hàm số
\pgfplotsset{compat=1.17}
%\usepackage{mathptmx}
\usepackage{mathpazo}
\ifXeTeX% XeLaTeX: font chữ có đủ dấu tiếng Việt, giữ nguyên font toán mathpazo
	\usepackage[no-math]{fontspec}
	\setmainfont{TeX Gyre Pagella}
	\newfontfamily\EZqag{TeX Gyre Adventor}[NFSSFamily=qag,Ligatures=TeX]
\fi
\usepackage[many]{tcolorbox}
\usetikzlibrary{math,through,calc,intersections,angles,quotes,shapes,shapes.geometric,arrows,patterns,matrix,chains,arrows.meta,decorations.shapes,decorations.fractals,decorations.markings,shadows,decorations.pathreplacing}
\usepackage[top=1.6cm, bottom=1.8cm, left=1.8cm, right=1.8cm]{geometry}% v2: top 1.6cm (bản cũ 1.5cm) để chữ không sát dải đầu trang
\providecommand{\CheDo}{loigiai}% loigiai | dethi | solcolor | color | book
\expandafter\usepackage\expandafter[\CheDo]{ex_test}
% \tenchuong: khai báo ở file chính (trước \input) hoặc \renewcommand sau \input
%============================================================
%  [GỘP TỪ eduzy-de.tex] EDUZY - MÀU NÂU ĐỒNG / VÀNG KIM
%  + ĐẦU TRANG, CHÂN TRANG KIỂU DẢI NGANG (v2), ĐỐI XỨNG TRANG CHẴN / TRANG LẺ.
%  (không cần file eduzy-de.tex riêng nữa)
%  Khối này phải đặt SAU \usepackage[...]{ex_test} và SAU \def\tenchuong{...}
%============================================================
\usepackage{lastpage}
\usepackage{needspace}
%--- DÀN TRANG ĐỀ THI: tách câu (THÊM 05/11/2026; chỉ thêm lệnh mới, KHÔNG đổi hành vi mặc định) ---
%  Hai công tắc, mặc định tắt (= preamble chạy y như cũ). Bật bằng cách đặt TRƯỚC %============================================================
%  PREAMBLE EDUZY v2 - dùng chung cho đề thi, bài học, chuyên đề (gói ex_test)
%  Khác bản preamble-eduzy.tex ở đầu trang, chân trang: dải nền nâu đồng tràn ngang,
%  logo + Eduzy.com.vn, ô chữ vàng kem cạnh xiên, chân trang "Trang (số)";
%  logo Eduzy chìm giữa mọi trang (tắt: \logochimtat, bật lại: \logochimbat).
%  Cách dùng ở file chính (xem mau-chuan.tex, mau-chuyen-de.tex):
%      \def\CheDo{loigiai}          % loigiai | dethi | solcolor | color | book
%      \def\tenchuong{<dòng chữ trên đầu trang>}
%      \def\thuonghieu{<chữ ở chân trang lẻ>}   % không khai báo thì dùng mặc định
%      \def\tacgia{<chữ ở chân trang chẵn>}      % không khai báo thì dùng mặc định
%      \def\khauhieu{<chữ cạnh logo ở đầu trang>}  % tùy chọn, mặc định Eduzy.com.vn
%      \input{preamble-eduzy-v2}
%      \TLdapso                      % tài liệu: \TLfull | \TLdapso | \TLan
%      \begin{document} ... \end{document}
%  Logo Eduzy vẽ bằng TikZ ngay trong file (\eduzylogo), không cần file ảnh.
%  Biên dịch: pdfLaTeX (máy có gói vietnam) hoặc XeLaTeX.
%============================================================
\documentclass[12pt,a4paper,twoside]{article}%<-- twoside: bat dau trang/chan trang chan-le
\usepackage{iftex}
\usepackage{amsmath,amssymb,tikz,mathrsfs,tkz-tab}
\usepackage{graphicx}
\usepackage{pgfplots}% đồ thị hàm số
\pgfplotsset{compat=1.17}
%\usepackage{mathptmx}
\usepackage{mathpazo}
\ifXeTeX% XeLaTeX: font chữ có đủ dấu tiếng Việt, giữ nguyên font toán mathpazo
	\usepackage[no-math]{fontspec}
	\setmainfont{TeX Gyre Pagella}
	\newfontfamily\EZqag{TeX Gyre Adventor}[NFSSFamily=qag,Ligatures=TeX]
\fi
\usepackage[many]{tcolorbox}
\usetikzlibrary{math,through,calc,intersections,angles,quotes,shapes,shapes.geometric,arrows,patterns,matrix,chains,arrows.meta,decorations.shapes,decorations.fractals,decorations.markings,shadows,decorations.pathreplacing}
\usepackage[top=1.6cm, bottom=1.8cm, left=1.8cm, right=1.8cm]{geometry}% v2: top 1.6cm (bản cũ 1.5cm) để chữ không sát dải đầu trang
\providecommand{\CheDo}{loigiai}% loigiai | dethi | solcolor | color | book
\expandafter\usepackage\expandafter[\CheDo]{ex_test}
% \tenchuong: khai báo ở file chính (trước \input) hoặc \renewcommand sau \input
%============================================================
%  [GỘP TỪ eduzy-de.tex] EDUZY - MÀU NÂU ĐỒNG / VÀNG KIM
%  + ĐẦU TRANG, CHÂN TRANG KIỂU DẢI NGANG (v2), ĐỐI XỨNG TRANG CHẴN / TRANG LẺ.
%  (không cần file eduzy-de.tex riêng nữa)
%  Khối này phải đặt SAU \usepackage[...]{ex_test} và SAU \def\tenchuong{...}
%============================================================
\usepackage{lastpage}
\usepackage{needspace}
%--- DÀN TRANG ĐỀ THI: tách câu (THÊM 05/11/2026; chỉ thêm lệnh mới, KHÔNG đổi hành vi mặc định) ---
%  Hai công tắc, mặc định tắt (= preamble chạy y như cũ). Bật bằng cách đặt TRƯỚC \input{preamble-eduzy-v2}
%  trong file đề (tools/dan_trang/dan_trang.py tự ghi khi đặt \NeedspaceCau):
%    \def\TachCauDS{1}   câu Đúng/Sai được tách: câu dẫn (kèm hình) ở cuối trang, bảng Phát biểu sang trang sau
%                        (câu dẫn và hình không tách nhau, bảng không tách giữa hai trang)
%    \def\TachCauLG{1}   bản lời giải (\CheDo khác dethi): phần đề của câu cao hơn 1/3 trang được phép tách
%  Lệnh dùng trong file đề (giá trị do tools/dan_trang/dan_trang.py đo ở bản dethi):
%    \NeedspaceCau[<cao câu dẫn + hình>]{<cao cả khối câu>}   ngay trên \begin{ex} (câu đầu phần: trên tiêu đề phần)
%        công tắc tắt: đúng bằng \Needspace{<cao cả khối câu>} như cũ;
%        \TachCauDS=1 và có [..] (câu Đúng/Sai được chọn tách): chỉ giữ câu dẫn + hình;
%        \TachCauLG=1 ở bản lời giải: đòi tối đa 1/3 \textheight.
%    \NeedspaceBang{<cao bảng Phát biểu>}%   ngay trước \choiceTFt: \TachCauDS=1 thì giữ trọn bảng trên một trang;
%        công tắc tắt: không làm gì.
\providecommand{\TachCauDS}{0}
\providecommand{\TachCauLG}{0}
\newlength{\EZnsCau}
\newcommand{\NeedspaceCau}[2][]{%
	\setlength{\EZnsCau}{#2}%
	\ifnum\TachCauDS=1 \if\relax\detokenize{#1}\relax\else\setlength{\EZnsCau}{#1}\fi\fi
	\ifnum\TachCauLG=1 \ifdefstring{\kindOPTN}{dethi}{}{\ifdim\EZnsCau>0.3333\textheight\setlength{\EZnsCau}{0.3333\textheight}\fi}\fi
	\Needspace{\EZnsCau}}
\newcommand{\NeedspaceBang}[1]{\ifnum\TachCauDS=1 \Needspace{#1}\fi}

%--- Thang màu: đậm -> nhạt
\definecolor{gold1}{RGB}{126,96,0}    % nâu đồng đậm  #7E6000
\definecolor{gold2}{RGB}{192,144,3}   % vàng kim      #C09003
\definecolor{gold3}{RGB}{255,215,99}  % vàng nhạt     #FFD763
\definecolor{gold4}{RGB}{255,231,153} % vàng rất nhạt #FFE799
\definecolor{gold5}{RGB}{255,243,205} % vàng kem      #FFF3CD
%--- Ghi đè tông xanh/đỏ cũ
\definecolor{myblue}{RGB}{126,96,0}
\colorlet{blue}{gold1}
\colorlet{cyan}{gold2}
\colorlet{teal}{gold2}

%--- Tên chương mặc định (file đề nên tự khai báo lại)
\providecommand{\tenchuong}{}
%--- Chân trang
\providecommand{\thuonghieu}{TOÁN 12 -- Eduzy.com.vn}
\providecommand{\tacgia}{Tô Gia Bảo -- 0353703783}

%--- Chữ cạnh logo ở đầu trang
\providecommand{\khauhieu}{Eduzy.com.vn}

%====================== THÔNG SỐ ĐẦU TRANG, CHÂN TRANG (v2) ======================
%  Dải nền đậm tràn ngang trang; ô chữ vàng kem có cạnh xiên, kèm 2 vạch xiên chuyển màu.
%  Đầu trang lẻ : logo + \khauhieu (trái)   | ô chữ \tenchuong (phải)
%  Chân trang lẻ: ô chữ \thuonghieu (trái)  | Trang (số) (phải)
%  Trang chẵn soi gương trang lẻ; chân trang chẵn ghi \tacgia.
%--- Màu: chỉ dùng thang vàng kim Eduzy
\colorlet{ezNen}{gold1}% dải nền đậm
\colorlet{ezVach1}{gold3}% vạch xiên sát dải nền (đậm hơn)
\colorlet{ezVach2}{gold4}% vạch xiên sát ô chữ (nhạt hơn)
\colorlet{ezO}{gold5}% ô chữ vàng kem
%--- Kích thước
\newlength{\EZdai}\setlength{\EZdai}{1.1cm}% chiều cao dải nền
\newlength{\EZo}\setlength{\EZo}{0.68cm}% chiều cao ô chữ
\newlength{\EZxien}\setlength{\EZxien}{0.55cm}% độ lệch ngang của cạnh xiên
\newlength{\EZvach}\setlength{\EZvach}{0.75cm}% bề rộng mỗi vạch xiên
\newlength{\EZkhe}\setlength{\EZkhe}{0.12cm}% khe trắng giữa vạch và ô chữ
\newlength{\EZdem}\setlength{\EZdem}{0.4cm}% khoảng đệm từ chữ tới cạnh xiên
\newlength{\EZle}\setlength{\EZle}{1.8cm}% lề trái/phải của chữ (bằng lề thân bài)
\newlength{\EZrongmin}\setlength{\EZrongmin}{4.5cm}% ô chữ rộng tối thiểu
\newlength{\EZrong}% bề rộng chữ trong ô (tính tự động)
%--- Font
\newcommand{\EZfontO}{\fontsize{10.5pt}{12pt}\selectfont\bfseries}% chữ trong ô
\newcommand{\EZfontKH}{\fontfamily{qag}\fontsize{11pt}{12pt}\selectfont\bfseries}% \khauhieu cạnh logo
\newcommand{\EZfontTrang}{\fontsize{11.5pt}{12pt}\selectfont\bfseries}% chữ "Trang" và số trang
%--- Đo bề rộng chữ trong ô
\newcommand{\EZdorong}[1]{%
	\settowidth{\EZrong}{\EZfontO #1}%
	\ifdim\EZrong<\EZrongmin\setlength{\EZrong}{\EZrongmin}\fi}
%--- Logo + chữ trắng \khauhieu (dựng sẵn một lần khi vào \begin{document})
\newsavebox{\EZlogobox}
\AtBeginDocument{\sbox{\EZlogobox}{%
		\raisebox{-0.5\height}{\eduzylogo[1cm]}\hspace{0.15cm}%
		\raisebox{-0.5\height}{\color{white}\EZfontKH\khauhieu}}}
%--- Hai vạch xiên + khe trắng. #1: điểm gốc ở mép dưới, #2: hướng (1 sang phải, -1 sang trái),
%    #3: hướng nghiêng (1 kiểu "/", -1 kiểu "\")
\newcommand{\EZvachxien}[3]{%
	\fill[white] (#1) --++({(#3)*\EZxien},\EZo) --++({(#2)*\EZkhe},0) --++({-(#3)*\EZxien},-\EZo) --cycle;
	\fill[ezVach2] ([xshift=(#2)*\EZkhe]#1) --++({(#3)*\EZxien},\EZo) --++({(#2)*\EZvach},0) --++({-(#3)*\EZxien},-\EZo) --cycle;
	\fill[ezVach1] ([xshift=(#2)*(\EZkhe+\EZvach)]#1) --++({(#3)*\EZxien},\EZo) --++({(#2)*\EZvach},0) --++({-(#3)*\EZxien},-\EZo) --cycle;
	\draw[white,line width=0.6pt] ([xshift=(#2)*(\EZkhe+\EZvach)]#1) --++({(#3)*\EZxien},\EZo);
}
%--- Chữ trong ô (căn giữa theo chiều cao ô, không lệch vì dấu tiếng Việt)
\tikzset{ezchu/.style={inner sep=0pt,text=ezNen,text height=0.72em,text depth=0.22em}}

%====================== LOGO CHÌM GIỮA TRANG ======================
%  Vẽ cùng đầu trang nên nằm dưới chữ của thân bài, có trên mọi trang.
%  Tắt: \logochimtat (từ trang đó trở đi), bật lại: \logochimbat.
%  Đổi cỡ, độ đậm: khai báo trước \input, ví dụ \def\EZchimcao{9cm} \def\EZchimdo{0.1}
\providecommand{\EZchimcao}{10cm}% chiều cao logo chìm
\providecommand{\EZchimdo}{0.12}% độ đậm (0 = trong suốt, 1 = rõ hẳn)
\newif\ifEZchim \EZchimtrue
\newcommand{\logochimtat}{\EZchimfalse}
\newcommand{\logochimbat}{\EZchimtrue}
%  Dựng sẵn logo mờ một lần khi vào \begin{document}; transparency group làm mờ cả logo
%  một lượt nên không lộ vệt ở chỗ các dải màu chồng nhau.
\newsavebox{\EZchimtho}
\newsavebox{\EZchimbox}
\AtBeginDocument{%
	\sbox{\EZchimtho}{\eduzylogo[\EZchimcao]}%
	\sbox{\EZchimbox}{\begin{tikzpicture}
			\begin{scope}[transparency group,opacity=\EZchimdo]
				\node[inner sep=0pt] {\usebox{\EZchimtho}};
			\end{scope}
	\end{tikzpicture}}}
\newcommand{\EZvechim}{%
	\ifEZchim
		\node[inner sep=0pt] at (current page.center) {\usebox{\EZchimbox}};
	\fi}

%====================== ĐẦU TRANG - TRANG LẺ ======================
%  (logo bên trái, ô chữ \tenchuong bên phải)
\newcommand{\myfancyheadOdd}{%
	\EZdorong{\tenchuong}%
	\begin{tikzpicture}[remember picture,overlay]
		\coordinate (T) at (current page.north west);
		\coordinate (P) at (current page.north east);
		%--- logo chìm giữa trang
		\EZvechim
		%--- dải nền
		\fill[ezNen] (T) rectangle ([yshift=-\EZdai]P);
		%--- ô chữ: X là góc dưới bên trái, cạnh trái xiên kiểu "/"
		\coordinate (X) at ([xshift=-\EZle-\EZrong-\EZdem-\EZxien,yshift=-\EZdai]P);
		\fill[ezO] (X) --++(\EZxien,\EZo) -- ([yshift=-\EZdai+\EZo]P) -- ([yshift=-\EZdai]P) --cycle;
		\EZvachxien{X}{-1}{1}
		\node[ezchu,anchor=east] at ([xshift=-\EZle,yshift=-\EZdai+0.5*\EZo]P) {\EZfontO\tenchuong};
		%--- logo + tên
		\node[inner sep=0pt,anchor=west] at ([xshift=\EZle,yshift=-0.5*\EZdai]T) {\usebox{\EZlogobox}};
	\end{tikzpicture}%
}

%====================== ĐẦU TRANG - TRANG CHẴN =====================
%  (soi gương: ô chữ \tenchuong bên trái, logo bên phải)
\newcommand{\myfancyheadEven}{%
	\EZdorong{\tenchuong}%
	\begin{tikzpicture}[remember picture,overlay]
		\coordinate (T) at (current page.north west);
		\coordinate (P) at (current page.north east);
		%--- logo chìm giữa trang
		\EZvechim
		%--- dải nền
		\fill[ezNen] (T) rectangle ([yshift=-\EZdai]P);
		%--- ô chữ: X là góc dưới bên phải, cạnh phải xiên kiểu "\"
		\coordinate (X) at ([xshift=\EZle+\EZrong+\EZdem+\EZxien,yshift=-\EZdai]T);
		\fill[ezO] (X) --++(-\EZxien,\EZo) -- ([yshift=-\EZdai+\EZo]T) -- ([yshift=-\EZdai]T) --cycle;
		\EZvachxien{X}{1}{-1}
		\node[ezchu,anchor=west] at ([xshift=\EZle,yshift=-\EZdai+0.5*\EZo]T) {\EZfontO\tenchuong};
		%--- logo + tên
		\node[inner sep=0pt,anchor=east] at ([xshift=-\EZle,yshift=-0.5*\EZdai]P) {\usebox{\EZlogobox}};
	\end{tikzpicture}%
}

%====================== CHÂN TRANG - TRANG LẺ ======================
%  (ô chữ \thuonghieu bên trái, "Trang (số)" bên phải)
\newcommand{\myfancyfootOdd}{%
	\EZdorong{\thuonghieu}%
	\begin{tikzpicture}[remember picture,overlay]
		\coordinate (T) at (current page.south west);
		\coordinate (P) at (current page.south east);
		%--- dải nền
		\fill[ezNen] (T) rectangle ([yshift=\EZdai]P);
		%--- ô chữ sát mép trên dải: X là góc dưới bên phải, cạnh phải xiên kiểu "/"
		\coordinate (X) at ([xshift=\EZle+\EZrong+\EZdem,yshift=\EZdai-\EZo]T);
		\fill[ezO] ([yshift=\EZdai-\EZo]T) -- (X) --++(\EZxien,\EZo) -- ([yshift=\EZdai]T) --cycle;
		\EZvachxien{X}{1}{1}
		\node[ezchu,anchor=west] at ([xshift=\EZle,yshift=\EZdai-0.5*\EZo]T) {\EZfontO\thuonghieu};
		%--- số trang
		\node[circle,draw=white,line width=0.9pt,text=white,inner sep=0pt,minimum size=0.62cm,anchor=east]
		(so) at ([xshift=-\EZle,yshift=0.5*\EZdai]P) {\EZfontTrang\thepage};
		\node[inner sep=0pt,text=white,anchor=base east] (tr) at ([xshift=-0.15cm]so.base west) {\EZfontTrang Trang};
	\end{tikzpicture}%
}

%====================== CHÂN TRANG - TRANG CHẴN ====================
%  (soi gương: "Trang (số)" bên trái, ô chữ \tacgia bên phải)
\newcommand{\myfancyfootEven}{%
	\EZdorong{\tacgia}%
	\begin{tikzpicture}[remember picture,overlay]
		\coordinate (T) at (current page.south west);
		\coordinate (P) at (current page.south east);
		%--- dải nền
		\fill[ezNen] (T) rectangle ([yshift=\EZdai]P);
		%--- ô chữ sát mép trên dải: X là góc dưới bên trái, cạnh trái xiên kiểu "\"
		\coordinate (X) at ([xshift=-\EZle-\EZrong-\EZdem,yshift=\EZdai-\EZo]P);
		\fill[ezO] ([yshift=\EZdai-\EZo]P) -- (X) --++(-\EZxien,\EZo) -- ([yshift=\EZdai]P) --cycle;
		\EZvachxien{X}{-1}{-1}
		\node[ezchu,anchor=east] at ([xshift=-\EZle,yshift=\EZdai-0.5*\EZo]P) {\EZfontO\tacgia};
		%--- số trang
		%    (vòng tròn ẩn so0 nằm giữa dải để lấy đường chân chữ, giống trang lẻ)
		\node[circle,inner sep=0pt,minimum size=0.62cm] (so0) at ([yshift=0.5*\EZdai]T) {\EZfontTrang\phantom{\thepage}};
		\node[inner sep=0pt,text=white,anchor=base west] (tr) at ($(so0.base -| T)+(\EZle,0)$) {\EZfontTrang Trang};
		\node[circle,draw=white,line width=0.9pt,text=white,inner sep=0pt,minimum size=0.62cm,anchor=base west]
		(so) at ([xshift=0.15cm]tr.base east) {\EZfontTrang\thepage};
	\end{tikzpicture}%
}

%--- Giữ tương thích với các file cũ đang gọi \myfancyhead, \myfancyfoot
\newcommand{\myfancyhead}{\myfancyheadOdd}
\newcommand{\myfancyfoot}{\myfancyfootOdd}

%======================= GẮN VÀO TRANG =============================
\pagestyle{fancy}
\fancyhf{}
\fancyhead[LO]{\myfancyheadOdd}
\fancyhead[RE]{\myfancyheadEven}
\fancyfoot[LO]{\myfancyfootOdd}
\fancyfoot[RE]{\myfancyfootEven}
\renewcommand{\headrulewidth}{0pt}
\renewcommand{\footrulewidth}{0pt}
\fancypagestyle{plain}{%
	\fancyhf{}
	\fancyhead[LO]{\myfancyheadOdd}
	\fancyhead[RE]{\myfancyheadEven}
	\fancyfoot[LO]{\myfancyfootOdd}
	\fancyfoot[RE]{\myfancyfootEven}
	\renewcommand{\headrulewidth}{0pt}
	\renewcommand{\footrulewidth}{0pt}
}
%=================== HẾT PHẦN GỘP TỪ eduzy-de.tex ===================
\newcommand{\indam}[1]{
	{\fontsize{12pt}{9pt}\sffamily \selectfont\bfseries\color{gold1} #1}
}
\newcommand{\ind}[1]{
	{\fontsize{14pt}{9pt}\fontfamily{qag} \selectfont\bfseries\color{gold1} #1}
}
\renewcommand{\FalseEX}{\stepcounter{dapan}\textbf{\selectfont\color{gold1}\Alph{dapan}.}\ignorespaces}
\renewtheorem{ex}{\color{gold1}Câu}
\renewcommand{\baselinestretch}{1.1}
\newcommand{\hoac}[1]{
	\left[\begin{aligned}#1\end{aligned}\right.}
\newcommand{\heva}[1]{
	\left\{\begin{aligned}#1\end{aligned}\right.}
\usepackage{esvect}
\def\vec{\vv}
\def\vt{\vv}
\def\overrightarrow{\vv}
%====================== LOGO EDUZY VẼ BẰNG TIKZ ======================
%  Đường viền dò từ logo gốc; dải màu vàng kim đo trực tiếp trên ảnh gốc
%  (nghiêng 21 độ, tô bằng 80 dải hẹp nên hiển thị giống nhau trên mọi trình xem PDF).
%  Dùng: \eduzylogo          -> cao 2.4cm (mặc định)
%        \eduzylogo[1.5cm]   -> cao 1.5cm
\definecolor{ezText}{RGB}{245,200,46}
\newcommand{\eduzylogo}[1][2.4cm]{%
	\resizebox{!}{#1}{%
		\begin{tikzpicture}[x=0.01cm,y=0.01cm]
			\path[use as bounding box] (0,0) rectangle (224,-224);
			% mũ tốt nghiệp + vương miện
			\begin{scope}
				\clip
				(119.4,-176.1)..controls(118.9,-176.1)and(116.8,-176)..(114.6,-176)
				..controls(101.4,-176)and(85.8,-175.9)..(83.7,-175.7)
				..controls(81.8,-175.6)and(81,-175.6)..(80.2,-175.7)
				..controls(78.8,-175.9)and(77.7,-175.9)..(76.2,-175.7)
				..controls(74.9,-175.5)and(70.8,-175.5)..(64.3,-175.7)
				..controls(59.1,-175.9)and(58.3,-175.7)..(58.6,-173.7)
				..controls(58.6,-172.8)and(59.6,-169.6)..(60.3,-168)
				..controls(60.8,-166.9)and(61.4,-165.8)..(62.3,-165)
				..controls(63.1,-164.2)and(64.3,-163.4)..(65,-163.2)
				..controls(65.5,-163.1)and(72.3,-163)..(76.6,-163.1)
				..controls(77.9,-163.1)and(80.1,-163.1)..(81.4,-163)
				..controls(82.8,-162.9)and(86.3,-162.9)..(89.2,-162.9)
				..controls(93,-162.9)and(94.9,-162.8)..(96,-162.7)
				..controls(97.2,-162.5)and(101,-162.5)..(114,-162.4)
				..controls(123,-162.3)and(134.9,-162.2)..(140.2,-162.1)
				..controls(150.1,-162)and(152.2,-162.1)..(153.6,-162.4)
				..controls(154.4,-162.6)and(157.1,-163.8)..(158.1,-164.4)
				..controls(158.8,-164.8)and(159.7,-165.7)..(159.9,-166.3)
				..controls(160,-166.5)and(160.2,-167.1)..(160.2,-167.6)
				..controls(160.3,-168)and(160.6,-169.4)..(160.9,-170.6)
				..controls(161.2,-171.7)and(161.5,-173.1)..(161.6,-173.6)
				..controls(161.8,-174.4)and(161.8,-174.5)..(161.5,-174.9)
				..controls(161.2,-175.6)and(160.8,-175.9)..(159.9,-176)
				..controls(159.2,-176.1)and(154.1,-176.2)..(151.5,-176.2)
				..controls(148.1,-176.1)and(122.6,-176.1)..(121.6,-176.2)
				..controls(120.8,-176.2)and(119.8,-176.2)..(119.4,-176.1)--cycle
				(130.4,-160)..controls(129,-160)and(119.8,-160)..(110.1,-159.9)
				..controls(100.3,-159.8)and(88.5,-159.7)..(83.9,-159.6)
				..controls(77.8,-159.5)and(75.3,-159.5)..(74.5,-159.6)
				..controls(73.9,-159.7)and(72.5,-159.8)..(71.3,-159.8)
				..controls(69.6,-159.8)and(69,-159.8)..(68.5,-159.6)
				..controls(67.6,-159.3)and(66.8,-158.6)..(66.2,-157.6)
				..controls(66,-157.3)and(65.6,-156.5)..(65.3,-156.1)
				..controls(64.6,-155.1)and(64.5,-154.9)..(63.7,-152.9)
				..controls(63.2,-151.6)and(62.9,-150.9)..(61.8,-149.6)
				..controls(61.5,-149.2)and(61.1,-148.5)..(60.9,-147.9)
				..controls(60.7,-147.4)and(60.3,-146.6)..(60.1,-146.1)
				..controls(59.9,-145.5)and(59.5,-144.9)..(59,-144.3)
				..controls(58.5,-143.7)and(58.2,-143.2)..(57.9,-142.5)
				..controls(57.7,-141.9)and(57.3,-141.1)..(57,-140.7)
				..controls(56.6,-140.3)and(56.1,-139.3)..(55.8,-138.5)
				..controls(55.5,-137.8)and(55.1,-136.8)..(54.9,-136.5)
				..controls(54.7,-136.1)and(54.3,-135.2)..(53.9,-134.5)
				..controls(53.6,-133.8)and(53,-132.7)..(52.6,-132)
				..controls(51.6,-130.5)and(51.4,-130.2)..(51.1,-129)
				..controls(50.6,-127.5)and(50.6,-127.4)..(49.7,-126.3)
				..controls(49.2,-125.7)and(48.6,-124.8)..(48.3,-124.2)
				..controls(48,-123.6)and(47.5,-122.7)..(47.3,-122.2)
				..controls(47,-121.8)and(46.5,-120.7)..(46.1,-119.8)
				..controls(45.8,-118.9)and(45.2,-117.8)..(44.9,-117.2)
				..controls(44.2,-116.1)and(43.3,-114.4)..(42.3,-112.2)
				..controls(41.8,-111.3)and(41.3,-110.4)..(41.1,-110.1)
				..controls(40.5,-109.4)and(40.3,-109)..(40,-107.8)
				..controls(39.8,-107)and(39.6,-106.5)..(39.2,-106)
				..controls(38.9,-105.6)and(38.4,-104.7)..(38,-104)
				..controls(37.7,-103.4)and(37.2,-102.5)..(37,-102.1)
				..controls(36.7,-101.7)and(36.3,-100.8)..(36,-100.1)
				..controls(35.6,-98.9)and(35.3,-98.4)..(34.6,-97.6)
				..controls(34.2,-97.1)and(34.2,-96.3)..(34.6,-96)
				..controls(35,-95.7)and(35.5,-95.9)..(36.2,-96.7)
				..controls(36.5,-97)and(37.2,-97.5)..(37.6,-97.7)..controls(38,-98)and(38.6,-98.4)..(39.1,-98.9)
				..controls(39.6,-99.4)and(40.3,-100)..(40.6,-100.3)
				..controls(41.3,-100.8)and(42.4,-101.8)..(43.7,-102.9)
				..controls(44,-103.2)and(44.6,-103.7)..(45,-103.9)
				..controls(45.4,-104.2)and(45.9,-104.6)..(46.1,-104.8)
				..controls(47.1,-105.9)and(48.8,-107.1)..(50.5,-107.9)
				..controls(51.8,-108.6)and(52.3,-109)..(53.4,-110)
				..controls(54.8,-111.2)and(55.6,-111.6)..(57,-112.1)
				..controls(57.5,-112.3)and(58.5,-112.8)..(59.4,-113.3)
				..controls(60.2,-113.7)and(61.3,-114.3)..(61.9,-114.5)
				..controls(62.4,-114.7)and(63.3,-115.1)..(63.8,-115.3)
				..controls(64.4,-115.6)and(65.4,-115.9)..(66.2,-116.1)
				..controls(67,-116.3)and(68.4,-116.6)..(69.2,-116.9)
				..controls(72.3,-117.9)and(72.8,-118)..(78,-118.1)
				..controls(84,-118.3)and(86.1,-118.1)..(88.1,-117.1)
				..controls(89.4,-116.5)and(93,-115)..(93.7,-114.9)
				..controls(94.8,-114.6)and(95.4,-114.2)..(95.9,-113.4)
				..controls(96.2,-113)and(96.7,-112.5)..(97.1,-112.2)
				..controls(98,-111.6)and(101.3,-108.4)..(102.1,-107.2)
				..controls(102.4,-106.8)and(102.9,-106.2)..(103.2,-105.8)
				..controls(103.5,-105.5)and(104,-104.8)..(104.2,-104.2)
				..controls(104.5,-103.6)and(105,-102.8)..(105.3,-102.4)
				..controls(105.6,-102)and(106.1,-101.2)..(106.5,-100.6)
				..controls(107.1,-99.5)and(108.6,-97.5)..(109.5,-96.5)
				..controls(109.9,-96.1)and(110.1,-96)..(110.4,-96)
				..controls(111.1,-96)and(111.6,-96.5)..(111.8,-97.6)
				..controls(112,-98.6)and(112.3,-99.3)..(112.8,-99.9)
				..controls(113.7,-100.9)and(115.2,-103.1)..(115.7,-104.2)
				..controls(116.1,-105)and(116.5,-105.5)..(116.9,-106)
				..controls(117.3,-106.4)and(117.8,-106.9)..(118,-107.3)
				..controls(118.2,-107.6)and(118.6,-108)..(119,-108.3)
				..controls(119.6,-108.8)and(121.1,-110.3)..(121.8,-111.2)
				..controls(122.1,-111.7)and(122.5,-112)..(123.4,-112.4)
				..controls(124.8,-113.2)and(125.8,-113.8)..(127,-114.6)
				..controls(127.5,-114.9)and(128.1,-115.3)..(128.2,-115.4)
				..controls(128.4,-115.5)and(129.3,-115.8)..(130.2,-116)
				..controls(131.1,-116.2)and(132.5,-116.6)..(133.3,-116.9)
				..controls(135.7,-117.9)and(136.1,-117.9)..(142.5,-117.8)
				..controls(145.8,-117.8)and(148.4,-117.7)..(148.8,-117.6)
				..controls(149.2,-117.6)and(150.2,-117.2)..(151,-116.8)
				..controls(152,-116.4)and(152.8,-116)..(153.6,-115.9)
				..controls(154.5,-115.7)and(155.2,-115.4)..(156.3,-114.9)
				..controls(157.1,-114.5)and(158.2,-114)..(158.7,-113.8)
				..controls(159.1,-113.7)and(159.8,-113.3)..(160.2,-113)
				..controls(160.6,-112.7)and(161.3,-112.4)..(161.7,-112.2)
				..controls(162.7,-111.9)and(163.2,-111.7)..(164.1,-111.1)
				..controls(164.5,-110.8)and(165.3,-110.4)..(165.9,-110)
				..controls(166.5,-109.7)and(167.3,-109.2)..(167.6,-108.9)
				..controls(167.9,-108.6)and(168.6,-108.2)..(169,-108)
				..controls(169.4,-107.8)and(170.2,-107.3)..(170.7,-106.9)
				..controls(171.2,-106.5)and(172.1,-105.8)..(172.7,-105.4)
				..controls(173.2,-105)and(174.1,-104.4)..(174.5,-104)
				..controls(174.9,-103.7)and(175.7,-103.1)..(176.2,-102.6)
				..controls(176.7,-102.2)and(177.4,-101.5)..(177.9,-101.2)
				..controls(178.3,-100.8)and(178.8,-100.4)..(179,-100.2)
				..controls(179.2,-100)and(179.6,-99.6)..(180,-99.3)
				..controls(180.4,-99.1)and(180.9,-98.6)..(181.2,-98.4)
				..controls(181.5,-98.1)and(182.1,-97.7)..(182.5,-97.5)
				..controls(183.2,-97.2)and(183.5,-96.9)..(184.3,-95.9)
				..controls(184.6,-95.6)and(184.6,-95.6)..(184.8,-95.7)
				..controls(185,-95.8)and(185.1,-96)..(185.1,-96.2)
				..controls(185.1,-96.5)and(185,-96.7)..(184.6,-97.2)
				..controls(184.1,-97.8)and(184,-97.9)..(183.9,-98.8)
				..controls(183.7,-100.3)and(183.3,-101.3)..(182.9,-101.8)
				..controls(182.6,-102.1)and(182.3,-102.6)..(182.2,-102.9)
				..controls(182,-103.3)and(181.6,-104)..(181.2,-104.6)
				..controls(180.9,-105.2)and(180.5,-106.1)..(180.4,-106.6)
				..controls(179.9,-108)and(179.7,-108.5)..(179,-109.4)
				..controls(178.6,-109.8)and(178.2,-110.5)..(177.9,-111)
				..controls(177.7,-111.5)and(177.3,-112.3)..(177,-112.7)
				..controls(176.6,-113.2)and(176.4,-113.7)..(176.3,-114.2)
				..controls(176.2,-114.7)and(175.9,-115.2)..(175.4,-116)
				..controls(175,-116.6)and(174.4,-117.7)..(174,-118.5)
				..controls(173.7,-119.2)and(173.3,-120.1)..(173,-120.5)
				..controls(172.8,-120.9)and(172.4,-121.8)..(172.1,-122.5)
				..controls(171.8,-123.2)and(171.3,-124.1)..(171,-124.6)
				..controls(170.8,-125.1)and(170.3,-126)..(169.9,-126.6)
				..controls(169.6,-127.3)and(169.2,-128.1)..(169,-128.3)
				..controls(168.8,-128.6)and(168.5,-129.3)..(168.3,-129.8)
				..controls(168.1,-130.3)and(167.8,-131.1)..(167.5,-131.6)
				..controls(166.9,-132.5)and(164.2,-137.9)..(164.2,-138.1)
				..controls(164.2,-138.1)and(163.9,-138.8)..(163.4,-139.7)
				..controls(162.9,-140.5)and(162.4,-141.5)..(162.3,-141.9)
				..controls(161.9,-143)and(161.6,-143.4)..(160.9,-144.1)
				..controls(160.4,-144.6)and(160.3,-144.8)..(160.1,-145.4)
				..controls(160,-145.8)and(159.6,-146.8)..(159.2,-147.8)
				..controls(158.7,-148.8)and(158.1,-150)..(157.9,-150.6)
				..controls(157.6,-151.1)and(157.2,-151.9)..(156.9,-152.4)
				..controls(156.5,-152.8)and(156.1,-153.7)..(155.8,-154.4)
				..controls(155.6,-155)and(155.2,-155.9)..(154.9,-156.3)
				..controls(154.7,-156.7)and(154.4,-157.4)..(154.2,-157.9)
				..controls(154.1,-158.4)and(153.8,-159)..(153.5,-159.3)
				..controls(152.8,-160.1)and(153,-160.1)..(142.4,-160.1)
				..controls(137.2,-160.1)and(131.8,-160.1)..(130.4,-160)--cycle
				(141.5,-103.7)..controls(140.8,-103)and(139.8,-102.1)..(139.2,-101.6)
				..controls(138.5,-101.2)and(137.9,-100.6)..(137.8,-100.4)
				..controls(137,-99.5)and(134,-96.5)..(133.2,-96)
				..controls(132.7,-95.7)and(131.8,-94.9)..(131.1,-94.3)
				..controls(129.4,-92.7)and(127.9,-91.6)..(126.6,-91.1)
				..controls(125.2,-90.5)and(124.1,-90)..(122.9,-89.1)
				..controls(122.1,-88.6)and(121.3,-88.2)..(119.9,-87.7)
				..controls(118.8,-87.3)and(117.5,-86.8)..(116.9,-86.5)
				..controls(114.3,-85.5)and(113.7,-85.4)..(109.7,-85.5)
				..controls(105.3,-85.7)and(103.9,-85.9)..(101.6,-87)
				..controls(100.7,-87.4)and(99.4,-87.8)..(97.9,-88)
				..controls(97,-88.2)and(96.6,-88.3)..(95.8,-88.9)
				..controls(95.4,-89.2)and(94.7,-89.6)..(94.2,-89.8)
				..controls(92.9,-90.3)and(92.3,-90.7)..(90.9,-92.1)
				..controls(90.2,-92.7)and(89.4,-93.3)..(89,-93.5)..controls(87,-94.5)and(85.2,-96)..(83.6,-97.9)
				..controls(83.3,-98.2)and(82.5,-98.9)..(81.8,-99.4)
				..controls(80.7,-100.1)and(80.5,-100.3)..(80,-101.1)
				..controls(78.7,-103.1)and(77.2,-103.9)..(76.4,-103)
				..controls(75.7,-102.2)and(75.7,-103.2)..(75.7,-83.6)
				..controls(75.6,-73.6)and(75.6,-64.8)..(75.5,-64)
				..controls(75.5,-63.2)and(75.5,-61.4)..(75.6,-59.8)--(75.7,-57.1)--(76,-56.8)
				..controls(76.8,-56.2)and(78,-56.3)..(79.8,-57.2)
				..controls(80.4,-57.5)and(81.3,-57.8)..(81.8,-57.9)
				..controls(82.3,-58.1)and(83,-58.4)..(83.6,-58.8)..controls(84.3,-59.2)and(85,-59.5)..(86,-59.8)
				..controls(86.8,-60.1)and(88,-60.6)..(88.7,-60.9)
				..controls(89.3,-61.2)and(90.6,-61.7)..(91.4,-62)
				..controls(92.2,-62.3)and(93.7,-62.9)..(94.6,-63.3)
				..controls(95.6,-63.7)and(96.8,-64.1)..(97.4,-64.3)
				..controls(98,-64.5)and(99,-64.9)..(99.7,-65.2)
				..controls(100.4,-65.6)and(101.5,-66)..(102.3,-66.2)
				..controls(104.2,-66.8)and(105,-67.1)..(106.2,-67.8)
				..controls(107.9,-68.7)and(110.2,-69.4)..(111.3,-69.2)
				..controls(112.5,-69)and(115,-68.2)..(115.6,-67.8)
				..controls(117.3,-66.7)and(117.6,-66.6)..(119.5,-66.2)
				..controls(121.9,-65.8)and(122.1,-65.7)..(123,-65.1)
				..controls(124.1,-64.4)and(124.4,-64.3)..(126,-63.9)
				..controls(127,-63.7)and(127.7,-63.4)..(128.4,-63.1)
				..controls(129,-62.8)and(129.7,-62.5)..(130.6,-62.4)
				..controls(132.1,-62.1)and(134,-61.5)..(135.2,-61)
				..controls(135.7,-60.8)and(136.7,-60.5)..(137.3,-60.3)
				..controls(138.4,-60)and(139.3,-59.7)..(141.1,-59)
				..controls(141.6,-58.7)and(142,-58.6)..(142.5,-58.6)
				..controls(143.2,-58.6)and(143.3,-58.7)..(143.6,-59)
				..controls(144,-59.5)and(144.1,-59.9)..(144,-60.8)
				..controls(144,-61.2)and(144,-71.1)..(144,-82.8)--(143.9,-104.1)--(143.7,-104.5)
				..controls(143.5,-104.7)and(143.3,-104.8)..(143,-104.9)
				..controls(142.6,-104.9)and(142.5,-104.8)..(141.5,-103.7)--cycle
				(27.8,-91.4)..controls(27.1,-90.9)and(26.3,-90)..(26.1,-89.5)
				..controls(26,-89.2)and(25.8,-88.6)..(25.7,-87.9)
				..controls(25.4,-86.2)and(25.6,-84.8)..(26.5,-83.5)
				..controls(27.4,-82.3)and(28.3,-82)..(29.4,-82.3)
				..controls(31.1,-82.8)and(32.1,-84.9)..(32,-87.4)
				..controls(31.9,-89.3)and(31.2,-90.6)..(29.9,-91.4)
				..controls(29.1,-91.9)and(28.6,-91.9)..(27.8,-91.4)--cycle
				(190.1,-91.3)..controls(189.7,-91.2)and(189.3,-90.9)..(189.1,-90.7)
				..controls(188.3,-89.8)and(187.6,-87.7)..(187.6,-86.2)
				..controls(187.6,-85.5)and(187.7,-85.2)..(188.2,-84.4)
				..controls(189.1,-82.6)and(190.1,-81.8)..(191.3,-82.1)
				..controls(192,-82.3)and(193.5,-83.7)..(193.8,-84.4)
				..controls(194.1,-85.1)and(194.2,-87.1)..(194,-88.3)
				..controls(193.8,-89.8)and(193.3,-90.5)..(192.3,-91.3)
				..controls(191.7,-91.7)and(190.9,-91.7)..(190.1,-91.3)--cycle
				(173.7,-83.4)..controls(172.2,-82.5)and(172,-82.5)..(170.6,-82.4)
				..controls(169.7,-82.4)and(168.8,-82.4)..(168.3,-82.5)
				..controls(167.2,-82.7)and(166.9,-82.7)..(166.6,-82.2)
				..controls(166.3,-81.8)and(166.3,-81)..(166.8,-80)
				..controls(167.2,-78.9)and(167.4,-77.4)..(167.7,-74.1)
				..controls(167.8,-72.8)and(167.9,-71.8)..(167.9,-71.7)
				..controls(168,-71.7)and(168,-70.2)..(168,-68.4)..controls(168,-65.4)and(168,-64.9)..(167.7,-63)
				..controls(167.3,-60.5)and(167.1,-59.9)..(166.4,-58.4)
				..controls(166.1,-57.7)and(165.8,-56.9)..(165.7,-56.6)
				..controls(165.7,-56.2)and(165.4,-55.6)..(165.2,-55.2)
				..controls(164.6,-53.8)and(161.9,-50.7)..(161,-50.3)
				..controls(159.8,-49.7)and(157.9,-50)..(155.2,-51.2)
				..controls(154.7,-51.4)and(153.4,-51.8)..(152.2,-52.1)
				..controls(150.8,-52.4)and(149.4,-52.9)..(148.6,-53.2)
				..controls(147.8,-53.6)and(146.7,-54)..(146,-54.2)
				..controls(145.4,-54.4)and(144.5,-54.7)..(144,-54.9)
				..controls(143.6,-55.2)and(142.7,-55.5)..(141.8,-55.7)
				..controls(141,-55.9)and(139.6,-56.3)..(138.7,-56.6)
				..controls(137.8,-57)and(136.3,-57.4)..(135.3,-57.6)
				..controls(133.8,-57.9)and(133.2,-58.1)..(132.1,-58.7)
				..controls(131.1,-59.2)and(130.3,-59.4)..(129.4,-59.6)
				..controls(128.7,-59.8)and(127.8,-60)..(127.4,-60.1)
				..controls(127,-60.3)and(126,-60.5)..(125.1,-60.7)
				..controls(123.2,-61.1)and(121.7,-61.5)..(120.2,-62.1)
				..controls(119.7,-62.3)and(118.1,-62.8)..(116.7,-63.2)
				..controls(114.6,-63.9)and(114,-64.1)..(113,-64.7)
				..controls(111.9,-65.4)and(111.8,-65.4)..(111,-65.5)
				..controls(109.9,-65.5)and(109.3,-65.4)..(107.7,-64.7)
				..controls(107,-64.4)and(106,-64)..(105.3,-63.9)
				..controls(103.5,-63.5)and(102.2,-63)..(101.2,-62.5)
				..controls(100.1,-61.8)and(99.3,-61.5)..(97.8,-61.1)
				..controls(97.2,-61)and(96.2,-60.7)..(95.7,-60.4)
				..controls(95.1,-60.1)and(94.2,-59.8)..(93.5,-59.6)
				..controls(92.9,-59.4)and(91.7,-59)..(91,-58.7)..controls(90.2,-58.4)and(88.9,-58)..(88,-57.7)
				..controls(87.2,-57.5)and(86.1,-57.1)..(85.5,-56.8)
				..controls(85,-56.6)and(84.1,-56.2)..(83.5,-56)..controls(82.9,-55.7)and(81.8,-55.2)..(81,-54.9)
				..controls(80.1,-54.5)and(79.1,-54.2)..(78.7,-54.1)
				..controls(77.9,-53.8)and(77.3,-53.6)..(76.2,-52.9)
				..controls(75.7,-52.6)and(75.2,-52.4)..(74.1,-52.2)
				..controls(72.5,-51.9)and(72.4,-51.8)..(71.3,-51.1)
				..controls(69.9,-50.1)and(69.4,-49.9)..(67.9,-49.8)
				..controls(67,-49.7)and(66.6,-49.6)..(66.3,-49.4)
				..controls(66.1,-49.3)and(65.7,-49)..(65.4,-48.8)
				..controls(65.1,-48.7)and(64.2,-48.4)..(63.6,-48.2)
				..controls(62.9,-48)and(61.7,-47.6)..(60.9,-47.2)
				..controls(60.1,-46.8)and(58.9,-46.4)..(58.2,-46.2)
				..controls(57.6,-46)and(56.8,-45.7)..(56.5,-45.5)
				..controls(54.9,-44.6)and(54.6,-44.5)..(53.2,-44.2)
				..controls(51.7,-43.8)and(51,-43.6)..(49.9,-42.9)
				..controls(49.5,-42.6)and(48.9,-42.4)..(47.5,-42)
				..controls(45.7,-41.6)and(45.6,-41.5)..(44.3,-40.6)
				..controls(43.5,-40.1)and(42.8,-39.6)..(42.7,-39.4)
				..controls(42.4,-38.7)and(43,-38.1)..(44.5,-37.6)..controls(44.9,-37.5)and(46,-37)..(46.9,-36.6)
				..controls(47.9,-36.1)and(48.8,-35.8)..(49.6,-35.6)
				..controls(50.5,-35.4)and(51.3,-35.2)..(52.4,-34.6)
				..controls(53.3,-34.1)and(54.5,-33.7)..(54.9,-33.5)
				..controls(55.4,-33.4)and(56.4,-33)..(57.2,-32.6)
				..controls(58.6,-32)and(60.1,-31.5)..(61.6,-31.2)
				..controls(62.3,-31.1)and(62.7,-31)..(63.2,-30.7)
				..controls(64.1,-30)and(64.3,-29.9)..(65.3,-29.7)
				..controls(65.7,-29.6)and(66.5,-29.2)..(67.1,-28.9)
				..controls(67.6,-28.6)and(68.8,-28.2)..(69.6,-27.9)
				..controls(70.5,-27.7)and(71.6,-27.3)..(72.1,-27)
				..controls(72.6,-26.7)and(73.8,-26.2)..(74.8,-25.9)
				..controls(75.7,-25.5)and(76.7,-25.1)..(77.1,-25)
				..controls(77.4,-24.8)and(78.5,-24.4)..(79.6,-24.1)
				..controls(80.6,-23.8)and(82,-23.2)..(82.8,-22.9)
				..controls(83.5,-22.5)and(84.8,-22)..(85.8,-21.8)
				..controls(86.7,-21.5)and(87.8,-21.1)..(88.1,-20.9)
				..controls(88.4,-20.7)and(89.3,-20.3)..(90.1,-20.1)
				..controls(90.8,-19.8)and(91.9,-19.4)..(92.4,-19.1)
				..controls(93.5,-18.5)and(94.4,-18.2)..(95.4,-18.1)
				..controls(95.7,-18)and(96.4,-17.8)..(96.9,-17.5)
				..controls(97.5,-17.3)and(98.4,-16.9)..(99,-16.6)
				..controls(99.7,-16.3)and(100.9,-15.9)..(101.9,-15.6)
				..controls(102.9,-15.3)and(104.4,-14.7)..(105.3,-14.3)
				..controls(108.1,-13.2)and(108.5,-13.1)..(110.1,-13.1)
				..controls(111.4,-13.1)and(111.5,-13.1)..(112.4,-13.5)
				..controls(113.2,-14)and(115,-14.8)..(116.6,-15.4)
				..controls(116.9,-15.5)and(117.9,-15.8)..(118.8,-16.1)
				..controls(120.9,-16.6)and(121.5,-16.9)..(122.6,-17.9)
				..controls(123,-18.3)and(123.7,-18.8)..(123.9,-18.9)
				..controls(124.6,-19.4)and(126,-19.8)..(127.9,-20.2)
				..controls(129.6,-20.5)and(129.8,-20.6)..(131.1,-21.4)
				..controls(131.5,-21.7)and(132.4,-22.1)..(133.1,-22.3)
				..controls(136.1,-23.2)and(138.3,-24.1)..(140.2,-25)
				..controls(141,-25.5)and(141.5,-25.6)..(142.5,-25.8)
				..controls(145.1,-26.3)and(146,-26.6)..(149.4,-28.1)
				..controls(151.3,-28.9)and(153.5,-29.8)..(154.3,-30.1)
				..controls(155.1,-30.4)and(156.2,-30.8)..(156.6,-31)
				..controls(157.1,-31.2)and(158.2,-31.6)..(159.1,-31.8)
				..controls(160.6,-32.2)and(160.9,-32.3)..(161.7,-32.9)
				..controls(162.4,-33.3)and(162.9,-33.6)..(163.7,-33.8)
				..controls(164.3,-34)and(164.8,-34.1)..(164.9,-34.1)
				..controls(165,-34.1)and(165.6,-34.4)..(166.4,-34.8)
				..controls(167.3,-35.2)and(168.1,-35.4)..(169,-35.6)
				..controls(170,-35.8)and(170.5,-36)..(171.7,-36.6)
				..controls(172.4,-37)and(173.3,-37.4)..(173.6,-37.4)
				..controls(175.6,-37.7)and(176.2,-38)..(176.7,-39.1)
				..controls(176.9,-39.6)and(176.9,-39.7)..(176.7,-40.4)
				..controls(176.4,-41.6)and(175.8,-42.7)..(175.2,-43.4)
				..controls(174.5,-44.1)and(172.6,-45.5)..(171.7,-45.8)
				..controls(170.6,-46.2)and(169.1,-46)..(168.2,-45.4)
				..controls(167.9,-45.2)and(167.3,-44.8)..(166.9,-44.4)
				..controls(166.4,-44.1)and(165.8,-43.6)..(165.4,-43.2)
				..controls(165.1,-42.9)and(164.4,-42.3)..(163.9,-41.9)
				..controls(163.3,-41.5)and(162.7,-40.9)..(162.4,-40.5)
				..controls(161.4,-39.2)and(159.9,-38)..(158.8,-37.8)
				..controls(158.4,-37.8)and(158.3,-37.7)..(158.2,-37.2)
				..controls(158,-36.8)and(157.4,-36.1)..(157.1,-36.1)
				..controls(157,-36.1)and(156.7,-36.1)..(156.5,-36)
				..controls(156.1,-35.9)and(155.9,-35.7)..(155.7,-35.2)
				..controls(155.3,-34.6)and(155.1,-34.4)..(154,-34.1)
				..controls(153.6,-34)and(153.2,-33.7)..(152.9,-33.5)
				..controls(152.4,-32.9)and(151.6,-32.4)..(151.3,-32.4)
				..controls(150.9,-32.4)and(150.4,-32.6)..(150.4,-32.9)
				..controls(150.4,-33.2)and(150.8,-33.4)..(151.4,-33.6)
				..controls(152.1,-33.7)and(152.8,-34.1)..(153.5,-34.9)
				..controls(154.1,-35.5)and(154.3,-35.6)..(155,-35.8)
				..controls(155.7,-36)and(155.7,-36)..(156,-36.6)
				..controls(156.3,-37.3)and(156.7,-37.7)..(157.5,-37.8)
				..controls(158.1,-37.9)and(158.2,-38)..(158.4,-38.8)
				..controls(158.5,-39.4)and(158.6,-39.6)..(159.6,-40.4)
				..controls(160.2,-41)and(160.9,-41.6)..(161.3,-41.8)
				..controls(161.7,-42)and(161.9,-42.3)..(162.2,-42.8)
				..controls(162.4,-43.2)and(162.8,-43.8)..(163.1,-44.1)
				..controls(163.9,-44.8)and(165.1,-46.7)..(165.7,-48)
				..controls(166.2,-49.3)and(166.4,-49.7)..(167.2,-50.5)
				..controls(167.5,-50.8)and(167.9,-51.4)..(168.1,-51.8)
				..controls(168.3,-52.2)and(168.8,-53)..(169.1,-53.6)
				..controls(169.6,-54.3)and(169.8,-54.9)..(169.9,-55.6)
				..controls(170.4,-57.8)and(170.5,-58.1)..(171.2,-59.5)
				..controls(171.6,-60.3)and(172,-61.5)..(172.2,-62.1)
				..controls(172.3,-62.8)and(172.7,-63.7)..(172.9,-64.1)
				..controls(173.2,-64.6)and(173.4,-65.3)..(173.4,-65.6)
				..controls(173.8,-68.8)and(174.2,-70.6)..(174.9,-72.6)
				..controls(175.5,-74.5)and(175.9,-77.1)..(176.1,-80.2)
				..controls(176.3,-83.4)and(175.5,-84.4)..(173.7,-83.4)--cycle;
				\definecolor{ezs}{RGB}{247,179,26}\fill[ezs](-5.7,-156.6)--(-3.4,-157.5)--(61.9,12.7)--(59.7,13.6)--cycle;
				\definecolor{ezs}{RGB}{247,179,22}\fill[ezs](-3.8,-157.3)--(-1.5,-158.2)--(63.9,12)--(61.6,12.9)--cycle;
				\definecolor{ezs}{RGB}{248,179,19}\fill[ezs](-1.8,-158.1)--(0.4,-158.9)--(65.8,11.2)--(63.5,12.1)--cycle;
				\definecolor{ezs}{RGB}{248,178,16}\fill[ezs](0.1,-158.8)--(2.4,-159.7)--(67.7,10.5)--(65.4,11.4)--cycle;
				\definecolor{ezs}{RGB}{249,178,15}\fill[ezs](2,-159.5)--(4.3,-160.4)--(69.6,9.8)--(67.3,10.7)--cycle;
				\definecolor{ezs}{RGB}{251,177,17}\fill[ezs](3.9,-160.3)--(6.2,-161.1)--(71.5,9)--(69.2,9.9)--cycle;
				\definecolor{ezs}{RGB}{252,179,18}\fill[ezs](5.8,-161)--(8.1,-161.9)--(73.4,8.3)--(71.1,9.2)--cycle;
				\definecolor{ezs}{RGB}{252,182,17}\fill[ezs](7.7,-161.7)--(10,-162.6)--(75.3,7.6)--(73,8.5)--cycle;
				\definecolor{ezs}{RGB}{252,185,17}\fill[ezs](9.6,-162.5)--(11.9,-163.4)--(77.2,6.8)--(75,7.7)--cycle;
				\definecolor{ezs}{RGB}{251,189,18}\fill[ezs](11.5,-163.2)--(13.8,-164.1)--(79.2,6.1)--(76.9,7)--cycle;
				\definecolor{ezs}{RGB}{251,191,22}\fill[ezs](13.4,-163.9)--(15.7,-164.8)--(81.1,5.4)--(78.8,6.2)--cycle;
				\definecolor{ezs}{RGB}{252,193,25}\fill[ezs](15.4,-164.7)--(17.6,-165.6)--(83,4.6)--(80.7,5.5)--cycle;
				\definecolor{ezs}{RGB}{253,195,27}\fill[ezs](17.3,-165.4)--(19.6,-166.3)--(84.9,3.9)--(82.6,4.8)--cycle;
				\definecolor{ezs}{RGB}{253,196,31}\fill[ezs](19.2,-166.1)--(21.5,-167)--(86.8,3.2)--(84.5,4)--cycle;
				\definecolor{ezs}{RGB}{254,199,35}\fill[ezs](21.1,-166.9)--(23.4,-167.8)--(88.7,2.4)--(86.4,3.3)--cycle;
				\definecolor{ezs}{RGB}{253,201,38}\fill[ezs](23,-167.6)--(25.3,-168.5)--(90.6,1.7)--(88.3,2.6)--cycle;
				\definecolor{ezs}{RGB}{253,203,40}\fill[ezs](24.9,-168.3)--(27.2,-169.2)--(92.5,1)--(90.3,1.8)--cycle;
				\definecolor{ezs}{RGB}{253,205,42}\fill[ezs](26.8,-169.1)--(29.1,-170)--(94.4,0.2)--(92.2,1.1)--cycle;
				\definecolor{ezs}{RGB}{253,207,45}\fill[ezs](28.7,-169.8)--(31,-170.7)--(96.4,-0.5)--(94.1,0.4)--cycle;
				\definecolor{ezs}{RGB}{253,210,49}\fill[ezs](30.7,-170.5)--(32.9,-171.4)--(98.3,-1.2)--(96,-0.4)--cycle;
				\definecolor{ezs}{RGB}{253,212,54}\fill[ezs](32.6,-171.3)--(34.9,-172.2)--(100.2,-2)--(97.9,-1.1)--cycle;
				\definecolor{ezs}{RGB}{253,214,59}\fill[ezs](34.5,-172)--(36.8,-172.9)--(102.1,-2.7)--(99.8,-1.8)--cycle;
				\definecolor{ezs}{RGB}{253,216,63}\fill[ezs](36.4,-172.7)--(38.7,-173.6)--(104,-3.4)--(101.7,-2.6)--cycle;
				\definecolor{ezs}{RGB}{253,218,65}\fill[ezs](38.3,-173.5)--(40.6,-174.4)--(105.9,-4.2)--(103.6,-3.3)--cycle;
				\definecolor{ezs}{RGB}{254,220,66}\fill[ezs](40.2,-174.2)--(42.5,-175.1)--(107.8,-4.9)--(105.5,-4)--cycle;
				\definecolor{ezs}{RGB}{254,222,68}\fill[ezs](42.1,-175)--(44.4,-175.8)--(109.7,-5.6)--(107.5,-4.8)--cycle;
				\definecolor{ezs}{RGB}{254,224,71}\fill[ezs](44,-175.7)--(46.3,-176.6)--(111.7,-6.4)--(109.4,-5.5)--cycle;
				\definecolor{ezs}{RGB}{254,226,75}\fill[ezs](46,-176.4)--(48.2,-177.3)--(113.6,-7.1)--(111.3,-6.2)--cycle;
				\definecolor{ezs}{RGB}{254,228,79}\fill[ezs](47.9,-177.2)--(50.1,-178)--(115.5,-7.8)--(113.2,-7)--cycle;
				\definecolor{ezs}{RGB}{254,230,82}\fill[ezs](49.8,-177.9)--(52.1,-178.8)--(117.4,-8.6)--(115.1,-7.7)--cycle;
				\definecolor{ezs}{RGB}{254,232,85}\fill[ezs](51.7,-178.6)--(54,-179.5)--(119.3,-9.3)--(117,-8.4)--cycle;
				\definecolor{ezs}{RGB}{254,234,88}\fill[ezs](53.6,-179.4)--(55.9,-180.2)--(121.2,-10)--(118.9,-9.2)--cycle;
				\definecolor{ezs}{RGB}{254,237,91}\fill[ezs](55.5,-180.1)--(57.8,-181)--(123.1,-10.8)--(120.8,-9.9)--cycle;
				\definecolor{ezs}{RGB}{254,239,94}\fill[ezs](57.4,-180.8)--(59.7,-181.7)--(125,-11.5)--(122.8,-10.6)--cycle;
				\definecolor{ezs}{RGB}{254,241,97}\fill[ezs](59.3,-181.6)--(61.6,-182.4)--(127,-12.2)--(124.7,-11.4)--cycle;
				\definecolor{ezs}{RGB}{254,243,101}\fill[ezs](61.2,-182.3)--(63.5,-183.2)--(128.9,-13)--(126.6,-12.1)--cycle;
				\definecolor{ezs}{RGB}{254,245,103}\fill[ezs](63.2,-183)--(65.4,-183.9)--(130.8,-13.7)--(128.5,-12.8)--cycle;
				\definecolor{ezs}{RGB}{254,245,105}\fill[ezs](65.1,-183.8)--(67.4,-184.6)--(132.7,-14.4)--(130.4,-13.6)--cycle;
				\definecolor{ezs}{RGB}{254,245,105}\fill[ezs](67,-184.5)--(69.3,-185.4)--(134.6,-15.2)--(132.3,-14.3)--cycle;
				\definecolor{ezs}{RGB}{254,244,103}\fill[ezs](68.9,-185.2)--(71.2,-186.1)--(136.5,-15.9)--(134.2,-15)--cycle;
				\definecolor{ezs}{RGB}{253,243,101}\fill[ezs](70.8,-186)--(73.1,-186.8)--(138.4,-16.6)--(136.1,-15.8)--cycle;
				\definecolor{ezs}{RGB}{253,241,98}\fill[ezs](72.7,-186.7)--(75,-187.6)--(140.3,-17.4)--(138,-16.5)--cycle;
				\definecolor{ezs}{RGB}{253,239,95}\fill[ezs](74.6,-187.4)--(76.9,-188.3)--(142.2,-18.1)--(140,-17.2)--cycle;
				\definecolor{ezs}{RGB}{253,236,92}\fill[ezs](76.5,-188.2)--(78.8,-189)--(144.2,-18.8)--(141.9,-18)--cycle;
				\definecolor{ezs}{RGB}{253,234,88}\fill[ezs](78.5,-188.9)--(80.7,-189.8)--(146.1,-19.6)--(143.8,-18.7)--cycle;
				\definecolor{ezs}{RGB}{253,232,85}\fill[ezs](80.4,-189.6)--(82.7,-190.5)--(148,-20.3)--(145.7,-19.4)--cycle;
				\definecolor{ezs}{RGB}{253,230,82}\fill[ezs](82.3,-190.4)--(84.6,-191.2)--(149.9,-21)--(147.6,-20.2)--cycle;
				\definecolor{ezs}{RGB}{253,228,78}\fill[ezs](84.2,-191.1)--(86.5,-192)--(151.8,-21.8)--(149.5,-20.9)--cycle;
				\definecolor{ezs}{RGB}{253,226,74}\fill[ezs](86.1,-191.8)--(88.4,-192.7)--(153.7,-22.5)--(151.4,-21.6)--cycle;
				\definecolor{ezs}{RGB}{253,224,70}\fill[ezs](88,-192.6)--(90.3,-193.4)--(155.6,-23.3)--(153.3,-22.4)--cycle;
				\definecolor{ezs}{RGB}{253,222,67}\fill[ezs](89.9,-193.3)--(92.2,-194.2)--(157.5,-24)--(155.3,-23.1)--cycle;
				\definecolor{ezs}{RGB}{253,220,65}\fill[ezs](91.8,-194)--(94.1,-194.9)--(159.5,-24.7)--(157.2,-23.8)--cycle;
				\definecolor{ezs}{RGB}{253,217,63}\fill[ezs](93.7,-194.8)--(96,-195.6)--(161.4,-25.5)--(159.1,-24.6)--cycle;
				\definecolor{ezs}{RGB}{253,215,60}\fill[ezs](95.7,-195.5)--(97.9,-196.4)--(163.3,-26.2)--(161,-25.3)--cycle;
				\definecolor{ezs}{RGB}{253,213,56}\fill[ezs](97.6,-196.2)--(99.9,-197.1)--(165.2,-26.9)--(162.9,-26)--cycle;
				\definecolor{ezs}{RGB}{252,211,53}\fill[ezs](99.5,-197)--(101.8,-197.8)--(167.1,-27.7)--(164.8,-26.8)--cycle;
				\definecolor{ezs}{RGB}{252,209,49}\fill[ezs](101.4,-197.7)--(103.7,-198.6)--(169,-28.4)--(166.7,-27.5)--cycle;
				\definecolor{ezs}{RGB}{252,207,45}\fill[ezs](103.3,-198.4)--(105.6,-199.3)--(170.9,-29.1)--(168.6,-28.2)--cycle;
				\definecolor{ezs}{RGB}{253,205,42}\fill[ezs](105.2,-199.2)--(107.5,-200)--(172.8,-29.9)--(170.6,-29)--cycle;
				\definecolor{ezs}{RGB}{253,202,40}\fill[ezs](107.1,-199.9)--(109.4,-200.8)--(174.7,-30.6)--(172.5,-29.7)--cycle;
				\definecolor{ezs}{RGB}{253,200,38}\fill[ezs](109,-200.6)--(111.3,-201.5)--(176.7,-31.3)--(174.4,-30.4)--cycle;
				\definecolor{ezs}{RGB}{253,198,35}\fill[ezs](111,-201.4)--(113.2,-202.2)--(178.6,-32.1)--(176.3,-31.2)--cycle;
				\definecolor{ezs}{RGB}{253,196,31}\fill[ezs](112.9,-202.1)--(115.2,-203)--(180.5,-32.8)--(178.2,-31.9)--cycle;
				\definecolor{ezs}{RGB}{253,194,28}\fill[ezs](114.8,-202.8)--(117.1,-203.7)--(182.4,-33.5)--(180.1,-32.6)--cycle;
				\definecolor{ezs}{RGB}{253,192,26}\fill[ezs](116.7,-203.6)--(119,-204.5)--(184.3,-34.3)--(182,-33.4)--cycle;
				\definecolor{ezs}{RGB}{253,189,24}\fill[ezs](118.6,-204.3)--(120.9,-205.2)--(186.2,-35)--(183.9,-34.1)--cycle;
				\definecolor{ezs}{RGB}{253,186,20}\fill[ezs](120.5,-205)--(122.8,-205.9)--(188.1,-35.7)--(185.8,-34.8)--cycle;
				\definecolor{ezs}{RGB}{253,184,17}\fill[ezs](122.4,-205.8)--(124.7,-206.7)--(190,-36.5)--(187.8,-35.6)--cycle;
				\definecolor{ezs}{RGB}{252,183,15}\fill[ezs](124.3,-206.5)--(126.6,-207.4)--(192,-37.2)--(189.7,-36.3)--cycle;
				\definecolor{ezs}{RGB}{252,181,14}\fill[ezs](126.3,-207.2)--(128.5,-208.1)--(193.9,-37.9)--(191.6,-37.1)--cycle;
				\definecolor{ezs}{RGB}{252,179,13}\fill[ezs](128.2,-208)--(130.4,-208.9)--(195.8,-38.7)--(193.5,-37.8)--cycle;
				\definecolor{ezs}{RGB}{252,178,12}\fill[ezs](130.1,-208.7)--(132.4,-209.6)--(197.7,-39.4)--(195.4,-38.5)--cycle;
				\definecolor{ezs}{RGB}{251,178,10}\fill[ezs](132,-209.4)--(134.3,-210.3)--(199.6,-40.1)--(197.3,-39.3)--cycle;
				\definecolor{ezs}{RGB}{251,178,10}\fill[ezs](133.9,-210.2)--(136.2,-211.1)--(201.5,-40.9)--(199.2,-40)--cycle;
				\definecolor{ezs}{RGB}{250,177,11}\fill[ezs](135.8,-210.9)--(138.1,-211.8)--(203.4,-41.6)--(201.1,-40.7)--cycle;
				\definecolor{ezs}{RGB}{250,178,11}\fill[ezs](137.7,-211.6)--(140,-212.5)--(205.3,-42.3)--(203.1,-41.5)--cycle;
				\definecolor{ezs}{RGB}{248,180,10}\fill[ezs](139.6,-212.4)--(141.9,-213.3)--(207.3,-43.1)--(205,-42.2)--cycle;
				\definecolor{ezs}{RGB}{247,181,9}\fill[ezs](141.5,-213.1)--(143.8,-214)--(209.2,-43.8)--(206.9,-42.9)--cycle;
				\definecolor{ezs}{RGB}{246,182,9}\fill[ezs](143.5,-213.8)--(145.7,-214.7)--(211.1,-44.5)--(208.8,-43.7)--cycle;
				\definecolor{ezs}{RGB}{246,182,10}\fill[ezs](145.4,-214.6)--(147.7,-215.5)--(213,-45.3)--(210.7,-44.4)--cycle;
			\end{scope}
			% chữ EDUZY (Times đậm, làm dày nét bằng 5 lớp lệch nhẹ; căn đúng khung chữ logo gốc)
			\foreach \dx/\dy in {-0.7/0,0.7/0,0/-0.3,0/0.3,0/0}
			\node[anchor=center,inner sep=0pt,outer sep=0pt,text=ezText,align=center] at (109.6+\dx,-193.6+\dy)
			{\resizebox{1.01cm}{0.22cm}{\fontencoding{T1}\fontfamily{ptm}\fontseries{b}\selectfont EDUZY}};
		\end{tikzpicture}%
	}%
}
%====================== THIẾT LẬP THÂN ĐỀ ======================
%--- Môi trường tự luận {bt}: ex_test 3.3.4 chưa khai báo sẵn nên tạo ở đây
%    (ex_test đã tự móc \fixshowans và \listenumerate cho bt)
\makeatletter
\@ifundefined{bt}{\newtheorem{bt}{\color{gold1}\namebt}}{}
\makeatother
\renewcommand{\sectionmark}[1]{\markright{\thesection\ #1}}
%--- Khung tiêu đề (logo + tối đa 4 dòng; dòng 3, 4 để trống {} thì bỏ qua)
%  Đề thi:    \khungtieude{THI THỬ THPT QUỐC GIA NĂM 2026}{Môn Toán -- Đề số 21}
%                         {(Đề thi gồm: ...)}{Thời gian làm bài: 90 phút (không kể thời gian phát đề).}
%  Chuyên đề: \khungtieude{CHUYÊN ĐỀ 3 -- DÃY SỐ VÀ CÔNG THỨC TRUY HỒI}
%                         {Bồi dưỡng học sinh giỏi Toán THPT}{...}{}
\newcommand{\khungtieude}[4]{%
	\begin{tcolorbox}[colframe=gold1,colback=gold5!45,boxrule=0.5mm]
		\begin{minipage}[c]{3cm}
			\centerline{\eduzylogo[2.4cm]}
		\end{minipage}\hspace{0.5cm}
		\begin{minipage}[c]{12.5cm}
			\vskip 0.15 cm
			\centerline{\ind{#1}}
			\vskip 0.15 cm
			\centerline{\textbf{#2}}
			\if\relax\detokenize{#3}\relax\else
				\vskip 0.15 cm
				\centerline{\textit{#3}}
			\fi
			\if\relax\detokenize{#4}\relax\else
				\vskip 0.15 cm
				\centerline{\textit{#4}}
			\fi
			\vskip 0.15 cm
		\end{minipage}
	\end{tcolorbox}
}
%====================== BỔ SUNG CHO BÀI HỌC, CHUYÊN ĐỀ ======================
\usepackage[shortlabels]{enumitem}
\usepackage{tabularx}
%--- Danh sách gọn: \begin{enumerate}[gon] ... \end{enumerate}
\SetEnumitemKey{gon}{leftmargin=*,itemsep=1pt,parsep=0pt,topsep=3pt}
%--- Khung lý thuyết: \begin{khung}{} ... \end{khung}; tiêu đề chỉ dùng khi thật ngắn ("Mẹo", "Lỗi thường gặp")
\newtcolorbox{khung}[1]{breakable,enhanced,colframe=gold1,colback=gold5!45,
	boxrule=0.4mm,fonttitle=\bfseries\sffamily,coltitle=white,colbacktitle=gold1,
	attach boxed title to top left={xshift=6mm,yshift=-3mm},
	boxed title style={arc=1mm,boxrule=0.3mm},title={#1}}
%--- Tiêu đề mục nhỏ trong một dạng: \muc{Dấu hiệu}, \muc{Phương pháp}, \muc{Bài tự luyện}
\makeatletter
\newcommand{\muc}[1]{\par\medskip\noindent\indam{#1}\par\nobreak\smallskip\@afterindentfalse\@afterheading}
\makeatother
%--- Ô công thức cần nhớ (căn giữa, viền vàng, vừa khít nội dung)
%     \begin{center}\congthuc{$Q(n)=(n-1)!$}\end{center}
\newtcbox{\congthuc}{on line,colframe=gold1,colback=gold5!45,boxrule=0.5pt,arc=2pt,
	left=8pt,right=8pt,top=4pt,bottom=4pt}
%--- Công tắc lời giải bài tự luyện (ví dụ mẫu luôn dùng \loigiai)
%     \loigiaiTL[{đáp số}]{lời giải đầy đủ}
%     \TLfull -> in lời giải đầy đủ;  \TLdapso -> chỉ in một dòng "Đáp số";  \TLan -> ẩn hoàn toàn
%     (\CheDo là dethi thì mọi lời giải và đáp số đều ẩn)
\makeatletter
\newif\ifTLhien   \TLhientrue
\newif\ifTLdapso  \TLdapsofalse
\newcommand{\TLfull}{\TLhientrue\TLdapsofalse}
\newcommand{\TLdapso}{\TLhientrue\TLdapsotrue}
\newcommand{\TLan}{\TLhienfalse}
\newcommand{\loigiaiTL}[2][]{%
	\ifTLhien
		\ifTLdapso
			\def\@tlarg{#1}%
			\ifx\@tlarg\@empty\else
				\ifdefstring{\kindOPTN}{dethi}{}{\par\noindent{\bfseries\color{gold1}Đáp số:} #1\par}%
			\fi
		\else
			\loigiai{#2}%
		\fi
	\fi}
\makeatother
\TLdapso% mặc định: chỉ in đáp số bài tự luyện
%====================== BẢNG ĐÚNG/SAI (\choiceTFt) ======================
%  Hàng tiêu đề nền nâu đồng chữ trắng (Phát biểu | Đ | S), các hàng tô xen kẽ kem / trắng,
%  đường kẻ vàng kim, nhãn a) b) c) d) màu nâu đồng. Bản có đáp án: cột được chọn có vòng tròn
%  nâu đồng với dấu tích trắng (câu đúng tích ở cột Đúng, câu sai tích ở cột Sai).
%  Tiêu đề ghi đủ "Đúng | Sai": \OPTN{viettat=0} sau \input (cột Đúng/Sai rộng ra, phát biểu hẹp lại).
%  Dùng bảng gốc của ex_test: \def\BangDSgoc{1} trước \input{preamble-eduzy-v2}.
\providecommand{\BangDSgoc}{0}
\ifnum\BangDSgoc=0
\usepackage{colortbl,pifont}
\colorlet{ezDSke}{gold2!60!gold1}% đường kẻ bảng (nét 0,4pt như ex_test để không đổi chiều cao bảng)
\colorlet{ezDSle}{gold5!60}% hàng a), c)
\colorlet{ezDSchan}{white}% hàng b), d)
%--- dấu chọn ở cột Đúng/Sai
%    (bề rộng 0, căn giữa ô: không làm đổi bề rộng cột của ex_test, không báo tràn lề)
\newcommand{\EZdauchon}{\makebox[0pt]{\tikz[baseline=-1ex]\node[circle,fill=gold1,inner sep=0pt,
	minimum size=1.2em,text=white,font=\small]{\ding{51}};}}
\def\tickT{\EZdauchon}\def\tickF{\EZdauchon}
\AtBeginDocument{\sbox{\ccTFTab}{\EZdauchon}}
%--- nhãn a) b) c) d) màu nâu đồng (ở bản lời giải, nhãn câu đúng vẫn đóng khung như ex_test)
\renewcommand{\FalseTF}{\stepcounter{dapan}\textbf{\color{gold1}\DapAnTF\sepTF}\ignorespaces}
%--- trước khi dựng bảng: kẻ màu vàng kim (bề rộng cột giữ nguyên như ex_test); sau bảng trả lại màu đen
\newcount\EZhangTF
\newcommand{\EZbatdauTF}{%
	\global\EZhangTF=0\relax
	\arrayrulecolor{ezDSke}}
\newcommand{\EZketthucTF}{\arrayrulecolor{black}}
%--- hàng tiêu đề nền nâu đồng, chữ trắng
\makeatletter
\newcommand{\EZo@TF}[1]{\centering\bfseries\color{white}#1}
\newcommand{\EZtieudeTF}{%
	\ifcase\cotDS
		\ifnum\cotSTT=1\relax
			\def\rowHeader{\hline\rowcolor{gold1}\EZo@TF{TT}&\EZo@TF{\phatbieu}\tabularnewline}%
		\else
			\def\rowHeader{\hline\rowcolor{gold1}\EZo@TF{\phatbieu}\tabularnewline}%
		\fi
	\or
		\ifnum\cotSTT=1\relax
			\def\rowHeader{\hline\rowcolor{gold1}\EZo@TF{TT}&\EZo@TF{\phatbieu}&\EZo@TF{\Dung}\tabularnewline}%
		\else
			\def\rowHeader{\hline\rowcolor{gold1}\EZo@TF{\phatbieu}&\EZo@TF{\Dung}\tabularnewline}%
		\fi
	\or
		\ifnum\cotSTT=1\relax
			\def\rowHeader{\hline\rowcolor{gold1}\EZo@TF{TT}&\EZo@TF{\phatbieu}&\EZo@TF{\Dung}&\EZo@TF{\Sai}\tabularnewline}%
		\else
			\def\rowHeader{\hline\rowcolor{gold1}\EZo@TF{\phatbieu}&\EZo@TF{\Dung}&\EZo@TF{\Sai}\tabularnewline}%
		\fi
	\fi}
%--- tô xen kẽ từng hàng (gọi ngay sau \hline đầu mỗi hàng)
\newcommand{\EZmauhangTF}{%
	\noalign{\global\advance\EZhangTF\@ne}%
	\ifodd\EZhangTF\rowcolor{ezDSle}\else\rowcolor{ezDSchan}\fi}
%--- gắn vào \choiceTFt / \writeChoiceTFt của ex_test
\xpretocmd{\choiceTFt}{\EZbatdauTF}{}{\typeout{[preamble-eduzy-v2] Không gắn được bảng Đúng/Sai (đầu bảng)}}
\xpatchcmd{\choiceTFt}{\ifnum\dongPB=0\def\rowHeader{}\fi}{\EZtieudeTF\ifnum\dongPB=0\def\rowHeader{}\fi}
	{}{\typeout{[preamble-eduzy-v2] Không gắn được bảng Đúng/Sai (hàng tiêu đề)}}
\xpatchcmd{\writeChoiceTFt}{\hline\xdef\dotEX{\sdotEX}}{\hline\EZmauhangTF\xdef\dotEX{\sdotEX}}
	{}{\typeout{[preamble-eduzy-v2] Không gắn được bảng Đúng/Sai (tô hàng)}}
\xpatchcmd{\writeChoiceTFt}{{\end{longtable}}}{{\end{longtable}\EZketthucTF}}
	{}{\typeout{[preamble-eduzy-v2] Không gắn được bảng Đúng/Sai (cuối longtable)}}
\xpatchcmd{\writeChoiceTFt}{{\end{xtabular}\end{center}\medskip\noindent}}{{\end{xtabular}\end{center}\EZketthucTF\medskip\noindent}}
	{}{\typeout{[preamble-eduzy-v2] Không gắn được bảng Đúng/Sai (cuối xtabular)}}
\fi
%=====================================================================
\renewcommand{\parallel}{\mathrel{/\!\!/}}

%  trong file đề (tools/dan_trang/dan_trang.py tự ghi khi đặt \NeedspaceCau):
%    \def\TachCauDS{1}   câu Đúng/Sai được tách: câu dẫn (kèm hình) ở cuối trang, bảng Phát biểu sang trang sau
%                        (câu dẫn và hình không tách nhau, bảng không tách giữa hai trang)
%    \def\TachCauLG{1}   bản lời giải (\CheDo khác dethi): phần đề của câu cao hơn 1/3 trang được phép tách
%  Lệnh dùng trong file đề (giá trị do tools/dan_trang/dan_trang.py đo ở bản dethi):
%    \NeedspaceCau[<cao câu dẫn + hình>]{<cao cả khối câu>}   ngay trên \begin{ex} (câu đầu phần: trên tiêu đề phần)
%        công tắc tắt: đúng bằng \Needspace{<cao cả khối câu>} như cũ;
%        \TachCauDS=1 và có [..] (câu Đúng/Sai được chọn tách): chỉ giữ câu dẫn + hình;
%        \TachCauLG=1 ở bản lời giải: đòi tối đa 1/3 \textheight.
%    \NeedspaceBang{<cao bảng Phát biểu>}%   ngay trước \choiceTFt: \TachCauDS=1 thì giữ trọn bảng trên một trang;
%        công tắc tắt: không làm gì.
\providecommand{\TachCauDS}{0}
\providecommand{\TachCauLG}{0}
\newlength{\EZnsCau}
\newcommand{\NeedspaceCau}[2][]{%
	\setlength{\EZnsCau}{#2}%
	\ifnum\TachCauDS=1 \if\relax\detokenize{#1}\relax\else\setlength{\EZnsCau}{#1}\fi\fi
	\ifnum\TachCauLG=1 \ifdefstring{\kindOPTN}{dethi}{}{\ifdim\EZnsCau>0.3333\textheight\setlength{\EZnsCau}{0.3333\textheight}\fi}\fi
	\Needspace{\EZnsCau}}
\newcommand{\NeedspaceBang}[1]{\ifnum\TachCauDS=1 \Needspace{#1}\fi}

%--- Thang màu: đậm -> nhạt
\definecolor{gold1}{RGB}{126,96,0}    % nâu đồng đậm  #7E6000
\definecolor{gold2}{RGB}{192,144,3}   % vàng kim      #C09003
\definecolor{gold3}{RGB}{255,215,99}  % vàng nhạt     #FFD763
\definecolor{gold4}{RGB}{255,231,153} % vàng rất nhạt #FFE799
\definecolor{gold5}{RGB}{255,243,205} % vàng kem      #FFF3CD
%--- Ghi đè tông xanh/đỏ cũ
\definecolor{myblue}{RGB}{126,96,0}
\colorlet{blue}{gold1}
\colorlet{cyan}{gold2}
\colorlet{teal}{gold2}

%--- Tên chương mặc định (file đề nên tự khai báo lại)
\providecommand{\tenchuong}{}
%--- Chân trang
\providecommand{\thuonghieu}{TOÁN 12 -- Eduzy.com.vn}
\providecommand{\tacgia}{Tô Gia Bảo -- 0353703783}

%--- Chữ cạnh logo ở đầu trang
\providecommand{\khauhieu}{Eduzy.com.vn}

%====================== THÔNG SỐ ĐẦU TRANG, CHÂN TRANG (v2) ======================
%  Dải nền đậm tràn ngang trang; ô chữ vàng kem có cạnh xiên, kèm 2 vạch xiên chuyển màu.
%  Đầu trang lẻ : logo + \khauhieu (trái)   | ô chữ \tenchuong (phải)
%  Chân trang lẻ: ô chữ \thuonghieu (trái)  | Trang (số) (phải)
%  Trang chẵn soi gương trang lẻ; chân trang chẵn ghi \tacgia.
%--- Màu: chỉ dùng thang vàng kim Eduzy
\colorlet{ezNen}{gold1}% dải nền đậm
\colorlet{ezVach1}{gold3}% vạch xiên sát dải nền (đậm hơn)
\colorlet{ezVach2}{gold4}% vạch xiên sát ô chữ (nhạt hơn)
\colorlet{ezO}{gold5}% ô chữ vàng kem
%--- Kích thước
\newlength{\EZdai}\setlength{\EZdai}{1.1cm}% chiều cao dải nền
\newlength{\EZo}\setlength{\EZo}{0.68cm}% chiều cao ô chữ
\newlength{\EZxien}\setlength{\EZxien}{0.55cm}% độ lệch ngang của cạnh xiên
\newlength{\EZvach}\setlength{\EZvach}{0.75cm}% bề rộng mỗi vạch xiên
\newlength{\EZkhe}\setlength{\EZkhe}{0.12cm}% khe trắng giữa vạch và ô chữ
\newlength{\EZdem}\setlength{\EZdem}{0.4cm}% khoảng đệm từ chữ tới cạnh xiên
\newlength{\EZle}\setlength{\EZle}{1.8cm}% lề trái/phải của chữ (bằng lề thân bài)
\newlength{\EZrongmin}\setlength{\EZrongmin}{4.5cm}% ô chữ rộng tối thiểu
\newlength{\EZrong}% bề rộng chữ trong ô (tính tự động)
%--- Font
\newcommand{\EZfontO}{\fontsize{10.5pt}{12pt}\selectfont\bfseries}% chữ trong ô
\newcommand{\EZfontKH}{\fontfamily{qag}\fontsize{11pt}{12pt}\selectfont\bfseries}% \khauhieu cạnh logo
\newcommand{\EZfontTrang}{\fontsize{11.5pt}{12pt}\selectfont\bfseries}% chữ "Trang" và số trang
%--- Đo bề rộng chữ trong ô
\newcommand{\EZdorong}[1]{%
	\settowidth{\EZrong}{\EZfontO #1}%
	\ifdim\EZrong<\EZrongmin\setlength{\EZrong}{\EZrongmin}\fi}
%--- Logo + chữ trắng \khauhieu (dựng sẵn một lần khi vào \begin{document})
\newsavebox{\EZlogobox}
\AtBeginDocument{\sbox{\EZlogobox}{%
		\raisebox{-0.5\height}{\eduzylogo[1cm]}\hspace{0.15cm}%
		\raisebox{-0.5\height}{\color{white}\EZfontKH\khauhieu}}}
%--- Hai vạch xiên + khe trắng. #1: điểm gốc ở mép dưới, #2: hướng (1 sang phải, -1 sang trái),
%    #3: hướng nghiêng (1 kiểu "/", -1 kiểu "\")
\newcommand{\EZvachxien}[3]{%
	\fill[white] (#1) --++({(#3)*\EZxien},\EZo) --++({(#2)*\EZkhe},0) --++({-(#3)*\EZxien},-\EZo) --cycle;
	\fill[ezVach2] ([xshift=(#2)*\EZkhe]#1) --++({(#3)*\EZxien},\EZo) --++({(#2)*\EZvach},0) --++({-(#3)*\EZxien},-\EZo) --cycle;
	\fill[ezVach1] ([xshift=(#2)*(\EZkhe+\EZvach)]#1) --++({(#3)*\EZxien},\EZo) --++({(#2)*\EZvach},0) --++({-(#3)*\EZxien},-\EZo) --cycle;
	\draw[white,line width=0.6pt] ([xshift=(#2)*(\EZkhe+\EZvach)]#1) --++({(#3)*\EZxien},\EZo);
}
%--- Chữ trong ô (căn giữa theo chiều cao ô, không lệch vì dấu tiếng Việt)
\tikzset{ezchu/.style={inner sep=0pt,text=ezNen,text height=0.72em,text depth=0.22em}}

%====================== LOGO CHÌM GIỮA TRANG ======================
%  Vẽ cùng đầu trang nên nằm dưới chữ của thân bài, có trên mọi trang.
%  Tắt: \logochimtat (từ trang đó trở đi), bật lại: \logochimbat.
%  Đổi cỡ, độ đậm: khai báo trước \input, ví dụ \def\EZchimcao{9cm} \def\EZchimdo{0.1}
\providecommand{\EZchimcao}{10cm}% chiều cao logo chìm
\providecommand{\EZchimdo}{0.12}% độ đậm (0 = trong suốt, 1 = rõ hẳn)
\newif\ifEZchim \EZchimtrue
\newcommand{\logochimtat}{\EZchimfalse}
\newcommand{\logochimbat}{\EZchimtrue}
%  Dựng sẵn logo mờ một lần khi vào \begin{document}; transparency group làm mờ cả logo
%  một lượt nên không lộ vệt ở chỗ các dải màu chồng nhau.
\newsavebox{\EZchimtho}
\newsavebox{\EZchimbox}
\AtBeginDocument{%
	\sbox{\EZchimtho}{\eduzylogo[\EZchimcao]}%
	\sbox{\EZchimbox}{\begin{tikzpicture}
			\begin{scope}[transparency group,opacity=\EZchimdo]
				\node[inner sep=0pt] {\usebox{\EZchimtho}};
			\end{scope}
	\end{tikzpicture}}}
\newcommand{\EZvechim}{%
	\ifEZchim
		\node[inner sep=0pt] at (current page.center) {\usebox{\EZchimbox}};
	\fi}

%====================== ĐẦU TRANG - TRANG LẺ ======================
%  (logo bên trái, ô chữ \tenchuong bên phải)
\newcommand{\myfancyheadOdd}{%
	\EZdorong{\tenchuong}%
	\begin{tikzpicture}[remember picture,overlay]
		\coordinate (T) at (current page.north west);
		\coordinate (P) at (current page.north east);
		%--- logo chìm giữa trang
		\EZvechim
		%--- dải nền
		\fill[ezNen] (T) rectangle ([yshift=-\EZdai]P);
		%--- ô chữ: X là góc dưới bên trái, cạnh trái xiên kiểu "/"
		\coordinate (X) at ([xshift=-\EZle-\EZrong-\EZdem-\EZxien,yshift=-\EZdai]P);
		\fill[ezO] (X) --++(\EZxien,\EZo) -- ([yshift=-\EZdai+\EZo]P) -- ([yshift=-\EZdai]P) --cycle;
		\EZvachxien{X}{-1}{1}
		\node[ezchu,anchor=east] at ([xshift=-\EZle,yshift=-\EZdai+0.5*\EZo]P) {\EZfontO\tenchuong};
		%--- logo + tên
		\node[inner sep=0pt,anchor=west] at ([xshift=\EZle,yshift=-0.5*\EZdai]T) {\usebox{\EZlogobox}};
	\end{tikzpicture}%
}

%====================== ĐẦU TRANG - TRANG CHẴN =====================
%  (soi gương: ô chữ \tenchuong bên trái, logo bên phải)
\newcommand{\myfancyheadEven}{%
	\EZdorong{\tenchuong}%
	\begin{tikzpicture}[remember picture,overlay]
		\coordinate (T) at (current page.north west);
		\coordinate (P) at (current page.north east);
		%--- logo chìm giữa trang
		\EZvechim
		%--- dải nền
		\fill[ezNen] (T) rectangle ([yshift=-\EZdai]P);
		%--- ô chữ: X là góc dưới bên phải, cạnh phải xiên kiểu "\"
		\coordinate (X) at ([xshift=\EZle+\EZrong+\EZdem+\EZxien,yshift=-\EZdai]T);
		\fill[ezO] (X) --++(-\EZxien,\EZo) -- ([yshift=-\EZdai+\EZo]T) -- ([yshift=-\EZdai]T) --cycle;
		\EZvachxien{X}{1}{-1}
		\node[ezchu,anchor=west] at ([xshift=\EZle,yshift=-\EZdai+0.5*\EZo]T) {\EZfontO\tenchuong};
		%--- logo + tên
		\node[inner sep=0pt,anchor=east] at ([xshift=-\EZle,yshift=-0.5*\EZdai]P) {\usebox{\EZlogobox}};
	\end{tikzpicture}%
}

%====================== CHÂN TRANG - TRANG LẺ ======================
%  (ô chữ \thuonghieu bên trái, "Trang (số)" bên phải)
\newcommand{\myfancyfootOdd}{%
	\EZdorong{\thuonghieu}%
	\begin{tikzpicture}[remember picture,overlay]
		\coordinate (T) at (current page.south west);
		\coordinate (P) at (current page.south east);
		%--- dải nền
		\fill[ezNen] (T) rectangle ([yshift=\EZdai]P);
		%--- ô chữ sát mép trên dải: X là góc dưới bên phải, cạnh phải xiên kiểu "/"
		\coordinate (X) at ([xshift=\EZle+\EZrong+\EZdem,yshift=\EZdai-\EZo]T);
		\fill[ezO] ([yshift=\EZdai-\EZo]T) -- (X) --++(\EZxien,\EZo) -- ([yshift=\EZdai]T) --cycle;
		\EZvachxien{X}{1}{1}
		\node[ezchu,anchor=west] at ([xshift=\EZle,yshift=\EZdai-0.5*\EZo]T) {\EZfontO\thuonghieu};
		%--- số trang
		\node[circle,draw=white,line width=0.9pt,text=white,inner sep=0pt,minimum size=0.62cm,anchor=east]
		(so) at ([xshift=-\EZle,yshift=0.5*\EZdai]P) {\EZfontTrang\thepage};
		\node[inner sep=0pt,text=white,anchor=base east] (tr) at ([xshift=-0.15cm]so.base west) {\EZfontTrang Trang};
	\end{tikzpicture}%
}

%====================== CHÂN TRANG - TRANG CHẴN ====================
%  (soi gương: "Trang (số)" bên trái, ô chữ \tacgia bên phải)
\newcommand{\myfancyfootEven}{%
	\EZdorong{\tacgia}%
	\begin{tikzpicture}[remember picture,overlay]
		\coordinate (T) at (current page.south west);
		\coordinate (P) at (current page.south east);
		%--- dải nền
		\fill[ezNen] (T) rectangle ([yshift=\EZdai]P);
		%--- ô chữ sát mép trên dải: X là góc dưới bên trái, cạnh trái xiên kiểu "\"
		\coordinate (X) at ([xshift=-\EZle-\EZrong-\EZdem,yshift=\EZdai-\EZo]P);
		\fill[ezO] ([yshift=\EZdai-\EZo]P) -- (X) --++(-\EZxien,\EZo) -- ([yshift=\EZdai]P) --cycle;
		\EZvachxien{X}{-1}{-1}
		\node[ezchu,anchor=east] at ([xshift=-\EZle,yshift=\EZdai-0.5*\EZo]P) {\EZfontO\tacgia};
		%--- số trang
		%    (vòng tròn ẩn so0 nằm giữa dải để lấy đường chân chữ, giống trang lẻ)
		\node[circle,inner sep=0pt,minimum size=0.62cm] (so0) at ([yshift=0.5*\EZdai]T) {\EZfontTrang\phantom{\thepage}};
		\node[inner sep=0pt,text=white,anchor=base west] (tr) at ($(so0.base -| T)+(\EZle,0)$) {\EZfontTrang Trang};
		\node[circle,draw=white,line width=0.9pt,text=white,inner sep=0pt,minimum size=0.62cm,anchor=base west]
		(so) at ([xshift=0.15cm]tr.base east) {\EZfontTrang\thepage};
	\end{tikzpicture}%
}

%--- Giữ tương thích với các file cũ đang gọi \myfancyhead, \myfancyfoot
\newcommand{\myfancyhead}{\myfancyheadOdd}
\newcommand{\myfancyfoot}{\myfancyfootOdd}

%======================= GẮN VÀO TRANG =============================
\pagestyle{fancy}
\fancyhf{}
\fancyhead[LO]{\myfancyheadOdd}
\fancyhead[RE]{\myfancyheadEven}
\fancyfoot[LO]{\myfancyfootOdd}
\fancyfoot[RE]{\myfancyfootEven}
\renewcommand{\headrulewidth}{0pt}
\renewcommand{\footrulewidth}{0pt}
\fancypagestyle{plain}{%
	\fancyhf{}
	\fancyhead[LO]{\myfancyheadOdd}
	\fancyhead[RE]{\myfancyheadEven}
	\fancyfoot[LO]{\myfancyfootOdd}
	\fancyfoot[RE]{\myfancyfootEven}
	\renewcommand{\headrulewidth}{0pt}
	\renewcommand{\footrulewidth}{0pt}
}
%=================== HẾT PHẦN GỘP TỪ eduzy-de.tex ===================
\newcommand{\indam}[1]{
	{\fontsize{12pt}{9pt}\sffamily \selectfont\bfseries\color{gold1} #1}
}
\newcommand{\ind}[1]{
	{\fontsize{14pt}{9pt}\fontfamily{qag} \selectfont\bfseries\color{gold1} #1}
}
\renewcommand{\FalseEX}{\stepcounter{dapan}\textbf{\selectfont\color{gold1}\Alph{dapan}.}\ignorespaces}
\renewtheorem{ex}{\color{gold1}Câu}
\renewcommand{\baselinestretch}{1.1}
\newcommand{\hoac}[1]{
	\left[\begin{aligned}#1\end{aligned}\right.}
\newcommand{\heva}[1]{
	\left\{\begin{aligned}#1\end{aligned}\right.}
\usepackage{esvect}
\def\vec{\vv}
\def\vt{\vv}
\def\overrightarrow{\vv}
%====================== LOGO EDUZY VẼ BẰNG TIKZ ======================
%  Đường viền dò từ logo gốc; dải màu vàng kim đo trực tiếp trên ảnh gốc
%  (nghiêng 21 độ, tô bằng 80 dải hẹp nên hiển thị giống nhau trên mọi trình xem PDF).
%  Dùng: \eduzylogo          -> cao 2.4cm (mặc định)
%        \eduzylogo[1.5cm]   -> cao 1.5cm
\definecolor{ezText}{RGB}{245,200,46}
\newcommand{\eduzylogo}[1][2.4cm]{%
	\resizebox{!}{#1}{%
		\begin{tikzpicture}[x=0.01cm,y=0.01cm]
			\path[use as bounding box] (0,0) rectangle (224,-224);
			% mũ tốt nghiệp + vương miện
			\begin{scope}
				\clip
				(119.4,-176.1)..controls(118.9,-176.1)and(116.8,-176)..(114.6,-176)
				..controls(101.4,-176)and(85.8,-175.9)..(83.7,-175.7)
				..controls(81.8,-175.6)and(81,-175.6)..(80.2,-175.7)
				..controls(78.8,-175.9)and(77.7,-175.9)..(76.2,-175.7)
				..controls(74.9,-175.5)and(70.8,-175.5)..(64.3,-175.7)
				..controls(59.1,-175.9)and(58.3,-175.7)..(58.6,-173.7)
				..controls(58.6,-172.8)and(59.6,-169.6)..(60.3,-168)
				..controls(60.8,-166.9)and(61.4,-165.8)..(62.3,-165)
				..controls(63.1,-164.2)and(64.3,-163.4)..(65,-163.2)
				..controls(65.5,-163.1)and(72.3,-163)..(76.6,-163.1)
				..controls(77.9,-163.1)and(80.1,-163.1)..(81.4,-163)
				..controls(82.8,-162.9)and(86.3,-162.9)..(89.2,-162.9)
				..controls(93,-162.9)and(94.9,-162.8)..(96,-162.7)
				..controls(97.2,-162.5)and(101,-162.5)..(114,-162.4)
				..controls(123,-162.3)and(134.9,-162.2)..(140.2,-162.1)
				..controls(150.1,-162)and(152.2,-162.1)..(153.6,-162.4)
				..controls(154.4,-162.6)and(157.1,-163.8)..(158.1,-164.4)
				..controls(158.8,-164.8)and(159.7,-165.7)..(159.9,-166.3)
				..controls(160,-166.5)and(160.2,-167.1)..(160.2,-167.6)
				..controls(160.3,-168)and(160.6,-169.4)..(160.9,-170.6)
				..controls(161.2,-171.7)and(161.5,-173.1)..(161.6,-173.6)
				..controls(161.8,-174.4)and(161.8,-174.5)..(161.5,-174.9)
				..controls(161.2,-175.6)and(160.8,-175.9)..(159.9,-176)
				..controls(159.2,-176.1)and(154.1,-176.2)..(151.5,-176.2)
				..controls(148.1,-176.1)and(122.6,-176.1)..(121.6,-176.2)
				..controls(120.8,-176.2)and(119.8,-176.2)..(119.4,-176.1)--cycle
				(130.4,-160)..controls(129,-160)and(119.8,-160)..(110.1,-159.9)
				..controls(100.3,-159.8)and(88.5,-159.7)..(83.9,-159.6)
				..controls(77.8,-159.5)and(75.3,-159.5)..(74.5,-159.6)
				..controls(73.9,-159.7)and(72.5,-159.8)..(71.3,-159.8)
				..controls(69.6,-159.8)and(69,-159.8)..(68.5,-159.6)
				..controls(67.6,-159.3)and(66.8,-158.6)..(66.2,-157.6)
				..controls(66,-157.3)and(65.6,-156.5)..(65.3,-156.1)
				..controls(64.6,-155.1)and(64.5,-154.9)..(63.7,-152.9)
				..controls(63.2,-151.6)and(62.9,-150.9)..(61.8,-149.6)
				..controls(61.5,-149.2)and(61.1,-148.5)..(60.9,-147.9)
				..controls(60.7,-147.4)and(60.3,-146.6)..(60.1,-146.1)
				..controls(59.9,-145.5)and(59.5,-144.9)..(59,-144.3)
				..controls(58.5,-143.7)and(58.2,-143.2)..(57.9,-142.5)
				..controls(57.7,-141.9)and(57.3,-141.1)..(57,-140.7)
				..controls(56.6,-140.3)and(56.1,-139.3)..(55.8,-138.5)
				..controls(55.5,-137.8)and(55.1,-136.8)..(54.9,-136.5)
				..controls(54.7,-136.1)and(54.3,-135.2)..(53.9,-134.5)
				..controls(53.6,-133.8)and(53,-132.7)..(52.6,-132)
				..controls(51.6,-130.5)and(51.4,-130.2)..(51.1,-129)
				..controls(50.6,-127.5)and(50.6,-127.4)..(49.7,-126.3)
				..controls(49.2,-125.7)and(48.6,-124.8)..(48.3,-124.2)
				..controls(48,-123.6)and(47.5,-122.7)..(47.3,-122.2)
				..controls(47,-121.8)and(46.5,-120.7)..(46.1,-119.8)
				..controls(45.8,-118.9)and(45.2,-117.8)..(44.9,-117.2)
				..controls(44.2,-116.1)and(43.3,-114.4)..(42.3,-112.2)
				..controls(41.8,-111.3)and(41.3,-110.4)..(41.1,-110.1)
				..controls(40.5,-109.4)and(40.3,-109)..(40,-107.8)
				..controls(39.8,-107)and(39.6,-106.5)..(39.2,-106)
				..controls(38.9,-105.6)and(38.4,-104.7)..(38,-104)
				..controls(37.7,-103.4)and(37.2,-102.5)..(37,-102.1)
				..controls(36.7,-101.7)and(36.3,-100.8)..(36,-100.1)
				..controls(35.6,-98.9)and(35.3,-98.4)..(34.6,-97.6)
				..controls(34.2,-97.1)and(34.2,-96.3)..(34.6,-96)
				..controls(35,-95.7)and(35.5,-95.9)..(36.2,-96.7)
				..controls(36.5,-97)and(37.2,-97.5)..(37.6,-97.7)..controls(38,-98)and(38.6,-98.4)..(39.1,-98.9)
				..controls(39.6,-99.4)and(40.3,-100)..(40.6,-100.3)
				..controls(41.3,-100.8)and(42.4,-101.8)..(43.7,-102.9)
				..controls(44,-103.2)and(44.6,-103.7)..(45,-103.9)
				..controls(45.4,-104.2)and(45.9,-104.6)..(46.1,-104.8)
				..controls(47.1,-105.9)and(48.8,-107.1)..(50.5,-107.9)
				..controls(51.8,-108.6)and(52.3,-109)..(53.4,-110)
				..controls(54.8,-111.2)and(55.6,-111.6)..(57,-112.1)
				..controls(57.5,-112.3)and(58.5,-112.8)..(59.4,-113.3)
				..controls(60.2,-113.7)and(61.3,-114.3)..(61.9,-114.5)
				..controls(62.4,-114.7)and(63.3,-115.1)..(63.8,-115.3)
				..controls(64.4,-115.6)and(65.4,-115.9)..(66.2,-116.1)
				..controls(67,-116.3)and(68.4,-116.6)..(69.2,-116.9)
				..controls(72.3,-117.9)and(72.8,-118)..(78,-118.1)
				..controls(84,-118.3)and(86.1,-118.1)..(88.1,-117.1)
				..controls(89.4,-116.5)and(93,-115)..(93.7,-114.9)
				..controls(94.8,-114.6)and(95.4,-114.2)..(95.9,-113.4)
				..controls(96.2,-113)and(96.7,-112.5)..(97.1,-112.2)
				..controls(98,-111.6)and(101.3,-108.4)..(102.1,-107.2)
				..controls(102.4,-106.8)and(102.9,-106.2)..(103.2,-105.8)
				..controls(103.5,-105.5)and(104,-104.8)..(104.2,-104.2)
				..controls(104.5,-103.6)and(105,-102.8)..(105.3,-102.4)
				..controls(105.6,-102)and(106.1,-101.2)..(106.5,-100.6)
				..controls(107.1,-99.5)and(108.6,-97.5)..(109.5,-96.5)
				..controls(109.9,-96.1)and(110.1,-96)..(110.4,-96)
				..controls(111.1,-96)and(111.6,-96.5)..(111.8,-97.6)
				..controls(112,-98.6)and(112.3,-99.3)..(112.8,-99.9)
				..controls(113.7,-100.9)and(115.2,-103.1)..(115.7,-104.2)
				..controls(116.1,-105)and(116.5,-105.5)..(116.9,-106)
				..controls(117.3,-106.4)and(117.8,-106.9)..(118,-107.3)
				..controls(118.2,-107.6)and(118.6,-108)..(119,-108.3)
				..controls(119.6,-108.8)and(121.1,-110.3)..(121.8,-111.2)
				..controls(122.1,-111.7)and(122.5,-112)..(123.4,-112.4)
				..controls(124.8,-113.2)and(125.8,-113.8)..(127,-114.6)
				..controls(127.5,-114.9)and(128.1,-115.3)..(128.2,-115.4)
				..controls(128.4,-115.5)and(129.3,-115.8)..(130.2,-116)
				..controls(131.1,-116.2)and(132.5,-116.6)..(133.3,-116.9)
				..controls(135.7,-117.9)and(136.1,-117.9)..(142.5,-117.8)
				..controls(145.8,-117.8)and(148.4,-117.7)..(148.8,-117.6)
				..controls(149.2,-117.6)and(150.2,-117.2)..(151,-116.8)
				..controls(152,-116.4)and(152.8,-116)..(153.6,-115.9)
				..controls(154.5,-115.7)and(155.2,-115.4)..(156.3,-114.9)
				..controls(157.1,-114.5)and(158.2,-114)..(158.7,-113.8)
				..controls(159.1,-113.7)and(159.8,-113.3)..(160.2,-113)
				..controls(160.6,-112.7)and(161.3,-112.4)..(161.7,-112.2)
				..controls(162.7,-111.9)and(163.2,-111.7)..(164.1,-111.1)
				..controls(164.5,-110.8)and(165.3,-110.4)..(165.9,-110)
				..controls(166.5,-109.7)and(167.3,-109.2)..(167.6,-108.9)
				..controls(167.9,-108.6)and(168.6,-108.2)..(169,-108)
				..controls(169.4,-107.8)and(170.2,-107.3)..(170.7,-106.9)
				..controls(171.2,-106.5)and(172.1,-105.8)..(172.7,-105.4)
				..controls(173.2,-105)and(174.1,-104.4)..(174.5,-104)
				..controls(174.9,-103.7)and(175.7,-103.1)..(176.2,-102.6)
				..controls(176.7,-102.2)and(177.4,-101.5)..(177.9,-101.2)
				..controls(178.3,-100.8)and(178.8,-100.4)..(179,-100.2)
				..controls(179.2,-100)and(179.6,-99.6)..(180,-99.3)
				..controls(180.4,-99.1)and(180.9,-98.6)..(181.2,-98.4)
				..controls(181.5,-98.1)and(182.1,-97.7)..(182.5,-97.5)
				..controls(183.2,-97.2)and(183.5,-96.9)..(184.3,-95.9)
				..controls(184.6,-95.6)and(184.6,-95.6)..(184.8,-95.7)
				..controls(185,-95.8)and(185.1,-96)..(185.1,-96.2)
				..controls(185.1,-96.5)and(185,-96.7)..(184.6,-97.2)
				..controls(184.1,-97.8)and(184,-97.9)..(183.9,-98.8)
				..controls(183.7,-100.3)and(183.3,-101.3)..(182.9,-101.8)
				..controls(182.6,-102.1)and(182.3,-102.6)..(182.2,-102.9)
				..controls(182,-103.3)and(181.6,-104)..(181.2,-104.6)
				..controls(180.9,-105.2)and(180.5,-106.1)..(180.4,-106.6)
				..controls(179.9,-108)and(179.7,-108.5)..(179,-109.4)
				..controls(178.6,-109.8)and(178.2,-110.5)..(177.9,-111)
				..controls(177.7,-111.5)and(177.3,-112.3)..(177,-112.7)
				..controls(176.6,-113.2)and(176.4,-113.7)..(176.3,-114.2)
				..controls(176.2,-114.7)and(175.9,-115.2)..(175.4,-116)
				..controls(175,-116.6)and(174.4,-117.7)..(174,-118.5)
				..controls(173.7,-119.2)and(173.3,-120.1)..(173,-120.5)
				..controls(172.8,-120.9)and(172.4,-121.8)..(172.1,-122.5)
				..controls(171.8,-123.2)and(171.3,-124.1)..(171,-124.6)
				..controls(170.8,-125.1)and(170.3,-126)..(169.9,-126.6)
				..controls(169.6,-127.3)and(169.2,-128.1)..(169,-128.3)
				..controls(168.8,-128.6)and(168.5,-129.3)..(168.3,-129.8)
				..controls(168.1,-130.3)and(167.8,-131.1)..(167.5,-131.6)
				..controls(166.9,-132.5)and(164.2,-137.9)..(164.2,-138.1)
				..controls(164.2,-138.1)and(163.9,-138.8)..(163.4,-139.7)
				..controls(162.9,-140.5)and(162.4,-141.5)..(162.3,-141.9)
				..controls(161.9,-143)and(161.6,-143.4)..(160.9,-144.1)
				..controls(160.4,-144.6)and(160.3,-144.8)..(160.1,-145.4)
				..controls(160,-145.8)and(159.6,-146.8)..(159.2,-147.8)
				..controls(158.7,-148.8)and(158.1,-150)..(157.9,-150.6)
				..controls(157.6,-151.1)and(157.2,-151.9)..(156.9,-152.4)
				..controls(156.5,-152.8)and(156.1,-153.7)..(155.8,-154.4)
				..controls(155.6,-155)and(155.2,-155.9)..(154.9,-156.3)
				..controls(154.7,-156.7)and(154.4,-157.4)..(154.2,-157.9)
				..controls(154.1,-158.4)and(153.8,-159)..(153.5,-159.3)
				..controls(152.8,-160.1)and(153,-160.1)..(142.4,-160.1)
				..controls(137.2,-160.1)and(131.8,-160.1)..(130.4,-160)--cycle
				(141.5,-103.7)..controls(140.8,-103)and(139.8,-102.1)..(139.2,-101.6)
				..controls(138.5,-101.2)and(137.9,-100.6)..(137.8,-100.4)
				..controls(137,-99.5)and(134,-96.5)..(133.2,-96)
				..controls(132.7,-95.7)and(131.8,-94.9)..(131.1,-94.3)
				..controls(129.4,-92.7)and(127.9,-91.6)..(126.6,-91.1)
				..controls(125.2,-90.5)and(124.1,-90)..(122.9,-89.1)
				..controls(122.1,-88.6)and(121.3,-88.2)..(119.9,-87.7)
				..controls(118.8,-87.3)and(117.5,-86.8)..(116.9,-86.5)
				..controls(114.3,-85.5)and(113.7,-85.4)..(109.7,-85.5)
				..controls(105.3,-85.7)and(103.9,-85.9)..(101.6,-87)
				..controls(100.7,-87.4)and(99.4,-87.8)..(97.9,-88)
				..controls(97,-88.2)and(96.6,-88.3)..(95.8,-88.9)
				..controls(95.4,-89.2)and(94.7,-89.6)..(94.2,-89.8)
				..controls(92.9,-90.3)and(92.3,-90.7)..(90.9,-92.1)
				..controls(90.2,-92.7)and(89.4,-93.3)..(89,-93.5)..controls(87,-94.5)and(85.2,-96)..(83.6,-97.9)
				..controls(83.3,-98.2)and(82.5,-98.9)..(81.8,-99.4)
				..controls(80.7,-100.1)and(80.5,-100.3)..(80,-101.1)
				..controls(78.7,-103.1)and(77.2,-103.9)..(76.4,-103)
				..controls(75.7,-102.2)and(75.7,-103.2)..(75.7,-83.6)
				..controls(75.6,-73.6)and(75.6,-64.8)..(75.5,-64)
				..controls(75.5,-63.2)and(75.5,-61.4)..(75.6,-59.8)--(75.7,-57.1)--(76,-56.8)
				..controls(76.8,-56.2)and(78,-56.3)..(79.8,-57.2)
				..controls(80.4,-57.5)and(81.3,-57.8)..(81.8,-57.9)
				..controls(82.3,-58.1)and(83,-58.4)..(83.6,-58.8)..controls(84.3,-59.2)and(85,-59.5)..(86,-59.8)
				..controls(86.8,-60.1)and(88,-60.6)..(88.7,-60.9)
				..controls(89.3,-61.2)and(90.6,-61.7)..(91.4,-62)
				..controls(92.2,-62.3)and(93.7,-62.9)..(94.6,-63.3)
				..controls(95.6,-63.7)and(96.8,-64.1)..(97.4,-64.3)
				..controls(98,-64.5)and(99,-64.9)..(99.7,-65.2)
				..controls(100.4,-65.6)and(101.5,-66)..(102.3,-66.2)
				..controls(104.2,-66.8)and(105,-67.1)..(106.2,-67.8)
				..controls(107.9,-68.7)and(110.2,-69.4)..(111.3,-69.2)
				..controls(112.5,-69)and(115,-68.2)..(115.6,-67.8)
				..controls(117.3,-66.7)and(117.6,-66.6)..(119.5,-66.2)
				..controls(121.9,-65.8)and(122.1,-65.7)..(123,-65.1)
				..controls(124.1,-64.4)and(124.4,-64.3)..(126,-63.9)
				..controls(127,-63.7)and(127.7,-63.4)..(128.4,-63.1)
				..controls(129,-62.8)and(129.7,-62.5)..(130.6,-62.4)
				..controls(132.1,-62.1)and(134,-61.5)..(135.2,-61)
				..controls(135.7,-60.8)and(136.7,-60.5)..(137.3,-60.3)
				..controls(138.4,-60)and(139.3,-59.7)..(141.1,-59)
				..controls(141.6,-58.7)and(142,-58.6)..(142.5,-58.6)
				..controls(143.2,-58.6)and(143.3,-58.7)..(143.6,-59)
				..controls(144,-59.5)and(144.1,-59.9)..(144,-60.8)
				..controls(144,-61.2)and(144,-71.1)..(144,-82.8)--(143.9,-104.1)--(143.7,-104.5)
				..controls(143.5,-104.7)and(143.3,-104.8)..(143,-104.9)
				..controls(142.6,-104.9)and(142.5,-104.8)..(141.5,-103.7)--cycle
				(27.8,-91.4)..controls(27.1,-90.9)and(26.3,-90)..(26.1,-89.5)
				..controls(26,-89.2)and(25.8,-88.6)..(25.7,-87.9)
				..controls(25.4,-86.2)and(25.6,-84.8)..(26.5,-83.5)
				..controls(27.4,-82.3)and(28.3,-82)..(29.4,-82.3)
				..controls(31.1,-82.8)and(32.1,-84.9)..(32,-87.4)
				..controls(31.9,-89.3)and(31.2,-90.6)..(29.9,-91.4)
				..controls(29.1,-91.9)and(28.6,-91.9)..(27.8,-91.4)--cycle
				(190.1,-91.3)..controls(189.7,-91.2)and(189.3,-90.9)..(189.1,-90.7)
				..controls(188.3,-89.8)and(187.6,-87.7)..(187.6,-86.2)
				..controls(187.6,-85.5)and(187.7,-85.2)..(188.2,-84.4)
				..controls(189.1,-82.6)and(190.1,-81.8)..(191.3,-82.1)
				..controls(192,-82.3)and(193.5,-83.7)..(193.8,-84.4)
				..controls(194.1,-85.1)and(194.2,-87.1)..(194,-88.3)
				..controls(193.8,-89.8)and(193.3,-90.5)..(192.3,-91.3)
				..controls(191.7,-91.7)and(190.9,-91.7)..(190.1,-91.3)--cycle
				(173.7,-83.4)..controls(172.2,-82.5)and(172,-82.5)..(170.6,-82.4)
				..controls(169.7,-82.4)and(168.8,-82.4)..(168.3,-82.5)
				..controls(167.2,-82.7)and(166.9,-82.7)..(166.6,-82.2)
				..controls(166.3,-81.8)and(166.3,-81)..(166.8,-80)
				..controls(167.2,-78.9)and(167.4,-77.4)..(167.7,-74.1)
				..controls(167.8,-72.8)and(167.9,-71.8)..(167.9,-71.7)
				..controls(168,-71.7)and(168,-70.2)..(168,-68.4)..controls(168,-65.4)and(168,-64.9)..(167.7,-63)
				..controls(167.3,-60.5)and(167.1,-59.9)..(166.4,-58.4)
				..controls(166.1,-57.7)and(165.8,-56.9)..(165.7,-56.6)
				..controls(165.7,-56.2)and(165.4,-55.6)..(165.2,-55.2)
				..controls(164.6,-53.8)and(161.9,-50.7)..(161,-50.3)
				..controls(159.8,-49.7)and(157.9,-50)..(155.2,-51.2)
				..controls(154.7,-51.4)and(153.4,-51.8)..(152.2,-52.1)
				..controls(150.8,-52.4)and(149.4,-52.9)..(148.6,-53.2)
				..controls(147.8,-53.6)and(146.7,-54)..(146,-54.2)
				..controls(145.4,-54.4)and(144.5,-54.7)..(144,-54.9)
				..controls(143.6,-55.2)and(142.7,-55.5)..(141.8,-55.7)
				..controls(141,-55.9)and(139.6,-56.3)..(138.7,-56.6)
				..controls(137.8,-57)and(136.3,-57.4)..(135.3,-57.6)
				..controls(133.8,-57.9)and(133.2,-58.1)..(132.1,-58.7)
				..controls(131.1,-59.2)and(130.3,-59.4)..(129.4,-59.6)
				..controls(128.7,-59.8)and(127.8,-60)..(127.4,-60.1)
				..controls(127,-60.3)and(126,-60.5)..(125.1,-60.7)
				..controls(123.2,-61.1)and(121.7,-61.5)..(120.2,-62.1)
				..controls(119.7,-62.3)and(118.1,-62.8)..(116.7,-63.2)
				..controls(114.6,-63.9)and(114,-64.1)..(113,-64.7)
				..controls(111.9,-65.4)and(111.8,-65.4)..(111,-65.5)
				..controls(109.9,-65.5)and(109.3,-65.4)..(107.7,-64.7)
				..controls(107,-64.4)and(106,-64)..(105.3,-63.9)
				..controls(103.5,-63.5)and(102.2,-63)..(101.2,-62.5)
				..controls(100.1,-61.8)and(99.3,-61.5)..(97.8,-61.1)
				..controls(97.2,-61)and(96.2,-60.7)..(95.7,-60.4)
				..controls(95.1,-60.1)and(94.2,-59.8)..(93.5,-59.6)
				..controls(92.9,-59.4)and(91.7,-59)..(91,-58.7)..controls(90.2,-58.4)and(88.9,-58)..(88,-57.7)
				..controls(87.2,-57.5)and(86.1,-57.1)..(85.5,-56.8)
				..controls(85,-56.6)and(84.1,-56.2)..(83.5,-56)..controls(82.9,-55.7)and(81.8,-55.2)..(81,-54.9)
				..controls(80.1,-54.5)and(79.1,-54.2)..(78.7,-54.1)
				..controls(77.9,-53.8)and(77.3,-53.6)..(76.2,-52.9)
				..controls(75.7,-52.6)and(75.2,-52.4)..(74.1,-52.2)
				..controls(72.5,-51.9)and(72.4,-51.8)..(71.3,-51.1)
				..controls(69.9,-50.1)and(69.4,-49.9)..(67.9,-49.8)
				..controls(67,-49.7)and(66.6,-49.6)..(66.3,-49.4)
				..controls(66.1,-49.3)and(65.7,-49)..(65.4,-48.8)
				..controls(65.1,-48.7)and(64.2,-48.4)..(63.6,-48.2)
				..controls(62.9,-48)and(61.7,-47.6)..(60.9,-47.2)
				..controls(60.1,-46.8)and(58.9,-46.4)..(58.2,-46.2)
				..controls(57.6,-46)and(56.8,-45.7)..(56.5,-45.5)
				..controls(54.9,-44.6)and(54.6,-44.5)..(53.2,-44.2)
				..controls(51.7,-43.8)and(51,-43.6)..(49.9,-42.9)
				..controls(49.5,-42.6)and(48.9,-42.4)..(47.5,-42)
				..controls(45.7,-41.6)and(45.6,-41.5)..(44.3,-40.6)
				..controls(43.5,-40.1)and(42.8,-39.6)..(42.7,-39.4)
				..controls(42.4,-38.7)and(43,-38.1)..(44.5,-37.6)..controls(44.9,-37.5)and(46,-37)..(46.9,-36.6)
				..controls(47.9,-36.1)and(48.8,-35.8)..(49.6,-35.6)
				..controls(50.5,-35.4)and(51.3,-35.2)..(52.4,-34.6)
				..controls(53.3,-34.1)and(54.5,-33.7)..(54.9,-33.5)
				..controls(55.4,-33.4)and(56.4,-33)..(57.2,-32.6)
				..controls(58.6,-32)and(60.1,-31.5)..(61.6,-31.2)
				..controls(62.3,-31.1)and(62.7,-31)..(63.2,-30.7)
				..controls(64.1,-30)and(64.3,-29.9)..(65.3,-29.7)
				..controls(65.7,-29.6)and(66.5,-29.2)..(67.1,-28.9)
				..controls(67.6,-28.6)and(68.8,-28.2)..(69.6,-27.9)
				..controls(70.5,-27.7)and(71.6,-27.3)..(72.1,-27)
				..controls(72.6,-26.7)and(73.8,-26.2)..(74.8,-25.9)
				..controls(75.7,-25.5)and(76.7,-25.1)..(77.1,-25)
				..controls(77.4,-24.8)and(78.5,-24.4)..(79.6,-24.1)
				..controls(80.6,-23.8)and(82,-23.2)..(82.8,-22.9)
				..controls(83.5,-22.5)and(84.8,-22)..(85.8,-21.8)
				..controls(86.7,-21.5)and(87.8,-21.1)..(88.1,-20.9)
				..controls(88.4,-20.7)and(89.3,-20.3)..(90.1,-20.1)
				..controls(90.8,-19.8)and(91.9,-19.4)..(92.4,-19.1)
				..controls(93.5,-18.5)and(94.4,-18.2)..(95.4,-18.1)
				..controls(95.7,-18)and(96.4,-17.8)..(96.9,-17.5)
				..controls(97.5,-17.3)and(98.4,-16.9)..(99,-16.6)
				..controls(99.7,-16.3)and(100.9,-15.9)..(101.9,-15.6)
				..controls(102.9,-15.3)and(104.4,-14.7)..(105.3,-14.3)
				..controls(108.1,-13.2)and(108.5,-13.1)..(110.1,-13.1)
				..controls(111.4,-13.1)and(111.5,-13.1)..(112.4,-13.5)
				..controls(113.2,-14)and(115,-14.8)..(116.6,-15.4)
				..controls(116.9,-15.5)and(117.9,-15.8)..(118.8,-16.1)
				..controls(120.9,-16.6)and(121.5,-16.9)..(122.6,-17.9)
				..controls(123,-18.3)and(123.7,-18.8)..(123.9,-18.9)
				..controls(124.6,-19.4)and(126,-19.8)..(127.9,-20.2)
				..controls(129.6,-20.5)and(129.8,-20.6)..(131.1,-21.4)
				..controls(131.5,-21.7)and(132.4,-22.1)..(133.1,-22.3)
				..controls(136.1,-23.2)and(138.3,-24.1)..(140.2,-25)
				..controls(141,-25.5)and(141.5,-25.6)..(142.5,-25.8)
				..controls(145.1,-26.3)and(146,-26.6)..(149.4,-28.1)
				..controls(151.3,-28.9)and(153.5,-29.8)..(154.3,-30.1)
				..controls(155.1,-30.4)and(156.2,-30.8)..(156.6,-31)
				..controls(157.1,-31.2)and(158.2,-31.6)..(159.1,-31.8)
				..controls(160.6,-32.2)and(160.9,-32.3)..(161.7,-32.9)
				..controls(162.4,-33.3)and(162.9,-33.6)..(163.7,-33.8)
				..controls(164.3,-34)and(164.8,-34.1)..(164.9,-34.1)
				..controls(165,-34.1)and(165.6,-34.4)..(166.4,-34.8)
				..controls(167.3,-35.2)and(168.1,-35.4)..(169,-35.6)
				..controls(170,-35.8)and(170.5,-36)..(171.7,-36.6)
				..controls(172.4,-37)and(173.3,-37.4)..(173.6,-37.4)
				..controls(175.6,-37.7)and(176.2,-38)..(176.7,-39.1)
				..controls(176.9,-39.6)and(176.9,-39.7)..(176.7,-40.4)
				..controls(176.4,-41.6)and(175.8,-42.7)..(175.2,-43.4)
				..controls(174.5,-44.1)and(172.6,-45.5)..(171.7,-45.8)
				..controls(170.6,-46.2)and(169.1,-46)..(168.2,-45.4)
				..controls(167.9,-45.2)and(167.3,-44.8)..(166.9,-44.4)
				..controls(166.4,-44.1)and(165.8,-43.6)..(165.4,-43.2)
				..controls(165.1,-42.9)and(164.4,-42.3)..(163.9,-41.9)
				..controls(163.3,-41.5)and(162.7,-40.9)..(162.4,-40.5)
				..controls(161.4,-39.2)and(159.9,-38)..(158.8,-37.8)
				..controls(158.4,-37.8)and(158.3,-37.7)..(158.2,-37.2)
				..controls(158,-36.8)and(157.4,-36.1)..(157.1,-36.1)
				..controls(157,-36.1)and(156.7,-36.1)..(156.5,-36)
				..controls(156.1,-35.9)and(155.9,-35.7)..(155.7,-35.2)
				..controls(155.3,-34.6)and(155.1,-34.4)..(154,-34.1)
				..controls(153.6,-34)and(153.2,-33.7)..(152.9,-33.5)
				..controls(152.4,-32.9)and(151.6,-32.4)..(151.3,-32.4)
				..controls(150.9,-32.4)and(150.4,-32.6)..(150.4,-32.9)
				..controls(150.4,-33.2)and(150.8,-33.4)..(151.4,-33.6)
				..controls(152.1,-33.7)and(152.8,-34.1)..(153.5,-34.9)
				..controls(154.1,-35.5)and(154.3,-35.6)..(155,-35.8)
				..controls(155.7,-36)and(155.7,-36)..(156,-36.6)
				..controls(156.3,-37.3)and(156.7,-37.7)..(157.5,-37.8)
				..controls(158.1,-37.9)and(158.2,-38)..(158.4,-38.8)
				..controls(158.5,-39.4)and(158.6,-39.6)..(159.6,-40.4)
				..controls(160.2,-41)and(160.9,-41.6)..(161.3,-41.8)
				..controls(161.7,-42)and(161.9,-42.3)..(162.2,-42.8)
				..controls(162.4,-43.2)and(162.8,-43.8)..(163.1,-44.1)
				..controls(163.9,-44.8)and(165.1,-46.7)..(165.7,-48)
				..controls(166.2,-49.3)and(166.4,-49.7)..(167.2,-50.5)
				..controls(167.5,-50.8)and(167.9,-51.4)..(168.1,-51.8)
				..controls(168.3,-52.2)and(168.8,-53)..(169.1,-53.6)
				..controls(169.6,-54.3)and(169.8,-54.9)..(169.9,-55.6)
				..controls(170.4,-57.8)and(170.5,-58.1)..(171.2,-59.5)
				..controls(171.6,-60.3)and(172,-61.5)..(172.2,-62.1)
				..controls(172.3,-62.8)and(172.7,-63.7)..(172.9,-64.1)
				..controls(173.2,-64.6)and(173.4,-65.3)..(173.4,-65.6)
				..controls(173.8,-68.8)and(174.2,-70.6)..(174.9,-72.6)
				..controls(175.5,-74.5)and(175.9,-77.1)..(176.1,-80.2)
				..controls(176.3,-83.4)and(175.5,-84.4)..(173.7,-83.4)--cycle;
				\definecolor{ezs}{RGB}{247,179,26}\fill[ezs](-5.7,-156.6)--(-3.4,-157.5)--(61.9,12.7)--(59.7,13.6)--cycle;
				\definecolor{ezs}{RGB}{247,179,22}\fill[ezs](-3.8,-157.3)--(-1.5,-158.2)--(63.9,12)--(61.6,12.9)--cycle;
				\definecolor{ezs}{RGB}{248,179,19}\fill[ezs](-1.8,-158.1)--(0.4,-158.9)--(65.8,11.2)--(63.5,12.1)--cycle;
				\definecolor{ezs}{RGB}{248,178,16}\fill[ezs](0.1,-158.8)--(2.4,-159.7)--(67.7,10.5)--(65.4,11.4)--cycle;
				\definecolor{ezs}{RGB}{249,178,15}\fill[ezs](2,-159.5)--(4.3,-160.4)--(69.6,9.8)--(67.3,10.7)--cycle;
				\definecolor{ezs}{RGB}{251,177,17}\fill[ezs](3.9,-160.3)--(6.2,-161.1)--(71.5,9)--(69.2,9.9)--cycle;
				\definecolor{ezs}{RGB}{252,179,18}\fill[ezs](5.8,-161)--(8.1,-161.9)--(73.4,8.3)--(71.1,9.2)--cycle;
				\definecolor{ezs}{RGB}{252,182,17}\fill[ezs](7.7,-161.7)--(10,-162.6)--(75.3,7.6)--(73,8.5)--cycle;
				\definecolor{ezs}{RGB}{252,185,17}\fill[ezs](9.6,-162.5)--(11.9,-163.4)--(77.2,6.8)--(75,7.7)--cycle;
				\definecolor{ezs}{RGB}{251,189,18}\fill[ezs](11.5,-163.2)--(13.8,-164.1)--(79.2,6.1)--(76.9,7)--cycle;
				\definecolor{ezs}{RGB}{251,191,22}\fill[ezs](13.4,-163.9)--(15.7,-164.8)--(81.1,5.4)--(78.8,6.2)--cycle;
				\definecolor{ezs}{RGB}{252,193,25}\fill[ezs](15.4,-164.7)--(17.6,-165.6)--(83,4.6)--(80.7,5.5)--cycle;
				\definecolor{ezs}{RGB}{253,195,27}\fill[ezs](17.3,-165.4)--(19.6,-166.3)--(84.9,3.9)--(82.6,4.8)--cycle;
				\definecolor{ezs}{RGB}{253,196,31}\fill[ezs](19.2,-166.1)--(21.5,-167)--(86.8,3.2)--(84.5,4)--cycle;
				\definecolor{ezs}{RGB}{254,199,35}\fill[ezs](21.1,-166.9)--(23.4,-167.8)--(88.7,2.4)--(86.4,3.3)--cycle;
				\definecolor{ezs}{RGB}{253,201,38}\fill[ezs](23,-167.6)--(25.3,-168.5)--(90.6,1.7)--(88.3,2.6)--cycle;
				\definecolor{ezs}{RGB}{253,203,40}\fill[ezs](24.9,-168.3)--(27.2,-169.2)--(92.5,1)--(90.3,1.8)--cycle;
				\definecolor{ezs}{RGB}{253,205,42}\fill[ezs](26.8,-169.1)--(29.1,-170)--(94.4,0.2)--(92.2,1.1)--cycle;
				\definecolor{ezs}{RGB}{253,207,45}\fill[ezs](28.7,-169.8)--(31,-170.7)--(96.4,-0.5)--(94.1,0.4)--cycle;
				\definecolor{ezs}{RGB}{253,210,49}\fill[ezs](30.7,-170.5)--(32.9,-171.4)--(98.3,-1.2)--(96,-0.4)--cycle;
				\definecolor{ezs}{RGB}{253,212,54}\fill[ezs](32.6,-171.3)--(34.9,-172.2)--(100.2,-2)--(97.9,-1.1)--cycle;
				\definecolor{ezs}{RGB}{253,214,59}\fill[ezs](34.5,-172)--(36.8,-172.9)--(102.1,-2.7)--(99.8,-1.8)--cycle;
				\definecolor{ezs}{RGB}{253,216,63}\fill[ezs](36.4,-172.7)--(38.7,-173.6)--(104,-3.4)--(101.7,-2.6)--cycle;
				\definecolor{ezs}{RGB}{253,218,65}\fill[ezs](38.3,-173.5)--(40.6,-174.4)--(105.9,-4.2)--(103.6,-3.3)--cycle;
				\definecolor{ezs}{RGB}{254,220,66}\fill[ezs](40.2,-174.2)--(42.5,-175.1)--(107.8,-4.9)--(105.5,-4)--cycle;
				\definecolor{ezs}{RGB}{254,222,68}\fill[ezs](42.1,-175)--(44.4,-175.8)--(109.7,-5.6)--(107.5,-4.8)--cycle;
				\definecolor{ezs}{RGB}{254,224,71}\fill[ezs](44,-175.7)--(46.3,-176.6)--(111.7,-6.4)--(109.4,-5.5)--cycle;
				\definecolor{ezs}{RGB}{254,226,75}\fill[ezs](46,-176.4)--(48.2,-177.3)--(113.6,-7.1)--(111.3,-6.2)--cycle;
				\definecolor{ezs}{RGB}{254,228,79}\fill[ezs](47.9,-177.2)--(50.1,-178)--(115.5,-7.8)--(113.2,-7)--cycle;
				\definecolor{ezs}{RGB}{254,230,82}\fill[ezs](49.8,-177.9)--(52.1,-178.8)--(117.4,-8.6)--(115.1,-7.7)--cycle;
				\definecolor{ezs}{RGB}{254,232,85}\fill[ezs](51.7,-178.6)--(54,-179.5)--(119.3,-9.3)--(117,-8.4)--cycle;
				\definecolor{ezs}{RGB}{254,234,88}\fill[ezs](53.6,-179.4)--(55.9,-180.2)--(121.2,-10)--(118.9,-9.2)--cycle;
				\definecolor{ezs}{RGB}{254,237,91}\fill[ezs](55.5,-180.1)--(57.8,-181)--(123.1,-10.8)--(120.8,-9.9)--cycle;
				\definecolor{ezs}{RGB}{254,239,94}\fill[ezs](57.4,-180.8)--(59.7,-181.7)--(125,-11.5)--(122.8,-10.6)--cycle;
				\definecolor{ezs}{RGB}{254,241,97}\fill[ezs](59.3,-181.6)--(61.6,-182.4)--(127,-12.2)--(124.7,-11.4)--cycle;
				\definecolor{ezs}{RGB}{254,243,101}\fill[ezs](61.2,-182.3)--(63.5,-183.2)--(128.9,-13)--(126.6,-12.1)--cycle;
				\definecolor{ezs}{RGB}{254,245,103}\fill[ezs](63.2,-183)--(65.4,-183.9)--(130.8,-13.7)--(128.5,-12.8)--cycle;
				\definecolor{ezs}{RGB}{254,245,105}\fill[ezs](65.1,-183.8)--(67.4,-184.6)--(132.7,-14.4)--(130.4,-13.6)--cycle;
				\definecolor{ezs}{RGB}{254,245,105}\fill[ezs](67,-184.5)--(69.3,-185.4)--(134.6,-15.2)--(132.3,-14.3)--cycle;
				\definecolor{ezs}{RGB}{254,244,103}\fill[ezs](68.9,-185.2)--(71.2,-186.1)--(136.5,-15.9)--(134.2,-15)--cycle;
				\definecolor{ezs}{RGB}{253,243,101}\fill[ezs](70.8,-186)--(73.1,-186.8)--(138.4,-16.6)--(136.1,-15.8)--cycle;
				\definecolor{ezs}{RGB}{253,241,98}\fill[ezs](72.7,-186.7)--(75,-187.6)--(140.3,-17.4)--(138,-16.5)--cycle;
				\definecolor{ezs}{RGB}{253,239,95}\fill[ezs](74.6,-187.4)--(76.9,-188.3)--(142.2,-18.1)--(140,-17.2)--cycle;
				\definecolor{ezs}{RGB}{253,236,92}\fill[ezs](76.5,-188.2)--(78.8,-189)--(144.2,-18.8)--(141.9,-18)--cycle;
				\definecolor{ezs}{RGB}{253,234,88}\fill[ezs](78.5,-188.9)--(80.7,-189.8)--(146.1,-19.6)--(143.8,-18.7)--cycle;
				\definecolor{ezs}{RGB}{253,232,85}\fill[ezs](80.4,-189.6)--(82.7,-190.5)--(148,-20.3)--(145.7,-19.4)--cycle;
				\definecolor{ezs}{RGB}{253,230,82}\fill[ezs](82.3,-190.4)--(84.6,-191.2)--(149.9,-21)--(147.6,-20.2)--cycle;
				\definecolor{ezs}{RGB}{253,228,78}\fill[ezs](84.2,-191.1)--(86.5,-192)--(151.8,-21.8)--(149.5,-20.9)--cycle;
				\definecolor{ezs}{RGB}{253,226,74}\fill[ezs](86.1,-191.8)--(88.4,-192.7)--(153.7,-22.5)--(151.4,-21.6)--cycle;
				\definecolor{ezs}{RGB}{253,224,70}\fill[ezs](88,-192.6)--(90.3,-193.4)--(155.6,-23.3)--(153.3,-22.4)--cycle;
				\definecolor{ezs}{RGB}{253,222,67}\fill[ezs](89.9,-193.3)--(92.2,-194.2)--(157.5,-24)--(155.3,-23.1)--cycle;
				\definecolor{ezs}{RGB}{253,220,65}\fill[ezs](91.8,-194)--(94.1,-194.9)--(159.5,-24.7)--(157.2,-23.8)--cycle;
				\definecolor{ezs}{RGB}{253,217,63}\fill[ezs](93.7,-194.8)--(96,-195.6)--(161.4,-25.5)--(159.1,-24.6)--cycle;
				\definecolor{ezs}{RGB}{253,215,60}\fill[ezs](95.7,-195.5)--(97.9,-196.4)--(163.3,-26.2)--(161,-25.3)--cycle;
				\definecolor{ezs}{RGB}{253,213,56}\fill[ezs](97.6,-196.2)--(99.9,-197.1)--(165.2,-26.9)--(162.9,-26)--cycle;
				\definecolor{ezs}{RGB}{252,211,53}\fill[ezs](99.5,-197)--(101.8,-197.8)--(167.1,-27.7)--(164.8,-26.8)--cycle;
				\definecolor{ezs}{RGB}{252,209,49}\fill[ezs](101.4,-197.7)--(103.7,-198.6)--(169,-28.4)--(166.7,-27.5)--cycle;
				\definecolor{ezs}{RGB}{252,207,45}\fill[ezs](103.3,-198.4)--(105.6,-199.3)--(170.9,-29.1)--(168.6,-28.2)--cycle;
				\definecolor{ezs}{RGB}{253,205,42}\fill[ezs](105.2,-199.2)--(107.5,-200)--(172.8,-29.9)--(170.6,-29)--cycle;
				\definecolor{ezs}{RGB}{253,202,40}\fill[ezs](107.1,-199.9)--(109.4,-200.8)--(174.7,-30.6)--(172.5,-29.7)--cycle;
				\definecolor{ezs}{RGB}{253,200,38}\fill[ezs](109,-200.6)--(111.3,-201.5)--(176.7,-31.3)--(174.4,-30.4)--cycle;
				\definecolor{ezs}{RGB}{253,198,35}\fill[ezs](111,-201.4)--(113.2,-202.2)--(178.6,-32.1)--(176.3,-31.2)--cycle;
				\definecolor{ezs}{RGB}{253,196,31}\fill[ezs](112.9,-202.1)--(115.2,-203)--(180.5,-32.8)--(178.2,-31.9)--cycle;
				\definecolor{ezs}{RGB}{253,194,28}\fill[ezs](114.8,-202.8)--(117.1,-203.7)--(182.4,-33.5)--(180.1,-32.6)--cycle;
				\definecolor{ezs}{RGB}{253,192,26}\fill[ezs](116.7,-203.6)--(119,-204.5)--(184.3,-34.3)--(182,-33.4)--cycle;
				\definecolor{ezs}{RGB}{253,189,24}\fill[ezs](118.6,-204.3)--(120.9,-205.2)--(186.2,-35)--(183.9,-34.1)--cycle;
				\definecolor{ezs}{RGB}{253,186,20}\fill[ezs](120.5,-205)--(122.8,-205.9)--(188.1,-35.7)--(185.8,-34.8)--cycle;
				\definecolor{ezs}{RGB}{253,184,17}\fill[ezs](122.4,-205.8)--(124.7,-206.7)--(190,-36.5)--(187.8,-35.6)--cycle;
				\definecolor{ezs}{RGB}{252,183,15}\fill[ezs](124.3,-206.5)--(126.6,-207.4)--(192,-37.2)--(189.7,-36.3)--cycle;
				\definecolor{ezs}{RGB}{252,181,14}\fill[ezs](126.3,-207.2)--(128.5,-208.1)--(193.9,-37.9)--(191.6,-37.1)--cycle;
				\definecolor{ezs}{RGB}{252,179,13}\fill[ezs](128.2,-208)--(130.4,-208.9)--(195.8,-38.7)--(193.5,-37.8)--cycle;
				\definecolor{ezs}{RGB}{252,178,12}\fill[ezs](130.1,-208.7)--(132.4,-209.6)--(197.7,-39.4)--(195.4,-38.5)--cycle;
				\definecolor{ezs}{RGB}{251,178,10}\fill[ezs](132,-209.4)--(134.3,-210.3)--(199.6,-40.1)--(197.3,-39.3)--cycle;
				\definecolor{ezs}{RGB}{251,178,10}\fill[ezs](133.9,-210.2)--(136.2,-211.1)--(201.5,-40.9)--(199.2,-40)--cycle;
				\definecolor{ezs}{RGB}{250,177,11}\fill[ezs](135.8,-210.9)--(138.1,-211.8)--(203.4,-41.6)--(201.1,-40.7)--cycle;
				\definecolor{ezs}{RGB}{250,178,11}\fill[ezs](137.7,-211.6)--(140,-212.5)--(205.3,-42.3)--(203.1,-41.5)--cycle;
				\definecolor{ezs}{RGB}{248,180,10}\fill[ezs](139.6,-212.4)--(141.9,-213.3)--(207.3,-43.1)--(205,-42.2)--cycle;
				\definecolor{ezs}{RGB}{247,181,9}\fill[ezs](141.5,-213.1)--(143.8,-214)--(209.2,-43.8)--(206.9,-42.9)--cycle;
				\definecolor{ezs}{RGB}{246,182,9}\fill[ezs](143.5,-213.8)--(145.7,-214.7)--(211.1,-44.5)--(208.8,-43.7)--cycle;
				\definecolor{ezs}{RGB}{246,182,10}\fill[ezs](145.4,-214.6)--(147.7,-215.5)--(213,-45.3)--(210.7,-44.4)--cycle;
			\end{scope}
			% chữ EDUZY (Times đậm, làm dày nét bằng 5 lớp lệch nhẹ; căn đúng khung chữ logo gốc)
			\foreach \dx/\dy in {-0.7/0,0.7/0,0/-0.3,0/0.3,0/0}
			\node[anchor=center,inner sep=0pt,outer sep=0pt,text=ezText,align=center] at (109.6+\dx,-193.6+\dy)
			{\resizebox{1.01cm}{0.22cm}{\fontencoding{T1}\fontfamily{ptm}\fontseries{b}\selectfont EDUZY}};
		\end{tikzpicture}%
	}%
}
%====================== THIẾT LẬP THÂN ĐỀ ======================
%--- Môi trường tự luận {bt}: ex_test 3.3.4 chưa khai báo sẵn nên tạo ở đây
%    (ex_test đã tự móc \fixshowans và \listenumerate cho bt)
\makeatletter
\@ifundefined{bt}{\newtheorem{bt}{\color{gold1}\namebt}}{}
\makeatother
\renewcommand{\sectionmark}[1]{\markright{\thesection\ #1}}
%--- Khung tiêu đề (logo + tối đa 4 dòng; dòng 3, 4 để trống {} thì bỏ qua)
%  Đề thi:    \khungtieude{THI THỬ THPT QUỐC GIA NĂM 2026}{Môn Toán -- Đề số 21}
%                         {(Đề thi gồm: ...)}{Thời gian làm bài: 90 phút (không kể thời gian phát đề).}
%  Chuyên đề: \khungtieude{CHUYÊN ĐỀ 3 -- DÃY SỐ VÀ CÔNG THỨC TRUY HỒI}
%                         {Bồi dưỡng học sinh giỏi Toán THPT}{...}{}
\newcommand{\khungtieude}[4]{%
	\begin{tcolorbox}[colframe=gold1,colback=gold5!45,boxrule=0.5mm]
		\begin{minipage}[c]{3cm}
			\centerline{\eduzylogo[2.4cm]}
		\end{minipage}\hspace{0.5cm}
		\begin{minipage}[c]{12.5cm}
			\vskip 0.15 cm
			\centerline{\ind{#1}}
			\vskip 0.15 cm
			\centerline{\textbf{#2}}
			\if\relax\detokenize{#3}\relax\else
				\vskip 0.15 cm
				\centerline{\textit{#3}}
			\fi
			\if\relax\detokenize{#4}\relax\else
				\vskip 0.15 cm
				\centerline{\textit{#4}}
			\fi
			\vskip 0.15 cm
		\end{minipage}
	\end{tcolorbox}
}
%====================== BỔ SUNG CHO BÀI HỌC, CHUYÊN ĐỀ ======================
\usepackage[shortlabels]{enumitem}
\usepackage{tabularx}
%--- Danh sách gọn: \begin{enumerate}[gon] ... \end{enumerate}
\SetEnumitemKey{gon}{leftmargin=*,itemsep=1pt,parsep=0pt,topsep=3pt}
%--- Khung lý thuyết: \begin{khung}{} ... \end{khung}; tiêu đề chỉ dùng khi thật ngắn ("Mẹo", "Lỗi thường gặp")
\newtcolorbox{khung}[1]{breakable,enhanced,colframe=gold1,colback=gold5!45,
	boxrule=0.4mm,fonttitle=\bfseries\sffamily,coltitle=white,colbacktitle=gold1,
	attach boxed title to top left={xshift=6mm,yshift=-3mm},
	boxed title style={arc=1mm,boxrule=0.3mm},title={#1}}
%--- Tiêu đề mục nhỏ trong một dạng: \muc{Dấu hiệu}, \muc{Phương pháp}, \muc{Bài tự luyện}
\makeatletter
\newcommand{\muc}[1]{\par\medskip\noindent\indam{#1}\par\nobreak\smallskip\@afterindentfalse\@afterheading}
\makeatother
%--- Ô công thức cần nhớ (căn giữa, viền vàng, vừa khít nội dung)
%     \begin{center}\congthuc{$Q(n)=(n-1)!$}\end{center}
\newtcbox{\congthuc}{on line,colframe=gold1,colback=gold5!45,boxrule=0.5pt,arc=2pt,
	left=8pt,right=8pt,top=4pt,bottom=4pt}
%--- Công tắc lời giải bài tự luyện (ví dụ mẫu luôn dùng \loigiai)
%     \loigiaiTL[{đáp số}]{lời giải đầy đủ}
%     \TLfull -> in lời giải đầy đủ;  \TLdapso -> chỉ in một dòng "Đáp số";  \TLan -> ẩn hoàn toàn
%     (\CheDo là dethi thì mọi lời giải và đáp số đều ẩn)
\makeatletter
\newif\ifTLhien   \TLhientrue
\newif\ifTLdapso  \TLdapsofalse
\newcommand{\TLfull}{\TLhientrue\TLdapsofalse}
\newcommand{\TLdapso}{\TLhientrue\TLdapsotrue}
\newcommand{\TLan}{\TLhienfalse}
\newcommand{\loigiaiTL}[2][]{%
	\ifTLhien
		\ifTLdapso
			\def\@tlarg{#1}%
			\ifx\@tlarg\@empty\else
				\ifdefstring{\kindOPTN}{dethi}{}{\par\noindent{\bfseries\color{gold1}Đáp số:} #1\par}%
			\fi
		\else
			\loigiai{#2}%
		\fi
	\fi}
\makeatother
\TLdapso% mặc định: chỉ in đáp số bài tự luyện
%====================== BẢNG ĐÚNG/SAI (\choiceTFt) ======================
%  Hàng tiêu đề nền nâu đồng chữ trắng (Phát biểu | Đ | S), các hàng tô xen kẽ kem / trắng,
%  đường kẻ vàng kim, nhãn a) b) c) d) màu nâu đồng. Bản có đáp án: cột được chọn có vòng tròn
%  nâu đồng với dấu tích trắng (câu đúng tích ở cột Đúng, câu sai tích ở cột Sai).
%  Tiêu đề ghi đủ "Đúng | Sai": \OPTN{viettat=0} sau \input (cột Đúng/Sai rộng ra, phát biểu hẹp lại).
%  Dùng bảng gốc của ex_test: \def\BangDSgoc{1} trước %============================================================
%  PREAMBLE EDUZY v2 - dùng chung cho đề thi, bài học, chuyên đề (gói ex_test)
%  Khác bản preamble-eduzy.tex ở đầu trang, chân trang: dải nền nâu đồng tràn ngang,
%  logo + Eduzy.com.vn, ô chữ vàng kem cạnh xiên, chân trang "Trang (số)";
%  logo Eduzy chìm giữa mọi trang (tắt: \logochimtat, bật lại: \logochimbat).
%  Cách dùng ở file chính (xem mau-chuan.tex, mau-chuyen-de.tex):
%      \def\CheDo{loigiai}          % loigiai | dethi | solcolor | color | book
%      \def\tenchuong{<dòng chữ trên đầu trang>}
%      \def\thuonghieu{<chữ ở chân trang lẻ>}   % không khai báo thì dùng mặc định
%      \def\tacgia{<chữ ở chân trang chẵn>}      % không khai báo thì dùng mặc định
%      \def\khauhieu{<chữ cạnh logo ở đầu trang>}  % tùy chọn, mặc định Eduzy.com.vn
%      \input{preamble-eduzy-v2}
%      \TLdapso                      % tài liệu: \TLfull | \TLdapso | \TLan
%      \begin{document} ... \end{document}
%  Logo Eduzy vẽ bằng TikZ ngay trong file (\eduzylogo), không cần file ảnh.
%  Biên dịch: pdfLaTeX (máy có gói vietnam) hoặc XeLaTeX.
%============================================================
\documentclass[12pt,a4paper,twoside]{article}%<-- twoside: bat dau trang/chan trang chan-le
\usepackage{iftex}
\usepackage{amsmath,amssymb,tikz,mathrsfs,tkz-tab}
\usepackage{graphicx}
\usepackage{pgfplots}% đồ thị hàm số
\pgfplotsset{compat=1.17}
%\usepackage{mathptmx}
\usepackage{mathpazo}
\ifXeTeX% XeLaTeX: font chữ có đủ dấu tiếng Việt, giữ nguyên font toán mathpazo
	\usepackage[no-math]{fontspec}
	\setmainfont{TeX Gyre Pagella}
	\newfontfamily\EZqag{TeX Gyre Adventor}[NFSSFamily=qag,Ligatures=TeX]
\fi
\usepackage[many]{tcolorbox}
\usetikzlibrary{math,through,calc,intersections,angles,quotes,shapes,shapes.geometric,arrows,patterns,matrix,chains,arrows.meta,decorations.shapes,decorations.fractals,decorations.markings,shadows,decorations.pathreplacing}
\usepackage[top=1.6cm, bottom=1.8cm, left=1.8cm, right=1.8cm]{geometry}% v2: top 1.6cm (bản cũ 1.5cm) để chữ không sát dải đầu trang
\providecommand{\CheDo}{loigiai}% loigiai | dethi | solcolor | color | book
\expandafter\usepackage\expandafter[\CheDo]{ex_test}
% \tenchuong: khai báo ở file chính (trước \input) hoặc \renewcommand sau \input
%============================================================
%  [GỘP TỪ eduzy-de.tex] EDUZY - MÀU NÂU ĐỒNG / VÀNG KIM
%  + ĐẦU TRANG, CHÂN TRANG KIỂU DẢI NGANG (v2), ĐỐI XỨNG TRANG CHẴN / TRANG LẺ.
%  (không cần file eduzy-de.tex riêng nữa)
%  Khối này phải đặt SAU \usepackage[...]{ex_test} và SAU \def\tenchuong{...}
%============================================================
\usepackage{lastpage}
\usepackage{needspace}
%--- DÀN TRANG ĐỀ THI: tách câu (THÊM 05/11/2026; chỉ thêm lệnh mới, KHÔNG đổi hành vi mặc định) ---
%  Hai công tắc, mặc định tắt (= preamble chạy y như cũ). Bật bằng cách đặt TRƯỚC \input{preamble-eduzy-v2}
%  trong file đề (tools/dan_trang/dan_trang.py tự ghi khi đặt \NeedspaceCau):
%    \def\TachCauDS{1}   câu Đúng/Sai được tách: câu dẫn (kèm hình) ở cuối trang, bảng Phát biểu sang trang sau
%                        (câu dẫn và hình không tách nhau, bảng không tách giữa hai trang)
%    \def\TachCauLG{1}   bản lời giải (\CheDo khác dethi): phần đề của câu cao hơn 1/3 trang được phép tách
%  Lệnh dùng trong file đề (giá trị do tools/dan_trang/dan_trang.py đo ở bản dethi):
%    \NeedspaceCau[<cao câu dẫn + hình>]{<cao cả khối câu>}   ngay trên \begin{ex} (câu đầu phần: trên tiêu đề phần)
%        công tắc tắt: đúng bằng \Needspace{<cao cả khối câu>} như cũ;
%        \TachCauDS=1 và có [..] (câu Đúng/Sai được chọn tách): chỉ giữ câu dẫn + hình;
%        \TachCauLG=1 ở bản lời giải: đòi tối đa 1/3 \textheight.
%    \NeedspaceBang{<cao bảng Phát biểu>}%   ngay trước \choiceTFt: \TachCauDS=1 thì giữ trọn bảng trên một trang;
%        công tắc tắt: không làm gì.
\providecommand{\TachCauDS}{0}
\providecommand{\TachCauLG}{0}
\newlength{\EZnsCau}
\newcommand{\NeedspaceCau}[2][]{%
	\setlength{\EZnsCau}{#2}%
	\ifnum\TachCauDS=1 \if\relax\detokenize{#1}\relax\else\setlength{\EZnsCau}{#1}\fi\fi
	\ifnum\TachCauLG=1 \ifdefstring{\kindOPTN}{dethi}{}{\ifdim\EZnsCau>0.3333\textheight\setlength{\EZnsCau}{0.3333\textheight}\fi}\fi
	\Needspace{\EZnsCau}}
\newcommand{\NeedspaceBang}[1]{\ifnum\TachCauDS=1 \Needspace{#1}\fi}

%--- Thang màu: đậm -> nhạt
\definecolor{gold1}{RGB}{126,96,0}    % nâu đồng đậm  #7E6000
\definecolor{gold2}{RGB}{192,144,3}   % vàng kim      #C09003
\definecolor{gold3}{RGB}{255,215,99}  % vàng nhạt     #FFD763
\definecolor{gold4}{RGB}{255,231,153} % vàng rất nhạt #FFE799
\definecolor{gold5}{RGB}{255,243,205} % vàng kem      #FFF3CD
%--- Ghi đè tông xanh/đỏ cũ
\definecolor{myblue}{RGB}{126,96,0}
\colorlet{blue}{gold1}
\colorlet{cyan}{gold2}
\colorlet{teal}{gold2}

%--- Tên chương mặc định (file đề nên tự khai báo lại)
\providecommand{\tenchuong}{}
%--- Chân trang
\providecommand{\thuonghieu}{TOÁN 12 -- Eduzy.com.vn}
\providecommand{\tacgia}{Tô Gia Bảo -- 0353703783}

%--- Chữ cạnh logo ở đầu trang
\providecommand{\khauhieu}{Eduzy.com.vn}

%====================== THÔNG SỐ ĐẦU TRANG, CHÂN TRANG (v2) ======================
%  Dải nền đậm tràn ngang trang; ô chữ vàng kem có cạnh xiên, kèm 2 vạch xiên chuyển màu.
%  Đầu trang lẻ : logo + \khauhieu (trái)   | ô chữ \tenchuong (phải)
%  Chân trang lẻ: ô chữ \thuonghieu (trái)  | Trang (số) (phải)
%  Trang chẵn soi gương trang lẻ; chân trang chẵn ghi \tacgia.
%--- Màu: chỉ dùng thang vàng kim Eduzy
\colorlet{ezNen}{gold1}% dải nền đậm
\colorlet{ezVach1}{gold3}% vạch xiên sát dải nền (đậm hơn)
\colorlet{ezVach2}{gold4}% vạch xiên sát ô chữ (nhạt hơn)
\colorlet{ezO}{gold5}% ô chữ vàng kem
%--- Kích thước
\newlength{\EZdai}\setlength{\EZdai}{1.1cm}% chiều cao dải nền
\newlength{\EZo}\setlength{\EZo}{0.68cm}% chiều cao ô chữ
\newlength{\EZxien}\setlength{\EZxien}{0.55cm}% độ lệch ngang của cạnh xiên
\newlength{\EZvach}\setlength{\EZvach}{0.75cm}% bề rộng mỗi vạch xiên
\newlength{\EZkhe}\setlength{\EZkhe}{0.12cm}% khe trắng giữa vạch và ô chữ
\newlength{\EZdem}\setlength{\EZdem}{0.4cm}% khoảng đệm từ chữ tới cạnh xiên
\newlength{\EZle}\setlength{\EZle}{1.8cm}% lề trái/phải của chữ (bằng lề thân bài)
\newlength{\EZrongmin}\setlength{\EZrongmin}{4.5cm}% ô chữ rộng tối thiểu
\newlength{\EZrong}% bề rộng chữ trong ô (tính tự động)
%--- Font
\newcommand{\EZfontO}{\fontsize{10.5pt}{12pt}\selectfont\bfseries}% chữ trong ô
\newcommand{\EZfontKH}{\fontfamily{qag}\fontsize{11pt}{12pt}\selectfont\bfseries}% \khauhieu cạnh logo
\newcommand{\EZfontTrang}{\fontsize{11.5pt}{12pt}\selectfont\bfseries}% chữ "Trang" và số trang
%--- Đo bề rộng chữ trong ô
\newcommand{\EZdorong}[1]{%
	\settowidth{\EZrong}{\EZfontO #1}%
	\ifdim\EZrong<\EZrongmin\setlength{\EZrong}{\EZrongmin}\fi}
%--- Logo + chữ trắng \khauhieu (dựng sẵn một lần khi vào \begin{document})
\newsavebox{\EZlogobox}
\AtBeginDocument{\sbox{\EZlogobox}{%
		\raisebox{-0.5\height}{\eduzylogo[1cm]}\hspace{0.15cm}%
		\raisebox{-0.5\height}{\color{white}\EZfontKH\khauhieu}}}
%--- Hai vạch xiên + khe trắng. #1: điểm gốc ở mép dưới, #2: hướng (1 sang phải, -1 sang trái),
%    #3: hướng nghiêng (1 kiểu "/", -1 kiểu "\")
\newcommand{\EZvachxien}[3]{%
	\fill[white] (#1) --++({(#3)*\EZxien},\EZo) --++({(#2)*\EZkhe},0) --++({-(#3)*\EZxien},-\EZo) --cycle;
	\fill[ezVach2] ([xshift=(#2)*\EZkhe]#1) --++({(#3)*\EZxien},\EZo) --++({(#2)*\EZvach},0) --++({-(#3)*\EZxien},-\EZo) --cycle;
	\fill[ezVach1] ([xshift=(#2)*(\EZkhe+\EZvach)]#1) --++({(#3)*\EZxien},\EZo) --++({(#2)*\EZvach},0) --++({-(#3)*\EZxien},-\EZo) --cycle;
	\draw[white,line width=0.6pt] ([xshift=(#2)*(\EZkhe+\EZvach)]#1) --++({(#3)*\EZxien},\EZo);
}
%--- Chữ trong ô (căn giữa theo chiều cao ô, không lệch vì dấu tiếng Việt)
\tikzset{ezchu/.style={inner sep=0pt,text=ezNen,text height=0.72em,text depth=0.22em}}

%====================== LOGO CHÌM GIỮA TRANG ======================
%  Vẽ cùng đầu trang nên nằm dưới chữ của thân bài, có trên mọi trang.
%  Tắt: \logochimtat (từ trang đó trở đi), bật lại: \logochimbat.
%  Đổi cỡ, độ đậm: khai báo trước \input, ví dụ \def\EZchimcao{9cm} \def\EZchimdo{0.1}
\providecommand{\EZchimcao}{10cm}% chiều cao logo chìm
\providecommand{\EZchimdo}{0.12}% độ đậm (0 = trong suốt, 1 = rõ hẳn)
\newif\ifEZchim \EZchimtrue
\newcommand{\logochimtat}{\EZchimfalse}
\newcommand{\logochimbat}{\EZchimtrue}
%  Dựng sẵn logo mờ một lần khi vào \begin{document}; transparency group làm mờ cả logo
%  một lượt nên không lộ vệt ở chỗ các dải màu chồng nhau.
\newsavebox{\EZchimtho}
\newsavebox{\EZchimbox}
\AtBeginDocument{%
	\sbox{\EZchimtho}{\eduzylogo[\EZchimcao]}%
	\sbox{\EZchimbox}{\begin{tikzpicture}
			\begin{scope}[transparency group,opacity=\EZchimdo]
				\node[inner sep=0pt] {\usebox{\EZchimtho}};
			\end{scope}
	\end{tikzpicture}}}
\newcommand{\EZvechim}{%
	\ifEZchim
		\node[inner sep=0pt] at (current page.center) {\usebox{\EZchimbox}};
	\fi}

%====================== ĐẦU TRANG - TRANG LẺ ======================
%  (logo bên trái, ô chữ \tenchuong bên phải)
\newcommand{\myfancyheadOdd}{%
	\EZdorong{\tenchuong}%
	\begin{tikzpicture}[remember picture,overlay]
		\coordinate (T) at (current page.north west);
		\coordinate (P) at (current page.north east);
		%--- logo chìm giữa trang
		\EZvechim
		%--- dải nền
		\fill[ezNen] (T) rectangle ([yshift=-\EZdai]P);
		%--- ô chữ: X là góc dưới bên trái, cạnh trái xiên kiểu "/"
		\coordinate (X) at ([xshift=-\EZle-\EZrong-\EZdem-\EZxien,yshift=-\EZdai]P);
		\fill[ezO] (X) --++(\EZxien,\EZo) -- ([yshift=-\EZdai+\EZo]P) -- ([yshift=-\EZdai]P) --cycle;
		\EZvachxien{X}{-1}{1}
		\node[ezchu,anchor=east] at ([xshift=-\EZle,yshift=-\EZdai+0.5*\EZo]P) {\EZfontO\tenchuong};
		%--- logo + tên
		\node[inner sep=0pt,anchor=west] at ([xshift=\EZle,yshift=-0.5*\EZdai]T) {\usebox{\EZlogobox}};
	\end{tikzpicture}%
}

%====================== ĐẦU TRANG - TRANG CHẴN =====================
%  (soi gương: ô chữ \tenchuong bên trái, logo bên phải)
\newcommand{\myfancyheadEven}{%
	\EZdorong{\tenchuong}%
	\begin{tikzpicture}[remember picture,overlay]
		\coordinate (T) at (current page.north west);
		\coordinate (P) at (current page.north east);
		%--- logo chìm giữa trang
		\EZvechim
		%--- dải nền
		\fill[ezNen] (T) rectangle ([yshift=-\EZdai]P);
		%--- ô chữ: X là góc dưới bên phải, cạnh phải xiên kiểu "\"
		\coordinate (X) at ([xshift=\EZle+\EZrong+\EZdem+\EZxien,yshift=-\EZdai]T);
		\fill[ezO] (X) --++(-\EZxien,\EZo) -- ([yshift=-\EZdai+\EZo]T) -- ([yshift=-\EZdai]T) --cycle;
		\EZvachxien{X}{1}{-1}
		\node[ezchu,anchor=west] at ([xshift=\EZle,yshift=-\EZdai+0.5*\EZo]T) {\EZfontO\tenchuong};
		%--- logo + tên
		\node[inner sep=0pt,anchor=east] at ([xshift=-\EZle,yshift=-0.5*\EZdai]P) {\usebox{\EZlogobox}};
	\end{tikzpicture}%
}

%====================== CHÂN TRANG - TRANG LẺ ======================
%  (ô chữ \thuonghieu bên trái, "Trang (số)" bên phải)
\newcommand{\myfancyfootOdd}{%
	\EZdorong{\thuonghieu}%
	\begin{tikzpicture}[remember picture,overlay]
		\coordinate (T) at (current page.south west);
		\coordinate (P) at (current page.south east);
		%--- dải nền
		\fill[ezNen] (T) rectangle ([yshift=\EZdai]P);
		%--- ô chữ sát mép trên dải: X là góc dưới bên phải, cạnh phải xiên kiểu "/"
		\coordinate (X) at ([xshift=\EZle+\EZrong+\EZdem,yshift=\EZdai-\EZo]T);
		\fill[ezO] ([yshift=\EZdai-\EZo]T) -- (X) --++(\EZxien,\EZo) -- ([yshift=\EZdai]T) --cycle;
		\EZvachxien{X}{1}{1}
		\node[ezchu,anchor=west] at ([xshift=\EZle,yshift=\EZdai-0.5*\EZo]T) {\EZfontO\thuonghieu};
		%--- số trang
		\node[circle,draw=white,line width=0.9pt,text=white,inner sep=0pt,minimum size=0.62cm,anchor=east]
		(so) at ([xshift=-\EZle,yshift=0.5*\EZdai]P) {\EZfontTrang\thepage};
		\node[inner sep=0pt,text=white,anchor=base east] (tr) at ([xshift=-0.15cm]so.base west) {\EZfontTrang Trang};
	\end{tikzpicture}%
}

%====================== CHÂN TRANG - TRANG CHẴN ====================
%  (soi gương: "Trang (số)" bên trái, ô chữ \tacgia bên phải)
\newcommand{\myfancyfootEven}{%
	\EZdorong{\tacgia}%
	\begin{tikzpicture}[remember picture,overlay]
		\coordinate (T) at (current page.south west);
		\coordinate (P) at (current page.south east);
		%--- dải nền
		\fill[ezNen] (T) rectangle ([yshift=\EZdai]P);
		%--- ô chữ sát mép trên dải: X là góc dưới bên trái, cạnh trái xiên kiểu "\"
		\coordinate (X) at ([xshift=-\EZle-\EZrong-\EZdem,yshift=\EZdai-\EZo]P);
		\fill[ezO] ([yshift=\EZdai-\EZo]P) -- (X) --++(-\EZxien,\EZo) -- ([yshift=\EZdai]P) --cycle;
		\EZvachxien{X}{-1}{-1}
		\node[ezchu,anchor=east] at ([xshift=-\EZle,yshift=\EZdai-0.5*\EZo]P) {\EZfontO\tacgia};
		%--- số trang
		%    (vòng tròn ẩn so0 nằm giữa dải để lấy đường chân chữ, giống trang lẻ)
		\node[circle,inner sep=0pt,minimum size=0.62cm] (so0) at ([yshift=0.5*\EZdai]T) {\EZfontTrang\phantom{\thepage}};
		\node[inner sep=0pt,text=white,anchor=base west] (tr) at ($(so0.base -| T)+(\EZle,0)$) {\EZfontTrang Trang};
		\node[circle,draw=white,line width=0.9pt,text=white,inner sep=0pt,minimum size=0.62cm,anchor=base west]
		(so) at ([xshift=0.15cm]tr.base east) {\EZfontTrang\thepage};
	\end{tikzpicture}%
}

%--- Giữ tương thích với các file cũ đang gọi \myfancyhead, \myfancyfoot
\newcommand{\myfancyhead}{\myfancyheadOdd}
\newcommand{\myfancyfoot}{\myfancyfootOdd}

%======================= GẮN VÀO TRANG =============================
\pagestyle{fancy}
\fancyhf{}
\fancyhead[LO]{\myfancyheadOdd}
\fancyhead[RE]{\myfancyheadEven}
\fancyfoot[LO]{\myfancyfootOdd}
\fancyfoot[RE]{\myfancyfootEven}
\renewcommand{\headrulewidth}{0pt}
\renewcommand{\footrulewidth}{0pt}
\fancypagestyle{plain}{%
	\fancyhf{}
	\fancyhead[LO]{\myfancyheadOdd}
	\fancyhead[RE]{\myfancyheadEven}
	\fancyfoot[LO]{\myfancyfootOdd}
	\fancyfoot[RE]{\myfancyfootEven}
	\renewcommand{\headrulewidth}{0pt}
	\renewcommand{\footrulewidth}{0pt}
}
%=================== HẾT PHẦN GỘP TỪ eduzy-de.tex ===================
\newcommand{\indam}[1]{
	{\fontsize{12pt}{9pt}\sffamily \selectfont\bfseries\color{gold1} #1}
}
\newcommand{\ind}[1]{
	{\fontsize{14pt}{9pt}\fontfamily{qag} \selectfont\bfseries\color{gold1} #1}
}
\renewcommand{\FalseEX}{\stepcounter{dapan}\textbf{\selectfont\color{gold1}\Alph{dapan}.}\ignorespaces}
\renewtheorem{ex}{\color{gold1}Câu}
\renewcommand{\baselinestretch}{1.1}
\newcommand{\hoac}[1]{
	\left[\begin{aligned}#1\end{aligned}\right.}
\newcommand{\heva}[1]{
	\left\{\begin{aligned}#1\end{aligned}\right.}
\usepackage{esvect}
\def\vec{\vv}
\def\vt{\vv}
\def\overrightarrow{\vv}
%====================== LOGO EDUZY VẼ BẰNG TIKZ ======================
%  Đường viền dò từ logo gốc; dải màu vàng kim đo trực tiếp trên ảnh gốc
%  (nghiêng 21 độ, tô bằng 80 dải hẹp nên hiển thị giống nhau trên mọi trình xem PDF).
%  Dùng: \eduzylogo          -> cao 2.4cm (mặc định)
%        \eduzylogo[1.5cm]   -> cao 1.5cm
\definecolor{ezText}{RGB}{245,200,46}
\newcommand{\eduzylogo}[1][2.4cm]{%
	\resizebox{!}{#1}{%
		\begin{tikzpicture}[x=0.01cm,y=0.01cm]
			\path[use as bounding box] (0,0) rectangle (224,-224);
			% mũ tốt nghiệp + vương miện
			\begin{scope}
				\clip
				(119.4,-176.1)..controls(118.9,-176.1)and(116.8,-176)..(114.6,-176)
				..controls(101.4,-176)and(85.8,-175.9)..(83.7,-175.7)
				..controls(81.8,-175.6)and(81,-175.6)..(80.2,-175.7)
				..controls(78.8,-175.9)and(77.7,-175.9)..(76.2,-175.7)
				..controls(74.9,-175.5)and(70.8,-175.5)..(64.3,-175.7)
				..controls(59.1,-175.9)and(58.3,-175.7)..(58.6,-173.7)
				..controls(58.6,-172.8)and(59.6,-169.6)..(60.3,-168)
				..controls(60.8,-166.9)and(61.4,-165.8)..(62.3,-165)
				..controls(63.1,-164.2)and(64.3,-163.4)..(65,-163.2)
				..controls(65.5,-163.1)and(72.3,-163)..(76.6,-163.1)
				..controls(77.9,-163.1)and(80.1,-163.1)..(81.4,-163)
				..controls(82.8,-162.9)and(86.3,-162.9)..(89.2,-162.9)
				..controls(93,-162.9)and(94.9,-162.8)..(96,-162.7)
				..controls(97.2,-162.5)and(101,-162.5)..(114,-162.4)
				..controls(123,-162.3)and(134.9,-162.2)..(140.2,-162.1)
				..controls(150.1,-162)and(152.2,-162.1)..(153.6,-162.4)
				..controls(154.4,-162.6)and(157.1,-163.8)..(158.1,-164.4)
				..controls(158.8,-164.8)and(159.7,-165.7)..(159.9,-166.3)
				..controls(160,-166.5)and(160.2,-167.1)..(160.2,-167.6)
				..controls(160.3,-168)and(160.6,-169.4)..(160.9,-170.6)
				..controls(161.2,-171.7)and(161.5,-173.1)..(161.6,-173.6)
				..controls(161.8,-174.4)and(161.8,-174.5)..(161.5,-174.9)
				..controls(161.2,-175.6)and(160.8,-175.9)..(159.9,-176)
				..controls(159.2,-176.1)and(154.1,-176.2)..(151.5,-176.2)
				..controls(148.1,-176.1)and(122.6,-176.1)..(121.6,-176.2)
				..controls(120.8,-176.2)and(119.8,-176.2)..(119.4,-176.1)--cycle
				(130.4,-160)..controls(129,-160)and(119.8,-160)..(110.1,-159.9)
				..controls(100.3,-159.8)and(88.5,-159.7)..(83.9,-159.6)
				..controls(77.8,-159.5)and(75.3,-159.5)..(74.5,-159.6)
				..controls(73.9,-159.7)and(72.5,-159.8)..(71.3,-159.8)
				..controls(69.6,-159.8)and(69,-159.8)..(68.5,-159.6)
				..controls(67.6,-159.3)and(66.8,-158.6)..(66.2,-157.6)
				..controls(66,-157.3)and(65.6,-156.5)..(65.3,-156.1)
				..controls(64.6,-155.1)and(64.5,-154.9)..(63.7,-152.9)
				..controls(63.2,-151.6)and(62.9,-150.9)..(61.8,-149.6)
				..controls(61.5,-149.2)and(61.1,-148.5)..(60.9,-147.9)
				..controls(60.7,-147.4)and(60.3,-146.6)..(60.1,-146.1)
				..controls(59.9,-145.5)and(59.5,-144.9)..(59,-144.3)
				..controls(58.5,-143.7)and(58.2,-143.2)..(57.9,-142.5)
				..controls(57.7,-141.9)and(57.3,-141.1)..(57,-140.7)
				..controls(56.6,-140.3)and(56.1,-139.3)..(55.8,-138.5)
				..controls(55.5,-137.8)and(55.1,-136.8)..(54.9,-136.5)
				..controls(54.7,-136.1)and(54.3,-135.2)..(53.9,-134.5)
				..controls(53.6,-133.8)and(53,-132.7)..(52.6,-132)
				..controls(51.6,-130.5)and(51.4,-130.2)..(51.1,-129)
				..controls(50.6,-127.5)and(50.6,-127.4)..(49.7,-126.3)
				..controls(49.2,-125.7)and(48.6,-124.8)..(48.3,-124.2)
				..controls(48,-123.6)and(47.5,-122.7)..(47.3,-122.2)
				..controls(47,-121.8)and(46.5,-120.7)..(46.1,-119.8)
				..controls(45.8,-118.9)and(45.2,-117.8)..(44.9,-117.2)
				..controls(44.2,-116.1)and(43.3,-114.4)..(42.3,-112.2)
				..controls(41.8,-111.3)and(41.3,-110.4)..(41.1,-110.1)
				..controls(40.5,-109.4)and(40.3,-109)..(40,-107.8)
				..controls(39.8,-107)and(39.6,-106.5)..(39.2,-106)
				..controls(38.9,-105.6)and(38.4,-104.7)..(38,-104)
				..controls(37.7,-103.4)and(37.2,-102.5)..(37,-102.1)
				..controls(36.7,-101.7)and(36.3,-100.8)..(36,-100.1)
				..controls(35.6,-98.9)and(35.3,-98.4)..(34.6,-97.6)
				..controls(34.2,-97.1)and(34.2,-96.3)..(34.6,-96)
				..controls(35,-95.7)and(35.5,-95.9)..(36.2,-96.7)
				..controls(36.5,-97)and(37.2,-97.5)..(37.6,-97.7)..controls(38,-98)and(38.6,-98.4)..(39.1,-98.9)
				..controls(39.6,-99.4)and(40.3,-100)..(40.6,-100.3)
				..controls(41.3,-100.8)and(42.4,-101.8)..(43.7,-102.9)
				..controls(44,-103.2)and(44.6,-103.7)..(45,-103.9)
				..controls(45.4,-104.2)and(45.9,-104.6)..(46.1,-104.8)
				..controls(47.1,-105.9)and(48.8,-107.1)..(50.5,-107.9)
				..controls(51.8,-108.6)and(52.3,-109)..(53.4,-110)
				..controls(54.8,-111.2)and(55.6,-111.6)..(57,-112.1)
				..controls(57.5,-112.3)and(58.5,-112.8)..(59.4,-113.3)
				..controls(60.2,-113.7)and(61.3,-114.3)..(61.9,-114.5)
				..controls(62.4,-114.7)and(63.3,-115.1)..(63.8,-115.3)
				..controls(64.4,-115.6)and(65.4,-115.9)..(66.2,-116.1)
				..controls(67,-116.3)and(68.4,-116.6)..(69.2,-116.9)
				..controls(72.3,-117.9)and(72.8,-118)..(78,-118.1)
				..controls(84,-118.3)and(86.1,-118.1)..(88.1,-117.1)
				..controls(89.4,-116.5)and(93,-115)..(93.7,-114.9)
				..controls(94.8,-114.6)and(95.4,-114.2)..(95.9,-113.4)
				..controls(96.2,-113)and(96.7,-112.5)..(97.1,-112.2)
				..controls(98,-111.6)and(101.3,-108.4)..(102.1,-107.2)
				..controls(102.4,-106.8)and(102.9,-106.2)..(103.2,-105.8)
				..controls(103.5,-105.5)and(104,-104.8)..(104.2,-104.2)
				..controls(104.5,-103.6)and(105,-102.8)..(105.3,-102.4)
				..controls(105.6,-102)and(106.1,-101.2)..(106.5,-100.6)
				..controls(107.1,-99.5)and(108.6,-97.5)..(109.5,-96.5)
				..controls(109.9,-96.1)and(110.1,-96)..(110.4,-96)
				..controls(111.1,-96)and(111.6,-96.5)..(111.8,-97.6)
				..controls(112,-98.6)and(112.3,-99.3)..(112.8,-99.9)
				..controls(113.7,-100.9)and(115.2,-103.1)..(115.7,-104.2)
				..controls(116.1,-105)and(116.5,-105.5)..(116.9,-106)
				..controls(117.3,-106.4)and(117.8,-106.9)..(118,-107.3)
				..controls(118.2,-107.6)and(118.6,-108)..(119,-108.3)
				..controls(119.6,-108.8)and(121.1,-110.3)..(121.8,-111.2)
				..controls(122.1,-111.7)and(122.5,-112)..(123.4,-112.4)
				..controls(124.8,-113.2)and(125.8,-113.8)..(127,-114.6)
				..controls(127.5,-114.9)and(128.1,-115.3)..(128.2,-115.4)
				..controls(128.4,-115.5)and(129.3,-115.8)..(130.2,-116)
				..controls(131.1,-116.2)and(132.5,-116.6)..(133.3,-116.9)
				..controls(135.7,-117.9)and(136.1,-117.9)..(142.5,-117.8)
				..controls(145.8,-117.8)and(148.4,-117.7)..(148.8,-117.6)
				..controls(149.2,-117.6)and(150.2,-117.2)..(151,-116.8)
				..controls(152,-116.4)and(152.8,-116)..(153.6,-115.9)
				..controls(154.5,-115.7)and(155.2,-115.4)..(156.3,-114.9)
				..controls(157.1,-114.5)and(158.2,-114)..(158.7,-113.8)
				..controls(159.1,-113.7)and(159.8,-113.3)..(160.2,-113)
				..controls(160.6,-112.7)and(161.3,-112.4)..(161.7,-112.2)
				..controls(162.7,-111.9)and(163.2,-111.7)..(164.1,-111.1)
				..controls(164.5,-110.8)and(165.3,-110.4)..(165.9,-110)
				..controls(166.5,-109.7)and(167.3,-109.2)..(167.6,-108.9)
				..controls(167.9,-108.6)and(168.6,-108.2)..(169,-108)
				..controls(169.4,-107.8)and(170.2,-107.3)..(170.7,-106.9)
				..controls(171.2,-106.5)and(172.1,-105.8)..(172.7,-105.4)
				..controls(173.2,-105)and(174.1,-104.4)..(174.5,-104)
				..controls(174.9,-103.7)and(175.7,-103.1)..(176.2,-102.6)
				..controls(176.7,-102.2)and(177.4,-101.5)..(177.9,-101.2)
				..controls(178.3,-100.8)and(178.8,-100.4)..(179,-100.2)
				..controls(179.2,-100)and(179.6,-99.6)..(180,-99.3)
				..controls(180.4,-99.1)and(180.9,-98.6)..(181.2,-98.4)
				..controls(181.5,-98.1)and(182.1,-97.7)..(182.5,-97.5)
				..controls(183.2,-97.2)and(183.5,-96.9)..(184.3,-95.9)
				..controls(184.6,-95.6)and(184.6,-95.6)..(184.8,-95.7)
				..controls(185,-95.8)and(185.1,-96)..(185.1,-96.2)
				..controls(185.1,-96.5)and(185,-96.7)..(184.6,-97.2)
				..controls(184.1,-97.8)and(184,-97.9)..(183.9,-98.8)
				..controls(183.7,-100.3)and(183.3,-101.3)..(182.9,-101.8)
				..controls(182.6,-102.1)and(182.3,-102.6)..(182.2,-102.9)
				..controls(182,-103.3)and(181.6,-104)..(181.2,-104.6)
				..controls(180.9,-105.2)and(180.5,-106.1)..(180.4,-106.6)
				..controls(179.9,-108)and(179.7,-108.5)..(179,-109.4)
				..controls(178.6,-109.8)and(178.2,-110.5)..(177.9,-111)
				..controls(177.7,-111.5)and(177.3,-112.3)..(177,-112.7)
				..controls(176.6,-113.2)and(176.4,-113.7)..(176.3,-114.2)
				..controls(176.2,-114.7)and(175.9,-115.2)..(175.4,-116)
				..controls(175,-116.6)and(174.4,-117.7)..(174,-118.5)
				..controls(173.7,-119.2)and(173.3,-120.1)..(173,-120.5)
				..controls(172.8,-120.9)and(172.4,-121.8)..(172.1,-122.5)
				..controls(171.8,-123.2)and(171.3,-124.1)..(171,-124.6)
				..controls(170.8,-125.1)and(170.3,-126)..(169.9,-126.6)
				..controls(169.6,-127.3)and(169.2,-128.1)..(169,-128.3)
				..controls(168.8,-128.6)and(168.5,-129.3)..(168.3,-129.8)
				..controls(168.1,-130.3)and(167.8,-131.1)..(167.5,-131.6)
				..controls(166.9,-132.5)and(164.2,-137.9)..(164.2,-138.1)
				..controls(164.2,-138.1)and(163.9,-138.8)..(163.4,-139.7)
				..controls(162.9,-140.5)and(162.4,-141.5)..(162.3,-141.9)
				..controls(161.9,-143)and(161.6,-143.4)..(160.9,-144.1)
				..controls(160.4,-144.6)and(160.3,-144.8)..(160.1,-145.4)
				..controls(160,-145.8)and(159.6,-146.8)..(159.2,-147.8)
				..controls(158.7,-148.8)and(158.1,-150)..(157.9,-150.6)
				..controls(157.6,-151.1)and(157.2,-151.9)..(156.9,-152.4)
				..controls(156.5,-152.8)and(156.1,-153.7)..(155.8,-154.4)
				..controls(155.6,-155)and(155.2,-155.9)..(154.9,-156.3)
				..controls(154.7,-156.7)and(154.4,-157.4)..(154.2,-157.9)
				..controls(154.1,-158.4)and(153.8,-159)..(153.5,-159.3)
				..controls(152.8,-160.1)and(153,-160.1)..(142.4,-160.1)
				..controls(137.2,-160.1)and(131.8,-160.1)..(130.4,-160)--cycle
				(141.5,-103.7)..controls(140.8,-103)and(139.8,-102.1)..(139.2,-101.6)
				..controls(138.5,-101.2)and(137.9,-100.6)..(137.8,-100.4)
				..controls(137,-99.5)and(134,-96.5)..(133.2,-96)
				..controls(132.7,-95.7)and(131.8,-94.9)..(131.1,-94.3)
				..controls(129.4,-92.7)and(127.9,-91.6)..(126.6,-91.1)
				..controls(125.2,-90.5)and(124.1,-90)..(122.9,-89.1)
				..controls(122.1,-88.6)and(121.3,-88.2)..(119.9,-87.7)
				..controls(118.8,-87.3)and(117.5,-86.8)..(116.9,-86.5)
				..controls(114.3,-85.5)and(113.7,-85.4)..(109.7,-85.5)
				..controls(105.3,-85.7)and(103.9,-85.9)..(101.6,-87)
				..controls(100.7,-87.4)and(99.4,-87.8)..(97.9,-88)
				..controls(97,-88.2)and(96.6,-88.3)..(95.8,-88.9)
				..controls(95.4,-89.2)and(94.7,-89.6)..(94.2,-89.8)
				..controls(92.9,-90.3)and(92.3,-90.7)..(90.9,-92.1)
				..controls(90.2,-92.7)and(89.4,-93.3)..(89,-93.5)..controls(87,-94.5)and(85.2,-96)..(83.6,-97.9)
				..controls(83.3,-98.2)and(82.5,-98.9)..(81.8,-99.4)
				..controls(80.7,-100.1)and(80.5,-100.3)..(80,-101.1)
				..controls(78.7,-103.1)and(77.2,-103.9)..(76.4,-103)
				..controls(75.7,-102.2)and(75.7,-103.2)..(75.7,-83.6)
				..controls(75.6,-73.6)and(75.6,-64.8)..(75.5,-64)
				..controls(75.5,-63.2)and(75.5,-61.4)..(75.6,-59.8)--(75.7,-57.1)--(76,-56.8)
				..controls(76.8,-56.2)and(78,-56.3)..(79.8,-57.2)
				..controls(80.4,-57.5)and(81.3,-57.8)..(81.8,-57.9)
				..controls(82.3,-58.1)and(83,-58.4)..(83.6,-58.8)..controls(84.3,-59.2)and(85,-59.5)..(86,-59.8)
				..controls(86.8,-60.1)and(88,-60.6)..(88.7,-60.9)
				..controls(89.3,-61.2)and(90.6,-61.7)..(91.4,-62)
				..controls(92.2,-62.3)and(93.7,-62.9)..(94.6,-63.3)
				..controls(95.6,-63.7)and(96.8,-64.1)..(97.4,-64.3)
				..controls(98,-64.5)and(99,-64.9)..(99.7,-65.2)
				..controls(100.4,-65.6)and(101.5,-66)..(102.3,-66.2)
				..controls(104.2,-66.8)and(105,-67.1)..(106.2,-67.8)
				..controls(107.9,-68.7)and(110.2,-69.4)..(111.3,-69.2)
				..controls(112.5,-69)and(115,-68.2)..(115.6,-67.8)
				..controls(117.3,-66.7)and(117.6,-66.6)..(119.5,-66.2)
				..controls(121.9,-65.8)and(122.1,-65.7)..(123,-65.1)
				..controls(124.1,-64.4)and(124.4,-64.3)..(126,-63.9)
				..controls(127,-63.7)and(127.7,-63.4)..(128.4,-63.1)
				..controls(129,-62.8)and(129.7,-62.5)..(130.6,-62.4)
				..controls(132.1,-62.1)and(134,-61.5)..(135.2,-61)
				..controls(135.7,-60.8)and(136.7,-60.5)..(137.3,-60.3)
				..controls(138.4,-60)and(139.3,-59.7)..(141.1,-59)
				..controls(141.6,-58.7)and(142,-58.6)..(142.5,-58.6)
				..controls(143.2,-58.6)and(143.3,-58.7)..(143.6,-59)
				..controls(144,-59.5)and(144.1,-59.9)..(144,-60.8)
				..controls(144,-61.2)and(144,-71.1)..(144,-82.8)--(143.9,-104.1)--(143.7,-104.5)
				..controls(143.5,-104.7)and(143.3,-104.8)..(143,-104.9)
				..controls(142.6,-104.9)and(142.5,-104.8)..(141.5,-103.7)--cycle
				(27.8,-91.4)..controls(27.1,-90.9)and(26.3,-90)..(26.1,-89.5)
				..controls(26,-89.2)and(25.8,-88.6)..(25.7,-87.9)
				..controls(25.4,-86.2)and(25.6,-84.8)..(26.5,-83.5)
				..controls(27.4,-82.3)and(28.3,-82)..(29.4,-82.3)
				..controls(31.1,-82.8)and(32.1,-84.9)..(32,-87.4)
				..controls(31.9,-89.3)and(31.2,-90.6)..(29.9,-91.4)
				..controls(29.1,-91.9)and(28.6,-91.9)..(27.8,-91.4)--cycle
				(190.1,-91.3)..controls(189.7,-91.2)and(189.3,-90.9)..(189.1,-90.7)
				..controls(188.3,-89.8)and(187.6,-87.7)..(187.6,-86.2)
				..controls(187.6,-85.5)and(187.7,-85.2)..(188.2,-84.4)
				..controls(189.1,-82.6)and(190.1,-81.8)..(191.3,-82.1)
				..controls(192,-82.3)and(193.5,-83.7)..(193.8,-84.4)
				..controls(194.1,-85.1)and(194.2,-87.1)..(194,-88.3)
				..controls(193.8,-89.8)and(193.3,-90.5)..(192.3,-91.3)
				..controls(191.7,-91.7)and(190.9,-91.7)..(190.1,-91.3)--cycle
				(173.7,-83.4)..controls(172.2,-82.5)and(172,-82.5)..(170.6,-82.4)
				..controls(169.7,-82.4)and(168.8,-82.4)..(168.3,-82.5)
				..controls(167.2,-82.7)and(166.9,-82.7)..(166.6,-82.2)
				..controls(166.3,-81.8)and(166.3,-81)..(166.8,-80)
				..controls(167.2,-78.9)and(167.4,-77.4)..(167.7,-74.1)
				..controls(167.8,-72.8)and(167.9,-71.8)..(167.9,-71.7)
				..controls(168,-71.7)and(168,-70.2)..(168,-68.4)..controls(168,-65.4)and(168,-64.9)..(167.7,-63)
				..controls(167.3,-60.5)and(167.1,-59.9)..(166.4,-58.4)
				..controls(166.1,-57.7)and(165.8,-56.9)..(165.7,-56.6)
				..controls(165.7,-56.2)and(165.4,-55.6)..(165.2,-55.2)
				..controls(164.6,-53.8)and(161.9,-50.7)..(161,-50.3)
				..controls(159.8,-49.7)and(157.9,-50)..(155.2,-51.2)
				..controls(154.7,-51.4)and(153.4,-51.8)..(152.2,-52.1)
				..controls(150.8,-52.4)and(149.4,-52.9)..(148.6,-53.2)
				..controls(147.8,-53.6)and(146.7,-54)..(146,-54.2)
				..controls(145.4,-54.4)and(144.5,-54.7)..(144,-54.9)
				..controls(143.6,-55.2)and(142.7,-55.5)..(141.8,-55.7)
				..controls(141,-55.9)and(139.6,-56.3)..(138.7,-56.6)
				..controls(137.8,-57)and(136.3,-57.4)..(135.3,-57.6)
				..controls(133.8,-57.9)and(133.2,-58.1)..(132.1,-58.7)
				..controls(131.1,-59.2)and(130.3,-59.4)..(129.4,-59.6)
				..controls(128.7,-59.8)and(127.8,-60)..(127.4,-60.1)
				..controls(127,-60.3)and(126,-60.5)..(125.1,-60.7)
				..controls(123.2,-61.1)and(121.7,-61.5)..(120.2,-62.1)
				..controls(119.7,-62.3)and(118.1,-62.8)..(116.7,-63.2)
				..controls(114.6,-63.9)and(114,-64.1)..(113,-64.7)
				..controls(111.9,-65.4)and(111.8,-65.4)..(111,-65.5)
				..controls(109.9,-65.5)and(109.3,-65.4)..(107.7,-64.7)
				..controls(107,-64.4)and(106,-64)..(105.3,-63.9)
				..controls(103.5,-63.5)and(102.2,-63)..(101.2,-62.5)
				..controls(100.1,-61.8)and(99.3,-61.5)..(97.8,-61.1)
				..controls(97.2,-61)and(96.2,-60.7)..(95.7,-60.4)
				..controls(95.1,-60.1)and(94.2,-59.8)..(93.5,-59.6)
				..controls(92.9,-59.4)and(91.7,-59)..(91,-58.7)..controls(90.2,-58.4)and(88.9,-58)..(88,-57.7)
				..controls(87.2,-57.5)and(86.1,-57.1)..(85.5,-56.8)
				..controls(85,-56.6)and(84.1,-56.2)..(83.5,-56)..controls(82.9,-55.7)and(81.8,-55.2)..(81,-54.9)
				..controls(80.1,-54.5)and(79.1,-54.2)..(78.7,-54.1)
				..controls(77.9,-53.8)and(77.3,-53.6)..(76.2,-52.9)
				..controls(75.7,-52.6)and(75.2,-52.4)..(74.1,-52.2)
				..controls(72.5,-51.9)and(72.4,-51.8)..(71.3,-51.1)
				..controls(69.9,-50.1)and(69.4,-49.9)..(67.9,-49.8)
				..controls(67,-49.7)and(66.6,-49.6)..(66.3,-49.4)
				..controls(66.1,-49.3)and(65.7,-49)..(65.4,-48.8)
				..controls(65.1,-48.7)and(64.2,-48.4)..(63.6,-48.2)
				..controls(62.9,-48)and(61.7,-47.6)..(60.9,-47.2)
				..controls(60.1,-46.8)and(58.9,-46.4)..(58.2,-46.2)
				..controls(57.6,-46)and(56.8,-45.7)..(56.5,-45.5)
				..controls(54.9,-44.6)and(54.6,-44.5)..(53.2,-44.2)
				..controls(51.7,-43.8)and(51,-43.6)..(49.9,-42.9)
				..controls(49.5,-42.6)and(48.9,-42.4)..(47.5,-42)
				..controls(45.7,-41.6)and(45.6,-41.5)..(44.3,-40.6)
				..controls(43.5,-40.1)and(42.8,-39.6)..(42.7,-39.4)
				..controls(42.4,-38.7)and(43,-38.1)..(44.5,-37.6)..controls(44.9,-37.5)and(46,-37)..(46.9,-36.6)
				..controls(47.9,-36.1)and(48.8,-35.8)..(49.6,-35.6)
				..controls(50.5,-35.4)and(51.3,-35.2)..(52.4,-34.6)
				..controls(53.3,-34.1)and(54.5,-33.7)..(54.9,-33.5)
				..controls(55.4,-33.4)and(56.4,-33)..(57.2,-32.6)
				..controls(58.6,-32)and(60.1,-31.5)..(61.6,-31.2)
				..controls(62.3,-31.1)and(62.7,-31)..(63.2,-30.7)
				..controls(64.1,-30)and(64.3,-29.9)..(65.3,-29.7)
				..controls(65.7,-29.6)and(66.5,-29.2)..(67.1,-28.9)
				..controls(67.6,-28.6)and(68.8,-28.2)..(69.6,-27.9)
				..controls(70.5,-27.7)and(71.6,-27.3)..(72.1,-27)
				..controls(72.6,-26.7)and(73.8,-26.2)..(74.8,-25.9)
				..controls(75.7,-25.5)and(76.7,-25.1)..(77.1,-25)
				..controls(77.4,-24.8)and(78.5,-24.4)..(79.6,-24.1)
				..controls(80.6,-23.8)and(82,-23.2)..(82.8,-22.9)
				..controls(83.5,-22.5)and(84.8,-22)..(85.8,-21.8)
				..controls(86.7,-21.5)and(87.8,-21.1)..(88.1,-20.9)
				..controls(88.4,-20.7)and(89.3,-20.3)..(90.1,-20.1)
				..controls(90.8,-19.8)and(91.9,-19.4)..(92.4,-19.1)
				..controls(93.5,-18.5)and(94.4,-18.2)..(95.4,-18.1)
				..controls(95.7,-18)and(96.4,-17.8)..(96.9,-17.5)
				..controls(97.5,-17.3)and(98.4,-16.9)..(99,-16.6)
				..controls(99.7,-16.3)and(100.9,-15.9)..(101.9,-15.6)
				..controls(102.9,-15.3)and(104.4,-14.7)..(105.3,-14.3)
				..controls(108.1,-13.2)and(108.5,-13.1)..(110.1,-13.1)
				..controls(111.4,-13.1)and(111.5,-13.1)..(112.4,-13.5)
				..controls(113.2,-14)and(115,-14.8)..(116.6,-15.4)
				..controls(116.9,-15.5)and(117.9,-15.8)..(118.8,-16.1)
				..controls(120.9,-16.6)and(121.5,-16.9)..(122.6,-17.9)
				..controls(123,-18.3)and(123.7,-18.8)..(123.9,-18.9)
				..controls(124.6,-19.4)and(126,-19.8)..(127.9,-20.2)
				..controls(129.6,-20.5)and(129.8,-20.6)..(131.1,-21.4)
				..controls(131.5,-21.7)and(132.4,-22.1)..(133.1,-22.3)
				..controls(136.1,-23.2)and(138.3,-24.1)..(140.2,-25)
				..controls(141,-25.5)and(141.5,-25.6)..(142.5,-25.8)
				..controls(145.1,-26.3)and(146,-26.6)..(149.4,-28.1)
				..controls(151.3,-28.9)and(153.5,-29.8)..(154.3,-30.1)
				..controls(155.1,-30.4)and(156.2,-30.8)..(156.6,-31)
				..controls(157.1,-31.2)and(158.2,-31.6)..(159.1,-31.8)
				..controls(160.6,-32.2)and(160.9,-32.3)..(161.7,-32.9)
				..controls(162.4,-33.3)and(162.9,-33.6)..(163.7,-33.8)
				..controls(164.3,-34)and(164.8,-34.1)..(164.9,-34.1)
				..controls(165,-34.1)and(165.6,-34.4)..(166.4,-34.8)
				..controls(167.3,-35.2)and(168.1,-35.4)..(169,-35.6)
				..controls(170,-35.8)and(170.5,-36)..(171.7,-36.6)
				..controls(172.4,-37)and(173.3,-37.4)..(173.6,-37.4)
				..controls(175.6,-37.7)and(176.2,-38)..(176.7,-39.1)
				..controls(176.9,-39.6)and(176.9,-39.7)..(176.7,-40.4)
				..controls(176.4,-41.6)and(175.8,-42.7)..(175.2,-43.4)
				..controls(174.5,-44.1)and(172.6,-45.5)..(171.7,-45.8)
				..controls(170.6,-46.2)and(169.1,-46)..(168.2,-45.4)
				..controls(167.9,-45.2)and(167.3,-44.8)..(166.9,-44.4)
				..controls(166.4,-44.1)and(165.8,-43.6)..(165.4,-43.2)
				..controls(165.1,-42.9)and(164.4,-42.3)..(163.9,-41.9)
				..controls(163.3,-41.5)and(162.7,-40.9)..(162.4,-40.5)
				..controls(161.4,-39.2)and(159.9,-38)..(158.8,-37.8)
				..controls(158.4,-37.8)and(158.3,-37.7)..(158.2,-37.2)
				..controls(158,-36.8)and(157.4,-36.1)..(157.1,-36.1)
				..controls(157,-36.1)and(156.7,-36.1)..(156.5,-36)
				..controls(156.1,-35.9)and(155.9,-35.7)..(155.7,-35.2)
				..controls(155.3,-34.6)and(155.1,-34.4)..(154,-34.1)
				..controls(153.6,-34)and(153.2,-33.7)..(152.9,-33.5)
				..controls(152.4,-32.9)and(151.6,-32.4)..(151.3,-32.4)
				..controls(150.9,-32.4)and(150.4,-32.6)..(150.4,-32.9)
				..controls(150.4,-33.2)and(150.8,-33.4)..(151.4,-33.6)
				..controls(152.1,-33.7)and(152.8,-34.1)..(153.5,-34.9)
				..controls(154.1,-35.5)and(154.3,-35.6)..(155,-35.8)
				..controls(155.7,-36)and(155.7,-36)..(156,-36.6)
				..controls(156.3,-37.3)and(156.7,-37.7)..(157.5,-37.8)
				..controls(158.1,-37.9)and(158.2,-38)..(158.4,-38.8)
				..controls(158.5,-39.4)and(158.6,-39.6)..(159.6,-40.4)
				..controls(160.2,-41)and(160.9,-41.6)..(161.3,-41.8)
				..controls(161.7,-42)and(161.9,-42.3)..(162.2,-42.8)
				..controls(162.4,-43.2)and(162.8,-43.8)..(163.1,-44.1)
				..controls(163.9,-44.8)and(165.1,-46.7)..(165.7,-48)
				..controls(166.2,-49.3)and(166.4,-49.7)..(167.2,-50.5)
				..controls(167.5,-50.8)and(167.9,-51.4)..(168.1,-51.8)
				..controls(168.3,-52.2)and(168.8,-53)..(169.1,-53.6)
				..controls(169.6,-54.3)and(169.8,-54.9)..(169.9,-55.6)
				..controls(170.4,-57.8)and(170.5,-58.1)..(171.2,-59.5)
				..controls(171.6,-60.3)and(172,-61.5)..(172.2,-62.1)
				..controls(172.3,-62.8)and(172.7,-63.7)..(172.9,-64.1)
				..controls(173.2,-64.6)and(173.4,-65.3)..(173.4,-65.6)
				..controls(173.8,-68.8)and(174.2,-70.6)..(174.9,-72.6)
				..controls(175.5,-74.5)and(175.9,-77.1)..(176.1,-80.2)
				..controls(176.3,-83.4)and(175.5,-84.4)..(173.7,-83.4)--cycle;
				\definecolor{ezs}{RGB}{247,179,26}\fill[ezs](-5.7,-156.6)--(-3.4,-157.5)--(61.9,12.7)--(59.7,13.6)--cycle;
				\definecolor{ezs}{RGB}{247,179,22}\fill[ezs](-3.8,-157.3)--(-1.5,-158.2)--(63.9,12)--(61.6,12.9)--cycle;
				\definecolor{ezs}{RGB}{248,179,19}\fill[ezs](-1.8,-158.1)--(0.4,-158.9)--(65.8,11.2)--(63.5,12.1)--cycle;
				\definecolor{ezs}{RGB}{248,178,16}\fill[ezs](0.1,-158.8)--(2.4,-159.7)--(67.7,10.5)--(65.4,11.4)--cycle;
				\definecolor{ezs}{RGB}{249,178,15}\fill[ezs](2,-159.5)--(4.3,-160.4)--(69.6,9.8)--(67.3,10.7)--cycle;
				\definecolor{ezs}{RGB}{251,177,17}\fill[ezs](3.9,-160.3)--(6.2,-161.1)--(71.5,9)--(69.2,9.9)--cycle;
				\definecolor{ezs}{RGB}{252,179,18}\fill[ezs](5.8,-161)--(8.1,-161.9)--(73.4,8.3)--(71.1,9.2)--cycle;
				\definecolor{ezs}{RGB}{252,182,17}\fill[ezs](7.7,-161.7)--(10,-162.6)--(75.3,7.6)--(73,8.5)--cycle;
				\definecolor{ezs}{RGB}{252,185,17}\fill[ezs](9.6,-162.5)--(11.9,-163.4)--(77.2,6.8)--(75,7.7)--cycle;
				\definecolor{ezs}{RGB}{251,189,18}\fill[ezs](11.5,-163.2)--(13.8,-164.1)--(79.2,6.1)--(76.9,7)--cycle;
				\definecolor{ezs}{RGB}{251,191,22}\fill[ezs](13.4,-163.9)--(15.7,-164.8)--(81.1,5.4)--(78.8,6.2)--cycle;
				\definecolor{ezs}{RGB}{252,193,25}\fill[ezs](15.4,-164.7)--(17.6,-165.6)--(83,4.6)--(80.7,5.5)--cycle;
				\definecolor{ezs}{RGB}{253,195,27}\fill[ezs](17.3,-165.4)--(19.6,-166.3)--(84.9,3.9)--(82.6,4.8)--cycle;
				\definecolor{ezs}{RGB}{253,196,31}\fill[ezs](19.2,-166.1)--(21.5,-167)--(86.8,3.2)--(84.5,4)--cycle;
				\definecolor{ezs}{RGB}{254,199,35}\fill[ezs](21.1,-166.9)--(23.4,-167.8)--(88.7,2.4)--(86.4,3.3)--cycle;
				\definecolor{ezs}{RGB}{253,201,38}\fill[ezs](23,-167.6)--(25.3,-168.5)--(90.6,1.7)--(88.3,2.6)--cycle;
				\definecolor{ezs}{RGB}{253,203,40}\fill[ezs](24.9,-168.3)--(27.2,-169.2)--(92.5,1)--(90.3,1.8)--cycle;
				\definecolor{ezs}{RGB}{253,205,42}\fill[ezs](26.8,-169.1)--(29.1,-170)--(94.4,0.2)--(92.2,1.1)--cycle;
				\definecolor{ezs}{RGB}{253,207,45}\fill[ezs](28.7,-169.8)--(31,-170.7)--(96.4,-0.5)--(94.1,0.4)--cycle;
				\definecolor{ezs}{RGB}{253,210,49}\fill[ezs](30.7,-170.5)--(32.9,-171.4)--(98.3,-1.2)--(96,-0.4)--cycle;
				\definecolor{ezs}{RGB}{253,212,54}\fill[ezs](32.6,-171.3)--(34.9,-172.2)--(100.2,-2)--(97.9,-1.1)--cycle;
				\definecolor{ezs}{RGB}{253,214,59}\fill[ezs](34.5,-172)--(36.8,-172.9)--(102.1,-2.7)--(99.8,-1.8)--cycle;
				\definecolor{ezs}{RGB}{253,216,63}\fill[ezs](36.4,-172.7)--(38.7,-173.6)--(104,-3.4)--(101.7,-2.6)--cycle;
				\definecolor{ezs}{RGB}{253,218,65}\fill[ezs](38.3,-173.5)--(40.6,-174.4)--(105.9,-4.2)--(103.6,-3.3)--cycle;
				\definecolor{ezs}{RGB}{254,220,66}\fill[ezs](40.2,-174.2)--(42.5,-175.1)--(107.8,-4.9)--(105.5,-4)--cycle;
				\definecolor{ezs}{RGB}{254,222,68}\fill[ezs](42.1,-175)--(44.4,-175.8)--(109.7,-5.6)--(107.5,-4.8)--cycle;
				\definecolor{ezs}{RGB}{254,224,71}\fill[ezs](44,-175.7)--(46.3,-176.6)--(111.7,-6.4)--(109.4,-5.5)--cycle;
				\definecolor{ezs}{RGB}{254,226,75}\fill[ezs](46,-176.4)--(48.2,-177.3)--(113.6,-7.1)--(111.3,-6.2)--cycle;
				\definecolor{ezs}{RGB}{254,228,79}\fill[ezs](47.9,-177.2)--(50.1,-178)--(115.5,-7.8)--(113.2,-7)--cycle;
				\definecolor{ezs}{RGB}{254,230,82}\fill[ezs](49.8,-177.9)--(52.1,-178.8)--(117.4,-8.6)--(115.1,-7.7)--cycle;
				\definecolor{ezs}{RGB}{254,232,85}\fill[ezs](51.7,-178.6)--(54,-179.5)--(119.3,-9.3)--(117,-8.4)--cycle;
				\definecolor{ezs}{RGB}{254,234,88}\fill[ezs](53.6,-179.4)--(55.9,-180.2)--(121.2,-10)--(118.9,-9.2)--cycle;
				\definecolor{ezs}{RGB}{254,237,91}\fill[ezs](55.5,-180.1)--(57.8,-181)--(123.1,-10.8)--(120.8,-9.9)--cycle;
				\definecolor{ezs}{RGB}{254,239,94}\fill[ezs](57.4,-180.8)--(59.7,-181.7)--(125,-11.5)--(122.8,-10.6)--cycle;
				\definecolor{ezs}{RGB}{254,241,97}\fill[ezs](59.3,-181.6)--(61.6,-182.4)--(127,-12.2)--(124.7,-11.4)--cycle;
				\definecolor{ezs}{RGB}{254,243,101}\fill[ezs](61.2,-182.3)--(63.5,-183.2)--(128.9,-13)--(126.6,-12.1)--cycle;
				\definecolor{ezs}{RGB}{254,245,103}\fill[ezs](63.2,-183)--(65.4,-183.9)--(130.8,-13.7)--(128.5,-12.8)--cycle;
				\definecolor{ezs}{RGB}{254,245,105}\fill[ezs](65.1,-183.8)--(67.4,-184.6)--(132.7,-14.4)--(130.4,-13.6)--cycle;
				\definecolor{ezs}{RGB}{254,245,105}\fill[ezs](67,-184.5)--(69.3,-185.4)--(134.6,-15.2)--(132.3,-14.3)--cycle;
				\definecolor{ezs}{RGB}{254,244,103}\fill[ezs](68.9,-185.2)--(71.2,-186.1)--(136.5,-15.9)--(134.2,-15)--cycle;
				\definecolor{ezs}{RGB}{253,243,101}\fill[ezs](70.8,-186)--(73.1,-186.8)--(138.4,-16.6)--(136.1,-15.8)--cycle;
				\definecolor{ezs}{RGB}{253,241,98}\fill[ezs](72.7,-186.7)--(75,-187.6)--(140.3,-17.4)--(138,-16.5)--cycle;
				\definecolor{ezs}{RGB}{253,239,95}\fill[ezs](74.6,-187.4)--(76.9,-188.3)--(142.2,-18.1)--(140,-17.2)--cycle;
				\definecolor{ezs}{RGB}{253,236,92}\fill[ezs](76.5,-188.2)--(78.8,-189)--(144.2,-18.8)--(141.9,-18)--cycle;
				\definecolor{ezs}{RGB}{253,234,88}\fill[ezs](78.5,-188.9)--(80.7,-189.8)--(146.1,-19.6)--(143.8,-18.7)--cycle;
				\definecolor{ezs}{RGB}{253,232,85}\fill[ezs](80.4,-189.6)--(82.7,-190.5)--(148,-20.3)--(145.7,-19.4)--cycle;
				\definecolor{ezs}{RGB}{253,230,82}\fill[ezs](82.3,-190.4)--(84.6,-191.2)--(149.9,-21)--(147.6,-20.2)--cycle;
				\definecolor{ezs}{RGB}{253,228,78}\fill[ezs](84.2,-191.1)--(86.5,-192)--(151.8,-21.8)--(149.5,-20.9)--cycle;
				\definecolor{ezs}{RGB}{253,226,74}\fill[ezs](86.1,-191.8)--(88.4,-192.7)--(153.7,-22.5)--(151.4,-21.6)--cycle;
				\definecolor{ezs}{RGB}{253,224,70}\fill[ezs](88,-192.6)--(90.3,-193.4)--(155.6,-23.3)--(153.3,-22.4)--cycle;
				\definecolor{ezs}{RGB}{253,222,67}\fill[ezs](89.9,-193.3)--(92.2,-194.2)--(157.5,-24)--(155.3,-23.1)--cycle;
				\definecolor{ezs}{RGB}{253,220,65}\fill[ezs](91.8,-194)--(94.1,-194.9)--(159.5,-24.7)--(157.2,-23.8)--cycle;
				\definecolor{ezs}{RGB}{253,217,63}\fill[ezs](93.7,-194.8)--(96,-195.6)--(161.4,-25.5)--(159.1,-24.6)--cycle;
				\definecolor{ezs}{RGB}{253,215,60}\fill[ezs](95.7,-195.5)--(97.9,-196.4)--(163.3,-26.2)--(161,-25.3)--cycle;
				\definecolor{ezs}{RGB}{253,213,56}\fill[ezs](97.6,-196.2)--(99.9,-197.1)--(165.2,-26.9)--(162.9,-26)--cycle;
				\definecolor{ezs}{RGB}{252,211,53}\fill[ezs](99.5,-197)--(101.8,-197.8)--(167.1,-27.7)--(164.8,-26.8)--cycle;
				\definecolor{ezs}{RGB}{252,209,49}\fill[ezs](101.4,-197.7)--(103.7,-198.6)--(169,-28.4)--(166.7,-27.5)--cycle;
				\definecolor{ezs}{RGB}{252,207,45}\fill[ezs](103.3,-198.4)--(105.6,-199.3)--(170.9,-29.1)--(168.6,-28.2)--cycle;
				\definecolor{ezs}{RGB}{253,205,42}\fill[ezs](105.2,-199.2)--(107.5,-200)--(172.8,-29.9)--(170.6,-29)--cycle;
				\definecolor{ezs}{RGB}{253,202,40}\fill[ezs](107.1,-199.9)--(109.4,-200.8)--(174.7,-30.6)--(172.5,-29.7)--cycle;
				\definecolor{ezs}{RGB}{253,200,38}\fill[ezs](109,-200.6)--(111.3,-201.5)--(176.7,-31.3)--(174.4,-30.4)--cycle;
				\definecolor{ezs}{RGB}{253,198,35}\fill[ezs](111,-201.4)--(113.2,-202.2)--(178.6,-32.1)--(176.3,-31.2)--cycle;
				\definecolor{ezs}{RGB}{253,196,31}\fill[ezs](112.9,-202.1)--(115.2,-203)--(180.5,-32.8)--(178.2,-31.9)--cycle;
				\definecolor{ezs}{RGB}{253,194,28}\fill[ezs](114.8,-202.8)--(117.1,-203.7)--(182.4,-33.5)--(180.1,-32.6)--cycle;
				\definecolor{ezs}{RGB}{253,192,26}\fill[ezs](116.7,-203.6)--(119,-204.5)--(184.3,-34.3)--(182,-33.4)--cycle;
				\definecolor{ezs}{RGB}{253,189,24}\fill[ezs](118.6,-204.3)--(120.9,-205.2)--(186.2,-35)--(183.9,-34.1)--cycle;
				\definecolor{ezs}{RGB}{253,186,20}\fill[ezs](120.5,-205)--(122.8,-205.9)--(188.1,-35.7)--(185.8,-34.8)--cycle;
				\definecolor{ezs}{RGB}{253,184,17}\fill[ezs](122.4,-205.8)--(124.7,-206.7)--(190,-36.5)--(187.8,-35.6)--cycle;
				\definecolor{ezs}{RGB}{252,183,15}\fill[ezs](124.3,-206.5)--(126.6,-207.4)--(192,-37.2)--(189.7,-36.3)--cycle;
				\definecolor{ezs}{RGB}{252,181,14}\fill[ezs](126.3,-207.2)--(128.5,-208.1)--(193.9,-37.9)--(191.6,-37.1)--cycle;
				\definecolor{ezs}{RGB}{252,179,13}\fill[ezs](128.2,-208)--(130.4,-208.9)--(195.8,-38.7)--(193.5,-37.8)--cycle;
				\definecolor{ezs}{RGB}{252,178,12}\fill[ezs](130.1,-208.7)--(132.4,-209.6)--(197.7,-39.4)--(195.4,-38.5)--cycle;
				\definecolor{ezs}{RGB}{251,178,10}\fill[ezs](132,-209.4)--(134.3,-210.3)--(199.6,-40.1)--(197.3,-39.3)--cycle;
				\definecolor{ezs}{RGB}{251,178,10}\fill[ezs](133.9,-210.2)--(136.2,-211.1)--(201.5,-40.9)--(199.2,-40)--cycle;
				\definecolor{ezs}{RGB}{250,177,11}\fill[ezs](135.8,-210.9)--(138.1,-211.8)--(203.4,-41.6)--(201.1,-40.7)--cycle;
				\definecolor{ezs}{RGB}{250,178,11}\fill[ezs](137.7,-211.6)--(140,-212.5)--(205.3,-42.3)--(203.1,-41.5)--cycle;
				\definecolor{ezs}{RGB}{248,180,10}\fill[ezs](139.6,-212.4)--(141.9,-213.3)--(207.3,-43.1)--(205,-42.2)--cycle;
				\definecolor{ezs}{RGB}{247,181,9}\fill[ezs](141.5,-213.1)--(143.8,-214)--(209.2,-43.8)--(206.9,-42.9)--cycle;
				\definecolor{ezs}{RGB}{246,182,9}\fill[ezs](143.5,-213.8)--(145.7,-214.7)--(211.1,-44.5)--(208.8,-43.7)--cycle;
				\definecolor{ezs}{RGB}{246,182,10}\fill[ezs](145.4,-214.6)--(147.7,-215.5)--(213,-45.3)--(210.7,-44.4)--cycle;
			\end{scope}
			% chữ EDUZY (Times đậm, làm dày nét bằng 5 lớp lệch nhẹ; căn đúng khung chữ logo gốc)
			\foreach \dx/\dy in {-0.7/0,0.7/0,0/-0.3,0/0.3,0/0}
			\node[anchor=center,inner sep=0pt,outer sep=0pt,text=ezText,align=center] at (109.6+\dx,-193.6+\dy)
			{\resizebox{1.01cm}{0.22cm}{\fontencoding{T1}\fontfamily{ptm}\fontseries{b}\selectfont EDUZY}};
		\end{tikzpicture}%
	}%
}
%====================== THIẾT LẬP THÂN ĐỀ ======================
%--- Môi trường tự luận {bt}: ex_test 3.3.4 chưa khai báo sẵn nên tạo ở đây
%    (ex_test đã tự móc \fixshowans và \listenumerate cho bt)
\makeatletter
\@ifundefined{bt}{\newtheorem{bt}{\color{gold1}\namebt}}{}
\makeatother
\renewcommand{\sectionmark}[1]{\markright{\thesection\ #1}}
%--- Khung tiêu đề (logo + tối đa 4 dòng; dòng 3, 4 để trống {} thì bỏ qua)
%  Đề thi:    \khungtieude{THI THỬ THPT QUỐC GIA NĂM 2026}{Môn Toán -- Đề số 21}
%                         {(Đề thi gồm: ...)}{Thời gian làm bài: 90 phút (không kể thời gian phát đề).}
%  Chuyên đề: \khungtieude{CHUYÊN ĐỀ 3 -- DÃY SỐ VÀ CÔNG THỨC TRUY HỒI}
%                         {Bồi dưỡng học sinh giỏi Toán THPT}{...}{}
\newcommand{\khungtieude}[4]{%
	\begin{tcolorbox}[colframe=gold1,colback=gold5!45,boxrule=0.5mm]
		\begin{minipage}[c]{3cm}
			\centerline{\eduzylogo[2.4cm]}
		\end{minipage}\hspace{0.5cm}
		\begin{minipage}[c]{12.5cm}
			\vskip 0.15 cm
			\centerline{\ind{#1}}
			\vskip 0.15 cm
			\centerline{\textbf{#2}}
			\if\relax\detokenize{#3}\relax\else
				\vskip 0.15 cm
				\centerline{\textit{#3}}
			\fi
			\if\relax\detokenize{#4}\relax\else
				\vskip 0.15 cm
				\centerline{\textit{#4}}
			\fi
			\vskip 0.15 cm
		\end{minipage}
	\end{tcolorbox}
}
%====================== BỔ SUNG CHO BÀI HỌC, CHUYÊN ĐỀ ======================
\usepackage[shortlabels]{enumitem}
\usepackage{tabularx}
%--- Danh sách gọn: \begin{enumerate}[gon] ... \end{enumerate}
\SetEnumitemKey{gon}{leftmargin=*,itemsep=1pt,parsep=0pt,topsep=3pt}
%--- Khung lý thuyết: \begin{khung}{} ... \end{khung}; tiêu đề chỉ dùng khi thật ngắn ("Mẹo", "Lỗi thường gặp")
\newtcolorbox{khung}[1]{breakable,enhanced,colframe=gold1,colback=gold5!45,
	boxrule=0.4mm,fonttitle=\bfseries\sffamily,coltitle=white,colbacktitle=gold1,
	attach boxed title to top left={xshift=6mm,yshift=-3mm},
	boxed title style={arc=1mm,boxrule=0.3mm},title={#1}}
%--- Tiêu đề mục nhỏ trong một dạng: \muc{Dấu hiệu}, \muc{Phương pháp}, \muc{Bài tự luyện}
\makeatletter
\newcommand{\muc}[1]{\par\medskip\noindent\indam{#1}\par\nobreak\smallskip\@afterindentfalse\@afterheading}
\makeatother
%--- Ô công thức cần nhớ (căn giữa, viền vàng, vừa khít nội dung)
%     \begin{center}\congthuc{$Q(n)=(n-1)!$}\end{center}
\newtcbox{\congthuc}{on line,colframe=gold1,colback=gold5!45,boxrule=0.5pt,arc=2pt,
	left=8pt,right=8pt,top=4pt,bottom=4pt}
%--- Công tắc lời giải bài tự luyện (ví dụ mẫu luôn dùng \loigiai)
%     \loigiaiTL[{đáp số}]{lời giải đầy đủ}
%     \TLfull -> in lời giải đầy đủ;  \TLdapso -> chỉ in một dòng "Đáp số";  \TLan -> ẩn hoàn toàn
%     (\CheDo là dethi thì mọi lời giải và đáp số đều ẩn)
\makeatletter
\newif\ifTLhien   \TLhientrue
\newif\ifTLdapso  \TLdapsofalse
\newcommand{\TLfull}{\TLhientrue\TLdapsofalse}
\newcommand{\TLdapso}{\TLhientrue\TLdapsotrue}
\newcommand{\TLan}{\TLhienfalse}
\newcommand{\loigiaiTL}[2][]{%
	\ifTLhien
		\ifTLdapso
			\def\@tlarg{#1}%
			\ifx\@tlarg\@empty\else
				\ifdefstring{\kindOPTN}{dethi}{}{\par\noindent{\bfseries\color{gold1}Đáp số:} #1\par}%
			\fi
		\else
			\loigiai{#2}%
		\fi
	\fi}
\makeatother
\TLdapso% mặc định: chỉ in đáp số bài tự luyện
%====================== BẢNG ĐÚNG/SAI (\choiceTFt) ======================
%  Hàng tiêu đề nền nâu đồng chữ trắng (Phát biểu | Đ | S), các hàng tô xen kẽ kem / trắng,
%  đường kẻ vàng kim, nhãn a) b) c) d) màu nâu đồng. Bản có đáp án: cột được chọn có vòng tròn
%  nâu đồng với dấu tích trắng (câu đúng tích ở cột Đúng, câu sai tích ở cột Sai).
%  Tiêu đề ghi đủ "Đúng | Sai": \OPTN{viettat=0} sau \input (cột Đúng/Sai rộng ra, phát biểu hẹp lại).
%  Dùng bảng gốc của ex_test: \def\BangDSgoc{1} trước \input{preamble-eduzy-v2}.
\providecommand{\BangDSgoc}{0}
\ifnum\BangDSgoc=0
\usepackage{colortbl,pifont}
\colorlet{ezDSke}{gold2!60!gold1}% đường kẻ bảng (nét 0,4pt như ex_test để không đổi chiều cao bảng)
\colorlet{ezDSle}{gold5!60}% hàng a), c)
\colorlet{ezDSchan}{white}% hàng b), d)
%--- dấu chọn ở cột Đúng/Sai
%    (bề rộng 0, căn giữa ô: không làm đổi bề rộng cột của ex_test, không báo tràn lề)
\newcommand{\EZdauchon}{\makebox[0pt]{\tikz[baseline=-1ex]\node[circle,fill=gold1,inner sep=0pt,
	minimum size=1.2em,text=white,font=\small]{\ding{51}};}}
\def\tickT{\EZdauchon}\def\tickF{\EZdauchon}
\AtBeginDocument{\sbox{\ccTFTab}{\EZdauchon}}
%--- nhãn a) b) c) d) màu nâu đồng (ở bản lời giải, nhãn câu đúng vẫn đóng khung như ex_test)
\renewcommand{\FalseTF}{\stepcounter{dapan}\textbf{\color{gold1}\DapAnTF\sepTF}\ignorespaces}
%--- trước khi dựng bảng: kẻ màu vàng kim (bề rộng cột giữ nguyên như ex_test); sau bảng trả lại màu đen
\newcount\EZhangTF
\newcommand{\EZbatdauTF}{%
	\global\EZhangTF=0\relax
	\arrayrulecolor{ezDSke}}
\newcommand{\EZketthucTF}{\arrayrulecolor{black}}
%--- hàng tiêu đề nền nâu đồng, chữ trắng
\makeatletter
\newcommand{\EZo@TF}[1]{\centering\bfseries\color{white}#1}
\newcommand{\EZtieudeTF}{%
	\ifcase\cotDS
		\ifnum\cotSTT=1\relax
			\def\rowHeader{\hline\rowcolor{gold1}\EZo@TF{TT}&\EZo@TF{\phatbieu}\tabularnewline}%
		\else
			\def\rowHeader{\hline\rowcolor{gold1}\EZo@TF{\phatbieu}\tabularnewline}%
		\fi
	\or
		\ifnum\cotSTT=1\relax
			\def\rowHeader{\hline\rowcolor{gold1}\EZo@TF{TT}&\EZo@TF{\phatbieu}&\EZo@TF{\Dung}\tabularnewline}%
		\else
			\def\rowHeader{\hline\rowcolor{gold1}\EZo@TF{\phatbieu}&\EZo@TF{\Dung}\tabularnewline}%
		\fi
	\or
		\ifnum\cotSTT=1\relax
			\def\rowHeader{\hline\rowcolor{gold1}\EZo@TF{TT}&\EZo@TF{\phatbieu}&\EZo@TF{\Dung}&\EZo@TF{\Sai}\tabularnewline}%
		\else
			\def\rowHeader{\hline\rowcolor{gold1}\EZo@TF{\phatbieu}&\EZo@TF{\Dung}&\EZo@TF{\Sai}\tabularnewline}%
		\fi
	\fi}
%--- tô xen kẽ từng hàng (gọi ngay sau \hline đầu mỗi hàng)
\newcommand{\EZmauhangTF}{%
	\noalign{\global\advance\EZhangTF\@ne}%
	\ifodd\EZhangTF\rowcolor{ezDSle}\else\rowcolor{ezDSchan}\fi}
%--- gắn vào \choiceTFt / \writeChoiceTFt của ex_test
\xpretocmd{\choiceTFt}{\EZbatdauTF}{}{\typeout{[preamble-eduzy-v2] Không gắn được bảng Đúng/Sai (đầu bảng)}}
\xpatchcmd{\choiceTFt}{\ifnum\dongPB=0\def\rowHeader{}\fi}{\EZtieudeTF\ifnum\dongPB=0\def\rowHeader{}\fi}
	{}{\typeout{[preamble-eduzy-v2] Không gắn được bảng Đúng/Sai (hàng tiêu đề)}}
\xpatchcmd{\writeChoiceTFt}{\hline\xdef\dotEX{\sdotEX}}{\hline\EZmauhangTF\xdef\dotEX{\sdotEX}}
	{}{\typeout{[preamble-eduzy-v2] Không gắn được bảng Đúng/Sai (tô hàng)}}
\xpatchcmd{\writeChoiceTFt}{{\end{longtable}}}{{\end{longtable}\EZketthucTF}}
	{}{\typeout{[preamble-eduzy-v2] Không gắn được bảng Đúng/Sai (cuối longtable)}}
\xpatchcmd{\writeChoiceTFt}{{\end{xtabular}\end{center}\medskip\noindent}}{{\end{xtabular}\end{center}\EZketthucTF\medskip\noindent}}
	{}{\typeout{[preamble-eduzy-v2] Không gắn được bảng Đúng/Sai (cuối xtabular)}}
\fi
%=====================================================================
\renewcommand{\parallel}{\mathrel{/\!\!/}}
.
\providecommand{\BangDSgoc}{0}
\ifnum\BangDSgoc=0
\usepackage{colortbl,pifont}
\colorlet{ezDSke}{gold2!60!gold1}% đường kẻ bảng (nét 0,4pt như ex_test để không đổi chiều cao bảng)
\colorlet{ezDSle}{gold5!60}% hàng a), c)
\colorlet{ezDSchan}{white}% hàng b), d)
%--- dấu chọn ở cột Đúng/Sai
%    (bề rộng 0, căn giữa ô: không làm đổi bề rộng cột của ex_test, không báo tràn lề)
\newcommand{\EZdauchon}{\makebox[0pt]{\tikz[baseline=-1ex]\node[circle,fill=gold1,inner sep=0pt,
	minimum size=1.2em,text=white,font=\small]{\ding{51}};}}
\def\tickT{\EZdauchon}\def\tickF{\EZdauchon}
\AtBeginDocument{\sbox{\ccTFTab}{\EZdauchon}}
%--- nhãn a) b) c) d) màu nâu đồng (ở bản lời giải, nhãn câu đúng vẫn đóng khung như ex_test)
\renewcommand{\FalseTF}{\stepcounter{dapan}\textbf{\color{gold1}\DapAnTF\sepTF}\ignorespaces}
%--- trước khi dựng bảng: kẻ màu vàng kim (bề rộng cột giữ nguyên như ex_test); sau bảng trả lại màu đen
\newcount\EZhangTF
\newcommand{\EZbatdauTF}{%
	\global\EZhangTF=0\relax
	\arrayrulecolor{ezDSke}}
\newcommand{\EZketthucTF}{\arrayrulecolor{black}}
%--- hàng tiêu đề nền nâu đồng, chữ trắng
\makeatletter
\newcommand{\EZo@TF}[1]{\centering\bfseries\color{white}#1}
\newcommand{\EZtieudeTF}{%
	\ifcase\cotDS
		\ifnum\cotSTT=1\relax
			\def\rowHeader{\hline\rowcolor{gold1}\EZo@TF{TT}&\EZo@TF{\phatbieu}\tabularnewline}%
		\else
			\def\rowHeader{\hline\rowcolor{gold1}\EZo@TF{\phatbieu}\tabularnewline}%
		\fi
	\or
		\ifnum\cotSTT=1\relax
			\def\rowHeader{\hline\rowcolor{gold1}\EZo@TF{TT}&\EZo@TF{\phatbieu}&\EZo@TF{\Dung}\tabularnewline}%
		\else
			\def\rowHeader{\hline\rowcolor{gold1}\EZo@TF{\phatbieu}&\EZo@TF{\Dung}\tabularnewline}%
		\fi
	\or
		\ifnum\cotSTT=1\relax
			\def\rowHeader{\hline\rowcolor{gold1}\EZo@TF{TT}&\EZo@TF{\phatbieu}&\EZo@TF{\Dung}&\EZo@TF{\Sai}\tabularnewline}%
		\else
			\def\rowHeader{\hline\rowcolor{gold1}\EZo@TF{\phatbieu}&\EZo@TF{\Dung}&\EZo@TF{\Sai}\tabularnewline}%
		\fi
	\fi}
%--- tô xen kẽ từng hàng (gọi ngay sau \hline đầu mỗi hàng)
\newcommand{\EZmauhangTF}{%
	\noalign{\global\advance\EZhangTF\@ne}%
	\ifodd\EZhangTF\rowcolor{ezDSle}\else\rowcolor{ezDSchan}\fi}
%--- gắn vào \choiceTFt / \writeChoiceTFt của ex_test
\xpretocmd{\choiceTFt}{\EZbatdauTF}{}{\typeout{[preamble-eduzy-v2] Không gắn được bảng Đúng/Sai (đầu bảng)}}
\xpatchcmd{\choiceTFt}{\ifnum\dongPB=0\def\rowHeader{}\fi}{\EZtieudeTF\ifnum\dongPB=0\def\rowHeader{}\fi}
	{}{\typeout{[preamble-eduzy-v2] Không gắn được bảng Đúng/Sai (hàng tiêu đề)}}
\xpatchcmd{\writeChoiceTFt}{\hline\xdef\dotEX{\sdotEX}}{\hline\EZmauhangTF\xdef\dotEX{\sdotEX}}
	{}{\typeout{[preamble-eduzy-v2] Không gắn được bảng Đúng/Sai (tô hàng)}}
\xpatchcmd{\writeChoiceTFt}{{\end{longtable}}}{{\end{longtable}\EZketthucTF}}
	{}{\typeout{[preamble-eduzy-v2] Không gắn được bảng Đúng/Sai (cuối longtable)}}
\xpatchcmd{\writeChoiceTFt}{{\end{xtabular}\end{center}\medskip\noindent}}{{\end{xtabular}\end{center}\EZketthucTF\medskip\noindent}}
	{}{\typeout{[preamble-eduzy-v2] Không gắn được bảng Đúng/Sai (cuối xtabular)}}
\fi
%=====================================================================
\renewcommand{\parallel}{\mathrel{/\!\!/}}

%     %=====================================================================
% preamble-vatli-hinh.tex -- BỘ MACRO HÌNH VẬT LÍ (dùng chung cho mọi đề Vật lí)
% Gọi ngay sau preamble chính:
%     %============================================================
%  PREAMBLE EDUZY v2 - dùng chung cho đề thi, bài học, chuyên đề (gói ex_test)
%  Khác bản preamble-eduzy.tex ở đầu trang, chân trang: dải nền nâu đồng tràn ngang,
%  logo + Eduzy.com.vn, ô chữ vàng kem cạnh xiên, chân trang "Trang (số)";
%  logo Eduzy chìm giữa mọi trang (tắt: \logochimtat, bật lại: \logochimbat).
%  Cách dùng ở file chính (xem mau-chuan.tex, mau-chuyen-de.tex):
%      \def\CheDo{loigiai}          % loigiai | dethi | solcolor | color | book
%      \def\tenchuong{<dòng chữ trên đầu trang>}
%      \def\thuonghieu{<chữ ở chân trang lẻ>}   % không khai báo thì dùng mặc định
%      \def\tacgia{<chữ ở chân trang chẵn>}      % không khai báo thì dùng mặc định
%      \def\khauhieu{<chữ cạnh logo ở đầu trang>}  % tùy chọn, mặc định Eduzy.com.vn
%      \input{preamble-eduzy-v2}
%      \TLdapso                      % tài liệu: \TLfull | \TLdapso | \TLan
%      \begin{document} ... \end{document}
%  Logo Eduzy vẽ bằng TikZ ngay trong file (\eduzylogo), không cần file ảnh.
%  Biên dịch: pdfLaTeX (máy có gói vietnam) hoặc XeLaTeX.
%============================================================
\documentclass[12pt,a4paper,twoside]{article}%<-- twoside: bat dau trang/chan trang chan-le
\usepackage{iftex}
\usepackage{amsmath,amssymb,tikz,mathrsfs,tkz-tab}
\usepackage{graphicx}
\usepackage{pgfplots}% đồ thị hàm số
\pgfplotsset{compat=1.17}
%\usepackage{mathptmx}
\usepackage{mathpazo}
\ifXeTeX% XeLaTeX: font chữ có đủ dấu tiếng Việt, giữ nguyên font toán mathpazo
	\usepackage[no-math]{fontspec}
	\setmainfont{TeX Gyre Pagella}
	\newfontfamily\EZqag{TeX Gyre Adventor}[NFSSFamily=qag,Ligatures=TeX]
\fi
\usepackage[many]{tcolorbox}
\usetikzlibrary{math,through,calc,intersections,angles,quotes,shapes,shapes.geometric,arrows,patterns,matrix,chains,arrows.meta,decorations.shapes,decorations.fractals,decorations.markings,shadows,decorations.pathreplacing}
\usepackage[top=1.6cm, bottom=1.8cm, left=1.8cm, right=1.8cm]{geometry}% v2: top 1.6cm (bản cũ 1.5cm) để chữ không sát dải đầu trang
\providecommand{\CheDo}{loigiai}% loigiai | dethi | solcolor | color | book
\expandafter\usepackage\expandafter[\CheDo]{ex_test}
% \tenchuong: khai báo ở file chính (trước \input) hoặc \renewcommand sau \input
%============================================================
%  [GỘP TỪ eduzy-de.tex] EDUZY - MÀU NÂU ĐỒNG / VÀNG KIM
%  + ĐẦU TRANG, CHÂN TRANG KIỂU DẢI NGANG (v2), ĐỐI XỨNG TRANG CHẴN / TRANG LẺ.
%  (không cần file eduzy-de.tex riêng nữa)
%  Khối này phải đặt SAU \usepackage[...]{ex_test} và SAU \def\tenchuong{...}
%============================================================
\usepackage{lastpage}
\usepackage{needspace}
%--- DÀN TRANG ĐỀ THI: tách câu (THÊM 05/11/2026; chỉ thêm lệnh mới, KHÔNG đổi hành vi mặc định) ---
%  Hai công tắc, mặc định tắt (= preamble chạy y như cũ). Bật bằng cách đặt TRƯỚC \input{preamble-eduzy-v2}
%  trong file đề (tools/dan_trang/dan_trang.py tự ghi khi đặt \NeedspaceCau):
%    \def\TachCauDS{1}   câu Đúng/Sai được tách: câu dẫn (kèm hình) ở cuối trang, bảng Phát biểu sang trang sau
%                        (câu dẫn và hình không tách nhau, bảng không tách giữa hai trang)
%    \def\TachCauLG{1}   bản lời giải (\CheDo khác dethi): phần đề của câu cao hơn 1/3 trang được phép tách
%  Lệnh dùng trong file đề (giá trị do tools/dan_trang/dan_trang.py đo ở bản dethi):
%    \NeedspaceCau[<cao câu dẫn + hình>]{<cao cả khối câu>}   ngay trên \begin{ex} (câu đầu phần: trên tiêu đề phần)
%        công tắc tắt: đúng bằng \Needspace{<cao cả khối câu>} như cũ;
%        \TachCauDS=1 và có [..] (câu Đúng/Sai được chọn tách): chỉ giữ câu dẫn + hình;
%        \TachCauLG=1 ở bản lời giải: đòi tối đa 1/3 \textheight.
%    \NeedspaceBang{<cao bảng Phát biểu>}%   ngay trước \choiceTFt: \TachCauDS=1 thì giữ trọn bảng trên một trang;
%        công tắc tắt: không làm gì.
\providecommand{\TachCauDS}{0}
\providecommand{\TachCauLG}{0}
\newlength{\EZnsCau}
\newcommand{\NeedspaceCau}[2][]{%
	\setlength{\EZnsCau}{#2}%
	\ifnum\TachCauDS=1 \if\relax\detokenize{#1}\relax\else\setlength{\EZnsCau}{#1}\fi\fi
	\ifnum\TachCauLG=1 \ifdefstring{\kindOPTN}{dethi}{}{\ifdim\EZnsCau>0.3333\textheight\setlength{\EZnsCau}{0.3333\textheight}\fi}\fi
	\Needspace{\EZnsCau}}
\newcommand{\NeedspaceBang}[1]{\ifnum\TachCauDS=1 \Needspace{#1}\fi}

%--- Thang màu: đậm -> nhạt
\definecolor{gold1}{RGB}{126,96,0}    % nâu đồng đậm  #7E6000
\definecolor{gold2}{RGB}{192,144,3}   % vàng kim      #C09003
\definecolor{gold3}{RGB}{255,215,99}  % vàng nhạt     #FFD763
\definecolor{gold4}{RGB}{255,231,153} % vàng rất nhạt #FFE799
\definecolor{gold5}{RGB}{255,243,205} % vàng kem      #FFF3CD
%--- Ghi đè tông xanh/đỏ cũ
\definecolor{myblue}{RGB}{126,96,0}
\colorlet{blue}{gold1}
\colorlet{cyan}{gold2}
\colorlet{teal}{gold2}

%--- Tên chương mặc định (file đề nên tự khai báo lại)
\providecommand{\tenchuong}{}
%--- Chân trang
\providecommand{\thuonghieu}{TOÁN 12 -- Eduzy.com.vn}
\providecommand{\tacgia}{Tô Gia Bảo -- 0353703783}

%--- Chữ cạnh logo ở đầu trang
\providecommand{\khauhieu}{Eduzy.com.vn}

%====================== THÔNG SỐ ĐẦU TRANG, CHÂN TRANG (v2) ======================
%  Dải nền đậm tràn ngang trang; ô chữ vàng kem có cạnh xiên, kèm 2 vạch xiên chuyển màu.
%  Đầu trang lẻ : logo + \khauhieu (trái)   | ô chữ \tenchuong (phải)
%  Chân trang lẻ: ô chữ \thuonghieu (trái)  | Trang (số) (phải)
%  Trang chẵn soi gương trang lẻ; chân trang chẵn ghi \tacgia.
%--- Màu: chỉ dùng thang vàng kim Eduzy
\colorlet{ezNen}{gold1}% dải nền đậm
\colorlet{ezVach1}{gold3}% vạch xiên sát dải nền (đậm hơn)
\colorlet{ezVach2}{gold4}% vạch xiên sát ô chữ (nhạt hơn)
\colorlet{ezO}{gold5}% ô chữ vàng kem
%--- Kích thước
\newlength{\EZdai}\setlength{\EZdai}{1.1cm}% chiều cao dải nền
\newlength{\EZo}\setlength{\EZo}{0.68cm}% chiều cao ô chữ
\newlength{\EZxien}\setlength{\EZxien}{0.55cm}% độ lệch ngang của cạnh xiên
\newlength{\EZvach}\setlength{\EZvach}{0.75cm}% bề rộng mỗi vạch xiên
\newlength{\EZkhe}\setlength{\EZkhe}{0.12cm}% khe trắng giữa vạch và ô chữ
\newlength{\EZdem}\setlength{\EZdem}{0.4cm}% khoảng đệm từ chữ tới cạnh xiên
\newlength{\EZle}\setlength{\EZle}{1.8cm}% lề trái/phải của chữ (bằng lề thân bài)
\newlength{\EZrongmin}\setlength{\EZrongmin}{4.5cm}% ô chữ rộng tối thiểu
\newlength{\EZrong}% bề rộng chữ trong ô (tính tự động)
%--- Font
\newcommand{\EZfontO}{\fontsize{10.5pt}{12pt}\selectfont\bfseries}% chữ trong ô
\newcommand{\EZfontKH}{\fontfamily{qag}\fontsize{11pt}{12pt}\selectfont\bfseries}% \khauhieu cạnh logo
\newcommand{\EZfontTrang}{\fontsize{11.5pt}{12pt}\selectfont\bfseries}% chữ "Trang" và số trang
%--- Đo bề rộng chữ trong ô
\newcommand{\EZdorong}[1]{%
	\settowidth{\EZrong}{\EZfontO #1}%
	\ifdim\EZrong<\EZrongmin\setlength{\EZrong}{\EZrongmin}\fi}
%--- Logo + chữ trắng \khauhieu (dựng sẵn một lần khi vào \begin{document})
\newsavebox{\EZlogobox}
\AtBeginDocument{\sbox{\EZlogobox}{%
		\raisebox{-0.5\height}{\eduzylogo[1cm]}\hspace{0.15cm}%
		\raisebox{-0.5\height}{\color{white}\EZfontKH\khauhieu}}}
%--- Hai vạch xiên + khe trắng. #1: điểm gốc ở mép dưới, #2: hướng (1 sang phải, -1 sang trái),
%    #3: hướng nghiêng (1 kiểu "/", -1 kiểu "\")
\newcommand{\EZvachxien}[3]{%
	\fill[white] (#1) --++({(#3)*\EZxien},\EZo) --++({(#2)*\EZkhe},0) --++({-(#3)*\EZxien},-\EZo) --cycle;
	\fill[ezVach2] ([xshift=(#2)*\EZkhe]#1) --++({(#3)*\EZxien},\EZo) --++({(#2)*\EZvach},0) --++({-(#3)*\EZxien},-\EZo) --cycle;
	\fill[ezVach1] ([xshift=(#2)*(\EZkhe+\EZvach)]#1) --++({(#3)*\EZxien},\EZo) --++({(#2)*\EZvach},0) --++({-(#3)*\EZxien},-\EZo) --cycle;
	\draw[white,line width=0.6pt] ([xshift=(#2)*(\EZkhe+\EZvach)]#1) --++({(#3)*\EZxien},\EZo);
}
%--- Chữ trong ô (căn giữa theo chiều cao ô, không lệch vì dấu tiếng Việt)
\tikzset{ezchu/.style={inner sep=0pt,text=ezNen,text height=0.72em,text depth=0.22em}}

%====================== LOGO CHÌM GIỮA TRANG ======================
%  Vẽ cùng đầu trang nên nằm dưới chữ của thân bài, có trên mọi trang.
%  Tắt: \logochimtat (từ trang đó trở đi), bật lại: \logochimbat.
%  Đổi cỡ, độ đậm: khai báo trước \input, ví dụ \def\EZchimcao{9cm} \def\EZchimdo{0.1}
\providecommand{\EZchimcao}{10cm}% chiều cao logo chìm
\providecommand{\EZchimdo}{0.12}% độ đậm (0 = trong suốt, 1 = rõ hẳn)
\newif\ifEZchim \EZchimtrue
\newcommand{\logochimtat}{\EZchimfalse}
\newcommand{\logochimbat}{\EZchimtrue}
%  Dựng sẵn logo mờ một lần khi vào \begin{document}; transparency group làm mờ cả logo
%  một lượt nên không lộ vệt ở chỗ các dải màu chồng nhau.
\newsavebox{\EZchimtho}
\newsavebox{\EZchimbox}
\AtBeginDocument{%
	\sbox{\EZchimtho}{\eduzylogo[\EZchimcao]}%
	\sbox{\EZchimbox}{\begin{tikzpicture}
			\begin{scope}[transparency group,opacity=\EZchimdo]
				\node[inner sep=0pt] {\usebox{\EZchimtho}};
			\end{scope}
	\end{tikzpicture}}}
\newcommand{\EZvechim}{%
	\ifEZchim
		\node[inner sep=0pt] at (current page.center) {\usebox{\EZchimbox}};
	\fi}

%====================== ĐẦU TRANG - TRANG LẺ ======================
%  (logo bên trái, ô chữ \tenchuong bên phải)
\newcommand{\myfancyheadOdd}{%
	\EZdorong{\tenchuong}%
	\begin{tikzpicture}[remember picture,overlay]
		\coordinate (T) at (current page.north west);
		\coordinate (P) at (current page.north east);
		%--- logo chìm giữa trang
		\EZvechim
		%--- dải nền
		\fill[ezNen] (T) rectangle ([yshift=-\EZdai]P);
		%--- ô chữ: X là góc dưới bên trái, cạnh trái xiên kiểu "/"
		\coordinate (X) at ([xshift=-\EZle-\EZrong-\EZdem-\EZxien,yshift=-\EZdai]P);
		\fill[ezO] (X) --++(\EZxien,\EZo) -- ([yshift=-\EZdai+\EZo]P) -- ([yshift=-\EZdai]P) --cycle;
		\EZvachxien{X}{-1}{1}
		\node[ezchu,anchor=east] at ([xshift=-\EZle,yshift=-\EZdai+0.5*\EZo]P) {\EZfontO\tenchuong};
		%--- logo + tên
		\node[inner sep=0pt,anchor=west] at ([xshift=\EZle,yshift=-0.5*\EZdai]T) {\usebox{\EZlogobox}};
	\end{tikzpicture}%
}

%====================== ĐẦU TRANG - TRANG CHẴN =====================
%  (soi gương: ô chữ \tenchuong bên trái, logo bên phải)
\newcommand{\myfancyheadEven}{%
	\EZdorong{\tenchuong}%
	\begin{tikzpicture}[remember picture,overlay]
		\coordinate (T) at (current page.north west);
		\coordinate (P) at (current page.north east);
		%--- logo chìm giữa trang
		\EZvechim
		%--- dải nền
		\fill[ezNen] (T) rectangle ([yshift=-\EZdai]P);
		%--- ô chữ: X là góc dưới bên phải, cạnh phải xiên kiểu "\"
		\coordinate (X) at ([xshift=\EZle+\EZrong+\EZdem+\EZxien,yshift=-\EZdai]T);
		\fill[ezO] (X) --++(-\EZxien,\EZo) -- ([yshift=-\EZdai+\EZo]T) -- ([yshift=-\EZdai]T) --cycle;
		\EZvachxien{X}{1}{-1}
		\node[ezchu,anchor=west] at ([xshift=\EZle,yshift=-\EZdai+0.5*\EZo]T) {\EZfontO\tenchuong};
		%--- logo + tên
		\node[inner sep=0pt,anchor=east] at ([xshift=-\EZle,yshift=-0.5*\EZdai]P) {\usebox{\EZlogobox}};
	\end{tikzpicture}%
}

%====================== CHÂN TRANG - TRANG LẺ ======================
%  (ô chữ \thuonghieu bên trái, "Trang (số)" bên phải)
\newcommand{\myfancyfootOdd}{%
	\EZdorong{\thuonghieu}%
	\begin{tikzpicture}[remember picture,overlay]
		\coordinate (T) at (current page.south west);
		\coordinate (P) at (current page.south east);
		%--- dải nền
		\fill[ezNen] (T) rectangle ([yshift=\EZdai]P);
		%--- ô chữ sát mép trên dải: X là góc dưới bên phải, cạnh phải xiên kiểu "/"
		\coordinate (X) at ([xshift=\EZle+\EZrong+\EZdem,yshift=\EZdai-\EZo]T);
		\fill[ezO] ([yshift=\EZdai-\EZo]T) -- (X) --++(\EZxien,\EZo) -- ([yshift=\EZdai]T) --cycle;
		\EZvachxien{X}{1}{1}
		\node[ezchu,anchor=west] at ([xshift=\EZle,yshift=\EZdai-0.5*\EZo]T) {\EZfontO\thuonghieu};
		%--- số trang
		\node[circle,draw=white,line width=0.9pt,text=white,inner sep=0pt,minimum size=0.62cm,anchor=east]
		(so) at ([xshift=-\EZle,yshift=0.5*\EZdai]P) {\EZfontTrang\thepage};
		\node[inner sep=0pt,text=white,anchor=base east] (tr) at ([xshift=-0.15cm]so.base west) {\EZfontTrang Trang};
	\end{tikzpicture}%
}

%====================== CHÂN TRANG - TRANG CHẴN ====================
%  (soi gương: "Trang (số)" bên trái, ô chữ \tacgia bên phải)
\newcommand{\myfancyfootEven}{%
	\EZdorong{\tacgia}%
	\begin{tikzpicture}[remember picture,overlay]
		\coordinate (T) at (current page.south west);
		\coordinate (P) at (current page.south east);
		%--- dải nền
		\fill[ezNen] (T) rectangle ([yshift=\EZdai]P);
		%--- ô chữ sát mép trên dải: X là góc dưới bên trái, cạnh trái xiên kiểu "\"
		\coordinate (X) at ([xshift=-\EZle-\EZrong-\EZdem,yshift=\EZdai-\EZo]P);
		\fill[ezO] ([yshift=\EZdai-\EZo]P) -- (X) --++(-\EZxien,\EZo) -- ([yshift=\EZdai]P) --cycle;
		\EZvachxien{X}{-1}{-1}
		\node[ezchu,anchor=east] at ([xshift=-\EZle,yshift=\EZdai-0.5*\EZo]P) {\EZfontO\tacgia};
		%--- số trang
		%    (vòng tròn ẩn so0 nằm giữa dải để lấy đường chân chữ, giống trang lẻ)
		\node[circle,inner sep=0pt,minimum size=0.62cm] (so0) at ([yshift=0.5*\EZdai]T) {\EZfontTrang\phantom{\thepage}};
		\node[inner sep=0pt,text=white,anchor=base west] (tr) at ($(so0.base -| T)+(\EZle,0)$) {\EZfontTrang Trang};
		\node[circle,draw=white,line width=0.9pt,text=white,inner sep=0pt,minimum size=0.62cm,anchor=base west]
		(so) at ([xshift=0.15cm]tr.base east) {\EZfontTrang\thepage};
	\end{tikzpicture}%
}

%--- Giữ tương thích với các file cũ đang gọi \myfancyhead, \myfancyfoot
\newcommand{\myfancyhead}{\myfancyheadOdd}
\newcommand{\myfancyfoot}{\myfancyfootOdd}

%======================= GẮN VÀO TRANG =============================
\pagestyle{fancy}
\fancyhf{}
\fancyhead[LO]{\myfancyheadOdd}
\fancyhead[RE]{\myfancyheadEven}
\fancyfoot[LO]{\myfancyfootOdd}
\fancyfoot[RE]{\myfancyfootEven}
\renewcommand{\headrulewidth}{0pt}
\renewcommand{\footrulewidth}{0pt}
\fancypagestyle{plain}{%
	\fancyhf{}
	\fancyhead[LO]{\myfancyheadOdd}
	\fancyhead[RE]{\myfancyheadEven}
	\fancyfoot[LO]{\myfancyfootOdd}
	\fancyfoot[RE]{\myfancyfootEven}
	\renewcommand{\headrulewidth}{0pt}
	\renewcommand{\footrulewidth}{0pt}
}
%=================== HẾT PHẦN GỘP TỪ eduzy-de.tex ===================
\newcommand{\indam}[1]{
	{\fontsize{12pt}{9pt}\sffamily \selectfont\bfseries\color{gold1} #1}
}
\newcommand{\ind}[1]{
	{\fontsize{14pt}{9pt}\fontfamily{qag} \selectfont\bfseries\color{gold1} #1}
}
\renewcommand{\FalseEX}{\stepcounter{dapan}\textbf{\selectfont\color{gold1}\Alph{dapan}.}\ignorespaces}
\renewtheorem{ex}{\color{gold1}Câu}
\renewcommand{\baselinestretch}{1.1}
\newcommand{\hoac}[1]{
	\left[\begin{aligned}#1\end{aligned}\right.}
\newcommand{\heva}[1]{
	\left\{\begin{aligned}#1\end{aligned}\right.}
\usepackage{esvect}
\def\vec{\vv}
\def\vt{\vv}
\def\overrightarrow{\vv}
%====================== LOGO EDUZY VẼ BẰNG TIKZ ======================
%  Đường viền dò từ logo gốc; dải màu vàng kim đo trực tiếp trên ảnh gốc
%  (nghiêng 21 độ, tô bằng 80 dải hẹp nên hiển thị giống nhau trên mọi trình xem PDF).
%  Dùng: \eduzylogo          -> cao 2.4cm (mặc định)
%        \eduzylogo[1.5cm]   -> cao 1.5cm
\definecolor{ezText}{RGB}{245,200,46}
\newcommand{\eduzylogo}[1][2.4cm]{%
	\resizebox{!}{#1}{%
		\begin{tikzpicture}[x=0.01cm,y=0.01cm]
			\path[use as bounding box] (0,0) rectangle (224,-224);
			% mũ tốt nghiệp + vương miện
			\begin{scope}
				\clip
				(119.4,-176.1)..controls(118.9,-176.1)and(116.8,-176)..(114.6,-176)
				..controls(101.4,-176)and(85.8,-175.9)..(83.7,-175.7)
				..controls(81.8,-175.6)and(81,-175.6)..(80.2,-175.7)
				..controls(78.8,-175.9)and(77.7,-175.9)..(76.2,-175.7)
				..controls(74.9,-175.5)and(70.8,-175.5)..(64.3,-175.7)
				..controls(59.1,-175.9)and(58.3,-175.7)..(58.6,-173.7)
				..controls(58.6,-172.8)and(59.6,-169.6)..(60.3,-168)
				..controls(60.8,-166.9)and(61.4,-165.8)..(62.3,-165)
				..controls(63.1,-164.2)and(64.3,-163.4)..(65,-163.2)
				..controls(65.5,-163.1)and(72.3,-163)..(76.6,-163.1)
				..controls(77.9,-163.1)and(80.1,-163.1)..(81.4,-163)
				..controls(82.8,-162.9)and(86.3,-162.9)..(89.2,-162.9)
				..controls(93,-162.9)and(94.9,-162.8)..(96,-162.7)
				..controls(97.2,-162.5)and(101,-162.5)..(114,-162.4)
				..controls(123,-162.3)and(134.9,-162.2)..(140.2,-162.1)
				..controls(150.1,-162)and(152.2,-162.1)..(153.6,-162.4)
				..controls(154.4,-162.6)and(157.1,-163.8)..(158.1,-164.4)
				..controls(158.8,-164.8)and(159.7,-165.7)..(159.9,-166.3)
				..controls(160,-166.5)and(160.2,-167.1)..(160.2,-167.6)
				..controls(160.3,-168)and(160.6,-169.4)..(160.9,-170.6)
				..controls(161.2,-171.7)and(161.5,-173.1)..(161.6,-173.6)
				..controls(161.8,-174.4)and(161.8,-174.5)..(161.5,-174.9)
				..controls(161.2,-175.6)and(160.8,-175.9)..(159.9,-176)
				..controls(159.2,-176.1)and(154.1,-176.2)..(151.5,-176.2)
				..controls(148.1,-176.1)and(122.6,-176.1)..(121.6,-176.2)
				..controls(120.8,-176.2)and(119.8,-176.2)..(119.4,-176.1)--cycle
				(130.4,-160)..controls(129,-160)and(119.8,-160)..(110.1,-159.9)
				..controls(100.3,-159.8)and(88.5,-159.7)..(83.9,-159.6)
				..controls(77.8,-159.5)and(75.3,-159.5)..(74.5,-159.6)
				..controls(73.9,-159.7)and(72.5,-159.8)..(71.3,-159.8)
				..controls(69.6,-159.8)and(69,-159.8)..(68.5,-159.6)
				..controls(67.6,-159.3)and(66.8,-158.6)..(66.2,-157.6)
				..controls(66,-157.3)and(65.6,-156.5)..(65.3,-156.1)
				..controls(64.6,-155.1)and(64.5,-154.9)..(63.7,-152.9)
				..controls(63.2,-151.6)and(62.9,-150.9)..(61.8,-149.6)
				..controls(61.5,-149.2)and(61.1,-148.5)..(60.9,-147.9)
				..controls(60.7,-147.4)and(60.3,-146.6)..(60.1,-146.1)
				..controls(59.9,-145.5)and(59.5,-144.9)..(59,-144.3)
				..controls(58.5,-143.7)and(58.2,-143.2)..(57.9,-142.5)
				..controls(57.7,-141.9)and(57.3,-141.1)..(57,-140.7)
				..controls(56.6,-140.3)and(56.1,-139.3)..(55.8,-138.5)
				..controls(55.5,-137.8)and(55.1,-136.8)..(54.9,-136.5)
				..controls(54.7,-136.1)and(54.3,-135.2)..(53.9,-134.5)
				..controls(53.6,-133.8)and(53,-132.7)..(52.6,-132)
				..controls(51.6,-130.5)and(51.4,-130.2)..(51.1,-129)
				..controls(50.6,-127.5)and(50.6,-127.4)..(49.7,-126.3)
				..controls(49.2,-125.7)and(48.6,-124.8)..(48.3,-124.2)
				..controls(48,-123.6)and(47.5,-122.7)..(47.3,-122.2)
				..controls(47,-121.8)and(46.5,-120.7)..(46.1,-119.8)
				..controls(45.8,-118.9)and(45.2,-117.8)..(44.9,-117.2)
				..controls(44.2,-116.1)and(43.3,-114.4)..(42.3,-112.2)
				..controls(41.8,-111.3)and(41.3,-110.4)..(41.1,-110.1)
				..controls(40.5,-109.4)and(40.3,-109)..(40,-107.8)
				..controls(39.8,-107)and(39.6,-106.5)..(39.2,-106)
				..controls(38.9,-105.6)and(38.4,-104.7)..(38,-104)
				..controls(37.7,-103.4)and(37.2,-102.5)..(37,-102.1)
				..controls(36.7,-101.7)and(36.3,-100.8)..(36,-100.1)
				..controls(35.6,-98.9)and(35.3,-98.4)..(34.6,-97.6)
				..controls(34.2,-97.1)and(34.2,-96.3)..(34.6,-96)
				..controls(35,-95.7)and(35.5,-95.9)..(36.2,-96.7)
				..controls(36.5,-97)and(37.2,-97.5)..(37.6,-97.7)..controls(38,-98)and(38.6,-98.4)..(39.1,-98.9)
				..controls(39.6,-99.4)and(40.3,-100)..(40.6,-100.3)
				..controls(41.3,-100.8)and(42.4,-101.8)..(43.7,-102.9)
				..controls(44,-103.2)and(44.6,-103.7)..(45,-103.9)
				..controls(45.4,-104.2)and(45.9,-104.6)..(46.1,-104.8)
				..controls(47.1,-105.9)and(48.8,-107.1)..(50.5,-107.9)
				..controls(51.8,-108.6)and(52.3,-109)..(53.4,-110)
				..controls(54.8,-111.2)and(55.6,-111.6)..(57,-112.1)
				..controls(57.5,-112.3)and(58.5,-112.8)..(59.4,-113.3)
				..controls(60.2,-113.7)and(61.3,-114.3)..(61.9,-114.5)
				..controls(62.4,-114.7)and(63.3,-115.1)..(63.8,-115.3)
				..controls(64.4,-115.6)and(65.4,-115.9)..(66.2,-116.1)
				..controls(67,-116.3)and(68.4,-116.6)..(69.2,-116.9)
				..controls(72.3,-117.9)and(72.8,-118)..(78,-118.1)
				..controls(84,-118.3)and(86.1,-118.1)..(88.1,-117.1)
				..controls(89.4,-116.5)and(93,-115)..(93.7,-114.9)
				..controls(94.8,-114.6)and(95.4,-114.2)..(95.9,-113.4)
				..controls(96.2,-113)and(96.7,-112.5)..(97.1,-112.2)
				..controls(98,-111.6)and(101.3,-108.4)..(102.1,-107.2)
				..controls(102.4,-106.8)and(102.9,-106.2)..(103.2,-105.8)
				..controls(103.5,-105.5)and(104,-104.8)..(104.2,-104.2)
				..controls(104.5,-103.6)and(105,-102.8)..(105.3,-102.4)
				..controls(105.6,-102)and(106.1,-101.2)..(106.5,-100.6)
				..controls(107.1,-99.5)and(108.6,-97.5)..(109.5,-96.5)
				..controls(109.9,-96.1)and(110.1,-96)..(110.4,-96)
				..controls(111.1,-96)and(111.6,-96.5)..(111.8,-97.6)
				..controls(112,-98.6)and(112.3,-99.3)..(112.8,-99.9)
				..controls(113.7,-100.9)and(115.2,-103.1)..(115.7,-104.2)
				..controls(116.1,-105)and(116.5,-105.5)..(116.9,-106)
				..controls(117.3,-106.4)and(117.8,-106.9)..(118,-107.3)
				..controls(118.2,-107.6)and(118.6,-108)..(119,-108.3)
				..controls(119.6,-108.8)and(121.1,-110.3)..(121.8,-111.2)
				..controls(122.1,-111.7)and(122.5,-112)..(123.4,-112.4)
				..controls(124.8,-113.2)and(125.8,-113.8)..(127,-114.6)
				..controls(127.5,-114.9)and(128.1,-115.3)..(128.2,-115.4)
				..controls(128.4,-115.5)and(129.3,-115.8)..(130.2,-116)
				..controls(131.1,-116.2)and(132.5,-116.6)..(133.3,-116.9)
				..controls(135.7,-117.9)and(136.1,-117.9)..(142.5,-117.8)
				..controls(145.8,-117.8)and(148.4,-117.7)..(148.8,-117.6)
				..controls(149.2,-117.6)and(150.2,-117.2)..(151,-116.8)
				..controls(152,-116.4)and(152.8,-116)..(153.6,-115.9)
				..controls(154.5,-115.7)and(155.2,-115.4)..(156.3,-114.9)
				..controls(157.1,-114.5)and(158.2,-114)..(158.7,-113.8)
				..controls(159.1,-113.7)and(159.8,-113.3)..(160.2,-113)
				..controls(160.6,-112.7)and(161.3,-112.4)..(161.7,-112.2)
				..controls(162.7,-111.9)and(163.2,-111.7)..(164.1,-111.1)
				..controls(164.5,-110.8)and(165.3,-110.4)..(165.9,-110)
				..controls(166.5,-109.7)and(167.3,-109.2)..(167.6,-108.9)
				..controls(167.9,-108.6)and(168.6,-108.2)..(169,-108)
				..controls(169.4,-107.8)and(170.2,-107.3)..(170.7,-106.9)
				..controls(171.2,-106.5)and(172.1,-105.8)..(172.7,-105.4)
				..controls(173.2,-105)and(174.1,-104.4)..(174.5,-104)
				..controls(174.9,-103.7)and(175.7,-103.1)..(176.2,-102.6)
				..controls(176.7,-102.2)and(177.4,-101.5)..(177.9,-101.2)
				..controls(178.3,-100.8)and(178.8,-100.4)..(179,-100.2)
				..controls(179.2,-100)and(179.6,-99.6)..(180,-99.3)
				..controls(180.4,-99.1)and(180.9,-98.6)..(181.2,-98.4)
				..controls(181.5,-98.1)and(182.1,-97.7)..(182.5,-97.5)
				..controls(183.2,-97.2)and(183.5,-96.9)..(184.3,-95.9)
				..controls(184.6,-95.6)and(184.6,-95.6)..(184.8,-95.7)
				..controls(185,-95.8)and(185.1,-96)..(185.1,-96.2)
				..controls(185.1,-96.5)and(185,-96.7)..(184.6,-97.2)
				..controls(184.1,-97.8)and(184,-97.9)..(183.9,-98.8)
				..controls(183.7,-100.3)and(183.3,-101.3)..(182.9,-101.8)
				..controls(182.6,-102.1)and(182.3,-102.6)..(182.2,-102.9)
				..controls(182,-103.3)and(181.6,-104)..(181.2,-104.6)
				..controls(180.9,-105.2)and(180.5,-106.1)..(180.4,-106.6)
				..controls(179.9,-108)and(179.7,-108.5)..(179,-109.4)
				..controls(178.6,-109.8)and(178.2,-110.5)..(177.9,-111)
				..controls(177.7,-111.5)and(177.3,-112.3)..(177,-112.7)
				..controls(176.6,-113.2)and(176.4,-113.7)..(176.3,-114.2)
				..controls(176.2,-114.7)and(175.9,-115.2)..(175.4,-116)
				..controls(175,-116.6)and(174.4,-117.7)..(174,-118.5)
				..controls(173.7,-119.2)and(173.3,-120.1)..(173,-120.5)
				..controls(172.8,-120.9)and(172.4,-121.8)..(172.1,-122.5)
				..controls(171.8,-123.2)and(171.3,-124.1)..(171,-124.6)
				..controls(170.8,-125.1)and(170.3,-126)..(169.9,-126.6)
				..controls(169.6,-127.3)and(169.2,-128.1)..(169,-128.3)
				..controls(168.8,-128.6)and(168.5,-129.3)..(168.3,-129.8)
				..controls(168.1,-130.3)and(167.8,-131.1)..(167.5,-131.6)
				..controls(166.9,-132.5)and(164.2,-137.9)..(164.2,-138.1)
				..controls(164.2,-138.1)and(163.9,-138.8)..(163.4,-139.7)
				..controls(162.9,-140.5)and(162.4,-141.5)..(162.3,-141.9)
				..controls(161.9,-143)and(161.6,-143.4)..(160.9,-144.1)
				..controls(160.4,-144.6)and(160.3,-144.8)..(160.1,-145.4)
				..controls(160,-145.8)and(159.6,-146.8)..(159.2,-147.8)
				..controls(158.7,-148.8)and(158.1,-150)..(157.9,-150.6)
				..controls(157.6,-151.1)and(157.2,-151.9)..(156.9,-152.4)
				..controls(156.5,-152.8)and(156.1,-153.7)..(155.8,-154.4)
				..controls(155.6,-155)and(155.2,-155.9)..(154.9,-156.3)
				..controls(154.7,-156.7)and(154.4,-157.4)..(154.2,-157.9)
				..controls(154.1,-158.4)and(153.8,-159)..(153.5,-159.3)
				..controls(152.8,-160.1)and(153,-160.1)..(142.4,-160.1)
				..controls(137.2,-160.1)and(131.8,-160.1)..(130.4,-160)--cycle
				(141.5,-103.7)..controls(140.8,-103)and(139.8,-102.1)..(139.2,-101.6)
				..controls(138.5,-101.2)and(137.9,-100.6)..(137.8,-100.4)
				..controls(137,-99.5)and(134,-96.5)..(133.2,-96)
				..controls(132.7,-95.7)and(131.8,-94.9)..(131.1,-94.3)
				..controls(129.4,-92.7)and(127.9,-91.6)..(126.6,-91.1)
				..controls(125.2,-90.5)and(124.1,-90)..(122.9,-89.1)
				..controls(122.1,-88.6)and(121.3,-88.2)..(119.9,-87.7)
				..controls(118.8,-87.3)and(117.5,-86.8)..(116.9,-86.5)
				..controls(114.3,-85.5)and(113.7,-85.4)..(109.7,-85.5)
				..controls(105.3,-85.7)and(103.9,-85.9)..(101.6,-87)
				..controls(100.7,-87.4)and(99.4,-87.8)..(97.9,-88)
				..controls(97,-88.2)and(96.6,-88.3)..(95.8,-88.9)
				..controls(95.4,-89.2)and(94.7,-89.6)..(94.2,-89.8)
				..controls(92.9,-90.3)and(92.3,-90.7)..(90.9,-92.1)
				..controls(90.2,-92.7)and(89.4,-93.3)..(89,-93.5)..controls(87,-94.5)and(85.2,-96)..(83.6,-97.9)
				..controls(83.3,-98.2)and(82.5,-98.9)..(81.8,-99.4)
				..controls(80.7,-100.1)and(80.5,-100.3)..(80,-101.1)
				..controls(78.7,-103.1)and(77.2,-103.9)..(76.4,-103)
				..controls(75.7,-102.2)and(75.7,-103.2)..(75.7,-83.6)
				..controls(75.6,-73.6)and(75.6,-64.8)..(75.5,-64)
				..controls(75.5,-63.2)and(75.5,-61.4)..(75.6,-59.8)--(75.7,-57.1)--(76,-56.8)
				..controls(76.8,-56.2)and(78,-56.3)..(79.8,-57.2)
				..controls(80.4,-57.5)and(81.3,-57.8)..(81.8,-57.9)
				..controls(82.3,-58.1)and(83,-58.4)..(83.6,-58.8)..controls(84.3,-59.2)and(85,-59.5)..(86,-59.8)
				..controls(86.8,-60.1)and(88,-60.6)..(88.7,-60.9)
				..controls(89.3,-61.2)and(90.6,-61.7)..(91.4,-62)
				..controls(92.2,-62.3)and(93.7,-62.9)..(94.6,-63.3)
				..controls(95.6,-63.7)and(96.8,-64.1)..(97.4,-64.3)
				..controls(98,-64.5)and(99,-64.9)..(99.7,-65.2)
				..controls(100.4,-65.6)and(101.5,-66)..(102.3,-66.2)
				..controls(104.2,-66.8)and(105,-67.1)..(106.2,-67.8)
				..controls(107.9,-68.7)and(110.2,-69.4)..(111.3,-69.2)
				..controls(112.5,-69)and(115,-68.2)..(115.6,-67.8)
				..controls(117.3,-66.7)and(117.6,-66.6)..(119.5,-66.2)
				..controls(121.9,-65.8)and(122.1,-65.7)..(123,-65.1)
				..controls(124.1,-64.4)and(124.4,-64.3)..(126,-63.9)
				..controls(127,-63.7)and(127.7,-63.4)..(128.4,-63.1)
				..controls(129,-62.8)and(129.7,-62.5)..(130.6,-62.4)
				..controls(132.1,-62.1)and(134,-61.5)..(135.2,-61)
				..controls(135.7,-60.8)and(136.7,-60.5)..(137.3,-60.3)
				..controls(138.4,-60)and(139.3,-59.7)..(141.1,-59)
				..controls(141.6,-58.7)and(142,-58.6)..(142.5,-58.6)
				..controls(143.2,-58.6)and(143.3,-58.7)..(143.6,-59)
				..controls(144,-59.5)and(144.1,-59.9)..(144,-60.8)
				..controls(144,-61.2)and(144,-71.1)..(144,-82.8)--(143.9,-104.1)--(143.7,-104.5)
				..controls(143.5,-104.7)and(143.3,-104.8)..(143,-104.9)
				..controls(142.6,-104.9)and(142.5,-104.8)..(141.5,-103.7)--cycle
				(27.8,-91.4)..controls(27.1,-90.9)and(26.3,-90)..(26.1,-89.5)
				..controls(26,-89.2)and(25.8,-88.6)..(25.7,-87.9)
				..controls(25.4,-86.2)and(25.6,-84.8)..(26.5,-83.5)
				..controls(27.4,-82.3)and(28.3,-82)..(29.4,-82.3)
				..controls(31.1,-82.8)and(32.1,-84.9)..(32,-87.4)
				..controls(31.9,-89.3)and(31.2,-90.6)..(29.9,-91.4)
				..controls(29.1,-91.9)and(28.6,-91.9)..(27.8,-91.4)--cycle
				(190.1,-91.3)..controls(189.7,-91.2)and(189.3,-90.9)..(189.1,-90.7)
				..controls(188.3,-89.8)and(187.6,-87.7)..(187.6,-86.2)
				..controls(187.6,-85.5)and(187.7,-85.2)..(188.2,-84.4)
				..controls(189.1,-82.6)and(190.1,-81.8)..(191.3,-82.1)
				..controls(192,-82.3)and(193.5,-83.7)..(193.8,-84.4)
				..controls(194.1,-85.1)and(194.2,-87.1)..(194,-88.3)
				..controls(193.8,-89.8)and(193.3,-90.5)..(192.3,-91.3)
				..controls(191.7,-91.7)and(190.9,-91.7)..(190.1,-91.3)--cycle
				(173.7,-83.4)..controls(172.2,-82.5)and(172,-82.5)..(170.6,-82.4)
				..controls(169.7,-82.4)and(168.8,-82.4)..(168.3,-82.5)
				..controls(167.2,-82.7)and(166.9,-82.7)..(166.6,-82.2)
				..controls(166.3,-81.8)and(166.3,-81)..(166.8,-80)
				..controls(167.2,-78.9)and(167.4,-77.4)..(167.7,-74.1)
				..controls(167.8,-72.8)and(167.9,-71.8)..(167.9,-71.7)
				..controls(168,-71.7)and(168,-70.2)..(168,-68.4)..controls(168,-65.4)and(168,-64.9)..(167.7,-63)
				..controls(167.3,-60.5)and(167.1,-59.9)..(166.4,-58.4)
				..controls(166.1,-57.7)and(165.8,-56.9)..(165.7,-56.6)
				..controls(165.7,-56.2)and(165.4,-55.6)..(165.2,-55.2)
				..controls(164.6,-53.8)and(161.9,-50.7)..(161,-50.3)
				..controls(159.8,-49.7)and(157.9,-50)..(155.2,-51.2)
				..controls(154.7,-51.4)and(153.4,-51.8)..(152.2,-52.1)
				..controls(150.8,-52.4)and(149.4,-52.9)..(148.6,-53.2)
				..controls(147.8,-53.6)and(146.7,-54)..(146,-54.2)
				..controls(145.4,-54.4)and(144.5,-54.7)..(144,-54.9)
				..controls(143.6,-55.2)and(142.7,-55.5)..(141.8,-55.7)
				..controls(141,-55.9)and(139.6,-56.3)..(138.7,-56.6)
				..controls(137.8,-57)and(136.3,-57.4)..(135.3,-57.6)
				..controls(133.8,-57.9)and(133.2,-58.1)..(132.1,-58.7)
				..controls(131.1,-59.2)and(130.3,-59.4)..(129.4,-59.6)
				..controls(128.7,-59.8)and(127.8,-60)..(127.4,-60.1)
				..controls(127,-60.3)and(126,-60.5)..(125.1,-60.7)
				..controls(123.2,-61.1)and(121.7,-61.5)..(120.2,-62.1)
				..controls(119.7,-62.3)and(118.1,-62.8)..(116.7,-63.2)
				..controls(114.6,-63.9)and(114,-64.1)..(113,-64.7)
				..controls(111.9,-65.4)and(111.8,-65.4)..(111,-65.5)
				..controls(109.9,-65.5)and(109.3,-65.4)..(107.7,-64.7)
				..controls(107,-64.4)and(106,-64)..(105.3,-63.9)
				..controls(103.5,-63.5)and(102.2,-63)..(101.2,-62.5)
				..controls(100.1,-61.8)and(99.3,-61.5)..(97.8,-61.1)
				..controls(97.2,-61)and(96.2,-60.7)..(95.7,-60.4)
				..controls(95.1,-60.1)and(94.2,-59.8)..(93.5,-59.6)
				..controls(92.9,-59.4)and(91.7,-59)..(91,-58.7)..controls(90.2,-58.4)and(88.9,-58)..(88,-57.7)
				..controls(87.2,-57.5)and(86.1,-57.1)..(85.5,-56.8)
				..controls(85,-56.6)and(84.1,-56.2)..(83.5,-56)..controls(82.9,-55.7)and(81.8,-55.2)..(81,-54.9)
				..controls(80.1,-54.5)and(79.1,-54.2)..(78.7,-54.1)
				..controls(77.9,-53.8)and(77.3,-53.6)..(76.2,-52.9)
				..controls(75.7,-52.6)and(75.2,-52.4)..(74.1,-52.2)
				..controls(72.5,-51.9)and(72.4,-51.8)..(71.3,-51.1)
				..controls(69.9,-50.1)and(69.4,-49.9)..(67.9,-49.8)
				..controls(67,-49.7)and(66.6,-49.6)..(66.3,-49.4)
				..controls(66.1,-49.3)and(65.7,-49)..(65.4,-48.8)
				..controls(65.1,-48.7)and(64.2,-48.4)..(63.6,-48.2)
				..controls(62.9,-48)and(61.7,-47.6)..(60.9,-47.2)
				..controls(60.1,-46.8)and(58.9,-46.4)..(58.2,-46.2)
				..controls(57.6,-46)and(56.8,-45.7)..(56.5,-45.5)
				..controls(54.9,-44.6)and(54.6,-44.5)..(53.2,-44.2)
				..controls(51.7,-43.8)and(51,-43.6)..(49.9,-42.9)
				..controls(49.5,-42.6)and(48.9,-42.4)..(47.5,-42)
				..controls(45.7,-41.6)and(45.6,-41.5)..(44.3,-40.6)
				..controls(43.5,-40.1)and(42.8,-39.6)..(42.7,-39.4)
				..controls(42.4,-38.7)and(43,-38.1)..(44.5,-37.6)..controls(44.9,-37.5)and(46,-37)..(46.9,-36.6)
				..controls(47.9,-36.1)and(48.8,-35.8)..(49.6,-35.6)
				..controls(50.5,-35.4)and(51.3,-35.2)..(52.4,-34.6)
				..controls(53.3,-34.1)and(54.5,-33.7)..(54.9,-33.5)
				..controls(55.4,-33.4)and(56.4,-33)..(57.2,-32.6)
				..controls(58.6,-32)and(60.1,-31.5)..(61.6,-31.2)
				..controls(62.3,-31.1)and(62.7,-31)..(63.2,-30.7)
				..controls(64.1,-30)and(64.3,-29.9)..(65.3,-29.7)
				..controls(65.7,-29.6)and(66.5,-29.2)..(67.1,-28.9)
				..controls(67.6,-28.6)and(68.8,-28.2)..(69.6,-27.9)
				..controls(70.5,-27.7)and(71.6,-27.3)..(72.1,-27)
				..controls(72.6,-26.7)and(73.8,-26.2)..(74.8,-25.9)
				..controls(75.7,-25.5)and(76.7,-25.1)..(77.1,-25)
				..controls(77.4,-24.8)and(78.5,-24.4)..(79.6,-24.1)
				..controls(80.6,-23.8)and(82,-23.2)..(82.8,-22.9)
				..controls(83.5,-22.5)and(84.8,-22)..(85.8,-21.8)
				..controls(86.7,-21.5)and(87.8,-21.1)..(88.1,-20.9)
				..controls(88.4,-20.7)and(89.3,-20.3)..(90.1,-20.1)
				..controls(90.8,-19.8)and(91.9,-19.4)..(92.4,-19.1)
				..controls(93.5,-18.5)and(94.4,-18.2)..(95.4,-18.1)
				..controls(95.7,-18)and(96.4,-17.8)..(96.9,-17.5)
				..controls(97.5,-17.3)and(98.4,-16.9)..(99,-16.6)
				..controls(99.7,-16.3)and(100.9,-15.9)..(101.9,-15.6)
				..controls(102.9,-15.3)and(104.4,-14.7)..(105.3,-14.3)
				..controls(108.1,-13.2)and(108.5,-13.1)..(110.1,-13.1)
				..controls(111.4,-13.1)and(111.5,-13.1)..(112.4,-13.5)
				..controls(113.2,-14)and(115,-14.8)..(116.6,-15.4)
				..controls(116.9,-15.5)and(117.9,-15.8)..(118.8,-16.1)
				..controls(120.9,-16.6)and(121.5,-16.9)..(122.6,-17.9)
				..controls(123,-18.3)and(123.7,-18.8)..(123.9,-18.9)
				..controls(124.6,-19.4)and(126,-19.8)..(127.9,-20.2)
				..controls(129.6,-20.5)and(129.8,-20.6)..(131.1,-21.4)
				..controls(131.5,-21.7)and(132.4,-22.1)..(133.1,-22.3)
				..controls(136.1,-23.2)and(138.3,-24.1)..(140.2,-25)
				..controls(141,-25.5)and(141.5,-25.6)..(142.5,-25.8)
				..controls(145.1,-26.3)and(146,-26.6)..(149.4,-28.1)
				..controls(151.3,-28.9)and(153.5,-29.8)..(154.3,-30.1)
				..controls(155.1,-30.4)and(156.2,-30.8)..(156.6,-31)
				..controls(157.1,-31.2)and(158.2,-31.6)..(159.1,-31.8)
				..controls(160.6,-32.2)and(160.9,-32.3)..(161.7,-32.9)
				..controls(162.4,-33.3)and(162.9,-33.6)..(163.7,-33.8)
				..controls(164.3,-34)and(164.8,-34.1)..(164.9,-34.1)
				..controls(165,-34.1)and(165.6,-34.4)..(166.4,-34.8)
				..controls(167.3,-35.2)and(168.1,-35.4)..(169,-35.6)
				..controls(170,-35.8)and(170.5,-36)..(171.7,-36.6)
				..controls(172.4,-37)and(173.3,-37.4)..(173.6,-37.4)
				..controls(175.6,-37.7)and(176.2,-38)..(176.7,-39.1)
				..controls(176.9,-39.6)and(176.9,-39.7)..(176.7,-40.4)
				..controls(176.4,-41.6)and(175.8,-42.7)..(175.2,-43.4)
				..controls(174.5,-44.1)and(172.6,-45.5)..(171.7,-45.8)
				..controls(170.6,-46.2)and(169.1,-46)..(168.2,-45.4)
				..controls(167.9,-45.2)and(167.3,-44.8)..(166.9,-44.4)
				..controls(166.4,-44.1)and(165.8,-43.6)..(165.4,-43.2)
				..controls(165.1,-42.9)and(164.4,-42.3)..(163.9,-41.9)
				..controls(163.3,-41.5)and(162.7,-40.9)..(162.4,-40.5)
				..controls(161.4,-39.2)and(159.9,-38)..(158.8,-37.8)
				..controls(158.4,-37.8)and(158.3,-37.7)..(158.2,-37.2)
				..controls(158,-36.8)and(157.4,-36.1)..(157.1,-36.1)
				..controls(157,-36.1)and(156.7,-36.1)..(156.5,-36)
				..controls(156.1,-35.9)and(155.9,-35.7)..(155.7,-35.2)
				..controls(155.3,-34.6)and(155.1,-34.4)..(154,-34.1)
				..controls(153.6,-34)and(153.2,-33.7)..(152.9,-33.5)
				..controls(152.4,-32.9)and(151.6,-32.4)..(151.3,-32.4)
				..controls(150.9,-32.4)and(150.4,-32.6)..(150.4,-32.9)
				..controls(150.4,-33.2)and(150.8,-33.4)..(151.4,-33.6)
				..controls(152.1,-33.7)and(152.8,-34.1)..(153.5,-34.9)
				..controls(154.1,-35.5)and(154.3,-35.6)..(155,-35.8)
				..controls(155.7,-36)and(155.7,-36)..(156,-36.6)
				..controls(156.3,-37.3)and(156.7,-37.7)..(157.5,-37.8)
				..controls(158.1,-37.9)and(158.2,-38)..(158.4,-38.8)
				..controls(158.5,-39.4)and(158.6,-39.6)..(159.6,-40.4)
				..controls(160.2,-41)and(160.9,-41.6)..(161.3,-41.8)
				..controls(161.7,-42)and(161.9,-42.3)..(162.2,-42.8)
				..controls(162.4,-43.2)and(162.8,-43.8)..(163.1,-44.1)
				..controls(163.9,-44.8)and(165.1,-46.7)..(165.7,-48)
				..controls(166.2,-49.3)and(166.4,-49.7)..(167.2,-50.5)
				..controls(167.5,-50.8)and(167.9,-51.4)..(168.1,-51.8)
				..controls(168.3,-52.2)and(168.8,-53)..(169.1,-53.6)
				..controls(169.6,-54.3)and(169.8,-54.9)..(169.9,-55.6)
				..controls(170.4,-57.8)and(170.5,-58.1)..(171.2,-59.5)
				..controls(171.6,-60.3)and(172,-61.5)..(172.2,-62.1)
				..controls(172.3,-62.8)and(172.7,-63.7)..(172.9,-64.1)
				..controls(173.2,-64.6)and(173.4,-65.3)..(173.4,-65.6)
				..controls(173.8,-68.8)and(174.2,-70.6)..(174.9,-72.6)
				..controls(175.5,-74.5)and(175.9,-77.1)..(176.1,-80.2)
				..controls(176.3,-83.4)and(175.5,-84.4)..(173.7,-83.4)--cycle;
				\definecolor{ezs}{RGB}{247,179,26}\fill[ezs](-5.7,-156.6)--(-3.4,-157.5)--(61.9,12.7)--(59.7,13.6)--cycle;
				\definecolor{ezs}{RGB}{247,179,22}\fill[ezs](-3.8,-157.3)--(-1.5,-158.2)--(63.9,12)--(61.6,12.9)--cycle;
				\definecolor{ezs}{RGB}{248,179,19}\fill[ezs](-1.8,-158.1)--(0.4,-158.9)--(65.8,11.2)--(63.5,12.1)--cycle;
				\definecolor{ezs}{RGB}{248,178,16}\fill[ezs](0.1,-158.8)--(2.4,-159.7)--(67.7,10.5)--(65.4,11.4)--cycle;
				\definecolor{ezs}{RGB}{249,178,15}\fill[ezs](2,-159.5)--(4.3,-160.4)--(69.6,9.8)--(67.3,10.7)--cycle;
				\definecolor{ezs}{RGB}{251,177,17}\fill[ezs](3.9,-160.3)--(6.2,-161.1)--(71.5,9)--(69.2,9.9)--cycle;
				\definecolor{ezs}{RGB}{252,179,18}\fill[ezs](5.8,-161)--(8.1,-161.9)--(73.4,8.3)--(71.1,9.2)--cycle;
				\definecolor{ezs}{RGB}{252,182,17}\fill[ezs](7.7,-161.7)--(10,-162.6)--(75.3,7.6)--(73,8.5)--cycle;
				\definecolor{ezs}{RGB}{252,185,17}\fill[ezs](9.6,-162.5)--(11.9,-163.4)--(77.2,6.8)--(75,7.7)--cycle;
				\definecolor{ezs}{RGB}{251,189,18}\fill[ezs](11.5,-163.2)--(13.8,-164.1)--(79.2,6.1)--(76.9,7)--cycle;
				\definecolor{ezs}{RGB}{251,191,22}\fill[ezs](13.4,-163.9)--(15.7,-164.8)--(81.1,5.4)--(78.8,6.2)--cycle;
				\definecolor{ezs}{RGB}{252,193,25}\fill[ezs](15.4,-164.7)--(17.6,-165.6)--(83,4.6)--(80.7,5.5)--cycle;
				\definecolor{ezs}{RGB}{253,195,27}\fill[ezs](17.3,-165.4)--(19.6,-166.3)--(84.9,3.9)--(82.6,4.8)--cycle;
				\definecolor{ezs}{RGB}{253,196,31}\fill[ezs](19.2,-166.1)--(21.5,-167)--(86.8,3.2)--(84.5,4)--cycle;
				\definecolor{ezs}{RGB}{254,199,35}\fill[ezs](21.1,-166.9)--(23.4,-167.8)--(88.7,2.4)--(86.4,3.3)--cycle;
				\definecolor{ezs}{RGB}{253,201,38}\fill[ezs](23,-167.6)--(25.3,-168.5)--(90.6,1.7)--(88.3,2.6)--cycle;
				\definecolor{ezs}{RGB}{253,203,40}\fill[ezs](24.9,-168.3)--(27.2,-169.2)--(92.5,1)--(90.3,1.8)--cycle;
				\definecolor{ezs}{RGB}{253,205,42}\fill[ezs](26.8,-169.1)--(29.1,-170)--(94.4,0.2)--(92.2,1.1)--cycle;
				\definecolor{ezs}{RGB}{253,207,45}\fill[ezs](28.7,-169.8)--(31,-170.7)--(96.4,-0.5)--(94.1,0.4)--cycle;
				\definecolor{ezs}{RGB}{253,210,49}\fill[ezs](30.7,-170.5)--(32.9,-171.4)--(98.3,-1.2)--(96,-0.4)--cycle;
				\definecolor{ezs}{RGB}{253,212,54}\fill[ezs](32.6,-171.3)--(34.9,-172.2)--(100.2,-2)--(97.9,-1.1)--cycle;
				\definecolor{ezs}{RGB}{253,214,59}\fill[ezs](34.5,-172)--(36.8,-172.9)--(102.1,-2.7)--(99.8,-1.8)--cycle;
				\definecolor{ezs}{RGB}{253,216,63}\fill[ezs](36.4,-172.7)--(38.7,-173.6)--(104,-3.4)--(101.7,-2.6)--cycle;
				\definecolor{ezs}{RGB}{253,218,65}\fill[ezs](38.3,-173.5)--(40.6,-174.4)--(105.9,-4.2)--(103.6,-3.3)--cycle;
				\definecolor{ezs}{RGB}{254,220,66}\fill[ezs](40.2,-174.2)--(42.5,-175.1)--(107.8,-4.9)--(105.5,-4)--cycle;
				\definecolor{ezs}{RGB}{254,222,68}\fill[ezs](42.1,-175)--(44.4,-175.8)--(109.7,-5.6)--(107.5,-4.8)--cycle;
				\definecolor{ezs}{RGB}{254,224,71}\fill[ezs](44,-175.7)--(46.3,-176.6)--(111.7,-6.4)--(109.4,-5.5)--cycle;
				\definecolor{ezs}{RGB}{254,226,75}\fill[ezs](46,-176.4)--(48.2,-177.3)--(113.6,-7.1)--(111.3,-6.2)--cycle;
				\definecolor{ezs}{RGB}{254,228,79}\fill[ezs](47.9,-177.2)--(50.1,-178)--(115.5,-7.8)--(113.2,-7)--cycle;
				\definecolor{ezs}{RGB}{254,230,82}\fill[ezs](49.8,-177.9)--(52.1,-178.8)--(117.4,-8.6)--(115.1,-7.7)--cycle;
				\definecolor{ezs}{RGB}{254,232,85}\fill[ezs](51.7,-178.6)--(54,-179.5)--(119.3,-9.3)--(117,-8.4)--cycle;
				\definecolor{ezs}{RGB}{254,234,88}\fill[ezs](53.6,-179.4)--(55.9,-180.2)--(121.2,-10)--(118.9,-9.2)--cycle;
				\definecolor{ezs}{RGB}{254,237,91}\fill[ezs](55.5,-180.1)--(57.8,-181)--(123.1,-10.8)--(120.8,-9.9)--cycle;
				\definecolor{ezs}{RGB}{254,239,94}\fill[ezs](57.4,-180.8)--(59.7,-181.7)--(125,-11.5)--(122.8,-10.6)--cycle;
				\definecolor{ezs}{RGB}{254,241,97}\fill[ezs](59.3,-181.6)--(61.6,-182.4)--(127,-12.2)--(124.7,-11.4)--cycle;
				\definecolor{ezs}{RGB}{254,243,101}\fill[ezs](61.2,-182.3)--(63.5,-183.2)--(128.9,-13)--(126.6,-12.1)--cycle;
				\definecolor{ezs}{RGB}{254,245,103}\fill[ezs](63.2,-183)--(65.4,-183.9)--(130.8,-13.7)--(128.5,-12.8)--cycle;
				\definecolor{ezs}{RGB}{254,245,105}\fill[ezs](65.1,-183.8)--(67.4,-184.6)--(132.7,-14.4)--(130.4,-13.6)--cycle;
				\definecolor{ezs}{RGB}{254,245,105}\fill[ezs](67,-184.5)--(69.3,-185.4)--(134.6,-15.2)--(132.3,-14.3)--cycle;
				\definecolor{ezs}{RGB}{254,244,103}\fill[ezs](68.9,-185.2)--(71.2,-186.1)--(136.5,-15.9)--(134.2,-15)--cycle;
				\definecolor{ezs}{RGB}{253,243,101}\fill[ezs](70.8,-186)--(73.1,-186.8)--(138.4,-16.6)--(136.1,-15.8)--cycle;
				\definecolor{ezs}{RGB}{253,241,98}\fill[ezs](72.7,-186.7)--(75,-187.6)--(140.3,-17.4)--(138,-16.5)--cycle;
				\definecolor{ezs}{RGB}{253,239,95}\fill[ezs](74.6,-187.4)--(76.9,-188.3)--(142.2,-18.1)--(140,-17.2)--cycle;
				\definecolor{ezs}{RGB}{253,236,92}\fill[ezs](76.5,-188.2)--(78.8,-189)--(144.2,-18.8)--(141.9,-18)--cycle;
				\definecolor{ezs}{RGB}{253,234,88}\fill[ezs](78.5,-188.9)--(80.7,-189.8)--(146.1,-19.6)--(143.8,-18.7)--cycle;
				\definecolor{ezs}{RGB}{253,232,85}\fill[ezs](80.4,-189.6)--(82.7,-190.5)--(148,-20.3)--(145.7,-19.4)--cycle;
				\definecolor{ezs}{RGB}{253,230,82}\fill[ezs](82.3,-190.4)--(84.6,-191.2)--(149.9,-21)--(147.6,-20.2)--cycle;
				\definecolor{ezs}{RGB}{253,228,78}\fill[ezs](84.2,-191.1)--(86.5,-192)--(151.8,-21.8)--(149.5,-20.9)--cycle;
				\definecolor{ezs}{RGB}{253,226,74}\fill[ezs](86.1,-191.8)--(88.4,-192.7)--(153.7,-22.5)--(151.4,-21.6)--cycle;
				\definecolor{ezs}{RGB}{253,224,70}\fill[ezs](88,-192.6)--(90.3,-193.4)--(155.6,-23.3)--(153.3,-22.4)--cycle;
				\definecolor{ezs}{RGB}{253,222,67}\fill[ezs](89.9,-193.3)--(92.2,-194.2)--(157.5,-24)--(155.3,-23.1)--cycle;
				\definecolor{ezs}{RGB}{253,220,65}\fill[ezs](91.8,-194)--(94.1,-194.9)--(159.5,-24.7)--(157.2,-23.8)--cycle;
				\definecolor{ezs}{RGB}{253,217,63}\fill[ezs](93.7,-194.8)--(96,-195.6)--(161.4,-25.5)--(159.1,-24.6)--cycle;
				\definecolor{ezs}{RGB}{253,215,60}\fill[ezs](95.7,-195.5)--(97.9,-196.4)--(163.3,-26.2)--(161,-25.3)--cycle;
				\definecolor{ezs}{RGB}{253,213,56}\fill[ezs](97.6,-196.2)--(99.9,-197.1)--(165.2,-26.9)--(162.9,-26)--cycle;
				\definecolor{ezs}{RGB}{252,211,53}\fill[ezs](99.5,-197)--(101.8,-197.8)--(167.1,-27.7)--(164.8,-26.8)--cycle;
				\definecolor{ezs}{RGB}{252,209,49}\fill[ezs](101.4,-197.7)--(103.7,-198.6)--(169,-28.4)--(166.7,-27.5)--cycle;
				\definecolor{ezs}{RGB}{252,207,45}\fill[ezs](103.3,-198.4)--(105.6,-199.3)--(170.9,-29.1)--(168.6,-28.2)--cycle;
				\definecolor{ezs}{RGB}{253,205,42}\fill[ezs](105.2,-199.2)--(107.5,-200)--(172.8,-29.9)--(170.6,-29)--cycle;
				\definecolor{ezs}{RGB}{253,202,40}\fill[ezs](107.1,-199.9)--(109.4,-200.8)--(174.7,-30.6)--(172.5,-29.7)--cycle;
				\definecolor{ezs}{RGB}{253,200,38}\fill[ezs](109,-200.6)--(111.3,-201.5)--(176.7,-31.3)--(174.4,-30.4)--cycle;
				\definecolor{ezs}{RGB}{253,198,35}\fill[ezs](111,-201.4)--(113.2,-202.2)--(178.6,-32.1)--(176.3,-31.2)--cycle;
				\definecolor{ezs}{RGB}{253,196,31}\fill[ezs](112.9,-202.1)--(115.2,-203)--(180.5,-32.8)--(178.2,-31.9)--cycle;
				\definecolor{ezs}{RGB}{253,194,28}\fill[ezs](114.8,-202.8)--(117.1,-203.7)--(182.4,-33.5)--(180.1,-32.6)--cycle;
				\definecolor{ezs}{RGB}{253,192,26}\fill[ezs](116.7,-203.6)--(119,-204.5)--(184.3,-34.3)--(182,-33.4)--cycle;
				\definecolor{ezs}{RGB}{253,189,24}\fill[ezs](118.6,-204.3)--(120.9,-205.2)--(186.2,-35)--(183.9,-34.1)--cycle;
				\definecolor{ezs}{RGB}{253,186,20}\fill[ezs](120.5,-205)--(122.8,-205.9)--(188.1,-35.7)--(185.8,-34.8)--cycle;
				\definecolor{ezs}{RGB}{253,184,17}\fill[ezs](122.4,-205.8)--(124.7,-206.7)--(190,-36.5)--(187.8,-35.6)--cycle;
				\definecolor{ezs}{RGB}{252,183,15}\fill[ezs](124.3,-206.5)--(126.6,-207.4)--(192,-37.2)--(189.7,-36.3)--cycle;
				\definecolor{ezs}{RGB}{252,181,14}\fill[ezs](126.3,-207.2)--(128.5,-208.1)--(193.9,-37.9)--(191.6,-37.1)--cycle;
				\definecolor{ezs}{RGB}{252,179,13}\fill[ezs](128.2,-208)--(130.4,-208.9)--(195.8,-38.7)--(193.5,-37.8)--cycle;
				\definecolor{ezs}{RGB}{252,178,12}\fill[ezs](130.1,-208.7)--(132.4,-209.6)--(197.7,-39.4)--(195.4,-38.5)--cycle;
				\definecolor{ezs}{RGB}{251,178,10}\fill[ezs](132,-209.4)--(134.3,-210.3)--(199.6,-40.1)--(197.3,-39.3)--cycle;
				\definecolor{ezs}{RGB}{251,178,10}\fill[ezs](133.9,-210.2)--(136.2,-211.1)--(201.5,-40.9)--(199.2,-40)--cycle;
				\definecolor{ezs}{RGB}{250,177,11}\fill[ezs](135.8,-210.9)--(138.1,-211.8)--(203.4,-41.6)--(201.1,-40.7)--cycle;
				\definecolor{ezs}{RGB}{250,178,11}\fill[ezs](137.7,-211.6)--(140,-212.5)--(205.3,-42.3)--(203.1,-41.5)--cycle;
				\definecolor{ezs}{RGB}{248,180,10}\fill[ezs](139.6,-212.4)--(141.9,-213.3)--(207.3,-43.1)--(205,-42.2)--cycle;
				\definecolor{ezs}{RGB}{247,181,9}\fill[ezs](141.5,-213.1)--(143.8,-214)--(209.2,-43.8)--(206.9,-42.9)--cycle;
				\definecolor{ezs}{RGB}{246,182,9}\fill[ezs](143.5,-213.8)--(145.7,-214.7)--(211.1,-44.5)--(208.8,-43.7)--cycle;
				\definecolor{ezs}{RGB}{246,182,10}\fill[ezs](145.4,-214.6)--(147.7,-215.5)--(213,-45.3)--(210.7,-44.4)--cycle;
			\end{scope}
			% chữ EDUZY (Times đậm, làm dày nét bằng 5 lớp lệch nhẹ; căn đúng khung chữ logo gốc)
			\foreach \dx/\dy in {-0.7/0,0.7/0,0/-0.3,0/0.3,0/0}
			\node[anchor=center,inner sep=0pt,outer sep=0pt,text=ezText,align=center] at (109.6+\dx,-193.6+\dy)
			{\resizebox{1.01cm}{0.22cm}{\fontencoding{T1}\fontfamily{ptm}\fontseries{b}\selectfont EDUZY}};
		\end{tikzpicture}%
	}%
}
%====================== THIẾT LẬP THÂN ĐỀ ======================
%--- Môi trường tự luận {bt}: ex_test 3.3.4 chưa khai báo sẵn nên tạo ở đây
%    (ex_test đã tự móc \fixshowans và \listenumerate cho bt)
\makeatletter
\@ifundefined{bt}{\newtheorem{bt}{\color{gold1}\namebt}}{}
\makeatother
\renewcommand{\sectionmark}[1]{\markright{\thesection\ #1}}
%--- Khung tiêu đề (logo + tối đa 4 dòng; dòng 3, 4 để trống {} thì bỏ qua)
%  Đề thi:    \khungtieude{THI THỬ THPT QUỐC GIA NĂM 2026}{Môn Toán -- Đề số 21}
%                         {(Đề thi gồm: ...)}{Thời gian làm bài: 90 phút (không kể thời gian phát đề).}
%  Chuyên đề: \khungtieude{CHUYÊN ĐỀ 3 -- DÃY SỐ VÀ CÔNG THỨC TRUY HỒI}
%                         {Bồi dưỡng học sinh giỏi Toán THPT}{...}{}
\newcommand{\khungtieude}[4]{%
	\begin{tcolorbox}[colframe=gold1,colback=gold5!45,boxrule=0.5mm]
		\begin{minipage}[c]{3cm}
			\centerline{\eduzylogo[2.4cm]}
		\end{minipage}\hspace{0.5cm}
		\begin{minipage}[c]{12.5cm}
			\vskip 0.15 cm
			\centerline{\ind{#1}}
			\vskip 0.15 cm
			\centerline{\textbf{#2}}
			\if\relax\detokenize{#3}\relax\else
				\vskip 0.15 cm
				\centerline{\textit{#3}}
			\fi
			\if\relax\detokenize{#4}\relax\else
				\vskip 0.15 cm
				\centerline{\textit{#4}}
			\fi
			\vskip 0.15 cm
		\end{minipage}
	\end{tcolorbox}
}
%====================== BỔ SUNG CHO BÀI HỌC, CHUYÊN ĐỀ ======================
\usepackage[shortlabels]{enumitem}
\usepackage{tabularx}
%--- Danh sách gọn: \begin{enumerate}[gon] ... \end{enumerate}
\SetEnumitemKey{gon}{leftmargin=*,itemsep=1pt,parsep=0pt,topsep=3pt}
%--- Khung lý thuyết: \begin{khung}{} ... \end{khung}; tiêu đề chỉ dùng khi thật ngắn ("Mẹo", "Lỗi thường gặp")
\newtcolorbox{khung}[1]{breakable,enhanced,colframe=gold1,colback=gold5!45,
	boxrule=0.4mm,fonttitle=\bfseries\sffamily,coltitle=white,colbacktitle=gold1,
	attach boxed title to top left={xshift=6mm,yshift=-3mm},
	boxed title style={arc=1mm,boxrule=0.3mm},title={#1}}
%--- Tiêu đề mục nhỏ trong một dạng: \muc{Dấu hiệu}, \muc{Phương pháp}, \muc{Bài tự luyện}
\makeatletter
\newcommand{\muc}[1]{\par\medskip\noindent\indam{#1}\par\nobreak\smallskip\@afterindentfalse\@afterheading}
\makeatother
%--- Ô công thức cần nhớ (căn giữa, viền vàng, vừa khít nội dung)
%     \begin{center}\congthuc{$Q(n)=(n-1)!$}\end{center}
\newtcbox{\congthuc}{on line,colframe=gold1,colback=gold5!45,boxrule=0.5pt,arc=2pt,
	left=8pt,right=8pt,top=4pt,bottom=4pt}
%--- Công tắc lời giải bài tự luyện (ví dụ mẫu luôn dùng \loigiai)
%     \loigiaiTL[{đáp số}]{lời giải đầy đủ}
%     \TLfull -> in lời giải đầy đủ;  \TLdapso -> chỉ in một dòng "Đáp số";  \TLan -> ẩn hoàn toàn
%     (\CheDo là dethi thì mọi lời giải và đáp số đều ẩn)
\makeatletter
\newif\ifTLhien   \TLhientrue
\newif\ifTLdapso  \TLdapsofalse
\newcommand{\TLfull}{\TLhientrue\TLdapsofalse}
\newcommand{\TLdapso}{\TLhientrue\TLdapsotrue}
\newcommand{\TLan}{\TLhienfalse}
\newcommand{\loigiaiTL}[2][]{%
	\ifTLhien
		\ifTLdapso
			\def\@tlarg{#1}%
			\ifx\@tlarg\@empty\else
				\ifdefstring{\kindOPTN}{dethi}{}{\par\noindent{\bfseries\color{gold1}Đáp số:} #1\par}%
			\fi
		\else
			\loigiai{#2}%
		\fi
	\fi}
\makeatother
\TLdapso% mặc định: chỉ in đáp số bài tự luyện
%====================== BẢNG ĐÚNG/SAI (\choiceTFt) ======================
%  Hàng tiêu đề nền nâu đồng chữ trắng (Phát biểu | Đ | S), các hàng tô xen kẽ kem / trắng,
%  đường kẻ vàng kim, nhãn a) b) c) d) màu nâu đồng. Bản có đáp án: cột được chọn có vòng tròn
%  nâu đồng với dấu tích trắng (câu đúng tích ở cột Đúng, câu sai tích ở cột Sai).
%  Tiêu đề ghi đủ "Đúng | Sai": \OPTN{viettat=0} sau \input (cột Đúng/Sai rộng ra, phát biểu hẹp lại).
%  Dùng bảng gốc của ex_test: \def\BangDSgoc{1} trước \input{preamble-eduzy-v2}.
\providecommand{\BangDSgoc}{0}
\ifnum\BangDSgoc=0
\usepackage{colortbl,pifont}
\colorlet{ezDSke}{gold2!60!gold1}% đường kẻ bảng (nét 0,4pt như ex_test để không đổi chiều cao bảng)
\colorlet{ezDSle}{gold5!60}% hàng a), c)
\colorlet{ezDSchan}{white}% hàng b), d)
%--- dấu chọn ở cột Đúng/Sai
%    (bề rộng 0, căn giữa ô: không làm đổi bề rộng cột của ex_test, không báo tràn lề)
\newcommand{\EZdauchon}{\makebox[0pt]{\tikz[baseline=-1ex]\node[circle,fill=gold1,inner sep=0pt,
	minimum size=1.2em,text=white,font=\small]{\ding{51}};}}
\def\tickT{\EZdauchon}\def\tickF{\EZdauchon}
\AtBeginDocument{\sbox{\ccTFTab}{\EZdauchon}}
%--- nhãn a) b) c) d) màu nâu đồng (ở bản lời giải, nhãn câu đúng vẫn đóng khung như ex_test)
\renewcommand{\FalseTF}{\stepcounter{dapan}\textbf{\color{gold1}\DapAnTF\sepTF}\ignorespaces}
%--- trước khi dựng bảng: kẻ màu vàng kim (bề rộng cột giữ nguyên như ex_test); sau bảng trả lại màu đen
\newcount\EZhangTF
\newcommand{\EZbatdauTF}{%
	\global\EZhangTF=0\relax
	\arrayrulecolor{ezDSke}}
\newcommand{\EZketthucTF}{\arrayrulecolor{black}}
%--- hàng tiêu đề nền nâu đồng, chữ trắng
\makeatletter
\newcommand{\EZo@TF}[1]{\centering\bfseries\color{white}#1}
\newcommand{\EZtieudeTF}{%
	\ifcase\cotDS
		\ifnum\cotSTT=1\relax
			\def\rowHeader{\hline\rowcolor{gold1}\EZo@TF{TT}&\EZo@TF{\phatbieu}\tabularnewline}%
		\else
			\def\rowHeader{\hline\rowcolor{gold1}\EZo@TF{\phatbieu}\tabularnewline}%
		\fi
	\or
		\ifnum\cotSTT=1\relax
			\def\rowHeader{\hline\rowcolor{gold1}\EZo@TF{TT}&\EZo@TF{\phatbieu}&\EZo@TF{\Dung}\tabularnewline}%
		\else
			\def\rowHeader{\hline\rowcolor{gold1}\EZo@TF{\phatbieu}&\EZo@TF{\Dung}\tabularnewline}%
		\fi
	\or
		\ifnum\cotSTT=1\relax
			\def\rowHeader{\hline\rowcolor{gold1}\EZo@TF{TT}&\EZo@TF{\phatbieu}&\EZo@TF{\Dung}&\EZo@TF{\Sai}\tabularnewline}%
		\else
			\def\rowHeader{\hline\rowcolor{gold1}\EZo@TF{\phatbieu}&\EZo@TF{\Dung}&\EZo@TF{\Sai}\tabularnewline}%
		\fi
	\fi}
%--- tô xen kẽ từng hàng (gọi ngay sau \hline đầu mỗi hàng)
\newcommand{\EZmauhangTF}{%
	\noalign{\global\advance\EZhangTF\@ne}%
	\ifodd\EZhangTF\rowcolor{ezDSle}\else\rowcolor{ezDSchan}\fi}
%--- gắn vào \choiceTFt / \writeChoiceTFt của ex_test
\xpretocmd{\choiceTFt}{\EZbatdauTF}{}{\typeout{[preamble-eduzy-v2] Không gắn được bảng Đúng/Sai (đầu bảng)}}
\xpatchcmd{\choiceTFt}{\ifnum\dongPB=0\def\rowHeader{}\fi}{\EZtieudeTF\ifnum\dongPB=0\def\rowHeader{}\fi}
	{}{\typeout{[preamble-eduzy-v2] Không gắn được bảng Đúng/Sai (hàng tiêu đề)}}
\xpatchcmd{\writeChoiceTFt}{\hline\xdef\dotEX{\sdotEX}}{\hline\EZmauhangTF\xdef\dotEX{\sdotEX}}
	{}{\typeout{[preamble-eduzy-v2] Không gắn được bảng Đúng/Sai (tô hàng)}}
\xpatchcmd{\writeChoiceTFt}{{\end{longtable}}}{{\end{longtable}\EZketthucTF}}
	{}{\typeout{[preamble-eduzy-v2] Không gắn được bảng Đúng/Sai (cuối longtable)}}
\xpatchcmd{\writeChoiceTFt}{{\end{xtabular}\end{center}\medskip\noindent}}{{\end{xtabular}\end{center}\EZketthucTF\medskip\noindent}}
	{}{\typeout{[preamble-eduzy-v2] Không gắn được bảng Đúng/Sai (cuối xtabular)}}
\fi
%=====================================================================
\renewcommand{\parallel}{\mathrel{/\!\!/}}

%     %=====================================================================
% preamble-vatli-hinh.tex -- BỘ MACRO HÌNH VẬT LÍ (dùng chung cho mọi đề Vật lí)
% Gọi ngay sau preamble chính:
%     \input{preamble-eduzy-v2}
%     \input{preamble-vatli-hinh}
% Tên file bắt đầu bằng "preamble" nên dong-goi.sh tự chép vào zip thành phẩm;
% tools/so_hinh.py đọc file này khi dựng ảnh so hình.
% Lưu ý màu: preamble-eduzy-v2 đổi blue, cyan, teal sang tông vàng đồng (gold1, gold2).
% Cần màu xanh thật thì dùng VLxanh (hoặc khai báo \definecolor cục bộ).
% Rút từ đề mẫu VatLi2026-CaoBang-SoGD-Lan1 (thầy duyệt 06/11/2026).
%=====================================================================

%--- Mạch điện: circuitikz kiểu châu Âu (điện trở hình chữ nhật...)
\usepackage[european]{circuitikz}

%--- Màu dùng chung
\definecolor{VLdo}{RGB}{214,72,60}
\definecolor{VLxanh}{RGB}{60,100,200}
\definecolor{VLdong}{RGB}{200,128,70}
\definecolor{VLgo}{RGB}{186,132,84}
\definecolor{VLla}{RGB}{150,200,150}

%--- Mũi tên: VLmui (mũi tên ở đầu đường), VLgiua=<vị trí> (mũi tên giữa đường, mặc định 0,55)
\tikzset{
	VLmui/.style={-{Stealth[length=4.5pt,width=3.5pt]}},
	VLgiua/.style={postaction={decorate},decoration={markings,mark=at position #1 with {\arrow{Stealth[length=6.5pt,width=5pt]}}}},
	VLgiua/.default=0.55
}
% cọc nối dây của điện kế: \VLcoc{x}{y}{màu}
\newcommand{\VLcoc}[3]{%
	\begin{scope}[shift={(#1,#2)}]
		\shade[left color=#3!45!black,right color=#3!55!black,middle color=#3!80!white] (-0.17,0) rectangle (0.17,0.28);
		\shade[left color=#3!55!black,right color=#3!65!black,middle color=#3!60!white] (-0.14,0.28) rectangle (0.14,0.5);
		\fill[#3!55!black] (0,0.5) ellipse (0.14 and 0.045);
		\fill[black!85] (0,0.5) ellipse (0.055 and 0.02);
	\end{scope}}
% phích cắm có rãnh: \VLphich{x}{y}{màu}
\newcommand{\VLphich}[3]{%
	\begin{scope}[shift={(#1,#2)}]
		\shade[left color=#3!50!black,right color=#3!60!black,middle color=#3!55!white] (-0.12,0) rectangle (0.12,0.62);
		\draw[#3!45!black,very thin] foreach \y in {0.1,0.18,...,0.5}{(-0.12,\y)--(0.12,\y)};
		\shade[left color=#3!50!black,right color=#3!60!black,middle color=#3!55!white] (-0.08,0.62) rectangle (0.08,0.85);
		\fill[#3!45!black] (0,0.85) ellipse (0.08 and 0.03);
	\end{scope}}
% điện kế nhìn chính diện, gốc ở góc trái dưới, rộng 4, cao 4: \VLdienke{1: thang số | 2: thang -, +}
\newcommand{\VLdienke}[1]{%
	\fill[black!82,rounded corners=2pt] (0,0) rectangle (4,1.35);
	\fill[black!60] (0.05,1.22) rectangle (3.95,1.35);
	\fill[VLla,draw=black!60,rounded corners=1pt] (0.28,0.25) rectangle (3.72,1.02);
	\fill[gray!8,draw=gray!65,rounded corners=1.5pt] (0.08,1.35) rectangle (3.92,4.05);
	\fill[white,draw=gray!45] (0.28,1.5) rectangle (3.72,3.85);
	\fill[gray!18] (1.25,1.52) rectangle (2.75,2.05);
	\draw[gray!40,very thin] foreach \x in {1.3,1.4,...,2.7}{(\x,1.55)--(\x,2.0)};
	\draw[black!80] (2,1.35) ++(127:1.95) arc[start angle=127,end angle=53,radius=1.95];
	\foreach \g in {53,55.4667,...,127.1}{\draw[black!80,very thin] (2,1.35)++(\g:1.95)--++(\g:0.12);}
	\foreach \g in {53,65.3333,...,127.1}{\draw[black!85] (2,1.35)++(\g:1.95)--++(\g:0.22);}
	\ifnum#1=1
		\foreach \g/\s in {127/15,114.6667/10,102.3333/5,90/0,77.6667/5,65.3333/10,53/15}{\node[font=\fontsize{5}{5}\selectfont,inner sep=0] at ($(2,1.35)+(\g:2.37)$) {\s};}
		\node[font=\scriptsize\bfseries] at (2,2.75) {G};
	\else
		\node[font=\scriptsize] at ($(2,1.35)+(127:1.7)$) {$-$};
		\node[font=\scriptsize] at ($(2,1.35)+(53:1.7)$) {$+$};
	\fi
	\draw[VLdo,semithick] (2,1.8)--(2,3.42);
	\fill[black!85] (2,1.8) circle (0.13);
	\fill[gray!55] (2,1.8) circle (0.045);
	\draw[white,line width=1.2pt,opacity=0.7] (0.2,3.95)--(0.2,1.6);
}
% nam châm thẳng 3D: \VLnamcham{dài}{cao}{màu trái}{màu phải}{chữ trái}{chữ phải}
\newcommand{\VLnamcham}[6]{%
	\fill[#3!55!white] (0,#2)--(0.5*#1,#2)--(0.5*#1+0.22,#2+0.13)--(0.22,#2+0.13)--cycle;
	\fill[#4!55!white] (0.5*#1,#2)--(#1,#2)--(#1+0.22,#2+0.13)--(0.5*#1+0.22,#2+0.13)--cycle;
	\fill[#4!70!black] (#1,0)--(#1+0.22,0.13)--(#1+0.22,#2+0.13)--(#1,#2)--cycle;
	\shade[top color=#3!75!white,bottom color=#3!95!black] (0,0) rectangle (0.5*#1,#2);
	\shade[top color=#4!75!white,bottom color=#4!95!black] (0.5*#1,0) rectangle (#1,#2);
	\node[white,font=\scriptsize\bfseries] at (0.25*#1,0.5*#2) {#5};
	\node[white,font=\scriptsize\bfseries] at (0.75*#1,0.5*#2) {#6};
}
% lưới chấm (từ trường hướng ra) / lưới dấu nhân (hướng vào)
\newcommand{\VLcham}[2]{\fill (#1,#2) circle (1.6pt);}
\newcommand{\VLnhan}[2]{\draw[semithick] (#1,#2)++(-0.06,-0.06)--++(0.12,0.12) (#1,#2)++(-0.06,0.06)--++(0.12,-0.12);}
% thanh dẫn hình trụ nằm ngang / thẳng đứng (tâm, nửa dài)
\newcommand{\VLthanhngang}[3]{%
	\shade[top color=gray!60,bottom color=gray!85,middle color=gray!35,draw=black!85] (#1-#3,#2-0.1) rectangle (#1+#3,#2+0.1);
	\filldraw[fill=gray!35,draw=black!80] (#1+#3,#2) ellipse (0.04 and 0.1);}
\newcommand{\VLthanhdung}[3]{%
	\shade[left color=gray!85,right color=gray!60,middle color=gray!35,draw=black!85] (#1-0.1,#2-#3) rectangle (#1+0.1,#2+#3);
	\filldraw[fill=gray!35,draw=black!80] (#1,#2+#3) ellipse (0.1 and 0.035);}

% Tên file bắt đầu bằng "preamble" nên dong-goi.sh tự chép vào zip thành phẩm;
% tools/so_hinh.py đọc file này khi dựng ảnh so hình.
% Lưu ý màu: preamble-eduzy-v2 đổi blue, cyan, teal sang tông vàng đồng (gold1, gold2).
% Cần màu xanh thật thì dùng VLxanh (hoặc khai báo \definecolor cục bộ).
% Rút từ đề mẫu VatLi2026-CaoBang-SoGD-Lan1 (thầy duyệt 06/11/2026).
%=====================================================================

%--- Mạch điện: circuitikz kiểu châu Âu (điện trở hình chữ nhật...)
\usepackage[european]{circuitikz}

%--- Màu dùng chung
\definecolor{VLdo}{RGB}{214,72,60}
\definecolor{VLxanh}{RGB}{60,100,200}
\definecolor{VLdong}{RGB}{200,128,70}
\definecolor{VLgo}{RGB}{186,132,84}
\definecolor{VLla}{RGB}{150,200,150}

%--- Mũi tên: VLmui (mũi tên ở đầu đường), VLgiua=<vị trí> (mũi tên giữa đường, mặc định 0,55)
\tikzset{
	VLmui/.style={-{Stealth[length=4.5pt,width=3.5pt]}},
	VLgiua/.style={postaction={decorate},decoration={markings,mark=at position #1 with {\arrow{Stealth[length=6.5pt,width=5pt]}}}},
	VLgiua/.default=0.55
}
% cọc nối dây của điện kế: \VLcoc{x}{y}{màu}
\newcommand{\VLcoc}[3]{%
	\begin{scope}[shift={(#1,#2)}]
		\shade[left color=#3!45!black,right color=#3!55!black,middle color=#3!80!white] (-0.17,0) rectangle (0.17,0.28);
		\shade[left color=#3!55!black,right color=#3!65!black,middle color=#3!60!white] (-0.14,0.28) rectangle (0.14,0.5);
		\fill[#3!55!black] (0,0.5) ellipse (0.14 and 0.045);
		\fill[black!85] (0,0.5) ellipse (0.055 and 0.02);
	\end{scope}}
% phích cắm có rãnh: \VLphich{x}{y}{màu}
\newcommand{\VLphich}[3]{%
	\begin{scope}[shift={(#1,#2)}]
		\shade[left color=#3!50!black,right color=#3!60!black,middle color=#3!55!white] (-0.12,0) rectangle (0.12,0.62);
		\draw[#3!45!black,very thin] foreach \y in {0.1,0.18,...,0.5}{(-0.12,\y)--(0.12,\y)};
		\shade[left color=#3!50!black,right color=#3!60!black,middle color=#3!55!white] (-0.08,0.62) rectangle (0.08,0.85);
		\fill[#3!45!black] (0,0.85) ellipse (0.08 and 0.03);
	\end{scope}}
% điện kế nhìn chính diện, gốc ở góc trái dưới, rộng 4, cao 4: \VLdienke{1: thang số | 2: thang -, +}
\newcommand{\VLdienke}[1]{%
	\fill[black!82,rounded corners=2pt] (0,0) rectangle (4,1.35);
	\fill[black!60] (0.05,1.22) rectangle (3.95,1.35);
	\fill[VLla,draw=black!60,rounded corners=1pt] (0.28,0.25) rectangle (3.72,1.02);
	\fill[gray!8,draw=gray!65,rounded corners=1.5pt] (0.08,1.35) rectangle (3.92,4.05);
	\fill[white,draw=gray!45] (0.28,1.5) rectangle (3.72,3.85);
	\fill[gray!18] (1.25,1.52) rectangle (2.75,2.05);
	\draw[gray!40,very thin] foreach \x in {1.3,1.4,...,2.7}{(\x,1.55)--(\x,2.0)};
	\draw[black!80] (2,1.35) ++(127:1.95) arc[start angle=127,end angle=53,radius=1.95];
	\foreach \g in {53,55.4667,...,127.1}{\draw[black!80,very thin] (2,1.35)++(\g:1.95)--++(\g:0.12);}
	\foreach \g in {53,65.3333,...,127.1}{\draw[black!85] (2,1.35)++(\g:1.95)--++(\g:0.22);}
	\ifnum#1=1
		\foreach \g/\s in {127/15,114.6667/10,102.3333/5,90/0,77.6667/5,65.3333/10,53/15}{\node[font=\fontsize{5}{5}\selectfont,inner sep=0] at ($(2,1.35)+(\g:2.37)$) {\s};}
		\node[font=\scriptsize\bfseries] at (2,2.75) {G};
	\else
		\node[font=\scriptsize] at ($(2,1.35)+(127:1.7)$) {$-$};
		\node[font=\scriptsize] at ($(2,1.35)+(53:1.7)$) {$+$};
	\fi
	\draw[VLdo,semithick] (2,1.8)--(2,3.42);
	\fill[black!85] (2,1.8) circle (0.13);
	\fill[gray!55] (2,1.8) circle (0.045);
	\draw[white,line width=1.2pt,opacity=0.7] (0.2,3.95)--(0.2,1.6);
}
% nam châm thẳng 3D: \VLnamcham{dài}{cao}{màu trái}{màu phải}{chữ trái}{chữ phải}
\newcommand{\VLnamcham}[6]{%
	\fill[#3!55!white] (0,#2)--(0.5*#1,#2)--(0.5*#1+0.22,#2+0.13)--(0.22,#2+0.13)--cycle;
	\fill[#4!55!white] (0.5*#1,#2)--(#1,#2)--(#1+0.22,#2+0.13)--(0.5*#1+0.22,#2+0.13)--cycle;
	\fill[#4!70!black] (#1,0)--(#1+0.22,0.13)--(#1+0.22,#2+0.13)--(#1,#2)--cycle;
	\shade[top color=#3!75!white,bottom color=#3!95!black] (0,0) rectangle (0.5*#1,#2);
	\shade[top color=#4!75!white,bottom color=#4!95!black] (0.5*#1,0) rectangle (#1,#2);
	\node[white,font=\scriptsize\bfseries] at (0.25*#1,0.5*#2) {#5};
	\node[white,font=\scriptsize\bfseries] at (0.75*#1,0.5*#2) {#6};
}
% lưới chấm (từ trường hướng ra) / lưới dấu nhân (hướng vào)
\newcommand{\VLcham}[2]{\fill (#1,#2) circle (1.6pt);}
\newcommand{\VLnhan}[2]{\draw[semithick] (#1,#2)++(-0.06,-0.06)--++(0.12,0.12) (#1,#2)++(-0.06,0.06)--++(0.12,-0.12);}
% thanh dẫn hình trụ nằm ngang / thẳng đứng (tâm, nửa dài)
\newcommand{\VLthanhngang}[3]{%
	\shade[top color=gray!60,bottom color=gray!85,middle color=gray!35,draw=black!85] (#1-#3,#2-0.1) rectangle (#1+#3,#2+0.1);
	\filldraw[fill=gray!35,draw=black!80] (#1+#3,#2) ellipse (0.04 and 0.1);}
\newcommand{\VLthanhdung}[3]{%
	\shade[left color=gray!85,right color=gray!60,middle color=gray!35,draw=black!85] (#1-0.1,#2-#3) rectangle (#1+0.1,#2+#3);
	\filldraw[fill=gray!35,draw=black!80] (#1,#2+#3) ellipse (0.1 and 0.035);}

% Tên file bắt đầu bằng "preamble" nên dong-goi.sh tự chép vào zip thành phẩm;
% tools/so_hinh.py đọc file này khi dựng ảnh so hình.
% Lưu ý màu: preamble-eduzy-v2 đổi blue, cyan, teal sang tông vàng đồng (gold1, gold2).
% Cần màu xanh thật thì dùng VLxanh (hoặc khai báo \definecolor cục bộ).
% Rút từ đề mẫu VatLi2026-CaoBang-SoGD-Lan1 (thầy duyệt 06/11/2026).
%=====================================================================

%--- Mạch điện: circuitikz kiểu châu Âu (điện trở hình chữ nhật...)
\usepackage[european]{circuitikz}

%--- Màu dùng chung
\definecolor{VLdo}{RGB}{214,72,60}
\definecolor{VLxanh}{RGB}{60,100,200}
\definecolor{VLdong}{RGB}{200,128,70}
\definecolor{VLgo}{RGB}{186,132,84}
\definecolor{VLla}{RGB}{150,200,150}

%--- Mũi tên: VLmui (mũi tên ở đầu đường), VLgiua=<vị trí> (mũi tên giữa đường, mặc định 0,55)
\tikzset{
	VLmui/.style={-{Stealth[length=4.5pt,width=3.5pt]}},
	VLgiua/.style={postaction={decorate},decoration={markings,mark=at position #1 with {\arrow{Stealth[length=6.5pt,width=5pt]}}}},
	VLgiua/.default=0.55
}
% cọc nối dây của điện kế: \VLcoc{x}{y}{màu}
\newcommand{\VLcoc}[3]{%
	\begin{scope}[shift={(#1,#2)}]
		\shade[left color=#3!45!black,right color=#3!55!black,middle color=#3!80!white] (-0.17,0) rectangle (0.17,0.28);
		\shade[left color=#3!55!black,right color=#3!65!black,middle color=#3!60!white] (-0.14,0.28) rectangle (0.14,0.5);
		\fill[#3!55!black] (0,0.5) ellipse (0.14 and 0.045);
		\fill[black!85] (0,0.5) ellipse (0.055 and 0.02);
	\end{scope}}
% phích cắm có rãnh: \VLphich{x}{y}{màu}
\newcommand{\VLphich}[3]{%
	\begin{scope}[shift={(#1,#2)}]
		\shade[left color=#3!50!black,right color=#3!60!black,middle color=#3!55!white] (-0.12,0) rectangle (0.12,0.62);
		\draw[#3!45!black,very thin] foreach \y in {0.1,0.18,...,0.5}{(-0.12,\y)--(0.12,\y)};
		\shade[left color=#3!50!black,right color=#3!60!black,middle color=#3!55!white] (-0.08,0.62) rectangle (0.08,0.85);
		\fill[#3!45!black] (0,0.85) ellipse (0.08 and 0.03);
	\end{scope}}
% điện kế nhìn chính diện, gốc ở góc trái dưới, rộng 4, cao 4: \VLdienke{1: thang số | 2: thang -, +}
\newcommand{\VLdienke}[1]{%
	\fill[black!82,rounded corners=2pt] (0,0) rectangle (4,1.35);
	\fill[black!60] (0.05,1.22) rectangle (3.95,1.35);
	\fill[VLla,draw=black!60,rounded corners=1pt] (0.28,0.25) rectangle (3.72,1.02);
	\fill[gray!8,draw=gray!65,rounded corners=1.5pt] (0.08,1.35) rectangle (3.92,4.05);
	\fill[white,draw=gray!45] (0.28,1.5) rectangle (3.72,3.85);
	\fill[gray!18] (1.25,1.52) rectangle (2.75,2.05);
	\draw[gray!40,very thin] foreach \x in {1.3,1.4,...,2.7}{(\x,1.55)--(\x,2.0)};
	\draw[black!80] (2,1.35) ++(127:1.95) arc[start angle=127,end angle=53,radius=1.95];
	\foreach \g in {53,55.4667,...,127.1}{\draw[black!80,very thin] (2,1.35)++(\g:1.95)--++(\g:0.12);}
	\foreach \g in {53,65.3333,...,127.1}{\draw[black!85] (2,1.35)++(\g:1.95)--++(\g:0.22);}
	\ifnum#1=1
		\foreach \g/\s in {127/15,114.6667/10,102.3333/5,90/0,77.6667/5,65.3333/10,53/15}{\node[font=\fontsize{5}{5}\selectfont,inner sep=0] at ($(2,1.35)+(\g:2.37)$) {\s};}
		\node[font=\scriptsize\bfseries] at (2,2.75) {G};
	\else
		\node[font=\scriptsize] at ($(2,1.35)+(127:1.7)$) {$-$};
		\node[font=\scriptsize] at ($(2,1.35)+(53:1.7)$) {$+$};
	\fi
	\draw[VLdo,semithick] (2,1.8)--(2,3.42);
	\fill[black!85] (2,1.8) circle (0.13);
	\fill[gray!55] (2,1.8) circle (0.045);
	\draw[white,line width=1.2pt,opacity=0.7] (0.2,3.95)--(0.2,1.6);
}
% nam châm thẳng 3D: \VLnamcham{dài}{cao}{màu trái}{màu phải}{chữ trái}{chữ phải}
\newcommand{\VLnamcham}[6]{%
	\fill[#3!55!white] (0,#2)--(0.5*#1,#2)--(0.5*#1+0.22,#2+0.13)--(0.22,#2+0.13)--cycle;
	\fill[#4!55!white] (0.5*#1,#2)--(#1,#2)--(#1+0.22,#2+0.13)--(0.5*#1+0.22,#2+0.13)--cycle;
	\fill[#4!70!black] (#1,0)--(#1+0.22,0.13)--(#1+0.22,#2+0.13)--(#1,#2)--cycle;
	\shade[top color=#3!75!white,bottom color=#3!95!black] (0,0) rectangle (0.5*#1,#2);
	\shade[top color=#4!75!white,bottom color=#4!95!black] (0.5*#1,0) rectangle (#1,#2);
	\node[white,font=\scriptsize\bfseries] at (0.25*#1,0.5*#2) {#5};
	\node[white,font=\scriptsize\bfseries] at (0.75*#1,0.5*#2) {#6};
}
% lưới chấm (từ trường hướng ra) / lưới dấu nhân (hướng vào)
\newcommand{\VLcham}[2]{\fill (#1,#2) circle (1.6pt);}
\newcommand{\VLnhan}[2]{\draw[semithick] (#1,#2)++(-0.06,-0.06)--++(0.12,0.12) (#1,#2)++(-0.06,0.06)--++(0.12,-0.12);}
% thanh dẫn hình trụ nằm ngang / thẳng đứng (tâm, nửa dài)
\newcommand{\VLthanhngang}[3]{%
	\shade[top color=gray!60,bottom color=gray!85,middle color=gray!35,draw=black!85] (#1-#3,#2-0.1) rectangle (#1+#3,#2+0.1);
	\filldraw[fill=gray!35,draw=black!80] (#1+#3,#2) ellipse (0.04 and 0.1);}
\newcommand{\VLthanhdung}[3]{%
	\shade[left color=gray!85,right color=gray!60,middle color=gray!35,draw=black!85] (#1-0.1,#2-#3) rectangle (#1+0.1,#2+#3);
	\filldraw[fill=gray!35,draw=black!80] (#1,#2+#3) ellipse (0.1 and 0.035);}

% nhiệt kế y tế thủy ngân nằm nghiêng (thân thủy tinh, thang 35–42, đầu bầu kim loại bên trái), gốc ở đầu bầu, dài 5,2:
% \VLnhietkeyte{góc nghiêng} (đặt trong scope có shift)
\newcommand{\VLnhietkeyte}[1]{%
	\begin{scope}[rotate=#1]
		\shade[left color=gray!70,right color=gray!40] (0,0)--(0.32,-0.07)--(0.32,0.07)--cycle;
		\shade[top color=gray!8,bottom color=gray!40,middle color=white,rounded corners=2pt] (0.3,-0.2) rectangle (5.2,0.2);
		\draw[gray!55!black,thin,rounded corners=2pt] (0.3,-0.2) rectangle (5.2,0.2);
		\fill[VLluc!80] (4.55,-0.2) rectangle (4.75,0.2);
		\fill[gray!65] (0.3,-0.035) rectangle (3.0,0.035);
		\foreach \VLk in {0,...,35}{\draw[black!70,very thin] ({1.2+0.09*\VLk},0.2)--++(0,{-0.05-0.04*(\VLk==int(\VLk/5)*5)});}
		\foreach \VLs/\VLx/\VLy in {35/1.2/-0.1,36/1.65/0.11,37/2.1/-0.1,38/2.55/0.11,39/3.0/-0.1,40/3.45/0.11,41/3.9/-0.1,42/4.35/0.11}{\node[font=\fontsize{4.5}{4.5}\selectfont,inner sep=0pt] at (\VLx,\VLy) {\VLs};}
	\end{scope}}
\begin{document}
	\khungtieude{ĐỀ THI THỬ THPT SỞ GD\&ĐT NINH BÌNH -- LẦN 3}{Môn Vật lí}{(Đề thi gồm: 18 câu trắc nghiệm, 04 câu đúng-sai, 06 câu trả lời ngắn)}{Thời gian làm bài: 50 phút (không kể thời gian phát đề).}

	\noindent\textbf{Cho biết:} $T(\text{K})=t\,(^\circ\text{C})+273$; $R=8{,}31$~J$\cdot$mol$^{-1}\cdot$K$^{-1}$; $N_A=6{,}02\cdot 10^{23}$~hạt/mol; $\pi=3{,}14$.

	\NeedspaceCau{10.3\baselineskip}
	\noindent
	\subsection*{\indam{PHẦN I. Câu trắc nghiệm nhiều phương án lựa chọn. Thí sinh trả lời từ câu 1 đến câu 18. Mỗi câu hỏi thí sinh chỉ chọn một phương án.}}
	\setcounter{ex}{0}
	\Opensolutionfile{ans}[ans/VatLi2026-NinhBinh-SoGD-Lan3-1]

	\begin{ex}%[VL12-C3-NB]
		Độ lớn suất điện động cảm ứng trong mạch kín tỉ lệ thuận với
		\choice
		{\True tốc độ biến thiên của từ thông qua mạch}
		{thời gian xảy ra sự biến thiên từ thông qua mạch}
		{độ lớn của từ thông qua mạch}
		{tổng trở của mạch kín}
		\loigiai{
			Định luật Faraday: $|e_\text{c}|=\left|\dfrac{\Delta\Phi}{\Delta t}\right|$, tỉ lệ với tốc độ biến thiên của từ thông qua mạch.
		}
	\end{ex}

	\NeedspaceCau{11.3\baselineskip}
	\begin{ex}%[VL12-C2-NB]
		Đồ thị nào sau đây biểu diễn quá trình biến đổi trạng thái của một lượng khí lí tưởng xác định khi áp suất không đổi?
		\begin{center}
			\begin{tikzpicture}[line join=round,line cap=round,font=\footnotesize]
				% Hình 1
				\begin{scope}
					\draw[VLmui,thick] (0,0)--(2.6,0) node[below,inner sep=3pt] {$T$ (K)};
					\draw[VLmui,thick] (0,0)--(0,2.4) node[right,inner sep=2pt] {$V$ (m$^3$)};
					\node[below left,inner sep=1.5pt] at (0,0) {O};
					\draw[thick,VLdo,dashed] (0,0)--(0.75,0.62);
					\draw[very thick,VLdo] (0.75,0.62)--(2.2,1.82);
					\node at (1.3,-0.55) {Hình 1};
				\end{scope}
				% Hình 2
				\begin{scope}[xshift=3.7cm]
					\draw[VLmui,thick] (0,0)--(2.6,0) node[below,inner sep=3pt] {$T$ (K)};
					\draw[VLmui,thick] (0,0)--(0,2.4) node[right,inner sep=2pt] {$p$ (Pa)};
					\node[below left,inner sep=1.5pt] at (0,0) {O};
					\draw[thick,VLdo,dashed] (0,0)--(0.75,0.62);
					\draw[very thick,VLdo] (0.75,0.62)--(2.2,1.82);
					\node at (1.3,-0.55) {Hình 2};
				\end{scope}
				% Hình 3
				\begin{scope}[xshift=7.4cm]
					\draw[VLmui,thick] (0,0)--(2.6,0) node[below,inner sep=3pt] {$V$ (m$^3$)};
					\draw[VLmui,thick] (0,0)--(0,2.4) node[right,inner sep=2pt] {$p$ (Pa)};
					\node[below left,inner sep=1.5pt] at (0,0) {O};
					\draw[very thick,VLdo] plot[domain=0.25:2.25,samples=60] (\x,{0.45/\x+0.03});
					\node at (1.3,-0.55) {Hình 3};
				\end{scope}
				% Hình 4
				\begin{scope}[xshift=11.1cm]
					\draw[VLmui,thick] (0,0)--(2.6,0) node[below,inner sep=3pt] {$V$ (m$^3$)};
					\draw[VLmui,thick] (0,0)--(0,2.4) node[right,inner sep=2pt] {$p$ (Pa)};
					\node[below left,inner sep=1.5pt] at (0,0) {O};
					\draw[thick,VLdo,dashed] (1.6,0)--(1.6,0.85);
					\draw[very thick,VLdo] (1.6,0.85)--(1.6,2.05);
					\node at (1.3,-0.55) {Hình 4};
				\end{scope}
			\end{tikzpicture}
		\end{center}
		\choice
		{Hình 3}
		{Hình 4}
		{Hình 2}
		{\True Hình 1}
		\loigiai{
			Đẳng áp: $\dfrac{V}{T}=\text{hằng số}$, đồ thị $V$ theo $T$ là đường thẳng có đường kéo dài qua gốc O (Hình 1).
		}
	\end{ex}

	\NeedspaceCau{5.1\baselineskip}
	\begin{ex}%[VL12-C1-NB]
		Một trong những thông số cần thiết để một chất được ứng dụng làm chất làm mát động cơ nhiệt là có
		\choice
		{nhiệt nóng chảy riêng lớn}
		{khối lượng riêng lớn}
		{\True nhiệt dung riêng lớn}
		{nhiệt độ nóng chảy lớn}
		\loigiai{
			Chất có nhiệt dung riêng lớn hấp thụ được nhiều nhiệt lượng mà nhiệt độ tăng ít nên làm mát tốt.
		}
	\end{ex}

	\NeedspaceCau{4.5\baselineskip}
	\begin{ex}%[VL12-C1-TH]
		Khi tiến hành thí nghiệm đo độ biến thiên nhiệt độ của một vật, một bạn học sinh thu được kết quả theo thang đo nhiệt độ Celsius và thang đo nhiệt độ Kelvin tương ứng là $\Delta t$ và $\Delta T$. Công thức đúng là
		\choice
		{$\Delta t=273\cdot\Delta T$}
		{$\Delta T=\Delta t+273$}
		{$\Delta T=273\cdot\Delta t$}
		{\True $\Delta t=\Delta T$}
		\loigiai{
			$T=t+273\Rightarrow \Delta T=\Delta t$ (mỗi độ chia $1\,^\circ\text{C}$ bằng $1$~K).
		}
	\end{ex}

	\NeedspaceCau{11.6\baselineskip}
	\begin{ex}%[VL12-C2-TH]
		Hình bên mô tả một lượng khí lí tưởng được nhốt trong xilanh, biết áp suất khí quyển là $1$~bar và áp kế chỉ mức ``0'' ứng với áp suất khí quyển. Khi áp kế chỉ $4$~bar thì áp suất của khí trong xilanh lúc này là
		% giuhinh
		\begin{center}
			\begin{tikzpicture}[scale=0.85,line join=round,line cap=round,font=\footnotesize]
				\pgfmathsetseed{5}
				% khí trong xilanh
				\fill[VLxanh!10] (0,0) rectangle (1.6,1.05);
				\foreach \i in {1,...,140}{\fill[gray!70] ({0.05+1.5*rnd},{0.05+0.95*rnd}) circle (0.35pt);}
				% thành xilanh
				\shade[left color=gray!35,right color=gray!10] (-0.1,-0.1) rectangle (0,1.85);
				\shade[left color=gray!10,right color=gray!35] (1.6,-0.1) rectangle (1.7,1.85);
				\fill[gray!40] (-0.1,-0.1) rectangle (1.7,0);
				\draw[thick] (-0.1,1.85)--(-0.1,-0.1)--(1.7,-0.1)--(1.7,1.85);
				\draw (0,1.85)--(0,0)--(1.6,0)--(1.6,1.85);
				% pit-tông, cán, núm
				\shade[top color=gray!55,bottom color=black!80] (0,1.05) rectangle (1.6,1.3);
				\draw (0,1.05) rectangle (1.6,1.3);
				\shade[left color=gray!70,right color=gray!40,middle color=gray!5] (0.74,1.3) rectangle (0.86,2.45);
				\draw[thin] (0.74,1.3) rectangle (0.86,2.45);
				\shade[top color=gray!20,bottom color=gray!65] (0.35,2.45) rectangle (1.25,2.56);
				\draw[thin] (0.35,2.45) rectangle (1.25,2.56);
				% bàn tay ấn
				\shade[top color=orange!12,bottom color=orange!35,draw=black!75] (2.1,2.58)--(0.8,2.58) .. controls (0.35,2.58) and (0.05,2.6) .. (0.02,2.7) .. controls (0.02,2.78) and (0.3,2.8) .. (0.65,2.8) .. controls (0.95,2.95) and (1.5,3.02) .. (2.1,2.98)--cycle;
				\draw[black!60,thin] (0.65,2.8) .. controls (0.8,2.75) and (0.95,2.72) .. (1.2,2.74);
				\shade[top color=gray!25,bottom color=gray!55,draw=black!75] (2.1,2.52) rectangle (2.4,3.04);
				% ống nối, áp kế nhỏ
				\fill[gray!55] (1.7,0.32) rectangle (2.22,0.4);
				\draw[thin] (1.7,0.32)--(2.22,0.32) (1.7,0.4)--(2.22,0.4);
				\filldraw[fill=white,draw=black,thick] (2.45,0.36) circle (0.23);
				\draw[VLmui,thin] (2.37,0.28)--(2.55,0.46);
				% khung phóng to
				\filldraw[fill=white,draw=black,thick,rounded corners=6pt] (2.75,0.65) rectangle (5.15,3.35);
				\fill[white] (2.76,0.95)--(2.62,0.52)--(2.76,1.4)--cycle;
				\draw[thick] (2.75,0.95)--(2.62,0.52)--(2.75,1.4);
				% mặt đồng hồ
				\coordinate (T) at (3.95,2.1);
				\fill[gray!55] ($(T)+(-0.12,-1.0)$) rectangle ($(T)+(0.12,-1.25)$);
				\draw ($(T)+(-0.12,-1.0)$) rectangle ($(T)+(0.12,-1.25)$);
				\filldraw[fill=gray!8,draw=black,thin] ($(T)+(-0.38,-1.25)$) rectangle ($(T)+(0.12,-1.4)$);
				\shade[inner color=gray!15,outer color=gray!60] (T) circle (1.05);
				\draw[thick] (T) circle (1.05);
				\filldraw[fill=white,draw=black] (T) circle (0.9);
				\foreach \g in {-40,-31,...,220}{\draw[very thin] ($(T)+(\g:0.8)$)--($(T)+(\g:0.88)$);}
				\foreach \g/\s in {180/0,145/1,112/2,80/3,48/4,14/5,-22/6}{\node[inner sep=0pt] at ($(T)+(\g:0.64)$) {\s};}
				\node[inner sep=0pt,gray] at ($(T)+(214:0.64)$) {1};
				\node[inner sep=0pt] at ($(T)+(0,-0.52)$) {Bar};
				\draw[VLmui,very thick] ($(T)+(228:0.12)$)--($(T)+(48:0.42)$);
				\fill (T) circle (0.06);
			\end{tikzpicture}
		\end{center}
		\choice
		{$4$~bar}
		{\True $5$~bar}
		{$1$~bar}
		{$3$~bar}
		\loigiai{
			Áp kế chỉ độ chênh so với áp suất khí quyển: $p=p_0+4=1+4=5$~bar.
		}
	\end{ex}

	\NeedspaceCau{4.1\baselineskip}
	\begin{ex}%[VL12-C3-NB]
		Sóng điện từ là quá trình lan truyền của
		\choice
		{từ trường trong không gian}
		{\True điện từ trường trong không gian}
		{điện trường trong môi trường vật chất}
		{từ trường trong môi trường vật chất}
		\loigiai{
			Sóng điện từ là điện từ trường lan truyền trong không gian (cả trong chân không).
		}
	\end{ex}

	\NeedspaceCau{4.3\baselineskip}
	\begin{ex}%[VL12-C2-NB]
		Theo mô hình động học phân tử của khí lí tưởng, hằng số Boltzmann là $k$, mối quan hệ giữa động năng tịnh tiến trung bình của phân tử $\overline{E}_{\text{đ}}$ và nhiệt độ tuyệt đối $T$ là
		\choice
		{\True $\overline{E}_{\text{đ}}=\dfrac{3}{2}kT$}
		{$\overline{E}_{\text{đ}}=\dfrac{1}{3}kT$}
		{$\overline{E}_{\text{đ}}=\dfrac{1}{2}kT$}
		{$\overline{E}_{\text{đ}}=\dfrac{2}{3}kT$}
		\loigiai{
			Động năng tịnh tiến trung bình của phân tử khí lí tưởng: $\overline{E}_{\text{đ}}=\dfrac{3}{2}kT$.
		}
	\end{ex}

	\NeedspaceCau{7.1\baselineskip}
	\begin{ex}%[VL12-C1-NB]
		Khi nói về mô hình động học phân tử về cấu tạo chất, phát biểu nào sau đây \textbf{không đúng}?
		\choice
		{Nhiệt độ càng cao thì tốc độ chuyển động của các phân tử càng lớn}
		{Các phân tử chuyển động không ngừng}
		{Giữa các phân tử có lực hút và lực đẩy}
		{\True Các chất được cấu tạo từ các hạt không riêng biệt gọi là phân tử}
		\loigiai{
			Các chất được cấu tạo từ các hạt riêng biệt là phân tử (nguyên tử).
		}
	\end{ex}

	\NeedspaceCau{2.7\baselineskip}
	\begin{ex}%[VL12-C4-NB]
		Phản ứng hạt nhân \textbf{không} tuân theo định luật bảo toàn
		\choice
		{động lượng}
		{số nucleon}
		{\True khối lượng}
		{điện tích}
		\loigiai{
			Phản ứng hạt nhân bảo toàn điện tích, số nucleon, động lượng, năng lượng toàn phần; khối lượng không bảo toàn.
		}
	\end{ex}

	\NeedspaceCau{4.7\baselineskip}
	\begin{ex}%[VL12-C3-TH]
		Một đoạn dây dẫn đặt nằm ngang có dòng điện chạy qua theo chiều từ Nam ra Bắc theo quy ước địa lí. Dây được đặt ở khu vực có từ trường đều các đường sức có chiều thẳng đứng từ trên xuống dưới. Lực từ do từ trường tác dụng lên dây có chiều hướng về phía
		\choice
		{Đông}
		{Nam}
		{\True Tây}
		{Bắc}
		\loigiai{
			Quy tắc bàn tay trái: đường sức hướng xuống đâm vào lòng bàn tay, chiều từ cổ tay đến ngón giữa hướng Bắc, ngón cái choãi ra chỉ hướng Tây.
		}
	\end{ex}

	\NeedspaceCau{4.9\baselineskip}
	\begin{ex}%[VL12-C2-NB]
		Định luật Boyle cho biết sự biến đổi trạng thái của một lượng khí xác định khi thông số
		\choice
		{áp suất của khí không đổi}
		{khối lượng riêng khí không đổi}
		{thể tích của khí không đổi}
		{\True nhiệt độ của khí không đổi}
		\loigiai{
			Định luật Boyle áp dụng cho quá trình đẳng nhiệt: $pV=\text{hằng số}$.
		}
	\end{ex}

	\NeedspaceCau{4.6\baselineskip}
	\begin{ex}%[VL12-C1-TH]
		Trong thời tiết mùa đông giá lạnh ở trong phòng học nếu chạm tay vào song sắt cửa sổ ta có cảm giác lạnh nhưng chạm tay vào bàn gỗ ta có cảm giác đỡ lạnh hơn. Gọi $T_1$, $T_2$ và $T_3$ lần lượt là nhiệt độ của bàn tay, song sắt cửa sổ và bàn gỗ, mối liên hệ giữa chúng là
		\choice
		{\True $T_1>T_2=T_3$}
		{$T_1>T_3>T_2$}
		{$T_1<T_2=T_3$}
		{$T_1>T_2>T_3$}
		\loigiai{
			Song sắt và bàn gỗ cùng ở nhiệt độ phòng ($T_2=T_3$), thấp hơn nhiệt độ bàn tay; sắt dẫn nhiệt tốt hơn gỗ nên ta thấy lạnh hơn.
		}
	\end{ex}

	\NeedspaceCau{4.6\baselineskip}
	\begin{ex}%[VL12-C2-NB]
		Với $p$ là áp suất, $V$ là thể tích, $T$ là nhiệt độ tuyệt đối. Phương trình nào sau đây là phương trình trạng thái của khí lí tưởng?
		\choice
		{$\dfrac{V}{T}=\text{hằng số}$}
		{\True $\dfrac{pV}{T}=\text{hằng số}$}
		{$\dfrac{VT}{p}=\text{hằng số}$}
		{$\dfrac{p}{T}=\text{hằng số}$}
		\loigiai{
			Phương trình trạng thái của khí lí tưởng: $\dfrac{pV}{T}=\text{hằng số}$.
		}
	\end{ex}

	\NeedspaceCau{2.6\baselineskip}
	\begin{ex}%[VL12-C4-NB]
		Trong hạt nhân oxygen ${}^{17}_{\phantom{1}8}\text{O}$ có
		\choice
		{$17$ hạt neutron}
		{\True $9$ hạt neutron}
		{$8$ hạt nucleon}
		{$9$ hạt proton}
		\loigiai{
			$Z=8$ proton, $A=17$ nucleon, $N=A-Z=17-8=9$ neutron.
		}
	\end{ex}

	\NeedspaceCau{7.4\baselineskip}
	\begin{ex}%[VL12-C3-TH]
		Hình vẽ bên cho biết hướng Bắc địa lí và chiều của dòng điện trong dây dẫn thẳng dài đi từ ngoài vào trong mặt phẳng hình vẽ. Vectơ cảm ứng từ do dòng điện này gây ra tại M có chiều hướng về phía
		\par\noindent
		\begin{minipage}[c]{0.64\linewidth}
		\choice
		{\True Đông}
		{Bắc}
		{Tây}
		{Nam}
		\end{minipage}\hfill
		\begin{minipage}[c]{0.34\linewidth}
			\centering
			\begin{tikzpicture}[line join=round,line cap=round,font=\footnotesize]
				\draw[gray,densely dotted] (-0.8,0)--(0.8,0) (0,-0.6)--(0,1.25);
				\fill (0,1.25) circle (1.8pt) node[left=3pt] {M};
				\filldraw[fill=white,draw=VLdo,semithick] (0,0) circle (0.13);
				\draw[VLdo,semithick] (0,0)++(-0.092,-0.092)--++(0.184,0.184) (0,0)++(-0.092,0.092)--++(0.184,-0.184);
				\node[VLdo,below left,inner sep=2pt] at (-0.08,-0.1) {$I$};
				\draw[VLmui,thick] (1.45,-0.55)--(1.45,1.2) node[above,inner sep=2pt] {Bắc};
			\end{tikzpicture}
		\end{minipage}
		\loigiai{
			Quy tắc nắm tay phải: đường sức là các đường tròn quanh dây, chiều kim đồng hồ (nhìn từ phía dòng điện đi vào). Tại M (phía Bắc của dây), $\vv{B}$ hướng về phía Đông.
		}
	\end{ex}

	\NeedspaceCau{14.2\baselineskip}
	\begin{ex}%[VL12-C4-TH]
		Hình bên là đồ thị đơn giản hóa phân bố năng lượng liên kết riêng $E_{\text{lkr}}$ theo số khối $A$ của hạt nhân. Từ đồ thị cho biết hạt nhân nào bền vững nhất trong các hạt nhân sau: ${}^{63}_{29}\text{Cu}$, ${}^{238}_{\phantom{2}92}\text{U}$, ${}^{206}_{\phantom{2}82}\text{Pb}$, ${}^4_2\text{He}$?
		\begin{center}
			\begin{tikzpicture}[x=0.032cm,y=0.33cm,line join=round,line cap=round,font=\footnotesize]
				\draw[gray!70,dashed,thin] (0,0) grid[xstep=20,ystep=2] (240,10);
				\draw[VLmui,thick] (0,0)--(262,0) node[right,inner sep=2pt] {$A$};
				\draw[VLmui,thick] (0,0)--(0,11.6) node[right,inner sep=2pt] {$E_{\text{lkr}}$ (MeV)};
				\foreach \x in {40,80,...,240}{\node[below,inner sep=2pt] at (\x,0) {$\x$};}
				\foreach \y in {2,4,...,10}{\node[left,inner sep=2pt] at (0,\y) {$\y$};}
				\node[below left,inner sep=2pt] at (0,0) {$0$};
				\draw[thick,VLdo] (2,1.1)--(3.7,2.9)--(3.9,7.15)--(7,5.6)--(12.4,7.7)--(15.3,7.45)--(16.6,8.0)--(19.5,7.7);
				\draw[thick,VLdo] plot[smooth] coordinates {(19.5,7.7) (31.1,8.4) (39.8,8.66) (56.3,8.77) (74.5,8.73) (90.3,8.62) (120.1,8.51) (153.2,8.37) (208.7,7.82) (235.6,7.46) (244,7.3)};
				\foreach \p in {(2,1.1),(3.7,2.9),(3.9,7.15),(7,5.6),(12.4,7.7),(15.3,7.45),(16.6,8.0),(19.5,7.7),(31.1,8.4),(39.8,8.66),(56.3,8.77),(74.5,8.73),(90.3,8.62),(120.1,8.51),(153.2,8.37),(208.7,7.82),(235.6,7.46)}{\fill \p circle (1.4pt);}
			\end{tikzpicture}
		\end{center}
		\choice
		{\True ${}^{63}_{29}\text{Cu}$}
		{${}^{238}_{\phantom{2}92}\text{U}$}
		{${}^4_2\text{He}$}
		{${}^{206}_{\phantom{2}82}\text{Pb}$}
		\loigiai{
			Hạt nhân có năng lượng liên kết riêng càng lớn càng bền. Theo đồ thị, $E_{\text{lkr}}$ lớn nhất ở vùng $A$ từ $50$ đến $80$ nên ${}^{63}_{29}\text{Cu}$ bền vững nhất.
		}
	\end{ex}

	\NeedspaceCau{3.8\baselineskip}
	\begin{ex}%[VL12-C1-NB]
		Tổng động năng và thế năng của các phân tử cấu tạo nên vật được gọi là
		\choice
		{nhiệt lượng của vật}
		{cơ năng của vật}
		{\True nội năng của vật}
		{động lượng của vật}
		\loigiai{
			Nội năng của vật là tổng động năng và thế năng của các phân tử cấu tạo nên vật.
		}
	\end{ex}

	\NeedspaceCau{4.9\baselineskip}
	\begin{ex}%[VL12-C3-NB]
		Thiết bị nào dưới đây \textbf{không} hoạt động dựa trên nguyên tắc của hiện tượng cảm ứng điện từ?
		\choice
		{Máy phát điện xoay chiều}
		{Máy biến áp}
		{\True Kim la bàn}
		{Đàn guitar điện}
		\loigiai{
			Kim la bàn hoạt động dựa trên tương tác giữa kim nam châm và từ trường Trái Đất, không dựa trên cảm ứng điện từ.
		}
	\end{ex}
	\Closesolutionfile{ans}

	\NeedspaceCau{16.4\baselineskip}
	\subsection*{\indam{PHẦN II. Câu trắc nghiệm đúng sai. Thí sinh trả lời từ câu 1 đến câu 4. Trong mỗi ý a), b), c), d) ở mỗi câu, thí sinh chọn đúng hoặc sai.}}
	\setcounter{ex}{0}
	\Opensolutionfile{ans}[ans/VatLi2026-NinhBinh-SoGD-Lan3-2]

	\begin{ex}%[VL12-C4-TH]
		Cho phản ứng hạt nhân ${}^{235}_{\phantom{2}92}\text{U}+{}^1_0\text{n}\to{}^{A_1}_{Z_1}\text{X}+{}^{A_2}_{Z_2}\text{Y}+k\,{}^1_0\text{n}$.
		\NeedspaceBang{11.2\baselineskip}%
		\choiceTFt
		{Phản ứng là phản ứng thu năng lượng}
		{Cơ chế của phản ứng trên được Otto Hahn đưa ra là dùng hạt neutron có động năng rất lớn (cỡ MeV) để bắn vào hạt nhân ${}^{A_1}_{Z_1}\text{X}$, ${}^{A_2}_{Z_2}\text{Y}$}
		{Các hạt nhân ${}^{A_1}_{Z_1}\text{X}$, ${}^{A_2}_{Z_2}\text{Y}$ sinh ra từ phản ứng phân hạch của ${}^{235}_{\phantom{2}92}\text{U}$ là hoàn toàn xác định}
		{\True Thực tế phản ứng trên không xảy ra trực tiếp mà phải qua trạng thái trung gian với phương trình là ${}^{235}_{\phantom{2}92}\text{U}+{}^1_0\text{n}\to{}^{236}_{\phantom{2}92}\text{U}^*\to{}^{A_1}_{Z_1}\text{X}+{}^{A_2}_{Z_2}\text{Y}+k\,{}^1_0\text{n}$}
		\loigiai{
			Đây là phản ứng phân hạch ${}^{235}_{\phantom{2}92}\text{U}$ do hấp thụ neutron.
			\begin{itemchoice}
				\itemch \textbf{Sai.} Phản ứng phân hạch tỏa năng lượng (cỡ $200$~MeV mỗi phân hạch).
				\itemch \textbf{Sai.} Neutron chậm (động năng nhỏ, cỡ $0{,}01$~eV) bắn vào hạt nhân ${}^{235}_{\phantom{2}92}\text{U}$, không phải vào các hạt nhân sản phẩm X, Y.
				\itemch \textbf{Sai.} Phân hạch ${}^{235}_{\phantom{2}92}\text{U}$ có thể tạo ra nhiều cặp hạt nhân sản phẩm khác nhau.
				\itemch \textbf{Đúng.} ${}^{235}_{\phantom{2}92}\text{U}$ hấp thụ neutron chậm thành hạt nhân kích thích ${}^{236}_{\phantom{2}92}\text{U}^*$, rồi hạt nhân này vỡ thành X, Y và $k$ neutron.
			\end{itemchoice}
		}
	\end{ex}

	\NeedspaceCau{25.9\baselineskip}
	\begin{ex}%[VL12-C3-VD]
		Sơ đồ hình bên mô tả các bộ phận chính của một máy đo rung chấn kiểu cảm ứng. Đầu ra nối tới mạch đo suất điện động cảm ứng. Đầu nhọn của máy đo rung chấn gắn chặt vào mặt đất. Ngay khi xảy ra động đất, vỏ và cuộn dây di chuyển lên xuống do chuyển động của mặt đất. Nam châm vẫn đứng yên do quán tính của nó. Trong khoảng thời gian $3{,}0$~ms, cuộn dây di chuyển từ vị trí có cảm ứng từ là $8{,}0$~mT đến vị trí có cảm ứng từ là $22{,}0$~mT (coi rằng các đường sức từ luôn vuông góc với mặt phẳng vòng dây). Cuộn dây của máy địa chấn có $250$ vòng và diện tích $18{,}0~\text{cm}^2$.
		\begin{center}
			\begin{tikzpicture}[scale=0.694,line join=round,line cap=round,font=\footnotesize]
				% vỏ
				\fill[gray!22] (0,0) rectangle (2.4,3.6);
				\fill[white] (0.12,0.12) rectangle (2.28,3.48);
				\draw[thick] (0,0) rectangle (2.4,3.6);
				\draw (0.12,0.12) rectangle (2.28,3.48);
				% đầu nhọn
				\shade[left color=gray!60,right color=gray!25] (0.95,0)--(1.2,-1.0)--(1.45,0)--cycle;
				\draw (0.95,0)--(1.2,-1.0)--(1.45,0);
				% lò xo
				\draw[gray!40!black,semithick,decorate,decoration={coil,aspect=0.45,segment length=2.6pt,amplitude=3.2pt}] (1.2,3.48)--(1.2,2.42);
				\draw[gray!40!black,semithick,decorate,decoration={coil,aspect=0.45,segment length=2.6pt,amplitude=3.2pt}] (1.2,1.42)--(1.2,0.12);
				% cuộn dây (nửa sau)
				\foreach \y in {1.82,1.9,1.98,2.06}{\draw[VLdong!70!black,thick] (0.38,\y) arc[start angle=180,end angle=0,x radius=0.82,y radius=0.16];}
				% nam châm
				\shade[left color=gray!70,right color=gray!45,middle color=gray!8] (1.07,1.42) rectangle (1.33,2.42);
				\draw[thin] (1.07,1.42) rectangle (1.33,2.42);
				\shade[left color=gray!70,right color=gray!45,middle color=gray!8] (0.98,2.42) rectangle (1.42,2.5);
				\draw[thin] (0.98,2.42) rectangle (1.42,2.5);
				% cuộn dây (nửa trước)
				\foreach \y in {1.82,1.9,1.98,2.06}{\draw[VLdong!85!black,thick] (0.38,\y) arc[start angle=180,end angle=360,x radius=0.82,y radius=0.16];}
				% dây ra
				\draw[thick] (0.75,3.6) .. controls (0.75,3.1) and (0.3,3.0) .. (0.45,2.6) .. controls (0.55,2.35) and (0.35,2.2) .. (0.38,2.06);
				\draw[thick] (1.7,3.6) .. controls (1.7,3.05) and (2.15,3.0) .. (2.0,2.6) .. controls (1.9,2.35) and (2.05,2.2) .. (2.02,2.06);
				\fill (0.75,3.6) circle (2.2pt) (1.7,3.6) circle (2.2pt);
				% nhãn
				\node[align=center,inner sep=1pt] (DR) at (1.25,4.45) {Đầu ra nối tới mạch đo\\suất điện động cảm ứng};
				\draw[thin] (DR.south)++(-0.2,0)--(0.78,3.68) (DR.south)++(0.2,0)--(1.67,3.68);
				\draw[thin] (1.38,2.3)--(2.75,2.85) node[right,inner sep=1pt] {Nam châm};
				\draw[thin] (1.95,1.82)--(2.75,1.15) node[right,inner sep=1pt] {Cuộn dây};
				\node[left,inner sep=1pt] (LX) at (-0.35,2.4) {Lò xo};
				\draw[thin] (LX.east)--(1.08,3.0) (LX.east)--(1.08,0.85);
				\draw[thin] (0.06,1.25)--(-0.45,1.25) node[left,inner sep=1pt] {Vỏ};
				\draw[thin] (1.08,-0.55)--(-0.15,-0.35) node[left,inner sep=1pt] {Đầu nhọn};
			\end{tikzpicture}
		\end{center}
		\NeedspaceBang{10\baselineskip}%
		\choiceTFt
		{Tín hiệu ở đầu ra của cuộn dây khi có rung chấn đo được là tín hiệu điện một chiều}
		{\True Nguyên tắc hoạt động của thiết bị dựa trên hiện tượng cảm ứng điện từ}
		{Khi tăng độ cứng của lò xo, thiết bị đo sẽ nhạy hơn với những rung chấn chậm}
		{\True Độ lớn suất điện động trung bình tạo ra trong khoảng thời gian $3{,}0$~ms như trên là $2{,}1$~V}
		\loigiai{
			\begin{itemchoice}
				\itemch \textbf{Sai.} Cuộn dây dao động lên xuống quanh nam châm, từ thông lúc tăng lúc giảm nên suất điện động đổi chiều: tín hiệu không phải một chiều.
				\itemch \textbf{Đúng.} Từ thông qua cuộn dây biến thiên làm xuất hiện suất điện động cảm ứng.
				\itemch \textbf{Sai.} Lò xo cứng hơn thì nam châm dễ chuyển động theo vỏ khi rung chấn chậm, chuyển động tương đối nhỏ nên thiết bị kém nhạy hơn.
				\itemch \textbf{Đúng.} $|e|=NS\dfrac{\Delta B}{\Delta t}=250\cdot 18\cdot 10^{-4}\cdot\dfrac{(22-8)\cdot 10^{-3}}{3\cdot 10^{-3}}=2{,}1$~V.
			\end{itemchoice}
		}
	\end{ex}

	\NeedspaceCau{21.4\baselineskip}
	\begin{ex}%[VL12-C2-TH]
		% KIỂM TRA: ý d) bản gốc Đúng / tính lại Sai (theo hình, (3) nằm trên đường thẳng qua O và (1) nên T3 = 4T2)
		Hình bên biểu diễn quá trình biến đổi trạng thái của một lượng khí lí tưởng xác định. Lúc đầu khối khí biến đổi đẳng nhiệt từ trạng thái (1) đến trạng thái (2) với $p_2=2p_1$, sau đó khối khí biến đổi đẳng áp đến trạng thái (3).
		\begin{center}
			\begin{tikzpicture}[x=0.78cm,y=1.05cm,line join=round,line cap=round,font=\footnotesize]
				\draw[VLmui,thick] (0,0)--(5.0,0) node[below,inner sep=2pt] {$V$};
				\draw[VLmui,thick] (0,0)--(0,3.0) node[left,inner sep=2pt] {$p$};
				\node[below left,inner sep=2pt] at (0,0) {O};
				\draw[gray,dashed,thin] plot[domain=0.72:4.4,samples=60] (\x,{2/\x});
				\draw[gray,dashed,thin] (0,0)--(4,2);
				\draw[gray,dashed,thin] (0,2)--(1,2)--(1,0) (0,1)--(2,1)--(2,0);
				\draw[very thick,VLdo,VLgiua=0.6] plot[domain=2:1,samples=40] (\x,{2/\x});
				\draw[very thick,VLxanh,VLgiua=0.55] (1,2)--(4,2);
				\foreach \p in {(1,2),(2,1),(4,2)}{\filldraw[fill=white,draw=black,thick] \p circle (2pt);}
				\node[above right,inner sep=1.5pt] at (1,2) {(2)};
				\node[right,inner sep=3pt] at (2,1) {(1)};
				\node[right,inner sep=3pt] at (4,2) {(3)};
				\node[left,inner sep=2pt] at (0,2) {$p_2$};
				\node[left,inner sep=2pt] at (0,1) {$p_1$};
				\node[below,inner sep=2pt] at (1,0) {$V_2$};
				\node[below,inner sep=2pt] at (2,0) {$V_1$};
			\end{tikzpicture}
		\end{center}
		\NeedspaceBang{10\baselineskip}%
		\choiceTFt
		{Khối khí thực hiện công khi biến đổi từ trạng thái (1) đến trạng thái (2)}
		{\True Quá trình biến đổi từ trạng thái (2) đến trạng thái (3) thì nhiệt độ của khối khí tăng}
		{\True Nội năng của khối khí chỉ phụ thuộc vào nhiệt độ}
		{Nhiệt độ tuyệt đối của khối khí ở trạng thái (3) gấp ba lần nhiệt độ tuyệt đối của khối khí ở trạng thái (2)}
		\loigiai{
			\begin{itemchoice}
				\itemch \textbf{Sai.} Từ (1) đến (2) đẳng nhiệt, áp suất tăng gấp đôi nên thể tích giảm một nửa: khí bị nén, nhận công.
				\itemch \textbf{Đúng.} Từ (2) đến (3) đẳng áp, thể tích tăng nên $T=\dfrac{pV}{nR}$ tăng.
				\itemch \textbf{Đúng.} Khí lí tưởng bỏ qua tương tác giữa các phân tử, nội năng chỉ là tổng động năng phân tử nên chỉ phụ thuộc nhiệt độ.
				\itemch \textbf{Sai.} $V_2=\dfrac{V_1}{2}$; (1) và (3) cùng nằm trên đường thẳng qua O nên $\dfrac{p_3}{V_3}=\dfrac{p_1}{V_1}\Rightarrow V_3=2V_1=4V_2$.\\
				Đẳng áp: $\dfrac{T_3}{T_2}=\dfrac{V_3}{V_2}=4$.
			\end{itemchoice}
		}
	\end{ex}

	\NeedspaceCau[16.4\baselineskip]{30\baselineskip}
	\begin{ex}%[VL12-C1-TH]
		Một nhóm học sinh tìm hiểu về sự truyền nhiệt. Các bạn có các dụng cụ và cách tiến hành như sau:\\
		\textit{Dụng cụ:}\\
		$\bullet$ Cốc nhôm đựng $200$~ml nước ở nhiệt độ $36\,^\circ\text{C}$ (1).\\
		$\bullet$ Bình cách nhiệt đựng $500$~ml nước ở nhiệt độ $41\,^\circ\text{C}$ (2).\\
		$\bullet$ Hai nhiệt kế thủy ngân (3).\\
		\textit{Tiến hành:} Đặt cốc nhôm vào trong lòng bình cách nhiệt như hình bên và quan sát số chỉ nhiệt kế để tìm hiểu về sự truyền nhiệt của chúng.
		% giuhinh
		\begin{center}
			\begin{tikzpicture}[scale=0.8,line join=round,line cap=round,font=\footnotesize]
				% (1) cốc nhôm
				\begin{scope}[shift={(0,0)}]
					\fill[VLxanh!25] (0.02,0.04) rectangle (0.58,0.55);
					\shade[left color=gray!45,right color=gray!20,middle color=white,opacity=0.55] (0,0) rectangle (0.6,0.8);
					\draw[gray!50!black] (0,0.8)--(0,0)--(0.6,0)--(0.6,0.8);
					\draw[gray!50!black] (-0.04,0.8)--(0.64,0.8);
					\VLnhietke{0.3}{0.15}{3.1}{2.5}
					\node[below,inner sep=2pt] at (0.3,0) {(1)};
				\end{scope}
				% (2) bình cách nhiệt
				\begin{scope}[shift={(1.1,0)}]
					\fill[VLxanh!25] (0.08,0.08) rectangle (1.52,0.85);
					\draw[gray!50!black,thick,rounded corners=1pt] (0.05,1.95)--(0.05,0.05)--(1.55,0.05)--(1.55,1.95);
					\draw[gray!50!black] (0,0) rectangle (1.6,2.0);
					\fill[gray!35] (-0.08,1.95) rectangle (1.68,2.08);
					\draw[gray!50!black] (-0.08,1.95) rectangle (1.68,2.08);
					\VLnhietke{0.8}{0.3}{3.35}{2.85}
					\node[below,inner sep=2pt] at (0.8,0) {(2)};
				\end{scope}
				% ghép: cốc trong bình
				\begin{scope}[shift={(3.25,0)}]
					\fill[VLxanh!25] (0.08,0.08) rectangle (1.52,1.12);
					\draw[gray!50!black,thick,rounded corners=1pt] (0.05,1.95)--(0.05,0.05)--(1.55,0.05)--(1.55,1.95);
					\draw[gray!50!black] (0,0) rectangle (1.6,2.0);
					\fill[gray!35] (-0.08,1.95) rectangle (1.68,2.08);
					\draw[gray!50!black] (-0.08,1.95) rectangle (1.68,2.08);
					\fill[VLxanh!25] (0.82,0.1) rectangle (1.38,1.0);
					\shade[left color=gray!45,right color=gray!20,middle color=white,opacity=0.55] (0.8,0.08) rectangle (1.4,1.35);
					\draw[gray!50!black] (0.8,1.35)--(0.8,0.08)--(1.4,0.08)--(1.4,1.35);
					\draw[gray!50!black] (0.76,1.35)--(1.44,1.35);
					\VLnhietke{0.45}{0.3}{3.35}{2.85}
					\VLnhietke{1.1}{0.25}{3.1}{2.5}
				\end{scope}
				% nhãn (3)
				\node (N3) at (0.95,3.75) {(3)};
				\draw[thin] (N3.west)--(0.37,3.0) (N3.south east)--(1.84,3.25);
				% nhiệt kế y tế
				\begin{scope}[shift={(5.6,1.0)}]
					\VLnhietkeyte{27}
					\draw[thin] (1.9,0.85)--(1.9,0.0) node[below,inner sep=1.5pt] {(3)};
				\end{scope}
			\end{tikzpicture}
		\end{center}
		\NeedspaceBang{13.9\baselineskip}%
		\choiceTFt
		{\True Dụng cụ (3) là nhiệt kế thủy ngân (nhiệt kế y tế) hoạt động dựa trên tính chất giãn nở vì nhiệt của thủy ngân}
		{Nội năng của nước ở bình (2) tăng do thế năng của cột nước tăng vì nó được dâng cao hơn so với ban đầu}
		{\True Thang chia nhiệt độ của nhiệt kế thủy ngân ghi nhiệt độ từ $35\,^\circ\text{C}$ đến $42\,^\circ\text{C}$ nên việc lựa chọn nhiệt độ trong thí nghiệm của nhóm là một trong những cách chọn hợp lí}
		{\True Sau một thời gian cả hai nhiệt kế chỉ giá trị không đổi và hai giá trị này bằng nhau chứng tỏ không có sự truyền nhiệt lượng giữa hai vật có nhiệt độ bằng nhau}
		\loigiai{
			\begin{itemchoice}
				\itemch \textbf{Đúng.} Thủy ngân nở ra khi nóng lên, cột thủy ngân dâng lên trong ống.
				\itemch \textbf{Sai.} Nước ở bình (2) nóng hơn nên truyền nhiệt cho cốc nhôm và nước trong cốc: nội năng của nó giảm. Thế năng của cả cột nước (cơ năng) không thuộc nội năng.
				\itemch \textbf{Đúng.} $36\,^\circ\text{C}$, $41\,^\circ\text{C}$ và nhiệt độ cân bằng (nằm giữa hai giá trị này) đều trong thang đo $35\,^\circ\text{C}$ -- $42\,^\circ\text{C}$.
				\itemch \textbf{Đúng.} Hai vật đạt cân bằng nhiệt; khi nhiệt độ bằng nhau thì không còn truyền nhiệt giữa chúng.
			\end{itemchoice}
		}
	\end{ex}
	\Closesolutionfile{ans}

	\NeedspaceCau{13.4\baselineskip}
	\subsection*{\indam{PHẦN III. Câu trắc nghiệm trả lời ngắn. Thí sinh trả lời từ câu 1 đến câu 6.}}
	\setcounter{ex}{0}
	\Opensolutionfile{ans}[ans/VatLi2026-NinhBinh-SoGD-Lan3-3]

	\begin{ex}%[VL12-C2-VD]
		Trong việc bảo quản thực phẩm người ta thường dùng máy hút chân không để hút không khí ra khỏi túi đựng thực phẩm. Một máy hút chân không đơn giản, cứ mỗi chu trình khép kín hoạt động bằng cách làm giãn khí trong túi vào xilanh của máy rồi đẩy khí từ xilanh ra ngoài (thời gian đẩy khí ra ngoài coi như rất nhỏ trong chu trình). Người ta dùng máy có dung tích xilanh là $1$ lít để hút chân không cho một túi đựng thực phẩm có dung tích $1$ lít, có áp suất ban đầu là $512$~mmHg. Biết cứ mỗi phút máy thực hiện được $6$ chu trình khép kín. Coi các chu trình máy thực hiện nhiệt độ của khí luôn không đổi và khí trong túi là khí lí tưởng. Thời gian để máy làm giảm áp suất trong túi xuống còn $1$~mmHg là bao nhiêu giây (làm tròn kết quả đến chữ số hàng đơn vị)?
		\par
		\shortans[oly]{90}
		\loigiai{
			Mỗi chu trình khí giãn đẳng nhiệt từ $1$ lít đến $1+1=2$ lít nên áp suất giảm một nửa: $p_n=\dfrac{p_0}{2^n}$.\\
			$\dfrac{512}{2^n}=1\Rightarrow n=9$ chu trình. Mỗi chu trình kéo dài $\dfrac{60}{6}=10$~s nên $t=9\cdot 10=90$~s.
		}
	\end{ex}

	\NeedspaceCau{16.7\baselineskip}
	\begin{ex}%[VL12-C2-VD]
		% KIỂM TRA: bản gốc 3,7 (đọc p2 = p0/3, giả thiết khí đơn nguyên tử U = 3pV/2) / tính lại 5,6 (theo lưới hình: p0 = 4 ô, p2 = 1 ô, V2 = 4V0 nên p1V1 = p2V2, T1 = T2, ΔU = 0)
		Một lượng khí lí tưởng xác định thực hiện quá trình biến đổi từ trạng thái (1) đến trạng thái (2) theo ba con đường biến đổi khác nhau (a, b, c) trên hệ tọa độ $(p-V)$ như hình bên. Trong đó $p$ là áp suất và $V$ là thể tích của khối khí. Gọi $Q_a$, $Q_b$, $Q_c$ lần lượt là nhiệt lượng mà khối khí nhận được trong quá trình đi từ trạng thái (1) đến trạng thái (2) theo các con đường (a), (b) và (c). Ta có $Q_a=xQ_b=yQ_c$. Giá trị của tổng $(x+y)$ bằng bao nhiêu (làm tròn kết quả đến chữ số hàng phần mười)?
		\begin{center}
			\begin{tikzpicture}[x=0.8cm,y=0.6cm,line join=round,line cap=round,font=\footnotesize]
				\draw[gray!55,densely dotted] (0,0) grid[step=1] (5,5);
				\draw[VLmui,thick] (0,0)--(6.0,0) node[below,inner sep=3pt] {$V$ (m$^3$)};
				\draw[VLmui,thick] (0,0)--(0,5.7) node[right,inner sep=2pt] {$p$ (Pa)};
				\node[below left,inner sep=2pt] at (0,0) {O};
				\node[left,inner sep=2pt] at (0,4) {$p_0$};
				\node[below,inner sep=2pt] at (1,0) {$V_0$};
				\draw[very thick,VLdo,VLgiua=0.5] (1,4)--(4,4);
				\draw[very thick,VLdo,VLgiua=0.5] (4,4)--(4,1);
				\draw[very thick,VLxanh,VLgiua=0.5] (1,4)--(1,1);
				\draw[very thick,VLxanh,VLgiua=0.5] (1,1)--(4,1);
				\draw[very thick,VLluc,VLgiua=0.55] (1,4)--(4,1);
				\foreach \p in {(1,4),(4,1)}{\filldraw[fill=white,draw=black,thick] \p circle (2pt);}
				\node[above,inner sep=3pt] at (1.1,4) {(1)};
				\node[below right,inner sep=2pt] at (4,1) {(2)};
				\node[above,inner sep=2pt] at (2.5,4) {a};
				\node[above right,inner sep=1pt] at (2.55,2.5) {b};
				\node[left,inner sep=3pt] at (1,2.5) {c};
			\end{tikzpicture}
		\end{center}
		\shortans[oly]{5{,}6}
		\loigiai{
			Theo hình: $p_1=p_0$, $V_1=V_0$; $p_2=\dfrac{p_0}{4}$, $V_2=4V_0\Rightarrow p_2V_2=p_1V_1\Rightarrow T_2=T_1$, $\Delta U=0$ nên $Q=A$ (công khí thực hiện, bằng diện tích dưới đường biến đổi).\\
			$Q_a=p_0\cdot 3V_0=3p_0V_0$; $Q_b=\dfrac{p_0+\frac{p_0}{4}}{2}\cdot 3V_0=1{,}875p_0V_0$; $Q_c=\dfrac{p_0}{4}\cdot 3V_0=0{,}75p_0V_0$.\\
			$x=\dfrac{Q_a}{Q_b}=1{,}6$; $y=\dfrac{Q_a}{Q_c}=4\Rightarrow x+y=5{,}6$.
		}
	\end{ex}

	\NeedspaceCau{7.6\baselineskip}
	\begin{ex}%[VL12-C1-VD]
		Một em bé bị sốt, thân nhiệt đang có nhiệt độ $38{,}9\,^\circ\text{C}$ được bác sĩ chỉ định dùng thuốc hạ sốt. Cơn sốt hạ xuống thân nhiệt còn $37{,}4\,^\circ\text{C}$ trong $20$ phút. Coi rằng cơ chế bay hơi của mồ hôi trên da là cách duy nhất để làm giảm nhiệt độ cơ thể, em bé có cân nặng là $15$~kg. Nhiệt dung riêng của cơ thể em bé lấy bằng nhiệt dung riêng của nước là $1000$~cal/(kg.$^\circ$C) và nhiệt hóa hơi riêng của nước ở điều kiện trên là $580$~cal/g. Tốc độ bay hơi trung bình trên da do thuốc gây ra là bao nhiêu gam/phút (làm tròn kết quả đến chữ số hàng phần trăm)?
		\par
		\shortans[oly]{1{,}94}
		\loigiai{
			Nhiệt lượng cơ thể tỏa ra: $Q=mc\Delta t=15\cdot 1000\cdot(38{,}9-37{,}4)=22\,500$~cal.\\
			Khối lượng mồ hôi bay hơi: $m'=\dfrac{Q}{L}=\dfrac{22\,500}{580}\approx 38{,}79$~g. Tốc độ bay hơi: $\dfrac{38{,}79}{20}\approx 1{,}94$~g/phút.
		}
	\end{ex}

	\NeedspaceCau{4.6\baselineskip}
	\begin{ex}%[VL12-C4-TH]
		Khối lượng của hạt nhân ${}^{30}_{15}\text{P}$, proton và neutron lần lượt là $29{,}97005$~amu, $1{,}00728$~amu và $1{,}00867$~amu. Biết rằng $1~\text{amu}=931{,}5~\text{MeV}/c^2$. Năng lượng liên kết riêng của hạt nhân ${}^{30}_{15}\text{P}$ bằng bao nhiêu MeV (làm tròn kết quả đến chữ số hàng phần trăm)?
		\par
		\shortans[oly]{8{,}36}
		\loigiai{
			$\Delta m=15\cdot 1{,}00728+15\cdot 1{,}00867-29{,}97005=0{,}2692$~amu.\\
			$W_{\text{lkr}}=\dfrac{\Delta m c^2}{A}=\dfrac{0{,}2692\cdot 931{,}5}{30}\approx 8{,}36$~MeV/nucleon.
		}
	\end{ex}

	\NeedspaceCau{5.6\baselineskip}
	\begin{ex}%[VL12-C3-TH]
		Một khung dây kim loại phẳng có diện tích $40~\text{cm}^2$ gồm $1000$ vòng dây, quay đều với tốc độ $3000$ vòng/phút quanh trục vuông góc với đường sức của một từ trường đều có cảm ứng từ $0{,}01$~T. Suất điện động hiệu dụng xuất hiện trong khung dây có giá trị bằng bao nhiêu vôn (lấy $\pi=3{,}14$) (làm tròn kết quả đến chữ số hàng phần mười)?
		\par
		\shortans[oly]{8{,}9}
		\loigiai{
			$\omega=2\pi\cdot\dfrac{3000}{60}=314$~rad/s; $E_0=NBS\omega=1000\cdot 0{,}01\cdot 40\cdot 10^{-4}\cdot 314=12{,}56$~V.\\
			$E=\dfrac{E_0}{\sqrt{2}}=\dfrac{12{,}56}{\sqrt{2}}\approx 8{,}9$~V.
		}
	\end{ex}

	\NeedspaceCau{13.8\baselineskip}
	\begin{ex}%[VL12-C3-VD]
		Một dây đồng điện trở $3~\Omega$ được uốn thành hình vuông cạnh $40$~cm, hai đầu dây đồng được nối với hai cực của một nguồn điện có suất điện động $6$~V, điện trở không đáng kể. Mạch điện đặt trong một từ trường có $\vv{B}$ cùng hướng với vectơ pháp tuyến $\vv{n}$ của mặt phẳng hình vuông như hình bên. Độ lớn cảm ứng từ tăng theo thời gian theo quy luật $B=2+15t$ (T), ($t$ đo bằng giây). Cho rằng từ trường do dòng điện gây ra không đáng kể. Cường độ dòng điện trong mạch là bao nhiêu ampe (làm tròn kết quả đến chữ số hàng phần mười)?
		% giuhinh
		\begin{center}
			\begin{tikzpicture}[line join=round,line cap=round,font=\footnotesize]
				\ctikzset{bipoles/length=0.9cm}
				\fill[VLxanh!8] (0,0) rectangle (3.4,2.6);
				\draw[VLdong!70!black,line width=1.1pt] (3.4,1.75)--(3.4,2.6)--(0,2.6)--(0,0)--(3.4,0)--(3.4,0.85);
				\draw[thick] (3.4,1.75) to[battery1] (3.4,0.85);
				\node[right,inner sep=2pt] at (3.6,1.3) {$E,\,r$};
				\draw[VLxanh!85!black,thick] (1.45,1.3) circle (0.22);
				\draw[VLxanh!85!black,thick] (1.45,1.3)++(-0.155,-0.155)--++(0.31,0.31) (1.45,1.3)++(-0.155,0.155)--++(0.31,-0.31);
				\node[VLxanh!85!black,right,inner sep=3pt] at (1.67,1.3) {$\vv{B}$};
			\end{tikzpicture}
		\end{center}
		\shortans[oly]{2{,}8}
		\loigiai{
			$e_\text{c}=S\left|\dfrac{\Delta B}{\Delta t}\right|=0{,}4^2\cdot 15=2{,}4$~V.\\
			$B$ hướng vào trong mặt phẳng hình và tăng nên dòng điện cảm ứng ngược chiều kim đồng hồ, cùng chiều dòng điện do nguồn tạo ra (đi ra từ cực dương ở trên): hai suất điện động mắc cùng chiều.\\
			$I=\dfrac{E+e_\text{c}}{R}=\dfrac{6+2{,}4}{3}=2{,}8$~A.
		}
	\end{ex}

	\centerline{--- HẾT ---}
	\Closesolutionfile{ans}

	%===================== BẢNG ĐÁP ÁN =====================
	%\newpage
	%\begin{center}\indam{BẢNG ĐÁP ÁN}\end{center}
	%\noindent\indam{PHẦN I.}\\
	%\inputansbox{6}{ans/VatLi2026-NinhBinh-SoGD-Lan3-1}
	%\noindent\indam{PHẦN II.}\\
	%\inputansbox{2}{ans/VatLi2026-NinhBinh-SoGD-Lan3-2}
	%\noindent\indam{PHẦN III.}\\
	%\inputansbox{6}{ans/VatLi2026-NinhBinh-SoGD-Lan3-3}
\end{document}
